% ======================================================================
%  F. BRUHAT and J. TITS
%  "Groupes reductifs sur un corps local, I: Donnees radicielles valuees"
%  Publications mathematiques de l'I.H.E.S. 41 (1972), 5-251
%
%  COMPLETE ENGLISH TRANSLATION -- single continuous edition
%
%  Assembled from twelve part-translations covering pp. 5-251 without gap.
%  The sectioning is that of the original; the original numbering of every
%  statement is preserved, nothing has been renumbered.  Every numbered
%  statement carries a PDF anchor, and every textual cross-reference of the
%  form (6.4.28) that resolves to a statement present in the document is a
%  hyperlink, as are the bibliography citations [n].
%
%  Compile:  pdflatex bruhat-tits.tex   (three times)
% ======================================================================
\documentclass[12pt,a4paper,oneside,english]{smfart}

\usepackage[T1]{fontenc}
\usepackage[utf8]{inputenc}
\usepackage{newtxtext}
\makeatletter
\let\Bbbk\relax
\makeatother
\usepackage{newtxmath}
\usepackage{mathrsfs}
\DeclareMathAlphabet{\mathsf}{OT1}{phv}{m}{n}
\usepackage{bm}
\usepackage{enumitem}
\usepackage{array,longtable,booktabs,multicol}
\usepackage{tikz}
\usetikzlibrary{arrows.meta,positioning,calc}
\usepackage[margin=2.5cm,marginparwidth=1.5cm,marginparsep=4mm]{geometry}
\usepackage{marginnote}
\usepackage{xcolor}
\usepackage[expansion=false]{microtype}
\usepackage{needspace}
\ifdefined\pdfmapfile\pdfmapfile{+newtx.map}\fi
\definecolor{btlink}{rgb}{0.10,0.25,0.55}
\usepackage[
  colorlinks=true,   linkcolor=btlink,   citecolor=btlink,   urlcolor=btlink,
  % the sources restart the original's local equation numbering and \tag does
  % not step the counter; letting hyperref name its own destinations keeps
  % them unique.  Our own targets are explicit \hypertarget names.
  hypertexnames=false,   bookmarks=false,   pdftitle={Reductive groups over a local field I -- English translation},   pdfauthor={F. Bruhat and J. Tits},   pdfsubject={Publ. Math. IHES 41 (1972), 5-251}, ]{hyperref}

\renewcommand*{\marginfont}{\scriptsize\color{gray!80!black}}
\setlength{\parindent}{1.4em}
\allowdisplaybreaks
\emergencystretch=2em
\tolerance=1200
\hbadness=10000
\vbadness=10000

\pagestyle{plain}

% ---------------------------------------------------------------- anchors
\makeatletter
\newcommand{\btanchor}[1]{\leavevmode\Hy@raisedlink{\hypertarget{bt:#1}{}}}
\newcommand{\btpageanchor}[1]{\Hy@raisedlink{\hypertarget{p:#1}{}}}
\makeatother
\newcommand{\btref}[1]{\hyperlink{bt:#1}{#1}}
\newcommand{\btbib}[1]{\hyperlink{bib:#1}{#1}}
\newcommand{\btunnum}[1]{\par\medskip\noindent\emph{#1}.\ ---\ \ignorespaces}

% original page numbers, in the margin
\newcommand{\btpage}[1]{\marginnote{p.~#1}\btpageanchor{#1}}
\newcommand{\pg}[1]{\btpage{#1}}
\newcommand{\pmark}[1]{\btpage{#1}}
\newcommand{\pagemark}[1]{\btpage{#1}}

% ------------------------------------------------------- sectioning
\newcommand{\btsecA}[2]{%
  \par\needspace{4\baselineskip}\phantomsection
  \section*{#2}\markboth{#2}{#2}%
  \ifx\relax#1\relax\else\btanchor{#1}\fi}
\newcommand{\btsecB}[2]{%
  \par\needspace{4\baselineskip}\phantomsection
  \subsection*{#2}%
  \ifx\relax#1\relax\else\btanchor{#1}\fi}
\newcommand{\btsecC}[2]{%
  \par\needspace{3\baselineskip}\phantomsection
  \subsubsection*{#2}%
  \ifx\relax#1\relax\else\btanchor{#1}\fi}

\setcounter{tocdepth}{2}
\setcounter{secnumdepth}{0}

% theorem machinery, hoisted out of the sources
\theoremstyle{definition}
\newtheorem*{defnstar}{Definition}
\newtheoremstyle{bt}{\medskipamount}{\medskipamount}%
  {\itshape}{0pt}{\itshape}{. ---}{ }%
  {\thmname{#1}\thmnote{ (#3)}}
\theoremstyle{bt}
\newtheorem{prop}{Proposition}
\newtheorem{lem}{Lemma}
\newtheorem{cor}{Corollary}
\newtheorem{defn}{Definition}
\newtheorem{rem}{Remark}
\newtheorem{rems}{Remarks}
\newtheorem{exam}{Example}
\newtheorem{exams}{Examples}

\title[Reductive Groups over a Local Field I]{Reductive Groups over a Local Field I. Valued Root Data}
\author{F. Bruhat and J. Tits}
\date{}

\begin{document}
\raggedbottom
\maketitle
\begin{center}\small English translation of \emph{Groupes r\'eductifs sur un corps local, I: Donn\'ees radicielles valu\'ees}, \emph{Publ.\ Math.\ I.H.\'E.S.} \textbf{41} (1972), 5--251.\end{center}
\tableofcontents

% ---------------------------------------- bruhattitsch1.tex
\begingroup\providecommand{\num}{}
\renewcommand{\num}[1]{\medskip\noindent\textbf{(#1)}\ \ignorespaces}
\providecommand{\Num}{}
\renewcommand{\Num}[1]{\medskip\noindent\textbf{#1}\ \ignorespaces}
\providecommand{\vv}{}
\renewcommand{\vv}[1]{{\vphantom{#1}}^{v}\mkern-2mu#1}
\providecommand{\vA}{}
\renewcommand{\vA}{\vv{\mathbf{A}}}
\providecommand{\vW}{}
\renewcommand{\vW}{\vv{\mathbf{W}}}
\providecommand{\vSig}{}
\renewcommand{\vSig}{\vv{\boldsymbol{\Sigma}}}
\providecommand{\Bg}{}
\renewcommand{\Bg}{\mathfrak{B}}
\providecommand{\Ps}{}
\renewcommand{\Ps}{\mathfrak{P}}
\providecommand{\Bldg}{}
\renewcommand{\Bldg}{\mathscr{I}}
\providecommand{\Fac}{}
\renewcommand{\Fac}{\mathscr{F}}
\providecommand{\Wall}{}
\renewcommand{\Wall}{\mathscr{H}}
\providecommand{\Ohat}{}
\renewcommand{\Ohat}{\mathcal{O}}
\providecommand{\Ghat}{}
\renewcommand{\Ghat}{\widehat{G}}
\providecommand{\Bhat}{}
\renewcommand{\Bhat}{\widehat{\mathrm{B}}}
\providecommand{\Phat}{}
\renewcommand{\Phat}{\widehat{\mathrm{P}}}
\providecommand{\Ntil}{}
\renewcommand{\Ntil}{\widetilde{\mathrm{N}}}
\providecommand{\Stab}{}\renewcommand{\Stab}{\operatorname{Stab}}
\providecommand{\Card}{}\renewcommand{\Card}{\operatorname{Card}}
\providecommand{\codim}{}\renewcommand{\codim}{\operatorname{codim}}
\providecommand{\Ker}{}\renewcommand{\Ker}{\operatorname{Ker}}
\providecommand{\Cox}{}\renewcommand{\Cox}{\operatorname{Cox}}
\providecommand{\Aut}{}\renewcommand{\Aut}{\operatorname{Aut}}
\setlength{\parindent}{1.5em}
% =====================================================================
\btsecA{}{Introduction}

\emph{The Introduction (pp.~7--12) is discursive expository prose rather than
mathematics; what follows is a summary of its content and of the notation it fixes, not a line-by-line rendering.}

\medskip

\noindent\textbf{Programme.} The memoir develops the theory of reductive
algebraic groups over a local field $K$ along the lines opened by Iwahori and Matsumoto; the main results had been announced in [\btbib{10}]--[16], with conjectures in [\btbib{9}]. The guiding picture: a semisimple group over $K$ --- or rather the group $G$ of its $K$-rational points --- is to be studied as an ``infinite-dimensional object'' over the residue field $k$. The analogy with the ordinary theory of algebraic groups over $k$ is systematic, and, as there, one separates the \emph{absolute} theory ($k$ algebraically closed) from the
\emph{relative} theory ($k$ arbitrary).

\medskip
\noindent\textbf{Iwahori and parahoric subgroups.} The role played in the
classical absolute theory by Borel subgroups (maximal connected solvable) and parabolic subgroups is played here, for $k$ algebraically closed, by
\emph{Iwahori subgroups} and \emph{parahoric subgroups}. Chapter~II shows they
carry a natural structure of connected prosolvable proalgebraic group, maximal for that property --- so an \emph{a priori} algebro-geometric account (not given here) would presumably define Iwahori subgroups as the maximal connected prosolvable subgroups of $G$. Here they are introduced otherwise: by the explicit Iwahori--Matsumoto construction [\btbib{28}] from the root structure when $G$ is split; and, in general (Steinberg: $G$ is quasi-split when $k$ is algebraically closed), either by extending that constructive method to the quasi-split case --- as Hijikata had done in cases [\btbib{26}], [\btbib{27}] --- or by Galois descent to the split case.

Just as a Borel subgroup together with the normalizer of a maximal torus forms a Tits system with finite Coxeter Weyl group, so for $G$ simply connected an Iwahori subgroup $B$ together with the rational points of the normalizer of a suitable maximal torus forms a Tits system, this time with infinite Weyl group; whence again that $B$ is self-normalizing, a Bruhat decomposition, and so on. \emph{Parahoric subgroups} are the subgroups containing an Iwahori $B$ and equal to a finite union of double cosets mod $B$. The non-simply-connected case behaves like the non-connected case classically; the Introduction restricts itself for simplicity to the simply connected case.

\medskip
\noindent\textbf{Descent and the relative case.} Passage to the relative case
goes by descent, using Galois descent techniques --- whence the frequent (and probably often superfluous) hypothesis that $k$ be perfect. As classically one replaces Borel subgroups, which need not exist over $k$, by minimal parabolic $k$-subgroups, so here one considers the minimal ``parahoric $k$-subgroups'', which are the ``$B$'' of Tits systems whose Weyl group is the affine Weyl group of a root system --- Tits systems \emph{of affine type}. That root system need not be homothetic to the relative root system of $G$, but has the same ordinary Weyl group.

\medskip
\noindent\textbf{Why parahoric subgroups matter.} (i)~For $G$ simply
connected, the maximal bounded subgroups are exactly the maximal parahoric subgroups; the maximal bounded subgroups of an arbitrary semisimple group over $K$ are determined completely, generalizing Hijikata's results for $\mathrm{PGL}_n$ and groups of nondegenerate sesquilinear forms. (ii)~Parahoric subgroups carry natural proalgebraic structures over $k$, so that questions about $G$ --- Galois cohomology in particular --- reduce ``modulo $p$'' to questions about finite-dimensional semisimple groups over $k$; this generalizes theorems of Kneser for $K$ locally compact of characteristic zero (anisotropic groups, vanishing of $H^1$, classification). (iii)~If the group splits over an unramified extension of $K$, the parahoric subgroups appear as groups of ``integral points'' of group schemes over the ring $\Ohat$ of integers of $K$ which give the group back under $\Ohat\rightarrow K$.

\medskip
\noindent\textbf{Affine roots and valuations (notation fixed here).} To $G$
one associates an \emph{affine root system}, absolute or relative, whose elements are the affine half-spaces
\[
  \alpha_{a,k}\;=\;\{x\in\mathbf{A}\mid a(x)+k\geqslant 0\},
  \qquad a\in\Phi,\ k\in\Gamma_a,
\]
of the dual $\mathbf{A}$ of the real vector space spanned by the relative root system $\Phi$ of $G$, where $\Gamma_a$ is a discrete subgroup of $\mathbf{R}$. Each affine root $\alpha$ has a subgroup $U_\alpha\leqslant G$; for fixed $a$ the $U_{\alpha_{a,k}}$, $k\in\Gamma_a$, form a filtration of the usual root subgroup $U_a$. These filtrations come from a \emph{valuation} of the root datum formed by the $U_a$ --- the guiding remark being that a valuation of a field is hardly more than a filtration of its additive group. Successive quotients are additive groups over $k$ in the absolute case, connected unipotent groups of class $\leqslant 2$ in the relative case.

\medskip
\noindent\textbf{Plan of Chapter I.} The chapter is the ``abstract'' part of
the theory: pure group theory, independent of any algebro-geometric structure, hence applicable to groups unrelated to algebraic groups. It is expounded from two successive points of view --- Tits systems of affine type (\S\S2--5), geometric at the outset and giving pride of place to parahoric subgroups; and valued root data (\S\S6--8), more analytic, resting on the $U_a$ and their valuations. The two are not strictly equivalent: there are Tits systems of affine type not arising from a valued root datum (automorphism groups of certain trees), while the valuations considered are not generally assumed discrete and only discrete ones yield such Tits systems. Consequently many statements appear twice in slightly different forms; the authors note the redundancy, decline to recast the memoir, and remark that the two accounts at least exhibit two aspects of the objects and call on quite different modes of reasoning.

Section by section:
\begin{itemize}[leftmargin=1.6em,itemsep=2pt]
\item \S1 --- recollections on Tits systems and affine Weyl groups.
\item \S2 --- to each Tits system of affine type $(B,N)$ is attached a
  \emph{building} $\Bldg$: a geometric polysimplicial complex, related to but
  not identical with the combinatorial object of [\btbib{39}] or [\btbib{5}] (chap.~IV, \S2,   exercises 10--17), carrying a distance for which it is complete. In rank $1$   (Weyl group infinite dihedral) $\Bldg$ is a \emph{tree}. It plays in part   the role of the Riemannian symmetric space of real semisimple Lie theory.
\item \S3 --- maximal bounded subgroups of $G$, for a bornology naturally
  attached to $(B,N)$ (in the algebraic case, the one defined by the valuation   of $K$), via a \emph{fixed point lemma} valid for these buildings and for   Riemannian symmetric spaces alike --- so usable also for maximal compact   subgroups of real Lie groups.
\item \S4 --- Iwasawa and Cartan decompositions, and relations between double
  cosets with consequences for representation theory over locally compact   local fields. From $(B,N)$ one builds $\Bg\leqslant G$, which in the   algebraic case is a minimal parabolic $K$-subgroup; the Iwasawa   decomposition is $G=\Bg.P$ with $P$ maximal bounded --- valid only for the   ``good'' maximal bounded subgroups and under an extra hypothesis on   $(B,N)$, satisfied e.g.\ whenever $G=\Bg.N.\Bg$. Unlike the real case one   cannot generally split off $R\leqslant\Bg$ with $G=R.P$, $R\cap P=\{1\}$.   Non-simply-connected groups are accommodated by carrying along a second   group $\Ghat$ with a \emph{B-N-adapted} homomorphism   $\varphi:G\rightarrow\Ghat$ ($\Ghat$, $G$ being in applications the rational   points of an arbitrary semisimple group and of the simply connected covering   of its identity component) --- whence the heaviness of the notation.
\item \S5 --- when $G=\Bg N\Bg$, $(\Bg,N)$ is a Tits system with finite Weyl
  group: a \emph{double Tits system}. This paragraph is an axiomatic study of   that situation; it is not used later and may be skipped.
\item \S6 --- valued root data. A root datum of type $\Phi$ in $G$ is a family
  $(T,(U_a)_{a\in\Phi})$ subject to (\btref{6.1.1}) and (DR1)--(DR6), modelled on   $G=$ rational points of a semisimple group over $K$, $T$ the centralizer of   a maximal split torus, $\Phi$ the relative roots, $U_a$ the unipotent root   subgroups ($U_{(a)}$ of [\btbib{3}]). A valuation is a family   $\varphi_a:U_a\rightarrow\mathbf{R}\cup\{\infty\}$ subject to (\btref{6.2.1}),   (V0)--(V5). A discrete valuation of a generating root datum yields a double   Tits system in a finite-index subgroup (\btref{6.5}), but the discrete case is not   assumed thereafter. Parahoric subgroups become special cases of the groups   $U_f$ of no.~6.4; over a local field with residue field $k$, every parahoric   group is an extension of a reductive $k$-group by a pronilpotent group.
\item \S\S7--8 --- \S\S2, 3, 4 transposed to valued root data: a building
  $\Bldg$ (that of the associated Tits system in the discrete case; in general   with chambers and facets but without the polysimplicial structure and   without completeness), the bornology of $G$, and the maximal bounded   subgroups. The order of proof is roughly reversed: subgroups generalizing   the parahorics, then Iwasawa and Bruhat decompositions, then the building.
\item \S9 --- technical; the descent theorem carrying valuations and affine
  Tits systems from split to quasi-split to arbitrary semisimple groups.
\item \S10 --- the classical groups by direct computation, independent of
  descent; all valuations of their natural root data are determined. Since   non-discrete valuations and skew fields of infinite rank over their center   are allowed, these results are not entirely recovered later.
\end{itemize}

\medskip
\noindent\textbf{Acknowledgements.} The work originated in discussions at the
1965 AMS Summer Institute ``Algebraic groups and discontinuous subgroups'' at Boulder, Colorado; the authors thank M.~Kneser, T.~A.~Springer and R.~Steinberg for conversations there, and Yale University, the Institute for Advanced Study, the IH\'ES and the Sonderforschungsbereich Theoretische Mathematik of the University of Bonn for the contacts that made the chapter possible so quickly.

\medskip
\noindent\emph{Terminological footnote (p.~8).} The authors adhere to
Bourbaki's terminology for ease of reference to [\btbib{5}]; hence what some call a (BN)-pair is here called a \emph{Tits system} --- one author having insisted, the other resigning himself.

% =====================================================================

\btsecA{1}{1. Recollections and notation}

\btsecB{1.1}{1.1. Polysimplicial complexes}

\num{1.1.1}\btanchor{1.1.1} Let $F$ be a set. By an \emph{affine structure} on $F$ we mean a
map $[0,1]\times F\times F\rightarrow F$, written $(t,x,y)\mapsto tx+(1-t)y$, such that there exists a bijection $j$ of $F$ onto a convex subset possessing an interior point of a finite-dimensional real affine space $E$, satisfying
\begin{equation}
  j\bigl(tx+(1-t)y\bigr)=t\,j(x)+(1-t)\,j(y)
\end{equation}
for all $t\in[0,1]$ and $x,y\in F$. Every convex subset of a finite-dimensional affine space carries a canonical affine structure. A set equipped with an affine structure is called an \emph{affine set}.

If $F$ is an affine set, the affine space $E$ and the bijection $j$ satisfying (1) are clearly determined up to affine isomorphism. One may therefore transport to $F$ the topology of $j(F)$, together with the notions of dimension, of open and closed segment, of parallelism of two segments, of convex subset and convex hull, of \emph{internal point}, of \emph{extreme point}, and so on. For each point $x$ of $F$, let $F_x$ be the set consisting of $x$ together with those points $y\in F$, $y\neq x$, for which there is an open segment of $F$ containing both $x$ and $y$. One checks at once that $F_x$ is a convex subset of $F$ and that the sets $F_x$, $x\in F$, form a partition of $F$. A subset of $F$ of the form $F_x$ is called a \emph{facet} (\emph{facette}) of $F$.

Let $F$ and $F'$ be affine sets. A \emph{morphism} of $F$ into $F'$ (also called an \emph{affine map}) is a map $j:F\rightarrow F'$ satisfying (1). This defines the category of affine sets; every bijective morphism is an isomorphism. The map $\bigl(t,(x,x'),(y,y')\bigr)\mapsto\bigl(tx+(1-t)y,\;tx'+(1-t)y'\bigr)$ endows $F\times F'$ with an affine structure, the \emph{product} of those of $F$ and $F'$.

\num{1.1.2}\btanchor{1.1.2} A \emph{closed} (resp.\ \emph{open}) \emph{simplex of dimension
$r$} is an affine set isomorphic to a closed (resp.\ open) simplex of an $r$-dimensional real vector space $E$, i.e.\ to the convex hull (resp.\ the interior of the convex hull) of $r+1$ affinely independent points of $E$. A closed simplex $F$ of dimension $r$ has $r+1$ extreme points; a facet of $F$ is an open simplex of dimension $\leqslant r$, namely the set of internal points of the convex hull of a nonempty subset of the set of extreme points of $F$.

\num{1.1.3}\btanchor{1.1.3} A \emph{closed polysimplex} is an affine set $F$ isomorphic to the
product of finitely many closed simplices $F_1,\dots,F_k$ ($k\geqslant0$) of dimension $>0$. One verifies easily that the simplices $F_j$ and the projections $F\rightarrow F_j$ are determined up to isomorphism. After identifying $F$ with $F_1\times\dots\times F_k$, the facets of $F$ are the subsets $G_1\times\dots\times G_k$ with $G_j$ a facet of $F_j$; they are therefore \emph{open polysimplices}, i.e.\ affine sets isomorphic to products of open simplices.

\num{1.1.4}\btanchor{1.1.4} Let $A$ be a set and $\Fac$ a set of subsets of $A$, whose
elements will be called the \emph{facets} of $A$. Assume:
\begin{enumerate}[label=(\arabic*),leftmargin=2.4em]
\item \emph{the elements of $\Fac$ form a partition of $A$.}
\end{enumerate}
Suppose given moreover an order relation on $\Fac$, written $F\prec F'$, and for $F\in\Fac$ let $\bar F$ denote the union of the elements of $\Fac$ bounded above by $F$ (such elements being called the \emph{facets of $F$}). Finally, for each $F\in\Fac$ suppose given an affine structure on $\bar F$ such that:
\begin{enumerate}[label=(\arabic*),leftmargin=2.4em,start=2]
\item \emph{$\bar F$ is a closed polysimplex;}
\item \emph{the set of $F'\in\Fac$ bounded above by $F$ is the set of facets
  of $\bar F$; the closure of such an $F'$ in $\bar F$ coincides with   $\bar F'$, and the affine structure of $\bar F'$ is induced by that of   $\bar F$.}
\end{enumerate}
It follows at once that each facet $F$ of $A$ is an open polysimplex, namely the set of internal points of $\bar F$.

\num{1.1.5}\btanchor{1.1.5} We assume henceforth that the dimensions of the facets of $A$ are
bounded above, and write $\dim A$ for their supremum. For $F\in\Fac$ put $\codim F=\dim A-\dim F$. The facets of codimension $0$ (resp.\ $1$) are called the \emph{chambers} (resp.\ the \emph{panels}, \emph{cloisons}) of $A$. Two chambers are \emph{adjacent} (\emph{mitoyennes}) if they are distinct and have a common panel. A \emph{gallery of length $n$} is a sequence $\Gamma=(C_0,C_1,\dots,C_n)$ of $n+1$ chambers such that $C_{i-1}$ and $C_i$ are adjacent or equal ($1\leqslant i\leqslant n$). One says $C_0$ is the
\emph{origin} and $C_n$ the \emph{extremity} of $\Gamma$ (or that $C_0$ and
$C_n$ are the extremities of $\Gamma$). One says $\Gamma$ is
\emph{non-stuttering} (\emph{sans b\'egaiement}) if $C_{i-1}\neq C_i$ for
$1\leqslant i\leqslant n$.

Let $F,F'$ be facets of $A$. A gallery $\Gamma=(C_0,\dots,C_n)$ \emph{joins} $F$ to $F'$ if $F\subset\bar C_0$ and $F'\subset\bar C_n$. It is \emph{taut} (\emph{tendue}) between $F$ and $F'$ if it joins them and there is no gallery joining $F$ to $F'$ of length strictly less than $n$. For $x,x'\in A$, a gallery is said to be taut between $x$ and $x'$ if it is taut between the facets $F,F'$ containing $x$ and $x'$ respectively. A gallery is
\emph{minimal} if it is taut between its extremities.

\begin{defnstar}[1.1.6]\btanchor{1.1.6}
One says that $A$ (equipped with the set $\Fac$, the order relation $F\prec F'$, and the affine structures on each $\bar F$, satisfying the preceding conditions) is a \emph{polysimplicial complex} if, whatever the facets $F$ and $F'$ of $A$, there exists a gallery joining $F$ and $F'$.
\end{defnstar}

\noindent If $A$ is a polysimplicial complex, the chambers of $A$ are none
other than the maximal facets. If the chambers of a polysimplicial complex are simplices, one says it is a \emph{simplicial complex}; all facets are then (open) simplices.

\num{1.1.7}\btanchor{1.1.7} Let $A,A'$ be polysimplicial complexes. A \emph{morphism} of $A$
into $A'$ is a map $f:A\rightarrow A'$ with the two properties:
\begin{enumerate}[label=(\roman*),leftmargin=2.4em]
\item the image under $f$ of a facet of $A$ is a facet of $A'$;
\item if $F$ is a facet of $A$, the restriction of $f$ to $\bar F$ is an
  affine map of $\bar F$ onto $\overline{f(F)}$.
\end{enumerate}
The morphism $f$ is \emph{chambered} (\emph{chambr\'e}) if for every chamber $C$ of $A$ the restriction of $f$ to $C$ is an isomorphism of $C$ onto a chamber of $A'$.

\num{1.1.8}\btanchor{1.1.8} Let $A_1,A_2$ be polysimplicial complexes. Define the facets of
$A_1\times A_2$ to be the products $F_1\times F_2$ with $F_i$ a facet of $A_i$ ($i=1,2$), and order them by
\[
  F_1'\times F_2'\prec F_1\times F_2
  \quad\Longleftrightarrow\quad
  F_1'\prec F_1\ \text{ and }\ F_2'\prec F_2 .
\]
One then has $\overline{F_1\times F_2}=\bar F_1\times\bar F_2$, and one may endow $\overline{F_1\times F_2}$ with the product affine structure. One verifies easily that these data make $A_1\times A_2$ a polysimplicial complex, the \emph{product} of those of $A_1$ and $A_2$; the canonical projections are morphisms. This extends in the obvious way to a product of any finite number of polysimplicial complexes.

\btsecB{1.2}{1.2. Coxeter systems and Tits systems}

In this number we recall some definitions and results concerning Coxeter systems and Tits systems. For the proofs we refer to ([\btbib{5}], chap.~IV), which we designate by \emph{l.c.} throughout 1.2.

\num{1.2.1}\btanchor{1.2.1} Recall that a \emph{Coxeter system} is a pair $(W,S)$, where $W$
is a group and $S$ a subset of $W$, satisfying the following axioms (\emph{l.c.}, \S1, no.~3):
\begin{enumerate}[label=(\roman*),leftmargin=2.4em]
\item $S$ generates $W$ and the elements of $S$ are of order $2$;
\item let $m(s,s')$ be the order of the element $ss'$ of $W$ (for
  $s,s'\in S$), and let $I$ be the set of pairs $(s,s')$ of elements of $S$   such that $m(s,s')$ is finite; the generating set $S$ together with the   relations $(ss')^{m(s,s')}=1$ for $(s,s')\in I$ form a presentation of $W$.
\end{enumerate}
When $\Card S=1$ one has $W=\mathbf{Z}/2\mathbf{Z}$. When $\Card S=2$, the group $W$ is a dihedral group of order $2m(s,s')$, with $S=\{s,s'\}$ (\emph{l.c.}, \S1, no.~2).

\num{1.2.2}\btanchor{1.2.2} Let $(W,S)$ be a Coxeter system. The \emph{length} of an element
$w\in W$, written $\ell(w)$, is the least integer $q\geqslant0$ such that $w$ is a product of a sequence of $q$ elements of $S$. A \emph{reduced decomposition} of $w$ is a sequence $(s_1,\dots,s_q)$ of elements of $S$ with $w=s_1\dots s_q$ and $\ell(w)=q$ (\emph{l.c.}, \S1, no.~1). The set $\{s_1,\dots,s_q\}$ depends only on $w$ and not on the chosen reduced decomposition (\emph{l.c.}, \S1, no.~8, prop.~7); it is written $S_w$ and called the \emph{support} of $w$. The elements
\[
  t_j=(s_1\dots s_{j-1})\,s_j\,(s_1\dots s_{j-1})^{-1},
  \qquad 1\leqslant j\leqslant q,
\]
are pairwise distinct and their set depends only on $w$ (\emph{l.c.}, \S1, no.~4, lemma~2).

\num{1.2.3}\btanchor{1.2.3} Let $(s_1,\dots,s_q)$ be a reduced decomposition of an element
$w\in W$ and let $s\in S$. If $\ell(sw)\leqslant\ell(w)$, there is an integer $j$ with $1\leqslant j\leqslant q$ such that $s\,s_1\dots s_{j-1}=s_1\dots s_j$, and $(s_1,\dots,s_{j-1},s_{j+1},\dots,s_q)$ is a reduced decomposition of $sw$, which is therefore of length $q-1$ (\emph{l.c.}, \S1, no.~4, lemma~3). Let $s_1,\dots,s_q$ be elements of $S$ and put $w=s_1\dots s_q$; there exists a sequence $1\leqslant j_1<j_2<\dots<j_k\leqslant q$ such that $(s_{j_1},\dots,s_{j_k})$ is a reduced decomposition of $w$, and one has $k\equiv q\pmod 2$.

\num{1.2.4}\btanchor{1.2.4} For every subset $X$ of $S$, write $W_X$ for the subgroup of $W$
generated by $X$. It is also the set of elements $w\in W$ whose support is contained in $X$ (\emph{l.c.}, \S1, no.~8, cor.~1 of prop.~7). In particular one has $W_X\cap S=X$ and $W_X\cap W_Y=W_{X\cap Y}$ for $X,Y\subset S$. The pair $(W_X,X)$ is a Coxeter system (\emph{l.c.}, \S1, no.~8, th.~2).

\num{1.2.5}\btanchor{1.2.5} The \emph{Coxeter graph} $\Cox(W,S)$ (or $\Cox W$) of the Coxeter
system $(W,S)$ is the pair $(G,f)$ obtained as follows: $G$ is the graph whose vertices are the elements of $S$, two distinct vertices $s$ and $s'$ being joined by an edge if and only if the order $m(s,s')$ of $ss'$ is $\geqslant3$ (which means that $s$ and $s'$ do not commute); $f$ is the map $\{s,s'\}\mapsto m(s,s')$ from the set of edges into the set consisting of $\infty$ and the integers $\geqslant3$ (\emph{l.c.}, \S1, no.~9).

\num{1.2.6}\btanchor{1.2.6} By a \emph{Tits system} we mean a quadruple $(G,B,N,S)$, where $G$
is a group, $B$ and $N$ two subgroups of $G$, and $S$ a subset of $N/(B\cap N)$, satisfying the following axioms (\emph{l.c.}, \S2, no.~1):
\begin{itemize}[leftmargin=3.2em]
\item[(T1)] \emph{The set $B\cup N$ generates $G$, and $B\cap N$ is a normal
  subgroup of $N$.}
\item[(T2)] \emph{The set $S$ generates the group $W=N/(B\cap N)$ and consists
  of elements of order $2$.}\footnote{This second hypothesis made on $S$ is in   fact a consequence of the remaining axioms ([\btbib{37}]).}
\item[(T3)] \emph{For every $s\in S$ and every $w\in W$ one has
  $sBw\subset BwB\cup BswB$.}
\item[(T4)] \emph{For every $s\in S$ one has $sBs\neq B$.}
\end{itemize}
Note that the elements of $W$ are cosets modulo $B\cap N$, which gives a meaning to the products $Bw$, $BwB$, etc.

The set $S$ is the set of elements $w\in W$ such that $B\cup BwB$ is a subgroup of $G$ (\emph{l.c.}, \S2, no.~5, cor.\ of th.~3). It follows that, if $G$, $B$ and $N$ are given, there exists at most one subset $S$ of $N/(B\cap N)$ such that $(G,B,N,S)$ is a Tits system. This will allow us to say, by abuse of language, that the triple $(G,B,N)$ is a Tits system, or again that the pair $(B,N)$ is a Tits system in $G$.

The group $W$ is called the \emph{Weyl group} of the Tits system. The pair $(W,S)$ is a Coxeter system (\emph{l.c.}, \S2, no.~4, th.~2).

In the remainder of this number, $(G,B,N,S)$ denotes a Tits system, $W$ its Weyl group, and $\nu:N\rightarrow W$ the canonical homomorphism.

\num{1.2.7}\btanchor{1.2.7} The map $w\mapsto BwB$ is a bijection of $W$ onto the set
$B\backslash G/B$ of double cosets of $G$ modulo $B$ (\emph{l.c.}, \S2, no.~3, th.~1).

\num{1.2.8}\btanchor{1.2.8} Let $s\in S$ and $w\in W$; the three relations ``$BsBwB$ contains
a single double coset modulo $B$'', $BswB=BsBwB$, and $\ell(sw)=\ell(w)+1$ are equivalent (\emph{l.c.}, \S2, no.~4, th.~2). Consequently, if $(s_1,\dots,s_q)$ is a reduced decomposition of an element $w\in W$, one has $BwB=Bs_1Bs_2B\dots Bs_qB$.

\num{1.2.9}\btanchor{1.2.9} For every subset $X$ of $S$, put $B_X=BW_XB$ (the union of the
double cosets $BwB$ for $w$ running through the subgroup $W_X$ of $W$ generated by $X$; this set is denoted $G_X$ in (\emph{l.c.}, \S2)). The map $X\mapsto B_X$ is a bijection of the set of subsets of $S$ onto the set of subgroups of $G$ containing $B$, and one has $B_X\cap B_Y=B_{X\cap Y}$ for $X,Y\subset S$ (\emph{l.c.}, \S2, no.~5, th.~3).

\num{1.2.10}\btanchor{1.2.10} One says that a subgroup $P$ of $G$ is a \emph{parabolic
subgroup} if it contains a conjugate of $B$. There then exists a unique subset $X$ of $S$ such that $P$ is conjugate to $B_X$ (\emph{l.c.}, \S2, no.~6, prop.~4). One says that $X$ is the \emph{type} of $P$. For every subset $Y$ of $S$ containing $X$ there is exactly one parabolic subgroup of type $Y$ containing $P$, and one obtains in this way all the subgroups containing $P$.

Every parabolic subgroup is its own normalizer; two distinct parabolic subgroups whose intersection is again a parabolic subgroup are not conjugate; two parabolic subgroups contained in one and the same subgroup $P$ and conjugate by an element of $G$ are conjugate by an element of $P$ (\emph{l.c.}, \S2, no.~6, th.~4).

\num{1.2.11}\btanchor{1.2.11} Let $(G',B',N',S')$ be another Tits system, with $G'=G$. One says
that the Tits systems $(G,B,N,S)$ and $(G,B',N',S')$ are \emph{equivalent} if the subgroups $B$ and $B'$ are conjugate, or what amounts to the same, if they possess the same parabolic subgroups. The Weyl groups $W$ and $W'$, or more precisely the Coxeter systems $(W,S)$ and $(W',S')$, are then canonically isomorphic. Indeed, let $g\in G$ be such that $B'=gBg^{-1}$; since $B$ is its own normalizer, such an element is determined up to right multiplication by an element of $B$. The map $x\mapsto gxg^{-1}$ defines a bijection, independent of the choice of $g$, of the set $B\backslash G/B$ onto $B'\backslash G/B'$, hence a bijection $\gamma$ of $W$ onto $W'$ (cf.~(\btref{1.2.7})). Since the elements $s$ of $S$ (resp.\ $S'$) are characterized by the fact that $B\cup BsB$ (resp.\ $B'\cup B'sB'$) is a subgroup (\btref{1.2.6}), one has $\gamma(S)=S'$. On the other hand, let $(s_1,\dots,s_q)$ be a reduced decomposition of an element $w\in W$. One deduces immediately from (\btref{1.2.8}) that $(\gamma(s_1),\dots,\gamma(s_q))$ is a reduced decomposition of $\gamma(w)$. It follows at once that, for $s,t\in S$, the order of $\gamma(s)\gamma(t)$ in $W'$ equals that of $st$ in $W$, and that consequently there exists exactly one homomorphism $\gamma'$ of $W$ into $W'$ with $\gamma'(s)=\gamma(s)$ for all $s\in S$. Since
\[
  \gamma(w)=\gamma(s_1)\dots\gamma(s_q)
          =\gamma'(s_1)\dots\gamma'(s_q)=\gamma'(w),
\]
one sees that $\gamma$ is indeed an isomorphism of $(W,S)$ onto $(W',S')$.

\num{1.2.12}\btanchor{1.2.12} Put $H=\bigcap_{n\in N}nBn^{-1}$. This is a subgroup of $G$,
normalized by $N$. Put $\Ntil=NH$. Then $H$ is a normal subgroup of $\Ntil$, one has $H=B\cap\Ntil$, and the injection of $N$ into $\Ntil$ defines an isomorphism of $W$ onto $\Ntil/(B\cap\Ntil)$, by means of which we shall identify these two groups. One says that the Tits system $(G,B,N,S)$ is
\emph{saturated} if $\Ntil=N$. In all cases the quadruple $(G,B,\Ntil,S)$ is a
saturated Tits system, equivalent to $(G,B,N,S)$, called the \emph{saturated Tits system associated with} $(G,B,N,S)$ (\emph{l.c.}, \S2, exerc.~5).

Two Tits systems will be called \emph{strongly equivalent} if the associated saturated Tits systems are the same, up to inner automorphism. Note that two Tits systems, even saturated, may be equivalent without being strongly equivalent (except however if their Weyl group is finite (cf.\ \emph{l.c.},
\S2, exerc.~15)); in particular the groups $N$ and $N'$ may fail to be
conjugate (cf.~(\btref{2.8.13})).

\num{1.2.13}\btanchor{1.2.13} A homomorphism $\varphi$ of the group $G$ into a group $\Ghat$ is
said to be \emph{B-adapted} (resp.\ \emph{B-N-adapted}) if the following two conditions are satisfied:
\begin{enumerate}[label=(\roman*),leftmargin=2.4em]
\item the kernel of $\varphi$ is contained in $B$;
\item for every $g\in\Ghat$ there exists an element $h\in G$ such that
  $\varphi(hBh^{-1})=g\varphi(B)g^{-1}$ (resp.\   $\varphi(hBh^{-1})=g\varphi(B)g^{-1}$ and   $\varphi(hNh^{-1})=g\varphi(N)g^{-1}$).
\end{enumerate}
In the remainder of this number, $\varphi:G\rightarrow\Ghat$ denotes a B-adapted homomorphism.

\num{1.2.14}\btanchor{1.2.14} Since $G$ is generated by the union of the conjugates of $B$
(\emph{l.c.}, \S2, no.~4, cor.~2 of th.~2), the image $\varphi(G)$ is a normal subgroup of $\Ghat$. On the other hand, (i) entails that the kernel of $\varphi$ is contained in every parabolic subgroup of $G$, hence in $H$.

\num{1.2.15}\btanchor{1.2.15} For every parabolic subgroup $P$ of $G$ and every $g\in\Ghat$,
the inverse image $\varphi^{-1}\bigl(g\varphi(P)g^{-1}\bigr)$ is a parabolic subgroup of $G$, which we shall denote by $^{g}P$. One thus defines an action of $\Ghat$ on the set of parabolic subgroups of $G$. For every parabolic subgroup $P$ of $G$ we shall write $\Stab_{\Ghat}P$, or simply $\Stab P$, for the set of $g\in\Ghat$ such that $^{g}P=P$. One has $\varphi^{-1}(\Stab P)=P$.

\num{1.2.16}\btanchor{1.2.16} For $g\in\Ghat$, the right coset modulo $B$ of the element
$h\in G$ satisfying condition (ii) of (\btref{1.2.13}) is well determined, since $B$ is its own normalizer and the kernel of $\varphi$ is contained in $B$. One deduces that there exists exactly one permutation $\xi(g)$ of $W$ such that
\[
  \varphi\bigl(B\cdot\xi(g)(w)\cdot B\bigr)
   =\varphi(h)^{-1}\cdot g\cdot\varphi(BwB)\cdot g^{-1}\cdot\varphi(h)
\]
for every $w\in W$ and every $h\in G$ satisfying (\btref{1.2.13})(ii). An argument analogous to that of (\btref{1.2.11}) shows that $\xi(g)$ is an automorphism of the Coxeter system $(W,S)$ and that \emph{$\xi$ is a homomorphism of $\Ghat$ into the group of automorphisms of $(W,S)$, or again into the group of automorphisms of the Coxeter graph of $(W,S)$}. We put $\Xi=\xi(\Ghat)$.

\num{1.2.17}\btanchor{1.2.17} The kernel $\Ghat_0$ of $\xi$ contains $\varphi(G)$. It follows
that the restriction of $\xi$ to $\Stab B$ is a surjective map of $\Stab B$ onto $\Xi$. We put $\Bhat=\Ghat_0\cap\Stab B$. The pair $(\Bhat,\varphi(N))$ is a Tits system in $\Ghat_0$, with Weyl group canonically isomorphic to $W$ (\emph{l.c.}, \S2, exerc.~8).

More generally, if $P$ is a parabolic subgroup of $G$, we put $\Phat=\Ghat_0\cap\Stab P$. If $P\supset B$, one has $\Phat=\varphi(P)\cdot\Bhat$. The map $P\mapsto\Phat$ is a bijection of the set of parabolic subgroups of $G$ onto that of $\Ghat_0$.

\num{1.2.18}\btanchor{1.2.18} If $P$ is a parabolic subgroup of type $X\subset S$, the
parabolic subgroup $^{g}P$ is of type $\xi(g)(X)$ (for $g\in\Ghat$).

For every subset $X$ of $S$ we write $\Xi_X$ for the set of $\xi\in\Xi$ such that $\xi(X)=X$. If $g\in\Stab B_X$, the minimal parabolic subgroups $B$ and $^{g}B$ are contained in $B_X$, hence are conjugate by an element of $B_X$ (\btref{1.2.10}). Consequently
\[
  \Stab B_X=\varphi(B_X)\cdot(\Stab B_X\cap\Stab B)
          =\varphi(B_X)\cdot\bigl(\Stab B\cap\xi^{-1}(\Xi_X)\bigr),
\]
since $^{g}B_X=B_{\xi(g)(X)}$ for $g\in\Stab B$.

\num{1.2.19}\btanchor{1.2.19} Let $X\subset S$ and let $\mathsf{H}$ be a subgroup of $\Xi_X$.
Put
\[
  Q(X,\mathsf{H})=\varphi(B_X)\cdot\bigl(\Stab B\cap\xi^{-1}(\mathsf{H})\bigr).
\]
The map $(X,\mathsf{H})\mapsto Q(X,\mathsf{H})$ is a bijection of the set of pairs $(X,\mathsf{H})\in\Ps(S)\times\Ps(\Xi)$ such that $\mathsf{H}$ is a subgroup of $\Xi_X$, onto the set of subgroups of $\Ghat$ containing $\Bhat$. One has $Q(X,\Xi_X)=\Stab B_X$; moreover $Q(X,\mathsf{H})\subset Q(X',\mathsf{H}')$ if and only if $X\subset X'$ and $\mathsf{H}\subset\mathsf{H}'$ (\emph{l.c.}, \S2, exerc.~8). The subgroups $Q(X,\mathsf{H})$ and $Q(X',\mathsf{H}')$ are conjugate in $\Ghat$ if and only if there exists $t\in\Xi$ such that $X'=t.X$ and $\mathsf{H}'=t.\mathsf{H}.t^{-1}$.

Let $P$ (resp.\ $P'$) be a parabolic subgroup of $G$, of type $X$ (resp.\ $X'$). In order that $\Stab P\supset\Stab P'$ it is necessary and sufficient that $P\supset P'$ (whence $X\supset X'$) and that $\Xi_X\supset\Xi_{X'}$: this follows from what precedes, by reducing through conjugation to the case $P'=B_{X'}$.

\num{1.2.20}\btanchor{1.2.20} Suppose $\varphi$ is B-N-adapted and let $E$ be the intersection
of the normalizers of $\varphi(B)$ and of $\varphi(N)$ in $\Ghat$. One has $\Ghat=E\cdot\varphi(G)$, whence $\Ghat=E\cdot\Ghat_0$. Let $e\in E$ and $n\in N$, and let $^{e}n\in N$ be such that $e\varphi(n)e^{-1}=\varphi(^{e}n)$. One has $^{e}(nBn^{-1})={}^{e}n\,.\,B\,.\,(^{e}n)^{-1}$ and $\nu(^{e}n)=\xi(e)\bigl(\nu(n)\bigr)$.

\btsecB{1.3}{1.3. Affine Weyl groups}

In this number we recall some definitions and results concerning ``Weyl groups of affine type''. For further details and for the proofs we refer to [\btbib{5}], designated by \emph{l.c.} throughout 1.3.

\num{1.3.1}\btanchor{1.3.1} Let $\mathbf{A}$ be a finite-dimensional real affine space. We
assume that the space of translations $\vA$ of $\mathbf{A}$ is endowed with a (nondegenerate) scalar product, hence with a norm, and we equip $\mathbf{A}$ with the corresponding (Euclidean) distance: one then says that $\mathbf{A}$ is a \emph{Euclidean affine space}.

Two subsets of $\mathbf{A}$ deduced from one another by a translation are said to be \emph{equipollent}.

We denote moreover by $\mathbf{W}$ a subgroup of the group of displacements of $\mathbf{A}$, generated by orthogonal reflections in (affine) hyperplanes of $\mathbf{A}$ and acting properly on $\mathbf{A}$ when the latter is given the discrete topology; this last condition amounts to saying that $\mathbf{W}$ is a discrete subgroup of the group of displacements of $\mathbf{A}$ (\emph{l.c.}, chap.~V, \S4, no.~9). For $w\in\mathbf{W}$ we write $\vv{w}$ for the corresponding automorphism of $\vA$.

\num{1.3.2}\btanchor{1.3.2} It is known (\emph{l.c.}, chap.~V, \S3, no.~8) that $\mathbf{A}$
decomposes canonically into a product of Euclidean affine spaces $\mathbf{A}_0\times\dots\times\mathbf{A}_m$, each of them a quotient of $\mathbf{A}$, in such a way that the following conditions are realized: the action of $\mathbf{W}$ passes to the quotient and defines a discrete group $\mathbf{W}_i$ of displacements of $\mathbf{A}_i$, generated by orthogonal reflections in hyperplanes of $\mathbf{A}_i$; the group $\mathbf{W}$ is identified with the direct product of the $\mathbf{W}_i$; one has $\mathbf{W}_0=\{1\}$, and each $\mathbf{W}_i$ for $1\leqslant i\leqslant m$ is
\emph{irreducible}, in the sense that $\mathbf{W}_i\neq\{1\}$ and that the
canonical linear representation of $\mathbf{W}_i$ in the space $\vv{\mathbf{A}_i}$ of translations of $\mathbf{A}_i$ is irreducible.

In everything that follows we shall assume that $\mathbf{A}_0$ is reduced to a point, which allows us not to consider it, and that \emph{each $\mathbf{W}_i$ for $1\leqslant i\leqslant m$ is infinite}, or, what amounts to the same (\emph{l.c.}, chap.~V, \S3, no.~9), \emph{has no fixed point in $\mathbf{A}_i$}. Under these conditions we shall say that $\mathbf{W}$ is an
\emph{affine Weyl group}; when moreover $m=1$, we shall say that $\mathbf{W}$
is \emph{irreducible}.

\num{1.3.3}\btanchor{1.3.3} If $L$ is an affine hyperplane of $\mathbf{A}$, we write $s_L$ for
the orthogonal reflection admitting $L$ as its set of fixed points. Conversely, if $s$ is an orthogonal reflection in a hyperplane, we write $L_s$ for that hyperplane, the set of fixed points of $s$.

By a \emph{wall} of $\mathbf{A}$ (relative to $\mathbf{W}$) we mean any hyperplane $L$ of $\mathbf{A}$ such that the reflection $s_L$ belongs to $\mathbf{W}$. It is known that the set $\Wall$ of walls of $\mathbf{A}$ is locally finite (\emph{l.c.}, chap.~V, \S3, no.~1).

By an \emph{affine root} of $\mathbf{A}$ we mean any closed half-space of $\mathbf{A}$ bounded by a wall, this wall being called the \emph{wall of the root}. The set of affine roots of $\mathbf{A}$ will be written $\boldsymbol{\Sigma}$. If $\alpha\in\boldsymbol{\Sigma}$, one puts $r_\alpha=s_{\partial\alpha}$, where $\partial\alpha$ is the wall of $\alpha$; one puts $\alpha^{*}=\overline{\mathbf{A}-\alpha}$, and one denotes by $\alpha_{+}$ the intersection of the affine roots containing a neighborhood of $\alpha$: it is an affine root equipollent to $\alpha$.

It is clear that the datum of $\boldsymbol{\Sigma}$ determines $\mathbf{W}$. A set of closed half-spaces of a Euclidean affine space $\mathbf{A}$ which is the set of affine roots of $\mathbf{A}$ for an affine Weyl group acting in $\mathbf{A}$ will be called an \emph{affine root system}.

By a \emph{facet} of $\mathbf{A}$ (relative to $\mathbf{W}$ or to $\boldsymbol{\Sigma}$) we mean an equivalence class in $\mathbf{A}$ for the equivalence relation ``$x$ and $y$ are contained in the same affine roots''.

A facet $F$ is an open convex subset of the affine subspace (called the
\emph{support} of $F$) that it spans. Let us endow the set $\Fac$ of facets
with the order relation ``$F'$ is contained in the closure of $F$''; the set denoted $\bar F$ in no.~1.1 coincides with the closure of $F$, and $\mathbf{A}$, equipped with $\Fac$, with this order relation and with the canonical affine structures on each of the sets $\bar F$, is a polysimplicial complex, whose structure is the product of the polysimplicial complex structures on each of the $\mathbf{A}_i$ (\emph{l.c.}, chap.~V, \S3, no.~8, prop.~6 and no.~9, prop.~8), which are moreover simplicial complexes.

The \emph{chambers} of $\mathbf{A}$ (relative to $\mathbf{W}$) are the connected components of the complement in $\mathbf{A}$ of the union of the walls, and the \emph{panels} of $\mathbf{A}$ are the facets whose support is a hyperplane (which is then a wall). It is known that $\mathbf{W}$ is simply transitive on the set of chambers (\emph{l.c.}, chap.~V, \S3, no.~2, th.~1).

\num{1.3.4}\btanchor{1.3.4} We denote by $\mathbf{C}$ a chamber of $\mathbf{A}$ fixed once and
for all. It is known (\emph{l.c.}, chap.~V, \S3, no.~3, th.~2) that
\emph{the closure $\bar{\mathbf{C}}$ of $\mathbf{C}$ is a fundamental domain
for $\mathbf{W}$ acting in $\mathbf{A}$}. Moreover, $\mathbf{W}$ is generated by the set $\mathbf{S}$ of reflections in the walls supporting the panels of $\mathbf{C}$ (called the \emph{walls of the chamber $\mathbf{C}$}), and
\emph{the pair $(\mathbf{W},\mathbf{S})$ is a Coxeter system} (\emph{l.c.},
chap.~V, \S3, no.~2, th.~1).

More precisely, for $1\leqslant i\leqslant m$, the projection $\mathbf{C}_i$ of $\mathbf{C}$ in $\mathbf{A}_i$ is a chamber of $\mathbf{A}_i$ (relative to $\mathbf{W}_i$). If $\mathbf{S}_i$ is the set of reflections in the walls of $\mathbf{C}_i$, the Coxeter system $(\mathbf{W}_i,\mathbf{S}_i)$ is irreducible and $(\mathbf{W},\mathbf{S})$ is identified with the product of the $(\mathbf{W}_i,\mathbf{S}_i)$; in other words, after identifying the $\mathbf{W}_i$ with subgroups of $\mathbf{W}$, one has $\mathbf{S}=\bigcup_{1\leqslant i\leqslant m}\mathbf{S}_i$.

The \emph{rank} of $(\mathbf{W},\mathbf{S})$, that is to say the cardinality of $\mathbf{S}$, is equal to $\sum_{1\leqslant i\leqslant m}(\dim\mathbf{A}_i+1)$. On the other hand, one calls \emph{rank} of $\mathbf{W}$ (or of $\boldsymbol{\Sigma}$) the dimension of $\mathbf{A}$.

The Coxeter system $(\mathbf{W},\mathbf{S})$ (and also its Coxeter graph) determines the affine Weyl group $\mathbf{W}$ essentially. More precisely, let $\mathbf{A}'$ be a Euclidean affine space, $\mathbf{W}'$ an affine Weyl group in $\mathbf{A}'$, and $\mathbf{S}'$ the set of reflections in the walls of a chamber of $\mathbf{A}'$ relative to $\mathbf{W}'$; then every isomorphism of the Coxeter system $(\mathbf{W},\mathbf{S})$ onto $(\mathbf{W}',\mathbf{S}')$ is induced by a unique isomorphism $\alpha:\mathbf{A}\rightarrow\mathbf{A}'$ of affine spaces; if moreover $\mathbf{W}$ is irreducible, $\alpha$ is a similarity, i.e.\ transforms the distance of $\mathbf{A}$ into a multiple of that of $\mathbf{A}'$.

\num{1.3.5}\btanchor{1.3.5} For every facet $F$ there exists exactly one facet of $\mathbf{C}$
which is the transform of $F$ by an element of $\mathbf{W}$. On the other hand, for every subset $X$ of $\mathbf{S}$ such that $X\cap\mathbf{S}_i\neq\mathbf{S}_i$ for $1\leqslant i\leqslant m$, the set $\mathbf{C}_X$ of points $a\in\bar{\mathbf{C}}$ such that $X$ is exactly the set of $s\in\mathbf{S}$ with $a\in L_s$, is a facet of $\mathbf{C}$. Moreover, the map $X\mapsto\mathbf{C}_X$ is a bijection of the set
\[
  \mathbf{T}=\{X\subset\mathbf{S}\mid
     X\cap\mathbf{S}_i\neq\mathbf{S}_i\ \text{for }1\leqslant i\leqslant m\}
   =\{X\subset\mathbf{S}\mid \mathbf{W}_X\ \text{is finite}\}
\]
onto the set of facets of $\mathbf{C}$. One has $\mathbf{C}=\mathbf{C}_{\varnothing}$. We shall say that $\mathbf{T}$ is the set of \emph{types} of facets, and that a facet $F$ is \emph{of type} $X\in\mathbf{T}$ if $F$ is the transform of $\mathbf{C}_X$ by an element of $\mathbf{W}$.

The stabilizer of $\mathbf{C}_X$, which is also the stabilizer of an arbitrary point of $\mathbf{C}_X$, is the subgroup $\mathbf{W}_X$ of $\mathbf{W}$ generated by $X$ (\emph{l.c.}, chap.~V, \S3, no.~3, prop.~1).

\num{1.3.6}\btanchor{1.3.6} In order that a facet $F$ be a panel it is necessary and
sufficient that its type reduce to a single element $s$; one also says that $F$ is \emph{of type $s$}. If two chambers $\mathbf{C}$ and $\mathbf{C}'$ are adjacent (\btref{1.1.5}), they have exactly one panel in common. Conversely, every panel of $\mathbf{A}$ is a facet of exactly two chambers.

If $\Gamma=(\mathbf{C}_0,\dots,\mathbf{C}_n)$ is a non-stuttering gallery, one calls \emph{type} of $\Gamma$ the sequence $\mathbf{s}(\Gamma)=(s_1,\dots,s_n)$ of elements of $\mathbf{S}$ such that the panel common to $\mathbf{C}_{i-1}$ and $\mathbf{C}_i$ is of type $s_i$ (for $1\leqslant i\leqslant n$). Two galleries of $\mathbf{A}$ with the same origin and the same type are identical.

\num{1.3.7}\btanchor{1.3.7} One says that a point $x$ of $\mathbf{A}$ is a \emph{special
point} if, for every wall $L$ of $\mathbf{A}$, there exists a wall equipollent to $L$ containing $x$. It is known that every special point is an extreme point of the closure of a chamber, and that for every chamber $\mathbf{C}$ there is at least one special point adherent to $\mathbf{C}$ (\emph{l.c.}, chap.~V, \S3, no.~10, cor.\ of prop.~11). If $x$ is a special point, the group $\mathbf{W}$ is the semidirect product of the stabilizer $\mathbf{W}_x$ of $x$ by the normal subgroup $\mathbf{V}$ of translations of $\mathbf{A}$ belonging to $\mathbf{W}$, and the canonical map of $\mathbf{W}_x$ into the group of automorphisms of the vector space $\vA$ is an isomorphism of $\mathbf{W}_x$ onto the image $\vW$ of the whole of $\mathbf{W}$ in $\Aut\vA$ (\emph{l.c.}, chap.~V, \S3, no.~10, prop.~9). Moreover the group $\mathbf{V}$ is a
\emph{lattice} in $\vA$, that is to say a discrete subgroup of maximal rank of
$\vA$ (\emph{l.c.}, chap.~VI, \S2, no.~5, remark). In order that $x\in\mathbf{A}$ be a special point it is necessary and sufficient that the projection of $x$ in $\mathbf{A}_i$ be, for every $i\in\{1,\dots,m\}$, a special point of $\mathbf{A}_i$ for $\mathbf{W}_i$.

Suppose $\mathbf{W}$ irreducible. A vertex $s$ of the Coxeter graph of $(\mathbf{W},\mathbf{S})$ is said to be \emph{special} if the corresponding extreme point of $\bar{\mathbf{C}}$ (that is to say the point of intersection of the hyperplanes $L_t$ for $t\in\mathbf{S}$, $t\neq s$) is a special point.

\num{1.3.8}\btanchor{1.3.8} There exists in the dual $(\vA)^{*}$ of $\vA$ one and only one
reduced root system $\vSig$ possessing the following property: make $\mathbf{A}$ into a vector space and identify it with $\vA$ by choosing as origin a special point of $\mathbf{A}$; for $a\in\vSig$ and $k\in\mathbf{Z}$, denote by $L_{a,k}$ the affine hyperplane consisting of the points $x\in\mathbf{A}$ such that $a(x)=-k$; then the set $\Wall$ of walls of $\mathbf{A}$ is none other than the set of hyperplanes $L_{a,k}$ for $a\in\vSig$ and $k\in\mathbf{Z}$ (\emph{l.c.}, chap.~VI, \S2, no.~5, prop.~8). The affine roots of $\mathbf{A}$ are therefore the closed half-spaces
\[
  \alpha_{a,k}=\{x\in\mathbf{A}\mid a(x)+k\geqslant0\}
  \qquad\text{for } a\in\vSig,\ k\in\mathbf{Z}.
\]
If $\alpha\in\boldsymbol{\Sigma}$ is of the form $\alpha_{a,k}$, one puts $\vv{\alpha}=a$ and $r_{a,k}=r_\alpha$.

In other words, once a special point has been chosen as origin, the group $\mathbf{W}$ is identified with the affine Weyl group of the root system $\vSig$ in the sense of (\emph{l.c.}, chap.~VI, \S2). In particular, for $a\in\vSig$, let us write $a^{\vee}$ for the unique element of $\vA$ such that the reflection $r_a$ associated with $a$ is given by $r_a(v)=v-a(v)a^{\vee}$ ($v\in\vA$). The set of the $a^{\vee}$ for $a\in\vSig$ forms the \emph{inverse root system} of $\vSig$ (\emph{l.c.}, chap.~VI, \S1, no.~1, prop.~1). The special points are the \emph{weights} of the root system $(\vSig)^{\vee}$ (\emph{l.c.}, chap.~VI, \S2, no.~2, prop.~3), the group $\mathbf{V}$ of translations of $\mathbf{W}$ is the group of \emph{radical weights} of $(\vSig)^{\vee}$ and is generated by $(\vSig)^{\vee}$ (\emph{l.c.}, chap.~VI, \S2, no.~1, prop.~1), and the chambers of $\mathbf{A}$ are the \emph{alcoves} of $\vSig$ in the sense of (\emph{l.c.}, chap.~VI, \S2, no.~1). The group $\vW$ (acting in $\vA$ and in its dual) is identified with the (ordinary) Weyl group of $\vSig$. It is generated by the reflections $r_a$ for $a\in\vSig$.

The elements of $\vSig$ will be called the \emph{vector roots} of $\mathbf{A}$ (relative to $\mathbf{W}$).

\num{1.3.9}\btanchor{1.3.9} To the decomposition
$\mathbf{A}=\mathbf{A}_1\times\dots\times\mathbf{A}_m$ of (\btref{1.3.2}) there corresponds a decomposition of the dual of $\vA$ into a direct sum of the $(\vv{\mathbf{A}_i})^{*}$, and the intersections $\vSig\cap(\vv{\mathbf{A}_i})^{*}$ are none other than the irreducible components of $\vSig$ (\emph{l.c.}, chap.~VI, \S1, no.~2).

\num{1.3.10}\btanchor{1.3.10} By a \emph{vector chamber} we mean a connected component of the
complement in $\vA$ of the union of the (vector) hyperplanes that are kernels of the vector roots, that is to say a chamber of the inverse root system $(\vSig)^{\vee}$ in the sense of (\emph{l.c.}, chap.~VI, \S1, no.~5). More generally, by a \emph{vector facet} we mean an equivalence class for the following equivalence relation in $\vA$: ``for every vector root $a$, $a(x)$ and $a(y)$ are simultaneously either zero, or strictly positive, or strictly negative''. A vector facet is a simplicial cone, a vector chamber is an open simplicial cone. The group $\vW$ is simply transitive on the set of vector chambers, and the closure of a vector chamber is a fundamental domain for $\vW$ acting on $\vA$ (\emph{l.c.}, chap.~VI, \S1, no.~5). The opposite of a vector chamber is a vector chamber (\emph{l.c.}, chap.~VI, \S1, cor.~3 of prop.~17).

\num{1.3.11}\btanchor{1.3.11} By a \emph{sector} (\emph{quartier}) of $\mathbf{A}$ we mean the
set $x+D$ of transforms of a given point $x\in\mathbf{A}$ (called the
\emph{apex} of the sector) by the translations belonging to a given vector
chamber $D$ (called the \emph{direction} of the sector). Two sectors of opposite direction are said to be \emph{opposite}. Let $x$ be a special point of $\mathbf{A}$; every sector with apex $x$ contains exactly one chamber $\mathbf{C}$ to which $x$ is adherent; the sector $Q$ is then the union of the transforms of $\mathbf{C}$ by the homotheties with center $x$ and ratio $>0$, and one obtains in this way a bijection of the set of sectors with apex $x$ (or again of the set of vector chambers) onto the set of chambers of $\mathbf{A}$ whose closure contains $x$.

\num{1.3.12}\btanchor{1.3.12} With every vector chamber $D$ is associated a \emph{basis} $B(D)$
of $\vSig$ (or system of simple roots), that is to say a linearly independent subset of $\vSig$ such that every element of $\vSig$ is a linear combination, with integral coefficients all of the same sign, of elements of $B(D)$: the vector chamber $D$ is the set of $x\in\vA$ such that $a(x)>0$ for all $a\in B(D)$. The map $D\mapsto B(D)$ is a bijection of the set of vector chambers onto the set of bases of $\vSig$ (\emph{l.c.}, chap.~VI, \S1, nos.~5 and 7). For every $a\in B(D)$, the hyperplane $\Ker a$ contains a face of the simplicial cone $D$, and one obtains in this way a bijection of $B(D)$ onto the set of faces of $D$.

\num{1.3.13}\btanchor{1.3.13} The vector roots which are linear combinations with positive
coefficients of elements of $B(D)$ are said to be \emph{positive} for the vector chamber $D$; they are moreover the positive roots for the order relation on the dual of $\vA$ defined by the convex cone $D$. Their set is written $\vSig^{+}_{D}$ or $\vv{\Sigma}^{+}(D)$. One has $\vSig=\vSig^{+}_{D}\cup(-\vSig^{+}_{D})$. If $\mathbf{W}$ is irreducible, or what amounts to the same, if the root system $\vSig$ is irreducible, there exists a greatest root $\tilde a$ for this order relation; if one identifies $\mathbf{A}$ and $\vA$ as above by taking a special point as origin, the chamber of $\mathbf{A}$ contained in the sector $D$ and containing the origin in its closure is the set of points $x\in\mathbf{A}$ such that $a(x)>0$ for all $a\in B(D)$ and $\tilde a(x)<1$ (\emph{l.c.}, chap.~VI, \S2, no.~2, prop.~5).

\num{1.3.14}\btanchor{1.3.14} Let $D$ be a vector chamber. It is known that the set
$R=R(D)$ of reflections $r_a$ for $a\in B(D)$ generates $\vW$, and that $(\vW,R)$ is a Coxeter system (\emph{l.c.}, chap.~VI, \S1, no.~5, th.~2); one has
\[
  \vv{\Sigma}^{+}\bigl(r_a(D)\bigr)
    =\{b\in\vv{\Sigma}^{+}(D)\mid b\neq a\}\cup\{-a\}
  \quad\text{and}\quad -a\in B\bigl(r_a(D)\bigr)
  \qquad(\textit{ibid.}).
\]

\num{1.3.15}\btanchor{1.3.15} Let $M$ be a subset of $\vSig$. We shall say that an element
$m\in M$ is \emph{extremal in $M$} if the half-line $\mathbf{R}_{+}m$ is an extremal generator of the convex cone with origin $0$ generated by $M$. An order on $M$ will be called \emph{nibbling} (\emph{grignotant}) if it is total and if every element $m\in M$ is extremal in the set of elements of $M$ greater than $m$. In order that there exist a nibbling order on a subset $M$ of $\vSig$, it is necessary and sufficient that the convex cone with origin $0$ generated by $M$ be strictly convex, or again that its polar have nonempty interior, or again that it meet a vector chamber $D$, which means that $M\subset\vv{\Sigma}^{+}(D)$. In this case there exists a nibbling order on $M$ for which the first element is a given element extremal in $M$. In particular, there exists a nibbling order on $\vv{\Sigma}^{+}(D)$ for which the first element is an arbitrary simple root $a\in B(D)$.

If $M'\subset M$, the restriction to $M'$ of a nibbling order on $M$ is a nibbling order on $M'$.

\num{1.3.16}\btanchor{1.3.16} For every vector chamber $D$ one defines an order relation, said
to be \emph{associated with $D$} and written $\leqslant_{D}$, on $\mathbf{A}$ (resp.\ $\vA$) in the following manner. Let $D^{*}$ be the chamber of the root system $\vSig$ associated with the basis $B(D)$ (\emph{l.c.}, chap.~VI, \S1, no.~5): it is the set of $t\in(\vA)^{*}$ whose scalar product with every element of $B(D)$ is strictly positive. By definition, the relation $x\leqslant_{D}x'$, for $x,x'\in\mathbf{A}$ (resp.\ $\vA$), is equivalent to $t(x'-x)\geqslant0$ for all $t\in D^{*}$ (or, what amounts to the same, for every dominant weight $t$ of $\vSig$ relative to $B(D)$). These order relations are compatible with the vector space structure of $\vA$ and with the action of $\vA$ on $\mathbf{A}$. For $y\in\bar D$ and $w\in\vW$ one has $w(y)\leqslant_{D}y$ (\emph{l.c.}, chap.~VI, \S1, no.~6, prop.~18).
\par\endgroup

% ---------------------------------------- bt.tex
\begingroup\usetikzlibrary{arrows.meta,positioning,calc}
\newenvironment{BT}[2]%
  {\medskip\par\noindent\textbf{#1 (#2).}\;---\;\itshape}%
  {\par\medskip}
\newenvironment{BTu}[1]%
  {\medskip\par\noindent\textbf{#1.}\;---\;\itshape}%
  {\par\medskip}
\providecommand{\para}{}
\renewcommand{\para}[1]{\medskip\par\noindent\textbf{(#1)}\ \ignorespaces}
\providecommand{\proofref}{}
\renewcommand{\proofref}[1]{\par\smallskip\noindent
  {\small\itshape Proof:} {\small see the original text, #1.}\par\medskip}
\providecommand{\bA}{}
\renewcommand{\bA}{\mathbf{A}}
\providecommand{\bC}{}
\renewcommand{\bC}{\mathbf{C}}
\providecommand{\bW}{}
\renewcommand{\bW}{\mathbf{W}}
\providecommand{\bS}{}
\renewcommand{\bS}{\mathbf{S}}
\providecommand{\bT}{}
\renewcommand{\bT}{\mathbf{T}}
\providecommand{\bV}{}
\renewcommand{\bV}{\mathbf{V}}
\providecommand{\bSig}{}
\renewcommand{\bSig}{\boldsymbol{\Sigma}}
\providecommand{\lsv}{}
\renewcommand{\lsv}[1]{{\vphantom{#1}}^{v}\!#1}
\providecommand{\vA}{}
\renewcommand{\vA}{\lsv{\bA}}
\providecommand{\vW}{}
\renewcommand{\vW}{\lsv{\bW}}
\providecommand{\vSig}{}
\renewcommand{\vSig}{\lsv{\bSig}}
\providecommand{\vw}{}
\renewcommand{\vw}{\lsv{w}}
\providecommand{\valpha}{}
\renewcommand{\valpha}{\lsv{\alpha}}
\providecommand{\Wt}{}
\renewcommand{\Wt}{\widetilde{\bW}}
\providecommand{\Vt}{}
\renewcommand{\Vt}{\widetilde{\bV}}
\providecommand{\Wd}{}
\renewcommand{\Wd}[1]{\Wt^{\dagger}_{#1}}
\providecommand{\Ech}{}
\renewcommand{\Ech}{\mathscr{E}}
\providecommand{\Imm}{}
\renewcommand{\Imm}{\mathscr{I}}
\providecommand{\Par}{}
\renewcommand{\Par}{\mathscr{P}}
\providecommand{\Fac}{}
\renewcommand{\Fac}{\mathscr{F}}
\providecommand{\Sep}{}
\renewcommand{\Sep}{\mathscr{S}}
\providecommand{\Sect}{}
\renewcommand{\Sect}{\mathfrak{C}}
\providecommand{\Dyn}{}\renewcommand{\Dyn}{\operatorname{Dyn}}
\providecommand{\Cox}{}\renewcommand{\Cox}{\operatorname{Cox}}
\providecommand{\cl}{}\renewcommand{\cl}{\operatorname{cl}}
\providecommand{\Card}{}\renewcommand{\Card}{\operatorname{Card}}
\providecommand{\pr}{}\renewcommand{\pr}{\operatorname{pr}}
\providecommand{\Id}{}\renewcommand{\Id}{\operatorname{Id}}
\providecommand{\RR}{}
\renewcommand{\RR}{\mathbf{R}}
\providecommand{\NN}{}
\renewcommand{\NN}{\mathbf{N}}
\providecommand{\Rpos}{}
\renewcommand{\Rpos}{\RR^{*}_{+}}
\tikzset{
  dynvertex/.style   = {circle, draw, line width=.5pt, inner sep=0pt,                         minimum size=4.2pt, fill=white},   dynedge/.style     = {line width=.5pt},   dyndouble/.style   = {line width=.5pt, double, double distance=1.7pt},   dynthick/.style    = {line width=1.9pt},
}
\newenvironment{dyngraph}[1][1]%
  {\begin{tikzpicture}[baseline=-2pt, scale=#1, x=7.2mm, y=7.2mm]}%
  {\end{tikzpicture}}
\providecommand{\dv}{}
\renewcommand{\dv}[3]{\node[dynvertex] (#3) at (#1,#2) {};}
\providecommand{\dvo}{}
\renewcommand{\dvo}[2]{\node[dynvertex] (#2) at (#1,0) {};}
\providecommand{\esingle}{}
\renewcommand{\esingle}[2]{\draw[dynedge]   (#1) -- (#2);}
\providecommand{\edouble}{}
\renewcommand{\edouble}[2]{\draw[dyndouble] (#1) -- (#2);}
\providecommand{\etriple}{}
\renewcommand{\etriple}[2]{\draw[dyndouble] (#1) -- (#2);
                         \draw[dynedge]   (#1) -- (#2);}
\providecommand{\ethick}{}
\renewcommand{\ethick}[2]{\draw[dynthick]   (#1) -- (#2);}
\providecommand{\gtr}{}
\renewcommand{\gtr}[1]{\draw[line width=.55pt, line cap=round, line join=round]
  ({#1-0.115},0.15) -- ({#1+0.10},0) -- ({#1-0.115},-0.15);}
\providecommand{\lss}{}
\renewcommand{\lss}[1]{\draw[line width=.55pt, line cap=round, line join=round]
  ({#1+0.115},0.15) -- ({#1-0.10},0) -- ({#1+0.115},-0.15);}
\providecommand{\spec}{}
\renewcommand{\spec}[1]{\node[below=1.2pt of #1, inner sep=1pt]
  {\scriptsize\itshape s};}
\providecommand{\specr}{}
\renewcommand{\specr}[1]{\node[right=2.2pt of #1, inner sep=1pt]
  {\scriptsize\itshape s};}
\providecommand{\mult}{}
\renewcommand{\mult}[1]{\node[above=0.5pt of #1, inner sep=.5pt]
  {\scriptsize$\times$};}
\providecommand{\edots}{}
\renewcommand{\edots}[3]{%
  \draw[dynedge] (#1) -- ++(0.34,0);
  \draw[dynedge] (#2) -- ++(-0.34,0);
  \node at #3 {\scriptsize$\cdots\cdots$};}
\providecommand{\axdots}{}
\renewcommand{\axdots}[1]{\node at (#1,0) {\scriptsize$\cdots\cdots$};}
\bigskip

\btanchor{1.3.18}

\para{1.3.18}
Let $\Wt$ be the normalizer of $\bW$ in the group of automorphisms of $\bA$; it is also the group of automorphisms of the polysimplicial complex $\bA$. Let $\Vt$ be the subgroup of translations belonging to $\Wt$, let $x$ be a special point of $\bA$ adherent to $\bC$, and let $\mathrm{D}$ be the vector chamber satisfying $\bC\subset x+\mathrm{D}$.  For every subset $\Omega$ of $\bA$ one puts
\[
  \Wd{\Omega}=\{\,w\in\Wt \mid w(\Omega)=\Omega\,\}.
\]
Then (\textit{loc.\ cit.}, ch.~VI, \S2):

\begin{enumerate}[label=(\arabic*),leftmargin=2.4em,itemsep=.35em]
\item $\Wt$ is the semidirect product of $\Wd{\bC}$ by the normal subgroup
      $\bW$.

\item $\Wt$ is the semidirect product of $\Wd{x}$ by $\Vt$, and the map
      $w\mapsto \vw$ is an isomorphism of $\Wd{x}$ onto $\lsv{\Wt}$.

\item $\Vt$ is the weight group of the root system inverse to $\vSig$, and it
      acts simply transitively on the set of special points of $\bA$.

\item $\Wd{x}$ is the semidirect product of $\Wd{(x+\mathrm{D})}$ by $\bW_x$.
      If one lets $\Wd{(x+\mathrm{D})}$ act on the dual of $\vA$ by the       contragredient of the natural action on $\vA$, then       $\Wd{(x+\mathrm{D})}$ fixes $\vSig$ and $\mathrm{B}(\mathrm{D})$, whence       a homomorphism of $\Wd{(x+\mathrm{D})}$ into the automorphism group of       the Dynkin graph of $\vSig$; this homomorphism is bijective.

\item The elements of $\Wd{\bC}$ permute the walls of $\bC$; one deduces a
      homomorphism of $\Wd{\bC}$ into the automorphism group of $(\bW,\bS)$,       equivalently of the Coxeter graph of $\bW$, and this homomorphism is       bijective.  If $\bW$ is irreducible, the image of       $\Wd{(x+\mathrm{D})}$ under this isomorphism is the fixer subgroup of       the special vertex of $\Cox\bW$ associated with $x$.

\item Put $\Wt_{\mathrm{int}}=\bW_x\cdot\Vt$.  This is a normal subgroup of
      $\Wt$, and $\Wt$ is the semidirect product of $\Wd{(x+\mathrm{D})}$ by       $\Wt_{\mathrm{int}}$.
\end{enumerate}
\bigskip

\btsecB{1.4}{1.4. \'Echelonnages}

\noindent\footnotesize
(Note added in proof.)\ \ There is an equivalence between the notion of
\emph{\'echelonnage} used here and what I.~G.~Macdonald calls an
\emph{affine root system}.
\normalsize

\medskip
The notation of no.~1.3 is retained.

\para{1.4.1}\btanchor{1.4.1}
Let $\Phi$ be a not necessarily reduced root system in the dual of $\vA$. The following two properties are equivalent:
\begin{enumerate}[label=(\roman*),leftmargin=2.4em,itemsep=.3em]
\item \textit{$\Phi$ and the reduced root system $\vSig$ have the same Weyl
      group $\vW$;}
\item \textit{for every $a\in\Phi$ there is a real number $\lambda(a)>0$ with
      $\lambda(a)a\in\vSig$, and the map $a\mapsto\lambda(a)a$ from $\Phi$ to       $\vSig$ is surjective.}
\end{enumerate}

\smallskip
In what follows $\Phi$ is always assumed to satisfy these conditions.  One sees readily that if $\vSig$ is irreducible and not of type $B_n$, $C_n$, $G_2$ or $F_4$, then $\Phi=\lambda\,\vSig$ for some $\lambda\in\Rpos$.  If on the other hand $\vSig$ is of type $B_n$, $C_n$, $F_4$ (resp.\ $G_2$), then $\Phi$ may be of type $C_n$, $B_n$, $F_4$ (resp.\ $G_2$), and one may have $\lambda(a)=2\lambda(b)$ (resp.\ $\lambda(a)=3\lambda(b)$) with $a$ running over the short roots and $b$ over the long roots of $\Phi$; moreover, if $\vSig$ is of type $B_n$ or $C_n$, then $\Phi$ may be of type $BC_n$.

\medskip
An \textbf{\'echelonnage of $\Phi$ by $\bSig$} is a subset $\Ech\subset\Phi\times\bSig$ with the following properties:

\begin{itemize}[leftmargin=3.2em,itemsep=.3em]
\item[(E\,1)] \textit{if $(a,\alpha)\in\Ech$, then $a\in\Rpos\cdot\valpha$;}
\item[(E\,2)] \textit{if $(a,\alpha)\in\Ech$ and $w\in\bW$, then
      $(\vw(a),w(\alpha))\in\Ech$;}
\item[(E\,3)] \textit{the restrictions to $\Ech$ of the two projections
      $\pr_1\colon\Phi\times\bSig\to\Phi$ and       $\pr_2\colon\Phi\times\bSig\to\bSig$ are surjective.}
\end{itemize}

\smallskip
Two roots $a\in\Phi$ and $\alpha\in\bSig$ with $(a,\alpha)\in\Ech$ are said to be \textbf{associated} by the \'echelonnage $\Ech$.

\para{1.4.2}\btanchor{1.4.2}
If $x\in\bA$ and $F$ is the facet of $\bA$ containing $x$, then the set $\Phi_x=\Phi_F$ of roots $a\in\Phi$ associated with the affine roots whose wall contains $x$ is a root system whose Weyl group is the canonical image in $\vW$ of the stabilizer of $x$.  It is called the \textbf{root system attached to $x$} (or to $F$) \textbf{by $\Ech$}.

\para{1.4.3}\btanchor{1.4.3}
Let $\Phi'\subset(\vA)^{*}$ be another root system.  Two \'echelonnages $\Ech\subset\Phi\times\bSig$ and $\Ech'\subset\Phi'\times\bSig$ are said to be
\textbf{homothetic} if there is a linear transformation
$\lambda\colon(\vA)^{*}\to(\vA)^{*}$ such that $(a,\alpha)\in\Ech$ if and only if $(\lambda(a),\alpha)\in\Ech'$.  In that case $\lambda$ induces on each $(\vA_i)^{*}$ a homothety of positive ratio.  An \'echelonnage is said to be
\textbf{similar} to $\Ech$ if it is isomorphic (in the evident sense) to an
\'echelonnage homothetic to $\Ech$.

\para{1.4.4}\btanchor{1.4.4}
In the remainder of this number one considers quadruples $(G,f,M,F)$ consisting of a graph $G$, a function $f$ from the set of edges of $G$ to $\{3,4,6,\infty\}$, a set $M$ of vertices of $G$, and a set $F$ of pairs $(r,s)$ of vertices such that $\{r,s\}$ is an edge.  The vertices belonging to $M$ are called \textbf{multipliable}.  The pair $(G,f)$ is the
\textbf{underlying Coxeter graph} of $(G,f,M,F)$.  Such a quadruple is drawn
according to the following conventions: an edge $\{r,s\}$ is drawn as a single, double, triple or thick stroke according as $f(\{r,s\})=3,\ 4,\ 6$ or $\infty$; it carries an inequality sign $>$ running from $r$ towards $s$ when $(r,s)\in F$; and multipliable vertices are marked with a cross.

\medskip
Let $\Pi$ be a finite subset of $(\vA)^{*}\smallsetminus\{0\}$ such that for every pair $(a,b)$ of elements of $\Pi$ one has
\[
  \frac{-2(a\mid b)}{(a\mid a)}\in\NN,
  \qquad\text{or}\quad a=2b,
  \qquad\text{or else}\quad b=2a,
\]
where $(\ \mid\ )$ denotes a scalar product on $(\vA)^{*}$ invariant under $\vW$.  The \textbf{Dynkin graph of $\Pi$} is the quadruple $(G,f,M,F)$ defined as follows: the vertices of $G$ are the elements $a$ of $\Pi$ with $a/2\notin\Pi$; a pair $\{a,b\}$ of vertices is an edge if $(a\mid b)\neq 0$; the function $f$ is defined by
\[
  \cos^{2}\frac{\pi}{f(\{a,b\})}
  =\frac{(a\mid b)^{2}}{(a\mid a)(b\mid b)}\,;
\]
a vertex $a$ belongs to $M$ if $2a\in\Pi$; and a pair $(a,b)$ of vertices belongs to $F$ if $(a\mid a)>(b\mid b)$.

\medskip
Let $\Phi$ be a root system in a subspace of $(\vA)^{*}$ whose Weyl group leaves invariant the restriction of $(\ \mid\ )$ to that subspace.  The
\textbf{Dynkin graph of $\Phi$}, written $\Dyn\Phi$, is by definition the
Dynkin graph of the set of elements of $\Phi$ that are positive multiples of the elements of a base.  Up to canonical isomorphism this graph does not depend on the chosen base.

\para{1.4.5}\btanchor{1.4.5}
The \textbf{Dynkin graph of an \'echelonnage} $\Ech\subset\Phi\times\bSig$, written $\Dyn\Ech$, is the Dynkin graph of the set of elements of $\Phi$ associated with the affine roots that contain $\bC$ and whose wall contains a panel of $\bC$.  Up to canonical isomorphism this graph depends only on $\Ech$, and not on the choice of fundamental chamber $\bC$; its vertex set is in canonical bijection with the set of panels of $\bC$, hence also with $\bS$.

\smallskip
Knowing $\Dyn\Ech$, one recovers the Coxeter graph of $\bW$ and the graphs $\Dyn\Phi$ and $\Dyn\Phi_x$ for every $x\in\bA$ (hence also the systems $\Phi$, $\bSig$ and $\Phi_x$ up to isomorphism) by the following rules, whose proofs are immediate.

\smallskip
\begin{enumerate}[label=(\arabic*),leftmargin=2.4em,itemsep=.5em]
\item \textit{The Coxeter graph of $\bW$ is the Coxeter graph underlying
      $\Dyn\Ech$.}

      \smallskip
      Suppose $\Dyn\Ech$ connected.  If its rank (that is, its number of       vertices minus one) equals $1$, then $\Dyn\Phi$ has a single vertex,       always multipliable except when $\Dyn\Ech$ is the graph
      \begin{dyngraph}\dvo{0}{a}\dvo{1}{b}\ethick{a}{b}\end{dyngraph}.
      If $\Dyn\Ech$ has rank $\geq 2$ and is not one of the following graphs:

\item \[
      \left\{
      \begin{array}{c}
      \begin{dyngraph}
        \dvo{0}{a}\dvo{1}{b}\dvo{2}{c}\dvo{3.5}{d}\dvo{4.5}{e}\dvo{5.5}{f}
        \edouble{a}{b}\gtr{0.5}
        \esingle{b}{c}\edots{c}{d}{(2.75,0)}\esingle{d}{e}
        \edouble{e}{f}\gtr{5.0}\mult{f}
      \end{dyngraph}\\[1.1em]
      \begin{dyngraph}
        \dvo{0}{a}\dvo{1}{b}\dvo{2}{c}\dvo{3.5}{d}\dvo{4.5}{e}\dvo{5.5}{f}
        \edouble{a}{b}\lss{0.5}
        \esingle{b}{c}\edots{c}{d}{(2.75,0)}\esingle{d}{e}
        \edouble{e}{f}\gtr{5.0}\mult{f}
      \end{dyngraph}\\[1.1em]
      \begin{dyngraph}
        \dvo{0}{a}\dvo{1}{b}\dvo{2}{c}\dvo{3.5}{d}\dvo{4.5}{e}\dvo{5.5}{f}
        \edouble{a}{b}\gtr{0.5}
        \esingle{b}{c}\edots{c}{d}{(2.75,0)}\esingle{d}{e}
        \edouble{e}{f}\gtr{5.0}
      \end{dyngraph}
      \end{array}
      \right\}
      \quad (n+1\ \text{vertices})
      \]
      then $\Dyn\Phi$ is obtained from $\Dyn\Ech$ by deleting a special
      vertex, and the result does not depend on which special vertex is
      chosen.  Finally, if $\Dyn\Ech$ is one of the graphs (2), then
      $\Dyn\Phi$ is the graph
      \[
      \begin{dyngraph}
        \dvo{0}{a}\dvo{1}{b}\dvo{2}{c}\dvo{3.5}{d}\dvo{4.5}{e}\dvo{5.5}{f}
        \esingle{a}{b}\esingle{b}{c}\edots{c}{d}{(2.75,0)}\esingle{d}{e}
        \edouble{e}{f}\gtr{5.0}\mult{f}
      \end{dyngraph}
      \qquad (n\ \text{vertices}).
      \]

      \smallskip
      When $\Dyn\Ech$ is not connected, $\Dyn\Phi$ is obtained by applying the
      preceding rules to the connected components of $\Dyn\Ech$.

\item \textit{If $x\in\bA$ lies in a facet of type $X\subset\bS$, the graph
      $\Dyn\Phi_x$ is the subgraph of $\Dyn\Ech$ whose vertices correspond to
      the elements of $X$.}
\end{enumerate}

\begin{BT}{Classification}{1.4.6}\btanchor{1.4.6}
An \'echelonnage is determined up to similarity by its Dynkin graph.  The
table below is the complete list of the Dynkin graphs of the irreducible
\'echelonnages.
\end{BT}

\noindent
The proof proceeds by classifying, for each irreducible affine Weyl group,
the corresponding \'echelonnages.  For $\bW$ of type $C_n$, for instance,
inspection of the classification of root systems shows that up to homothety
$\Phi$ is one of: the system $\Phi_0=\vSig$ of type $C_n$; the system
$\Phi_1$ of type $B_n$ obtained by keeping the long roots of $\vSig$ and
replacing the short roots by their doubles; or the system
$\Phi_2=\Phi_0\cup\Phi_1$ of type $BC_n$ --- these being the only systems
admitting $\vW$ as Weyl group.  For $\Phi=\Phi_0$ and for $\Phi=\Phi_1$ there
is exactly one \'echelonnage of $\Phi$ by $\bSig$, since for each
$\alpha\in\bSig$ there is exactly one root $a\in\Phi$ with
$\valpha\in\RR_{+}\!\cdot a$; for $\Phi=\Phi_2$ there are four pairwise
non-similar ones.  The remaining types are handled in the same way, more
easily.  (See the original text, pp.~28--29, for the details.)

\bigskip
\hrule
\bigskip

\btsecC{}{Table of the irreducible \'echelonnages}

\renewcommand{\arraystretch}{1.9}
\begin{longtable}{@{}l@{\hspace{1.1em}}c@{\hspace{1.1em}}c@{\hspace{1.1em}}c@{\hspace{2.2em}}l@{}}
\toprule[.9pt]
\textbf{Name} & \textbf{Rank} & \textbf{Type of $\vSig$} & \textbf{Type of $\Phi$}
  & \textbf{Dynkin graph and special vertices}\\
\midrule[.6pt]
\endfirsthead
\toprule[.9pt]
\textbf{Name} & \textbf{Rank} & \textbf{Type of $\vSig$} & \textbf{Type of $\Phi$}
  & \textbf{Dynkin graph and special vertices}\\
\midrule[.6pt]
\endhead
\bottomrule[.9pt]
\endfoot

$A_1$ & $1$ & $A_1$ & $A_1$ &
\begin{dyngraph}
  \dvo{0}{a}\dvo{1}{b}\ethick{a}{b}\spec{a}\spec{b}
\end{dyngraph}\\

$C\text{-}BC_1^{\mathrm{I}}$ & $1$ & $A_1$ & $BC_1$ &
\begin{dyngraph}
  \dvo{0}{a}\dvo{1}{b}\ethick{a}{b}\spec{a}\spec{b}\mult{b}
\end{dyngraph}\\

$C\text{-}BC_1^{\mathrm{II}}$ & $1$ & $A_1$ & $BC_1$ &
\begin{dyngraph}
  \dvo{0}{a}\dvo{1}{b}\ethick{a}{b}\spec{a}\spec{b}\mult{a}\mult{b}
\end{dyngraph}\\

$C\text{-}BC_1^{\mathrm{III}}$ & $1$ & $A_1$ & $BC_1$ &
\begin{dyngraph}
  \dvo{0}{a}\dvo{1}{b}\ethick{a}{b}\gtr{0.5}\spec{a}\spec{b}
\end{dyngraph}\\

$C\text{-}BC_1^{\mathrm{IV}}$ & $1$ & $A_1$ & $BC_1$ &
\begin{dyngraph}
  \dvo{0}{a}\dvo{1}{b}\ethick{a}{b}\gtr{0.5}\spec{a}\spec{b}\mult{b}
\end{dyngraph}\\

$A_n$ & $n\geq 2$ & $A_n$ & $A_n$ &
\begin{dyngraph}
  \dv{0}{0}{L}\spec{L}
  \dv{1}{0.62}{u1}\dv{2}{0.62}{u2}\dv{3.5}{0.62}{u3}\dv{4.5}{0.62}{u4}
  \dv{1}{-0.62}{d1}\dv{2}{-0.62}{d2}\dv{3.5}{-0.62}{d3}\dv{4.5}{-0.62}{d4}
  \dv{5.5}{0}{R}\spec{R}
  \esingle{L}{u1}\esingle{u1}{u2}\esingle{u3}{u4}\esingle{u4}{R}
  \edots{u2}{u3}{(2.75,0.62)}
  \esingle{L}{d1}\esingle{d1}{d2}\esingle{d3}{d4}\esingle{d4}{R}
  \edots{d2}{d3}{(2.75,-0.62)}
  \spec{u1}\spec{u2}\spec{u3}\spec{u4}
  \spec{d1}\spec{d2}\spec{d3}\spec{d4}
\end{dyngraph}\\

$B_n$ & $n\geq 3$ & $B_n$ & $B_n$ &
\begin{dyngraph}
  \dv{0}{0.55}{f1}\dv{0}{-0.55}{f2}\spec{f1}\spec{f2}
  \dvo{0.95}{v1}\dvo{1.95}{v2}\dvo{3.35}{v3}\dvo{4.35}{v4}\dvo{5.35}{v5}
  \esingle{f1}{v1}\esingle{f2}{v1}\esingle{v1}{v2}\esingle{v3}{v4}
  \edots{v2}{v3}{(2.65,0)}
  \edouble{v4}{v5}\gtr{4.85}
\end{dyngraph}\\

$B\text{-}C_n$ & $n\geq 3$ & $B_n$ & $C_n$ &
\begin{dyngraph}
  \dv{0}{0.55}{f1}\dv{0}{-0.55}{f2}\spec{f1}\spec{f2}
  \dvo{0.95}{v1}\dvo{1.95}{v2}\dvo{3.35}{v3}\dvo{4.35}{v4}\dvo{5.35}{v5}
  \esingle{f1}{v1}\esingle{f2}{v1}\esingle{v1}{v2}\esingle{v3}{v4}
  \edots{v2}{v3}{(2.65,0)}
  \edouble{v4}{v5}\lss{4.85}
\end{dyngraph}\\

$B\text{-}BC_n$ & $n\geq 3$ & $B_n$ & $BC_n$ &
\begin{dyngraph}
  \dv{0}{0.55}{f1}\dv{0}{-0.55}{f2}\spec{f1}\spec{f2}
  \dvo{0.95}{v1}\dvo{1.95}{v2}\dvo{3.35}{v3}\dvo{4.35}{v4}\dvo{5.35}{v5}
  \esingle{f1}{v1}\esingle{f2}{v1}\esingle{v1}{v2}\esingle{v3}{v4}
  \edots{v2}{v3}{(2.65,0)}
  \edouble{v4}{v5}\gtr{4.85}\mult{v5}
\end{dyngraph}\\

$C_n$ & $n\geq 2$ & $C_n$ & $C_n$ &
\begin{dyngraph}
  \dvo{0}{v0}\dvo{1}{v1}\dvo{2}{v2}\dvo{3.5}{v3}\dvo{4.5}{v4}\dvo{5.5}{v5}
  \edouble{v0}{v1}\gtr{0.5}
  \esingle{v1}{v2}\edots{v2}{v3}{(2.75,0)}\esingle{v3}{v4}
  \edouble{v4}{v5}\lss{5.0}
  \spec{v0}\spec{v5}
\end{dyngraph}\\

$C\text{-}B_n$ & $n\geq 2$ & $C_n$ & $B_n$ &
\begin{dyngraph}
  \dvo{0}{v0}\dvo{1}{v1}\dvo{2}{v2}\dvo{3.5}{v3}\dvo{4.5}{v4}\dvo{5.5}{v5}
  \edouble{v0}{v1}\lss{0.5}
  \esingle{v1}{v2}\edots{v2}{v3}{(2.75,0)}\esingle{v3}{v4}
  \edouble{v4}{v5}\gtr{5.0}
  \spec{v0}\spec{v5}
\end{dyngraph}\\

$C\text{-}BC_n^{\mathrm{I}}$ & $n\geq 2$ & $C_n$ & $BC_n$ &
\begin{dyngraph}
  \dvo{0}{v0}\dvo{1}{v1}\dvo{2}{v2}\dvo{3.5}{v3}\dvo{4.5}{v4}\dvo{5.5}{v5}
  \edouble{v0}{v1}\lss{0.5}
  \esingle{v1}{v2}\edots{v2}{v3}{(2.75,0)}\esingle{v3}{v4}
  \edouble{v4}{v5}\gtr{5.0}\mult{v5}
  \spec{v0}\spec{v5}
\end{dyngraph}\\

$C\text{-}BC_n^{\mathrm{II}}$ & $n\geq 2$ & $C_n$ & $BC_n$ &
\begin{dyngraph}
  \dvo{0}{v0}\dvo{1}{v1}\dvo{2}{v2}\dvo{3.5}{v3}\dvo{4.5}{v4}\dvo{5.5}{v5}
  \edouble{v0}{v1}\lss{0.5}
  \esingle{v1}{v2}\edots{v2}{v3}{(2.75,0)}\esingle{v3}{v4}
  \edouble{v4}{v5}\gtr{5.0}
  \mult{v0}\mult{v5}\spec{v0}\spec{v5}
\end{dyngraph}\\

$C\text{-}BC_n^{\mathrm{III}}$ & $n\geq 2$ & $C_n$ & $BC_n$ &
\begin{dyngraph}
  \dvo{0}{v0}\dvo{1}{v1}\dvo{2}{v2}\dvo{3.5}{v3}\dvo{4.5}{v4}\dvo{5.5}{v5}
  \edouble{v0}{v1}\gtr{0.5}
  \esingle{v1}{v2}\edots{v2}{v3}{(2.75,0)}\esingle{v3}{v4}
  \edouble{v4}{v5}\gtr{5.0}
  \spec{v0}\spec{v5}
\end{dyngraph}\\

$C\text{-}BC_n^{\mathrm{IV}}$ & $n\geq 2$ & $C_n$ & $BC_n$ &
\begin{dyngraph}
  \dvo{0}{v0}\dvo{1}{v1}\dvo{2}{v2}\dvo{3.5}{v3}\dvo{4.5}{v4}\dvo{5.5}{v5}
  \edouble{v0}{v1}\gtr{0.5}
  \esingle{v1}{v2}\edots{v2}{v3}{(2.75,0)}\esingle{v3}{v4}
  \edouble{v4}{v5}\gtr{5.0}\mult{v5}
  \spec{v0}\spec{v5}
\end{dyngraph}\\

$D_n$ & $n\geq 4$ & $D_n$ & $D_n$ &
\begin{dyngraph}
  \dv{0}{0.55}{f1}\dv{0}{-0.55}{f2}\spec{f1}\spec{f2}
  \dvo{0.95}{v1}\dvo{1.95}{v2}\dvo{3.35}{v3}\dvo{4.35}{v4}
  \dv{5.3}{0.55}{g1}\dv{5.3}{-0.55}{g2}\spec{g1}\spec{g2}
  \esingle{f1}{v1}\esingle{f2}{v1}\esingle{v1}{v2}\esingle{v3}{v4}
  \edots{v2}{v3}{(2.65,0)}
  \esingle{v4}{g1}\esingle{v4}{g2}
\end{dyngraph}\\

$E_6$ & $6$ & $E_6$ & $E_6$ &
\begin{dyngraph}
  \dvo{0}{e1}\dvo{1}{e2}\dvo{2}{e3}\dvo{3}{e4}\dvo{4}{e5}
  \esingle{e1}{e2}\esingle{e2}{e3}\esingle{e3}{e4}\esingle{e4}{e5}
  \dv{2}{1}{m1}\dv{2}{2}{m2}
  \esingle{e3}{m1}\esingle{m1}{m2}
  \spec{e1}\spec{e5}\specr{m2}
\end{dyngraph}\\

$E_7$ & $7$ & $E_7$ & $E_7$ &
\begin{dyngraph}
  \dvo{0}{e1}\dvo{1}{e2}\dvo{2}{e3}\dvo{3}{e4}\dvo{4}{e5}\dvo{5}{e6}\dvo{6}{e7}
  \esingle{e1}{e2}\esingle{e2}{e3}\esingle{e3}{e4}
  \esingle{e4}{e5}\esingle{e5}{e6}\esingle{e6}{e7}
  \dv{3}{1}{m1}\esingle{e4}{m1}
  \spec{e1}\spec{e7}
\end{dyngraph}\\

$E_8$ & $8$ & $E_8$ & $E_8$ &
\begin{dyngraph}
  \dvo{0}{e1}\dvo{1}{e2}\dvo{2}{e3}\dvo{3}{e4}
  \dvo{4}{e5}\dvo{5}{e6}\dvo{6}{e7}\dvo{7}{e8}
  \esingle{e1}{e2}\esingle{e2}{e3}\esingle{e3}{e4}\esingle{e4}{e5}
  \esingle{e5}{e6}\esingle{e6}{e7}\esingle{e7}{e8}
  \dv{5}{1}{m1}\esingle{e6}{m1}
  \spec{e1}
\end{dyngraph}\\

$F_4$ & $4$ & $F_4$ & $F_4$ &
\begin{dyngraph}
  \dvo{0}{a}\dvo{1}{b}\dvo{2}{c}\dvo{3}{d}\dvo{4}{e}
  \esingle{a}{b}\esingle{b}{c}\edouble{c}{d}\gtr{2.5}\esingle{d}{e}
  \spec{a}
\end{dyngraph}\\

$F_4^{\mathrm{I}}$ & $4$ & $F_4$ & $F_4$ &
\begin{dyngraph}
  \dvo{0}{a}\dvo{1}{b}\dvo{2}{c}\dvo{3}{d}\dvo{4}{e}
  \esingle{a}{b}\esingle{b}{c}\edouble{c}{d}\lss{2.5}\esingle{d}{e}
  \spec{a}
\end{dyngraph}\\

$G_2$ & $2$ & $G_2$ & $G_2$ &
\begin{dyngraph}
  \dvo{0}{a}\dvo{1}{b}\dvo{2.1}{c}
  \esingle{a}{b}\etriple{b}{c}\gtr{1.55}
  \spec{a}
\end{dyngraph}\\

$G_2^{\mathrm{I}}$ & $2$ & $G_2$ & $G_2$ &
\begin{dyngraph}
  \dvo{0}{a}\dvo{1}{b}\dvo{2.1}{c}
  \esingle{a}{b}\etriple{b}{c}\lss{1.55}
  \spec{a}
\end{dyngraph}\\
\end{longtable}
\renewcommand{\arraystretch}{1.0}

\bigskip
\hrule
\bigskip

\btsecB{1.5}{1.5. Conventions and notation}

\para{1.5.1}\btanchor{1.5.1}
A Tits system $(G,B,N,S)$ with Weyl group $W$ is said to be \textbf{of affine
type} if there exist an affine Weyl group $\bW$ acting on a euclidean affine
space $\bA$, a chamber $\bC$ of $\bA$, and an isomorphism of the Coxeter
system $(W,S)$ onto the Coxeter system $(\bW,\bS)$, where $\bS$ is the set of
reflections in the walls of $\bC$ (\btref{1.3.4}).  The \textbf{rank} of such a Tits
system is the rank of $\bW$, that is, the dimension of $\bA$ (\btref{1.3.4}).

\smallskip
A \textbf{parahoric subgroup} of $G$ is then a parabolic subgroup of $G$ whose
type $X$ is, after the identification of $W$ with $\bW$, an element of the set
$\bT$ of (\btref{1.3.5}) --- that is, is such that $\bW_X$ is \emph{finite}.  A
subgroup $P$ containing $B$ is therefore parahoric if and only if the set of
double cosets $B\backslash P/B$ is finite.  When $\bW$ is irreducible the
parahoric subgroups are exactly the proper parabolic subgroups, i.e.\ those
different from $G$.  Every parabolic subgroup contained in a parahoric
subgroup is parahoric; the minimal parabolic subgroups, i.e.\ the conjugates
of $B$, are parahoric.

\para{1.5.2}\btanchor{1.5.2}
\textit{Throughout the rest of the chapter the hypotheses and notation of
\emph{1.3} are retained.}  In particular $\bW$ denotes an affine Weyl group
acting on a euclidean affine space $\bA$, and the symbols $\vA$, $\bC$,
$\bSig$, etc.\ have the same meaning as in 1.3.

\smallskip
\textit{In \S\S2--4, $(G,B,N,S)$ denotes a saturated Tits system of affine
type.}  The Weyl group $W$ of $G$ is identified with $\bW$ as in (\btref{1.5.1}), and
the symbols $B_X$, $H$, etc.\ have the same meaning as in 1.2.

\bigskip
\hrule
\bigskip

\btsecA{2}{2. The building of a Tits system of affine type}

With the Tits system $(G,B,N,S)$ of (\btref{1.5.2}) one associates a set $\Imm$,
equipped with: a structure of polysimplicial complex (\btref{2.1}); a set of chambered
morphisms of polysimplicial complexes from $\bA$ into $\Imm$, called the
\emph{structural maps} of $\Imm$ (\btref{2.2}); a distance making it a complete metric
space, for which the structural maps are isometries (\btref{2.5}); and an action of
$G$ respecting all of these structures (\btref{2.1.4}).  The object $\Imm$ so defined
is called the \textbf{building} of the Tits system $(G,B,N,S)$, or, by abuse
of language, of $G$.  It depends only on the conjugacy class of the pair
$(B,N)$ (\btref{2.7}).

\btsecB{2.1}{2.1. The polysimplicial complex \texorpdfstring{$\mathscr{I}$}{I}}

\para{2.1.1}\btanchor{2.1.1}
Let $\Par$ be the set of parahoric subgroups of $G$.  For $P\in\Par$ let
$\tau(P)\in\bT$ denote the \emph{type} of $P$ (\btref{1.2.10}).  Let $\Imm$ be the set
of pairs $(P,x)$ with $P\in\Par$ and $x\in\bC_{\tau(P)}$ (\btref{1.3.5}).  For
$P\in\Par$ the set
\[
  F=F(P)=\{(P,x)\mid x\in\bC_{\tau(P)}\}
\]
is called a \textbf{facet} of $\Imm$, of \emph{type} $\tau(F)=\tau(P)$ and of
\emph{codimension} $\Card\tau(P)$.  If $P'\in\Par$ and $P\subset P'$, one says
that $F(P')$ is a facet of $F(P)$; one then has $\tau(P')\supset\tau(P)$.
Conversely $F(P)$ possesses, for each type $Y\supset\tau(P)$, one and only one
facet of type $Y$ (\btref{1.2.10}).  We write $\overline{F}$ for the union of the
facets of a given facet $F$.

\para{2.1.2}\btanchor{2.1.2}
The map $(P,x)\mapsto x$ is called the \textbf{canonical map} of $\Imm$ into
$\overline{\bC}$.  If $F$ is a facet of $\Imm$ of type $X$, the restriction of
this map to $F$ (resp.\ $\overline{F}$) is a bijection of $F$ (resp.\
$\overline{F}$) onto $\bC_X$ (resp.\ $\overline{\bC}_X$), whose inverse is
called the canonical map of $\bC_X$ (resp.\ $\overline{\bC}_X$) onto $F$
(resp.\ $\overline{F}$).  \textit{We endow $F$ (resp.\ $\overline{F}$) with
the structure of open (resp.\ closed) affine polysimplex obtained by transport
of structure from that of $\bC_X$ (resp.\ $\overline{\bC}_X$) via this
canonical map.}  Writing $\Fac$ for the set of facets of $\Imm$ and
$F'\prec F$ for the incidence relation $F'\subset\overline{F}$, the conditions
of (\btref{1.1.4}) are satisfied.  The \textbf{chambers} of $\Imm$ are the facets of
type $\emptyset$, associated with the minimal parahoric subgroups.  At the
other extreme, if $P$ is a maximal parahoric subgroup, the facet $F(P)$
contains a single point, written $a_P$ and called the \textbf{vertex} of
$\Imm$ associated with $P$.

\para{2.1.3}\btanchor{2.1.3}
If two chambers of $\Imm$ are \emph{adjacent}, that is, distinct with a panel
in common (\btref{1.1.5}), this panel is well determined: indeed, if $F(P_1)$ and
$F(P_2)$ are two distinct panels of one chamber $F(P)$, then $P=P_1\cap P_2$.

\smallskip
This permits the following definition.  If $\Gamma=(C_0,\dots,C_n)$ is a
non-stuttering gallery of $\Imm$ (\btref{1.1.5}), the \textbf{type} of $\Gamma$ is the
sequence $\mathbf{s}(\Gamma)=(s_1,\dots,s_n)$ of elements of $\bS$ such that
the panel common to the adjacent chambers $C_{i-1}$ and $C_i$ is of type $s_i$
for $1\leq i\leq n$.

\para{2.1.4}\btanchor{2.1.4}
The group $G$ acts on $\Imm$ by
\begin{equation}
  g\cdot(P,x)=(gPg^{-1},x)
  \qquad (P\in\Par,\ x\in\bC_{\tau(P)}).
\end{equation}
This is indeed an action, since $gPg^{-1}$ is a parahoric subgroup of the same
type as $P$.

\begin{BTu}{Proposition}
\begin{enumerate}[label=(\roman*),leftmargin=2.4em,itemsep=.25em]
\item For every type $X\in\bT$, the group $G$ permutes the facets of type $X$
      transitively.
\item For every parahoric subgroup $P$, the restriction to $F(P)$ (resp.\
      $\overline{F(P)}$) of the action of $g\in G$ is an affine bijection of
      $F(P)$ (resp.\ $\overline{F(P)}$) onto $F(gPg^{-1})$ (resp.\
      $\overline{F(gPg^{-1})}$), whose composite with the canonical map of
      $\bC_{\tau(P)}$ onto $F(P)$ is the canonical map of $\bC_{\tau(P)}$ onto
      $F(gPg^{-1})$.
\end{enumerate}
\end{BTu}
\proofref{(i) follows from the conjugacy of parahoric subgroups of a given
type (\btref{1.2.10}); (ii) is immediate}

\begin{BT}{Proposition}{2.1.5}\btanchor{2.1.5}
Let $Q$ be a parahoric subgroup of $G$.
\begin{enumerate}[label=(\roman*),leftmargin=2.4em,itemsep=.25em]
\item $Q$ is the stabilizer in $G$ of the facet $F(Q)$ and of each of its
      points;
\item $\overline{F(Q)}$ is the set of points of $\Imm$ fixed by $Q$.
\end{enumerate}
\end{BT}
\proofref{both follow at once from the fact that a parahoric subgroup is its
own normalizer}

\begin{BT}{Corollary}{2.1.6}\btanchor{2.1.6}
Let $C$ be a chamber of $\Imm$.  Then $\overline{C}$ is a fundamental domain
for $G$ in $\Imm$.
\end{BT}
\proofref{every facet of $\Imm$ is a transform of exactly one facet of $C$;
(\btref{2.1.5}) then shows each orbit meets $\overline{C}$ in exactly one point}

\begin{BT}{Proposition}{2.1.7}\btanchor{2.1.7}
Let $C$ and $C'$ be two chambers of $\Imm$.  There exists one and only one
element $w=w(C,C')\in\bW$ such that $g^{-1}g'\in BwB$ for all $g,g'\in G$ with
$C=g\cdot F(B)$ and $C'=g'\cdot F(B)$.  One has $w(h\cdot C,h\cdot C')=w(C,C')$
for every $h\in G$, and $w(C',C)=w(C,C')^{-1}$.  Setting $w(C)=w(F(B),C)$: for
every $n\in\nu^{-1}(w(C))$ there exists $b\in B$ with $C=bn\cdot F(B)$.
\end{BT}
\proofref{p.~33}

\begin{BT}{Lemma}{2.1.8}\btanchor{2.1.8}
In order that a chamber $C=bn\cdot F(B)$ be adjacent to $F(B)$ it is necessary
and sufficient that $\nu(n)\in\bS$.  The panel common to $C$ and $F(B)$ is
then of type $\nu(n)$.
\end{BT}
\proofref{pp.~33--34}

\begin{BT}{Proposition}{2.1.9}\btanchor{2.1.9}
Let $\Gamma=(C_0,\dots,C_m)$ be a sequence of chambers with $C_0=F(B)$.  The
following conditions are equivalent:
\begin{enumerate}[label=(\roman*),leftmargin=2.4em,itemsep=.25em]
\item $\Gamma$ is a non-stuttering gallery;
\item there exist $b_j\in B$, $s_j\in\bS$ and
      $\bar{s}_j\in\nu^{-1}(s_j)$ \emph{(}$j=1,\dots,m$\emph{)} such that
      \[
        C_j=b_1\bar{s}_1\cdots b_j\bar{s}_j\cdot F(B)
        \qquad\text{for } j=1,\dots,m.
      \]
\end{enumerate}
Under these conditions $\Gamma$ is of type $(s_1,\dots,s_m)$.
\end{BT}
\proofref{p.~34; one direction uses (\btref{2.1.8}), the other is an
induction on $j$}

\begin{BT}{Corollary}{2.1.10}\btanchor{2.1.10}
Let $C$ be a chamber of $\Imm$ and put $w=w(F(B),C)$.  For every reduced
decomposition $\mathbf{s}=(s_1,\dots,s_m)$ of $w$ there exists a minimal
gallery with origin $F(B)$, extremity $C$ and type $\mathbf{s}$.
\end{BT}
\proofref{p.~34.  The result is made more precise in (\btref{2.3.9})}

\begin{BT}{Corollary}{2.1.11}\btanchor{2.1.11}
A non-stuttering gallery is minimal if and only if its type
$\mathbf{s}=(s_1,\dots,s_m)$ is a reduced decomposition of $s_1\cdots s_m$.
\end{BT}

\begin{BT}{Proposition}{2.1.12}\btanchor{2.1.12}
The set $\Imm$, equipped with the family of facets, the incidence relation
between facets, and the affine structures on each of the sets
$\overline{F(P)}$ for $P\in\Par$, is a polysimplicial complex, and $G$ acts on
$\Imm$ by automorphisms of polysimplicial complex.
\end{BT}
\proofref{pp.~34--35}

\begin{BT}{Example}{2.1.13}\btanchor{2.1.13}
Suppose $\bA$ of dimension $1$.  Then $\Card\bS=2$ and $\bW$ is an infinite
dihedral group.  Every finite sequence $(s_j)$ of elements of $\bS$ with
$s_j\neq s_{j+1}$ for all $j$ is a reduced decomposition, and every
non-stuttering gallery $(C_0,\dots,C_m)$ in which the panels
$\overline{C}_{j-1}\cap\overline{C}_j$ and
$\overline{C}_j\cap\overline{C}_{j+1}$ are distinct for $1\leq j\leq m-1$ is
minimal.  The building $\Imm$ is a simplicial complex of dimension $1$, i.e.\
a connected graph, and the above shows that this graph contains no circuit
(which would furnish a minimal gallery of length $>0$ with origin and
extremity equal).  \emph{Hence $\Imm$ is a tree.}
\end{BT}

\btsecB{2.2}{2.2. Structural maps and apartments}

\begin{BT}{Proposition}{2.2.1}\btanchor{2.2.1}
There exists one and only one map $j\colon\bA\to\Imm$ with the following two
properties:
\begin{enumerate}[label=(\roman*),leftmargin=2.4em,itemsep=.25em]
\item the restriction of $j$ to $\overline{\bC}$ is the canonical bijection
      \emph{(\btref{2.1.2})} of $\overline{\bC}$ onto $\overline{F(B)}$;
\item $j(\nu(n)\cdot x)=n\cdot j(x)$ for all $n\in N$ and $x\in\bA$.
\end{enumerate}
The map $j$ is injective.  If $F$ is a facet of $\bA$, then $j(F)$ is a facet
of $\Imm$ of the same type, one has $j(\overline{F})=\overline{j(F)}$, and the
restriction of $j$ to $\overline{F}$ is an affine bijection of $\overline{F}$
onto $\overline{j(F)}$.
\end{BT}
\proofref{p.~35}

\begin{BT}{Definition}{2.2.2}\btanchor{2.2.2}
The map $j$ of \emph{(\btref{2.2.1})} is called the \textbf{canonical map} of $\bA$
into $\Imm$.  A map $\varphi$ of $\bA$ into $\Imm$ is called a
\textbf{structural map} of $\Imm$ if there is $g\in G$ with
$\varphi(x)=g\cdot j(x)$ for all $x\in\bA$.  An \textbf{apartment} (resp.\
\textbf{half-apartment}, \textbf{sector}, \textbf{wall}) of $\Imm$ is a subset
of $\Imm$ which is the image of $\bA$ (resp.\ of an affine root, of a sector,
of a wall of $\bA$) under a structural map.
\end{BT}

\para{2.2.3}\btanchor{2.2.3}
It follows at once from (\btref{2.2.1}) that a structural map $\varphi$ is a chambered
morphism (\btref{1.1.7}) of polysimplicial complexes from $\bA$ into $\Imm$, and is
even an isomorphism of polysimplicial complexes of $\bA$ onto the apartment
$\varphi(\bA)$, which is a ``polysimplicial subcomplex'' of $\Imm$ in the
sense that an apartment containing a facet $F$ contains $\overline{F}$, and
that every apartment is the union of the facets of the chambers it contains.

\para{2.2.4}\btanchor{2.2.4}
Let $M$ be a half-apartment, the image of an affine root $\alpha$ of $\bA$
under a structural map $\varphi$.  The wall $\varphi(\partial\alpha)$ is the
union of the sets $\overline{F}$ for $F$ running over the panels of $\Imm$
belonging to a single chamber contained in $M$.  This wall therefore depends
only on $M$ and not on the choice of $\varphi$: it is called \textbf{the wall
of $M$} and written $\partial M$.

\begin{BT}{Proposition}{2.2.5}\btanchor{2.2.5}
With the preceding notation:
\begin{enumerate}[label=(\roman*),leftmargin=2.4em,itemsep=.25em]
\item $\displaystyle H=\bigcap_{n\in N}nBn^{-1}
      =\{g\in G\mid g\cdot j(x)=j(x)\ \text{for all}\ x\in\bA\}$;
\item $N=\{g\in G\mid g\cdot j(\bA)=j(\bA)\}$.
\end{enumerate}
\end{BT}
\proofref{p.~36; (ii) uses that the Tits system $(G,B,N)$ is
saturated}

\begin{BT}{Corollary}{2.2.6}\btanchor{2.2.6}
The group $G$ permutes transitively the pairs $(F,A)$ consisting of a facet
$F$ of a given type $X$ and an apartment $A$ containing $F$.
\end{BT}

\begin{BT}{Corollary}{2.2.7}\btanchor{2.2.7}
Let $\varphi$ be a structural map and $\psi$ a map of $\bA$ into the apartment
$\varphi(\bA)$.  In order that $\psi$ be a structural map it is necessary and
sufficient that there exist $w\in\bW$ with $\psi=\varphi\circ w$.
\end{BT}

\para{2.2.8}\btanchor{2.2.8}
By (\btref{2.2.7}) there exists on every apartment $A$ one and only one structure of
euclidean affine space for which every structural map with image $A$ is an
isomorphism of euclidean affine spaces from $\bA$.  \textit{We henceforth
regard $A$ as endowed with this euclidean structure, and in particular with
the corresponding distance, written $d_A$.}  For every $g\in G$ the map
$x\mapsto g\cdot x$ is then an isomorphism of $A$ onto the apartment
$g\cdot A$.

\btsecB{2.3}{2.3. Retractions}

\begin{BT}{Proposition}{2.3.1}\btanchor{2.3.1}
Two facets of $\Imm$ are contained in a common apartment.
\end{BT}
\proofref{p.~37; reduce to two chambers via (\btref{2.2.3}), then use
(\btref{2.1.7})}

\begin{BT}{Proposition}{2.3.2}\btanchor{2.3.2}
Let $A$ and $A'$ be two apartments containing a common chamber $C$.  There
exists one and only one map $\rho=\rho_{A',A;C}$ of $A'$ onto $A$ which is the
restriction to $A'$ of an action of $G$ and satisfies $\rho(C)=C$.  The
restriction of $\rho$ to $A\cap A'$ is the identity, and $\rho$ is an isometry
of $A'$ onto $A$.
\end{BT}
\proofref{p.~37}

\begin{BT}{Corollary}{2.3.3}\btanchor{2.3.3}
Let $F_1,F_2$ be two facets and $A_1,A_2$ two apartments each containing both
$F_1$ and $F_2$.  There exists $g\in G$ with $g\cdot A_1=A_2$ and
$g\cdot F_i=F_i$ for $i=1,2$.
\end{BT}

\begin{BT}{Theorem}{2.3.4}\btanchor{2.3.4}
Let $A$ be an apartment and $C$ a chamber contained in $A$.  There exists one
and only one map $\rho$ of $\Imm$ into $A$ with the following properties:
\begin{enumerate}[label=(\roman*),leftmargin=2.4em,itemsep=.25em]
\item $\rho(C)=C$;
\item for every apartment $A'$ containing $C$ there exists $g\in G$ with
      $\rho(x)=g\cdot x$ for all $x\in A'$.
\end{enumerate}
The restriction of $\rho$ to $A$ is the identity.  If $x\in\overline{C}$ then
$\rho^{-1}(x)=\{x\}$.  If $F$ is a facet of type $X$, so is $\rho(F)$; one has
$\rho(\overline{F})=\overline{\rho(F)}$, and the restriction of $\rho$ to
$\overline{F}$ is an affine bijection of $\overline{F}$ onto
$\overline{\rho(F)}$.
\end{BT}
\proofref{pp.~37--38; the key point is that $\rho_{A',A;C}(x)$
depends only on $x$}

\begin{BT}{Definition}{2.3.5}\btanchor{2.3.5}
The map $\rho$ of \emph{(\btref{2.3.4})} is called the \textbf{retraction of $\Imm$
onto $A$ with center $C$}, and is written $\rho_{A;C}$.
\end{BT}

\begin{BT}{Proposition}{2.3.6}\btanchor{2.3.6}
Let $C$ be a chamber, $F$ a facet and $\Gamma$ a taut gallery between $C$ and
$F$.  Every apartment containing $C$ and $F$ contains each term of $\Gamma$.
\end{BT}
\proofref{p.~38}

\begin{BT}{Corollary}{2.3.7}\btanchor{2.3.7}
\begin{enumerate}[label=(\roman*),leftmargin=2.4em,itemsep=.25em]
\item Let $\Gamma$ and $\Gamma'$ be two minimal galleries of the same type.
      There exists $g\in G$ with $\Gamma=g\cdot\Gamma'$.
\item Let $C,C',C''$ be three pairwise distinct chambers possessing a common
      panel.  There exists $g\in G$ with $g\cdot C=C$ and $g\cdot C'=C''$.
\end{enumerate}
\end{BT}

\begin{BT}{Lemma}{2.3.8}\btanchor{2.3.8}
Let $A$ be an apartment and $C$ a chamber of $A$.
\begin{enumerate}[label=(\roman*),leftmargin=2.4em,itemsep=.25em]
\item For every $w\in\bW$ there exists one and only one chamber $C'$ contained
      in $A$ with $w(C,C')=w$ \emph{(}notation of \emph{(\btref{2.1.7})}\emph{)}.
\item Two non-stuttering galleries of the same type and the same origin,
      contained in $A$, coincide.
\end{enumerate}
\end{BT}

\begin{BT}{Corollary}{2.3.9}\btanchor{2.3.9}
Let $C$ and $C'$ be two chambers of $\Imm$ and put $w=w(C',C)$
\emph{(\btref{2.1.7})}.  Let $E$ be the set of minimal galleries with origin $C'$ and
extremity $C$.  The map assigning to $\Gamma\in E$ its type
$\mathbf{s}(\Gamma)$ \emph{(\btref{2.1.3})} is a bijection of $E$ onto the set of
reduced decompositions of $w$.
\end{BT}
\proofref{p.~39}

\para{2.3.10}\btanchor{2.3.10}
Let $w\in\bW$.  Applying (\btref{2.3.9}) to $j(\bC)$ and $j(w\cdot\bC)$ shows that the
set of minimal galleries with origin $\bC$ and extremity $w(\bC)$ is in
bijection with the reduced decompositions of $w$.  If
$\Gamma=(C_0=\bC,\dots,C_m=w(\bC))$ is such a gallery, of type
$\mathbf{s}=(s_1,\dots,s_m)$, the reflection in the wall $L_i$ supporting the
panel common to $C_{i-1}$ and $C_i$ is
\[
  t_i=(s_1\cdots s_{i-1})\,s_i\,(s_1\cdots s_{i-1})^{-1},
\]
and $w=s_1\cdots s_m=t_m\cdots t_1$.  It was recalled in (\btref{1.2.2}) that the
$t_i$ are pairwise distinct and that their set --- written $T_w$ in what
follows --- is independent of the reduced decomposition chosen.  The $t_i$ are
called the \textbf{reflections associated with $w$}.

\smallskip
\textit{In order that the reflection $r$ in a wall $L$ of $\bA$ belong to
$T_w$, it is necessary and sufficient that $\bC$ and $w(\bC)$ lie on opposite
sides of $L$.}  In particular, \textit{the length $\ell(w)$ equals the number
of walls of $\bA$ separating $\bC$ and $w\cdot\bC$}, and a minimal gallery of
$\bA$ is entirely contained in every affine root containing its extremities.
More generally:

\begin{BT}{Proposition}{2.3.11}\btanchor{2.3.11}
Every minimal gallery of $\Imm$ is contained in every half-apartment of
$\Imm$ containing its extremities.
\end{BT}

\begin{BT}{Proposition}{2.3.12}\btanchor{2.3.12}
Let $A$ be an apartment of $\Imm$, $C$ a chamber of $A$, and
$\rho=\rho_{A;C}$.  Let $F$ be a panel of $\Imm$; let $C_1$ be the chamber of
$A$ admitting $\rho(F)$ as a panel and lying on the same side as $C$ of the
wall of $A$ containing $\rho(F)$, and let $C_2$ be the other chamber of $A$
admitting $\rho(F)$ as a panel.  Then there exists one and only one chamber
$C'$ of $\Imm$ admitting $F$ as a panel with $\rho(C')=C_1$; and for every
chamber $C''$ of $\Imm$ admitting $F$ as a panel and distinct from $C'$, one
has $\rho(C'')=C_2$.
\end{BT}
\proofref{pp.~39--40}

\btsecB{2.4}{2.4. Closed subsets}

\begin{BT}{Definition}{2.4.1}\btanchor{2.4.1}
Let $M$ be a subset of $\Imm$ \emph{(}resp.\ $\bA$\emph{)} meeting at least
one chamber.  $M$ is said to be \textbf{closed} if, for every gallery
$(C_0,\dots,C_m)$ taut between a chamber $C_0$ meeting $M$ and a point of $M$,
one has $\overline{C}_k\subset M$ for $0\leq k\leq m$.
\end{BT}

\noindent
Every closed subset meeting a chamber is the union of the sets $\overline{C}$
for the chambers $C$ meeting $M$.  An apartment, a half-apartment and an
affine root are closed ((\btref{2.3.6}) and (\btref{2.3.11})).  The intersection of a family
of closed subsets containing a common chamber is again closed.  Hence:

\para{2.4.2}\btanchor{2.4.2}
If $M$ is a subset of $\Imm$ (resp.\ $\bA$) meeting at least one chamber,
there is a smallest closed subset of $\Imm$ (resp.\ $\bA$) containing $M$.  It
is called the \textbf{enclosure} of $M$ and written $\cl(M)$.

\para{2.4.3}\btanchor{2.4.3}
The notions of closed subset and of enclosure depend only on the
polysimplicial structure of $\Imm$ (resp.\ $\bA$).  If $M$ meets at least one
chamber, then
\[
\begin{aligned}
  j(\cl(M))&=\cl(j(M)) &&\text{for } M\subset\bA,\\
  g\cdot\cl(M)&=\cl(g\cdot M) &&\text{for } M\subset\Imm,\ g\in G.
\end{aligned}
\]

\begin{BT}{Lemma}{2.4.4}\btanchor{2.4.4}
Let $C$ and $C'$ be two chambers of $\bA$.  The enclosure of $C\cup C'$
coincides with the intersection of the affine roots of $\bA$ containing
$C\cup C'$, and with the union of the closures of the chambers belonging to at
least one minimal gallery with extremities $C$ and $C'$.
\end{BT}
\proofref{pp.~40--41; the argument compares
$\Card\Sep(C,C')$ with $\Card\Sep(C,C'')+\Card\Sep(C',C'')$, where
$\Sep(X,Y)$ denotes the set of walls of $\bA$ separating $X$ and $Y$}

\begin{BT}{Proposition}{2.4.5}\btanchor{2.4.5}
Let $M$ be a subset of $\bA$ meeting at least one chamber.  The enclosure of
$M$ coincides with the intersection of the affine roots containing $M$, and is
the convex hull of the union of finitely many points and half-lines.
\end{BT}
\proofref{p.~41}

\para{2.4.6}\btanchor{2.4.6}
Proposition (\btref{2.4.5}) allows definitions (\btref{2.4.1}) and (\btref{2.4.2}) to be extended to
arbitrary subsets of $\bA$:

\begin{BTu}{Definition}
For every subset $M$ of $\bA$, the \textbf{enclosure} of $M$, written
$\cl(M)$, is the intersection of the affine roots of $\bA$ containing $M$.
$M$ is said to be \textbf{closed} if $\cl(M)=M$, equivalently if $M$ is an
intersection of affine roots.
\end{BTu}

\noindent
\textit{A subset $M$ of $\bA$ is closed if and only if it is convex and a
union of closures of facets.}

\begin{BT}{Proposition}{2.4.7}\btanchor{2.4.7}
Let $M$ be a subset of $\Imm$ contained in the union of finitely many facets.
In order that $M$ be contained in an apartment it is necessary and sufficient
that there exist two chambers $C$ and $C'$ with $M$ contained in the enclosure
of $C\cup C'$.
\end{BT}

\begin{BT}{Proposition}{2.4.8}\btanchor{2.4.8}
Let $\mathrm{D}$ be a vector chamber.
\begin{enumerate}[label=(\roman*),leftmargin=2.4em,itemsep=.25em]
\item For every bounded subset $M$ of $\bA$ there exist $x,x'\in\bA$ with
      $M\subset(x+\mathrm{D})\cap(x'-\mathrm{D})$.
\item For all $x,x',y,y'\in\bA$ with $y\in x-\overline{\mathrm{D}}$ and
      $y'\in x'+\overline{\mathrm{D}}$, one has
      \[
        (x+\overline{\mathrm{D}})\cap(x'-\overline{\mathrm{D}})
        \subset\cl(\{y,y'\}).
      \]
\end{enumerate}
\end{BT}
\proofref{pp.~41--42}

\begin{BT}{Corollary}{2.4.9}\btanchor{2.4.9}
Let $\mathrm{D}$ be a vector chamber and $x$ a point of $\bA$.  For any points
$y\in x+\mathrm{D}$ and $y'\in x-\mathrm{D}$, the enclosure of $\{y,y'\}$ is a
neighborhood of $x$.
\end{BT}

\begin{BT}{Proposition}{2.4.10}\btanchor{2.4.10}
Let $F$ be a facet of $\bA$, $L$ the affine subspace generated by $F$, and
$\Delta$ a connected component of the complement in $\bA$ of the union of the
walls containing $F$.  For every bounded subset $M$ of $\overline{\Delta}$
there exist a point $y\in L$ and a sector $\Sect\subset\Delta$ such that
$M\subset\cl(\{y,z\})$ for every point $z\in\Sect$.
\end{BT}
\proofref{p.~42}

\begin{BT}{Proposition}{2.4.11}\btanchor{2.4.11}
With the notation of \emph{(\btref{2.4.10})}: for every bounded subset $M$ of
$\overline{\Delta}$ and every chamber $C$ of $\Imm$ admitting $j(F)$ as a
facet, there exists an apartment of $\Imm$ containing $C$ and $j(M)$.
\end{BT}
\proofref{p.~42}

\begin{BT}{Corollary}{2.4.12}\btanchor{2.4.12}
Let $\alpha$ be a half-apartment, $M$ a bounded subset of $\alpha$, and $C$ a
chamber of $\Imm$ possessing a panel contained in the wall $\partial\alpha$ of
$\alpha$.  Then there exists an apartment containing $C$ and $M$.
\end{BT}

\begin{BT}{Proposition}{2.4.13}\btanchor{2.4.13}
Let $M$ be a subset of $\Imm$ meeting at least one chamber, and let $g\in G$
satisfy $g\cdot x=x$ for all $x\in M$.  Then $g\cdot x=x$ for all
$x\in\cl(M)$.  In other words, \textbf{the fixers of $M$ and of $\cl(M)$
coincide}.
\end{BT}
\proofref{p.~43; it suffices to show that the fixed-point set
$\{x\in\Imm\mid g\cdot x=x\}$ is closed}

\btsecB{2.5}{2.5. The metric of the building}

We saw in (\btref{2.2.8}) that each apartment $A$ of $\Imm$ carries a structure of
euclidean affine space, and in particular a distance $d_A$.  We show that
there exists on $\Imm$ one and only one distance whose restriction to each
apartment $A$ is $d_A$.  This distance is $G$-invariant and makes $\Imm$ a
complete, contractible metric space in which any two points are joined by a
unique geodesic.

\begin{BT}{Lemma}{2.5.1}\btanchor{2.5.1}
There exists one and only one function $d\colon\Imm\times\Imm\to\RR_{+}$ such
that for every apartment $A$ of $\Imm$ the restriction of $d$ to $A\times A$
is the function $d_A$.
\end{BT}
\proofref{p.~43; uniqueness from (\btref{2.3.1}), existence from (\btref{2.3.3})
and (\btref{2.2.8})}

\para{2.5.2}\btanchor{2.5.2}
The function $d$ is clearly invariant under $G$.

\begin{BT}{Proposition}{2.5.3}\btanchor{2.5.3}
Let $C$ be a chamber of $\Imm$ and $A$ an apartment containing $C$.  Let
$\rho=\rho_{A;C}$ be the retraction onto $A$ with center $C$ \emph{(\btref{2.3.5})},
and let $x,y\in\Imm$.  Then:
\begin{enumerate}[label=(\roman*),leftmargin=2.4em,itemsep=.25em]
\item $d(\rho(x),\rho(y))\leq d(x,y)$;
\item if $x\in\overline{C}$, then $d(\rho(x),\rho(y))=d(x,y)$;
\item if there is a facet $F$ of $\Imm$ with $x,y\in\overline{F}$, then
      $d(\rho(x),\rho(y))=d(x,y)$.
\end{enumerate}
\end{BT}
\proofref{p.~43}

\begin{BT}{Proposition}{2.5.4}\btanchor{2.5.4}
\begin{enumerate}[label=(\roman*),leftmargin=2.4em,itemsep=.25em]
\item The function $d$ is a distance on $\Imm$.
\item Let $x,y\in\Imm$ and
      $D=\{z\in\Imm\mid d(x,y)=d(x,z)+d(z,y)\}$.  Then $D$ is contained in
      every apartment $A$ containing $x$ and $y$, and coincides with the
      segment $[xy]$ of the affine space $A$.
\item If an isometry of $\Imm$ --- in particular an action of $G$ --- fixes
      the two points $x$ and $y$, it fixes every point of $D$.
\end{enumerate}
\end{BT}
\proofref{p.~44}

\begin{BT}{Definition}{2.5.5}\btanchor{2.5.5}
The \textbf{building} of the Tits system $(G,B,N,S)$ is the set $\Imm$
equipped with its structure of polysimplicial complex defined in \emph{2.1},
with the family of structural maps defined in \emph{2.2}, and with the
distance $d$.
\end{BT}

\begin{BT}{Definition}{2.5.6}\btanchor{2.5.6}
With the notation of \emph{(\btref{2.5.4})}, $D$ is called the \textbf{closed segment
with extremities $x$ and $y$}, written $[xy]$.  A subset $M$ of $\Imm$ is
\textbf{convex} if $x,y\in M$ implies $[xy]\subset M$.
\end{BT}

\noindent
The intersection of a family of convex subsets of $\Imm$ is convex, so one may
speak of the \textbf{convex hull} of a subset of $\Imm$.

\para{2.5.7}\btanchor{2.5.7}
For a subset of $\Imm$ contained in an apartment $A$ to be convex in $\Imm$ in
the sense of (\btref{2.5.6}), it is necessary and sufficient that it be convex for the
affine structure of $A$.  In particular every apartment is convex in $\Imm$.
The intersection $A\cap A'$ of two apartments is convex in each of them, and
the affine structures induced on $A\cap A'$ by that of $A$ and by that of $A'$
are identical; since $A\cap A'$ is a union of closures of facets, $A\cap A'$
is closed in $A$ in the sense of (\btref{2.4.6}).

\smallskip
\textit{If $M$ is a subset of $\Imm$ meeting at least one chamber, then $M$ is
closed if and only if it is convex and a union of closures of chambers.}

\smallskip
More generally, a subset $M$ of $\Imm$ is called \textbf{closed} if it is
convex and a union of closures of facets (cf.\ (\btref{2.4.6})).  For $M\subset\Imm$
one calls \emph{enclosure} \dots
\par\endgroup

% ---------------------------------------- bruhattitsen.tex
\begingroup\microtypesetup{expansion=false}
\providecommand{\Imm}{}
\renewcommand{\Imm}{\mathscr{I}}
\providecommand{\Sect}{}
\renewcommand{\Sect}{\mathfrak{C}}
\providecommand{\bA}{}
\renewcommand{\bA}{\mathbf{A}}
\providecommand{\bC}{}
\renewcommand{\bC}{\mathbf{C}}
\providecommand{\bW}{}
\renewcommand{\bW}{\mathbf{W}}
\providecommand{\bS}{}
\renewcommand{\bS}{\mathbf{S}}
\providecommand{\bT}{}
\renewcommand{\bT}{\mathbf{T}}
\providecommand{\bR}{}
\renewcommand{\bR}{\mathbf{R}}
\providecommand{\bQ}{}
\renewcommand{\bQ}{\mathbf{Q}}
\providecommand{\bN}{}
\renewcommand{\bN}{\mathbf{N}}
\providecommand{\bSig}{}
\renewcommand{\bSig}{\mathbf{\Sigma}}
\providecommand{\Par}{}
\renewcommand{\Par}{\mathscr{P}}
\providecommand{\Fac}{}
\renewcommand{\Fac}{\mathscr{F}}
\providecommand{\Bor}{}
\renewcommand{\Bor}{\mathscr{B}}
\providecommand{\Jset}{}
\renewcommand{\Jset}{\mathscr{J}}
\providecommand{\Chm}{}
\renewcommand{\Chm}{\mathscr{C}}
\providecommand{\Lch}{}
\renewcommand{\Lch}{\mathscr{L}}
\providecommand{\cl}{}
\renewcommand{\cl}{\operatorname{cl}}
\providecommand{\Stab}{}
\renewcommand{\Stab}{\operatorname{Stab}}
\providecommand{\Isom}{}
\renewcommand{\Isom}{\operatorname{Isom}}
\providecommand{\diam}{}
\renewcommand{\diam}{\operatorname{diam}}
\providecommand{\Aut}{}
\renewcommand{\Aut}{\operatorname{Aut}}
\providecommand{\Ker}{}
\renewcommand{\Ker}{\operatorname{Ker}}
\providecommand{\Card}{}
\renewcommand{\Card}{\operatorname{Card}}
\providecommand{\Ghat}{}
\renewcommand{\Ghat}{\widehat{G}}
\providecommand{\Bhat}{}
\renewcommand{\Bhat}{\widehat{B}}
\providecommand{\stmt}{}
\renewcommand{\stmt}[2]{\medskip\noindent\textit{#1}\ \textbf{(#2)}.\ ---\ }
\providecommand{\num}{}
\renewcommand{\num}[1]{\medskip\noindent\textbf{(#1)}\ }
\bigskip

%==================================================================

%==================================================================

\pagemark{45}
\noindent
\dots of $M$, and we shall denote by $\cl(M)$ the intersection of the closed
subsets of $\Imm$ containing $M$; it is the smallest closed subset containing
$M$.

\stmt{Proposition}{2.5.8}\btanchor{2.5.8}
\emph{Let $A$ and $A'$ be two apartments. There exists $g\in G$ such that
$g.A=A'$ and $g.x=x$ for every $x\in A\cap A'$.}

\medskip
Let $F$ be a facet open in $A\cap A'$ and let $g\in G$ be such that $g.A=A'$
and $g.F=F$ (\btref{2.2.6}). Let $y\in F$ and let $x\in A\cap A'$; let $D$ be the
segment $[xy]$, which is contained in $A\cap A'$. The image of $D$ under $g$
is a segment of $A'$ with endpoint $y$ and of length $d(x,y)$. Since $F$ is
open in $A\cap A'$, $D\cap F$ is a subsegment not reduced to $\{y\}$ and fixed
pointwise by $g$. One deduces that $g.D=D$ and that $g$ fixes all the points
of $D$, and in particular $x$.

\stmt{Corollary}{2.5.9}\btanchor{2.5.9}
\emph{Let $g\in G$; set $A=j(\bA)$. There exists $n\in N$ such that $g.x=n.x$
for every $x\in A\cap g^{-1}.A$. If moreover $g.x=x$ for $x$ belonging to a
nonempty open subset of $A$, then $g.x=x$ for every $x\in A\cap g^{-1}.A$.}

\medskip
Indeed, there exists $h\in G$ such that $hg.A=A$ and $h.y=y$ for
$y\in A\cap g.A$. One then has $hg=n\in N$ and $g.x=n.x$ for
$x\in A\cap g^{-1}.A$. The second assertion follows, since an element of $N$
fixing the points of a nonempty open subset of $A$ belongs to $H$.

\stmt{Proposition}{2.5.10}\btanchor{2.5.10}
\emph{If $F$ is a facet of $\Imm$, the closure of $F$ in the metric space
$\Imm$ is none other than the union $\overline{F}$ of the facets of $F$.}

\medskip
This follows at once from the compactness of $\overline{F}$ in any apartment
containing $F$.

\stmt{Lemma}{2.5.11}\btanchor{2.5.11}
\emph{Let $F$ be a facet of $\Imm$ and $x\in F$. There exists an open ball $D$
with center $x$ such that, if $F'$ is a facet of $\Imm$ meeting $D$, then
$F\subset\overline{F'}$. If $D$ is such a ball and if $g\in G$ is such that
$g.x\in D$, then $g.x=x$.}

\medskip
Let $x_0$ and $F_0$ be the images of $x$ and $F$ under the canonical map of
$\overline{F}$ into $\overline{\bC}$ (\btref{2.1.2}). There exists $\varepsilon>0$
such that every wall of $\bW$ meeting the open ball $D_0$ with center $x_0$
and radius $\varepsilon$ contains $x_0$. Let $F_1$ be a facet of $\bA$ meeting
$D_0$ and let $y\in F_1\cap D_0$. Every wall meeting the open segment
$\left]yx_0\right[$ contains $x_0$, hence also $[yx_0]$. This entails that
$\left]yx_0\right[\,\subset F_1$ and $x_0\in\overline{F_1}$. We therefore have
$F_0\subset\overline{F_1}$.

Now let $D$ be the open ball of $\Imm$ with center $x$ and radius
$\varepsilon$, and let $F'$ be a facet of $\Imm$ meeting $D$. Let $\alpha$ be
a structural map such that $\alpha(\overline{\bC})\supset F$ and
$\alpha(\bA)\supset F'$. Since $\alpha$ is an isometry, and since
$\alpha^{-1}|F$ is the canonical map of $F$ into $\overline{\bC}$, one has
$\alpha^{-1}(x)=x_0$, $\alpha^{-1}(F)=F_0$ and $\alpha^{-1}(D)=D_0$. Set
$F'_0=\alpha^{-1}(F')$. Since $F'_0\cap D_0\neq\emptyset$, one has
$F_0\subset\overline{F'_0}$ by the first part of the proof, whence
$F\subset\overline{F'}$.

Finally, let $g\in G$ be such that $g.x\in D$ and let $F'$ be the facet
containing $g.x$. One then has $x,\,g.x\in\overline{F'}$, whence $g.x=x$ by
(\btref{2.1.6}).

\stmt{Theorem}{2.5.12}\btanchor{2.5.12}
\emph{(i) The metric space $\Imm$ is complete.\\
(ii) Let $\Fac$ be a set of facets. Then $\bigcup_{F\in\Fac}\overline{F}$ is
closed in $\Imm$.}

\medskip
It suffices to show that, with the notation of (ii),
$M=\bigcup_{F\in\Fac}\overline{F}$ is \emph{complete}, that is to say that a
Cauchy sequence $(y_q)$ of $M$ possesses an accumulation point in $M$. Let $C$
\pagemark{46}
be a chamber of $\Imm$ and let $g_q\in G$ be such that
$x_q=g_q^{-1}.y_q\in\overline{C}$ (for $q\in\bN$). Up to replacing the
sequence $(y_q)$ by a subsequence, one may suppose that the sequence $(x_q)$
converges to $x\in\overline{C}$. Let $D$ be a ball with center $x$ possessing
the property of Lemma (\btref{2.5.11}). Since
\begin{align*}
d(g_p.x,g_q.x)&\le d(g_p.x,g_p.x_p)+d(g_p.x_p,g_q.x_q)+d(g_q.x_q,g_q.x)\\
&\le d(x,x_p)+d(y_p,y_q)+d(x_q,x)
\end{align*}
there exists an integer $m$ such that $g_q^{-1}g_p.x\in D$ for $p,q\ge m$, and
one has $g_p.x=g_q.x$ for $p,q\ge m$. For $p\ge m$ one has
$d(y_p,g_m.x)=d(x_p,x)$, and $y_p$ tends to $y=g_m.x$. Moreover, Lemma
(\btref{2.5.11}) shows that $y$ is adherent to the facet containing $y_p$ for $p$
large enough, and hence belongs to $M$.

\num{2.5.13}\btanchor{2.5.13}
\emph{Whatever the points $x,y\in\Imm$, the segment $[xy]$ is the unique
geodesic joining $x$ to $y$.}

\medskip
This follows at once from (\btref{2.5.4})(ii).

For every real number $t$ of the interval $[0,1]$ we shall denote by
$tx+(1-t)y$ the unique point $z\in[xy]$ satisfying $d(x,z)=(1-t)d(x,y)$.

\stmt{Lemma}{2.5.14}\btanchor{2.5.14}
\emph{Let $x,y,z,z'$ be four points of $\Imm$ such that $z\in[xy]$, and let
$\varepsilon\in\bR_+$ be such that
$d(x,z')\le d(x,z)+\varepsilon\, d(x,y)$ and
$d(y,z')\le d(y,z)+\varepsilon\, d(x,y)$. One then has
$d(z,z')^2\le 4\varepsilon\, d(x,z)\,d(y,z)+\varepsilon^2 d(x,y)^2$.}

\medskip
Let $A$ be an apartment containing $[xy]$ and $C$ a chamber of $A$ such that
$z\in\overline{C}$. Set $z''=\rho_{A;C}(z)$. One has $d(z,z')=d(z,z'')$ and
the four points $x,y,z,z''$ satisfy the same hypotheses as $x,y,z,z'$. One
sees that it suffices to prove the lemma when the four points under
consideration belong to one and the same apartment, and it then results from
an elementary computation of plane Euclidean geometry.

\stmt{Lemma}{2.5.15}\btanchor{2.5.15}
\emph{The map $(t,x,y)\mapsto tx+(1-t)y$ from $[0,1]\times\Imm\times\Imm$ into
$\Imm$ is continuous.}

\medskip
This follows easily from Lemma (\btref{2.5.14}).

\stmt{Proposition}{2.5.16}\btanchor{2.5.16}
\emph{The metric space $\Imm$ is contractible.}

\medskip
This follows from Lemma (\btref{2.5.15}).

\num{2.5.17}\btanchor{2.5.17}
Let $x,y\in\Imm$ and let $\Lch(x,y)$ be the set of finite sequences
$s=(x_0,\dots,x_m)$ of points of $\Imm$ such that $x_0=x$, $x_m=y$, and such
that $x_{j-1}$ and $x_j$ belong to the closure of one and the same chamber for
$1\le j\le m$. One has
\[
d(x,y)=\inf_{s\in\Lch(x,y)}\ \sum_{j=1}^{m} d(x_{j-1},x_j).
\]
It follows at once that \emph{the metric of $\Imm$ is completely determined by
the data of the polysimplicial complex structure of $\Imm$ and of the metric
space structures on the subsets of $\Imm$ of the form $\overline{C}$, where
$C$ is a chamber of $\Imm$.}

%==================================================================

\btsecB{2.6}{2.6. Irreducible components}

%==================================================================
\pagemark{47}

\stmt{Lemma}{2.6.1}\btanchor{2.6.1}
\emph{Let $(\bS',\bS'')$ be a partition of $\bS$ such that every element of
$\bS'$ commutes with every element of $\bS''$. Then $\bW_{\bS'}$ is a normal
subgroup of $\bW$ and $(G,B_{\bS'},N)$ is a Tits system whose Weyl group is
canonically isomorphic to $\bW/\bW_{\bS'}$ and to $\bW_{\bS''}$.}

\medskip
Since $\bW_{\bS'}\cap\bW_{\bS''}=\bW_{\bS'\cap\bS''}=\{e\}$, the group $\bW$
is the direct product of $\bW_{\bS'}$ and $\bW_{\bS''}$. Set
$H'=B_{\bS'}\cap N$: the image of $H'$ in $\bW$ is $\bW_{\bS'}$, and $N/H'$ is
identified with $\bW/\bW_{\bS'}$, or again with $\bW_{\bS''}$. Let
$s\in\bS''$ and $w\in\bW_{\bS''}$. One has:
\begin{align*}
sB_{\bS'}w &= sB\bW_{\bS'}Bw \subset sB\bW_{\bS'}wB = sBw\bW_{\bS'}B
\subset (Bsw\bW_{\bS'}B)\cup(Bw\bW_{\bS'}B)\\
&\subset (B_{\bS'}wB_{\bS'})\cup(B_{\bS'}swB_{\bS'}).
\end{align*}
Finally $sB_{\bS'}s\neq B_{\bS'}$, since $B_{\bS'}$ is its own normalizer. We
have thus verified the axioms of a Tits system.

Let us remark that this lemma is valid for an arbitrary Tits system, not
necessarily of affine type.

\num{2.6.2}\btanchor{2.6.2}
Let $\bA=\bA_1\times\dots\times\bA_m$ be the canonical decomposition of $\bA$
into a product of its irreducible components under $\bW$, and let
$\bW=\bW_1\times\dots\times\bW_m$, $\bC=\bC_1\times\dots\times\bC_m$ and
$\bS=\bS_1\cup\dots\cup\bS_m$ be the corresponding decompositions of $\bW$,
$\bC$ and $\bS$ ((\btref{1.3.2}) and (\btref{1.3.4})).

For every subset $J$ of $\{1,\dots,m\}$ let $p_J$ be the canonical projection
of $\bA$ onto $\bA_J=\prod_{i\in J}\bA_i$, let $\bC_J=\prod_{i\in J}\bC_i$,
and let $\bW_J=\prod_{i\in J}\bW_i$, regarded as a group of affine
transformations of $\bA_J$. Set $\bS_J=\bigcup_{i\in J}\bS_i$,
$B_J=B_{\bS_J}=B\bW_JB$ and $B^J=B_{\{1,\dots,m\}-J}$. It follows from Lemma
(\btref{2.6.1}) that \emph{$(G,B^J,N)$ is a Tits system with Weyl group $\bW_J$},
hence of affine type. Let $\Imm_J$ be the building of $(G,B^J,N)$ (constructed
from $\bA_J$ and $\bC_J$), $d_J$ its distance, $j_{0,J}$ the canonical map of
$\overline{\bC_J}$ onto $\overline{F(B^J)}$, etc. The conjugacy class of
$(B^J,N)$, hence also the building $\Imm_J$, depends only on the conjugacy
class of $(B,N)$.

\stmt{Proposition}{2.6.3}\btanchor{2.6.3}
\emph{(i) Let $\Par$ (resp. $\Par_J$) be the set of parahoric subgroups of
$(G,B,N)$ (resp. of $(G,B^J,N)$). For every $P\in\Par$ there exists a smallest
element, denoted $\pi_J(P)$, in the set of the elements of $\Par_J$ which
contain $P$.\\
(ii) There exists one and only one map $\varpi_J$ from $\Imm$ into $\Imm_J$
such that
\[
\varpi_J(F(P))=F(\pi_J(P))\quad\text{for every }P\in\Par;\qquad
p_J\circ\lambda=\lambda_J\circ\varpi_J
\]
where $\lambda$ (resp. $\lambda_J$) denotes the canonical map of $\Imm$
(resp. $\Imm_J$) into $\overline{\bC}$ (resp. $\overline{\bC_J}$).\\
(iii) For every $g\in G$ and $a\in\Imm$ one has
$g.\varpi_J(a)=\varpi_J(g.a)$, and $\varpi_J$ is surjective.\\
(iv) For every structural map $\varphi:\bA\to\Imm$ there exists one and only
one structural map $\varphi_J:\bA_J\to\Imm_J$ such that
$\varpi_J\circ\varphi=\varphi_J\circ p_J$.}

\medskip
To prove (i) one may suppose $P=B_X$ with $X\in\bT$. The subgroups of $G$
containing $P$ are then the $B_Y$ with $X\subset Y\subset\bS$, and if such a
subgroup belongs to $\Par_J$,
\pagemark{48}
it must contain a conjugate of $B^J$ and $Y$ must contain $\bS-\bS_J$. One
thus sees that $\pi_J(P)=B_{X\cup(\bS-\bS_J)}$ possesses the desired property.

Let us prove (ii). If $\varpi_J$ exists, it sends the point $(P,a)\in\Imm$
(with $a\in\lambda(F(P))$) to the point $(\pi_J(P),p_J(a))\in\Imm_J$. It
therefore suffices to show that
\begin{equation}
p_J(\lambda(F(P)))\subset \lambda_J(F(\pi_J(P))).
\end{equation}
By an operation of $G$ one reduces to the case $P=B_X$. One then has
$\lambda(F(P))=\bC_X$ and
$\lambda_J(F(\pi_J(P)))=\lambda_J((B^J)_{X\cap\bS_J})=(\bC_J)_{X\cap\bS_J}$
(the facet of type $X\cap\bS_J$ of $\bC_J$), whence (1).

The first assertion of (iii) follows from the uniqueness of $\varpi_J$. It
entails
$\varpi_J(\Imm)=G.\varpi_J(\overline{F(B)})=G.\overline{F(B^J)}=\Imm_J$.

Let us prove (iv). Taking (iii) into account, (iv) will follow from the
relation
\begin{equation}
\varpi_J(j(x))=j_J(p_J(x))\qquad\text{for }x\in\bA
\end{equation}
(where $j:\bA\to\Imm$ and $j_J:\bA_J\to\Imm_J$ are defined as in (\btref{2.2.1})). It
suffices, again by virtue of (iii) and of (\btref{2.2.1}), to establish (2) for
$x\in\bC_X$. Now in this case one has:
\[
\varpi_J(j(x))=\varpi_J(B_X,x)=\bigl((B^J)_{X\cap\bS_J},\,p_J(x)\bigr)=j_J(p_J(x)).
\]

\stmt{Proposition}{2.6.4}\btanchor{2.6.4}
\emph{Let $\{J\mid J\in\Jset\}$ be a partition of the set $\{1,\dots,m\}$. Let
$\Imm'$ be the product of the buildings $\Imm_J$ for $J\in\Jset$, endowed with
the product polysimplicial complex structure and with the distance
$\bigl(\sum_J d_J^2\bigr)^{1/2}$; let $\Par'$ be the product of the sets
$\Par_J$ endowed with the product order structure, let
$\varpi=(\varpi_J):\Imm\to\Imm'$ be the product map of the $\varpi_J$ and
$\pi=(\pi_J):\Par\to\Par'$ the product map of the $\pi_J$.\\
(i) $\varpi$ is an isometry and an isomorphism of polysimplicial complexes
from $\Imm$ onto $\Imm'$.\\
(ii) For every apartment $A$ of $\Imm$, the restriction of $\varpi$ to $A$ is
an isomorphism of Euclidean affine spaces and of polysimplicial complexes from
$A$ onto the product of the apartments $\varpi_J(A)$.\\
(iii) $\pi$ is an isomorphism of ordered sets from $\Par$ onto $\Par'$. For
every $P\in\Par$ one has $P=\bigcap_{J\in\Jset}\pi_J(P)$.}

\medskip
Assertion (ii) follows from (\btref{2.6.2}) and from (\btref{2.6.3})(iii). It entails that
$\varpi$ is isometric and is a morphism of polysimplicial complexes. In order
to establish (i) it therefore only remains to show that $\varpi$ is
surjective. By induction on $\Card\Jset$ one reduces to the case where
$\Jset=(J_1,J_2)$. Let then $x_k\in\Imm_{J_k}$ and $y_k\in\Imm$ with image
$x_k$ (for $k=1,2$) ((\btref{2.6.3})(iii)). Let $A$ be an apartment of $\Imm$
containing $y_1$ and $y_2$. One then has $x_k\in\varpi_{J_k}(A)$ for $k=1,2$,
whence $(x_1,x_2)\in\varpi(A)$ by (ii).

Finally, if $P\in\Par$, one has $\varpi(F(P))=\prod_J F(\pi_J(P))$. Since
$\varpi$ commutes with the action of $G$, one deduces that
$P=\bigcap_{J\in\Jset}\pi_J(P)$ (\btref{2.1.5}). It follows that $\pi$ is injective.
Let us show that $\pi$ is surjective: indeed, if $P_J\in\Par_J$ for
$J\in\Jset$, then $\prod_J F(P_J)$ is a facet $F'$ of $\Imm'$ and there exists
$P\in\Par$ such that $F'=\varpi(F(P))$, or again $P_J=\pi_J(P)$ for every $J$.
Finally, it is clear that $P\subset Q$ is equivalent to
$\pi_J(P)\subset\pi_J(Q)$ for every $J\in\Jset$ ($P,Q\in\Par$).

\num{2.6.5}\btanchor{2.6.5}
When $J$ is a subset with a single element $i$, we agree to write $i$ instead
of $\{i\}$ in the preceding notation. The Tits systems $(G,B^i,N)$ are
irreducible:
\pagemark{49}
they are called the \emph{irreducible components} of $(G,B,N)$. Proposition
(\btref{2.6.4}) thus furnishes a \emph{canonical decomposition of the building $\Imm$
into a product of simplicial complexes}.

\num{2.6.6}\btanchor{2.6.6}
Let $J\subset\{1,\dots,m\}$. It is clear that $(B_J,B,N\cap B_J)$ is also a
Tits system with Weyl group $\bW_J$. Its building is canonically isomorphic to
$\Imm_J$, the correspondence $P\mapsto P'$ between parahoric subgroups of
$(G,B^J,N)$ and of $(B_J,B,N\cap B_J)$ being describable as follows:
$P'=P\cap B_J$, and $P$ is the unique parahoric subgroup of $(G,B,N)$ of type
$\tau(P)\cup(\bS-\bS_J)$ containing $P'$.

%==================================================================

\btsecB{2.7}{2.7. \texorpdfstring{$B$}{B}-adapted homomorphisms}

%==================================================================

Let $\varphi:G\to\Ghat$ be a $B$-adapted homomorphism ((\btref{1.2.13}) sqq., whose
notation we take up again). We saw in 1.2 that $\Ghat$ operates on the set
$\Par$ of parahoric subgroups of $G$, and we defined a homomorphism $\xi$ of
$\Ghat$ into the group of automorphisms of the Coxeter system $(\bW,\bS)$, a
group which is identified with $\widetilde{\bW}^{\dagger}_{\bC}$, with the
notation of (\btref{1.3.18}), that is to say with the group of automorphisms of the
Euclidean space $\bA$ which preserve the chamber $\bC$.

\num{2.7.1}\btanchor{2.7.1}
Let then $y=(P,x)\in\Imm$ (with $P\in\Par$ and $x\in\bC_{\tau(P)}$). Since
${}^{g}P$ is a parahoric subgroup of $G$ of type $\xi(g)(\tau(P))$ (\btref{1.2.18}),
the pair $g.y=({}^{g}P,\xi(g).x)$ is again a point of $\Imm$, and one thus
defines a \emph{law of operation} of $\Ghat$ on the set $\Imm$. For this law
the stabilizer of a facet $F(P)$ is none other than the subgroup $\Stab P$
(\btref{1.2.15}), and $\Bhat=\Stab B\cap\xi^{-1}(1)$ is the fixer of the chamber
$F(B)$.

\stmt{Proposition}{2.7.2}\btanchor{2.7.2}
\emph{(i) If $g\in G$ and $y\in\Imm$, one has $\varphi(g).y=g.y$. In other
words, the laws of operation of $G$ and of $\Ghat$ on $\Imm$ are compatible
with the homomorphism $\varphi$. One has
\[
\Ker\varphi\subset\bigcap_{g\in G} gHg^{-1}=\bigcap_{g\in G} gBg^{-1}.
\]
(ii) For every $g\in\Ghat$ the map $y\mapsto g.y$ is an isometric automorphism
of polysimplicial complex of $\Imm$.\\
(iii) If, moreover, the homomorphism $\varphi$ is $B$-$N$-adapted (\btref{1.2.13}),
this map permutes among themselves the apartments (resp. the half-apartments,
the sectors, the walls) of $\Imm$.}

\medskip
The first assertion of (i) is immediate, taking into account that for $g\in G$
and $P\in\Par$ one has ${}^{\varphi(g)}P=gPg^{-1}$ and $\xi(\varphi(g))=1$.
The second follows, since if $g\in\Ker\varphi$ then $\varphi(g)$, hence also
$g$, operates trivially on $\Imm$.

Let us prove (ii). Let $g\in\Ghat$ and let $\gamma$ be the map $y\mapsto g.y$
of $\Imm$ into itself. Since the action of $\Ghat$ on $\Par$ respects the
inclusion relations in $\Par$, one sees that, for $P\in\Par$, the image under
$\gamma$ of the facet $F(P)$ is the facet $F({}^{g}P)$ and that
$\gamma(\overline{F(P)})=\overline{F({}^{g}P)}$. Moreover, the restriction of
$\gamma$ to $\overline{F(P)}$ is obtained by composing the canonical map of
$\overline{F(P)}$ onto $\overline{\bC}_{\tau(P)}$, the map $\xi(g)$, and the
canonical map of $\overline{\bC}_{\xi(g)(\tau(P))}$ onto
$\overline{F({}^{g}P)}$. Since $\xi(g)$ is an automorphism of the Euclidean
space $\bA$, one deduces that $\gamma$ is an automorphism
\pagemark{50}
of polysimplicial complex and that the restriction of $\gamma$ to
$\overline{F(P)}$ is isometric. Consequently $\gamma$ itself is isometric by
(\btref{2.5.17}).

Let us prove (iii). Up to multiplying $g$ on the left by an element of
$\varphi(G)$, one may suppose that $g$ normalizes $\varphi(B)$ and
$\varphi(N)$. One then has, for every $w\in\bW$ (\btref{1.2.20})
\[
{}^{g}(wBw^{-1})=w'Bw'^{-1}\quad\text{with }w'=\xi(g)(w)=\xi(g)w\xi(g)^{-1}
\]
and more generally
\[
{}^{g}(wB_Xw^{-1})=w'B_{\xi(g)(X)}w'^{-1}\qquad\text{for }X\in\bT.
\]
Let then $a\in\bA$ and let $X\in\bT$ and $w\in\bW$ be such that
$a\in w.\bC_X$. One has $j(a)=(wB_Xw^{-1},w^{-1}(a))$ and
$g.j(a)=(w'B_{\xi(g)(X)}w'^{-1},\,\xi(g).w^{-1}.a)$, whence
\[
g.j(a)=(w'B_{\xi(g)(X)}w'^{-1},\,w'^{-1}.\xi(g).a)=j(\xi(g).a).
\]
Consequently $\gamma$ preserves the apartment $j(\bA)$ and reduces on
$j(\bA)$ to the map $j\circ\xi(g)\circ j^{-1}$, and hence permutes among
themselves the walls (resp. the sectors, the half-apartments) of $j(\bA)$. One
deduces (iii) by using (i) and the transitivity of $G$ on the apartments.

\num{2.7.3}\btanchor{2.7.3}
On the other hand, $\Ghat$ \emph{does not permute among themselves the
structural maps}, unless $\xi(\Ghat)=\{1\}$.

\num{2.7.4}\btanchor{2.7.4}
What precedes applies in particular if one takes for $\Ghat$ the group
$\Aut_B G$ (resp. $\Aut_{(B,N)}G$) of the automorphisms of $G$ which preserve
the conjugacy class of $B$ (resp. of the pair $(B,N)$), and for $\varphi$ the
homomorphism which associates to $g\in G$ the inner automorphism defined by
$g$. One then has $\gamma.g.x=\gamma(g).\gamma.x$ for $\gamma\in\Ghat$,
$g\in G$ and $x\in\Imm$. Proposition (\btref{2.7.2}) then shows that \emph{the
building $\Imm$ with all its structures depends} (up to an automorphism of
$\bA$ preserving $\bC$, as far as the structural maps are concerned)
\emph{only on the conjugacy class of the pair $(B,N)$}. If one regards $\Imm$
as endowed only with its metric and its polysimplicial structure, it depends
only on the conjugacy class of $B$.

%==================================================================

\btsecB{2.8}{2.8. Characterization of apartments}

%==================================================================

We saw (\btref{2.4.7}) that, in order for a subset $M$ of $\Imm$ to be contained in an
apartment and bounded therein, it is necessary and sufficient that there exist
two chambers $C$ and $C'$ such that $M$ is contained in the enclosure of
$C\cup C'$. In particular, the family of the subsets of $\Imm$ which are
\emph{bounded subsets of an apartment} is completely determined by the data of
the polysimplicial complex structure of $\Imm$.

The same is not true in general of the family of the apartments: there exist
saturated equivalent Tits systems of affine type $(G,B,N,\bS)$ and
$(G,B,N',\bS')$ for which $N$ and $N'$ are not conjugate. The associated
polysimplicial complex $\Imm$ is the same for the two systems, but the
families of apartments are distinct: otherwise there would exist an element
$g\in G$ such that $j'(\bA)=g.j(\bA)$, denoting by $j$ and $j'$ the canonical
maps
\pagemark{51}
of $\bA$ into $\Imm$ corresponding to these two Tits systems with Weyl group
$\bW$, and one would have $N'=gNg^{-1}$ by (\btref{2.2.5}).

In this number we shall first give another characterization of the bounded
subsets of apartments, then, in the case where the group $G$ is
\emph{complete} for a suitable topology, of the apartments themselves.
Finally, we shall show that the apartment $j(\bA)$ is entirely determined by
the polysimplicial structure of $\Imm$ and by the group $N$ regarded as a
group of permutations of $\Imm$.

\emph{Henceforth we shall identify $\bA$ and $j(\bA)$ by means of $j$, and
shall therefore regard $\bA$ as an apartment of $\Imm$.}

\stmt{Proposition}{2.8.1}\btanchor{2.8.1}
\emph{Let $M$ be a bounded subset of $\Imm$ possessing one of the two
following properties:\\
a) the interior of $M$ is not empty;\\
b) $M$ is convex.\\
In order that there exist an apartment containing $M$, it is necessary and
sufficient that there exist an isometry of $M$ onto a subset of $\bA$.}

\medskip
That the condition is necessary is evident. Let us suppose it satisfied, and
suppose first that the interior of $M$ is not empty. There then exists a
chamber $C_0$ such that $C_0\cap M$ is of nonempty interior. On the other
hand, every isometry $\varphi:X\to X'$ of one subset of $\bA$ onto another
extends to an automorphism of $\bA$, which is unique as soon as $X$ contains a
nonempty open subset of $\bA$. One deduces easily that, if $A$ is an apartment
containing $C_0$, there exists one and only one isometry $\tau$ of $M$ onto a
subset of $A$ whose restriction to $C_0\cap M$ is the identity. Moreover,
$\tau$ is the restriction to $M$ of the retraction $\rho$ of center $C_0$ of
$\Imm$ onto $A$, since one has $d(\rho(x),y)=d(x,y)=d(\tau(x),y)$ for every
$x\in M$ and for every $y$ belonging to $C_0\cap M$, which is of nonempty
interior.

Let $E$ be the enclosure of $\rho(M)$, and let $\Chm$ be the set of the
chambers $C\subset E$ such that the restriction of $\rho$ to
$M\cup\overline{C}$ is again isometric. The restriction of $\rho$ to the set
$M'$, the union of $M$ and of the closures $\overline{C}$ of the chambers
belonging to $\Chm$, is also isometric. We shall then argue by induction on
the number $k$ of chambers contained in $E$ and not belonging to $\Chm$; what
precedes shows that if $k=0$ the restriction of $\rho$ to $M\cup E$ is
isometric, which entails $M\subset A$.

So suppose $k>0$. Considering a minimal gallery between $C_0$ (which
obviously belongs to $\Chm$) and a chamber contained in $E$ and not belonging
to $\Chm$, one sees that there exist two adjacent chambers $C$ and $C'$
contained in $E$ such that $C\in\Chm$ and $C'\notin\Chm$. The first part of
the proof shows that one may suppose $C=C_0$.

Let then $F$ be the panel common to $C$ and $C'$ and let $a\in F$. Let us
distinguish two cases:

\smallskip\noindent
1) \emph{For every $x\in M'$, $x\neq a$, the intersection
$[ax]\cap(\overline{C}\cup\overline{C'})$ is a subsegment of $[ax]$ not
reduced to $\{a\}$.} Let then $x\in M'$, $x\neq a$, $z\in\overline{C'}$, and
let $A'$ be an apartment containing $C'$ and $x$. The triangle $T$, the union
of the segments $[az]$, $[zx]$ and $[ax]$, is contained in $A'$, hence is
``Euclidean''. The triangle $T'=\rho(T)$ is then isometric to it, since on the
one hand $\rho$ is
\pagemark{52}
isometric on $[ax]$ and on the other hand the angles at the vertex $a$ of $T$
and $T'$ are equal, $\rho$ being the identity on
$\overline{C}\cup\overline{C'}$. One therefore has $d(z,x)=d(z,\rho(x))$, and
the restriction of $\rho$ to $M\cup\overline{C'}$ is isometric, contrary to
the hypothesis according to which $C'\notin\Chm$.

\smallskip\noindent
2) \emph{There exists $x\in M'$, $x\neq a$, such that
$[ax]\cap(\overline{C}\cup\overline{C'})=\{a\}$.} There then exists a chamber
$C''$ admitting $F$ as panel, distinct from $C$ and from $C'$, such that
$[ax]\cap C''\neq\emptyset$. Let then $A'$ be an apartment containing $C$ and
$x$, hence $C''$, and let $\rho'$ (resp. $\rho''$) be the retraction of $\Imm$
onto $A'$ of center $C$ (resp. $C''$). Let $y\in M'$; since
$\rho=\rho\circ\rho'$, one has
\begin{align*}
d(x,y)&=d(\rho(x),\rho(y))=d(\rho(\rho'(x)),\rho(\rho'(y)))\\
&\le d(\rho'(x),\rho'(y))\le d(x,y).
\end{align*}
Since $x=\rho'(x)=\rho''(x)$, one deduces
$d(x,\rho'(y))=d(x,y)\ge d(x,\rho''(y))$. Let $R'$ (resp. $R''$) be the
half-apartment of $A'$ whose wall contains $F$ and which contains $C$ (resp.
$C''$); one has $x\in R''$. Let us show that
$[ay]\cap(\overline{C}\cup\overline{C''})\neq\{a\}$ as soon as $y\neq a$.
Otherwise there would exist a chamber $C'''$ admitting $F$ as panel, distinct
from $C$ and from $C''$, with $[ay]\cap C'''\neq\emptyset$. Since
$\rho'(C''')=C''$ (resp. $\rho''(C''')=C$) and since $\rho'$ (resp. $\rho''$)
restricted to $[ay]$ is an isometry, one would have
$\rho'(y)\in\mathring{R}''$ and $\rho''(y)\in\mathring{R}'$, with moreover
$\rho''(y)=r.\rho'(y)$, where $r$ is the reflection in $A'$ with respect to
the hyperplane which is the common wall of $R'$ and $R''$. But this would
entail $d(x,\rho''(y))>d(x,\rho'(y))$, contrary to what we saw above.

In particular, if $y\in M'\cap A$ with $y\neq a$, one has
$[ay]\cap C\neq\{a\}$ since $C''\cap A=\emptyset$ and $y$ belongs to the
half-apartment $R$ of $A$ containing $C$ and whose wall contains $F$. By
(\btref{2.4.12}) there therefore exists an element $g\in G$ such that $g.y=y$ for
every $y\in M'\cap A$ and $g.C''=C'$. One then has $g.C=C$, and this entails
$\rho(g.z)=\rho(z)$ for every $z\in\Imm$. The restriction of $\rho$ to
$M''=g.M'$ is therefore again an isometry, and the enclosure of
$\rho(M'')=\rho(M')$ is still equal to $E$. Moreover, for every $y\in M''$
with $y\neq a$, one has $[ay]\cap(\overline{C}\cup\overline{C'})\neq\{a\}$ and
the study made above in 1) shows that the restriction of $\rho$ to
$M''\cup\overline{C'}$ is isometric. We thus find ourselves in the same
situation, with $M''$ instead of $M$, but the family $\Chm$ is increased by at
least one chamber, namely $C'$, and the number $k$ decreased by at least one
unit.

The induction hypothesis therefore entails that there exists an apartment
$A''$ containing $M''$, and the apartment $g^{-1}.A''$ contains $M$, which
completes the proof in case a).

Let us now suppose that $M$ is convex. Let $F$ be a facet of $\Imm$ meeting
$M$, of the greatest possible dimension, and let $x\in F\cap M$. For every
$y\in M$, $y\neq x$, the intersection $F\cap[xy]$ is not reduced to $\{x\}$.
Let $A$ be an apartment containing $F$ and let $C$ be a chamber of $A$ whose
closure contains $F$. It suffices for us, in order to reduce to the case
studied previously, to show that the restriction to $C\cup M$ of the
retraction $\rho=\rho_{A;C}$ is isometric, or again that
$d(\rho(y),\rho(z))=d(y,z)$ for $y,z\in M$. Now, the two triangles which are
the convex hulls of $\{x,y,z\}$ and of $\{x,\rho(y),\rho(z)\}$ are both
isometric to Euclidean triangles; their angles at the vertex $x$ are equal
since $[xy]\cap[x\rho(y)]\neq\{x\}\neq[xz]\cap[x\rho(z)]$, and the sides
$[xy]$ and $[x\rho(y)]$ (resp. $[xz]$ and $[x\rho(z)]$) are
\pagemark{53}
of equal lengths. Consequently one has $d(y,z)=d(\rho(y),\rho(z))$ ([\btbib{22}],
book I, prop. IV), which terminates the proof.

\num{2.8.2}\btanchor{2.8.2}
Let us suppose that $G$ is endowed with a structure of \emph{complete}
topological group for which $B$ is a closed subgroup. Let us suppose moreover
that there exists a decreasing sequence of subgroups $(V_n)_{n\ge 1}$ forming
a fundamental system of neighborhoods of the neutral element, and an
increasing sequence of bounded subsets $\Omega_n$ of $\bA$ such that, denoting
by $P_n$ the fixer of $j(\Omega_n)$ in $G$, one has
\begin{equation}
P_n=(P_n\cap V_n).H\qquad\text{for every }n\ge 1.
\end{equation}
It is immediate that (1) remains valid if one replaces $\Omega_n$ by a larger
subset of $\bA$. One may therefore suppose moreover that
$\bC\subset\Omega_1$, that the $\Omega_n$ are of the form
$\cl(C_n\cup C'_n)$, where $C_n$ and $C'_n$ are chambers of $\bA$ (\btref{2.4.8}), and
have as union the whole of $\bA$. Let us remark that $P_n$, being then an
intersection of conjugates of $B$, is a closed subgroup of $G$, and hence
$P_n\cap V_n$ is also closed.

\stmt{Proposition}{2.8.3}\btanchor{2.8.3}
\emph{Under the hypotheses of (\btref{2.8.2}), every subset $M$ of $\Imm$ which is the
union of an increasing sequence of subsets $M_n$ $(n\ge 0)$, each contained in
an apartment $A_n$, is itself contained in an apartment.}

\medskip
One may obviously suppose that the $M_n$ are bounded and are convex unions of
closures of facets, and that $M_0$ is open in $M_n$. Let us choose $u_n\in G$
such that $A_{n+1}=u_n.A_n$ and $u_n.x=x$ for every
$x\in M_n\subset A_n\cap A_{n+1}$ (\btref{2.5.8}). Set $g_n=u_{n-1}\cdots u_0$ and
$M'_n=g_n^{-1}.M_n\subset A_0$. One thus obtains an increasing sequence of
bounded subsets of $A=A_0$, whose union $M'$ is a convex union of closures of
facets.

\emph{Let us show that one may reduce to the case where $M'=A$.} Suppose
$M'\neq A$ and let $\Gamma=(C_0,\dots)$ be an ``infinite gallery'' having as
terms all the chambers of $A$ (each of them being able to figure there several
times) and such that $C_0$ contains a facet of maximal dimension of $M'_0$.
Set $D_{-1}=\emptyset$ and define $D_i$ and the chamber $C'_i$ by induction on
$i\ge 0$ in the following way: $C'_i$ is the first term of $\Gamma$ not
contained in $\cl(M'\cup D_{i-1})$, and $D_i=D_{i-1}\cup C'_i$. Up to
replacing the sequence $(M_n)$ by an extracted subsequence, one may suppose
that $C'_i$ is also the first term of $\Gamma$ not contained in
$\cl(M'_i\cup D_{i-1})$. On the other hand, let us denote by $C''_i$ the term
preceding $C'_i$ in $\Gamma$; when $C'_0=C_0$, the chamber $C'_0$ has no
predecessor in $\Gamma$ and one then designates by $C''_0$ a chamber adjacent
to $C'_0$ such that the panel common to $C'_0$ and $C''_0$ contains a facet of
maximal dimension of $M'$. We shall show by induction on the integer $n$ that
there exist, for $k\ge n\ge -1$, elements $h_{n,k}\in G$ such that, setting
$h_n=h_{n,n}$:
\begin{itemize}\itemsep2pt
\item[(1)] $h_{n,k}$ coincides with $g_k$ on $M'_k$;
\item[(2)] $h_{n,k}$ coincides with $h_n$ on $\cl(M'_n\cup D_n)$ for $n\ge 0$;
\item[(3)] $h_n$ coincides with $h_{n-1}$ on $\cl(M'_{n-1}\cup D_{n-1})$ for
$n\ge 1$.
\end{itemize}
One takes $h_{-1}=\mathrm{id}$ and $h_{-1,k}=g_k$ for $k\ge 0$. Let $n\ge 0$
and suppose the $h_{i,k}$ constructed for $i<n$. Let $k\ge n$. By definition
of $C'_n$ and $C''_n$, the enclosure of $M'_k\cup D_{n-1}$ is
\pagemark{54}
contained in the half-apartment of $A_0$ containing $C''_n$ and whose wall
contains the panel $F$ common to $C'_n$ and $C''_n$. In view of the induction
hypothesis, the facet $F'=h_{n-1,k}.F$ does not depend on $k$ and one has
$h_{n-1,k}.(M'_k\cup D_{n-1})=M_k\cup h_{n-1}.D_{n-1}$. By (\btref{2.4.12}) there
exists an apartment $A'_k$ containing $h_{n-1,k}.\cl(M'_k\cup D_{n-1})$ and
$h_{n-1}.C'_n$. By (\btref{2.5.8}) there exists $h_{n,k}\in G$ such that
$h_{n,k}.A_0=A'_k$ and $h_{n,k}.x=h_{n-1,k}.x$ for every
$x\in\cl(M'_k\cup D_{n-1})$. It is then immediate that (1) and (3) are
satisfied. Condition (2) is so as well, for the restrictions of $h_{n,k}$ and
of $h_n$ to $\cl(M'_n\cup D_n)$ have the same image
$\cl(M_n\cup h_{n-1}.D_n)$ and coincide on the chamber $C_0$.

Replacing $M_n$ by $\cl(M_n\cup h_n.C'_n)$ and $g_n$ by $h_n$, one is indeed
reduced to the case where $M'=A$.

So let us suppose henceforth that $M'=A$. Since each $M'_n$ is a convex union
of closures of facets, one sees that $A$ is also the union of the interiors of
the $M'_n$ and, up to extracting a subsequence, one may suppose that
$\Omega_n\subset M'_n$. Since $g_{n+1}.x=g_n.x$ for $x\in M'_n$, one has
$g_n^{-1}g_{n+1}\in P_n$, and up to replacing $g_{n+1}$ by an element of the
form $g'_{n+1}=g_{n+1}h_{n+1}$ with $h_{n+1}\in H$ (or again $u_n$ by
$u'_n=g_nh_{n+1}g_n^{-1}$), one may suppose that
$g_n^{-1}g_{n+1}\in P_n\cap V_n$.

The sequence $g_n$ then converges to an element $g\in G$ and one has
$g.x=g_n.x$ for every $x\in M'_n$, whence $M_n\subset g.A$, which completes
the proof.

\stmt{Corollary}{2.8.4}\btanchor{2.8.4}
\emph{Let us suppose the hypotheses of (\btref{2.8.2}) satisfied and let $M$ be a
subset of $\Imm$ of nonempty interior.\\
(i) In order that $M$ be contained in an apartment, it is necessary and
sufficient that it be isometric to a subset of $\bA$.\\
(ii) In order that $M$ be an apartment, it is necessary and sufficient that it
be isometric to $\bA$.}

\medskip
This follows from (\btref{2.8.1}) and (\btref{2.8.3}).

\stmt{Proposition}{2.8.5}\btanchor{2.8.5}
\emph{Under the hypotheses of (\btref{2.8.2}), in order that a subset $M$ of $\Imm$ be
an apartment, it is necessary and sufficient that $M$ be a closed subset
containing at least one chamber of $\Imm$ and that, for every panel $F$ of
$\Imm$ contained in $M$, there exist exactly two chambers of $\Imm$ contained
in $M$ and admitting $F$ as panel.}

\medskip
That the condition is necessary is quite clear (and independent of (\btref{2.8.2})!).
Let us suppose it realized. For every chamber $C\subset M$ and every
$s\in\bS$ there then exists one and only one chamber $C'\subset M$, adjacent
to $C$ and such that the common panel is of type $s$. One concludes
immediately, using (\btref{2.3.9}) and the fact that $M$ is closed, that for every
$w\in\bW$ and every chamber $C\subset M$ there exists one and only one chamber
$C'\subset M$ such that $w(C,C')=w$ (\btref{2.1.7}).

Let $C_n$ and $C'_n$ be as in (\btref{2.8.2}). Let us choose a chamber
$C^{*}\subset M$, set $w_n=w(\bC,C_n)$ and $w'_n=w(\bC,C'_n)$, and let
$C^{*}_n$ (resp. $C'^{*}_n$) be the chamber of $M$ such that
$w(C^{*},C^{*}_n)=w_n$ (resp. $w(C^{*},C'^{*}_n)=w'_n$). Since
$\bC\subset\cl(C_n\cup C'_n)$, one has
\[
\ell(w_nw'^{-1}_n)=\ell(w_n)+\ell(w'_n)\qquad\text{and}\qquad
w(C'_n,C_n)=w_nw'^{-1}_n.
\]
\pagemark{55}
One deduces at once that $C^{*}\subset M_n=\cl(C^{*}_n\cup C'^{*}_n)$, that
$w(C'^{*}_n,C^{*}_n)=w_nw'^{-1}_n$, and hence that there exists $g_n\in G$
such that $g_n.C_n=C^{*}_n$, $g_n.C'_n=C'^{*}_n$. For every $k<n$ one has
$g_n.C_k\subset\cl(C^{*}_n\cup C'^{*}_n)$ and $w(C^{*},g_n.C_k)=w_k$, whence
$C^{*}_k=g_n.C_k$. Similarly $C'^{*}_k=g_n.C'_k$, which entails
$M_k\subset M_n$. In view of (\btref{2.8.3}) there therefore exists $g\in G$ such that
$g.\bC=C^{*}$ and $M_n\subset g.\bA$ for every $n$, whence $g.C_n=C^{*}_n$ and
$g.C'_n=C'^{*}_n$, and consequently $g.\bA=\bigcup_n M_n\subset M$. If
moreover one had $M\neq\bA$, there would be at least one panel of $M$
belonging to three distinct chambers contained in $M$, contrary to the
hypothesis. Consequently $M=g.\bA$ is indeed an apartment of $\Imm$.

\stmt{Corollary}{2.8.6}\btanchor{2.8.6}
\emph{Let $\varphi$ be an injection of $\bA$ into $\Imm$. Suppose that
$\varphi$ is a chambered morphism (\btref{1.1.7}) of polysimplicial complexes and
sends each panel of $\bC$ onto a panel of the same type. Then $\varphi(\bA)$
is an apartment and $\varphi$ is a structural map.}

\medskip
One shows immediately by induction on $\ell(w(\bC,C))$ that $\varphi$ sends a
panel of an arbitrary chamber $C$ of $\bA$ onto a panel of the same type. One
deduces that the image of a minimal gallery is a minimal gallery of the same
type, which entails that $\varphi(\bA)$ is closed, and hence is an apartment
by (\btref{2.8.5}). There therefore exists a structural map $\psi$ such that
$\psi(\bA)=\varphi(\bA)$ and $\psi(\bC)=\varphi(\bC)$. But
$\psi^{-1}\circ\varphi$ is then an automorphism of polysimplicial complex of
$\bA$ whose restriction to $\bC$ is an affine automorphism of $\bC$ preserving
each panel of $\bC$, which entails that $\psi^{-1}\circ\varphi$ is the
identity, whence the corollary.

\num{2.8.7}\btanchor{2.8.7}
One thus sees that, under the hypotheses of (\btref{2.8.2}), \emph{the family of the
apartments of $\Imm$ depends only on the polysimplicial structure of $\Imm$}.
One deduces in particular that \emph{every $B$-adapted homomorphism
$\varphi:G\to\Ghat$ is automatically $B$-$N$-adapted}.

Moreover, one shows easily that \emph{the metric of $\Imm$ determines its
polysimplicial structure}: the points of $\Imm$ interior to a chamber are
those which possess a neighborhood isometric to an open subset of $\bA$. It
then follows from (\btref{2.8.4}) that, under the hypotheses of (\btref{2.8.2}), every
isometry of $\Imm$ onto itself is an automorphism of the building $\Imm$, up
to an automorphism of $(\bA,\bW)$.

\stmt{Lemma}{2.8.8}\btanchor{2.8.8}
\emph{Let $A$ be an apartment, $C$ a chamber of $A$, and $y,y'$ two points of
$\Imm$. If $C\subset\cl(\{\rho_{A;C}(y),\rho_{A;C}(y')\})$, then
$C\subset\cl(\{y,y'\})$.}

\medskip
Set $\rho=\rho_{A;C}$, $z=\rho(y)$ and $z'=\rho(y')$. Let $(\Gamma,C)$ (resp.
$\Gamma'$) be a taut gallery between $y$ and $C$ (resp. between $C$ and
$y'$), $C'$ the origin of $\Gamma$ and $C''$ the extremity of $\Gamma'$. In
view of (\btref{2.4.4}) the gallery $(\rho(\Gamma),\rho(\Gamma'))$ is minimal. The
gallery $(\Gamma,\Gamma')$, which is of the same type, is therefore minimal as
well (\btref{2.1.11}), and $C$ is contained in $\cl(C'\cup C'')$. Consequently there
exists $g\in G$ such that $g.C=C$, $g.C'=\rho(C')$ and $g.C''=\rho(C'')$. One
then has
\[
\cl(\{y,y'\})=g^{-1}.\cl(\{g.y,g.y'\})=g^{-1}.\cl(\{z,z'\})\supset g^{-1}.C=C.
\]

\stmt{Proposition}{2.8.9}\btanchor{2.8.9}
\emph{Let $D$ be a vector chamber, $x,x'\in\bA$ and $y,y'\in\Imm$. Suppose
that there exist $y_1\in x+D$ and $y'_1\in x'-D$ such that $d(y,y_1)$ (resp.
$d(y',y'_1)$) is less than the distance from $y_1$ (resp. $y'_1$) to the
complement of $x+D$ (resp. of $x'-D$) in $\bA$. Then
$\cl(\{y,y'\})\supset (x'+D)\cap(x-D)$.}

\medskip
\pagemark{56}
Let $C$ be a chamber meeting $(x-D)\cap(x'+D)$ and let $\rho$ be the
retraction of $\Imm$ onto $\bA$ of center $C$. One has
$\rho(y)\in x+\overline{D}$ and $\rho(y')\in x'-\overline{D}$. In view of
(\btref{2.4.8}) and (\btref{2.8.8}) one therefore has $C\subset\cl(\{y,y'\})$, whence our
assertion.

\stmt{Corollary}{2.8.10}\btanchor{2.8.10}
\emph{Let $D$ be a vector chamber, $v$ an element of $D$ and $\lambda$ a real
number such that $0\le\lambda<1$. There exists a real number $\varepsilon>0$
such that, for $x,x'\in\bA$ and $y,y'\in\Imm$, the relations $x'-x\in\bR v$,
$d(x,y)<\varepsilon\, d(x,x')$ and $d(x',y')<\varepsilon\, d(x,x')$ entail the
existence of $z,z'\in[yy']\cap\bA$ with $d(z,z')\ge\lambda\, d(x,x')$.}

\medskip
Let $x_i$ $(i=1,\dots,6)$ be six distinct points of a line $L$ of $\bA$, such
that $x_i\in[x_{i-1}x_{i+1}]$ for $i=2,\dots,5$, that $x_6-x_1\in\bR_+v$ and
that $d(x_3,x_4)>\lambda\, d(x_1,x_6)$. Let $r\in\bR^{*}_+$ and let $X_i$ be
the ball of $\bA$ with center $x_i$ and radius $r$ $(1\le i\le 6)$. Suppose
$r$ small enough that $X_1\subset x_2-D$,
$X_3\cup X_4\subset(x_2+D)\cap(x_5-D)$ and $X_6\subset x_5+D$, and set
$\varepsilon=r\, d(x_1,x_6)^{-1}$.

Let $x,x'\in\bA$ and $y,y'\in\Imm$ satisfying the conditions of the statement.
Up to transforming the system of the $x_i$ by a homothety or a translation and
multiplying $r$, $v$ and $D$ by the ratio of the homothety, which does not
change $\varepsilon$, one may suppose that $x=x_1$ and $x'=x_6$. Let $A$ be an
apartment containing $\{y,y'\}$. One has $d(x,y)<\varepsilon\, d(x,x')=r$ and,
similarly, $d(x',y')<r$. From (\btref{2.8.9}) one then deduces that
$A\supset(x_2+D)\cap(x_5-D)$. Let $C$ be a chamber of $\bA$ meeting this
intersection and $\rho$ the restriction to $A$ of the retraction
$\rho_{\bA;C}$. One has $C\subset A\cap\bA$, and $\rho$ is an isometry of $A$
onto $\bA$ leaving fixed all the points of $A\cap\bA$ (\btref{2.3.4}). Since
$\rho_{\bA;C}$ decreases distances (\btref{2.5.3}), one also has $\rho(y)\in X_1$ and
$\rho(y')\in X_6$. One deduces at once that there exist
$z\in[\rho(y)\rho(y')]\cap X_3$ and $z'\in[\rho(y)\rho(y')]\cap X_4$ such that
$d(z,z')\ge\lambda\, d(x,x')$. Since
$X_3\cup X_4\subset(x_2+D)\cap(x_5-D)\subset A\cap\bA$, one has
$z=\rho^{-1}(z)$ and $z'=\rho^{-1}(z')$, whence $z,z'\in[yy']\cap\bA$.

\stmt{Proposition}{2.8.11}\btanchor{2.8.11}
\emph{The apartment $\bA$ is the smallest nonempty convex subset of $\Imm$
invariant under $\nu^{-1}(V)$. It is also the smallest nonempty closed subset
of $\Imm$ invariant under $\nu^{-1}(V)$.}

\medskip
The second assertion is an immediate consequence of the first.

Let $M$ be a nonempty convex subset of $\Imm$ invariant under $\nu^{-1}(V)$,
$y\in M$, $x\in\bA$, and $g\in\nu^{-1}(V)$ such that $\nu(g)$ belongs to the
kernel of no vector root. Set $v=\nu(g)$ and let $\lambda,\varepsilon$ be real
numbers possessing the properties of Corollary (\btref{2.8.10}). For $n\in\bN$
sufficiently large one has
$d(x,y)=d(g^n.x,g^n.y)<\varepsilon\, d(x,g^n.x)$, and hence there exists
$z\in[y(g^n.y)]\cap\bA=M\cap\bA$. Since the convex hull of $\nu^{-1}(V)(z)$ is
manifestly $\bA$, one has $\bA\subset M$.

\num{2.8.12}\btanchor{2.8.12}
We shall show by an example that assertions (\btref{2.8.3}) to (\btref{2.8.7}) do not remain
valid if one suppresses the hypotheses of (\btref{2.8.2}).

Suppose $\operatorname{card}\bS=2$, set $\bS=\{s_1,s_2\}$, and let $X$ be a
subgroup of $B$ and $n'_i\in Bs_iB$ $(i=1,2)$ such that $n'^2_i\in X$ and
$n'_iXn'^{-1}_i=X$; let $N'$ be the group generated by $n'_1$, $n'_2$ and $X$.
It is clear that $X$ is normal in $N'$ and that there exists a surjective
homomorphism $\eta:\bW\to N'/X$, and only one, such that $\eta(s_i)=n'_iX$.
From (\btref{1.2.8}) it follows that if $w\in\bW$ one has $\eta(w)\subset BwB$.
Consequently $\eta$ is bijective. One deduces at once that $(G,B,N')$ is a
\pagemark{57}
Tits system with Weyl group $\bW$, equivalent to $(G,B,N)$ (\btref{1.2.11}).
Consequently the building $\Imm'$ of $(G,B,N')$ is identified, as a metric
space and as a polysimplicial complex, with the building $\Imm$ of $(G,B,N)$
(\btref{2.7.4}), and the apartments of $\Imm'$ are those of $\Imm$ if and only if
there exists $b\in B$ with $N'\subset bNb^{-1}$ (\btref{2.2.6}), that is to say
$\eta(w)\subset b.w H.b^{-1}$ for every $w\in\bW$.

\stmt{Example}{2.8.13}\btanchor{2.8.13}
Let $p$ be a prime number, and let $\omega_p$ be the $p$-adic valuation of
$\bQ$. Take $G=\mathrm{SL}_2(\bQ)$,
\[
B=\left\{\begin{pmatrix}a&b\\c&d\end{pmatrix}\in G\ \Big|\
\omega_p(a)=\omega_p(d)=0,\ \omega_p(b)\ge 0,\ \omega_p(c)>0\right\},
\]
\[
N=\left\{\begin{pmatrix}a&0\\0&d\end{pmatrix}\in G\right\}\cup
\left\{\begin{pmatrix}0&b\\c&0\end{pmatrix}\in G\right\}.
\]
By [\btbib{28}], $(G,B,N)$ is a Tits system of affine type of rank 1 (see also
(\btref{6.2.3}), a)). One has
\[
H=\left\{\begin{pmatrix}a&0\\0&a^{-1}\end{pmatrix}\ \Big|\ a\in\bQ,\
\omega_p(a)=0\right\},\qquad
s_1=\begin{pmatrix}0&1\\-1&0\end{pmatrix}.H,\qquad s_2=\begin{pmatrix}0&p\\-p^{-1}&0\end{pmatrix}.H.
\]
Set
\[
n'_1=\begin{pmatrix}0&1\\-1&0\end{pmatrix},\qquad n'_2=\begin{pmatrix}-1&2p\\-p^{-1}&1\end{pmatrix},\qquad X=\left\{\pm\begin{pmatrix}1&0\\0&1\end{pmatrix}\right\}.
\]
One verifies easily that the hypotheses of (\btref{2.8.12}) are satisfied and that the
eigenvalues of $n'_1n'_2$ are not rational. It follows that $n'_1n'_2$ is not
conjugate to an element of $s_1s_2H$, and hence that $\Imm$ and $\Imm'$ do not
have the same apartments.

%==================================================================

\btsecB{2.9}{2.9. Sectors}

%==================================================================

\stmt{Proposition}{2.9.1}\btanchor{2.9.1}
\emph{Let $A$ be an apartment of $\Imm$ and $\Sect$ a sector of $A$. There
exists one and only one map $\rho_{A;\Sect}$ from $\Imm$ into $A$ possessing
the following property: for every bounded subset $M$ of $\Imm$ there exists a
sector $\Sect'\subset\Sect$ such that, for every chamber $C$ meeting
$\Sect'$, one has $\rho_{A;\Sect}(x)=\rho_{A;C}(x)$ for every $x\in M$.}

\medskip
We must show that, given a bounded subset $M$ of $\Imm$, there exists a sector
$\Sect'\subset\Sect$ such that, if $C$ denotes a chamber meeting $\Sect'$, the
restriction of $\rho_{A;C}$ to $M$ does not depend on the choice of $C$. One
may obviously suppose that $M$ is a union of closures of chambers. Let us take
a point $a$ of $A$ and set $d=\sup\{d(a,x)\mid x\in M\}$. Let $M'$ be the ball
with center $a$ and radius $d$ in $A$ and let $\Sect_1$ be a sector opposite
to $\Sect$ and
\pagemark{58}
containing $M'$ ((\btref{2.4.8})(i)). Let $b$ be the vertex of $\Sect_1$ and let
$\Sect''$ be the sector of $A$ opposite to $\Sect_1$ (hence parallel to
$\Sect$) with vertex $b$. We shall show that the sector
$\Sect'=\Sect\cap\Sect''$ answers the question.

Let $C_0$ be a chamber of $A$ to which $b$ is adherent; set
$\rho=\rho_{A;C_0}$. Since a retraction decreases distances, one has
$\rho(C')\subset M'\subset\Sect_1$ for every chamber $C'$ contained in $M$.

Let $C$ be a chamber meeting $\Sect''$. By (\btref{2.4.9}), the enclosure of
$C\cup\rho(C')$ contains $C_0$. There therefore exists a minimal gallery
$\Gamma$ of the form
\[
(C,C_1,\dots,C_m,C_0,C_{m+1},\dots,C_n,\rho(C'))
\]
and there exists a minimal gallery of the form $(C_0,C'_{m+1},\dots,C'_n,C')$
with $\rho(C'_j)=C_j$ for $m+1\le j\le n$. The gallery
$(C,C_1,\dots,C_m,C_0,C'_{m+1},\dots,C'_n,C')$ is non-stuttering and has
the same type as $\Gamma$. It is therefore minimal, and its image under the
retraction $\rho_{A;C}$ is $\Gamma$, the unique minimal gallery of type
$\mathbf{s}(\Gamma)$ and origin $C$. Consequently one has
\[
\rho_{A;C}(C')=\rho_{A;C_0}(C')
\]
for every chamber $C'$ contained in $M$ and every chamber $C$ meeting
$\Sect''$, whence the proposition.

\stmt{Definition}{2.9.2}\btanchor{2.9.2}
\emph{The map $\rho_{A;\Sect}$ defined by (\btref{2.9.1}) is called the retraction of
$\Imm$ onto the apartment $A$ relative to the sector $\Sect$.}

\medskip
It is clear that $\rho=\rho_{A;\Sect}$ does not change if one replaces
$\Sect$ by a sector of $A$ parallel to $\Sect$, and therefore depends only on
the ``sector germ'' of $A$ defined by $\Sect$ (cf. (7.2.3), (\btref{7.4.12})). It
is also clear that $\rho$ decreases distances and that its restriction to the
closure of a chamber is isometric (\btref{2.5.3}).

\stmt{Corollary}{2.9.3}\btanchor{2.9.3}
\emph{Let $A$ be an apartment of $\Imm$, $\Sect$ a sector of $A$ and $C$ a
chamber of $\Imm$. There exists a sector $\Sect'\subset\Sect$ such that the
restriction of $\rho_{A;\Sect}$ to $\Sect'\cup C$ is isometric.}

\medskip
In view of (\btref{2.5.3})(ii), it suffices to take a sector $\Sect'\subset\Sect$ such
that $\rho_{A;\Sect}(C)=\rho_{A;C'}(C)$ for every chamber $C'$ meeting
$\Sect'$.

\stmt{Corollary}{2.9.4}\btanchor{2.9.4}
\emph{Suppose $G$ endowed with a structure of complete topological group
satisfying the hypotheses of (\btref{2.8.2}). Let $\Sect$ be a sector and $C$ a
chamber of $\Imm$. There exist a sector $\Sect'\subset\Sect$ and an apartment
containing $C$ and $\Sect'$.}

\medskip
This follows from (\btref{2.9.3}) and (\btref{2.8.4}).

\stmt{Proposition}{2.9.5}\btanchor{2.9.5}
\emph{Let $A$ and $A'$ be two apartments, $\Sect$ (resp. $\Sect'$) a sector of
$A$ (resp. of $A'$). There exist sectors $\Sect_1\subset\Sect$ and
$\Sect'_1\subset\Sect'$ such that the restriction to
$\Sect_1\cup\Sect'_1$ of the retraction of $\Imm$ onto $A$ relative to
$\Sect$ is isometric.}

\medskip
One may suppose that $A=\bA$ (identified with $j(\bA)$). Let us denote by
$\rho$ the retraction of $\Imm$ onto $\bA$ relative to the sector $\Sect$. For
every chamber $C$ of $A'$, let $m(C)$ be the center of gravity of $C$ and let
$Q(C)$ be the intersection of $C$ with the sector of $A'$ with vertex $m(C)$
parallel
\pagemark{59}
to $\Sect'$. One verifies easily that there exists one and only one element
$w(C)$ of the Weyl group ${}^{v}\bW$ of the root system ${}^{v}\bSig$ such
that $\rho(Q(C))\subset\rho(m(C))+w.D$, where $D$ is the vector chamber which
is the direction of $\Sect$. Let us then choose a chamber $C_0$ of $A'$ in
such a way that the length of $w(C_0)$ in ${}^{v}\bW$ (relative to the base
$B(D)$ of ${}^{v}\bSig$ associated with $D$) is the largest possible. One may
suppose that $\Sect'$ has vertex $m(C_0)$, that $\Sect$ is contained in
$\rho(m(C_0))+D$ and that one has $\rho(C_0)=\rho_{\bA;C}(C_0)$ for every
chamber $C$ meeting $\Sect$.

Let us make $\bA$ and $A'$ into vector spaces by choosing as origin the points
$m(C_0)$ and $\rho(m(C_0))$ respectively. One then has
$\Sect'=\bR_+.Q(C_0)$.

We shall show that, for every $v\in Q(C_0)$, every $t\in\bR_+$ and every
chamber $C$ of $\bA$ meeting $\Sect$, one has
\begin{equation}
\rho_{\bA;C}(tv)=t\rho(v)=\rho(tv).
\end{equation}
By continuity it suffices to prove (1) when the line $\bR v$ meets no facet of
codimension 2, which we suppose henceforth. There then exists a strictly
increasing sequence $(t_0,\dots,t_{n+1})$ of points of the interval $[0,t]$,
with $t_0=0$ and $t_{n+1}=t$, and a minimal gallery $(C_0,\dots,C_n)$
contained in $A'$ such that
$C_i\cap[0,tv]\supset\,\left]t_iv,\,t_{i+1}v\right[$ for $0\le i\le n$. Let us
argue by induction on the integer $n$, relation (1) being true when $n=0$.
Suppose $n>0$; in view of the induction hypothesis, relation (1), after
substitution of $s$ for $t$, is true for $0\le s\le t_n$. One deduces that
$\rho(C_{n-1})=\rho_{\bA;C}(C_{n-1})$.

It suffices to show that $\rho_{\bA;C}(C_n)$ is the chamber $C'_n$ adjacent to
$C'_{n-1}=\rho_{\bA;C}(C_{n-1})$ along the panel $F$, the image under $\rho$
of the panel common to $C_{n-1}$ and $C_n$. Indeed, this implies that the
restriction of $\rho_{\bA;C}$ to $\overline{C}_{n-1}\cup\overline{C}_n$ is
isometric, and hence that the map $t\mapsto\rho_{\bA;C}(tv)$ from
$[t_{n-1},t_{n+1}]$ into $\bA$ is affine. The first equality (1) follows from
it. The second reduces to the first by choosing the chamber $C$ ``sufficiently
far'' in $\Sect''$ so that $\rho$ and $\rho_{\bA;C}$ coincide on
$C_0\cup\dots\cup C_n$.

Let then $L$ be the wall of $\bA$ containing $F$ and let $\alpha$ be the
affine root with wall $L$ containing $C'_{n-1}$. One then has
$\rho([t_{n-1}v,t_nv])\subset\alpha$, whence, in view of the induction
hypothesis, $\rho([0,t_nv])\subset\alpha$ and
$C'_0=\rho(C_0)\subset\alpha$. Moreover $\alpha^{*}+w(C_0).D=\alpha^{*}$. If
$\alpha\supset C$, the sought relation $\rho_{\bA;C}(C_n)=C'_n$ follows
immediately from (\btref{2.3.12}). If $\alpha$ contains a sector parallel to $\Sect$,
one has $\alpha\supset\rho(C_0)+D$, and $\alpha$ contains $\Sect$, hence $C$.
There finally remains to examine the case where $\alpha^{*}$ contains $C$ and
a sector of direction $D$. If $\rho(C_n)\neq C'_n$, one has
$\rho(C_n)=C'_{n-1}$ and the segment $[\rho(t_nv),\rho(t_{n+1}v)]$ is
contained in the symmetric, with respect to $L$, of the half-line
$[t_n,\infty[\,.v$. One deduces that $w(C_n)=s.w(C_0)$, where $s$ is the image
in ${}^{v}\bW$ of the reflection with respect to $L$. But $D$ and $w(C_0).D$
are on the same side of the hyperplane of the fixed points of $s$, which
entails that the length of $sw(C_0)$ is strictly greater than that of
$w(C_0)$, whence a contradiction in view of the choice of $C_0$. Consequently
one has $\rho(C_n)=C'_n$. Let us now show that $\rho_{\bA;C}(C_n)=C'_n$.
Suppose on the contrary that $\rho_{\bA;C}(C_n)=C'_{n-1}$; there then exists a
chamber $C''$ adjacent to $C_n$ and to $C_{n-1}$ such that
$\rho_{\bA;C}(C'')=C'_n$ (\btref{2.3.12}). The length $h$ of a taut gallery
between $C$ and $C''$ is then equal to that of a taut gallery between $C$
and $C'_n$ and, since $C$ and $C'_{n-1}$ are on either side of $L$, the length
of a taut gallery between $C$ and $C'_{n-1}$ is equal to $h+1$. One
deduces that $\rho(C'')=C'_n$; indeed,
\pagemark{60}
$\rho(C'')$ admits $F$ as panel, hence is equal to $C'_n$ or to $C'_{n-1}$,
and cannot be equal to $C'_{n-1}$ since the image under $\rho$ of a gallery
stretched between $C$ and $C''$ is a gallery of length $h$ joining $C$ to
$\rho(C'')$. But this contradicts (\btref{2.3.12}): taking a chamber $C'''$ contained
in a sufficiently small sector contained in $\Sect$, one deduces indeed that
the two adjacent chambers $C_n$ and $C''$ are sent by the retraction
$\rho_{\bA;C'''}$ onto the same chamber $C'_n$, which is situated on the same
side of $L$ as $C'''$. Consequently we have indeed shown that
$\rho_{\bA;C}(C_n)=C'_n$, which completes the proof of (1).

Let us then show that the restriction of $\rho$ to $\Sect\cup\Sect'$ is
isometric. Relation (1) entails that, for $x\in\Sect$ and $y\in\Sect'$, one
has
\[
d(\rho(x),\rho(y))=d(x,\rho_{\bA;C}(y))=d(x,y)
\]
on taking for $C$ a chamber containing $x$ in its closure and applying
(\btref{2.5.3})(ii). If $x,y\in\Sect'$, there exists $t\in\bR^{*}_+$ such that $tx$
and $ty$ belong to $Q(C_0)$, and one then has, in view of (1) and of the fact
that the restriction of $\rho$ to $Q(C_0)$ is isometric,
\begin{align*}
d(\rho(x),\rho(y))&=d(t^{-1}\rho(tx),\,t^{-1}\rho(ty))=t^{-1}d(\rho(tx),\rho(ty))\\
&=t^{-1}d(tx,ty)=d(x,y).
\end{align*}
Since $\rho$ is the identity on $\Sect$, this completes the proof.

\stmt{Corollary}{2.9.6}\btanchor{2.9.6}
\emph{Let us take up again the hypotheses of (\btref{2.8.2}), and let $\Sect$ and
$\Sect'$ be two sectors of $\Imm$. There exist sectors $\Sect_1\subset\Sect$
and $\Sect'_1\subset\Sect'$ contained in one and the same apartment.}

\medskip
This follows from (\btref{2.9.5}) and (\btref{2.8.4}).

%==================================================================

\btsecA{3}{3. Bounded subgroups}

%==================================================================
\pagemark{61}

We keep the notation of the preceding paragraphs.

\btsecB{3.1}{3.1. Bornology defined by a Tits system}

\stmt{Definition}{3.1.1}\btanchor{3.1.1}
\emph{A \textbf{bornology} on a set $X$ is a set $\Bor$ of subsets of $X$,
called \textbf{bounded subsets}, stable under finite union, containing the
finite subsets of $X$, and such that $M\in\Bor$ and $M'\subset M$ entail
$M'\in\Bor$. If $X$ is a group, one says that $\Bor$ is \textbf{compatible
with the group law} of $X$, or makes $X$ a \textbf{bornological group}, if
$M,M'\in\Bor$ entails $M^{-1}M'\in\Bor$.}

\stmt{Example}{3.1.2}\btanchor{3.1.2}
a) The set of the relatively compact subsets of a Hausdorff topological group
$X$ makes $X$ a bornological group.

b) Let $E$ be a metric space. The group $\Isom E$ of the isometries of $E$
possesses a natural bornology formed of the sets $M$ possessing the following
equivalent properties:
\begin{itemize}\itemsep2pt
\item[(i)] there exists a point $x\in E$ such that the set of the $g(x)$ for
$g\in M$ is bounded in $E$;
\item[(ii)] for every bounded subset $F$ of $E$, the set of the $g(x)$ for
$g\in M$ and $x\in F$ is bounded.
\end{itemize}

c) Let $X'$ be a bornological group and $\varphi:X\to X'$ a homomorphism of
groups. The set of the subsets of $X$ whose image is a bounded subset of $X'$
forms a bornology on $X$, compatible with the group law of $X$, and called the
\emph{inverse image of the bornology of $X'$ by $\varphi$}.

\stmt{Proposition}{3.1.3}\btanchor{3.1.3}
\emph{Let $\Bor$ be the set of the subsets of $G$ whose canonical image in
$B\backslash G/B$ is finite. Then $\Bor$ is a bornology on $G$ compatible with
the group law of $G$.}

\medskip
That $\Bor$ be a bornology is evident. On the other hand it is clear that
$M\in\Bor$ entails $M^{-1}\in\Bor$, and it suffices to show that
$M,M'\in\Bor$ entails $MM'\in\Bor$, which is immediate since the product of
two double cosets modulo $B$ belongs to $\Bor$ ([\btbib{5}], chap. IV, \S\,2, n$^{o}$
1, lemma 1).

\stmt{Definition}{3.1.4}\btanchor{3.1.4}
\emph{The bornology $\Bor$ introduced in (\btref{3.1.3}) is said to be \textbf{defined
by the Tits system} $(G,B,N)$.}

\medskip
Since $M\in\Bor$ is equivalent to $gMg^{-1}\in\Bor$ (whatever $g\in G$), it is
clear that $\Bor$ depends only on the conjugacy class of $B$. \emph{Two
equivalent Tits systems define the same bornology}: we shall see later the
converse (\btref{3.5.1}).

\pagemark{62}
\stmt{Proposition}{3.1.5}\btanchor{3.1.5}
\emph{Let $\Bor$ be a bornology on $G$, compatible with the group law and such
that $B\in\Bor$. Let $\Bor_{\bW}$ be the set of the subsets $M$ of $\bW$ such
that $BMB\in\Bor$. Then $\Bor_{\bW}$ makes $\bW$ a bornological group. If $X$
is a subset of $\bW$ such that the union of the $BxBx^{-1}B$, for $x\in X$, is
bounded for $\Bor$, the set of the reflections associated with the various
elements of $X$ is bounded in $\bW$.}

\medskip
The first assertion is evident. The second follows from the fact that, if $t$
belongs to the set $T_w$ of the reflections associated with $w\in\bW$
(\btref{2.3.10}), one has $BtB\subset BwBw^{-1}B$ ([\btbib{5}], chap. IV, \S\,2, n$^{o}$ 4,
cor. 2 of th. 2).

\stmt{Corollary}{3.1.6}\btanchor{3.1.6}
\emph{Let us endow $G$ with the bornology $\Bor$ defined by the Tits system
$(G,B,N)$ and let $M$ be a subset of $G$. In order that $M$ be bounded, it is
necessary and sufficient that the union of the $gBg^{-1}$ for $g\in M$ be so.}

\medskip
Set $M'=\bigcup_{g\in M} gBg^{-1}$. If $M$ is bounded, one has
$M'\subset MBM^{-1}\in\Bor$. Conversely, suppose $M'$ bounded and let
$X\subset\bW$ be such that $BXB=BMB$. In view of (\btref{3.1.5}), the set of the
reflections associated with the various elements of $X$ is bounded in $\bW$,
that is to say finite. Since every $w\in\bW$ is a product of the reflections
associated with it, taken once and only once in a suitable order, $X$ itself
is finite and $M$ is bounded.

\stmt{Remark}{3.1.7}\btanchor{3.1.7}
The preceding results are valid for an arbitrary Tits system, not necessarily
of affine type.

\stmt{Proposition}{3.1.8}\btanchor{3.1.8}
\emph{Let $\Isom\Imm$ be the group of the isometries of the building $\Imm$ of
$(G,B,N)$. The bornology of $G$ defined by the Tits system $(G,B,N)$ is the
inverse image of the natural bornology of $\Isom\Imm$ by the canonical
homomorphism of $G$ into $\Isom\Imm$.}

\medskip
Let $M\subset G$ and let $X$ be the canonical image of $M$ in
$B\backslash G/B$. To say that the image of $M$ in $\Isom\Imm$ is bounded
amounts to saying that
\[
\sup_{g\in M} d\bigl(F(B),F(gBg^{-1})\bigr)<+\infty.
\]
Now, if $g\in BxB$ with $x\in\bW$, one has
$d(F(B),F(gBg^{-1}))=d(\bC,x(\bC))$. Since $\bW$ operates properly on $\bA$,
one has $\sup_{x\in X} d(\bC,x(\bC))<+\infty$ if and only if $X$ is finite,
whence the proposition.

\num{3.1.9}\btanchor{3.1.9}
Let $\varphi:G\to\Ghat$ be a $B$-adapted homomorphism and let
$B_1=\Stab B$ be the normalizer of $\varphi(B)$ in $\Ghat$. We defined in 2.7
a homomorphism of $\Ghat$ into $\Isom\Imm$ ``prolonging'' the canonical
homomorphism of $G$ into $\Isom\Imm$. We shall again call \emph{bornology
defined by the Tits system $(G,B,N)$} in $\Ghat$ the inverse image by this
homomorphism of the natural bornology of $\Isom\Imm$.

\medskip\noindent
\textit{Proposition}.\ ---\ \emph{Let us endow $G$ and $\Ghat$ with the
bornologies defined by the Tits system $(G,B,N)$. Let $M\subset\Ghat$; the
following conditions are equivalent:
\begin{itemize}\itemsep2pt
\item[(i)] $M$ is bounded in $\Ghat$;
\item[(ii)] $\varphi^{-1}(MB_1)$ is bounded in $G$;
\item[(iii)] the image of $M$ in $B_1\backslash\Ghat/B_1$ is finite.
\end{itemize}}

\pagemark{63}
Let us show that (i) entails (ii). Since $B_1$ is bounded in $\Ghat$, as
stabilizer of the chamber $F(B)$, $MB_1$ is bounded as soon as $M$ is, and the
image of $MB_1$ is bounded in $\Isom\Imm$, whence (ii).

If $\varphi^{-1}(MB_1)\subset BXB$ with $X\subset\bW$, one has
$M\subset B_1XB_1$; one deduces that (ii) entails (iii).

Finally, (iii) entails (i) since $B_1$ is bounded.

\btsecB{3.2}{3.2. A fixed point lemma}

\num{3.2.1}\btanchor{3.2.1}
Let $x,y\in\Imm$. By (\btref{2.5.4}) there exists a unique point $m\in\Imm$ such that
$d(x,m)=d(y,m)=\tfrac12 d(x,y)$: it is called the \emph{midpoint} of
$\{x,y\}$.

\medskip\noindent
\textit{Lemma}.\ ---\ \emph{Let $x,y,z\in\Imm$ and let $m$ be the midpoint of
$\{x,y\}$. One has}
\begin{equation}
d(x,z)^2+d(y,z)^2\ \ge\ 2d(m,z)^2+\tfrac12 d(x,y)^2 .
\end{equation}

\medskip
Let $A$ be an apartment of $\Imm$ containing $x$ and $y$, let $C$ be a chamber
of $A$ whose closure contains $m$, and set $z'=\rho_{A;C}(z)$ (\btref{2.3.5}). In view
of (\btref{2.5.3}) one has $d(x,z)\ge d(x,z')$, $d(y,z)\ge d(y,z')$ and
$d(m,z)=d(m,z')$. Relation (1) then follows from the equality
\[
d(x,z')^2+d(y,z')^2=2d(m,z')^2+\tfrac12 d(x,y)^2
\]
valid for every point $z'$ of the Euclidean space $A$ ([\btbib{31}], book VII,
prop. 122).

\stmt{Remark}{3.2.2}\btanchor{3.2.2}
Let $E$ be a metric space and let $x,y,m$ be three points of $E$ such that
relation (1) is satisfied for every point $z\in E$. Making successively $z=x$
and $z=y$, one sees that $d(m,x)=d(m,y)=\tfrac12 d(x,y)$. Moreover, if
$m'\in E$ is such that $d(m',x)=d(m',y)=\tfrac12 d(x,y)$, one has $m=m'$, as
one sees on making $z=m'$ in (1). In particular, for $x,y\in\Imm$, the
midpoint $m$ of $\{x,y\}$ is the unique point of $\Imm$ satisfying (1) for
every $z\in\Imm$.

\stmt{Lemma}{3.2.3}\btanchor{3.2.3}
\emph{Let $E$ be a complete metric space and $E'$ a subset of $E$ possessing
the following property:
\begin{itemize}
\item[(CN)] whatever the points $x$ and $y$ of $E'$, there exists a point
$m\in E'$ such that relation (1) of (\btref{3.2.1}) is satisfied for every point $z$
of $E'$.
\end{itemize}
If $M$ is a nonempty bounded subset of $E'$, the stabilizer of $M$ in
$\Isom E$ possesses a fixed point in the closure of $E'$ in $E$.}

\medskip
If $X$ and $Y$ are two subsets of $E$, we shall set
\[
\diam(X,Y)=\sup_{x\in X,\,y\in Y} d(x,y),\qquad \diam X=\diam(X,X).
\]
Let $k$ be a real number such that $0<k<1$. For every subset $X$ of $E'$, let
us denote by $f(X)$ the set of the points $m\in E'$ for which there exist
$x,y\in X$ such that relation (1) is
\pagemark{64}
satisfied for every $z\in E'$ and that $d(x,y)\ge k\,\diam X$; in view of
(CN), one has $f(X)\neq\emptyset$ if $X\neq\emptyset$.

For points $x,y,m$ satisfying these conditions, and for $z\in E'$, one has:
\begin{equation}
d(m,z)^2\ \le\ \tfrac12\bigl(d(x,z)^2+d(y,z)^2\bigr)-\tfrac14 d(x,y)^2
\ \le\ \diam(X,\{z\})^2-\tfrac{k^2}{4}(\diam X)^2 .
\end{equation}
For $z\in X$ one draws from this:
\begin{equation}
\diam(f(X),X)\ \le\ k_1\,\diam X\qquad\text{with }
k_1=\Bigl(1-\tfrac{k^2}{4}\Bigr)^{1/2}<1 .
\end{equation}
Moreover, taking $z\in f(X)$ in (2), one draws from (2) and (3):
\begin{equation}
\diam f(X)\ \le\ k_2\,\diam X\qquad\text{with }
k_2=\Bigl(1-\tfrac{k^2}{2}\Bigr)^{1/2}<1 .
\end{equation}
Let $M$ be a nonempty bounded subset of $E'$; by (4) one has
\begin{equation}
\diam f^{q}(M)\ \le\ k_2^{\,q}\,\diam M\qquad\text{for every integer }q\ge 1.
\end{equation}
Let us then choose a point $x_q\in f^{q}(M)$ for $q\in\bN^{*}$, which is
permissible since $f^{q}(M)$ is not empty. It follows from (3) and (5) that
$d(x_q,x_{q+1})\le k_1k_2^{\,q}\,\diam M$, and the Cauchy sequence $(x_q)$
converges to a point $x\in\overline{E'}$, independent of the choice of the
$x_q$ by (5). Finally, it is clear that the stabilizer of $M$ in $\Isom E$
leaves invariant each of the $f^{q}(M)$, and hence also $x$. This completes
the proof.

One will note that a priori the point $x$ thus constructed depends on the
choice of the constant $k$.

\stmt{Proposition}{3.2.4}\btanchor{3.2.4}
\emph{Let $M$ be a nonempty bounded subset of the building $\Imm$ of $G$. The
stabilizer of $M$ in $\Isom\Imm$ possesses a fixed point belonging to the
closure of the convex hull of $M$.}

\medskip
This follows from lemmas (\btref{3.2.1}) and (\btref{3.2.3}), taking in the latter for $E$ the
building $\Imm$ and for $E'$ the convex hull of $M$.

\stmt{Remark}{3.2.5}\btanchor{3.2.5}
Let $E$ be a metric space and $x,y,m$ three points of $E$. The proof of lemma
(\btref{3.2.1}) shows that if there exists a map $\rho$ from $E$ into a Euclidean
space, decreasing distances, such that $d(\rho(m),\rho(z))=d(m,z)$ for every
$z\in E$ and such that $\rho(m)$ is the midpoint of the segment
$[\rho(x)\rho(y)]$, then relation (1) is satisfied for every $z\in E$. This
applies in particular to the case where $E$ is a \emph{simply connected
Riemannian space of negative curvature} and where $m$ is the midpoint of the
geodesic segment $[xy]$: it suffices to take for $\rho$ the map of $E$ onto
its tangent space at $m$, inverse of the exponential map ([\btbib{24}], chap. I,
\S\,13). Lemma (\btref{3.2.3}) then furnishes in this case a simple proof of the
classical result on the existence of a fixed point for a compact group of
isometries of such a Riemannian manifold. It is rather curious that this same
lemma, which will serve us to determine the conjugacy classes of maximal
bounded subgroups of semisimple algebraic groups over a local field, can thus
be used to prove the conjugacy of maximal compact subgroups of real Lie
groups.
\par\endgroup

% ---------------------------------------- bruhattitsen_1.tex
\begingroup\providecommand{\imm}{}
\renewcommand{\imm}{\mathscr{I}}
\providecommand{\borno}{}
\renewcommand{\borno}{\mathscr{B}}
\providecommand{\Bfr}{}
\renewcommand{\Bfr}{\mathfrak{B}}
\providecommand{\Bfrh}{}
\renewcommand{\Bfrh}{\widehat{\mathfrak{B}}}
\providecommand{\Cfr}{}
\renewcommand{\Cfr}{\mathfrak{C}}
\providecommand{\Gh}{}
\renewcommand{\Gh}{\widehat{G}}
\providecommand{\Bh}{}
\renewcommand{\Bh}{\widehat{B}}
\providecommand{\Nh}{}
\renewcommand{\Nh}{\widehat{N}}
\providecommand{\Hh}{}
\renewcommand{\Hh}{\widehat{H}}
\providecommand{\Ph}{}
\renewcommand{\Ph}{\widehat{P}}
\providecommand{\Wh}{}
\renewcommand{\Wh}{\widehat{\mathbf{W}}}
\providecommand{\Vh}{}
\renewcommand{\Vh}{\widehat{\mathbf{V}}}
\providecommand{\nuh}{}
\renewcommand{\nuh}{\widehat{\nu}}
\providecommand{\bA}{}
\renewcommand{\bA}{\mathbf{A}}
\providecommand{\bW}{}
\renewcommand{\bW}{\mathbf{W}}
\providecommand{\bV}{}
\renewcommand{\bV}{\mathbf{V}}
\providecommand{\bS}{}
\renewcommand{\bS}{\mathbf{S}}
\providecommand{\bC}{}
\renewcommand{\bC}{\mathbf{C}}
\providecommand{\bD}{}
\renewcommand{\bD}{\mathbf{D}}
\providecommand{\bSig}{}
\renewcommand{\bSig}{\mathbf{\Sigma}}
\providecommand{\vp}{}
\renewcommand{\vp}[1]{{}^{v}\!#1}
\providecommand{\dgr}{}
\renewcommand{\dgr}{\dagger}
\providecommand{\dle}{}
\renewcommand{\dle}{\leqslant_{\bD}}
\providecommand{\Stab}{}\renewcommand{\Stab}{\operatorname{Stab}}
\providecommand{\cl}{}\renewcommand{\cl}{\operatorname{cl}}
\providecommand{\Card}{}\renewcommand{\Card}{\operatorname{Card}}
\providecommand{\Cox}{}\renewcommand{\Cox}{\operatorname{Cox}}
\providecommand{\Aut}{}\renewcommand{\Aut}{\operatorname{Aut}}
\providecommand{\Ker}{}\renewcommand{\Ker}{\operatorname{Ker}}
\providecommand{\Id}{}\renewcommand{\Id}{\operatorname{Id}}
\newenvironment{state}[2]%
  {\par\medskip\noindent\textit{#1 \textup{(#2)}}. --- \begingroup\itshape}%
  {\endgroup\par\medskip}
\newenvironment{states}[1]%
  {\par\medskip\noindent\textit{#1}. --- \begingroup\itshape}%
  {\endgroup\par\medskip}
\providecommand{\num}{}
\renewcommand{\num}[1]{\par\medskip\noindent\textbf{(#1)}\ \ignorespaces}
\setlength{\parskip}{2pt}
\bigskip

\noindent
Moreover, the building $\imm$ plays for $G$ a role somewhat analogous to the
one played, for a real semisimple Lie group, by the Riemannian symmetric
space obtained as the quotient by a maximal compact subgroup: it is a
complete contractible metric space on which the group acts, the stabilizers
of each point being bounded, and lemma (\btref{3.2.1}) expresses a kind of
``negative curvature'' property. For other examples of this analogy, see
[\btbib{35}].

%=====================================================================

\btsecB{3.3}{3.3. Maximal bounded subgroups}

%=====================================================================

Let $\varphi : G \to \Gh$ be a $B$-adapted homomorphism; endow $\Gh$ with the
bornology defined by the Tits system $(G,B,N,S)$ as in (\btref{3.1.9}).

\begin{state}{Theorem}{3.3.1}\btanchor{3.3.1}
In order that a subgroup of $\Gh$ be bounded, it is necessary and sufficient
that it stabilize a parahoric subgroup of $G$.
\end{state}

If a subgroup $\Gamma$ of $\Gh$ stabilizes a parahoric subgroup $P$ of $G$,
it fixes the facet $F(P)$ of $\imm$ and is consequently bounded, by the very
definition of the bornology of $\Gh$.

Conversely, if $\Gamma$ is bounded, every orbit of $\Gamma$ in $\imm$ is
bounded, and (\btref{3.2.4}) implies that $\Gamma$ possesses a fixed point $x$ in
$\imm$. Let $P$ be the parahoric subgroup of $G$ such that $x \in F(P)$.
Then $\Gamma$ preserves $F(P)$ and one has ${}^{g}P = P$ for all
$g \in \Gamma$ (\btref{2.7.1}), which completes the proof.

\begin{state}{Corollary}{3.3.2}\btanchor{3.3.2}
Every bounded subgroup of $\Gh$ is contained in a maximal bounded subgroup.
The maximal bounded subgroups are the maximal elements of the set of
stabilizers of parahoric subgroups of $G$.
\end{state}

This follows immediately from the theorem and from (\btref{1.2.19}), which shows
that the stabilizer of a parahoric subgroup is contained in only finitely
many stabilizers of parahoric subgroups.

\begin{state}{Corollary}{3.3.3}\btanchor{3.3.3}
The maximal bounded subgroups of $G$ are the maximal parahoric subgroups. In
particular, they form $\prod_{1\leqslant i\leqslant m}(\ell_i+1)$ conjugacy
classes (where the integers $\ell_i$, for $1\leqslant i\leqslant m$, are the
dimensions of the irreducible components of $\bA$).
\end{state}

It suffices to apply (\btref{3.3.2}) to the case $\varphi = \Id_G$.

\num{3.3.4}\btanchor{3.3.4}
Recall that in (\btref{1.2.16}) we defined a homomorphism $\xi$ from $\Gh$ into the
group of automorphisms of the Coxeter graph of $\bW$, such that, if $P$ is a
parahoric subgroup of type $X$ of $G$, then ${}^{g}P$ is of type $\xi(g)(X)$
for all $g\in\Gh$. Put $\Xi = \xi(\Gh)$. The following proposition is an
``abstract'' version of a result of Iwahori--Matsumoto:

\begin{states}{Proposition}
Let $P$ and $P'$ be two parahoric subgroups of $G$, let $X$ and $X'$ be their
types and let $\Xi_X$ be the stabilizer of $X$ in $\Xi$. In order that
$\Stab P$ be a maximal bounded subgroup of $\Gh$, it is necessary and
sufficient that $X$ be maximal among the types of parahoric subgroups
invariant under $\Xi_X$. In order that $\Stab P$ and $\Stab P'$ be conjugate
in $\Gh$, it is necessary and sufficient that there exist $t\in\Xi$ such that
$X' = t(X)$.
\end{states}

One may assume $P = B_X$. We recalled (\btref{1.2.19}) that
$\Stab P \subset \Stab P'$ is equivalent to $P' = B_{X'}$ with
$X \subset X'$ and $\Xi_X \subset \Xi_{X'}$. This last condition means that
$X'$ is invariant under $\Xi_X$, whence the first assertion.

If there exists $g\in\Gh$ such that $\Stab P' = g(\Stab P)g^{-1}$, one has
$P' = {}^{g}P$, whence $X' = \xi(g)(X)$. Conversely, if there exists
$t\in\Xi$ such that $X' = t(X)$, and if $g \in \xi^{-1}(t)$, the subgroups
${}^{g}P$ and $P'$ are of the same type, hence conjugate, and the same holds
for $\Stab P'$ and $\Stab P = g^{-1}(\Stab {}^{g}P)g$.

\num{3.3.5}\btanchor{3.3.5}
Suppose the Tits system $(G,B,N)$ irreducible. Then \emph{$\Stab P$ is
maximal if and only if the type $X$ of $P$ is the complement of an orbit of a
subgroup of $\Xi$ in $\bS$.}

If moreover $\Xi$ is cyclic of order $q$ (which is frequently the case in the
applications to simple algebraic groups), the number $c$ of conjugacy classes
of maximal bounded subgroups of $G$ satisfies the inequality $c < \ell + 1$
(where $\ell$ is the dimension of $\bA$), except when $q = 1$ or $2$, in
which case one has $c = \ell + 1$. Indeed, one sees easily that if $\Xi$
permutes cyclically a subset $Y$ of $\bS$, the number of conjugacy classes of
maximal bounded subgroups of $\Gh$ of the form $\Stab B_X$ with
$X \supset \bS - Y$ is equal to the number of divisors of $\Card Y$.

\num{3.3.6}\btanchor{3.3.6}
For examples of the explicit determination of the conjugacy classes of
maximal subgroups $\Stab P$, see [\btbib{28}].

%=====================================================================

\btsecB{3.4}{3.4. Characterization of the bornology defined by a Tits system of affine type}

%=====================================================================

We resume the notation of 2.6.

\begin{state}{Theorem}{3.4.1}\btanchor{3.4.1}
Let $\borno$ be a bornology on $G$, compatible with the group law of $G$ and
such that $B$ is bounded. There exists one and only one subset
$J \subset \{1,\dots,m\}$ such that $\borno$ is the bornology defined by the
Tits system $(G, B^{J}, N)$.
\end{state}

We shall first establish two lemmas.

\begin{state}{Lemma}{3.4.2}\btanchor{3.4.2}
Suppose $\bW$ irreducible and let $\bW^{0}$ be the stabilizer in $\bW$ of a
special point $a \in \overline{\bC}$. If $X$ is an infinite subset of $\bW$,
every reflection of $\bW$ is associated with at least one element of
$\bW^{0}\!\cdot\! X$.
\end{state}

Make $\bA$ into a vector space by taking $a$ as origin. Let $r$ be a
reflection of $\bW$ and let $\varphi$ be a linear form on $\bA$ and
$k \in \mathbf{R}$ be such that the equation $\varphi(x) = k$ defines the wall
$L$ of $\bA$, the set of fixed points of $r$, and such that $\bC$ is
contained in $\{x\in\bA \mid \varphi(x) < k\}$. In view of the irreducibility
of $\bW$, the set $\Phi$ of the transforms of $\varphi$ by $\vp{\bW}$
generates the dual of $\bA$, and $\Phi = -\Phi$, since the reflection in the
wall $\Ker\varphi$ transforms $\varphi$ into $-\varphi$. It follows that the
set
\[
E \;=\; \bigcap_{w \in \vp{\bW}} w\cdot\{x\in\bA \mid \varphi(x) < k\}
\]
is bounded, and that there exists $x \in X$ such that $x(\bC) \not\subset E$.
There then exists $w \in \bW^{0}$ such that
$wx(\bC) \subset \{x\in\bA \mid \varphi(x) > k\}$, and the reflection $r$ is
associated with $wx$ by (\btref{2.3.10}).

\begin{state}{Lemma}{3.4.3}\btanchor{3.4.3}
There exists an integer $N$ such that every element of $\bW$ is a product of
at most $N$ reflections of $\bW$.
\end{state}

Make $\bA$ into a vector space by taking a special point as origin, and
resume the notation of (\btref{1.3.8}). Let $\{a_1,\dots,a_\ell\}$ be a basis of the
root system associated with $\bW$. For every root $a \in \vp{\bSig}$ and
every integer $n \in \mathbf{Z}$, the composite $r_{a,n}\circ r_{a,0}$ is the
translation by the vector $-n\cdot a^{\vee}$; since the translations by the
vectors $a_i^{\vee}$ for $1\leqslant i\leqslant \ell$ form a basis of the
group $\bV$ of translations of $\bW$, and since $\bW$ is the semidirect
product of the stabilizer $\bW^{0}$ of the origin by $\bV$ (\btref{1.3.7}), one
deduces that one may take $N = N' + 2\ell$, where $N'$ is the smallest
integer such that every element of $\bW^{0}$ (which is finite and isomorphic
to $\vp{\bW}$) can be written as a product of at most $N'$ reflections.

\begin{state}{Remark}{3.4.4}\btanchor{3.4.4}
\upshape
Put $M = \bigcup_{g\in G} gBg^{-1}$ and let $N$ be an integer satisfying the
conditions of (\btref{3.4.3}). One then has $G = M^{2N}$. Indeed, for every
reflection $r \in \bW$ there exists $x \in \bW$ such that $r \in T_x$ (for
example $r$ itself). One then has
$BrB \subset BxBx^{-1}B \subset B\cdot M \subset M^{2}$. If $w \in \bW$ is
the product of $k \leqslant N$ reflections $r_j$, one therefore has
$BwB \subset \prod_j Br_jB \subset M^{2N}$.
\end{state}

\num{3.4.5}\btanchor{3.4.5}
Let us now prove theorem (\btref{3.4.1}). Let $\borno$ be a bornology on $G$,
compatible with the group law and such that $B$ is bounded. For each
$j = 1,\dots,m$, consider the bornology $\borno_j$ induced by $\borno$ on
$B_j = B\bW_j B$, and let $\borno_{\bW}$ (resp. $\borno_{\bW_j}$) be the
bornology on $\bW$ (resp. $\bW_j$) deduced from $\borno$ (resp. $\borno_j$)
as in (\btref{3.1.5}). If there exists a subset $X$ of $\bW_j$ which is infinite and
bounded for $\borno_{\bW_j}$, lemma (\btref{3.4.2}) together with (\btref{3.1.5}) shows that
the set of reflections of $\bW_j$ is bounded, hence, by (\btref{3.4.3}), that
$\bW_j$ itself is bounded for $\borno_j$. In other words, either
$B_j \in \borno$, or $\borno_j$ is the bornology defined by the Tits system
$(B_j, B, B_j\cap N)$.

Let then $J = \{\,j \in \{1,\dots,m\} \mid B_j \in \borno\,\}$. We shall show
that $\borno$ is the bornology defined by $(G, B^{J}, N)$. Replacing $B$ by
$B^{J}$, we may suppose that $J = \varnothing$. Let then
$X \in \borno_{\bW}$ and let $Y$ be the set of reflections associated with
the various points of $X$. We have seen (\btref{3.1.5}) that $Y$ is bounded for
$\borno_{\bW}$; consequently, by the first part of the proof, $Y \cap \bW_i$
is \emph{finite} for every $i$. Since every reflection of $\bW$ is contained
in one of the $\bW_i$, it follows that $Y$ itself is finite. Hence $X$ is
finite, whence our assertion.

Finally the uniqueness of $J$ is clear: it is the largest subset of
$\{1,\dots,m\}$ such that $B^{J}$ is bounded.

\begin{state}{Corollary}{3.4.6}\btanchor{3.4.6}
The bornology defined by $(G,B,N)$ is contained in every bornology compatible
with the group law of $G$ and for which $B$ is bounded. The bornologies
defined by the irreducible components of $(G,B,N)$ are the maximal
bornologies among the bornologies compatible with the group law of $G$ for
which $B$ is bounded and $G$ is not. If the system $(G,B,N)$ is irreducible,
the bornology it defines is the unique bornology compatible with the group
law of $G$ for which $B$ is bounded and $G$ is not.
\end{state}

\begin{state}{Corollary}{3.4.7}\btanchor{3.4.7}
Let $\varphi : G \to \Gh$ be a $B$-adapted homomorphism. The bornology
defined on $\Gh$ by the Tits system $(G,B,N)$ is the unique bornology on
$\Gh$ which is compatible with the group structure of $\Gh$ and for which
$\Bh$ is bounded and $\varphi(B^{J})$ is unbounded for every non-empty subset
$J$ of $\{1,\dots,m\}$.

If in particular the Tits system $(G,B,N)$ is irreducible, the bornology it
defines in $\Gh$ is the unique bornology compatible with the group structure
of $\Gh$ for which $\Bh$ is bounded and $\Gh$ is not.
\end{state}

If a bornology on $\Gh$ satisfies the conditions of the first part of the
statement, it induces on the subgroup $\Gh_0 = \Ker\xi$ (with the notation of
(\btref{1.2.17})) the bornology defined by the Tits system $(\Bh, \varphi(N))$, which
is obviously the same as the bornology defined in $\Gh_0$ by $(G,B,N)$ and
the homomorphism $\varphi : G \to \Gh_0$. Now $\Gh_0$ is of finite index in
$\Gh$, whence the first assertion of the corollary. The second follows.

%=====================================================================

\btsecB{3.5}{3.5. Characterization of a Tits system of affine type by its bornology}

%=====================================================================

Theorem (\btref{3.3.1}) shows that the maximal parahoric subgroups of the Tits system
$(G,B,N)$ are well determined by the datum of the bornology defined by this
Tits system. More precisely:

\begin{state}{Theorem}{3.5.1}\btanchor{3.5.1}
Two Tits systems of affine type in the same group are equivalent if and only
if they define the same bornology.
\end{state}

Let us begin by establishing two lemmas.

\begin{state}{Lemma}{3.5.2}\btanchor{3.5.2}
Suppose $\bW$ irreducible and let $P$, $P'$ be two distinct maximal parahoric
subgroups. In order that $P\cap P'$ be a parahoric subgroup, it is necessary
and sufficient that it be a maximal subgroup of $P$.
\end{state}

If $P \cap P'$ is a parahoric subgroup, one may assume that
$B \subset P\cap P'$. Since $\bW$ is irreducible, there exist two distinct
elements $s, s' \in \bS$ such that $P = B_{\bS - \{s\}}$ and
$P' = B_{\bS-\{s'\}}$, and $P\cap P' = B_{\bS-\{s,s'\}}$ is maximal in $P$.

Conversely, suppose $P\cap P'$ maximal in $P$. Let $D$ be a ball with center
$x = a_P$ (\btref{2.1.2}) possessing the properties of lemma (\btref{2.5.11}) and let
$a \in D \cap [a_P a_{P'}]$ with $a \neq a_P$. Let $Q$ be the parahoric
subgroup such that $a \in F(Q)$. One has $P\cap P' \subset Q$ by (\btref{2.5.4})(iii)
and $a_P \in \overline{F(Q)}$ by (\btref{2.5.11}), whence $Q \subset P$. One deduces
$P\cap P' = Q$.

\begin{state}{Lemma}{3.5.3}\btanchor{3.5.3}
Let $P_1,\dots,P_q$ be parahoric subgroups. The intersection
$P_1\cap\dots\cap P_q$ is a parahoric subgroup if and only if $P_h \cap P_k$
is a parahoric subgroup for $1\leqslant h,k\leqslant q$.
\end{state}

That the condition is necessary is evident. Let us show that it is sufficient
by induction on $q$. Suppose that $Q = P_1\cap\dots\cap P_{q-1}$ is a
parahoric subgroup and let $A$ be an apartment containing $F(Q)$ and
$F(P_q)$. Let $\Gamma$ be the convex hull of the union of the $F(P_h)$ for
$1\leqslant h\leqslant q$ and let $L$ be a wall of $A$. Since $F(P_h)$ and
$F(P_k)$ are two facets of $A$ situated on the same side of $L$ (for
$1\leqslant h,k \leqslant q$), the set $\Gamma$ lies entirely on one side of
$L$. This being true for every wall, $\Gamma$ is therefore contained in the
closure of a facet $F(Q')$, and one has $P_h \supset Q'$ for all $h$, which
entails that the intersection of the $P_h$ is indeed a parahoric subgroup.

\num{3.5.4}\btanchor{3.5.4}
Let us now pass to the \emph{proof of theorem} (\btref{3.5.1}). We have already seen
that the condition is necessary (\btref{3.1.4}). Let us show that it is sufficient.
Let $(B,N)$ and $(B',N')$ be two Tits systems of affine type in the same
group $G$, defining the same bornology. By (\btref{3.4.6}), there exists a bijection
from the set of irreducible components of $(G,B,N)$ onto that of the
irreducible components of $(G,B',N')$ such that the bornologies defined by
two irreducible components corresponding under this bijection are the same.
Taking (\btref{2.6.4}) into account, we see that it suffices to prove our assertion
in the case where $(B,N)$ and $(B',N')$ are irreducible, that is, in view of
(\btref{3.4.6}), when the bornology they define is maximal among the bornologies
compatible with the group law of $G$ for which $G$ is not bounded. Since
$(B,N)$ and $(B',N')$ have the same maximal parahoric subgroups, lemmas
(\btref{3.5.2}) and (\btref{3.5.3}) show that they have the same parahoric subgroups, that
is, they are equivalent.

\begin{state}{Remark}{3.5.5}\btanchor{3.5.5}
\upshape
The bornology defined by a Tits system of affine type is completely
determined by the datum of the maximal parahoric subgroups, and even by the
datum of a single parahoric subgroup, since the bounded subsets are those
which are contained in the union of finitely many double cosets modulo a
given parahoric subgroup. We thus see that two Tits systems of affine type in
the same group having the same maximal parahoric subgroups, or having a
common parahoric subgroup, are equivalent.

Let us point out that one can prove that two (arbitrary) Tits systems in the
same group having the same maximal \emph{parabolic} subgroups are equivalent.
\end{state}

\begin{state}{Proposition}{3.5.6}\btanchor{3.5.6}
Let $\varphi : G\to\Gh$ be a $B$-adapted homomorphism. Let $\Gh_0$ be the
kernel of the corresponding homomorphism $\xi : \Gh \to \Aut\Cox(\bW,\bS)$
(\btref{1.2.16}). The group $\Gh_0$ is the intersection of the subgroups of finite
index of $\Gh$ which contain a maximal bounded subgroup of $\Gh$ (for the
bornology defined by $\varphi$).
\end{state}

One may suppose $G = \Gh_0$ and $\varphi = \mathrm{id}$ (\btref{1.2.17}). If $L$ is a
subgroup of finite index of $\Gh$ containing a maximal bounded subgroup, then
$L\cap G$ contains a parahoric subgroup of $G$ (\btref{3.3.4}) and consequently
$L\cap G = G$, since a parabolic subgroup of the Tits system $(G,B,N)$ is not
of finite index in $G$ unless it is equal to $G$.

Now let $g \in \Gh - G$ and let $X$ be an orbit of $\xi(g)$ in
$\Cox(\bW,\bS)$ not reduced to a point. There exists a subset
$Y \subset \bS$ which is a type of maximal parahoric subgroup of $G$ and
which is such that $X \not\subset Y$ and $X \not\subset \bS - Y$: if
$\Card(X\cap\bS_i) \leqslant 1$ for every irreducible component $\bS_i$ of
$\Cox(\bW,\bS)$ --- which implies that $(\bW,\bS)$ is not irreducible --- it
suffices to choose a point $x_i$ in each $\bS_i$ in such a way that for one
index $i_1$ one has $x_{i_1} \in X$ and for an index $i_2$ one has
$x_{i_2}\notin X$ (recall that $\Card \bS_i \geqslant 2$ for all $i$), and to
take for $Y$ the complement in $\bS$ of the set of the $x_i$; if there exists
an index $i_0$ such that $\Card(X\cap \bS_{i_0}) \geqslant 2$, it suffices to
take $Y$ such that $\bS_{i_0} - (\bS_{i_0}\cap Y)$ is reduced to a point of
$X$. Let then $\Xi_Y$ be the stabilizer of $Y$ in $\Xi = \xi(\Gh)$ and let
$L = \xi^{-1}(\Xi_Y)$. Then $L$ contains the maximal bounded subgroup
$\Stab B_Y$, does not contain $g$, and is of finite index in $\Gh$.

\num{3.5.7}\btanchor{3.5.7}
Proposition (\btref{3.5.6}) shows that the datum of the bornology of $\Gh$ determines
$\Gh_0$ and also, by (\btref{3.5.1}), the Tits system of $\Gh_0$ up to equivalence,
hence the building of $\Gh$, etc., that is to say, everything that the
$B$-adapted homomorphism $\varphi$ brings in substance to the study of $\Gh$.

%=====================================================================

\btsecA{4}{4. Iwasawa and Cartan decompositions}

%=====================================================================

We keep the notation of the preceding sections.

We denote by $\varphi : G \to \Gh$ a $B$-$N$-adapted homomorphism. Recall
that in 2.7 we defined an action of $\Gh$ on the building $\imm$, for which
the operations of $\Gh$ are isometric automorphisms of polysimplicial complex
and permute among themselves the apartments (resp. half-apartments, walls,
sectors) of $\imm$.

\btsecB{4.1}{4.1. Fixers and stabilizers}

\num{4.1.1}\btanchor{4.1.1}
Let $\Omega$ be a subset of $\imm$; we set
\begin{align*}
\Ph_{\Omega} &= \{\, g\in\Gh \mid g.x = x \text{ for all } x\in\Omega \,\}\\
\Ph^{\dgr}_{\Omega} &= \{\, g\in\Gh \mid g.\Omega = \Omega \,\}\\
P_{\Omega} &= \varphi^{-1}(\Ph_{\Omega})
   = \{\, g\in G \mid g.x = x \text{ for all } x\in\Omega \,\}\\
P^{\dgr}_{\Omega} &= \varphi^{-1}(\Ph^{\dgr}_{\Omega})
   = \{\, g\in G \mid g.\Omega = \Omega \,\}.
\end{align*}
When $\Omega$ reduces to a point $x$, we write $P_x$, etc., instead of
$P_{\{x\}}$, etc. When $\Omega = \bC$ (identified with $j(\bC)$), one has
$P_{\bC} = P^{\dgr}_{\bC} = B$, the subgroup $\Ph_{\bC}$ is none other than
the subgroup $\Bh$ introduced in (\btref{2.7.1}), and $\Ph^{\dgr}_{\bC} = \Stab B$.

The subgroups $P_{\Omega}$ and $\Ph_{\Omega}$ do not change if $\Omega$ is
replaced by its closed convex hull ((\btref{2.5.4})(iii)).

It is clear that $g\Ph_{\Omega}g^{-1} = \Ph_{g.\Omega}$, etc., for
$\Omega\subset\imm$ and $g\in\Gh$.

Since the set of fixed points of an element $g\in G$ in $\bA$ is a convex
union of closures of facets, one always has
\begin{equation}
P_{\Omega} = P_{\cl(\Omega)} . \tag{1}
\end{equation}
If $\Omega$ contains a chamber $C$, one also has
\begin{equation}
\Ph_{\Omega} = \Ph_{\cl(\Omega)} . \tag{2}
\end{equation}
Indeed, one may suppose $C = \bC$. One then has
$\Ph_{\Omega}\subset\Bh\subset\Gh_0 = \Ker\xi$, and (2) follows from (1)
applied to the Tits system $(\Gh_0, \Bh, \varphi(N))$ (\btref{1.2.17}).

\num{4.1.2}\btanchor{4.1.2}
We denote by $\Nh$ the stabilizer $\Ph^{\dgr}_{\bA}$ of the apartment $\bA$
(identified with $j(\bA)$) in $\Gh$. One has, by (\btref{2.2.5}),
$N = P^{\dgr}_{\bA} = \varphi^{-1}(\Nh)$ and $\Gh = \Nh\cdot\varphi(G)$
(\btref{2.7.2}). We denote by $\nuh$ the obvious homomorphism from $\Nh$ into the
group of automorphisms of the Euclidean affine space $\bA$, and we set
$\Wh = \nuh(\Nh)$. We write $\Vh$ for the subgroup of translations of $\Wh$.
The group $\Wh$ is contained in the normalizer $\widetilde{\bW}$ of $\bW$ in
the group of automorphisms of $\bA$. Consequently, the subgroups $\bW$, $\bV$
and $\Vh$ are normal in $\Wh$.

We set $\Hh = \Ker\nuh$; it is clear that $\Hh = \Ph_{\bA}$ and that
$H = \varphi^{-1}(\Hh)$. For every chamber $C$ contained in $\bA$, one has
$\Nh\cap\Ph_{C} = \Hh$.

\begin{state}{Definition}{4.1.3}\btanchor{4.1.3}
We shall say that the $B$-$N$-adapted homomorphism $\varphi$ is
\emph{of connected type} if the images $\vp{\bW}$ and $\vp{\Wh}$ of $\bW$ and
$\Wh$ in the group of automorphisms of $\vp{\bA}$ coincide, in other words if
$\Wh \subset \widetilde{\bW}_{\mathrm{int}}$ ((\btref{1.3.18})(6)).
\end{state}

If $\varphi$ is of connected type and if $x$ is a special point of $\bA$, one
has $\Wh_x = \bW_x$ and $\Wh = \bW_x\cdot\Vh$.

\begin{state}{Proposition}{4.1.4}\btanchor{4.1.4}
For every subset $\Omega$ contained in $\bA$, put
$\Nh^{\dgr}_{\Omega} = \Ph^{\dgr}_{\Omega}\cap\Nh$ and
$\Nh_{\Omega} = \Ph_{\Omega}\cap\Nh$. One has
\[
\Ph^{\dgr}_{\Omega} = \Nh^{\dgr}_{\Omega}\cdot\varphi(P^{\dgr}_{\Omega})
\qquad\text{and}\qquad
\Ph_{\Omega} = \Nh_{\Omega}\cdot\varphi(P_{\Omega}).
\]
If $\Omega$ contains a non-empty open set, one has
$\Ph_{\Omega} = \Hh\cdot\varphi(P_{\Omega})$.
\end{state}

Let $g\in\Ph^{\dgr}_{\Omega}$; by (\btref{2.5.8}), there exists
$g'\in P^{\dgr}_{\Omega}$ such that $g'.\bA = g.\bA$. One then has
$\varphi(g'^{-1})g \in \Nh^{\dgr}_{\Omega}$, and even
$\varphi(g'^{-1})g\in\Nh_{\Omega}$ if $g\in\Ph_{\Omega}$; whence the first
assertion. If $\Omega$ contains a non-empty open set, one has $\nuh(n) = 1$
for all $n\in\Nh_{\Omega}$, which completes the proof of the proposition.

\num{4.1.5}\btanchor{4.1.5}
If $\bD$ is a vector chamber (\btref{1.3.10}), we shall denote by $\Vh_{\bD}$ the set
of elements of $\Vh$ belonging to $\overline{\bD}$, and by $E_{\bD}$ the set
of sectors of $\bA$ of direction $\bD$ (\btref{1.3.11}). If
$\Cfr_1, \Cfr_2 \in E_{\bD}$, one also has $\Cfr_1\cap\Cfr_2 \in E_{\bD}$:
this follows immediately from the fact that $\bD$ is a simplicial cone, the
set of points with strictly positive coordinates for a suitable basis of
$\vp{\bA}$. Since
$\Ph_{\Cfr_1}\cup\Ph_{\Cfr_2} \subset \Ph_{\Cfr_1\cap\Cfr_2}$, the subgroups
$\Ph_{\Cfr}$ for $\Cfr\in E_{\bD}$ form an increasing filtered family and
their \emph{union} is a subgroup of $G$ which we shall denote by
$\Bfrh^{0}_{\bD}$.

For $g\in\nuh^{-1}(\Vh)$ and $\Cfr\in E_{\bD}$, one has
$\nuh(g)\cdot\Cfr \in E_{\bD}$ and
$g\Ph_{\Cfr}g^{-1} = \Ph_{\nuh(g).\Cfr}$. Consequently, the subgroup
$\nuh^{-1}(\Vh)$ normalizes $\Bfrh^{0}_{\bD}$ and
$\nuh^{-1}(\Vh)\cdot\Bfrh^{0}_{\bD}$ is a \emph{subgroup} of $\Gh$, which we
shall denote by $\Bfrh_{\bD}$. We set
$\Bfr^{0}_{\bD} = \varphi^{-1}(\Bfrh^{0}_{\bD})$ and
$\Bfr_{\bD} = \varphi^{-1}(\Bfrh_{\bD})$. It is immediate that
$\Bfr_{\bD} = \nu^{-1}(\bV)\cdot\Bfr^{0}_{\bD}$, that
$\Bfrh_{\bD} = \nuh^{-1}(\Vh)\cdot\varphi(\Bfr_{\bD})$ and that
$\Bfrh^{0}_{\bD} = \Hh\cdot\varphi(\Bfr^{0}_{\bD})$.

\emph{In the remainder of this section, we fix a vector chamber $\bD$ and we
shall in general omit the index $\bD$ in the notation $\Bfrh^{0}_{\bD}$,
$\Bfrh_{\bD}$, $\Bfr^{0}_{\bD}$ and $\Bfr_{\bD}$.}

\btsecB{4.2}{4.2. Double cosets}

Let $Q$ and $Q'$ be two subgroups of $\Gh$ containing $\Hh$. Put
\[
\Wh_{Q} = \nuh(\Nh\cap Q), \qquad \Wh_{Q'} = \nuh(\Nh\cap Q').
\]
Since $\Hh = \Ker\nuh \subset Q\cap Q'\cap\Nh$, the inverse image under
$\nuh$ of a double coset $\Wh_{Q}w\Wh_{Q'}$ in $\Wh$ is a double coset
$(\Nh\cap Q)n(\Nh\cap Q')$ in $\Nh$ (where $n$ is any element of
$\nuh^{-1}(w)$) and is therefore contained in a single double coset $QnQ'$ of
$G$, contained in $Q\Nh Q'$. One thus obtains a map, called canonical and
denoted $\lambda_{Q,Q'}$, from $\Wh_{Q}\backslash\Wh/\Wh_{Q'}$ into
$Q\backslash Q\Nh Q'/Q'$, a map which is obviously surjective.

\begin{state}{Proposition}{4.2.1}\btanchor{4.2.1}
Let $\Omega$ and $\Omega'$ be two subsets of $\bA$ and let $Q$ and $Q'$ be
two subgroups of $\Gh$ such that
\[
\Hh\cdot\varphi(P_{\Omega}) \subset Q \subset \Ph^{\dgr}_{\Omega}
\qquad\text{and}\qquad
\Hh\cdot\varphi(P_{\Omega'}) \subset Q' \subset \Ph^{\dgr}_{\Omega'} .
\]
The map $\lambda_{Q,Q'}$ from $\Wh_{Q}\backslash\Wh/\Wh_{Q'}$ into
$Q\backslash Q\Nh Q'/Q'$ is bijective.
\end{state}

It suffices to show that if one has $qn = n'q'$, with $n, n'\in\Nh$,
$q\in Q$ and $q'\in Q'$, then $n' \in (\Nh\cap Q)n(\Nh\cap Q')$. Now, one has
$\Omega \subset \bA\cap q(\bA)$ and
$qn(\Omega') = n'(\Omega')\subset \bA\cap q(\bA)$. Moreover, there exists an
element $g_0 \in G$ such that $q(\bA) = g_0(\bA)$ and $g_0.x = x$ for all
$x \in \bA\cap q(\bA)$ (\btref{2.5.8}). One has
$g_0 \in P_{\Omega}\cap P_{n'(\Omega')}$ and
$g = \varphi(g_0)\in Q\cap n'Q'n'^{-1}$, whence
\[
m = g^{-1}qn = g^{-1}n'q' \in Qn \cap n'Q' \cap \Nh .
\]
Consequently, $n \in (\Nh\cap Q)m$ and $n' \in m(\Nh\cap Q')$, q.e.d.

\num{4.2.2}\btanchor{4.2.2}
Let us now take for the subgroup $Q'$ one of the subgroups $\Bfrh^{0}$ or
$\Bfrh$. One checks at once that $\Wh_{\Bfrh^{0}} = \{e\}$ and
$\Wh_{\Bfrh} = \Vh$.

\begin{states}{Corollary}
Let $\Omega$ be a subset of $\bA$ and $Q$ a subgroup of $\Gh$ such that
$\Hh\cdot\varphi(P_{\Omega}) \subset Q \subset \Ph^{\dgr}_{\Omega}$. The
following canonical maps are bijective:
\begin{align*}
\text{(i)} && \Wh_{Q}\backslash\Wh &\to Q\backslash Q\Nh\Bfrh^{0}/\Bfrh^{0};\\
\text{(ii)} && \Wh_{Q}\backslash\Wh/\Vh &\to Q\backslash Q\Nh\Bfrh/\Bfrh;\\
\text{(iii)} && \Wh &\to \Bfrh^{0}\backslash\Bfrh^{0}\Nh\Bfrh^{0}/\Bfrh^{0};\\
\text{(iv)} && \Wh/\Vh &\to \Bfrh\backslash\Bfrh\Nh\Bfrh/\Bfrh .
\end{align*}
\end{states}

Here again, only the injectivity of these maps has to be proved. Let us prove
that of (i); if $q_1n_1b_1 = q_2n_2b_2$ (with $q_i\in Q$, $n_i\in\Nh$ and
$b_i\in\Bfrh^{0}$), there exists a sector $\Cfr\in E_{\bD}$ such that
$b_i \in \Ph_{\Cfr}$. One then has $Qn_1\Ph_{\Cfr} = Qn_2\Ph_{\Cfr}$. Since
$\Wh_{\Ph_{\Cfr}} = \{e\}$, one deduces from proposition (\btref{4.2.1}) that
$\Wh_{Q}\nuh(n_1) = \Wh_{Q}\nuh(n_2)$, whence the injectivity of the map (i).

Since $\Bfrh = \nuh^{-1}(\Vh)\cdot\Bfrh^{0}$, the injectivity of (ii) follows
from that of (i). One proves in the same way the injectivity of (iii) and
(iv) (taking into account, for the definition of the map (iv), that $\Vh$ is
normal in $\Wh$, whence $\Vh\backslash\Wh/\Vh = \Wh/\Vh$).

\begin{state}{Corollary}{4.2.3}\btanchor{4.2.3}
Keep the hypotheses of corollary (\btref{4.2.2}) and suppose moreover that $Q$
contains the subgroup $\Bh$. Then the following two conditions are
equivalent:
\begin{itemize}
\item[(i)] $\Bh\Nh\Bfrh \subset Q\Bfrh$;
\item[(ii)] $\Wh = \Wh_{Q}\cdot\Vh$.
\end{itemize}
\end{state}

Indeed, condition (i) amounts to $Q\Nh\Bfrh = Q\Bfrh$ and is therefore
equivalent to (ii) by corollary (\btref{4.2.2})(ii).

\num{4.2.4}\btanchor{4.2.4}
Proposition (\btref{4.2.1}) will interest us above all when $\Gh = Q\Nh Q'$. Let us
note moreover that one has
\begin{equation}
Q\Nh Q' = \varphi(P_{\Omega})\Nh\varphi(P_{\Omega'})
        = \Ph_{\Omega}\Nh\Ph_{\Omega'}
        = \Ph^{\dgr}_{\Omega}\Nh\Ph^{\dgr}_{\Omega'} . \tag{1}
\end{equation}
Indeed, let $q\in Q$; one has $\Omega\subset\bA\cap q(\bA)$ and there exists
an element $g \in P_{\Omega}$ such that $\varphi(g)q.\bA = \bA$ (\btref{2.5.8}),
whence $\varphi(g)q \in \Nh\cap Q$. Consequently, one has
$Q = \varphi(P_{\Omega})(\Nh\cap Q)$, whence the first of the equalities (1).
The other two follow, taking for example $Q = \Ph_{\Omega}$ or
$Q = \Ph^{\dgr}_{\Omega}$.

The relation $\Gh = Q\Nh Q'$ has a simple geometric interpretation: it means
that \emph{whatever $g, g'\in\Gh$, there exists an apartment containing
$g(\Omega)$ and $g'(\Omega')$}. Indeed, suppose $\Gh = Q\Nh Q'$ and put
$g^{-1}g' = qnq'$ (with $g,g'\in\Gh$, $q\in Q$, $q'\in Q'$ and $n\in\Nh$).
One then has $g'(\Omega') = gqn(\Omega')\subset gq(\bA)$ and
$g(\Omega) = gq(\Omega)\subset gq(\bA)$. Conversely, if
$g(\Omega)\cup\Omega' \subset h(\bA)$ (with $g,h\in\Gh$), there exists, by
(\btref{2.5.8}), an element $q'\in Q'$ such that $q'h(\bA) = \bA$, whence
$q'g(\Omega)\subset\bA$ and $\Omega \subset \bA\cap g^{-1}q'^{-1}(\bA)$.
Again by (\btref{2.5.8}), there exists an element $q\in Q$ such that
$qg^{-1}q'^{-1}(\bA) = \bA$, that is such that $qg^{-1}q'^{-1}\in\Nh$, whence
finally $g^{-1} \in Q\Nh Q'$.

\num{4.2.5}\btanchor{4.2.5}
In an analogous way, the relation $\Gh = \Bfrh\Nh\Bfrh$ (which is equivalent
to $\Gh = \Bfrh^{0}\Nh\Bfrh^{0}$) \emph{means that, whatever the sectors
$\Cfr$ and $\Cfr'$ of the building $\imm$, there exist sectors
$\Cfr_1\subset\Cfr$ and $\Cfr'_1\subset\Cfr'$ contained in one and the same
apartment.} Indeed, suppose first that $\Gh = \Bfrh^{0}\Nh\Bfrh^{0}$ and let
$\Cfr$ and $\Cfr'$ be two sectors of $\imm$. In view of the transitivity of
$\vp{\bW}$ on the vector chambers and in view of the very definition of the
sectors of $\imm$ (\btref{2.2.2}), there exist elements $g, g'\in\Gh$ such that
$\Cfr = g.\Cfr_0$ and $\Cfr' = g'.\Cfr'_0$, where $\Cfr_0$, $\Cfr'_0$ are
sectors of $\bA$ of direction $\bD$. Put then $g^{-1}g' = bnb'$ with
$n\in\Nh$, $b, b'\in\Bfrh^{0}$. Replacing $\Cfr$ and $\Cfr'$ by smaller
sectors, we may suppose $\Cfr'_0 = \Cfr_0$ and $b, b'\in\Ph_{\Cfr_0}$. One
then has $\Cfr' = gbnb'.\Cfr_0 \subset gb(\bA)$ and
$\Cfr = gb.\Cfr_0\subset gb(\bA)$.

Conversely, let $g\in\Gh$; let $\Cfr\in E_{\bD}$ and suppose that there exist
sectors $\Cfr_1\subset\Cfr$ and $\Cfr'_1\subset g(\Cfr)$ contained in one and
the same apartment. Replacing $\Cfr$ by a smaller sector, we may suppose
$\Cfr_1 = \Cfr$ and $\Cfr'_1 = g(\Cfr)$. Let then $h\in\Gh$ be such that
$\Cfr\cup g(\Cfr)\subset h(\bA)$. By (\btref{2.5.8}), there exists
$b\in\Ph_{\Cfr}\subset\Bfrh^{0}$ such that $bh(\bA) = \bA$, and one then has
$\Cfr \subset \bA\cap g^{-1}b^{-1}(\bA)$, whence the existence of
$b'\in\Bfrh^{0}$ such that $b'g^{-1}b^{-1}(\bA) = \bA$, which entails
$g \in \Bfrh^{0}\Nh\Bfrh^{0}$.

Entirely analogous reasoning shows that the relation $\Gh = Q\Nh\Bfrh$ (or
$\Gh = Q\Nh\Bfrh^{0}$) \emph{means that, whatever the sector $\Cfr$ of
$\imm$, there exist a sector $\Cfr'\subset\Cfr$ and an apartment containing
$\Omega$ and $\Cfr'$}. In particular, the relation $G = BN\Bfr$ \emph{means
that for every chamber $C$ and every sector $\Cfr$ of $\imm$, there exist a
sector $\Cfr'\subset\Cfr$ and an apartment containing $C$ and $\Cfr'$.}

\btsecB{4.3}{4.3. Relations between double cosets}

\begin{state}{Proposition}{4.3.1}\btanchor{4.3.1}
Let $\Omega$, $\Omega'$ and $\Omega''$ be three subsets of $\bA$. Suppose
that one of them has non-empty interior and that one of the two following
conditions is satisfied:
\begin{itemize}
\item[a)] $\Omega$, $\Omega'$ and $\Omega''$ are bounded;
\item[b)] $G$ is complete for a topology satisfying the conditions of
(\btref{2.8.2}).
\end{itemize}
Then one has
\[
\Ph_{\Omega'}\cdot\Ph_{\Omega} \,\cap\, \Ph_{\Omega''}\cdot\Ph_{\Omega}
= (\Ph_{\Omega'}\cap\Ph_{\Omega''})\cdot\Ph_{\Omega}.
\]
\end{state}

Let $g = p'p = p''q$ with $p'\in\Ph_{\Omega'}$, $p''\in\Ph_{\Omega''}$,
$p, q\in\Ph_{\Omega}$. Put $M = \Omega'\cup\Omega''\cup g.\Omega$ and let $f$
be the map from $\Omega'\cup\Omega''\cup\Omega$ onto $M$ which is the
identity on $\Omega'\cup\Omega''$ and coincides with the action of $g$ on
$\Omega$ (note that $g.x = x$ for $x\in\Omega\cap(\Omega'\cup\Omega'')$).
Then $f$ is an isometry: if $x\in\Omega$ and $y\in\Omega'$ for example, one
has $d(g.x, y) = d(p'.x, p'.y) = d(x,y)$. Since $M$ has non-empty interior,
there exists an apartment $A$ containing $M$ ((\btref{2.8.1}) and (\btref{2.8.4})) and there
exists $h\in\Ph_{\Omega'}\cap\Ph_{\Omega''}$ with $h.A = \bA$ (\btref{2.5.8}), whence
$hg.\Omega\subset\bA$, and, in view of (\btref{2.5.9}) combined with the relation
$\Gh = \varphi(G)\cdot\Nh$, there exists $n\in\Nh$ such that $hg.x = n.x$ for
all $x\in\Omega$. For $x\in\Omega$ and $y\in\Omega'\cup\Omega''$, one then
has
\[
d(n^{-1}y, x) = d(y, n.x) = d(y, hg.x) = d(y, g.x) = d(y,x).
\]
If $\Omega$ (resp. $\Omega'\cup\Omega''$) contains a chamber, one deduces
$n^{-1}.y = y$ and $n \in \Ph_{\Omega'}\cap\Ph_{\Omega''}$ (resp. $n.x = x$
and $n\in\Ph_{\Omega}$). In both cases, one has
$g \in h^{-1}n\Ph_{\Omega} \subset
(\Ph_{\Omega'}\cap\Ph_{\Omega''})\cdot\Ph_{\Omega}$.

\begin{state}{Corollary}{4.3.2}\btanchor{4.3.2}
Under the hypotheses of (\btref{4.3.1}), one has
\[
\Ph_{\Omega}\cap(\Ph_{\Omega'}\cdot\Ph_{\Omega''})
= (\Ph_{\Omega}\cap\Ph_{\Omega'})\cdot(\Ph_{\Omega}\cap\Ph_{\Omega''}).
\]
\end{state}

Indeed, if $g = p'p''\in\Ph_{\Omega}$ with $p'\in\Ph_{\Omega'}$ and
$p''\in\Ph_{\Omega''}$, proposition (\btref{4.3.1}) shows that
$p'' = p'^{-1}g \in
(\Ph_{\Omega'}\cap\Ph_{\Omega''})\cdot(\Ph_{\Omega}\cap\Ph_{\Omega''})$. One
may therefore suppose $p''\in\Ph_{\Omega}$, and one then has
$p'\in\Ph_{\Omega}$.

\num{4.3.3}\btanchor{4.3.3}
It is known that proposition (\btref{4.3.1}) and corollary (\btref{4.3.2}) are also true when
$\Omega$, $\Omega'$ and $\Omega''$ are three facets of one and the same
chamber ([\btbib{5}], chap. IV, \S~2, exercise 9).

\begin{state}{Corollary}{4.3.4}\btanchor{4.3.4}
Let $g\in\Gh$ and $x\in\bA$. Suppose that the set $Y$ of the points $y$ of
$\bA$ such that $g\in\Ph_{y}\Ph_{x}$ contains a chamber. Then $Y$ is closed.
\end{state}

It suffices to show that, if $C$ is a chamber contained in $Y$ and $y\in Y$,
one has $\cl(C\cup\{y\}) \subset Y$. Now, one has
\[
g \in \Ph_{C}\Ph_{x}\cap\Ph_{y}\Ph_{x}     = (\Ph_{C}\cap\Ph_{y})\cdot\Ph_{x}     = \Ph_{\cl(C\cup\{y\})}\cdot\Ph_{x}.
\]

\begin{state}{Corollary}{4.3.5}\btanchor{4.3.5}
Let $\Omega$, $\Omega'$ and $\Omega''$ be three non-empty subsets of $\bA$.
Suppose that $\Omega\cup\Omega'\cup\Omega''$ (resp.
$\Omega'\cup\Omega''$) contains a chamber and that $\Omega$ is contained in
the enclosure of $\Omega'\cup\Omega''$. Then one has
\[
P_{\Omega'}P_{\Omega}\cap P_{\Omega''}P_{\Omega} = P_{\Omega}
\qquad(\text{resp. }
\Ph_{\Omega'}\Ph_{\Omega}\cap\Ph_{\Omega''}\Ph_{\Omega} = \Ph_{\Omega}).
\]
\end{state}

Indeed, one can write $\Omega$, $\Omega'$ and $\Omega''$ as unions of bounded
subsets $\Omega_n$, $\Omega'_n$ and $\Omega''_n$ in such a way that
$\Omega_n\cup\Omega'_n\cup\Omega''_n$ (resp. $\Omega'_n\cup\Omega''_n$)
contains a chamber and $\Omega_n \subset \cl(\Omega'_n\cup\Omega''_n)$. One
then has, by (\btref{4.1.1})(1),
\[
P_{\Omega'}P_{\Omega}\cap P_{\Omega''}P_{\Omega}
\subset P_{\Omega'_n}P_{\Omega_n}\cap P_{\Omega''_n}P_{\Omega_n}
= P_{\cl(\Omega'_n\cup\Omega''_n)}P_{\Omega_n} = P_{\Omega_n}
\]
for all $n$, whence the first equality. The second is proved in the same way,
using (\btref{4.1.1})(2) instead of (\btref{4.1.1})(1).

\begin{state}{Corollary}{4.3.6}\btanchor{4.3.6}
Let $\Omega$ and $\Omega'$ be two non-empty convex subsets of $\bA$ and let
$n\in\Nh$. One then has
\[
\Ph_{\Omega'}n\Ph_{\Omega}\cap\Bfrh^{0}n\Ph_{\Omega} = n\Ph_{\Omega}
\]
whenever $n.\Omega\subset\cl(\Omega'+\bD)$, and in particular when
$\Omega' = \Omega$ and $\nuh(n)\in\Vh_{\bD}$.
\end{state}

Replacing $\Omega$ by $n.\Omega$, hence $\Ph_{\Omega}$ by
$n\Ph_{\Omega}n^{-1}$, we may suppose that $n = 1$. One then has
\[
\Ph_{\Omega'}\Ph_{\Omega}\cap\Bfrh^{0}\Ph_{\Omega}
= \bigcup_{\Cfr\in E_{\bD}}
  \bigl(\Ph_{\Omega'}\Ph_{\Omega}\cap\Ph_{\Cfr}\Ph_{\Omega}\bigr).
\]
But, for every sector $\Cfr\in E_{\bD}$, the convex hull of
$\Cfr\cup\Omega'$ contains $\Omega'+\bD$, and the relation
$\Omega\subset\cl(\Omega'+\bD)$ entails $\Omega\subset\cl(\Omega'\cup\Cfr)$.

\begin{state}{Corollary}{4.3.7}\btanchor{4.3.7}
Suppose $G$ complete for a topology satisfying the conditions of (\btref{2.8.2}). Let
$\Omega\subset\bA$ and $n\in\Nh_{\Omega}$. Then one has
\[
\Bfrh^{0}n\Bfrh^{0}\cap\Ph_{\Omega}
= \Ph_{\Omega+\bD}\,n\,\Ph_{\Omega+\bD}.
\]
\end{state}

It suffices to show that, for $\Cfr\in E_{\bD}$, one has
\[
\Ph_{\Cfr}n\Ph_{\Cfr}\cap\Ph_{\Omega}
\subset \Ph_{\Omega+\bD}\,n\,\Ph_{\Omega+\bD}
\]
or again that
\[
(\Ph_{\Cfr}\cdot\Ph_{n.\Cfr})\cap\Ph_{\Omega}
\subset \Ph_{\Omega+\bD}\cdot\Ph_{\Omega+n\bD}
\]
which follows from (\btref{4.3.2}).

\begin{state}{Proposition}{4.3.8}\btanchor{4.3.8}
Let $x,y,z,u$ be four points of $\bA$ such that $y$ and $z$ belong to the
segment $[xu]$ and $d(x,y)\leqslant d(x,z)$. One has
\[
\Ph_{x}\Ph_{u}\cap\Ph_{y}\Ph_{z}
= (\Ph_{x}\cap\Ph_{y})\cdot(\Ph_{z}\cap\Ph_{u}).
\]
More precisely, if $p\in\Ph_{x}$ and $p'\in\Ph_{u}$ are such that
$pp'\in\Ph_{y}\Ph_{z}$, then $p\in\Ph_{y}$ and $p'\in\Ph_{z}$.
\end{state}

Put $g = pp' = qq'$ with $q\in\Ph_{y}$ and $q'\in\Ph_{z}$. One has
$z' = g.z = q.z$ and $u' = g.u = p.u$, whence $d(y,z') = d(y,z)$ and
$d(x,u') = d(x,u)$. One deduces:
\begin{align}
d(y,u') &\leqslant d(y, g.z) + d(g.z, g.u) = d(y,z)+d(z,u) = d(y,u)\\
d(x,u') &\leqslant d(x,y)+d(y,u') \leqslant d(x,y)+d(y,u) = d(x,u) = d(x,u')\\
d(x,z') &\leqslant d(x,y)+d(y,z') = d(x,y)+d(y,z) = d(x,z) = d(x, p.z)\\
d(x,u') &\leqslant d(x,z')+d(z',u') \leqslant d(x,z)+d(z,u) = d(x,u) = d(x,u')
\end{align}
and one sees that all the inequalities appearing in (1) to (4) are in fact
equalities. From (2) one then deduces that $y \in [xu'] = [x(p.u)]$, hence,
in view of (1), that $y = p.y$, that is $p\in\Ph_{y}$. Similarly, (4) shows
that $z'\in[xu'] = [x(p.u)]$ and, in view of (3), one has $z' = p.z$, whence
$p' = p^{-1}g \in \Ph_{z}$.

\begin{state}{Corollary}{4.3.9}\btanchor{4.3.9}
Let $y, z\in\bA$ and $t\in\nuh^{-1}(\Vh)$ be such that $t.z \in y+\bD$. One
has
\[
\Bfrh^{0}_{-\bD}\,t\,\Bfrh^{0}_{\bD}\cap\Ph_{y}t\Ph_{z}
= \Ph_{y-\bD}\,t\,\Ph_{z+\bD}.
\]
\end{state}

Multiplying on the right by $t^{-1}$, one reduces to the case $t = 1$. If
$pp'\in\Ph_{y}\Ph_{z}$, with $p\in\Bfrh^{0}_{-\bD}$ and
$p'\in\Bfrh^{0}_{\bD}$, there exists a sector $\Cfr'$ (resp. $\Cfr$) of
direction $\bD$ (resp. $-\bD$) such that $p'\in\Ph_{\Cfr'}$ (resp.
$p\in\Ph_{\Cfr}$), and there exist $x\in\Cfr$ and $u\in\Cfr'$ such that the
hypotheses of (\btref{4.3.8}) are satisfied. One then has
$p \in \Ph_{y}\cap\Ph_{\Cfr} = \Ph_{y-\bD}$ and
$p' \in \Ph_{z}\cap\Ph_{\Cfr'} = \Ph_{z+\bD}$.

\begin{state}{Remark}{4.3.10}\btanchor{4.3.10}
\upshape
One can show that the preceding corollary remains valid if one supposes only
that $t.z \in y+\overline{\bD}$.
\end{state}

\begin{state}{Lemma}{4.3.11}\btanchor{4.3.11}
Suppose $\varphi : G\to\Gh$ of connected type. Let $x, y\in\bA$ be such that
$y \in x+\overline{\bD}$. Let $\rho$ be the retraction of $\imm$ onto $\bA$
having as center a chamber of $\bA$. For every $g\in\Gh$, one has
$\rho(g.y) - \rho(g.x) \dle y - x$ (cf. (\btref{1.3.16})).
\end{state}

Let $(x_0 = x, \dots, x_k = y)$ be a monotone sequence of points of the
segment $[xy]$, such that $[x_i x_{i+1}]$ is contained in the closure of a
chamber $C_i$ ($0\leqslant i < k$). There exists $h_i \in\Gh$ such that the
restriction of $\rho$ to $g.\overline{C}_i$ coincides with the action of
$h_i$, and consequently there exists $n_i\in\Nh$ such that the restriction to
$\overline{C}_i$ of the map $z\mapsto\rho(g.z)$ coincides with the action of
$n_i$. Put $w_i = \vp{\nuh}(n_i)$. By (\btref{1.3.16}), one has
\[
\rho(g.x_{i+1}) - \rho(g.x_i) = w_i(x_{i+1}-x_i) \dle x_{i+1}-x_i
\]
since $x_{i+1}-x_i \in\overline{\bD}$ and $w_i\in\vp{\bW}$. Adding these
inequalities, one obtains the lemma.

\begin{state}{Proposition}{4.3.12}\btanchor{4.3.12}
Let $x,y,z,u\in\bA$ be such that $y-x$, $z-y$ and $u-z$ belong to
$\overline{\bD}$. If $\varphi$ is of connected type, one has
\[
\Ph_{x}\Ph_{y}\cap\Ph_{u}\Ph_{z} = \Ph_{y}\cap\Ph_{z}.
\]
\end{state}

Let $g \in p\Ph_{y}\cap q\Ph_{z}$, with $p\in\Ph_{x}$ and $q\in\Ph_{u}$. By
(\btref{4.3.11}), one has, for every retraction $\rho$ of $\imm$ onto $\bA$ with
center an arbitrary chamber $C$ of $\bA$:
\begin{align}
\rho(g.y) - x &= \rho(p.y)-\rho(p.x) \dle y-x\\
\rho(g.z) - \rho(g.y) &\dle z-y\\
u - \rho(g.z) &= \rho(q.u)-\rho(q.z) \dle u-z .
\end{align}
Adding, one sees that these inequalities are in fact equalities. One
therefore has $\rho(g.y) = y$ and $\rho(g.z) = z$. Choosing $C$ in such a way
that $y\in\overline{C}$ (resp. $z\in\overline{C}$), one obtains $g.y = y$
(resp. $g.z = z$).

\begin{state}{Corollary}{4.3.13}\btanchor{4.3.13}
Let $\Omega$, $\Omega'$ and $\Omega''$ be three non-empty convex subsets of
$\bA$, such that $\Omega$ is contained in the union of the intersections
$(\Omega'+\overline{D})\cap(\Omega''-\overline{D})$, for $D$ running over the
set of vector chambers of $\vp{\bA}$. One has
\begin{itemize}
\item[(i)] $\Ph_{\Omega'}\Ph_{\Omega}\cap\Ph_{\Omega''}\Ph_{\Omega}
             = \Ph_{\Omega}$
\item[(ii)] $\Ph_{\Omega}\cap\Ph_{\Omega'}\Ph_{\Omega''}
             = (\Ph_{\Omega}\cap\Ph_{\Omega'})\cdot
               (\Ph_{\Omega}\cap\Ph_{\Omega''})$.
\end{itemize}
\end{state}

Let $g\in\Ph_{\Omega'}\Ph_{\Omega}\cap\Ph_{\Omega''}\Ph_{\Omega}$ and let
$a\in\Omega$; there exist a vector chamber $D$ and elements $b\in\Omega'$ and
$c\in\Omega''$ such that $a-b$ and $c-a$ belong to $\overline{D}$. Applying
then (\btref{4.3.12}) with $\bD$ replaced by $D$, $x$ by $b$, $y$ and $z$ by $a$ and
$u$ by $c$, one obtains $g\in\Ph_{a}$, whence (i). Relation (ii) is deduced
as in (\btref{4.3.2}).

\begin{state}{Remark}{4.3.14}\btanchor{4.3.14}
\upshape
The hypothesis of (\btref{4.3.13}) is satisfied when $\Omega$ is contained in the
convex hull of $\Omega'\cup\Omega''$.
\end{state}

\begin{state}{Proposition}{4.3.15}\btanchor{4.3.15}
Let $x,y,z,u\in\bA$, $n\in\Nh$ and $w = \nuh(n)$. If
$\Ph_{x}\Ph_{u}\cap\Ph_{z}n\Ph_{y} \neq \varnothing$, there exists
$w'\in\Wh_{z}$ such that $d(x, w'w.y) \leqslant d(x,u)+d(u,y)$. If moreover
$\varphi$ is of connected type and if
$y \in u+\overline{\bD} \subset x+\overline{\bD}$, one may choose $w'$ such
that $w'w.y \dle y$. If in addition $x \in z+\overline{\bD}$, then
$w.y \dle y$.
\end{state}

Let $g = g'g'' \in hn\Ph_{y}$, with $g'\in\Ph_{x}$, $g''\in\Ph_{u}$ and
$h\in\Ph_{z}$. Let $C$ be a chamber of $\bA$ such that $z\in\overline{C}$ and
let $\rho$ be the retraction of $\imm$ onto $\bA$ with center $C$. There
exists $h'\in\Ph_{C}$ such that $\rho(g.y) = h'g.y = h'hn.y$. On the other
hand, since $z$ and $n.y$ belong to $\bA\cap(h'h)^{-1}.\bA$, there exists
$h''\in\Gh$ with $\bA = h''(h'h)^{-1}.\bA$, $h''.z = z$ and $h''n.y = y$
(\btref{2.5.8}). Put $n' = h'hh''^{-1}$: one has $n'\in\Nh$ and
$w' = \nuh(n')\in\Wh_{z}$. Moreover,
$\rho(g.y) = h'hn.y = h'hh''^{-1}n.y = n'n.y = w'w.y$. Consequently,
\begin{align*}
d(w'w.y, x) = d(\rho(g.y), x) &\leqslant d(x, g.y) = d(x, g''.y)\\
 &\leqslant d(x,u)+d(u, g''.y) = d(x,u)+d(u,y).
\end{align*}

Suppose now $\vp{\Wh} = \vp{\bW}$ and
$y \in u+\overline{\bD}\subset x+\overline{\bD}$. In view of (\btref{4.3.11}), one
has
\begin{align*}
w'w.y - x &= \rho(g.y)-\rho(g.u)+\rho(g.u)-\rho(x)
 = \rho(g.y)-\rho(g.u)+\rho(g'.u)-\rho(g'.x)\\
 &\dle y-u+u-x = y-x .
\end{align*}
Finally, if $x\in z+\overline{\bD}$, one has, again by (\btref{4.3.11}),
\begin{align*}
wy - z &= \rho(n'^{-1}g.y)-\rho(n'^{-1}g.u)+\rho(n'^{-1}g'.u)
        -\rho(n'^{-1}g'.x)+\rho(n'^{-1}.x)-\rho(n'^{-1}.z)\\
 &\dle y-u+u-x+x-z = y-z .
\end{align*}

\begin{state}{Corollary}{4.3.16}\btanchor{4.3.16}
Let $x,y\in\bA$, let $C$ be a chamber of $\bA$ and let $n\in\Nh$. If
$\Ph_{x}\Ph_{y}\cap\Ph_{C}n\Ph_{y}\neq\varnothing$, then
$d(x, n.y)\leqslant d(x,y)$. If moreover $\varphi$ is of connected type and
$y\in x+\overline{\bD}$, one has $n.y \dle y$.
\end{state}

It suffices to apply the preceding result, taking a point $z\in C$ such that
$\Ph_{C} = \Ph_{z}$ and noting that $\Wh_{z}$ is then reduced to $\{1\}$.

\begin{state}{Corollary}{4.3.17}\btanchor{4.3.17}
Suppose $\varphi$ of connected type. Let $x\in\bA$ and $n, n'\in\Nh$ be such
that $n.x \in x+\overline{\bD}$. Let $D$ be a vector chamber. If
$\Ph_{x}n\Ph_{x}\cap\Bfrh^{0}_{D}n'\Ph_{x}\neq\varnothing$, then one has
$n'n^{-1}.x \leqslant_{D} x$.
\end{state}

Indeed, there then exists a chamber $C$ such that
\[
\Ph_{x}n\Ph_{x}n^{-1}\cap\Ph_{C}n'n^{-1}\!\cdot n\Ph_{x}n^{-1}
\neq\varnothing
\]
and it suffices to apply (\btref{4.3.16}) taking $y = n.x$.

\btsecB{4.4}{4.4. The good maximal bounded subgroups}

In this section, we shall make explicit some particular cases of (\btref{4.2.3}),
(\btref{4.3.6}) and (\btref{4.3.17}). We keep the hypotheses and notation of the preceding
numbers and we assume moreover that
\begin{equation}
G = \Bfr.N.B . \tag{1}
\end{equation}
We shall see in \S\S~5 and 7 sufficient conditions for this relation to hold.
Recall (\btref{2.9.4}) that it holds as soon as $G$ is equipped with a topology
satisfying the conditions of (\btref{2.8.2}); it will also hold in all the
applications to the case of algebraic groups. Finally, note that relation (1)
is equivalent to
\begin{equation}
\Gh = \Bfrh.\Nh.\Bh \tag{1 bis}
\end{equation}
as follows from the geometric translation of (1) or of (1 bis) given in
(\btref{4.2.5}).

\num{4.4.1}\btanchor{4.4.1}
We saw (\btref{3.3.2}) that every maximal bounded subgroup $K$ of $\Gh$ is the
stabilizer of a well-determined facet $F$ of the building $\imm$. We shall
say that $K$ is a \emph{special} maximal bounded subgroup if this facet $F$
reduces to a special vertex of $\imm$.

\begin{states}{Definition}
We say that a maximal bounded subgroup $K$ of $\Gh$ is \emph{good} if one has
$\Gh = \Bfrh.K$.
\end{states}

Let $K$ be a good maximal bounded subgroup and let $g\in\Gh$. Put
$g^{-1} = bk$ with $b\in\Bfrh$ and $k\in K$. One then has
\[
g\Bfrh g^{-1}K = g\Bfrh bkK = g\Bfrh K = \Gh .
\]
We thus see that the notion of good maximal bounded subgroup is independent
of the choice of $\Bfrh$ in its conjugacy class (i.e. of the choice of the
apartment $\bA$ and of the vector chamber $\bD$) and that \emph{every
conjugate of a good maximal bounded subgroup is a good maximal bounded
subgroup}.

\begin{state}{Proposition}{4.4.2}\btanchor{4.4.2}
Let $K$ be a bounded subgroup of $\Gh$ containing $\Bh$. The following
conditions are equivalent:
\begin{itemize}
\item[(i)] $K$ is a good maximal bounded subgroup;
\item[(ii)] one has $\Wh = \Wh_{K}\cdot\Vh$;
\item[(iii)] the canonical map from $\Wh_{K}$ into $\vp{\Wh}$ is surjective;
\item[(iv)] $\Card\Wh_{K}\geqslant\Card\vp{\Wh}$.
\end{itemize}
\end{state}

It is immediate that (ii), (iii) and (iv) are equivalent. The equivalence of
(i) and (ii) follows from (\btref{4.2.3}) and from the fact that, if (ii) is
satisfied, then $K$ is a maximal bounded subgroup of $\Gh$. To establish this
assertion, consider a bounded subgroup $K'$ containing $K$. One has
$\Wh_{K'}\supset\Wh_{K}$; since the restriction of the map
$w\mapsto \vp{w}$ to $\Wh_{K'}$ is injective, one deduces from (iii)
(equivalent to (ii)) that $\Wh_{K'} = \Wh_{K}$, whence
$K' = B.\Wh_{K'}.B = K$.

\begin{state}{Proposition}{4.4.3}\btanchor{4.4.3}
Let $K$ be a good maximal bounded subgroup of $\Gh$ containing $\Bh$. One has
\begin{equation}
\Gh = \Bfrh^{0}\cdot\Vh\cdot K
\qquad(\text{``Iwasawa decomposition of } \Gh\text{''}) \tag{1}
\end{equation}
and the canonical map from $\Vh$ into $\Bfrh^{0}\backslash\Gh/K$ is
bijective. Moreover one has
\begin{equation}
\Gh = K\cdot\Vh_{\bD}\cdot K
\qquad(\text{``Cartan decomposition of } \Gh\text{''}) \tag{2}
\end{equation}
and the canonical map from $\Vh_{\bD}$ into $K\backslash\Gh/K$ is bijective
when the homomorphism $\varphi : G\to\Gh$ is of connected type.
\end{state}

The first assertion is only a particular case of (\btref{4.2.2})(i). Relation (2)
follows from (\btref{4.2.1}), taking into account that $\Vh_{\bD}$ is a fundamental
domain for $\vp{\bW}$ acting on $\Vh$: this indeed entails that the canonical
map from $\Vh_{\bD}$ into $\Wh_{K}\backslash\Wh/\Wh_{K}$ is surjective, and
is bijective when $\vp{\Wh} = \vp{\bW}$, that is when $\varphi$ is of
connected type.

\num{4.4.4}\btanchor{4.4.4}
Let us now make explicit some of the relations of 4.3 in the case of interest
to us, which will give us relations between the Iwasawa and Cartan
decompositions of $\Gh$:

\begin{states}{Proposition}
Suppose $\varphi$ of connected type. Let $K$ be a good maximal bounded
subgroup of $\Gh$ containing $\Bh$ and let $t\in\Vh_{\bD}$, $t'\in\Vh$,
$t''\in\Vh_{\bD}$.
\begin{itemize}
\item[(i)] If $K.t.K\cap\Bfrh^{0}.t'.K \neq\varnothing$, one has
$t' \dle t$ (in other words $p(t-t')\geqslant 0$ for every dominant weight
$p$ (relative to $\bD$) of $\vp{\bSig}$).
\item[(ii)] One has $K.t.K\cap\Bfrh^{0}.t.K = t.K$.
\item[(iii)] If $t'\in\Vh_{\bD}$ and if $KtKt'K\cap Kt''K\neq\varnothing$,
then one has $t'' \dle t+t'$ (writing here additively the composition law of
$\Vh$ identified with a subgroup of $\vp{\bA}$).
\item[(iv)] If $t'\in\Vh_{\bD}$, one has
\[
t^{-1}t''^{-1}Kt''Kt\cap t'Kt'^{-1}K = t^{-1}Kt\cap K .
\]
In particular, one has
\[
t^{-1}t''^{-1}Kt''Kt\cap K = t^{-1}Kt\cap K .
\]
\end{itemize}
\end{states}

Assertion (i) follows from (\btref{4.3.17}) and (ii) from (\btref{4.3.7}). The hypothesis of
(iii) entails
$K\cap t^{-1}Kt\cdot t^{-1}t''t'^{-1}\cdot t'Kt'^{-1}\neq\varnothing$;
setting $K = \Ph_{x}$, $y = u = t.x$ and $z = t^{-1}.x$ and applying
(\btref{4.3.15}), one deduces (iii). Finally, (iv) follows from (\btref{4.3.12}), taking
$K = \Ph_{z}$, $x = (t''t)^{-1}.z$, $y = t^{-1}.z$ and $u = t'.z$.

\begin{state}{Remark}{4.4.5}\btanchor{4.4.5}
\upshape
Suppose $\varphi$ of connected type and let $G_1$ be a normal subgroup of
$\Gh$ containing $\varphi(G)$. Then $\varphi : G\to G_1$ is
$B$-$N$-adapted of connected type. Let $F$ be a facet of $\bC$ whose
stabilizer $K_1$ in $G_1$ is a good maximal bounded subgroup of $G_1$. It is
then immediate that the stabilizer $K$ of $F$ in $\Gh$ is a good maximal
bounded subgroup of $\Gh$. If $g\in K$, there exists $h\in K_1$ such that
$g.\bA = h.\bA$ and, since $\vp{\nuh}(K_1\cap\Nh) = \vp{\Wh} = \vp{\bW}$, one
sees that
\[
K = K_1\cdot(\nuh^{-1}(\Vh)\cap K).
\]
One then deduces from (\btref{4.4.3}) that
\begin{align}
\Gh &= \Bfrh^{0}\cdot\nuh^{-1}(\Vh)\cdot K_1\\
\Gh &= K_1\cdot\nuh^{-1}(\Vh)\cdot K_1 .
\end{align}

Let us now put on $\Gh$ the bornology $\borno$ for which the bounded sets are
the finite unions of translates of bounded subsets of $G_1$. Since the
bornology of $G_1$ is invariant under $\Gh$ acting on $G_1$ by inner
automorphisms, this bornology is compatible with the group law of $\Gh$. If
one supposes moreover that $\Gh/G_1$ has no elements of finite order, a
subgroup of $\Gh$ is bounded for $\borno$ if and only if it is contained and
bounded in $G_1$. Consequently, $K_1$ is a maximal bounded subgroup of $\Gh$
(for $\borno$), giving rise to an ``Iwasawa decomposition'' (1) and to a
``Cartan decomposition'' (2). We shall see later an example of this
situation: $\Gh$ will be the group of rational points of a connected
reductive algebraic group $\mathscr{G}$ defined over a local field $k$ (for a
valuation $\omega$), $G$ the group of rational points over $k$ of the derived
group of $\mathscr{G}$, $\varphi$ the natural homomorphism and $G_1$ the set
of $g\in\Gh$ such that $\omega(\chi(g)) = 0$ for every character $\chi$ of
$\mathscr{G}$ rational over $k$, the bornology $\borno$ of $\Gh$ being none
other than the natural bornology of $\Gh$.
\end{state}

\num{4.4.6}\btanchor{4.4.6}
Let us now study the existence of good maximal bounded subgroups:

\begin{states}{Proposition}
(i) If $\varphi$ is of connected type, every special maximal bounded subgroup
is a good maximal bounded subgroup.

(ii) If $\varphi$ is the identity, the only good maximal bounded subgroups
are the special maximal bounded subgroups.
\end{states}

Taking (\btref{4.4.2}) into account, the first assertion is evident, since, if $x$ is
a special point of $\bA$, one has
$\vp{\Wh}_{x} \supset \vp{\bW}_{x} = \vp{\bW} = \vp{\Wh}$; the second is
evident as well, in view of the very definition of special points.

Recall that the special points are given by the table of no.~(\btref{1.3.12}).

\num{4.4.7}\btanchor{4.4.7}
Suppose still $\varphi$ \emph{of connected type}. If the homomorphism $\xi$
from $\Gh$ into $\Aut\Cox\bW$ is not trivial, there may exist good maximal
bounded subgroups which are not special. To determine them, one reduces at
once to the irreducible case, for, $\varphi$ being of connected type,
$\xi(\Gh)$ fixes each connected component of $\Cox\bW$. In the irreducible
case, one shows easily, by counting the orders of the two groups $\Wh_{K}$
and $\vp{\Wh}$ and using (\btref{4.4.2})(iv), that the only cases of non-special good
maximal bounded subgroups are the three following ones:
\begin{itemize}
\item $G$ of type $A_1$, \quad $\xi(\Gh) = \mathbf{Z}/2\mathbf{Z}$, \quad
$K = \Stab B$;
\item $G$ of type $C_2$, \quad $\xi(\Gh) = \mathbf{Z}/2\mathbf{Z}$, \quad
$K = \Stab B_X$, where $X$ is the set of the two special points of $\Cox\bW$
(which is $\circ\!\overset{4}{-\!\!-}\!\circ\!\overset{4}{-\!\!-}\!\circ$,
the two extreme vertices being the special ones);
\item $G$ of type $B_n$ ($n\geqslant 3$), \quad
$\xi(\Gh) = \mathbf{Z}/2\mathbf{Z}$, \quad $K = \Stab B_X$, where $X$ is the
set of points of $\bS$ distinct from the non-special extremal vertex of
$\Cox\bW$ (the diagram being the one with a fork of two vertices at one end
and a bond marked $4$ at the other end).
\end{itemize}

\num{4.4.8}\btanchor{4.4.8}
On the other hand, when $\varphi$ is not of connected type, there may be no
good maximal bounded subgroups. This is the case for example when $\bW$ is of
type $A_{2n+1}$, with $\xi(\Gh) = \mathbf{Z}/2\mathbf{Z}$ acting by symmetry
without fixed points on $\Cox\bW$, or when $\bW$ is of type $D_{2n+1}$, with
$\xi(\Gh) = \mathbf{Z}/2\mathbf{Z}$ acting without fixed points on
$\Cox\bW$.

On the other hand, one shows easily that, if $\varphi$ is not of connected
type and if $\bW$ is irreducible, a good maximal bounded subgroup is always
special.

\num{4.4.9}\btanchor{4.4.9}
The preceding relations have interesting consequences when $\Gh$ is a locally
compact topological group, the subgroup $\Bh$ being an open and closed
subgroup of $\Gh$, which happens for example when $\Gh$ is a semisimple
algebraic group over a locally compact local field. The bounded subsets of
$\Gh$ are then the relatively compact subsets, and the maximal bounded
subgroups are the maximal compact subgroups. Propositions (\btref{4.4.3}) and (\btref{4.4.4})
then imply for example that, if $K$ is a good maximal compact subgroup of
$\Gh$, the convolution algebra of continuous complex functions with compact
support on $\Gh$, invariant on the left and on the right under $K$, is
\emph{commutative}: the proof given in [\btbib{7}] adapts immediately; one may also
consult [\btbib{29}], [\btbib{30}] and [\btbib{32}].

%=====================================================================

\btsecA{5}{5. Double Tits systems}

%=====================================================================

In this section, we begin by giving necessary and sufficient conditions in
order that, with the notation of \S~4, one have $G = \Bfr N\Bfr$. We shall
see that this relation is equivalent to saying that $(G,\Bfr,N)$ is a Tits
system with finite Weyl group isomorphic to $\vp{\bW}$. We shall then say
that $(G,B,N,S)$ is a \emph{double Tits system}. It is systems of this type
that we shall construct later in the group of rational points of a simply
connected semisimple algebraic group over a local field.

In 5.2, we shall give necessary and sufficient conditions in order that a
family of subgroups $P_{\alpha}$ (for $\alpha\in\bSig$) be the family of
subgroups $P_{\alpha}$ associated with a double Tits system with Weyl group
$\bW$. Along the way, we shall establish results on the structure of the
subgroups $P_{\Omega}$ of a double Tits system.

Recall that the results of this section are not used in the remainder of the
chapter.

\btsecB{5.1}{5.1. Double Tits systems}

\emph{Throughout 5.1, we keep the notation of \S~4.} In particular, we assume
chosen a vector chamber $\bD$, and the letters $\Bfr$ and $\Bfr^{0}$ denote
the subgroups introduced in (\btref{4.1.5}).

\num{5.1.1}\btanchor{5.1.1}
Since $\Bfr = \Bfr^{0}\cdot\nu^{-1}(\bV)$ and
$\Bfr^{0}\cap N = H \subset \nu^{-1}(\bV)$, one has
$\Bfr\cap N = \nu^{-1}(\bV)$. Consequently, $\Bfr\cap N$ is a normal subgroup
of $N$ and the quotient group $N/\Bfr\cap N$ is canonically identified with
$\bW/\bV$, or again with $\vp{\bW}$. This last group is generated by the set
$R$ of the reflections with respect to the walls of the vector chamber $\bD$
(that is, with respect to the hyperplanes $\Ker a$, for $a$ running over the
basis of $\vp{\bSig}$ associated with $\bD$). Consequently, \emph{the
quadruple $(G,\Bfr,N,R)$ satisfies axiom} (T 2) \emph{of Tits systems}
(\btref{1.2.6}).

\begin{states}{Definition}
We say that the Tits system of affine type $(G,B,N,S)$ is a \emph{double Tits
system} if the quadruple $(G,\Bfr,N,R)$ is a Tits system.
\end{states}

\begin{state}{Definition}{5.1.2}\btanchor{5.1.2}
We say that two closed subsets $\Omega'$ and $\Omega''$ of $\bA$ are
\emph{transverse} if, for every affine root $\alpha$ of $\bA$ containing the
intersection $\Omega'\cap\Omega''$, one has either $\alpha\supset\Omega'$, or
$\alpha\supset\Omega''$.
\end{state}

\begin{state}{Theorem}{5.1.3}\btanchor{5.1.3}
The following conditions are equivalent:
\begin{itemize}
\item[(i)] $(G,B,N,S)$ is a double Tits system;
\item[(ii)] one has $G = \Bfr.N.\Bfr$;
\item[(ii bis)] whatever the sectors $\Cfr$ and $\Cfr'$ of $\imm$, there
exist sectors $\Cfr_1\subset\Cfr$ and $\Cfr'_1\subset\Cfr'$ contained in one
and the same apartment;
\item[(iii)] whatever the transverse closed subsets $\Omega'$ and $\Omega''$
of $\bA$ such that the intersection $\Omega = \Omega'\cap\Omega''$ contains a
chamber, one has
\[ P_{\Omega} = P_{\Omega'}\cdot P_{\Omega''} ; \]
\item[(iv)] whatever the chamber $C$ of $\bA$ and the vector chamber $D$, one
has
\[ P_{C} = P_{C+D}\cdot P_{C-D} ; \]
\item[(v)] there exist a chamber $C$ of $\bA$ and a vector chamber $D$ such
that $P_{C} = P_{C+D}\cdot P_{C-D}$.
\end{itemize}
Moreover, these conditions imply:
\begin{itemize}
\item[(vi)] $G = BN\Bfr$;
\item[(vi bis)] whatever the chamber $C$ and the sector $\Cfr$ of $\imm$,
there exist a sector $\Cfr'\subset\Cfr$ and an apartment $A$ containing $C$
and $\Cfr'$;
\item[(vii)] whatever the affine roots $\alpha$ and $\beta$ of $\bA$ such
that $\beta^{*}\subsetneq\alpha$, one has
\[ P_{\alpha\cap\beta} = P_{\alpha}\cdot P_{\beta} ; \]
\item[(viii)] for every affine root $\alpha$ of $\bA$, one has
\[ P_{\partial\alpha}
   = (P_{\alpha}r_{\alpha}P_{\alpha})\cup(P_{\alpha}P_{(\alpha^{*})_{+}}). \]
\end{itemize}
\end{state}

\num{5.1.4}\btanchor{5.1.4}
The proof of theorem (\btref{5.1.3}) will occupy the following numbers, up to
(\btref{5.1.21}). We already know that (i) implies (ii) (\btref{1.2.7}) and that (ii) is
equivalent to (ii bis) (\btref{4.2.5}).

\num{5.1.5}\btanchor{5.1.5}
\emph{Proof of the implication} (ii bis) $\Rightarrow$ (iii).

Suppose condition (ii bis) realized. Let $\Omega'$ and $\Omega''$ be two
transverse closed subsets of $\bA$, such that
$\Omega = \Omega'\cap\Omega''$ contains a chamber. We wish to prove that
\begin{equation}
P_{\Omega} = P_{\Omega'}\cdot P_{\Omega''} . \tag{1}
\end{equation}
It is clear that we may suppose that $\Omega$ contains the chamber $\bC$. Let
us first prove the following lemma:

\begin{state}{Lemma}{5.1.6}\btanchor{5.1.6}
Let $M$ be a subset of $\Omega'$. Then:
\begin{itemize}
\item[a)] $\Omega'$ and $\cl(\Omega''\cup M)$ are transverse and
$\cl(\Omega\cup M) = \Omega'\cap\cl(\Omega''\cup M)$;
\item[b)] $\Omega''$ and $\cl(\Omega\cup M)$ are transverse and
$\Omega = \Omega''\cap\cl(\Omega\cup M)$.
\end{itemize}
\end{state}

Since the intersection of two closed subsets is closed, one has
$\cl(M\cup\Omega)\subset\Omega'\cap\cl(M\cup\Omega'')$. On the other hand, if
an affine root $\alpha$ contains $\Omega\cup M$, then, since $\Omega'$ and
$\Omega''$ are transverse, either $\alpha\supset\Omega'$, or
$\alpha\supset\Omega''$ and $\alpha\supset\cl(M\cup\Omega'')$. Since
$\cl(M\cup\Omega)$ is the intersection of the affine roots containing
$M\cup\Omega$ (\btref{2.4.5}), one has
$\cl(M\cup\Omega)\supset\Omega'\cap\cl(M\cup\Omega'')$, whence a). Similarly,
one has $\Omega\subset\Omega''\cap\cl(M\cup\Omega)$ and if an affine root
$\alpha$ contains $\Omega$, either $\alpha\supset\Omega''$, or
$\alpha\supset\Omega'\supset M$ and $\alpha\supset\cl(M\cup\Omega)$. One
deduces b) at once.

\begin{state}{Lemma}{5.1.7}\btanchor{5.1.7}
Suppose moreover that $\Omega'$ contains a sector of direction $D'$. One then
has $P_{\Omega}\subset P_{\cl(\Omega+D')}\cdot P_{\Omega''}$.
\end{state}
\par\endgroup

% ---------------------------------------- bruhattits5.tex
\begingroup\providecommand{\marginfont}{}
\renewcommand*{\marginfont}{\scriptsize\itshape\color{gray}}
\allowdisplaybreaks
\setlength{\parskip}{0pt}
\providecommand{\bA}{}
\renewcommand{\bA}{\mathbf{A}}
\providecommand{\bC}{}
\renewcommand{\bC}{\mathbf{C}}
\providecommand{\bD}{}
\renewcommand{\bD}{\mathbf{D}}
\providecommand{\bW}{}
\renewcommand{\bW}{\mathbf{W}}
\providecommand{\bV}{}
\renewcommand{\bV}{\mathbf{V}}
\providecommand{\bS}{}
\renewcommand{\bS}{\mathbf{S}}
\providecommand{\bR}{}
\renewcommand{\bR}{\mathbf{R}}
\providecommand{\bSi}{}
\renewcommand{\bSi}{\bm{\Sigma}}
\providecommand{\bQ}{}
\renewcommand{\bQ}{\mathbf{Q}}
\providecommand{\Cq}{}
\renewcommand{\Cq}{\mathfrak{C}}
\providecommand{\Cqb}{}
\renewcommand{\Cqb}{\overline{\mathfrak{C}}}
\providecommand{\Bg}{}
\renewcommand{\Bg}{\mathfrak{B}}
\providecommand{\Pg}{}
\renewcommand{\Pg}{\mathfrak{P}}
\providecommand{\Ug}{}
\renewcommand{\Ug}{\mathfrak{U}}
\providecommand{\Imm}{}
\renewcommand{\Imm}{\mathscr{I}}
\providecommand{\Lsc}{}
\renewcommand{\Lsc}{\mathscr{L}}
\providecommand{\vs}{}
\renewcommand{\vs}[1]{{}^{v}\mkern-1.5mu #1}
\providecommand{\hG}{}
\renewcommand{\hG}{\widehat{G}}
\providecommand{\hN}{}
\renewcommand{\hN}{\widehat{N}}
\providecommand{\hP}{}
\renewcommand{\hP}{\widehat{P}}
\providecommand{\hH}{}
\renewcommand{\hH}{\widehat{H}}
\providecommand{\hW}{}
\renewcommand{\hW}{\widehat{\mathbf{W}}}
\providecommand{\hV}{}
\renewcommand{\hV}{\widehat{\mathbf{V}}}
\providecommand{\hnu}{}
\renewcommand{\hnu}{\hat{\nu}}
\providecommand{\hBg}{}
\renewcommand{\hBg}{\widehat{\mathfrak{B}}}
\providecommand{\dg}{}
\renewcommand{\dg}{^{\dagger}}
\providecommand{\cl}{}\renewcommand{\cl}{\operatorname{cl}}
\providecommand{\Norm}{}\renewcommand{\Norm}{\operatorname{Norm}}
\providecommand{\Stab}{}\renewcommand{\Stab}{\operatorname{Stab}}
\providecommand{\Aut}{}\renewcommand{\Aut}{\operatorname{Aut}}
\providecommand{\Ot}{}
\renewcommand{\Ot}{\widetilde{\Omega}}
\providecommand{\Pt}{}
\renewcommand{\Pt}{\widetilde{P}}
\providecommand{\Nt}{}
\renewcommand{\Nt}{\widetilde{N}}
\providecommand{\Sit}{}
\renewcommand{\Sit}{\widetilde{\bm{\Sigma}}}
\providecommand{\qedfr}{}
\renewcommand{\qedfr}{\hfill\textsc{q.e.d.}}
\newenvironment{btstat}[2]
  {\par\medskip\noindent\textit{#1} \textup{\textbf{(#2)}}.\ ---\ \itshape}
  {\par\medskip}
\newenvironment{btstatu}[1]
  {\par\medskip\noindent\textit{#1}.\ ---\ \itshape}
  {\par\medskip}
\providecommand{\btpar}{}
\renewcommand{\btpar}[1]{\par\medskip\noindent\textbf{(#1)}\ }
\bigskip
\pg{85}
\noindent
Let $g\in P_{\Omega}$. Let $\Cq'$ be a sector of direction $\mathrm{D}'$
contained in $\Omega'$ and let $\Cq''$ be a sector of $\bA$ opposite to
$\Cq'$. By (ii\,bis) we may assume, after shrinking $\Cq'$ and $\Cq''$ if
necessary, that the two sectors $\Cq'$ and $g.\Cq''$ are contained in one and
the same apartment $\mathrm{A}$. Let us show that $\Omega\subset \mathrm{A}$.
It suffices to show that every chamber $\mathrm{C}_{0}\subset\Omega$ is
contained in $\mathrm{A}$. Let $x\in\Cq'$ and $y\in\Cq''$ be such that
$\mathrm{C}_{0}\subset\cl(\{x,y\})$ (\btref{2.4.9}). It follows from (\btref{2.8.8}) that
$\mathrm{C}_{0}\subset\cl(\{x,g.y\})\subset \mathrm{A}$, whence our assertion.

Let moreover $h\in G$ be such that $\bA=h.\mathrm{A}$ and $h.z=z$ for every
$z\in \mathrm{A}\cap\bA$. The restrictions of $h$ and of $\rho_{\bA;\bC}$
(resp.\ of $g^{-1}$ and of $\rho_{\bA;\bC}$) to $\mathrm{A}$ (resp.\ to
$g.\bA$) coincide, since these are isometries and they agree on $\bC$.
Consequently the restrictions of $h$ and of $g^{-1}$ to $g.\Cq''$ are equal,
and one has $hg\in P_{\Cq''}$.

On the other hand, one has $\Omega''\subset\cl(\Omega\cup\Cq'')$. Indeed, if an
affine root $\alpha$ contains $\Omega$ and $\Cq''$, it cannot contain $\Cq'$,
hence a fortiori it does not contain $\Omega'$; since $\Omega'$ and $\Omega''$
are transverse, it contains $\Omega''$, whence our assertion (\btref{2.4.5}). As $h$
and $g$ belong to $P_{\Omega}$, we conclude that
$hg\in P_{\cl(\Omega\cup\Cq'')}\subset P_{\Omega''}$.

We have thus shown that $g=h^{-1}hg$ belongs to
$P_{\cl(\Omega\cup\Cq')}\cdot P_{\Omega''}$. To finish the proof it suffices to
remark that $\Omega+\mathrm{D}'\subset\cl(\Omega\cup\Cq')$.

\begin{btstat}{Lemma}{5.1.8}\btanchor{5.1.8}
Let $\mathrm{D}_{1},\dots,\mathrm{D}_{k}$ be vector chambers such that
$\Omega'+\mathrm{D}_{i}\subset\Omega'$ for $1\leqslant i\leqslant k$. Then
\[
P_{\Omega}\subset P_{\cl(\Omega+\mathrm{D}_{1}+\dots+\mathrm{D}_{k})}\cdot P_{\Omega''}.
\]
\end{btstat}

Lemma (\btref{5.1.6}) (applied to $\mathrm{M}=\Omega+\mathrm{D}_{1}+\dots+\mathrm{D}_{i}$)
shows that, for $0\leqslant i<k$, the closed sets $\Omega'$ and
$\Omega''_{i}=\cl(\Omega''+\mathrm{D}_{1}+\dots+\mathrm{D}_{i})$ are transverse
and that their intersection is
$\cl(\Omega+\mathrm{D}_{1}+\dots+\mathrm{D}_{i})$. Lemma (\btref{5.1.7}) then shows,
since $\Omega'$ contains a sector of direction $\mathrm{D}_{i+1}$, that
\begin{align*}
P_{\cl(\Omega+\mathrm{D}_{1}+\dots+\mathrm{D}_{i})}
&\subset P_{\cl(\Omega+\mathrm{D}_{1}+\dots+\mathrm{D}_{i+1})}\cdot
   P_{\cl(\Omega''+\mathrm{D}_{1}+\dots+\mathrm{D}_{i})}\\
&\subset P_{\cl(\Omega+\mathrm{D}_{1}+\dots+\mathrm{D}_{i+1})}\cdot P_{\Omega''}
\end{align*}
whence the lemma.

\begin{btstat}{Proposition}{5.1.9}\btanchor{5.1.9}
Assume that condition \textup{(ii\,bis)} still holds. Let $\alpha$ be a
half-apartment and $\mathrm{C}$ a chamber possessing a panel $\mathrm{F}$
contained in the wall $\partial\alpha$ of $\alpha$. Then there exists an
apartment containing $\alpha$ and $\mathrm{C}$.
\end{btstat}

We may assume $\alpha\subset\bA$ and $\mathrm{C}\not\subset\bA$. Let
$\mathrm{C}'$ (resp.\ $\mathrm{C}''$) be the chamber of $\bA$ admitting
$\mathrm{F}$ as a panel and contained in $\alpha$ (resp.\ not contained in
$\alpha$). Take $\Omega'=\alpha$, $\Omega''=\overline{\mathrm{C}'\cup\mathrm{C}''}$,
whence $\Omega=\overline{\mathrm{C}'}$. One sees at once that $\Omega'$ and
$\Omega''$ are closed and transverse, since $\alpha$ is the unique affine root
containing $\mathrm{C}'$ and not $\mathrm{C}''$. Let
$\mathrm{D}_{1},\dots,\mathrm{D}_{k}$ be the vector chambers on which the root
$\vs{\alpha}$ (\btref{1.3.8}) takes negative values. One checks easily that
$\alpha=\mathrm{C}'+\mathrm{D}_{1}+\dots+\mathrm{D}_{k}$, and Lemma (\btref{5.1.8})
shows that
\begin{equation*}
P_{\mathrm{C}'}=P_{\alpha}\cdot P_{\mathrm{C}'\cup\mathrm{C}''}.\tag{1}
\end{equation*}

Now the proposition is nothing but the geometric translation of (1). Indeed,
there exists an element $g\in P_{\mathrm{C}'}$ such that
$\mathrm{C}=g.\mathrm{C}''$ (\btref{2.3.7}). Write $g=h_{1}h_{2}$, with
$h_{1}\in P_{\alpha}$ \pg{86} and $h_{2}\in P_{\mathrm{C}'\cup\mathrm{C}''}$.
One has $\mathrm{C}=g.\mathrm{C}''=h_{1}.\mathrm{C}''$, and $\mathrm{C}$, like
$\alpha$, is contained in the apartment $h_{1}.\bA$. We leave it to the reader
to deduce (1) from (\btref{5.1.9}).

\btpar{5.1.10}\btanchor{5.1.10}
One may also formulate (\btref{5.1.9}) as follows: let $\alpha$ be a half-apartment and
let $\mathrm{F}$ be a panel contained in the wall of $\alpha$; then \emph{the
group $P_{\alpha}$ is transitive on the set of chambers not contained in
$\alpha$ and admitting $\mathrm{F}$ as a panel.}

If, with the notation of (\btref{5.1.9}), one assumes moreover that
$\mathrm{C}'=\bC$, then (1) reads $B=P_{\alpha}\cdot(B\cap r_{\alpha}Br_{\alpha})$,
from which one draws $Br_{\alpha}B\subset P_{\alpha}r_{\alpha}Br_{\alpha}r_{\alpha}B$
and finally
\begin{equation*}
Br_{\alpha}B=P_{\alpha}r_{\alpha}B.\tag{1\,bis}
\end{equation*}
Conversely, (1\,bis) implies $Br_{\alpha}\subset P_{\alpha}r_{\alpha}B$, whence
$B\subset P_{\alpha}r_{\alpha}Br_{\alpha}$ and
$B=P_{\alpha}\cdot(B\cap r_{\alpha}Br_{\alpha})$, that is to say (\btref{5.1.9})\,(1).

\begin{btstat}{Corollary}{5.1.11}\btanchor{5.1.11}
Assume \textup{(ii\,bis)} holds. For every closed subset $\Omega_{1}$ of $\bA$,
the set of points of $\bA$ left fixed by all the elements of $P_{\Omega_{1}}$
reduces to $\Omega_{1}$. If $\Omega_{1}$ and $\Omega'_{1}$ are two closed
subsets of $\bA$, one has $P_{\Omega_{1}}\subset P_{\Omega'_{1}}$ if and only
if $\Omega'_{1}\subset\Omega_{1}$.
\end{btstat}

Let $\alpha\in\bSi$ and $x\in\bA-\alpha$. Let $\mathrm{F}$ be a panel contained
in $\partial\alpha$ and let $\Gamma=(\mathrm{C},\dots,\mathrm{C}')$ be a taut
gallery between $\mathrm{F}$ and $x$. One has $\mathrm{C}\subset\alpha^{*}$. If
$x$ were left fixed by all the elements of $P_{\alpha}$, the same would hold for
$\mathrm{C}$ (\btref{2.4.13}), which would contradict (\btref{5.1.10}).

Now let $\Omega_{1}$ be a closed subset of $\bA$ and let $x\in\bA-\Omega_{1}$.
By definition (\btref{2.4.6}) there exists $\alpha\in\bSi$ with $\alpha\supset\Omega_{1}$
and $x\notin\alpha$. Then $x$ is not left fixed by all the elements of
$P_{\alpha}$, hence a fortiori not by all those of $P_{\Omega_{1}}$, which
establishes the first assertion of the corollary. The second follows at once.

\medskip
One should not believe that in general $\Omega$ is the set of points of the
building $\Imm$ left fixed by all the elements of $P_{\Omega}$: this may already
fail when $\Omega$ is an affine root $\alpha\in\bSi$, or when $\Omega$ equals
the whole of $\bA$. We shall see an example of this later (the case of
$SL_{2}(\bQ_{2})$).

\begin{btstat}{Lemma}{5.1.12}\btanchor{5.1.12}
Let $\mathrm{C}$ be a chamber of $\Omega'$. Then
\[
P_{\Omega}\subset P_{\cl(\Omega\cup \mathrm{C})}\cdot P_{\Omega''}.
\]
\end{btstat}

We shall argue by induction on the length $m$ of a minimal gallery
$\Gamma=(\mathrm{C}_{0},\mathrm{C}_{1},\dots,\mathrm{C}_{m}=\mathrm{C})$ of
smallest possible length among the galleries with end $\mathrm{C}$ and origin
contained in $\Omega$. The lemma is obvious if $m=0$. Since
$\mathrm{C}_{i}\subset\Omega'$ for $0\leqslant i\leqslant m$, $\Omega'$ being
closed, the induction hypothesis implies that
$P_{\Omega}\subset P_{\cl(\Omega\cup \mathrm{C}_{m-1})}\cdot P_{\Omega''}$, and
Lemma (\btref{5.1.6}) shows that one may assume $m=1$. Let then $\alpha$ be the affine
root containing $\mathrm{C}_{0}$ and not containing $\mathrm{C}_{1}$. One has
$\Omega''\subset\alpha$; otherwise $\Omega''$ would contain a chamber
$\mathrm{C}'\subset\bA-\alpha$ and one would have
$\mathrm{C}_{1}\subset\cl(\mathrm{C}_{0}\cup \mathrm{C}')\subset\Omega''$,
whence $\mathrm{C}_{1}\subset\Omega$ and $m=0$. Let then $g\in P_{\Omega}$; the
chamber $g^{-1}.\mathrm{C}_{1}$ admits as a panel the common panel of
$\mathrm{C}_{0}$ and $\mathrm{C}_{1}$, and we may apply Proposition (\btref{5.1.9}).
There therefore exists an apartment $\mathrm{A}'$ containing $\alpha$ and
$g^{-1}.\mathrm{C}_{1}$, and there exists an element
$h\in P_{\alpha}\subset P_{\Omega''}$ such that $h.\mathrm{A}'=\bA$. One then
has $hg^{-1}\in P_{\Omega}\cap P_{\mathrm{C}_{1}}$ and
$g=gh^{-1}.h\in P_{\cl(\Omega\cup \mathrm{C}_{1})}\cdot P_{\Omega''}$, \qedfr

\pg{87}
\btpar{5.1.13}\btanchor{5.1.13}
\emph{Let us now prove \textup{(\btref{5.1.5})\,(1)} in the case where $\Omega'$
contains a sector.}

One knows (\btref{2.4.5}) that $\Omega'$ is the convex hull of finitely many points
$x_{i}$ and half-lines $\mathrm{X}_{j}$. Now one sees easily that, if
$\mathrm{X}$ is a half-line and $\mathrm{D}'$ a vector chamber, the enclosure of
$\mathrm{X}+\mathrm{D}'$ is the closure of a finite union of sectors. One
deduces that there exist chambers $\mathrm{C}_{1},\dots,\mathrm{C}_{k}$ and
vector chambers $\mathrm{D}_{1},\dots,\mathrm{D}_{k}$ such that
\[
\Omega'=\cl\Bigl(\bigcup_{1\leqslant i\leqslant k}(\mathrm{C}_{i}+\mathrm{D}_{i})\Bigr).
\]

We already know that
$P_{\Omega}\subset P_{\cl(\Omega+\mathrm{D}_{1}+\dots+\mathrm{D}_{k})}\cdot P_{\Omega''}$
(Lemma (\btref{5.1.8})). It then suffices to apply Lemma (\btref{5.1.12}) a certain number of
times in order to obtain the desired relation:
\[
P_{\Omega}\subset P_{\cl((\Omega+\mathrm{D}_{1}+\dots+\mathrm{D}_{k})\cup \mathrm{C}_{1}\cup\dots\cup \mathrm{C}_{k})}\cdot P_{\Omega''}
\subset P_{\Omega'}\cdot P_{\Omega''}.
\]

\btpar{5.1.14}\btanchor{5.1.14}
Let us now complete the proof of (\btref{5.1.5})\,(1) in the general case. Let
$\mathrm{M}$ be a segment or a closed half-line with origin $x\in\Omega$,
contained in $\Omega'$. Let $\mathrm{X}$ be the set of affine roots containing
$\Omega$ and not containing $\mathrm{M}$, and put
$\Omega'''=\bigcap_{\alpha\in \mathrm{X}}\alpha$. Every affine root
$\alpha\in \mathrm{X}$ contains in its interior the open half-line with origin
$x$ opposite to $\mathrm{M}$. One deduces at once that \emph{$\Omega'''$
contains a sector}. Moreover, an affine root $\alpha\in \mathrm{X}$ contains
$\Omega$ and does not contain $\Omega'$, hence contains $\Omega''$.
Consequently one has $\Omega'''\supset\Omega''$.

Let us show that $\Omega=\Omega'''\cap\cl(\Omega\cup \mathrm{M})$ and that
$\Omega'''$ and $\cl(\Omega\cup \mathrm{M})$ are transverse. Indeed, if an
affine root contains $\Omega$, then either it contains $\mathrm{M}$ and hence
$\cl(\Omega\cup \mathrm{M})$, or it belongs to $\mathrm{X}$ and contains
$\Omega'''$. Since $\Omega'''$ contains a sector, one deduces from (\btref{5.1.13}) that
\[
P_{\Omega}\subset P_{\cl(\Omega\cup \mathrm{M})}\cdot P_{\Omega'''}
\subset P_{\cl(\Omega\cup \mathrm{M})}\cdot P_{\Omega''}.
\]

Taking Lemma (\btref{5.1.6}) into account, we are reduced to proving (\btref{5.1.5})\,(1) with
$\Omega''$ replaced by $\cl(\Omega''\cup \mathrm{M})$. Now $\Omega'$ is the
convex hull of finitely many points $x_{i}$ (for $1\leqslant i\leqslant a$) and
closed half-lines $\mathrm{X}_{j}$ with origin $y_{j}$ (for
$1\leqslant j\leqslant b$). Let $x$ be a point of $\Omega$; put
$\mathrm{M}_{i}=[xx_{i}]$ for $1\leqslant i\leqslant a$,
$\mathrm{M}_{i}=[xy_{i-a}]$ for $a<i\leqslant a+b$, and
$\mathrm{M}_{i}=\mathrm{X}_{i-(a+b)}$ for $a+b<i\leqslant a+2b$. Applying the
preceding argument successively to the pairs $\Omega'$,
$\Omega''_{i}=\cl(\Omega\cup \mathrm{M}_{1}\cup\dots\cup \mathrm{M}_{i})$ (for
$1\leqslant i\leqslant a+2b$), everything finally reduces to proving
(\btref{5.1.5})\,(1) in the case where $\Omega'\subset\Omega''$, which is trivial.

The proof of the implication (ii\,bis)\,$\Rightarrow$\,(iii) is complete.

\btpar{5.1.15}\btanchor{5.1.15}
\emph{Proof of the implications \textup{(iii)}\,$\Rightarrow$\,\textup{(iv)} and
\textup{(iv)}\,$\Rightarrow$\,\textup{(v)}.}

It suffices to remark that, for every chamber $\mathrm{C}$ of $\bA$, one has
$\mathrm{C}=\cl(\mathrm{C}+\mathrm{D})\cap\cl(\mathrm{C}-\mathrm{D})$ and that
every affine root containing $\mathrm{C}$ contains either
$\mathrm{C}+\mathrm{D}$ or $\mathrm{C}-\mathrm{D}$, i.e.\ that
$\cl(\mathrm{C}+\mathrm{D})$ and $\cl(\mathrm{C}-\mathrm{D})$ are transverse.

\btpar{5.1.16}\btanchor{5.1.16}
\emph{Throughout the rest of \textup{5.1}, except in \textup{(\btref{5.1.18})}, we
assume that condition \textup{(v)} is satisfied, and we denote by $\mathrm{C}$
(resp.\ $\mathrm{D}$) a chamber (resp.\ a vector chamber) such that
$P_{\mathrm{C}}=P_{\mathrm{C}+\mathrm{D}}\cdot P_{\mathrm{C}-\mathrm{D}}$.}

\pg{88}
\begin{btstatu}{Lemma}
For every sector $\Cq$ of $\bA$, of direction $\mathrm{D}'$, there exists a
chamber $\mathrm{C}'\subset\Cq$ such that
$P_{\mathrm{C}'}=P_{\mathrm{C}'+\mathrm{D}'}\cdot P_{\mathrm{C}'-\mathrm{D}'}$.
\end{btstatu}

Indeed, there exists $w\in\bW$ such that $\mathrm{D}'=\vs{w}(\mathrm{D})$. On the
other hand, since $\bV$ is a lattice in $\vs{\bA}$ and $\Cq$ is an open cone of
$\bA$, there is a translation $t\in\bV$ such that $tw(\mathrm{C})\subset\Cq$.
Take then $\mathrm{C}'=tw(\mathrm{C})$: one indeed has
$P_{\mathrm{C}'}=P_{tw(\mathrm{C}+\mathrm{D})}\cdot P_{tw(\mathrm{C}-\mathrm{D})}
=P_{\mathrm{C}'+\mathrm{D}'}\cdot P_{\mathrm{C}'-\mathrm{D}'}$.

\btpar{5.1.17}\btanchor{5.1.17}
\emph{Proof of \textup{(vi)} and \textup{(vi\,bis)}.}

We have already seen (\btref{4.2.5}) that (vi) and (vi\,bis) are equivalent. Let $\Cq$
be a sector of $\bA$, of direction $\mathrm{D}'$, and let $\mathrm{C}_{1}$ be a
chamber of $\Imm$. Let $a$ be a point of $\bA$, put
$d=\sup_{x\in \mathrm{C}_{1}}d(a,x)$ and let $\mathrm{S}$ be the ball of radius
$d$ and center $a$ in $\bA$. Let $b$ be a point of $\bA$ such that
$\mathrm{S}\subset b-\mathrm{D}'$ and let $\mathrm{C}'$ be a chamber contained
in the sector $(b+\mathrm{D}')\cap\Cq$ such that
$P_{\mathrm{C}'}=P_{\mathrm{C}'+\mathrm{D}'}\cdot P_{\mathrm{C}'-\mathrm{D}'}$.
Let $\rho$ be the retraction of $\Imm$ onto $\bA$ with center $\mathrm{C}'$
(\btref{2.3.5}). Since $\rho$ is the identity on $\bA$ and decreases distances (\btref{2.5.3}),
one has
$\rho(\mathrm{C}_{1})\subset \mathrm{S}\subset b-\mathrm{D}'\subset \mathrm{C}'-\mathrm{D}'$.
On the other hand, there exists an element $g\in P_{\mathrm{C}'}$ such that
$\rho(\mathrm{C}_{1})=g(\mathrm{C}_{1})$. Write then $g=xy$, with
$x\in P_{\mathrm{C}'-\mathrm{D}'}$ and $y\in P_{\mathrm{C}'+\mathrm{D}'}$. One
has $x^{-1}g(\mathrm{C}_{1})=g(\mathrm{C}_{1})$ since
$g(\mathrm{C}_{1})\subset \mathrm{C}'-\mathrm{D}'$, and consequently
$\mathrm{C}_{1}\subset g^{-1}(\bA)=y^{-1}(\bA)$. Since
$y\in P_{\mathrm{C}'+\mathrm{D}'}$, this implies that the chamber
$\mathrm{C}_{1}$ and the sector $c+\mathrm{D}'\subset\Cq$ (where $c$ is any
point of $\mathrm{C}'$) are both contained in the apartment $y^{-1}(\bA)$. One
passes from this to the general case of (vi\,bis) by transforming $\Cq$ by an
element of $G$.

\begin{btstat}{Remark}{5.1.18}\btanchor{5.1.18}
\upshape
On examining the preceding proof one sees easily that condition (vi)
$G=BN\Bg$ is implied by the following condition, weaker than (v):

\medskip
\noindent
(v\,bis) \emph{there exist a chamber $\mathrm{C}$ and a vector chamber
$\mathrm{D}$ such that $P_{\mathrm{C}}=P_{\mathrm{C}+\mathrm{D}}\cdot P_{\Omega}$
for every bounded subset $\Omega$ of $\mathrm{C}-\mathrm{D}$.}

\medskip
Moreover, one can prove that conversely (vi) implies (v\,bis), and even implies

\medskip
\noindent
(iii\,bis) \emph{Let $\Omega'$ and $\Omega''$ be two transverse closed subsets
such that $\Omega=\Omega'\cap\Omega''$ contains a chamber. If $\Omega'$ is
bounded and if $\Omega''$ is contained in a sector, then
$P_{\Omega}=P_{\Omega'}\cdot P_{\Omega''}$.}

\medskip
Finally, one can prove that, without making any supplementary hypothesis on the
Tits system $(G,B,N,S)$, the condition (iii\,ter), obtained by adding to (iii)
or to (iii\,bis) the hypothesis that $\Omega'$ and $\Omega''$ are both
\emph{bounded}, is always satisfied. For this it suffices to prove the analogue
of (\btref{5.1.12}) when $\Omega'$ and $\Omega''$ are bounded. This is done as in
(\btref{5.1.12}), but replacing the use of Proposition (\btref{5.1.9}) by that of
Corollary (\btref{2.4.12}).
\end{btstat}

\btpar{5.1.19}\btanchor{5.1.19}
\emph{Proof of \textup{(vii)}.}

Let $\alpha$ and $\beta$ be two affine roots of $\bA$ such that
$\beta^{*}\subsetneq\alpha$. There exists $w\in\bW$ such that
$w(\mathrm{C})\subset\alpha\cap\beta$. Let $a$ be a special point of
$\overline{w(\mathrm{C})}$. There exists an element $w'$ of the stabilizer
$\bW_{a}$ of $a$ in $\bW$ such that one of the faces of the vector chamber
$\mathrm{D}'=\vs{(w'w)}(\mathrm{D})$ is parallel to $\partial\alpha$ and to
$\partial\beta$, and one may assume that $\mathrm{C}'=w'w(\mathrm{C})$ is
contained in $\alpha\cap\beta$: \pg{89} indeed, if for instance
$w'w(\mathrm{C})\not\subset\alpha$, one necessarily has $a\in\partial\alpha$,
$w'w(\mathrm{C})\subset\alpha^{*}$, whence
$r_{\alpha}w'w(\mathrm{C})\subset\alpha$. Moreover one has
$P_{\mathrm{C}'}=P_{\mathrm{C}'+\mathrm{D}'}\cdot P_{\mathrm{C}'-\mathrm{D}'}$.
On the other hand, either the affine root $\alpha$ contains
$\mathrm{C}'+\mathrm{D}'$ and $\beta\supset \mathrm{C}'-\mathrm{D}'$, or
$\alpha\supset \mathrm{C}'-\mathrm{D}'$ and $\beta\supset \mathrm{C}'+\mathrm{D}'$.
Interchanging $\alpha$ and $\beta$ if necessary, we may assume that $\alpha$
contains $\mathrm{C}'+\mathrm{D}'$ and that $\beta$ contains
$\mathrm{C}'-\mathrm{D}'$.

Let $g\in P_{\alpha\cap\beta}\subset P_{\mathrm{C}'}$ and write $g=hk$, with
$h\in P_{\mathrm{C}'+\mathrm{D}'}$ and $k\in P_{\mathrm{C}'-\mathrm{D}'}$. One
has $h=gk^{-1}\in P_{\partial\alpha\cap(\mathrm{C}'-\mathrm{D}')}$. Now the
closure of the convex hull of $\partial\alpha\cap(\mathrm{C}'-\mathrm{D}')$ and
of $\mathrm{C}'+\mathrm{D}'$ is the whole affine root $\alpha$. Consequently one
has $h\in P_{\alpha}$ and likewise $k\in P_{\beta}$, which proves (vii).

\btpar{5.1.20}\btanchor{5.1.20}
\emph{Proof of \textup{(viii)}.}

Let us resume the notation of (\btref{5.1.19}), with $\beta=(\alpha^{*})_{+}$. Let $x$
be a point of $\mathrm{C}'$ and let $\delta_{0}$ be a half-line with origin $x$,
contained in the open cone $x-\mathrm{D}'$. Since $\partial\alpha$ is parallel
to a face of $\mathrm{D}'$ and since $\alpha\supset x+\mathrm{D}'$, hence
$\alpha\not\supset x-\mathrm{D}'$, this half-line meets $\partial\alpha$ in a
point $y$, which we may assume to be interior to a panel
$\mathrm{F}\subset\partial\alpha$. Let $\delta$ be the half-line carried by
$\delta_{0}$ with origin $y$ and let $\mathrm{C}_{1}$ (resp.\ $\mathrm{C}_{2}$)
be the chamber contained in $\alpha$ (resp.\ $\alpha^{*}$) admitting
$\mathrm{F}$ as a panel. For $z\in\delta$, let
$(\mathrm{X}_{0},\dots,\mathrm{X}_{k}=\mathrm{C}_{2})$ be a taut gallery between
$z$ and $\mathrm{C}_{2}$. Since $\mathrm{C}_{1}$ meets $[yx]$, one has
$\mathrm{C}_{1}\subset\cl(\mathrm{C}_{2}\cup \mathrm{C}')$ and there exists a
minimal gallery
$(\mathrm{X}_{k}=\mathrm{C}_{2},\mathrm{X}_{k+1}=\mathrm{C}_{1},\dots,\mathrm{X}_{n}=\mathrm{C}')$
(\btref{2.4.4}). Since $\mathrm{C}_{2}\subset\cl(\{z\}\cup \mathrm{C}')$, the gallery
$\Gamma=(\mathrm{X}_{0},\dots,\mathrm{X}_{n})$ is minimal, again by (\btref{2.4.4}).

Let $g\in P_{\partial\alpha}$. Suppose first that
$g(\mathrm{C}_{1})=\mathrm{C}_{1}$. Since the enclosure of
$\mathrm{C}_{1}\cup\partial\alpha$ is the strip $\alpha\cap(\alpha^{*})_{+}$, it
follows from (vii) that $g\in P_{\alpha}P_{(\alpha^{*})_{+}}$.

Suppose now that $\mathrm{C}'_{2}=g(\mathrm{C}_{1})\neq \mathrm{C}_{1}$. Put
$\delta'=gr_{\alpha}(\delta)$ and $\Cq=gr_{\alpha}(y-\mathrm{D}')$. By (vi),
there exist a sector $\Cq'\subset\Cq$ and an apartment $\mathrm{A}$ containing
$\Cq'$ and $\mathrm{C}'$. Since $\delta'$ is contained in $\Cq$, one has
$\delta'\cap\Cq'\neq\emptyset$. Take then a point $z\in\delta$ contained in a
chamber and such that $gr_{\alpha}(z)\in\Cq'$, and resume the minimal gallery
$\Gamma$ considered above. The sequence
\[
\Gamma'=(gr_{\alpha}(\mathrm{X}_{0}),\dots,gr_{\alpha}(\mathrm{X}_{k})=gr_{\alpha}(\mathrm{C}_{2}),
\mathrm{X}_{k+1}=\mathrm{C}_{1},\dots,\mathrm{X}_{n}=\mathrm{C}')
\]
is a gallery non-stuttering, since
$gr_{\alpha}(\mathrm{C}_{2})=\mathrm{C}'_{2}\neq \mathrm{C}_{1}$, and is of the
same type as $\Gamma$. It is therefore also a minimal gallery. Since
$gr_{\alpha}(\mathrm{X}_{0})$ contains $gr_{\alpha}(z)$, hence is contained in
$\mathrm{A}$ as is $\mathrm{C}'$, one deduces that the whole of $\Gamma'$ is
contained in $\mathrm{A}$, and in particular
$\mathrm{C}_{1}\cup \mathrm{C}'_{2}\subset \mathrm{A}$. Consequently
$\mathrm{A}$ contains the convex hull of $\{y\}\cup\Cq'$, a convex hull which is
nothing but $\Cq$. In particular, $\mathrm{A}$ contains the cone
$\mathrm{Y}=(y-\overline{\mathrm{D}'})\cap\partial\alpha$.

Let then $h$ be an element of $G$ such that $h(\bA)=\mathrm{A}$ and such that
$h$ leaves fixed all the points of $\bA\cap \mathrm{A}$. One has
$h\in P_{\mathrm{C}'}\cap P_{\mathrm{Y}}\cap P_{\mathrm{C}_{1}}$ and
$h(\mathrm{C}_{2})=\mathrm{C}'_{2}$. Write $h=pq$, with
$p\in P_{\mathrm{C}'+\mathrm{D}'}$ and $q\in P_{\mathrm{C}'-\mathrm{D}'}$. Since
$\mathrm{C}'-\overline{\mathrm{D}'}$ meets $\mathrm{C}_{1}$ and
$\mathrm{C}_{2}$ and contains $\mathrm{Y}$, one has
$\mathrm{C}'_{2}=p(\mathrm{C}_{2})$ and
$p\in P_{\mathrm{C}'+\mathrm{D}'}\cap P_{\mathrm{Y}}$. But the closure of the
convex hull of $\mathrm{Y}\cup(\mathrm{C}'+\mathrm{D}')$ is equal to $\alpha$,
hence $p\in P_{\alpha}$. Finally, one has
$r_{\alpha}p^{-1}g(\mathrm{C}_{1})=\mathrm{C}_{1}$ and
$r_{\alpha}p^{-1}g\in P_{\partial\alpha}$. By what we saw above, this implies
$r_{\alpha}p^{-1}g\in P_{(\alpha^{*})_{+}}P_{\alpha}$. One deduces
\[
g\in pr_{\alpha}P_{(\alpha^{*})_{+}}P_{\alpha}
\subset P_{\alpha}r_{\alpha}P_{(\alpha^{*})_{+}}r_{\alpha}^{-1}r_{\alpha}P_{\alpha}
=P_{\alpha}r_{\alpha}P_{\alpha}
\]
(since $r_{\alpha}P_{(\alpha^{*})_{+}}r_{\alpha}^{-1}\subset P_{\alpha}$).

\pg{90}
We have indeed shown that
$P_{\partial\alpha}\subset P_{\alpha}P_{(\alpha^{*})_{+}}\cup P_{\alpha}r_{\alpha}P_{\alpha}$,
whence (viii), the opposite inclusion being obvious.

\btpar{5.1.21}\btanchor{5.1.21}
\emph{Proof of \textup{(i)}.}

In view of (\btref{5.1.16}) we may assume $\mathrm{D}=\bD$. On the other hand, we have
already seen (\btref{5.1.1}) that axiom (T\,2) of Tits systems was satisfied.

\smallskip
\emph{Proof of} (T\,1). One knows that there exists an element $n\in N$ such
that $\vs{\nu}(n)(\bD)=-\bD$. One then has $P_{\mathrm{C}+\bD}\subset\Bg$ and
$P_{\mathrm{C}-\bD}\subset n\Bg n^{-1}$. This, together with (v), shows that the
group generated by $\Bg$ and $N$ contains $P_{\mathrm{C}}$, hence contains $B$,
which proves (T\,1).

\smallskip
\emph{Proof of} (T\,3). Let us make $\bA$ into a vector space by choosing as
origin a special point $o$, and let $\Cq$ be the sector $o+\bD$. This is a
simplicial cone in $\bA$. Let $\Lsc$ be the set of hyperplanes of $\bA$ which
contain a face of $\Cq$; the images of the reflections $s_{\mathrm{L}}$ for
$\mathrm{L}\in\Lsc$ form exactly the subset $R$ of $\vs{\bW}$. Let
$\mathrm{L}\in\Lsc$ and let $r\in N$ be such that $\nu(r)=s_{\mathrm{L}}$. We
have to show that
\begin{equation*}
\Bg r\Bg w\Bg\subset(\Bg rw\Bg)\cup(\Bg w\Bg)\tag{1}
\end{equation*}
for every $w\in N$. We shall first show that
\begin{equation*}
\Bg\cup\Bg r\Bg\ \text{is a group.}\tag{2}
\end{equation*}

Let $b\in\Bg^{0}$. Since $b$ leaves fixed all the points of a sufficiently small
sector contained in $\Cq$, there exists a point $x\in r.\Cq$ such that
$rbr^{-1}\in P_{x+r.\Cqb}$. Replacing $x$ by a point of $x+r.\Cq$ if necessary,
we may assume that $x$ is the transform of a point of $\mathrm{C}$ by an element
$t$ of $\nu^{-1}(\bV)$. One then has
\[
P_{x+r.\Cqb}\subset P_{t.\mathrm{C}}=tP_{\mathrm{C}}t^{-1} =tP_{\mathrm{C}+\Cqb}t^{-1}\cdot tP_{\mathrm{C}-\Cqb}t^{-1} =P_{x+\Cqb}\cdot P_{x-\Cqb}.
\]
There therefore exist $b_{1}\in P_{x+\Cqb}\subset\Bg^{0}$ and
$c\in P_{x-\Cqb}$ such that
\begin{equation*}
rbr^{-1}=b_{1}c.\tag{3}
\end{equation*}
Since $rbr^{-1}$ and $b_{1}$ belong to $P_{x+(\Cqb\cap r.\Cqb)}$, one also has
\begin{equation*}
c\in P_{x-\Cqb}\cap P_{x+(\Cqb\cap r.\Cqb)}.\tag{4}
\end{equation*}
Now $\Cqb\cap r.\Cqb$ is a cone with nonempty interior in the hyperplane
$\mathrm{L}$, and $\mathrm{L}$ is consequently the convex hull of the union of
$\Cqb\cap r.\Cqb$ and of $-(\Cqb\cap r.\Cqb)=(-\Cqb)\cap \mathrm{L}$. One
therefore draws from (4) that
$c\in P_{x+\mathrm{L}}\cap P_{x-\Cq}$. But the closed half-space $\mathrm{E}$ of
$\bA$ with wall $x+\mathrm{L}$ and not containing $x+\Cq$ is the convex hull of
the union of $x+\mathrm{L}$ and of $x-\Cq$, and one has $c\in P_{\mathrm{E}}$.
Let $\alpha$ be the affine root with boundary $\mathrm{L}$ containing
$\mathrm{C}$. What precedes shows that there exist affine roots $\gamma$
equipollent to $\alpha$ such that
\begin{equation*}
rbr^{-1}\in\Bg^{0}\cdot P_{\gamma^{*}}.\tag{5}
\end{equation*}
It suffices indeed to choose $\gamma$ such that $\mathrm{E}\supset\gamma^{*}$.

We shall now distinguish two cases:

\smallskip
\noindent
\emph{1st case: \textup{(5)} is satisfied for every affine root $\gamma$
equipollent to $\alpha$.}

Put then, for every affine root $\gamma$ contained in $\alpha$,
\[
rbr^{-1}=b_{\gamma}c_{\gamma}\quad\text{with } b_{\gamma}\in\Bg^{0}\text{ and } c_{\gamma}\in P_{\gamma^{*}}.
\]
\pg{91}
One has $c_{\gamma}c_{\alpha}^{-1}\in P_{\alpha^{*}}\cap\Bg^{0}$. But the convex
hull of a sector of direction $\bD$ and of $\alpha^{*}$ is the whole of $\bA$;
consequently one has $P_{\alpha^{*}}\cap\Bg^{0}=H$ and
$c_{\alpha}\in HP_{\gamma^{*}}=P_{\gamma^{*}}$. Since the intersection of the
subgroups $P_{\gamma^{*}}$ for $\gamma\subset\alpha$ reduces to $H$, one draws
$c_{\alpha}\in H$ and finally $rbr^{-1}\in\Bg^{0}$.

\smallskip
\noindent
\emph{2nd case: there exists a smallest affine root $\gamma$ equipollent to
$\alpha$ such that \textup{(5)} is satisfied.}

Let us then use (viii). One has
\begin{equation*}
P_{\gamma^{*}}=r_{\gamma}P_{\gamma}r_{\gamma}
\subset(P_{\gamma}r_{\gamma}P_{\gamma})\cup(P_{\gamma}P_{(\gamma^{*})_{+}}).\tag{6}
\end{equation*}
Since $P_{\gamma}\subset\Bg^{0}$ and since
$rbr^{-1}\notin\Bg^{0}P_{(\gamma^{*})_{+}}$, one sees that
$c_{\gamma}\in P_{\gamma}r_{\gamma}P_{\gamma}$, whence
\begin{equation*}
rbr^{-1}\in\Bg^{0}P_{\gamma}r_{\gamma}P_{\gamma}
\subset\Bg^{0}r_{\gamma}\Bg^{0}=\Bg^{0}rr^{-1}r_{\gamma}\Bg^{0}\subset\Bg^{0}r\Bg
\tag{7}
\end{equation*}
since $r^{-1}r_{\gamma}\in\nu^{-1}(\bV)\subset\Bg$.

In both cases we have therefore shown that
$rbr^{-1}\in\Bg^{0}\cup\Bg^{0}r\Bg$. One thus has
$r\Bg^{0}r^{-1}\subset\Bg^{0}\cup\Bg^{0}r\Bg$ and
\[
r\Bg r=r\Bg^{0}\nu^{-1}(\bV)r=r\Bg^{0}r^{-1}r\nu^{-1}(\bV)r
\subset(\Bg^{0}\cup\Bg^{0}r\Bg)\nu^{-1}(\bV)=\Bg\cup\Bg r\Bg
\]
which implies at once that $\Bg\cup\Bg r\Bg$ is a group.

Let now $w\in N$ and suppose first that \emph{the length of $\vs{\nu}(rw)$
(with respect to the generating system $R$ of $\vs{\bW}$) is greater than that
of $\vs{\nu}(w)$}. By ([\btbib{5}], chap.~V, \S\,3, n\textsuperscript{o}~2, th.~1), this
means that the vector chambers $\bD$ and $\vs{\nu}(rw).\bD$ lie on either side
of the hyperplane of $\vs{\bA}$ parallel to $\mathrm{L}$; consequently, for
every affine root $\gamma$ equipollent to $\alpha$, there exists a sector
$\Cq'\subset\Cq$ such that $rw.\Cq'\subset\gamma^{*}$, whence
$P_{\gamma^{*}}\subset P_{rw.\Cq'}$ and
\[
(rw)^{-1}P_{\gamma^{*}}rw\subset P_{\Cq'}\subset\Bg^{0}.
\]
Let then $b\in\Bg^{0}$. By (5), there exist an affine root $\gamma$ equipollent
to $\alpha$, an element $b_{1}\in\Bg$ and an element $c\in P_{\gamma^{*}}$ such
that $rbw=rbr^{-1}rw=b_{1}crw$. One deduces
\[
rbw=b_{1}rw(rw)^{-1}crw\in\Bg rw\cdot(rw)^{-1}P_{\gamma^{*}}rw\subset\Bg rw\Bg .
\]
Since $\nu^{-1}(\bV)$ is normal in $N$, it follows that
\begin{equation*}
\Bg r\Bg w\Bg\subset\Bg rw\Bg\quad\text{when }
\ell(\vs{\nu}(rw))>\ell(\vs{\nu}(w)).\tag{8}
\end{equation*}

Suppose now that the length of $\vs{\nu}(rw)$ is smaller than that of
$\vs{\nu}(w)$. Substituting $rw$ for $w$ in (8), one draws
$\Bg r\Bg rw\Bg\subset\Bg w\Bg$, whence, using (2),
\begin{equation*}
\Bg r\Bg w\Bg\subset\Bg r\Bg r^{-1}\Bg rw\Bg
\subset(\Bg\cup\Bg r\Bg)\Bg rw\Bg=\Bg rw\Bg\cup\Bg r\Bg rw\Bg
\subset\Bg rw\Bg\cup\Bg w\Bg.\tag{9}
\end{equation*}

We have thus indeed proved (1), that is to say (T\,3).

\medskip
The quadruple $(G,\Bg,N,R)$ satisfying axioms (T\,1), (T\,2) and (T\,3) of Tits
systems, one knows ([\btbib{5}], chap.~IV, \S\,2, n\textsuperscript{o}~3, remark) that
$G=\Bg N\Bg$, that is to say that (ii) is satisfied. It follows that
\emph{condition} (iv) \emph{is satisfied}.

\smallskip
Let us finally prove axiom (T\,4):
\[
r\Bg r^{-1}\neq\Bg\qquad(\text{for } r=s_{\mathrm{L}}\text{ with }\mathrm{L}\in\Lsc).
\]
Suppose there exists a hyperplane $\mathrm{L}\in\Lsc$ such that
$r\Bg r^{-1}=\Bg$ (with $r=s_{\mathrm{L}}$). If $\Cq=x+\bD$ is a sector of
direction $\bD$, one has
\[
P_{\Cq}\subset\Bg^{0}\subset r\Bg r^{-1}=\nu^{-1}(\bV)r\Bg^{0}r^{-1}.
\]
\pg{92}
For every element $g\in P_{\Cq}$ there therefore exist a sector $\Cq'$ of
direction $\mathrm{D}'=\vs{\nu}(r).\bD$, an element $h\in P_{\Cq'}$ and an
element $t\in\nu^{-1}(\bV)$ such that $g=th$. Let $\mathrm{E}$ (resp.\
$\mathrm{E}'$) be the enclosure of $\Cq\cup\Cq'$ (resp.\ $\Cq\cup t.\Cq'$); it
is a convex subset of $\bA$ and one has $g.\mathrm{E}=\mathrm{E}'$. Moreover,
the map $y\mapsto g.y$ from $\mathrm{E}$ onto $\mathrm{E}'$ is affine, reduces
to the identity on $\Cq$ and to the translation $y\mapsto t.y$ on $\Cq'$. One
concludes that $\nu(t)=1$ and that $g\in P_{\Cq'}$. In other words,
\emph{every element which leaves fixed the points of a sector of direction
$\bD$ also leaves fixed the points of a sector of direction $\mathrm{D}'$.}

Let then $\mathrm{C}'$ be the chamber contained in $\Cq$ to which $o$ is
adherent. What precedes implies that
$P_{\mathrm{C}'+\bD}\subset P_{\mathrm{C}'+\mathrm{D}'}$. Using (iv), one draws
\begin{equation*}
P_{\mathrm{C}'}=P_{\mathrm{C}'+\bD}\cdot P_{\mathrm{C}'-\bD}
\subset P_{\mathrm{C}'+\mathrm{D}'}\cdot P_{\mathrm{C}'-\bD}
\subset P_{(\mathrm{C}'+\mathrm{D}')\cap(\mathrm{C}'-\bD)}.\tag{10}
\end{equation*}
Now the only affine root which contains $\mathrm{C}'$ and does not contain
$r(\mathrm{C}')$ is the affine root $\alpha$ with boundary $\mathrm{L}$
containing $\mathrm{C}'$. It follows that the chamber $r(\mathrm{C}')$ is
contained in the intersection
$(\mathrm{C}'+\mathrm{D}')\cap(\mathrm{C}'-\bD)$, and one deduces from (10)
that $P_{\mathrm{C}'}\subset P_{r(\mathrm{C}')}$, which is impossible since the
stabilizers of two distinct chambers are two distinct parahoric subgroups. We
have thus reached a contradiction and axiom (T\,4) is satisfied.

This completes the proof of (i) and consequently of Theorem (\btref{5.1.3}).

\begin{btstat}{Remark}{5.1.22}\btanchor{5.1.22}
\upshape
Condition (viii) implies the following condition:

\medskip
\noindent
(viii\,bis) \emph{The union of two half-apartments with the same wall having no
chamber in common is an apartment.}

\medskip
Indeed, suppose (viii) satisfied and let $\mathrm{M}$ and $\mathrm{M}'$ be two
half-apartments with the same wall, having no chamber in common. We may
obviously restrict ourselves to the case where $\mathrm{M}'$ is an affine root
$\alpha$ of $\bA$. There then exists an element $g\in P_{\partial\alpha}$ such
that $g(\bA)$ is an apartment containing $\mathrm{M}$ (\btref{2.5.8}). Replacing $g$ by
$gr_{\alpha}$ if necessary, we may assume that $g(\alpha)=\mathrm{M}$. Since
$\mathrm{M}\cap\alpha$ contains no chamber, one has
$g\notin P_{\alpha}\cdot P_{(\alpha^{*})_{+}}$, whence, by (viii),
$g=hr_{\alpha}h'$, with $h,h'\in P_{\alpha}$. Then
$\mathrm{M}=hr_{\alpha}(\alpha)=h(\alpha^{*})$ and
$\mathrm{M}\cup\alpha=h(\bA)$, which proves (viii\,bis).

Conversely, one can show that condition (viii\,bis) implies (viii).
\end{btstat}

\btpar{5.1.23}\btanchor{5.1.23}
\emph{Until the end of \textup{5.1} we assume that $(G,B,N,S)$ is a double Tits
system.} For every affine root $\alpha$, the subgroup $P_{\partial\alpha}$ is
then generated by $P_{\alpha}$ and $P_{\alpha^{*}}$. Indeed, by (\btref{5.1.10}) there
exists an element $u\in P_{\alpha^{*}}$ which does not leave fixed a chamber
$\mathrm{C}$ contained in $\alpha$ and bordered by $\partial\alpha$. One then
has $u\notin P_{\alpha}\cdot P_{(\alpha^{*})_{+}}$, whence
$P_{\alpha}uP_{\alpha}=P_{\alpha}r_{\alpha}P_{\alpha}$ by (viii).

On the other hand, \emph{the set of points of $\Imm$ left fixed by every element
of $P_{\partial\alpha}$ reduces to $\partial\alpha$.} Indeed, let
$x\in\Imm$ not belonging to $\partial\alpha$ and let $y$ be a point of a panel
$\mathrm{F}$ contained in $\partial\alpha$. If $x$ is invariant under
$P_{\partial\alpha}$, the segment $[xy]$ is likewise left fixed pointwise by
$P_{\partial\alpha}$. Now there is a chamber $\mathrm{C}$ admitting $\mathrm{F}$
as a panel and containing a point of $[xy]$. This chamber would then be
invariant under $P_{\partial\alpha}$, which contradicts (\btref{5.1.10}).

\begin{btstat}{Lemma}{5.1.24}\btanchor{5.1.24}
Let $\alpha$ and $\beta$ be two affine roots of $\bA$. In order that
$\alpha^{*}=\beta$ it is necessary and sufficient that the group generated by
$P_{\alpha}$ and $P_{\beta}$ leave fixed no chamber, but leave fixed a panel of
$\Imm$.
\end{btstat}

\pg{93}
That the condition is necessary follows from (\btref{5.1.23}). Conversely, if the affine
root $\beta$ is not equipollent to $\alpha^{*}$, the group generated by
$P_{\alpha}$ and $P_{\beta}$ leaves fixed a chamber, since $\alpha\cap\beta$ has
nonempty interior. The same holds when $\beta\supsetneq\alpha^{*}$. If
$\beta\subset\alpha^{*}$, the fixed points of the group generated by
$P_{\alpha}$ and $P_{\beta}$ are contained in the set of fixed points of the
group generated by $P_{\alpha}$ and $P_{\alpha^{*}}$ (resp.\ $P_{\beta}$ and
$P_{\beta^{*}}$), that is to say, by (\btref{5.1.23}), in $\partial\alpha$ (resp.\
$\partial\beta$). If moreover $\beta\neq\alpha^{*}$, one has
$\partial\alpha\cap\partial\beta=\emptyset$, which proves that the condition is
sufficient.

\btpar{5.1.25}\btanchor{5.1.25}
It follows from (\btref{5.1.24}) that the datum of the set of subgroups $P_{\alpha}$,
for $\alpha\in\bSi$, allows one to determine the pairs
$(P_{\alpha},P_{\alpha^{*}})$, hence also the walls $\partial\alpha$ (\btref{5.1.23}),
and consequently the apartment $\bA$, which is the convex hull of the union of
the $\partial\alpha$. Consequently:

\begin{btstatu}{Proposition}
\upshape
\begin{itemize}\itshape
\item[(i)] In order that an element $g$ of $G$ belong to $N$, it is necessary
and sufficient that the inner automorphism defined by $g$ permute among
themselves the subgroups $P_{\alpha}$ for $\alpha\in\bSi$.
\item[(ii)] One has
$\displaystyle H=\bigcap_{\alpha\in\bSi}P_{\alpha}=\bigcap_{\alpha\in\bSi}\Norm P_{\alpha}$.
\end{itemize}
\end{btstatu}

\btpar{5.1.26}\btanchor{5.1.26}
For every affine root $\alpha$ of $\bA$, let us give ourselves a subset
$U_{\alpha}$ of $P_{\alpha}$ such that $P_{\alpha}=H.U_{\alpha}$. One then sees
easily, using (\btref{5.1.10}), that for every chamber $\mathrm{C}$ bordered by
$\partial\alpha$ and not contained in $\alpha$, there exists an element
$u\in U_{\alpha}$ such that $u(\mathrm{C})\neq \mathrm{C}$. One deduces, as in
(\btref{5.1.23}) and (\btref{5.1.24}), that the set of points of $\Imm$ left fixed by
$U_{\alpha}\cup U_{\alpha^{*}}$ reduces to $\partial\alpha$, and that, if
$\alpha$ and $\beta$ are two affine roots, one has $\beta=\alpha^{*}$ if and
only if $U_{\alpha}\cup U_{\beta}$ leaves fixed a panel and leaves fixed no
chamber of $\Imm$. One may then generalize Proposition (\btref{5.1.25}) as follows:
\emph{in order that an element $g\in G$ belong to the stabilizer $N$ of the
apartment $\bA$, it is necessary and sufficient that for every affine root
$\alpha$ of $\bA$ there exist an affine root $\beta$ of $\bA$ such that
$P_{\beta}=HgU_{\alpha}g^{-1}$.}

\btpar{5.1.27}\btanchor{5.1.27}
Let $\Omega$ be a nonempty closed subset of $\bA$ and let $\mathrm{L}$ be the
affine subspace of $\bA$ generated by $\Omega$. Since $\Omega$ is a convex union
of facets, one may choose a facet $\mathrm{F}\subset\Omega$ which is
\emph{open} in $\mathrm{L}$. Let $\Delta$ be a connected component of the
complement $\mathrm{X}$ in $\bA$ of the union of the walls of $\bA$ containing
$\mathrm{L}$. The orthogonal symmetry with respect to $\mathrm{L}$ transforms
$\Delta$ into another connected component of $\mathrm{X}$, denoted $\Delta'$.
One denotes by $\mathrm{C}_{0}$ (resp.\ $\mathrm{C}'_{0}$) the unique chamber
contained in $\Delta$ (resp.\ $\Delta'$) and containing $\mathrm{F}$ in its
closure.

\begin{btstatu}{Lemma}
Let $\mathrm{C}$ be a chamber of $\Imm$ containing $\mathrm{F}$ in its closure.
There exists an apartment containing $\mathrm{C}$ and $\Delta$.
\end{btstatu}

Let $\Gamma=(\mathrm{C}_{0},\mathrm{C}_{1},\dots,\mathrm{C}_{m}=\mathrm{C})$ be a
taut gallery between $\mathrm{C}_{0}$ and $\mathrm{C}$. Let us first show, by
induction on $i$, that $\mathrm{F}\subset\overline{\mathrm{C}}_{i}$ for
$0\leqslant i\leqslant m$. This is true by hypothesis for $i=0$. Let $\mathrm{A}$
be an apartment containing $\mathrm{C}_{0}$ and $\mathrm{C}$, hence $\Gamma$.
Suppose $i>0$ and let $\mathrm{M}$ be the wall of $\mathrm{A}$ separating
$\mathrm{C}_{i-1}$ and $\mathrm{C}_{i}$. Since $\mathrm{M}$ separates
$\mathrm{C}_{0}$ and $\mathrm{C}$ (\btref{2.3.10}) and since
$\mathrm{F}\subset\overline{\mathrm{C}}_{0}\cap\overline{\mathrm{C}}$, one has
$\mathrm{F}\subset \mathrm{M}$. By the induction hypothesis,
$\mathrm{F}\subset\overline{\mathrm{C}}_{i-1}$, hence
$\mathrm{F}\subset \mathrm{M}\cap\overline{\mathrm{C}}_{i-1}
=\overline{\mathrm{C}}_{i-1}\cap\overline{\mathrm{C}}_{i}$ and
$\mathrm{F}\subset\overline{\mathrm{C}}_{i}$.

Let us now prove the lemma by induction on $m$. If $m=0$, there is nothing to
prove. If $m\geqslant1$, there exists by the induction hypothesis an apartment
$\mathrm{A}$ con-\pg{94}taining $\Delta$ and $\mathrm{C}_{m-1}$. By (\btref{2.5.8}),
there exists $g\in P_{\Delta}$ with $g.\mathrm{A}=\bA$. Replacing $\mathrm{C}$
by $g.\mathrm{C}$, one sees that one may assume $\mathrm{C}_{m-1}\subset\bA$.
Let then $\alpha$ be the affine root containing $\mathrm{C}_{m-1}$ and whose
boundary $\partial\alpha$ contains
$\overline{\mathrm{C}}_{m-1}\cap\overline{\mathrm{C}}$. Then $\partial\alpha$
separates $\mathrm{C}_{0}$ and $\mathrm{C}$ (\btref{2.3.10}) and one has
$\mathrm{C}_{0}\subset\alpha$. Since $\partial\alpha\supset \mathrm{F}$, one has
$\Delta\subset\alpha$. But (\btref{5.1.10}) implies that there exists an apartment
containing $\mathrm{C}$ and $\alpha$, whence the lemma.

\begin{btstat}{Proposition}{5.1.28}\btanchor{5.1.28}
Let us keep the notation of \textup{(\btref{5.1.27})}. Put $N_{\Omega}=P_{\Omega}\cap N$
and let $\Ot$ be the intersection of the affine roots $\alpha$ such that
$\alpha\supset\Omega$ and $\Omega\not\subset\alpha^{*}$. Then $\Ot$ is a closed
subset with nonempty interior and one has
\[
P_{\Omega}=P_{\Delta}N_{\Omega}P_{\Delta}P_{\Ot}=P_{\mathrm{L}}\cdot P_{\Ot}.
\]
\end{btstat}

The affine roots $\alpha$ satisfying $\Omega\subset\alpha$ and
$\Omega\not\subset\alpha^{*}$ are those satisfying $\Omega\subset\alpha$ and
$\Omega\not\subset\partial\alpha$, or again $\Omega\subset\alpha$ and
$\mathrm{F}\not\subset\partial\alpha$. Let us show that
\begin{equation*}
\Ot=\cl(\Omega\cup \mathrm{C}_{0}\cup \mathrm{C}'_{0}).\tag{1}
\end{equation*}
Indeed, if $\alpha\supset\Omega$ and $\partial\alpha\not\supset \mathrm{F}$, the
facet $\mathrm{F}$ is interior to $\alpha$; since
$\mathrm{F}\subset\overline{\mathrm{C}}_{0}\cap\overline{\mathrm{C}}'_{0}$, one
concludes that $\mathrm{C}_{0}\cup \mathrm{C}'_{0}\subset\alpha$. Conversely, if
$\alpha\supset\Omega\cup \mathrm{C}_{0}\cup \mathrm{C}'_{0}$, then
$\partial\alpha\not\supset \mathrm{F}$, since every wall containing $\mathrm{F}$
separates $\mathrm{C}_{0}$ and $\mathrm{C}'_{0}$. Whence (1). It is then clear
that $\Ot$ is a closed subset with nonempty interior.

Let now $g\in P_{\Omega}$. Put $\mathrm{C}=g.\mathrm{C}_{0}$. By Lemma (\btref{5.1.27}),
there exists an apartment $\mathrm{A}$ containing $\mathrm{C}$ and $\Delta$, and
consequently an element $u\in P_{\Delta}$ such that $u.\mathrm{A}=\bA$ (\btref{2.5.8}),
whence $ug.\mathrm{C}_{0}\subset\bA$. There then exists $n\in N$ with
$nug.\mathrm{C}_{0}=\mathrm{C}_{0}$, and $n$ is the identity on
$\overline{\mathrm{C}}_{0}\cap ug.\overline{\mathrm{C}}_{0}\supset \mathrm{F}$.
Hence $n$ is the identity on $\mathrm{L}$ and $n\in N_{\Omega}$.

Put $\mathrm{C}'=nug.\mathrm{C}'_{0}$. Again by Lemma (\btref{5.1.27}), there exists an
apartment $\mathrm{A}'$ containing $\Delta$ and $\mathrm{C}'$, whence an element
$u'\in P_{\Delta}$ such that $g'.\mathrm{C}'_{0}\subset\bA$, on putting
$g'=u'nug$. Since $g'.\mathrm{C}_{0}=\mathrm{C}_{0}$, one has
$w(\mathrm{C}_{0},\mathrm{C}'_{0})=w(\mathrm{C}_{0},g'.\mathrm{C}'_{0})$ (with
the notation of (\btref{2.1.7})), whence $g'.\mathrm{C}'_{0}=\mathrm{C}'_{0}$ by
(\btref{2.3.8}). Hence $g'\in P_{\Omega\cup \mathrm{C}_{0}\cup \mathrm{C}'_{0}}$ and
$g'\in P_{\Ot}$ by (1). Finally one has indeed
$g=u^{-1}n^{-1}u'^{-1}g'\in P_{\Delta}N_{\Omega}P_{\Delta}P_{\Ot}$, whence the
proposition.

\begin{btstat}{Corollary}{5.1.29}\btanchor{5.1.29}
Let $\varphi:G\to\hG$ be a $B$-$N$-adapted homomorphism. Let us resume the
notation of \textup{(\btref{4.1.1})} and put $\hN_{\Omega}=\hN\cap\hP_{\Omega}$ and
$\hN\dg_{\Omega}=\hN\cap\hP\dg_{\Omega}$. One has
\begin{align*}
&a)\qquad \hP_{\Omega}=(\hN\dg_{\Delta}\cap\hN_{\Omega})\varphi(P_{\Omega})
=\varphi(P_{\Delta})\hN_{\Omega}\varphi(P_{\Delta}P_{\Ot})
=\hP_{\Delta}\hN_{\Omega}\hP_{\Delta}\hP_{\Ot};\\
&b)\qquad \hP\dg_{\Omega}=(\hN\dg_{\Delta}\cap\hN\dg_{\Omega})\varphi(P_{\Omega})
=\varphi(P_{\Delta})\hN\dg_{\Omega}\varphi(P_{\Delta}P_{\Ot})
=\hP_{\Delta}\hN\dg_{\Omega}\hP_{\Delta}\hP_{\Ot}.
\end{align*}
\end{btstat}

Let $g\in\hP\dg_{\Omega}$. By (\btref{2.5.8}), there exists $h\in P_{\Omega}$ with
$h.g^{-1}.\bA=\bA$, whence
$\varphi(h)g^{-1}\in\hN\cap\hP\dg_{\Omega}=\hN\dg_{\Omega}$. Hence
$\hP\dg_{\Omega}=\hN\dg_{\Omega}\varphi(P_{\Omega})$. But, if
$n\in\hN_{\Omega}$, $n.\Delta$ is a connected component of $\mathrm{X}$ and
there exists $n'\in N_{\Omega}$ such that $n'.n.\Delta=\Delta$, whence
$\hN\dg_{\Omega}=(\hN\dg_{\Delta}\cap\hN\dg_{\Omega})\varphi(N_{\Omega})$. This
proves the first of the equalities $b)$. On the other hand, if
$n\in\hN\dg_{\Delta}$, one has $n\hP_{\Delta}n^{-1}=\hP_{\Delta}$, whence
$n\varphi(P_{\Delta})n^{-1}=\varphi(P_{\Delta})$ since $\varphi(G)$ is normal in
$\hG$ and since $\varphi(P_{\Delta})=\varphi(G)\cap\hP_{\Delta}$. One deduces
the other equalities $b)$, and $a)$ follows at once from $b)$.

\pg{95}
\begin{btstat}{Remark}{5.1.30}\btanchor{5.1.30}
\upshape
It is clear that $\hN_{\Omega}=\hN_{\mathrm{L}}$. On the other hand, one may
have $\hN\dg_{\Omega}\neq\hN\dg_{\mathrm{L}}$. Moreover, one has
\begin{equation*}
\hN\dg_{\Delta}\cap\hN_{\Omega}
=\hnu^{-1}(\hW\dg_{\mathrm{C}_{0}}\cap\hW_{\mathrm{F}})\tag{1}
\end{equation*}
since $\mathrm{C}_{0}$ is the only chamber contained in $\Delta$ and admitting
$\mathrm{F}$ as a facet. One concludes that \emph{$\hH\varphi(P_{\Omega})$ is a
normal subgroup of finite index of $\hP_{\Omega}$}, the quotient
$\hP_{\Omega}/\hH\varphi(P_{\Omega})$ being isomorphic to
$\hW\dg_{\mathrm{C}_{0}}\cap\hW_{\mathrm{F}}$.
\end{btstat}

\begin{btstat}{Proposition}{5.1.31}\btanchor{5.1.31}
Every $B$-$N$-adapted homomorphism is also $\Bg$-$N$-adapted.
\end{btstat}

Let $\varphi:G\to\hG$ be a $B$-$N$-adapted homomorphism. One knows that
$\hG=\varphi(G).\hN$. Since $\hN$ normalizes $\varphi(N)$ and
$\varphi(\nu^{-1}(\bV))$, it suffices to show that, for every $n\in\hN$, there
exists $n'\in N$ such that $\varphi(n')n$ normalizes $\varphi(\Bg^{0})$. One may
write $w=\hnu(n)=w'w''t$, with $w'\in\bW_{x}$,
$w''\in\widetilde{\bW}\dg_{x+\bD}$ and $t\in\widetilde{\bV}$ (where $x$ is a
special point of $\bA$). Let $n'\in\nu^{-1}(w')$. For every sector $\Cq$ of
direction $\bD$ of $\bA$, one then has
\[
n\varphi(P_{\Cq})n^{-1}=\varphi(P_{w'w''t(\Cq)})=\varphi(n'P_{\Cq'}n'^{-1})
\]
where $\Cq'=w''t(\Cq)$ is a sector of $\bA$, of direction $\bD$. One therefore
has $n\varphi(\Bg^{0})n^{-1}=\varphi(n'\Bg^{0}n'^{-1})$, whence the proposition.

\medskip
Let moreover $\vs{\omega}$ be the homomorphism from $\hG$ into
$\Aut(\vs{\bW},R)$ associated with $\varphi$ (regarded as a $\Bg$-$N$-adapted
homomorphism). One knows that $\vs{\omega}$ is trivial on $\varphi(G)$ and we
have just shown that $\hG=\varphi(G).\mathrm{E}$, with
$\mathrm{E}=\{g\in\hN\mid \vs{\hnu}(g).\bD=\bD\}$. If $g\in \mathrm{E}$ and
$n\in N$, one has (\btref{1.2.20})
\[
\vs{\omega}(g)(\vs{\nu}(n))=\vs{\nu}(g\varphi(n)g^{-1})
=\vs{\hnu}(g)\,\vs{\nu}(n)\,\vs{\hnu}(g)^{-1}
\]
which shows that $\vs{\omega}$ restricted to $\mathrm{E}$ is identified with the
composite of $\vs{\hnu}$ and the canonical isomorphism of
$\vs{\hW}\dg_{\bD}$ onto $\Aut(\vs{\bW},R)$ ((\btref{1.3.18})\,(4)). It follows that
\emph{the kernel $\vs{\hG}_{0}$ of $\vs{\omega}$ is equal to
$\varphi(G)\hnu^{-1}(\hV)$}. Moreover, since $\hnu^{-1}(\hV)$ normalizes
$\varphi(\Bg)$ and since $\varphi(\Bg)$ is its own normalizer in $\varphi(G)$,
one deduces that
$\Stab\Bg\cap\vs{\hG}_{0}=\varphi(\Bg)\cdot\hnu^{-1}(\hV)=\hBg$, which shows
that the notation $\hBg$ is coherent with (\btref{1.2.17}).

\begin{btstat}{Corollary}{5.1.32}\btanchor{5.1.32}
If $\varphi:G\to\hG$ is a $B$-$N$-adapted homomorphism of connected type
\textup{(\btref{4.1.3})}, then $\hG=\vs{\hG}_{0}$, one has
$\hN\cap\hBg=\hnu^{-1}(\hV)$, the quotient group $\hN/\hBg\cap\hN$ is
identified with $\vs{\bW}$, and the quadruple $(\hG,\hBg,\hN,R)$ is a Tits
system with Weyl group $\vs{\bW}$.
\end{btstat}

Since $\vs{\hW}=\vs{\bW}$, one has $\mathrm{E}=\hnu^{-1}(\hV)$, whence
$\hG=\vs{\hG}_{0}$, and the corollary follows from (\btref{1.2.17}).

\btpar{5.1.33}\btanchor{5.1.33}
One knows ([\btbib{5}], [\btbib{39}]) that to every Tits system one may associate an abstract
simplicial complex, which we shall here call the \emph{combinatorial building}
of the system, whose simplices represent the parabolic subgroups, the face
relation corresponding to the inverse of the inclusion relation. In what follows
we shall briefly indicate the relations between the building $\Imm$ and the
combinatorial building of the Tits \pg{96}system $(G,\Bg,N,R)$ with Weyl group
$\vs{\bW}$, which we shall denote by $\vs{\Imm}$. By definition, the chambers of
$\vs{\Imm}$ are in bijective correspondence with the minimal parabolic
subgroups, that is to say with the conjugates of $\Bg$, this correspondence
being compatible with the action of $G$ on $\vs{\Imm}$ on the one hand and on
the set of conjugates of $\Bg$ (by inner automorphisms) on the other.

Moreover, the group $\Bg$ is the stabilizer of the sector germ
$\widetilde{\Cq}_{\Bg}$ (\btref{2.9.2}) containing the sectors $x+\bD$ for $x\in\bA$,
and it is easy to see that $G$ operates transitively on the set of sector germs
of $\Imm$, the stabilizer of $g.\widetilde{\Cq}_{\Bg}$ being $g\Bg g^{-1}$.
Since $\Bg$ is its own normalizer, one deduces a bijection compatible with the
action of $G$ between the set of sector germs of $\Imm$ and the set of minimal
parabolic subgroups, or again the set of chambers of $\vs{\Imm}$: the chamber
$f(\widetilde{\Cq})$ associated with the sector germ $\widetilde{\Cq}$ is the
one whose stabilizer coincides with that of $\widetilde{\Cq}$.

Moreover, \emph{sector germs $\widetilde{\Cq}_{i}$ $(1\leqslant i\leqslant m)$
are contained in one and the same apartment of $\Imm$ if and only if the
chambers $f(\widetilde{\Cq}_{i})$ are contained in one and the same apartment of
$\vs{\Imm}$} (recall that an apartment of $\vs{\Imm}$ is the transform by an
element of $G$ of the subcomplex $\vs{\bA}$ of $\vs{\Imm}$ whose simplices are
those associated with the parabolic subgroups $n\Pg n^{-1}$ for
$\Pg\supset\Bg$ and $n\in N$). To show this, it suffices to see that
$\widetilde{\Cq}\subset\bA$ is equivalent to $f(\widetilde{\Cq})\subset\vs{\bA}$,
which is obvious, since $\widetilde{\Cq}\subset\bA$ is equivalent to the
existence of an $n\in N$ with $\widetilde{\Cq}=n.\widetilde{\Cq}_{\Bg}$. One
deduces at once that \emph{there exists one and only one bijection, again
denoted $f$, from the set of apartments of $\Imm$ onto the set of apartments of
$\vs{\Imm}$, such that a sector germ $\widetilde{\Cq}$ is contained in an
apartment $\mathrm{A}$ if and only if the chamber $f(\widetilde{\Cq})$ is
contained in the apartment $f(\mathrm{A})$.} It is clear that this bijection is
compatible with the action of $G$ on $\Imm$ and on $\vs{\Imm}$.

It follows in particular from what precedes that to every automorphism $\theta$
of the building $\Imm$ there corresponds one and only one automorphism
$v(\theta)$ of $\vs{\Imm}$ such that
$v(\theta).f(\widetilde{\Cq})=f(\theta.\widetilde{\Cq})$ for every sector germ
$\widetilde{\Cq}$ of $\Imm$. One thus obtains a homomorphism $v$ from
$\Aut\Imm$ into $\Aut\vs{\Imm}$, a homomorphism which is \emph{injective}.
Indeed, if $v(\theta)=1$, then $\theta$ preserves all the apartments of $\Imm$.
Now one sees easily that each facet $\mathrm{F}$ of $\Imm$ has as closure the
intersection of the apartments of $\Imm$ which contain it; one deduces that
$\theta.\mathrm{F}=\mathrm{F}$ for every facet $\mathrm{F}$ of $\Imm$, whence
$\theta=1$. This makes it possible to use, for the study of $\Aut\Imm$, the
results established in [\btbib{39}] on the automorphism group of a combinatorial
building associated with a Tits system with finite Weyl group.

\btsecB{5.2}{5.2. Axiomatic characterization of the family of the \texorpdfstring{$P_{\Omega}$}{P\_Omega}}

In n\textsuperscript{os} (\btref{5.2.1}) to (\btref{5.2.27}) we abandon the hypotheses and
notation adopted in \S\S\,2 to 4, retaining only those of (\btref{1.5.2}). We denote by
$G$ a group, by $N$ a subgroup of $G$, by $\nu$ a surjective homomorphism from
$N$ onto $\bW$, and by $H$ the kernel of $\nu$. We put $W=N/H$ and we denote by
$S$ the inverse image of $\bS$ in $W$ under the isomorphism of $W$ onto $\bW$
deduced from $\nu$ (an isomorphism by means of which one will \pg{97}sometimes
identify $W$ and $\bW$). One is given, for every affine root $\alpha\in\bSi$, a
subgroup $P_{\alpha}$ of $G$. For every subset $\Omega$ of $\bA$, one denotes by
$P_{\Omega}$ the subgroup of $G$ generated by $H$ and the $P_{\alpha}$ for
$\alpha\in\bSi$ and $\alpha\supset\Omega$. It is clear that
\[
P_{\Omega}=P_{\cl(\Omega)}.
\]

We shall impose on these data conditions ((\btref{5.2.1})\,(1) to (6)) which will imply
that $(G,P_{\mathrm{C}},N,S)$ is a double Tits system. Later (\btref{5.2.27}) we shall
show that, conversely, the subgroups $P_{\alpha}$ associated as in \S\,4 with a
double Tits system satisfy these conditions, the notations $P_{\Omega}$ being
concordant. This will make it possible to apply to the case of double Tits
systems the results on the structure of the groups $P_{\Omega}$ which we shall
prove in (\btref{5.2.6}) et seq.

\btpar{5.2.1}\btanchor{5.2.1}
We assume, up to (\btref{5.2.27}), that the preceding data satisfy the following six
conditions:

\medskip
\begin{itemize}\itshape
\item[(A\,1)] For every $n\in N$ and every $\alpha\in\bSi$, one has
$nP_{\alpha}n^{-1}=P_{\nu(n).\alpha}$.
\item[(A\,2)] Let $\alpha\in\bSi$. For every $\beta\in\bSi$ with
$\beta\supsetneq\alpha$, one has $P_{\beta}\subsetneq P_{\alpha}$, and the
intersection of the $P_{\beta}$ for $\beta\supset\alpha$ is equal to $H$.
\item[(A\,3)] Let $\alpha,\beta\in\bSi$ be such that $\alpha\cap\beta$ has
nonempty interior, and let $\Omega$ be the intersection of the affine roots
containing $\alpha\cap\beta$ and not containing $\alpha$. Then
$P_{\alpha}P_{\Omega}$ is a group.
\item[(A\,4)] For every $\alpha\in\bSi$, the subset
$P_{\alpha}P_{(\alpha^{*})_{+}}\cup P_{\alpha}\nu^{-1}(r_{\alpha})P_{\alpha}$
is a subgroup of $G$.
\item[(A\,5)] For every $x\in\bA$ and every vector chamber $\mathrm{D}$, one has
$P_{x+\mathrm{D}}\cap P_{x-\mathrm{D}}=H$.
\item[(A\,6)] The group $G$ is generated by the union of the $P_{\alpha}$ for
$\alpha\in\bSi$.
\end{itemize}

\btpar{5.2.2}\btanchor{5.2.2}
It is clear that (A\,1) implies
\begin{equation*}
nP_{\Omega}n^{-1}=P_{\nu(n).\Omega}\qquad\text{for } n\in N\text{ and }\Omega\subset\bA.
\tag{A\,1\,bis}
\end{equation*}
Using then the transitivity of $\vs{\bW}$ on the vector chambers, one sees at
once that, granting (A\,1), it suffices to assume (A\,5) for \emph{one} vector
chamber $\mathrm{D}$. Moreover, (A\,2) shows that for $\Omega=\alpha$ the
notation $P_{\Omega}$ is unambiguous.

On the other hand, let us place ourselves in the hypotheses of (A\,3). Since
every affine root containing $\alpha\cap\beta$ contains either $\alpha$ or
$\Omega$, one deduces at once from (A\,2) and (A\,3) that one has
\begin{equation*}
P_{\alpha\cap\beta}=P_{\alpha}P_{\Omega}.\tag{1}
\end{equation*}
Suppose moreover that $\beta\supsetneq\alpha^{*}$. Then one has $\Omega=\beta$ and
\begin{equation*}
P_{\alpha\cap\beta}=P_{\alpha}P_{\beta}.\tag{2}
\end{equation*}
In particular, one has
\begin{equation*}
P_{\alpha\cap(\alpha^{*})_{+}}=P_{\alpha}P_{(\alpha^{*})_{+}}.\tag{3}
\end{equation*}

\btpar{5.2.3}\btanchor{5.2.3}
By (A\,2) one has $H\subset P_{\alpha}$ for every $\alpha\in\bSi$. In what
follows one denotes by $U_{\alpha}$ a system of representatives of $P_{\alpha}$
modulo $H$, i.e.\ a subset of $P_{\alpha}$ such that \pg{98}the map
$(u,h)\mapsto uh$ from $U_{\alpha}\times H$ into $P_{\alpha}$ is bijective. One
assumes the $U_{\alpha}$ chosen in such a way that $U_{\alpha}\supset U_{\beta}$
for $\alpha,\beta\in\bSi$ with $\alpha\subset\beta$. In the applications we have
in view, the $U_{\alpha}$ will be subgroups of $P_{\alpha}$. One also puts
$U_{\bA}=\{1\}$ and $\Ug_{a}=\bigcup_{\vs{\alpha}=a}U_{\alpha}$ for
$a\in\vs{\bSi}$.

\btpar{5.2.4}\btanchor{5.2.4}
Let $\Omega$ be a nonempty subset of $\bA$. One denotes by $\Sit(\Omega)$ the
set of affine roots containing $\Omega$, and by $\bSi(\Omega)$ the set of
minimal elements of $\Sit(\Omega)$, ordered by inclusion. Since
$\Omega\neq\emptyset$, every element of $\Sit(\Omega)$ contains one and only one
element of $\bSi(\Omega)$. The map $\alpha\mapsto\vs{\alpha}$ restricted to
$\bSi(\Omega)$ is evidently injective. Let $\vs{\bSi}(\Omega)$ be its image; for
$a\in\vs{\bSi}(\Omega)$ one will write $a(\Omega)$ for the unique element of
$\bSi(\Omega)$ satisfying $\vs{a(\Omega)}=a$. If
$a\in\vs{\bSi}-\vs{\bSi}(\Omega)$, one puts $a(\Omega)=\bA$. One has
\[
\cl(\Omega)=\bigcap_{\alpha\in\Sit(\Omega)}\alpha
=\bigcap_{\alpha\in\bSi(\Omega)}\alpha =\bigcap_{a\in\vs{\bSi}(\Omega)}a(\Omega) =\bigcap_{a\in\vs{\bSi}}a(\Omega).
\]

\btpar{5.2.5}\btanchor{5.2.5}
Let $\Omega$ be a subset of $\bA$ containing a sector of direction $\mathrm{D}$.
Then one has $\vs{\bSi}(\Omega)\subset\vs{\bSi}^{+}(\mathrm{D})$. Indeed, for
every $\alpha\in\bSi$ there exists $x\in\bA$ such that
$\alpha=\{x+t\mid t\in\vs{\bA},\ \vs{\alpha}(t)\geqslant0\}$ (\btref{1.3.8}). If
$\vs{\alpha}\notin\vs{\bSi}^{+}(\mathrm{D})$ and $z\in \mathrm{D}$, one has
$\vs{\alpha}(z)<0$. Let then $y\in\Omega$ and $\mu\in\bR_{+}$ be such that
$\mu\vs{\alpha}(z)<\vs{\alpha}(x-y)$. One has
$y+\mu z\in\Omega+\mathrm{D}\subset\cl(\Omega)$ and
$\vs{\alpha}(y+\mu z-x)<0$, whence $y+\mu z\notin\alpha$ and
$\alpha\notin\Sit(\Omega)$.

It follows (\btref{1.3.15}) that \emph{one may order $\vs{\bSi}(\Omega)$ according to a
nibbling order}.

\begin{btstat}{Proposition}{5.2.6}\btanchor{5.2.6}
Let $\Omega$ be a subset of $\bA$ containing a sector and let
$(a_{1},\dots,a_{n})$ be the elements of $\vs{\bSi}(\Omega)$ arranged according
to a nibbling order. The product map
$\pi:(u_{1},\dots,u_{n},h)\mapsto u_{1}\dots u_{n}h$ is a bijection of
$\prod_{1\leqslant j\leqslant n}U_{a_{j}(\Omega)}\times H$ onto $P_{\Omega}$.
\end{btstat}

We shall prove this proposition by induction on $n$, the assertion being obvious
for $n=0$ and $n=1$. Suppose $n>1$ and put
$\Omega'=\bigcap_{2\leqslant j\leqslant n}a_{j}(\Omega)$, whence
$\cl(\Omega)=a_{1}(\Omega)\cap\Omega'$. Since $a_{1}$ is extremal in
$\{a_{1},\dots,a_{n}\}$, the half-space
$\{t\in\vs{\bA}\mid a_{1}(t)<0\}$ does not contain the intersection of the
half-spaces $\{t\in\vs{\bA}\mid a_{j}(t)<0\}$ for $2\leqslant j\leqslant n$, and
one concludes immediately that no affine root equipollent to $a_{1}(\Omega)$
contains $\Omega'$, which implies that $\vs{\bSi}(\Omega')=\{a_{2},\dots,a_{n}\}$.
Since $(a_{2},\dots,a_{n})$ is a nibbling order, one has by the induction
hypothesis
\[
P_{\Omega'}=U_{a_{2}(\Omega)}\dots U_{a_{n}(\Omega)}H.
\]

Let us now show that $Q=U_{a_{1}(\Omega)}P_{\Omega'}$ is equal to $P_{\Omega}$.
For this it suffices to show that $Q$ is stable under left multiplication by a
generating system of $P_{\Omega}$. Now it is obvious that $HQ\subset Q$ and that
$U_{a_{1}(\Omega)}Q\subset Q$, since
$HU_{a_{1}(\Omega)}\cup U^{2}_{a_{1}(\Omega)}\subset P_{a_{1}(\Omega)}=U_{a_{1}(\Omega)}H$
and since $H\subset P_{\Omega'}$. If $j\geqslant2$, one has by (\btref{5.2.2})\,(1)
\[
P_{a_{1}(\Omega)\cap a_{j}(\Omega)}=P_{a_{1}(\Omega)}P_{\Omega_{j}} =U_{a_{1}(\Omega)}P_{\Omega_{j}}
\]
where $\Omega_{j}$ denotes the intersection of the affine roots containing
$a_{1}(\Omega)\cap a_{j}(\Omega)$ and not containing $a_{1}(\Omega)$. But these
affine roots contain $\Omega$ and do not contain $a_{1}(\Omega)$, hence they
contain $\Omega'$. One therefore has $\Omega_{j}\supset\Omega'$ and
$P_{\Omega_{j}}\subset P_{\Omega'}$. Whence \pg{99}
\begin{align*}
U_{a_{j}(\Omega)}Q&=U_{a_{j}(\Omega)}U_{a_{1}(\Omega)}P_{\Omega'}
\subset P_{a_{1}(\Omega)\cap a_{j}(\Omega)}P_{\Omega'}
=U_{a_{1}(\Omega)}P_{\Omega_{j}}P_{\Omega'}\\
&\subset U_{a_{1}(\Omega)}P_{\Omega'}=Q.
\end{align*}
We have thus shown that $\pi$ is \emph{surjective}.

Let us now show that $\pi$ is \emph{injective}. Let $h,h'\in H$ and
$u_{j},u'_{j}\in U_{a_{j}(\Omega)}$ be such that
$u_{1}\dots u_{n}h=u'_{1}\dots u'_{n}h'$. Since $a_{1}$ is extremal in
$\{a_{1},\dots,a_{n}\}$, there exists a vector chamber $\mathrm{D}$ such that
$a_{1}\in\vs{\bSi}^{+}(\mathrm{D})$ and $a_{j}\in-\vs{\bSi}^{+}(\mathrm{D})$ for
$2\leqslant j\leqslant n$. If $x$ is a point of $\Omega$, one then has
$a_{1}(\Omega)\supset x-\mathrm{D}$ and $\Omega'\supset x+\mathrm{D}$, whence
\[
u'^{-1}_{1}u_{1}\in P_{a_{1}(\Omega)}\cap P_{\Omega'}
\subset P_{x-\mathrm{D}}\cap P_{x+\mathrm{D}}=H
\]
and $u_{1}=u'_{1}$. The induction hypothesis then implies $u_{j}=u'_{j}$ for
every $j$, whence the injectivity of $\pi$.

\btpar{5.2.7}\btanchor{5.2.7}
One may express Proposition (\btref{5.2.6}) in a slightly different form. Let $\Omega$
be a subset of $\bA$ containing a sector of direction $\mathrm{D}$. Let us
arrange the elements of $\vs{\bSi}^{+}(\mathrm{D})$ in a nibbling order
$(a_{1},\dots,a_{m})$. Then \emph{the product map
$(u_{1},\dots,u_{m},h)\mapsto u_{1}\dots u_{m}h$ is a bijection of
$\prod_{1\leqslant j\leqslant m}U_{a_{j}(\Omega)}\times H$ onto $P_{\Omega}$}:
this follows immediately from the fact that the order induced on
$\vs{\bSi}(\Omega)$ is nibbling (\btref{1.3.15}), and from the definitions of
$a_{j}(\Omega)$ (\btref{5.2.4}) and of $U_{\bA}$ (\btref{5.2.3}).

\btpar{5.2.8}\btanchor{5.2.8}
For every $\alpha\in\bSi$, the product map $(h,u)\mapsto hu$ is a bijection of
$H\times U_{\alpha}^{-1}$ onto $P_{\alpha}$. Consequently one deduces from
(\btref{5.2.6}) the following assertion, which generalizes (\btref{5.2.6}): with the notation of
(\btref{5.2.6}) (or of (\btref{5.2.7})), the product map
$(u_{1},\dots,u_{k},h,u_{k+1},\dots,u_{n})\mapsto u_{1}\dots u_{k}hu_{k+1}\dots u_{n}$
is a bijection of
$\prod_{1\leqslant j\leqslant k}U_{a_{j}(\Omega)}\times H\times
\prod_{k+1\leqslant j\leqslant n}U^{-1}_{a_{j}(\Omega)}$ onto $P_{\Omega}$, for
any $k$ with $0\leqslant k\leqslant n$.

\begin{btstat}{Proposition}{5.2.9}\btanchor{5.2.9}
Let $\Omega$ be a nonempty subset of $\bA$ and let $\mathrm{D}$ be a vector
chamber. Let us arrange the elements of
$\vs{\bSi}(\Omega)\cap\vs{\bSi}^{+}(\mathrm{D})$ (resp.\
$\vs{\bSi}(\Omega)\cap\vs{\bSi}^{+}(-\mathrm{D})$) in a nibbling order
$(a_{1},\dots,a_{n})$ (resp.\ $(b_{1},\dots,b_{m})$). The product map
$(u_{1},\dots,u_{n},h,v_{1},\dots,v_{m})\mapsto u_{1}\dots u_{n}hv_{1}\dots v_{m}$
from $\prod_{1\leqslant j\leqslant n}U_{a_{j}(\Omega)}\times H\times
\prod_{1\leqslant k\leqslant m}U^{-1}_{b_{k}(\Omega)}$ into $P_{\Omega}$ is
injective.
\end{btstat}

One has $\cl(\Omega+\mathrm{D})=\bigcap_{1\leqslant j\leqslant n}a_{j}(\Omega)$
and $\cl(\Omega-\mathrm{D})=\bigcap_{1\leqslant k\leqslant m}b_{k}(\Omega)$,
whence
\[
U_{a_{1}(\Omega)}\dots U_{a_{n}(\Omega)}H\subset P_{\Omega+\mathrm{D}}
\qquad\text{and}\qquad
U^{-1}_{b_{1}(\Omega)}\dots U^{-1}_{b_{m}(\Omega)}\subset P_{\Omega-\mathrm{D}}.
\]
If, with obvious notation,
\[
u_{1}\dots u_{n}hv_{1}\dots v_{m}=u'_{1}\dots u'_{n}h'v'_{1}\dots v'_{m},
\]
then
\[
g=(u'_{1}\dots u'_{n})^{-1}u_{1}\dots u_{n}h =h'v'_{1}\dots v'_{m}(v_{1}\dots v_{m})^{-1}
\in P_{\Omega+\mathrm{D}}\cap P_{\Omega-\mathrm{D}}.
\]
But, if $x\in\Omega$, one has
$P_{\Omega+\mathrm{D}}\cap P_{\Omega-\mathrm{D}}\subset
P_{x+\mathrm{D}}\cap P_{x-\mathrm{D}}=H$, whence $g\in H$. It then suffices to
apply (\btref{5.2.6}) and (\btref{5.2.8}) to $\cl(\Omega+\mathrm{D})$ and to
$\cl(\Omega-\mathrm{D})$ in order to draw $u_{j}=u'_{j}$ and $v_{k}=v'_{k}$ for
$1\leqslant j\leqslant n$ and $1\leqslant k\leqslant m$.

\begin{btstat}{Lemma}{5.2.10}\btanchor{5.2.10}
Let $\Omega$ be a subset of $\bA$ with nonempty interior. For every vector
chamber $\mathrm{D}$, put
$Q_{\mathrm{D}}=P_{\Omega+\mathrm{D}}\cdot P_{\Omega-\mathrm{D}}$. The subset
$Q_{\mathrm{D}}$ is independent of the choice of $\mathrm{D}$.
\end{btstat}

\pg{100}
It suffices to prove that if $r$ is the reflection with respect to a wall of a
vector chamber $\mathrm{D}$, one has
$Q_{\mathrm{D}}\subset Q_{r(\mathrm{D})}$. Now one may arrange the elements of
$\vs{\bSi}^{+}(\mathrm{D})$ in a nibbling order $(a_{1},\dots,a_{m})$ in such a
way that $r$ is the reflection with respect to the hyperplane $a_{1}=0$
(\btref{1.3.15}). Then $(-a_{1},\dots,-a_{m})$ is a nibbling order on
$\vs{\bSi}^{+}(-\mathrm{D})$, and one has (\btref{5.2.8})
\[
P_{\Omega+\mathrm{D}}=HU^{-1}_{a_{1}(\Omega)}\dots U^{-1}_{a_{m}(\Omega)} =U_{a_{m}(\Omega)}\dots U_{a_{1}(\Omega)}H
\]
whence
\begin{equation*}
Q_{\mathrm{D}}=U_{a_{m}(\Omega)}\dots U_{a_{1}(\Omega)}
HU^{-1}_{-a_{1}(\Omega)}\dots U^{-1}_{-a_{m}(\Omega)}.\tag{1}
\end{equation*}
Since $a_{1}(\Omega)\cap(-a_{1}(\Omega))$ contains the nonempty interior of
$\Omega$, one has ((\btref{5.2.2})\,(2))
\begin{equation*}
U_{a_{1}(\Omega)}HU^{-1}_{-a_{1}(\Omega)}
\subset P_{a_{1}(\Omega)\cap-a_{1}(\Omega)}
=P_{-a_{1}(\Omega)}P_{a_{1}(\Omega)}
=U_{-a_{1}(\Omega)}HU^{-1}_{a_{1}(\Omega)}.\tag{2}
\end{equation*}
But $\cl(\Omega-r(\mathrm{D}))=(-a_{1}(\Omega))\cap
\bigcap_{2\leqslant j\leqslant m}a_{j}(\Omega)$ and
$\cl(\Omega+r(\mathrm{D}))=a_{1}(\Omega)\cap
\bigcap_{2\leqslant j\leqslant m}(-a_{j}(\Omega))$. Hence (1) and (2) imply
\[
Q_{\mathrm{D}}\subset P_{\Omega+r(\mathrm{D})}P_{\Omega-r(\mathrm{D})} =Q_{r(\mathrm{D})}.
\]

\begin{btstat}{Proposition}{5.2.11}\btanchor{5.2.11}
Let $\Omega$ be a subset of $\bA$ with nonempty interior, and let $\mathrm{D}$
be a vector chamber. One has
$P_{\Omega}=P_{\Omega+\mathrm{D}}P_{\Omega-\mathrm{D}}$. If one arranges the
elements of $\vs{\bSi}(\Omega)\cap\vs{\bSi}^{+}(\mathrm{D})$ (resp.\
$\vs{\bSi}(\Omega)\cap\vs{\bSi}^{+}(-\mathrm{D})$) in a nibbling order
$(a_{1},\dots,a_{n})$ (resp.\ $(b_{1},\dots,b_{m})$), the product map
\[
(u_{1},\dots,u_{n},h,v_{1},\dots,v_{m})\mapsto u_{1}\dots u_{n}hv_{1}\dots v_{m}
\]
is a bijection of $\prod_{1\leqslant j\leqslant n}U_{a_{j}(\Omega)}\times H
\times\prod_{1\leqslant k\leqslant m}U^{-1}_{b_{k}(\Omega)}$ onto $P_{\Omega}$.
\end{btstat}

For every root $a\in\vs{\bSi}(\Omega)$ there exists a vector chamber
$\mathrm{D}'$ contained in the half-space $a>0$. One then has
$U_{a(\Omega)}P_{\Omega+\mathrm{D}'}\subset P_{\Omega+\mathrm{D}'}$. Since
$P_{\Omega+\mathrm{D}}P_{\Omega-\mathrm{D}}
=P_{\Omega+\mathrm{D}'}P_{\Omega-\mathrm{D}'}$ (Lemma (\btref{5.2.10})), one deduces
that $P_{\Omega+\mathrm{D}}P_{\Omega-\mathrm{D}}$ is stable under left
multiplication by $H$ and by the $U_{a(\Omega)}$ for $a\in\vs{\bSi}(\Omega)$,
hence by $P_{\Omega}$, whence the first assertion. The second then follows from
(\btref{5.2.6}) and (\btref{5.2.9}).

\begin{btstat}{Remark}{5.2.12}\btanchor{5.2.12}
\upshape
As in (\btref{5.2.8}), one may move the factor $H$ within the product considered,
provided one multiplies by $-1$ the exponent of those $U_{\alpha}$ which one
moves from one side of $H$ to the other.
\end{btstat}

\begin{btstat}{Remark}{5.2.13}\btanchor{5.2.13}
Let $a\in\vs{\bSi}$. One may arrange the elements of $\vs{\bSi}$ in an order
$(a_{1},\dots,a_{m})$ such that $a=a_{1}$ and such that, with the notation of
\textup{(\btref{5.2.4})}, the product map from
$\prod_{1\leqslant j\leqslant k}U_{a_{j}(\Omega)}\times H\times
\prod_{k<j\leqslant m}U^{-1}_{a_{j}(\Omega)}$ into $P_{\Omega}$ is injective
(resp.\ bijective) for every subset $\Omega$ of $\bA$ (resp.\ for every subset
$\Omega$ of $\bA$ with nonempty interior), and every $k$ with
$0\leqslant k\leqslant m$.
\end{btstat}
It suffices indeed to choose the vector chamber $\mathrm{D}$ and the nibbling
order on $\vs{\bSi}^{+}(\mathrm{D})$ in such a way that $a$ is its first
element, then to choose a nibbling order on $\vs{\bSi}^{+}(-\mathrm{D})$ and to
apply (\btref{5.2.9}) (resp.\ (\btref{5.2.11})).

Such an order on $\vs{\bSi}$ will be called \emph{admissible}.

\begin{btstat}{Corollary}{5.2.14}\btanchor{5.2.14}
Let $\alpha\in\bSi$ and let $\Omega'$ be a closed subset of $\bA$ such that
$\Omega=\alpha\cap\Omega'$ has nonempty interior and such that $\alpha$ and
$\Omega'$ are transverse \textup{(\btref{5.1.2})}. One has
$P_{\Omega}=P_{\alpha}\cdot P_{\Omega'}$.
\end{btstat}

\pg{101}
Let us arrange the elements of $\vs{\bSi}$ in an admissible order
$(a_{1},\dots,a_{m})$ in such a way that $a_{1}=\vs{\alpha}$. One then has, by
(\btref{5.2.13}):
\begin{equation*}
P_{\Omega}=U_{a_{1}(\Omega)}U_{a_{2}(\Omega)}\dots U_{a_{m}(\Omega)}H.\tag{1}
\end{equation*}
But the affine roots $a_{j}(\Omega)$ for $a_{j}\in\vs{\bSi}(\Omega)$ and
$j\geqslant2$ do not contain $a_{1}(\Omega)$, hence contain $\Omega'$. On the
other hand, either $a_{1}(\Omega)\supset\Omega'$, or $a_{1}(\Omega)\supset\alpha$,
whence $a_{1}(\Omega)=\alpha$. In both cases, (1) indeed implies
$P_{\Omega}=P_{\alpha}\cdot P_{\Omega'}$.

\begin{btstat}{Corollary}{5.2.15}\btanchor{5.2.15}
Let $\alpha\in\bSi$ and let $\mathrm{C}$ be a chamber contained in $\alpha$ and
having a panel contained in $\partial\alpha$. One has
$P_{\mathrm{C}}=P_{\alpha}\cdot P_{\mathrm{C}\cup r_{\alpha}(\mathrm{C})}$.
\end{btstat}
It suffices to apply (\btref{5.2.14}) to
$\Omega'=\mathrm{C}\cup r_{\alpha}(\mathrm{C})$.

\begin{btstat}{Proposition}{5.2.16}\btanchor{5.2.16}
Let $\Omega$ and $\Omega'$ be two closed subsets of $\bA$, with nonempty
interior and such that $\Omega\cap\Omega'\neq\emptyset$. Then one has
\[
P_{\Omega}\cap P_{\Omega'}=P_{\cl(\Omega\cup\Omega')}.
\]
If moreover $P_{\Omega}\subset P_{\Omega'}$, one has $\Omega\supset\Omega'$.
\end{btstat}

It is clear that $P_{\cl(\Omega\cup\Omega')}\subset P_{\Omega}\cap P_{\Omega'}$.
On the other hand, let $x\in\Omega\cap\Omega'$ and let
$g\in P_{\Omega}\cap P_{\Omega'}\subset P_{\{x\}}$. Let us arrange the elements
of $\vs{\bSi}$ in an admissible order $(a_{1},\dots,a_{m})$. By (\btref{5.2.13}) one has
\[
g=u_{1}\dots u_{m}h=v_{1}\dots v_{m}k
\]
with $h,k\in H$, $u_{j}\in U_{a_{j}(\Omega)}$ and $v_{j}\in U_{a_{j}(\Omega')}$
for $1\leqslant j\leqslant m$, whence $u_{j},v_{j}\in U_{a_{j}(\{x\})}$.
Applying (\btref{5.2.9}) to $P_{\{x\}}$, one concludes that $u_{j}=v_{j}$ for every $j$,
whence
\[
u_{j}=v_{j}\in U_{a_{j}(\Omega)}\cap U_{a_{j}(\Omega')} =U_{a_{j}(\Omega)\cup a_{j}(\Omega')}\subset U_{a_{j}(\Omega\cup\Omega')}
\]
and finally $g\in P_{\Omega\cup\Omega'}$, which proves the first assertion of the
proposition.

If moreover $P_{\Omega}\subset P_{\Omega'}$, what precedes shows that
$U_{a_{j}(\Omega)}\subset U_{a_{j}(\Omega')}$, whence (by (A\,2))
$a_{j}(\Omega)\supset a_{j}(\Omega')$ for $1\leqslant j\leqslant m$, and
$\Omega\supset\Omega'$.

\begin{btstat}{Corollary}{5.2.17}\btanchor{5.2.17}
Let $\alpha\in\bSi$ and let $\mathrm{C}$ be a chamber contained in $\alpha$ and
having a panel contained in $\partial\alpha$. One has
$P_{\alpha^{*}}\cap P_{\mathrm{C}}=P_{(\alpha^{*})_{+}}$.
\end{btstat}
It suffices to apply (\btref{5.2.16}) with $\Omega=\alpha^{*}$ and
$\Omega'=\mathrm{C}$.

\begin{btstat}{Lemma}{5.2.18}\btanchor{5.2.18}
Let $\mathrm{C}$ and $\mathrm{C}'$ be two chambers. If
$P_{\mathrm{C}}\supset P_{\mathrm{C}'}$, one has $\mathrm{C}=\mathrm{C}'$.
\end{btstat}

Suppose $P_{\mathrm{C}}\supset P_{\mathrm{C}'}$ with
$\mathrm{C}\neq \mathrm{C}'$. One can find $x\in \mathrm{C}$ and
$y\in \mathrm{C}'$ such that $y-x$ belongs to a vector chamber $\mathrm{D}$ and
such that the segment $[xy]$ meets no facet of codimension $\geqslant2$. Let
$\mathrm{C}_{1}$ be the chamber adjacent to $\mathrm{C}$ meeting $[xy]$. One has
\[
P_{\mathrm{C}_{1}}=P_{\mathrm{C}_{1}+\mathrm{D}}\cdot P_{\mathrm{C}_{1}-\mathrm{D}}
\subset P_{x+\mathrm{D}}\cdot P_{y-\mathrm{D}}
\subset P_{\mathrm{C}}\cdot P_{\mathrm{C}'}=P_{\mathrm{C}}.
\]
Since $\overline{\mathrm{C}}_{1}\cap\overline{\mathrm{C}}\neq\emptyset$,
Proposition (\btref{5.2.16}) then implies
$\mathrm{C}\subset\overline{\mathrm{C}}_{1}$, which is absurd. One therefore
indeed has $\mathrm{C}=\mathrm{C}'$.

\begin{btstat}{Theorem}{5.2.19}\btanchor{5.2.19}
Put $B=P_{\mathrm{C}}$. One has $B\cap N=H$ and the quadruple $(G,B,N,S)$ is a
saturated Tits system of affine type, the homomorphism $\nu$ defining by passage
to the quotient \pg{102}an isomorphism of the Coxeter system $(W,S)$ onto
$(\bW,\bS)$. Moreover, $(G,B,N,S)$ is a double Tits system \textup{(\btref{5.1.1})} and
the notations $P_{\Omega}$ are coherent with those of \S\,\textup{4}: if one
identifies $\bA$ with its canonical image in the building $\Imm$ of
$(G,B,N,S)$, the subgroup $P_{\Omega}$ is, for every subset $\Omega$ of $\bA$,
the fixer of $\Omega$ in $G$.
\end{btstat}

The proof of this theorem will occupy n\textsuperscript{os} (\btref{5.2.20}) to (\btref{5.2.26}).

\btpar{5.2.20}\btanchor{5.2.20}
For every affine root $\alpha\in\bSi$, there exists $w\in\bW$ such that
$w.\bC\subset\alpha$, whence $P_{\alpha}\subset P_{w.\bC}=nBn^{-1}$ for
$n\in\nu^{-1}(w)$. Condition (A\,6) therefore implies that \emph{$G$ is
generated by $B\cup N$}.

On the other hand, let $n\in B\cap N$. One then has $nBn^{-1}=B$, whence
$P_{\nu(n).\bC}=P_{\bC}$ and $\nu(n).\bC=\bC$ by Lemma (\btref{5.2.18}). One therefore
has $\nu(n)=1$ and $n\in H$. This shows that $B\cap N=H$ and that $B\cap N$ is
normal in $N$, the quotient group being $W$. Consequently, axioms (T\,1) and
(T\,2) of Tits systems are satisfied.

\btpar{5.2.21}\btanchor{5.2.21}
Let us show that $H=\bigcap_{n\in N}nBn^{-1}$ (which will imply that the Tits
system is \emph{saturated} (\btref{1.2.12})). Indeed, let
$g\in\bigcap_{n\in N}nBn^{-1}$. For every bounded closed subset $\Omega$ of
$\bA$ containing $\bC$, one can find a sequence of chambers
$(\mathrm{C}_{i})_{0\leqslant i\leqslant n}$ with $\mathrm{C}_{0}=\bC$, such
that, if one puts $\Omega_{j}=\bigcup_{0\leqslant i\leqslant j}\overline{\mathrm{C}}_{i}$,
one has $\Omega_{j}\cap\overline{\mathrm{C}}_{j+1}\neq\emptyset$ for every $j$
and $\Omega_{n}=\Omega$. Since $g\in P_{\mathrm{C}}$ for every chamber
$\mathrm{C}$, one sees by induction on $j$ that $g\in P_{\Omega_{j}}$ for every
$j$: it suffices to use (\btref{5.2.16}). Hence $g\in P_{\Omega}$. Let us then arrange
the elements of $\vs{\bSi}$ in an admissible order $(a_{1},\dots,a_{m})$ and
write $g=u_{1}\dots u_{m}h$, with $h\in H$ and $u_{i}\in U_{a_{i}(\bC)}$. Since
$g\in P_{\Omega}$, one has $u_{i}\in U_{a_{i}(\Omega)}$ for every bounded closed
subset $\Omega$ of $\bA$ containing $\bC$, whence $u_{i}\in H$ by (A\,2) and
finally $g\in H$, which was to be proved.

\btpar{5.2.22}\btanchor{5.2.22}
Let us now prove axiom (T\,3). Let $r\in S$ and let $\alpha$ be the affine root
containing $\bC$ such that $r=r_{\alpha}$. By (\btref{5.2.15}) one has
$B=P_{\alpha}\cdot P_{\bC\cup r(\bC)}$. Let $n\in N$; one therefore has
\[
rBnB=rP_{\bC\cup r(\bC)}P_{\alpha}nB=P_{\bC\cup r(\bC)}rP_{\alpha}nB.
\]
If $\nu(n)^{-1}(\alpha)\supset\bC$, one has $n^{-1}P_{\alpha}n\subset B$, whence
\begin{equation*}
rBnB\subset P_{\bC\cup r(\bC)}rn(n^{-1}P_{\alpha}n)B\subset BrnB.\tag{1}
\end{equation*}
If $\nu(n)^{-1}(\alpha)\not\supset\bC$, one has $\nu(rn)^{-1}(\alpha)\supset\bC$
and $(rn)^{-1}P_{\alpha}rn\subset B$. On the other hand one has, by (A\,4),
\[
rP_{\alpha}r\subset(P_{(\alpha^{*})_{+}}P_{\alpha}\cup P_{\alpha}rP_{\alpha})
\subset(B\cup BrP_{\alpha})
\]
whence
\begin{equation*}
\begin{split}
rBnB=P_{\bC\cup r(\bC)}rP_{\alpha}rrnB&\subset(BrnB\cup BrP_{\alpha}rnB)\\
&=BrnB\cup Bn(rn)^{-1}P_{\alpha}rnB=BrnB\cup BnB
\end{split}\tag{2}
\end{equation*}
and formulas (1) and (2) prove (T\,3).

\btpar{5.2.23}\btanchor{5.2.23}
Finally, Lemma (\btref{5.2.18}) implies $rBr^{-1}\neq B$ for every $r\in S$, that is to
say axiom (T\,4). We have therefore indeed proved that \emph{$(G,B,N,S)$ is a
Tits system of affine type}, and \emph{saturated} by (\btref{5.2.21}).

\pg{103}
We may therefore construct the associated building $\Imm$, the canonical map
$j:\bA\to\Imm$, and consider the fixer in $G$ of the image $j(\Omega)$ of a
subset $\Omega$ of $\bA$, a fixer which we shall provisionally denote by
$\Pt_{\Omega}$. One knows that $\Pt_{\mathrm{C}}=P_{\mathrm{C}}$ for every
chamber $\mathrm{C}$ (\btref{2.1.5}). If $\alpha\in\bSi$ and $g\in P_{\alpha}$, one has
$g\in P_{\mathrm{C}}$ for every chamber $\mathrm{C}\subset\alpha$, whence
$g\in\Pt_{\alpha}$. Since $P_{\Omega}$ is by definition generated by $H$ and the
$P_{\alpha}$ for $\alpha\supset\Omega$, this implies
$P_{\Omega}\subset\Pt_{\Omega}$ for every $\Omega\subset\bA$.

\btpar{5.2.24}\btanchor{5.2.24}
Consequently, by (\btref{5.2.11}), one has
\[
\Pt_{\mathrm{C}}=P_{\mathrm{C}}=P_{\mathrm{C}+\mathrm{D}}\cdot P_{\mathrm{C}-\mathrm{D}}
\subset\Pt_{\mathrm{C}+\mathrm{D}}\cdot\Pt_{\mathrm{C}-\mathrm{D}}
\subset\Pt_{\mathrm{C}}
\]
for every chamber $\mathrm{C}$ and every vector chamber $\mathrm{D}$, which
shows that condition (iv) of Theorem (\btref{5.1.3}) is satisfied. In other words,
\emph{$(G,B,N,S)$ is indeed a double Tits system}.

\btpar{5.2.25}\btanchor{5.2.25}
To complete the proof of Theorem (\btref{5.2.19}) it only remains to prove that
$\Pt_{\Omega}=P_{\Omega}$ for every subset $\Omega$ of $\bA$.

Suppose first $\Omega$ closed and with nonempty interior. One may assume
moreover that $\Omega\supset\bC$. Let us arrange the elements of $\vs{\bSi}$ in
an admissible order $(a_{1},\dots,a_{m})$. If
$g\in\Pt_{\Omega}\subset\Pt_{\bC}=P_{\bC}$, one may write
$g=u_{1}\dots u_{m}h$, with $u_{j}\in U_{a_{j}(\bC)}$ and $h\in H$. Let
$j\in\{1,\dots,m\}$. There exists a chamber $\mathrm{C}\subset\Omega$ with
$a_{j}(\mathrm{C})=a_{j}(\Omega)$. Let
$\Gamma=(\bC,\mathrm{C}_{1},\dots,\mathrm{C}_{q}=\mathrm{C})$ be a taut gallery
between $\bC$ and $\mathrm{C}$. One has $\Gamma\subset\Omega$ and
$g\in\Pt_{\mathrm{C}_{i}}=P_{\mathrm{C}_{i}}$ for every $i=1,\dots,q$. One may
therefore write $g=u^{i}_{1}\dots u^{i}_{m}h^{i}$ with
$u^{i}_{k}\in U_{a_{k}(\mathrm{C}_{i})}$ for $1\leqslant k\leqslant m$ and
$h^{i}\in H$. But
$\overline{\mathrm{C}}_{i}\cap\overline{\mathrm{C}}_{i+1}\neq\emptyset$ and one
sees, on placing oneself in
$P_{\overline{\mathrm{C}}_{i}\cap\overline{\mathrm{C}}_{i+1}}$ and using
(\btref{5.2.13}), that $u^{i+1}_{k}=u^{i}_{k}$ for $1\leqslant k\leqslant m$ and
$0\leqslant i<q$. It follows that $u^{i}_{k}$ is independent of $i$ and that
$u_{k}\in U_{a_{k}(\mathrm{C})}$ for $1\leqslant k\leqslant m$. In particular,
one has $u_{j}\in U_{a_{j}(\mathrm{C})}=U_{a_{j}(\Omega)}$. Applying this for
$j=1,\dots,m$, one concludes $g\in P_{\Omega}$, whence
$\Pt_{\Omega}=P_{\Omega}$.

Let now $\Omega$ be an arbitrary nonempty closed subset of $\bA$ (one has
$\Pt_{\emptyset}=G=P_{\emptyset}$) and let us resume the notation $\mathrm{L}$,
$\mathrm{F}$ and $\Delta$ of (\btref{5.1.27}) and (\btref{5.1.28}); by (\btref{5.1.28}) one has
$\Pt_{\Omega}=\Pt_{\Delta}\Nt_{\Omega}\Pt_{\Delta}\Pt_{\Ot}$ with
$\Nt_{\Omega}=\Pt_{\Omega}\cap N$. Since $\Delta$ and $\Ot$ have nonempty
interior, one has $\Pt_{\Delta}=P_{\Delta}\subset P_{\Omega}$ and
$\Pt_{\Ot}=P_{\Ot}\subset P_{\Omega}$, and it suffices to prove that
$\Nt_{\Omega}\subset P_{\Omega}$. Now $\Nt_{\Omega}/H=\bW_{\mathrm{F}}$ is
generated by the reflections $r_{\alpha}$ associated with the affine roots
$\alpha$ such that $\mathrm{F}\subset\partial\alpha$, or again
$\Omega\subset\alpha\cap\alpha^{*}$. Our assertion therefore follows from the
following lemma:

\begin{btstat}{Lemma}{5.2.26}\btanchor{5.2.26}
Let $\alpha\in\bSi$. The subset $\nu^{-1}(r_{\alpha})$ of $N$ is contained in
the subgroup generated by $P_{\alpha}\cup P_{\alpha^{*}}$.
\end{btstat}

If this assertion were false, one would have
$P_{\alpha^{*}}\cap P_{\alpha}r_{\alpha}P_{\alpha}=\emptyset$ and consequently
$P_{\alpha^{*}}\subset P_{\alpha}\cdot P_{(\alpha^{*})_{+}}$ by (A\,4) (or by
condition (viii) of Theorem (\btref{5.1.3})). Using (\btref{5.2.2})\,(3), this would imply that
$P_{\alpha^{*}}$ leaves fixed every chamber contained in
$\alpha\cap(\alpha^{*})_{+}$, which contradicts (\btref{5.1.11}).

The proof of Theorem (\btref{5.2.19}) is complete.

\pg{104}
\btpar{5.2.27}\btanchor{5.2.27}
We shall now prove a ``converse'' of Theorem (\btref{5.2.19}). For this, let us resume
the hypotheses and notation of 5.1, except that we provisionally denote by
$\Pt_{\Omega}$ the fixer of $\Omega$ in $G$ and that we reserve the notation
$P_{\Omega}$ for the subgroup generated by the $\Pt_{\alpha}=P_{\alpha}$ for
$\alpha\in\bSi$ and $\alpha\supset\Omega$.

\begin{btstatu}{Proposition}
With the hypotheses and notation specified above, conditions
\textup{(A\,1)} to \textup{(A\,6)} of \textup{(\btref{5.2.1})} are satisfied as soon as
$(G,B,N,S)$ is a saturated double Tits system.
\end{btstatu}

The verification of (A\,1) is immediate; (A\,2) follows from (\btref{5.1.10}) and from
(\btref{2.2.5})\,(i); (A\,4) follows from condition (viii) of (\btref{5.1.3}); and (A\,5) is
verified immediately: if $g\in P_{x+\mathrm{D}}\cap P_{x-\mathrm{D}}$, then $g$
leaves fixed all the points of the sector $x+\mathrm{D}$ and of the opposite
sector $x-\mathrm{D}$, hence of their convex hull, which is nothing but $\bA$,
and $g\in\Pt_{\bA}=H$.

To establish (A\,3) and (A\,6), we shall first prove that if $\Omega$ is a
closed subset of $\bA$ containing a sector, one has $\Pt_{\Omega}=P_{\Omega}$.
For this we shall argue as in (\btref{5.2.6}). Let $(a_{1},\dots,a_{m})$ be the elements
of $\vs{\bSi}(\Omega)$ arranged in a nibbling order. Let us argue by induction
on $m$, the cases $m=0$ and $m=1$ being trivial. Put
$\Omega'=\bigcap_{2\leqslant j\leqslant m}a_{j}(\Omega)$; we saw in (\btref{5.2.6}) that
$\vs{\bSi}(\Omega')=\{a_{2},\dots,a_{m}\}$. By the induction hypothesis one
therefore has $\Pt_{\Omega'}=P_{\Omega'}$. Moreover, $a_{1}(\Omega)$ and
$\Omega'$ are transverse and, by (\btref{5.1.3})\,(iii), one has
\[
\Pt_{\Omega}=\Pt_{a_{1}(\Omega)}\cdot\Pt_{\Omega'}
=P_{a_{1}(\Omega)}\cdot P_{\Omega'}\subset P_{\Omega}
\]
whence $\Pt_{\Omega}=P_{\Omega}$.

Let us then place ourselves in the hypotheses of (A\,3). The closed subsets
$\alpha$ and $\Omega$ are then transverse and one has, by (\btref{5.1.3})\,(iii),
$P_{\alpha\cap\beta}\subset\Pt_{\alpha\cap\beta}=\Pt_{\alpha\cap\Omega}
=\Pt_{\alpha}\cdot\Pt_{\Omega}=P_{\alpha}\cdot\Pt_{\Omega}$. But $\Omega$
contains a sector, hence $\Pt_{\Omega}=P_{\Omega}$, whence
$P_{\alpha\cap\beta}\subset P_{\alpha}\cdot P_{\Omega}\subset P_{\alpha\cap\beta}$,
which proves (A\,3).

Finally, (\btref{5.1.3})\,(iv) shows that
\[
B=\Pt_{\bC+\mathrm{D}}\cdot\Pt_{\bC-\mathrm{D}} =P_{\bC+\mathrm{D}}\cdot P_{\bC-\mathrm{D}}\subset P_{\bC}
\]
for every vector chamber $\mathrm{D}$. Consequently the subgroup generated by
the $P_{\alpha}$ contains $B$. One proves exactly as in (\btref{5.2.26}) that it
contains the subsets $\nu^{-1}(r)$ for $r\in S$, hence $N$. This establishes
(A\,6) and completes the proof of the proposition.

\btpar{5.2.28}\btanchor{5.2.28}
It then follows from Theorem (\btref{5.2.19}) that $\Pt_{\Omega}=P_{\Omega}$ for every
subset $\Omega$ of $\bA$. In other words, returning completely to the notation
of 5.1, \emph{the fixer $P_{\Omega}$ of an arbitrary subset $\Omega$ of $\bA$
is, when $(G,B,N,S)$ is a double Tits system, generated by the subgroups
$P_{\alpha}$ for $\alpha\in\bSi$ and $\alpha\supset\Omega$, and the results
which we established above on the structure of the subgroups $P_{\Omega}$, and
in particular those of \textup{(\btref{5.2.6})}, \textup{(\btref{5.2.9})}, \textup{(\btref{5.2.11})},
\textup{(\btref{5.2.13})} and \textup{(\btref{5.2.16})}, are valid} ($U_{\alpha}$ being a subset
of $P_{\alpha}$ satisfying the conditions of (\btref{5.2.3})).

\btpar{5.2.29}\btanchor{5.2.29}
This implies that \emph{the subgroup $\Bg^{0}_{\mathrm{D}}$ (the union of the
subgroups $P_{\Cq}$ as $\Cq$ ranges over the set of sectors of direction
$\mathrm{D}$ \textup{(\btref{4.1.5})}) is generated by the union of the subgroups
$P_{\alpha}$ for $\vs{\alpha}\in\vs{\bSi}^{+}(\mathrm{D})$}. More precisely, if
one arranges the elements of $\vs{\bSi}^{+}(\mathrm{D})$ in

\medskip
\noindent\hrulefill
\par\endgroup

% ---------------------------------------- bruhattitsen_2.tex
\begingroup\providecommand{\vv}{}
\renewcommand{\vv}[1]{{}^{v}\!#1}
\providecommand{\vW}{}
\renewcommand{\vW}{\vv{W}}
\providecommand{\vnu}{}
\renewcommand{\vnu}{\vv{\nu}}
\providecommand{\vSig}{}
\renewcommand{\vSig}{\vv{\Sigma}}
\providecommand{\vA}{}
\renewcommand{\vA}{\vv{\mathbf{A}}}
\providecommand{\N}{}
\renewcommand{\N}{\mathcal{N}}
\providecommand{\red}{}
\renewcommand{\red}{\mathrm{red}}
\providecommand{\nm}{}
\renewcommand{\nm}{\mathrm{nm}}
\providecommand{\Phired}{}
\renewcommand{\Phired}{\Phi^{\red}}
\providecommand{\Bs}{}
\renewcommand{\Bs}{\mathfrak{B}}
\providecommand{\Us}{}
\renewcommand{\Us}{\mathfrak{U}}
\providecommand{\Cs}{}
\renewcommand{\Cs}{\mathfrak{C}}
\providecommand{\Ps}{}
\renewcommand{\Ps}{\mathscr{P}}
\providecommand{\co}{}
\renewcommand{\co}[1]{#1^{\vee}}
\providecommand{\bA}{}
\renewcommand{\bA}{\mathbf{A}}
\providecommand{\bV}{}
\renewcommand{\bV}{\mathbf{V}}
\providecommand{\bW}{}
\renewcommand{\bW}{\mathbf{W}}
\providecommand{\bS}{}
\renewcommand{\bS}{\mathbf{S}}
\providecommand{\bSig}{}
\renewcommand{\bSig}{\boldsymbol{\Sigma}}
\providecommand{\R}{}
\renewcommand{\R}{\mathbf{R}}
\providecommand{\Z}{}
\renewcommand{\Z}{\mathbf{Z}}
\providecommand{\Q}{}
\renewcommand{\Q}{\mathbf{Q}}
\providecommand{\Nstar}{}
\renewcommand{\Nstar}{\mathbf{N}^{*}}
\providecommand{\Add}{}
\renewcommand{\Add}{\mathscr{A}\!dd}
\providecommand{\Gsc}{}
\renewcommand{\Gsc}{\mathscr{G}}
\providecommand{\Tsc}{}
\renewcommand{\Tsc}{\mathscr{T}}
\providecommand{\Usc}{}
\renewcommand{\Usc}{\mathscr{U}}
\providecommand{\Ssc}{}
\renewcommand{\Ssc}{\mathscr{S}}
\providecommand{\Zsc}{}
\renewcommand{\Zsc}{\mathscr{Z}}
\providecommand{\para}{}
\renewcommand{\para}[1]{\par\medskip\noindent\textbf{(#1)}\ \ignorespaces}
\newenvironment{stmt}[2]{%
  \par\medskip\noindent\textbf{#1 (#2)}. --- \itshape\ignorespaces}{\par\medskip}
\newenvironment{sketch}{\par\smallskip\noindent\small\textit{Proof sketch.}\ \ignorespaces}{\par\smallskip}
\setlength{\parindent}{1.5em}

\para{5.2.29, end} With the elements of $\vSig^{+}(D)$ arranged in a nibbling
order $(a_1,\dots,a_m)$, the product map
$(u_1,\dots,u_m,h)\mapsto u_1\cdots u_m h$ is a bijection of
$\prod_{1\le j\le m}\Us_{a_j}\times H$ onto $\Bs^{0}_{D}$; here $\Us_a$ denotes
the union of the $U_\alpha$ with $\vv{\alpha}=a$. Likewise
$(u_1,\dots,u_m,t)\mapsto u_1\cdots u_m t$ is a bijection of
$\prod_{1\le j\le m}\Us_{a_j}\times\nu^{-1}(\bV)$ onto $\Bs_{D}$.

\para{5.2.30}\btanchor{5.2.30} We now study the case in which each $P_\alpha$ decomposes as a
semidirect product of $H$ by a normal subgroup $U_\alpha$, these decompositions
being suitably coherent. Precisely, until the end of \S\,5.2 we are given: a
group $G$; a subgroup $N$ of $G$; a surjective homomorphism $\nu$ of $N$ onto
$\bW$ with kernel denoted $H$; and, for each $\alpha\in\bSig$, a subgroup
$U_\alpha$ of $G$. These data are required to satisfy conditions
(B\,1)--(B\,6) below, where for any subset $\Omega\subset\bA$ we write
$U_\Omega$ for the subgroup of $G$ generated by the $U_\alpha$ with
$\alpha\in\bSig$ and $\alpha\supset\Omega$ (condition (B\,2) shows this notation
is unambiguous when $\Omega=\alpha\in\bSig$).

\begin{description}
\item[\rm (B\,1)] For all $\alpha\in\bSig$ and all $n\in N$,
  $\;nU_\alpha n^{-1}=U_{\nu(n).\alpha}$.
\item[\rm (B\,2)] Let $\alpha\in\bSig$. For every $\beta\in\bSig$ with
  $\beta\supsetneq\alpha$ one has $U_\beta\subsetneq U_\alpha$, and the
  intersection of the $U_\beta$ for $\beta\in\bSig$, $\beta\supset\alpha$, is
  reduced to the identity.
\item[\rm (B$'$3)] Let $\alpha,\beta\in\bSig$ have non-parallel boundaries, and
  let $\Omega$ be the intersection of those $\gamma\in\bSig$ with
  $\gamma\supset\alpha\cap\beta$ and $\gamma\not\supset\alpha$. Then $U_\alpha$
  normalizes $U_\Omega$.
\item[\rm (B$''$3)] Let $\alpha,\beta\in\bSig$ be such that
  $\beta\supsetneq\alpha^{*}$. Then $U_\alpha H U_\beta$ is a subgroup of $G$.
\item[\rm (B\,4)] Let $\alpha\in\bSig$. The subset
  $U_\alpha H U_{(\alpha^{*})_{+}}\cup U_\alpha\nu^{-1}(r_\alpha)U_\alpha$
  is a subgroup of $G$.
\item[\rm (B\,5)] For every $x\in\bA$ and every vector chamber $D$,
  $\;(H.U_{x+D})\cap U_{x-D}=\{1\}$.
\item[\rm (B\,6)] $G$ is generated by the union of $H$ and the $U_\alpha$,
  $\alpha\in\bSig$.
\end{description}

Because of (B\,1) it suffices to verify (B\,5) for one vector chamber fixed in
advance. Also (B\,1) implies that $H$ normalizes each $U_\alpha$, hence each
$U_\Omega$. One sets $P_\Omega=H.U_\Omega$: this is the subgroup of $G$
generated by $H$ and the $P_\alpha$ for $\alpha\in\bSig$, $\alpha\supset\Omega$.
By (B\,5), $P_\Omega$ is the semidirect product of $H$ by $U_\Omega$ as soon as
$\Omega$ contains a sector.

\begin{stmt}{Proposition}{5.2.31}\btanchor{5.2.31}
$N$, $\nu$ and the family of the $P_\alpha$ ($\alpha\in\bSig$) satisfy
conditions (A\,1)--(A\,6) of (\btref{5.2.1}).
\end{stmt}
\noindent This is immediate. Moreover the subgroups $U_\alpha$,
$\alpha\in\bSig$, satisfy the conditions of (\btref{5.2.3}), and $\Us_a$ is then a
subgroup for every $a\in\vSig$.

\para{5.2.32}\btanchor{5.2.32} The previous results therefore apply. In particular the quadruple
$(G,B,N,S)$ --- with $B=P_{\Cs}$ and $S$ the inverse image of $\bS$ in $N/H$ ---
is a saturated double Tits system. Furthermore (\btref{5.2.6}) and (\btref{5.2.11}) can be made
precise:

\begin{stmt}{Proposition}{5.2.32}
Let $\Omega$ be a non-empty subset of $\bA$.
\begin{enumerate}
\item[(i)] Suppose $\Omega$ contains a sector, and arrange the elements of
$\vSig(\Omega)$ in a nibbling order $(a_1,\dots,a_m)$. Then
$(u_1,\dots,u_m)\mapsto u_1\cdots u_m$ is a bijection of
$\prod_{1\le j\le m}U_{a_j(\Omega)}$ onto the subgroup $U_\Omega$.
\item[(ii)] Suppose $\Omega$ has non-empty interior and let $D$ be a vector
chamber. Then $(u,h,v)\mapsto uhv$ is a bijection of
$U_{\Omega+D}\times H\times U_{\Omega-D}$ onto $P_\Omega$.
\end{enumerate}
\end{stmt}
\noindent The proof of (i) runs exactly as that of (\btref{5.2.6}), using (B$'$3) in
place of (A\,3); (ii) follows from (i) together with (\btref{5.2.11}).

\para{5.2.33}\btanchor{5.2.33} Proposition (\btref{5.1.28}) may be sharpened similarly: in the notation
of (\btref{5.1.28}) one has $P_\Omega=U_\Delta N_\Omega U_\Delta U_{\widetilde\Delta}$.

\para{5.2.34}\btanchor{5.2.34} Let $D$ be a vector chamber. If $\Cs_1,\Cs_2$ are two sectors of
$\bA$ of direction $D$, then
$U_{\Cs_1}\cup U_{\Cs_2}\subset U_{\Cs_1\cap\Cs_2}$, and the union of the
subgroups $U_{\Cs}$, as $\Cs$ runs over the sectors of direction $D$, is a
subgroup, denoted $\Us_{D}$. Arranging the elements of $\vSig^{+}(D)$ in a
nibbling order $(a_1,\dots,a_m)$, it follows from (\btref{5.2.32}) that
$(u_1,\dots,u_m)\mapsto u_1\cdots u_m$ is a bijection of
$\prod_{1\le j\le m}\Us_{a_j}$ onto $\Us_{D}$. On the other hand (cf.\ (\btref{5.2.29}))
$\Bs^{0}_{D}$ is the semidirect product of $H$ by the normal subgroup $\Us_D$,
and $\Bs_D$ is the semidirect product of $\nu^{-1}(\bV)$ by the normal subgroup
$\Us_D$. Finally $G=\Bs_D.N.\Bs_D=\Us_D.N.\Us_D$, and one can even show that the
evident map is a bijection of $N$ onto $\Us_D\backslash G/\Us_D$.

\para{5.2.35}\btanchor{5.2.35} Let $\alpha\in\bSig$ and let $\Omega$ be a closed subset of $\bA$.
Assume $\alpha$ and $\Omega$ are transverse and that $\Omega'=\alpha\cap\Omega$
contains a chamber of $\bA$. Then $U_\alpha$ normalizes $U_\Omega$ and
$U_{\Omega'}=U_\alpha.U_\Omega$.

\begin{sketch}
Let $\beta\in\bSig$ with $\beta\supset\Omega$. For $\gamma\in\bSig$ with
$\gamma\supset\alpha\cap\beta$ and $\gamma\not\supset\alpha$ one has
$\gamma\supset\alpha\cap\Omega$ and $\gamma\not\supset\alpha$, hence
$\gamma\supset\Omega$. So for $u\in U_\alpha$, condition (B$'$3) gives
$uU_\beta u^{-1}\in U_\Omega$, which is the assertion.
\end{sketch}

\begin{stmt}{Remark}{5.2.36}\btanchor{5.2.36}
Condition (B$'$3) is equivalent to:
\begin{description}
\item[\rm (B$'$3 bis)] For $\alpha,\beta\in\bSig$ with non-parallel boundaries,
the commutator group $(U_\alpha,U_\beta)$ is contained in the subgroup generated
by the $U_\gamma$ for $\gamma\in\bSig$ with $\gamma\supset\alpha\cap\beta$,
$\gamma\not\supset\alpha$ and $\gamma\not\supset\beta$.
\end{description}
\end{stmt}
\noindent Indeed (B$'$3 bis) trivially implies (B$'$3), and the converse follows
from (\btref{5.2.6}).\footnote{In [\btbib{12}] and [\btbib{16}] a triple $(N,\nu,(U_\alpha)_{\alpha\in\bSig})$
satisfying (B\,1), (B\,2), (B$'$3 bis), (B$''$3), (B\,4), (B\,6) together with a
condition (DR\,6) analogous to but distinct from (B\,5) was called an
\emph{affine root datum of type $\bSig$}. In fact that condition (DR\,6), although
in general satisfied in the applications we have in view, does not seem to be a
consequence of the conditions of (\btref{5.2.30}); and moreover it is perhaps
insufficient to establish axiom (T\,4) of Tits systems, contrary to what we
asserted in [\btbib{12}] and [\btbib{16}].}

\bigskip\hrule\bigskip

Throughout \S\,6, $V$ denotes a real vector space, $V^{*}$ its dual, and $\Phi$
a root system in $V^{*}$ with Weyl group $\vW$ (acting on $V$ and on $V^{*}$).
A $\vW$-invariant inner product on $V$, hence on $V^{*}$, is given. For
$a\in\Phi$, $r_a$ is the reflection attached to $a$, and $\co{a}$ is the
\emph{coroot} attached to $a$, i.e.\ the element of $V$ with
$r_a(v)=v-a(v)\co{a}$ for all $v\in V$; the $\co{a}$, $a\in\Phi$, form the root
system dual to $\Phi$. A chamber $D_0$ of $\Phi$ in $V$ is fixed; $\Pi$ is the
corresponding basis of $\Phi$ (``system of simple roots''), and $\Phi^{+}$
(resp.\ $\Phi^{-}$) the set of positive (resp.\ negative) roots for $D_0$
(cf.\ (\btref{1.3.10}), (\btref{1.3.12})). We write $\Phired$ (resp.\ $\Phi^{\nm}$) for the set
of non-divisible (resp.\ non-multipliable) roots, i.e.\ the set of $a\in\Phi$
with $a/2\notin\Phi$ (resp.\ $2a\notin\Phi$), and we put
$\Phi^{+\red}=\Phired\cap\Phi^{+}$, $\Phi^{-\red}=\Phired\cap\Phi^{-}$.

In a group $X$, the normalizer of a subset $Y\subset X$ is written $\N_X(Y)$, or
simply $\N(Y)$ when no confusion is possible.

\btsecA{6}{6. Valued root data}
\btsecB{6.1}{6.1. Root data}

\begin{stmt}{Definition}{6.1.1}\btanchor{6.1.1}
A \emph{root datum of type $\Phi$} in a group $G$ is a system
$(T,(U_a,M_a)_{a\in\Phi})$ with the following properties.
\end{stmt}

\begin{description}
\item[\rm (DR\,1)] $T$ is a subgroup of $G$ and, for every $a\in\Phi$, $U_a$ is
a subgroup of $G$ not reduced to the identity element.
\item[\rm (DR\,2)] For $a,b\in\Phi$, the commutator group $(U_a,U_b)$ is
contained in the group generated by the $U_{pa+qb}$ for $p,q$ strictly positive
integers with $pa+qb\in\Phi$.
\item[\rm (DR\,3)] If $a$ and $2a$ belong to $\Phi$, then $U_{2a}\subsetneq U_a$.
\item[\rm (DR\,4)] For $a\in\Phi$, $M_a$ is a right coset of $T$ and
$U_{-a}-\{1\}\subset U_aM_aU_a$.
\item[\rm (DR\,5)] For $a,b\in\Phi$ and $n\in M_a$,
$\;nU_bn^{-1}=U_{r_a(b)}$.
\item[\rm (DR\,6)] If $U^{+}$ (resp.\ $U^{-}$) denotes the group generated by
the $U_a$ for $a\in\Phi^{+}$ (resp.\ $\Phi^{-}$), then $TU^{+}\cap U^{-}=\{1\}$.
\end{description}

\noindent The root datum is called \emph{generating} if $T$ and the $U_a$
generate $G$.

\para{6.1.2}\btanchor{6.1.2} Some immediate consequences of these axioms. For $a\in\Phi$ we put
$U_a^{*}=U_a-\{1\}$.

\medskip
\noindent\textbf{(1)} For $a\in\Phi$ one has $U_a\ne U_{-a}$ and
$U_aM_aU_a\cap\N(U_a)=\emptyset$; in particular $U_{-a}\cap\N(U_a)=\{1\}$.

\begin{sketch}
The first relation follows from (DR\,1) and (DR\,6). For the second: if
$g\in U_amU_a\cap\N(U_a)$ with $m\in M_a$, then $gU_ag^{-1}=U_a$, whence by
(DR\,5) $U_a=mU_am^{-1}=U_{-a}$.
\end{sketch}

\noindent\textbf{(2)} For $a\in\Phi$ and $u\in U_{-a}^{*}$ there is one and only
one element of $M_a$, written $m(u)$, with $u\in U_a.m(u).U_a$. More precisely
there is exactly one triple $(u',m,u'')\in U_a\times G\times U_a$ with
\[ u=u'mu'',\qquad mU_am^{-1}=U_{-a},\qquad mU_{-a}m^{-1}=U_a, \]
and one then has $m\in M_a$ and $u'\ne 1$.

\begin{sketch}
Given (DR\,4) and (DR\,5) it suffices to show that two triples with the stated
properties coincide. Setting $u_1'^{-1}u'=u_2'\in U_a$ and
$mu_1''u''^{-1}m^{-1}=u_2''\in U_{-a}$, one finds
$u_2'u_2''^{-1}=m_1m^{-1}\in\N(U_a)\cap\N(U_{-a})$, so $u_2'\in\N(U_{-a})$ and
$u_2''\in\N(U_a)$; by (1) both equal $1$. If $u'$ were $1$ one would get
$m=mu''^{-1}m^{-1}.u\in U_{-a}$, hence $U_a=mU_{-a}m^{-1}=U_{-a}$, contradicting
(1).
\end{sketch}

\noindent For $a\in\Phi$ we put $M_a^{0}=m(U_a^{*})=\{m(u)\mid u\in U_a^{*}\}$.

\medskip
\noindent\textbf{(3)} For $a\in\Phi$, $T$ normalizes $U_a$ and $M_a$.

\begin{sketch}
Let $t\in T$ and take $u,u',m,u''$ as in (2). Since $m\in M_a$, also $tm\in M_a$
by (DR\,4), and (DR\,5) shows that $t=(tm)m^{-1}$ normalizes $U_a$; similarly
for $U_{-a}$. Applying (2) to
${}^{t}u={}^{t}u'\cdot{}^{t}m\cdot{}^{t}u''$ gives ${}^{t}m\in M_a$; since $M_a$
is a right coset of $T$, $t$ normalizes $M_a$.
\end{sketch}

\noindent\textbf{(4)} One has $M_a=M_a^{-1}=M_{-a}$, and $M_a=M_{a/2}$ if
$a/2\in\Phi$.

\begin{sketch}
With $u,u',m,u''$ as in (2): if $a/2\in\Phi$ then (2) and (DR\,5) give
$m\in M_{a/2}$, hence $M_a=T.m=M_{a/2}$. In all cases
\[ u^{-1}=u''^{-1}m^{-1}u'^{-1},\qquad u'^{-1}=(mu''m^{-1}).m.u^{-1}, \]
so by (2) $m^{-1}\in M_a$ and $m\in M_{-a}$; with (3) and (DR\,4) this gives
$M_a^{-1}=m^{-1}T\subset M_aT=M_a$ and $M_{-a}=T.m=M_a$.
\end{sketch}

\noindent\textbf{(5)} $T\cup M_a$ is a subgroup of $G$ admitting $T$ as a normal
subgroup of index $2$. (Immediate from (3) and (4).)

\medskip
\noindent\textbf{(6)} For $a\in\Phi$, $\;U_a^{*}.M_a.U_a=U_{-a}^{*}.T.U_a$.
This follows from (DR\,4), (2) and (5) via the chain
\[ U_a^{*}M_aU_a\subset U_{-a}^{*}M_aU_{-a}M_aU_a\subset U_{-a}^{*}TU_a
   \subset U_a^{*}M_aU_aTU_a=U_a^{*}M_aU_a. \]

\medskip
\noindent\textbf{(7)} If $L_a$ denotes the group generated by $U_a$, $U_{-a}$
and $T$, then
\[ L_a=TU_a\cup U_aM_aU_a=M_aU_a\cup U_{-a}TU_a\ne U_{-a}TU_a. \]

\begin{sketch}
The second equality is (6). The first holds because the second and third members
are visibly contained in $L_a$ and invariant under right multiplication by $T$
(by (3)), the second being invariant under right multiplication by $U_a$ and the
third under right multiplication by $U_{-a}$. Finally $L_a=U_{-a}TU_a$ is
impossible: otherwise there would exist $u\in U_a$, $u'\in U_{-a}$, $t\in T$ with
$m=u'tu\in M_a$, whence $U_{-a}=mU_am^{-1}=u'U_au'^{-1}$ and $U_{-a}=U_a$,
contradicting (1).
\end{sketch}

\noindent\textbf{(8)} $\N(U_a)\cap L_a=TU_a$ and
$\N(U_a)\cap\N(U_{-a})\cap L_a=T$. (The first from (7) and (1); the second then
follows using (DR\,6).)

\medskip
\noindent From (8), (DR\,4) and (DR\,5) one gets at once
\begin{equation}
M_a=\{x\in L_a\mid xU_ax^{-1}=U_{-a}\ \text{and}\ xU_{-a}x^{-1}=U_a\}.
\end{equation}
In particular $M_a$ is entirely determined by the data of $T$, $U_a$ and
$U_{-a}$, which allows us to speak, by abuse of language, of \emph{the root datum
$(T,(U_a)_{a\in\Phi})$}.

\medskip
\noindent\textbf{(10)} Let $N$ be the group generated by $T$ and the union of
the $M_a$, $a\in\Phi$. If $\Phi\ne\emptyset$ this is also the group generated by
the $M_a$. By (DR\,5) there is a unique epimorphism $\vnu:N\to\vW$ such that for
all $a\in\Phi$ and all $n\in N$ one has $nU_an^{-1}=U_b$ with $b=\vnu(n)(a)$.
Moreover $\vnu(M_a)=\{r_a\}$ for every $a\in\Phi$.

\medskip
\noindent\textbf{(11)} Since $N$ normalizes $T$ by (8), transitivity of $\vW$ on
the chambers of $\Phi$ shows that condition (DR\,6) remains satisfied if the
chamber used to define $\Phi^{+}$ and $\Phi^{-}$ is replaced by another chamber.

\medskip
\noindent\textbf{(12)} Let $T^{0}$ be the intersection of $T$ with the group
generated by the union of the $M_a^{0}$, $a\in\Phi$ (cf.\ (2)); it is obviously a
normal subgroup of $T$, and even of $N$. Then
$(T^{0},(U_a,T^{0}M_a^{0})_{a\in\Phi})$ is a generating root datum of the group
$G'$ generated by the union of the $U_a$, $a\in\Phi$. More generally, if $X$ is a
subgroup of $T$ normalized by $M_a^{0}$, the pair
$(XT^{0},(U_a,XT^{0}M_a^{0})_{a\in\Phi})$ is a generating root datum of the group
generated by $X$ and $G'$. These assertions are evident.

\begin{stmt}{Examples}{6.1.3}\btanchor{6.1.3}
Let $K$ be a field.
\end{stmt}

\noindent\textbf{a)} Let $\Phi=\{a,-a\}$ and $G=SL_2(K)$. Let $T$ be the subgroup
of diagonal matrices, $U_a$ (resp.\ $U_{-a}$) the subgroup of ``unipotent'' upper
(resp.\ lower) triangular matrices,
$m=\begin{pmatrix}0&-1\\1&0\end{pmatrix}$ and $M_a=M_{-a}=Tm$. One obtains in
this way a generating root datum in $G$. Verification of the axioms is immediate;
for (DR\,4) one uses the relation
\begin{equation}
\begin{pmatrix}1&u\\0&1\end{pmatrix}
=\begin{pmatrix}1&0\\-u^{-1}&1\end{pmatrix}
 \begin{pmatrix}0&u\\-u^{-1}&0\end{pmatrix}
 \begin{pmatrix}1&0\\-u^{-1}&1\end{pmatrix}\qquad(u\in K^{*}),
\end{equation}
from which it also follows, in the notation of (\btref{6.1.2})(2), that
\[ m\!\left(\begin{pmatrix}1&u\\0&1\end{pmatrix}\right)
   =\begin{pmatrix}0&u\\-u^{-1}&0\end{pmatrix}. \]

\medskip
\noindent\textbf{b)} Assume from now on that $K$ is commutative; let $\Gsc$ be a
connected simple algebraic group defined and split over $K$, and $\Tsc$ a maximal
torus of $\Gsc$ defined and split over $K$. Let $\Phi$ be the root system of
$\Gsc$ relative to $\Tsc$. It is known that for each $a\in\Phi$ there is a unique
algebraic subgroup $\Usc_a$ of $\Gsc$ admitting an isomorphism $x_a$ of the
additive group $\Add$ onto $\Usc_a$ satisfying
$t.x_a(u).t^{-1}=x_a(a(t).u)$ for $u\in\Add$, $t\in\Tsc$; choosing such an
isomorphism amounts to choosing a generator $X_a$ of the Lie algebra of $\Usc_a$.
It is known ([\btbib{17}], [\btbib{19}]) that the $x_a$ (or the $X_a$) may be chosen so that:

\begin{enumerate}
\item[(2)] for every $a\in\Phi$ there is a unique epimorphism $\zeta_a$ defined
over $K$ from $SL_2$ onto the algebraic subgroup generated by $\Usc_a$ and
$\Usc_{-a}$, with
$x_a(u)=\zeta_a\!\left(\begin{smallmatrix}1&u\\0&1\end{smallmatrix}\right)$ and
$x_{-a}(u)=\zeta_a\!\left(\begin{smallmatrix}1&0\\u&1\end{smallmatrix}\right)$
for all $u\in\Add$;
\item[(3)] for $a,b\in\Phi$ with $b\ne -a$, and $u,v\in\Add$,
\[ (x_a(u),x_b(v))=\prod_{pa+qb\in\Phi,\;p,q\in\Nstar}
   x_{pa+qb}\bigl(C_{a,b;p,q}\,u^{p}v^{q}\bigr), \]
where the coefficients $C_{a,b;p,q}$ are \emph{integers} (the roots $pa+qb$ being
arranged in a prescribed order).
\end{enumerate}

\noindent This choice, which we assume made, corresponds to the choice of a
``Chevalley basis'' of the Lie algebra of $\Gsc$ in the sense of [\btbib{8}], or again of
a ``pinning'' of $\Gsc$ in the sense of [\btbib{19}].

Now let $G=\Gsc(K)$ be the set of $K$-rational points of $\Gsc$, and likewise
$T=\Tsc(K)$, $U_a=\Usc_a(K)$ and
$M_a=T.\zeta_a\!\left(\begin{smallmatrix}0&-1\\1&0\end{smallmatrix}\right)$
($a\in\Phi$). Then $(T,(U_a,M_a))$ is a generating root datum of type $\Phi$ in
$G$. Conditions (DR\,1) and (DR\,3) are obviously satisfied, (DR\,2) follows from
(3), and (DR\,4) follows from example a). Finally (DR\,5) and (DR\,6) are well
known (\emph{loc.\ cit.}). Note that the group $N$ is then equal to
$\N(\Tsc)(K)$, where $\N(\Tsc)$ denotes the normalizer of $\Tsc$ in $\Gsc$.

\medskip
\noindent\textbf{c)} More generally, let $\Gsc$ be a connected semisimple
algebraic group defined over $K$, $\Ssc$ a maximal $K$-split torus of $\Gsc$
defined over $K$, and $\Ssc'$ a maximal torus of $\Gsc$ defined over $K$ and
containing $\Ssc$; let $\Phi$ (resp.\ $\Phi'$) be the root system of $\Gsc$
relative to $\Ssc$ (resp.\ $\Ssc'$). For $a\in\Phi$, let $\Usc_a$ be the
algebraic subgroup of $\Gsc$ generated by the subgroups $\Usc_{a'}$ for those
$a'\in\Phi'$ whose restriction to $\Ssc$ equals $a$ or $2a$ (here $\Usc_{a'}$ is
the subgroup defined as in b) above, replacing $K$ by an algebraic closure of
$K$, or by an extension of $K$ splitting $\Ssc'$). Let $\Zsc(\Ssc)$ (resp.\
$\N(\Ssc)$) be the centralizer (resp.\ normalizer) of $\Ssc$ in $\Gsc$. Put
$G=\Gsc(K)$, $T=\Zsc(\Ssc)(K)$ and $U_a=\Usc_a(K)$ (for $a\in\Phi$). The results
of [\btbib{3}] then show that $(T,(U_a))$ is a generating root datum of type $\Phi$ in
$G$, for which the group $N$ equals $\N(\Ssc)(K)$.

\para{6.1.4}\btanchor{6.1.4} In all the sequel of \S\,6, $(T,(U_a,M_a))$ denotes a generating
root datum of type $\Phi$ in a group $G$, and the symbols $L_a$, $U^{+}$,
$U^{-}$, $N$, $\vnu$ have the meaning defined in (\btref{6.1.1}) or (\btref{6.1.2}). To simplify
notation we set $U_{2a}=\{1\}$ when $a\in\Phi$ and $2a\notin\Phi$.

\para{6.1.5}\btanchor{6.1.5} Let $\Phi=\bigcup_{i\in I}\Phi_i$ be the decomposition of $\Phi$
into irreducible components. For $i\in I$ let $G_i^{0}$ be the subgroup of $G$
generated by the $U_a$, $a\in\Phi_i$. It is immediate that
$(T\cap G_i^{0},(U_a,M_a\cap G_i^{0})_{a\in\Phi_i})$ is a generating root datum
of type $\Phi_i$ in $G_i^{0}$. The subgroup $T$ normalizes each $G_i^{0}$ and, by
(DR\,2), every element of $G_j^{0}$ for $j\ne i$ centralizes $G_i^{0}$. Since
$(T,(U_a))$ is generating, one sees that each $G_i^{0}$ is a normal subgroup of
$G$ and that $G_i^{0}\cap G_j^{0}$, for $i\ne j$, is contained in the center of
the normal subgroup $G^{0}$ of $G$ generated by the $U_a$ (or by the $G_i^{0}$)
--- a center which is moreover contained in $T$, as we shall see later. Note that
$G=T.G^{0}$.

\begin{stmt}{Proposition}{6.1.6}\btanchor{6.1.6}
The group $U^{+}$ is nilpotent. Moreover, let $\Psi$ be a subset of $\Phi^{+}$
and $\Psi^{\red}=\{a\in\Psi\mid a/2\notin\Psi\}$. For each $a\in\Psi$ let $Y_a$
be a subgroup of $U_a$, and put $X_a=Y_a$ or $Y_a.Y_{2a}$ according as
$2a\notin\Psi$ or $2a\in\Psi$. Let $X$ be the group generated by the $Y_a$, and
assume they have the following property:
\begin{enumerate}
\item[(i)] for $a,b\in\Psi$, $(Y_a,Y_b)$ is contained in the group generated by
the $Y_{ma+nb}$ for $m,n\in\Nstar$ with $ma+nb\in\Psi$.
\end{enumerate}
Then the product map $\prod_{a\in\Psi^{\red}}X_a\to X$ is bijective, whatever the
order in which the factors of the product are arranged.
\end{stmt}
\noindent This is an immediate consequence of (\btref{6.1.2})(11) and of the following
lemma.

\begin{stmt}{Lemma}{6.1.7}\btanchor{6.1.7}
Let $X$ be a group, and $\Psi$, $\Psi^{\red}$ as in (\btref{6.1.6}). For each $a\in\Psi$
let $Y_a$ be a subgroup of $X$ and put $X_a=Y_a$ or $Y_a.Y_{2a}$ according as
$2a\notin\Psi$ or $2a\in\Psi$. Assume the $Y_a$ generate $X$ and possess property
(i) of (\btref{6.1.6}) as well as the following:
\begin{enumerate}
\item[(ii)] for $v\in V$, the intersection of the group generated by the $Y_a$
($a\in\Psi$, $a(v)\le 0$) with the group generated by the $Y_a$ ($a\in\Psi$,
$a(v)>0$) is reduced to the identity element.
\end{enumerate}
Then $X$ is nilpotent and the product map
\begin{equation}
\prod_{a\in\Psi^{\red}}X_a\longrightarrow X
\end{equation}
is bijective, whatever the order in which the factors are arranged.
\end{stmt}

\begin{sketch}
Induction on $\#\Psi$. A maximal element $a_0$ of $\Psi$ for the order attached
to $D_0$ satisfies $ma_0+nb\notin\Psi$, so $Y_{a_0}$ is central by (i); pass to
$X'=X/Y_{a_0}$ and apply the induction hypothesis. Surjectivity of (1) follows
since $Y_{a_0}$ is central. Injectivity: if $\prod x_a=\prod y_a$, the image in
$X'$ gives $x_aX_{a_0}=y_aX_{a_0}$; for $a$ not proportional to $a_0$, hypothesis
(ii) applied at a point $v$ with $a(v)\le0<a_0(v)$ forces $x_a=y_a$, and the
remaining factor then agrees as well. \emph{(Original: p.~111--112.)}
\end{sketch}

\begin{stmt}{Proposition}{6.1.8}\btanchor{6.1.8}
Suppose $\Phi$ has rank $2$. Let $a_1,\dots,a_{2k}$ be the elements of $\Phired$
arranged in ``circular order'', i.e.\ in such a way that
\[ \Phired\cap(\Q_{+}a_{i-1}+\Q_{+}a_{i+1})=\{a_{i-1},a_i,a_{i+1}\}
   \qquad\text{for }1<i<2k. \]
The roots $a=a_1$ and $b=a_k$ then form a basis of $\Phi$, and the corresponding
positive elements of $\Phired$ are the $a_i$ for $1\le i\le k$. Define $a_i$ for
every $i\in\Z$ by the condition $a_{i+2k}=a_i$ for all $i\in\Z$. Let further
$u\in U_a$ and $u'\in U_b$ with $u,u'\ne 1$; put $m=m(u)$, $m'=m(u')^{-1}$ and
$n=mm'$. There exists one and only one sequence $(u_i)_{i\in\Z}$ with
$u_i\in U_{a_i}$, $u_1=u^{-1}$ and $u_k=u'$, such that the two following
conditions hold:
\begin{align}
(u,u')&=u_2\cdots u_{k-1},\\
u_i&\in U_{-a_i}.M_{a_i}.u_{k+i}^{-1}\qquad\text{for all }i\in\Z.
\end{align}
Put $m_i=m(u_i)^{-1}$, so that $m_1=m$ and $m_k=m'$. Then the following relations
hold for all $i\in\Z$:
\begin{align}
m_{k+i}&=m_i;\\
(u_i^{-1},u_{k+i-1})&=u_{i+1}\cdots u_{k+i-2};\\
u_{i+1}&=m_iu_{k+i-1}m_i^{-1};\\
u_2&=mu'm^{-1};\\
m_{i+1}m_i&=n;\\
m_{2i}&=n^{i}m'n^{-i}\quad\text{and}\quad m_{2i+1}=n^{i}mn^{-i};\\
mm'mm'\cdots&=m'mm'm\cdots,
\end{align}
each of the two members of the last relation having $k$ factors.
\end{stmt}

\begin{sketch}
The heart of the argument is the induction on $j\ge0$ over the statement
$(P_j)$: \emph{there is a unique finite sequence $(u_i)_{1\le i\le k+j}$ with
$u_i\in U_{a_i}$, $u_1=u^{-1}$, $u_k=u'$, with $u_i\ne1$ for $1\le i\le j+1$ and
$k\le i\le k+j$, satisfying $(4)$ for $1\le i\le j+1$ and $(2)$, $(5)$ for
$1\le i\le j$.} $(P_0)$ comes from (DR\,2) and (\btref{6.1.6}). The inductive step writes
$u_j=v_{k+j}m_j^{-1}u_{k+j}^{-1}$ and sets
$u'_{j+1}=m_ju_{k+j-1}m_j^{-1}$; commuting $U_{a_{k+j}}$ past $U_{a_{k+j-1}}$
turns the commutator $(u_j^{-1},u_{k+j-1})$ into
$u'_{j+1}\cdot(u'^{-1}_{j+1},u_{k+j})\cdot u_{k+j-1}^{-1}$, and (DR\,2) with
(\btref{6.1.6}) forces $u_{j+1}=u'_{j+1}\ne1$, yielding (5) for $i=j$ and (4) for
$i=j+1$. The symmetric argument (swap $a,b$ and $u,u'$) produces the sequence for
$i\le0$, and the two halves match up. Relations (7)--(9) follow from
$m_{i+1}=m_im_{i-1}m_i^{-1}=nm_{i-1}n^{-1}$, using $m_1=m$, $m_0=m_k=m'$ and
taking $i$ with $2i=k$ or $2i+1=k$. \emph{(Original: pp.~113--114.)}
\end{sketch}

\begin{stmt}{Remarks}{6.1.9}\btanchor{6.1.9}
\end{stmt}
\noindent\textbf{a)} It is easy to see that the proof of (\btref{6.1.8}) remains valid if
one takes a weaker definition of root data, for example that of [\btbib{36}], which
includes the case of the Ree groups of type ${}^{2}F_4$.

\medskip
\noindent\textbf{b)} Let $a\in\Phi$; for $x\in U_a-\{1\}$ write $f(x)$ for the
unique element of $U_a$ with $xM_af(x)\cap U_{-a}\ne\emptyset$. In many cases the
map $f$ is the identity, for example for the canonical root datum of $SL_2(K)$
((\btref{6.1.3})(1)), and more generally for the root datum of the group of rational
points of a connected semisimple algebraic group described in (\btref{6.1.3})~c), when
$a$ is a non-multipliable root ([\btbib{3}], \S\,7).

In all cases, from the relation $u_i=f(u_{k+i})^{-1}m_i^{-1}u_{k+i}^{-1}$ one
draws
\[ f(u_{k+i})^{-1}=u_i.m_i.m_i^{-1}u_{k+i}m_i,\qquad\text{whence}\qquad
   f(u_i)=m_i^{-1}u_{k+i}m_i. \]
From (5) and (7) one then deduces
\[ u_{i+2}=n.(m_i^{-1}u_{k+i}m_i).n^{-1}=n.f(u_i).n^{-1}, \]
whence
\[ u_{2i}=n^{i-1}mf^{\,i-1}(u')m^{-1}n^{1-i}\qquad\text{and}\qquad
   u_{2i+1}=n^{i}f^{\,i}(u^{-1})n^{-i}. \]
If in particular $f$ is the identity for each of the $a_i$, one has
\[ u_{2i+1}=n^{i}u^{-1}n^{-i}\qquad\text{and}\qquad
   u_{2i}=n^{i-1}mu'm^{-1}n^{1-i}. \]

\medskip
\noindent\textbf{c)} Relation (9) is not the only identity between $m$ and $m'$
whose image in $\vW$ gives the well-known relation $(\vnu(mm'))^{k}=1$. For
example, if $\Phi$ is of type $C_2$ and $a_1$ is a long root, then $a_1$ and
$a_3$ are strongly orthogonal and $m_1$, $m_3$ commute. One deduces that
\[ mm'mm'^{-1}=m'mm'^{-1}m. \]
Likewise, if $\Phi$ is of type $G_2$ and $a_1$ is long, then $a_1$ and $a_4$ are
strongly orthogonal, and one deduces that
\[ mm'mm'm^{-1}m'^{-1}=m'mm'mm^{-1}m'^{-1}m. \]

\begin{stmt}{Corollary}{6.1.10}\btanchor{6.1.10}
Let $\Psi$ be a root subsystem of $\Phi$ and $\Delta$ a basis of $\Psi$. For
$a\in\Delta$ let $N_a$ be a non-empty subset of $M_a^{0}$, and let $N_1$ be the
group generated by the $N_a$, $a\in\Delta$. Let $X$ be a subgroup of $T$
normalized by $N_1$ and containing the $N_a^{2}$ for $a\in\Delta$. Then the map
$\lambda$ sending $n\in N_1X$ to the restriction of $\vnu(n)$ to the subspace
generated by $\Psi$ is an epimorphism of $N_1X$ onto the Weyl group $\vW_1$ of
$\Psi$, with kernel $X$.
\end{stmt}

\begin{sketch}
It is immediate that $\lambda^{-1}(1)\supset X$ and that $\lambda(N_1X)=\vW_1$
(since $\vnu(N_a)=\{r_a\}$ for $a\in\Delta$). For every $a\in\Delta$ the image of
$N_a$ in the quotient group $N_1X/X$ reduces to a single element
$\overline{m}_a$ of order $2$, and the $\overline{m}_a$ generate $N_1X/X$. Let
$a,b\in\Delta$ with $a\ne b$; applying (\btref{6.1.8})(9) to the root datum
$(T,(U_c))$, where $c$ runs over the root system formed by the $pa+qb$ with
$p,q\in\Z$ and $pa+qb\in\Phi$, one sees, taking into account the relations
$\overline{m}_a^{2}=\overline{m}_b^{2}=1$, that
$(\overline{m}_a\overline{m}_b)^{k}=1$ with $k$ the order of $r_ar_b$.
Consequently the $\overline{m}_a$, $a\in\Delta$, satisfy the defining relations
of the Weyl group $\vW_1$, and the canonical map $N_1X/X\to\vW_1$ is an
isomorphism.
\end{sketch}

\begin{stmt}{Corollary}{6.1.11}\btanchor{6.1.11}
{\rm(i)} If $\Phi\ne\emptyset$, then $N$ is generated by the union of the $M_a$
for $a\in\Pi$. {\rm(ii)} One has $\vnu^{-1}(1)=T=N\cap TU^{+}$.
\end{stmt}

\begin{sketch}
Let $N'$ be the group generated by the $M_a$, $a\in\Pi$; then $\vnu(N')=\vW$. For
any root $b\in\Phi$ there is $w\in\vW$ such that $w(b)$ is proportional to a
simple root $a\in\Pi$ ([\btbib{5}], ch.~VI, \S\,1, prop.~15); if
$n\in\vnu^{-1}(w)\cap N'$, then by (\btref{6.1.2})(9) and (10)
\[ M_b=n^{-1}M_an\subset N', \]
whence (i). The relation $\vnu^{-1}(1)=T$ follows from (\btref{6.1.10}). Finally, if
$n\in N\cap TU^{+}$, then $w=\vnu(n)$ preserves $\Phi^{+}$ (by (DR\,1), (DR\,6)
and (\btref{6.1.2})(10)), hence is the identity element, and $n\in T$.
\end{sketch}

\begin{stmt}{Proposition}{6.1.12}\btanchor{6.1.12}
Put $R=\{M_a\mid a\in\Pi\}\subset N/T$. Then $(G,TU^{+},N,R)$ is a saturated Tits
system with Weyl group $N/T$ isomorphic to $\vW$.
\end{stmt}

\begin{sketch}
By (\btref{6.1.11}) it suffices to prove the following relations, for $n\in N$, $a\in\Pi$
and $m\in M_a$:
\begin{align}
mTU^{+}m^{-1}&\ne TU^{+},\\
TU^{+}n.\{1,m\}.TU^{+}m&=TU^{+}n.\{1,m\}.TU^{+},\\
\bigcap_{x\in N}xTU^{+}x^{-1}&\subset T.
\end{align}
The first follows immediately from (DR\,1), (DR\,6) and the fact that
$U_{-a}\subset mTU^{+}m^{-1}$. The second being invariant when $n$ is replaced by
$nm$, we need only establish it in the case where $a'=\vnu(n)(a)\in\Phi^{+}$. Let
$U'$ be the group generated by the $U_b$ for $b\in\Phi^{+}$ and $b\notin\R a$. By
(\btref{6.1.6}) one has $U^{+}=U_aU'$, and with the notation of (\btref{6.1.2})(7),
\[ TU^{+}n.\{1,m\}.TU^{+}=TU^{+}n.\{1,m\}.U_aTU'\subset TU^{+}nL_aTU'. \]
But by (\btref{6.1.2})(7),
\[ nL_a\subset nTU^{+}\cup nU_amU_aT=nTU^{+}\cup U_{a'}nmU_aT
   \subset TU^{+}.\{n,nm\}.TU^{+}, \]
whence (2). Finally (3) is a consequence of (DR\,6).
\end{sketch}

\begin{stmt}{Corollary}{6.1.13}\btanchor{6.1.13}
Keep the notation of (\btref{6.1.12}). Then
\begin{align}
\N(U^{+})&=TU^{+},\\
T&=\bigcap_{a\in\Phi}\N(U_a)=\N(U^{+})\cap\N(U^{-}).
\end{align}
\end{stmt}

\begin{sketch}
The normalizer of $U^{+}$ contains $TU^{+}$, so it has the form $TU^{+}N'TU^{+}$
with $N'\subset N$ (\btref{1.2.9}); but if an element $n$ of $N$ normalizes $U^{+}$, its
image under $\vnu$ preserves $\Phi^{+}$ by (DR\,1) and (DR\,6), hence is the
identity element, so that $n\in\vnu^{-1}(1)=T$. Thus $N'\subset T$ and
$\N(U^{+})=TU^{+}$. Relation (2) now follows from the inclusions (for the first,
cf.\ (\btref{6.1.2})(3))
\[ T\subset\bigcap_{a\in\Phi}\N(U_a)\subset\N(U^{+})\cap\N(U^{-})
   =TU^{+}\cap TU^{-}=T. \]
\end{sketch}

\begin{stmt}{Remark}{6.1.14}\btanchor{6.1.14}
The preceding corollary shows that a \emph{generating} root datum is entirely
determined by the system of the $U_a$.
\end{stmt}

\para{6.1.15}\btanchor{6.1.15} Let us note some further consequences of (\btref{6.1.12}).

\medskip
\noindent\textbf{a)} \emph{The injection of $N$ into $G$ defines a bijection of
$N$ (resp.\ $\vW$) onto the set of double cosets $U^{+}\backslash G/U^{+}$
(resp.\ $TU^{+}\backslash G/TU^{+}$).} Indeed, by (\btref{1.2.7}) it suffices to prove
that if $n\in N$ and $t\in T$ are such that $tn\in U^{+}nU^{+}$, then $t=1$. Now
by (\btref{6.1.6}) one has $nU^{+}n^{-1}\subset U^{+}.U^{-}$, and our assertion follows
from (DR\,6).

\medskip
\noindent\textbf{b)} Let $w\in\vW$; the subgroup $nU^{+}n^{-1}\cap U^{-}$ is
independent of $n\in\vnu^{-1}(w)$: it follows from (\btref{6.1.6}) and (DR\,6) that it is
the subgroup generated by the $U_a$ for $a\in\Phi^{-}\cap w(\Phi^{+})$. Denote it
$U_w^{+}$. One then checks easily that, for every $n\in\vnu^{-1}(w)$, the map
$(u,t,u')\mapsto utnu'$ (resp.\ $(u,u')\mapsto unu'$) is a bijection of
$U^{+}\times T\times U^{+}_{w}$ (resp.\ $U^{+}\times U^{+}_{w}$) onto
$U^{+}TnU^{+}$ (resp.\ $U^{+}nU^{+}$).

\medskip
\noindent\textbf{c)} There exists $n\in N$ with $nU^{+}n^{-1}=U^{-}$ and
$nU^{-}n^{-1}=U^{+}$: it suffices to choose $n$ so that $\vnu(n)$ is the element
of $\vW$ carrying the chamber $D$ to $-D$. Consequently, \emph{the injection of
$N$ into $G$ also defines a bijection of $N$ (resp.\ $\vW$) onto
$U^{+}\backslash G/U^{-}$ (resp.\ $TU^{+}\backslash G/TU^{-}$).}

\medskip
\noindent\textbf{d)} Let $X$ be a subset of $\Pi$, let $\vW_X$ be the subgroup of
$\vW$ generated by the $r_a$ for $a\in X$, and put $G_X=U^{+}T\vW_XU^{+}$. Then
$G_X$ is a parabolic subgroup of $G$, and the map $X\mapsto G_X$ is a bijection of
$\Ps(\Pi)$ onto the set of subgroups of $G$ containing $TU^{+}$: this is only a
restatement of (\btref{1.2.9}).

\btsecB{6.2}{6.2. Valuations of a root datum}

The notation of (\btref{6.1.4}) is retained.

\begin{stmt}{Definition}{6.2.1}\btanchor{6.2.1}
A \emph{valuation} of the root datum $(T,(U_a,M_a)_{a\in\Phi})$ is a family
$\varphi=(\varphi_a)_{a\in\Phi}$, where $\varphi_a$ is a map of $U_a$ into
$\R\cup\{\infty\}$, having the following properties.
\end{stmt}

\begin{description}
\item[\rm (V\,0)] for every $a\in\Phi$, the image of $\varphi_a$ contains at
least three elements;
\item[\rm (V\,1)] for $a\in\Phi$ and $k\in\R\cup\{\infty\}$, the set
$U_{a,k}=\varphi_a^{-1}([k,\infty])$ is a subgroup of $U_a$, and
$U_{a,\infty}=\{1\}$;
\item[\rm (V\,2)] for every $a\in\Phi$ and every $m\in M_a$, the function
$u\mapsto\varphi_{-a}(u)-\varphi_a(mum^{-1})$ is constant on
$U^{*}_{-a}=U_{-a}-\{1\}$;
\item[\rm (V\,3)] let $a,b\in\Phi$ and $k,\ell\in\R$; if $b\notin-\R_{+}a$, the
commutator group $(U_{a,k},U_{b,\ell})$ is contained in the group generated by
the $U_{pa+qb,\;pk+q\ell}$ for $p,q\in\Nstar$ with $pa+qb\in\Phi$;
\item[\rm (V\,4)] if $a$ and $2a$ belong to $\Phi$, then $\varphi_{2a}$ is the
restriction of $2\varphi_a$ to $U_{2a}$;
\item[\rm (V\,5)] let $a\in\Phi$, $u\in U_a$ and $u',u''\in U_{-a}$; if
$u'uu''\in M_a$, then $\varphi_{-a}(u')=-\varphi_a(u)$.
\end{description}

\para{6.2.2}\btanchor{6.2.2} Let $\varphi=(\varphi_a)$ be a valuation of $(T,(U_a))$. By
convention we set $U_{2a,k}=\{1\}$ when $a\in\Phi$ and $2a\notin\Phi$
($k\in\R\cup\{\infty\}$). For $a\in\Phi$ and $u\in U_a$,
\begin{equation}
\varphi_a(u)=\sup\{k\in\R\cup\{\infty\}\mid u\in U_{a,k}\}.
\end{equation}
By (V\,1),
\begin{equation}
\varphi_a(u^{-1})=\varphi_a(u).
\end{equation}
It follows that, when (V\,1) holds, axiom (V\,5) is equivalent to
\begin{description}
\item[\rm (V\,5 bis)] let $a\in\Phi$, $u\in U_a$ and $u',u''\in U_{-a}$; if
$u'uu''\in M_a$, then $\varphi_{-a}(u'')=-\varphi_a(u)$.
\end{description}

For $a\in\Phi$ and $k\in\R$ one puts
\[ M_{a,k}=M_a\cap U_{-a}\varphi_a^{-1}(k)U_{-a}. \]
By (\btref{6.1.2})(2), $M_{a,k}$ is non-empty if and only if $k\in\varphi_a(U_a^{*})$,
and
\begin{equation}
\varphi_a^{-1}(k)\subset U_{-a,-k}M_{a,k}U_{-a,-k}\qquad(a\in\Phi,\ k\in\R).
\end{equation}
When $2a\in\Phi$ one has $M_{2a,2k}\subset M_{a,k}$.

It follows from (V\,1) and (V\,3) that the $U_{a,k}$ form a basis of the filter of
neighborhoods of the identity for a topological group structure on $U_a$.

For $a\in\Phi$ one puts
\[ \Gamma_a=\varphi_a(U_a^{*}),\qquad
   \Gamma'_a=\{\varphi_a(u)\mid u\in U_a^{*},\ \varphi_a(u)=\sup\varphi_a(uU_{2a})\}. \]
By (V\,5) and (\btref{6.1.2})(2) one has $\Gamma_{-a}=-\Gamma_a$. If $2a\notin\Phi$ then
$\Gamma'_a=\Gamma_a$. By (V\,1) and (V\,4),
$\Gamma_a=\Gamma'_a\cup(1/2)\Gamma_{2a}$. Note that $\Gamma'_a$ may be empty: for
this it suffices that $U_{2a}$ be dense in $U_a$ (for an example of this
situation, cf.\ (\btref{6.2.3})~d)); this condition is also necessary when $\Gamma_a$ is
discrete in $\R$.

\begin{stmt}{Examples}{6.2.3}\btanchor{6.2.3}
Let $K$ be a field equipped with a non-improper real valuation $\omega$.
\end{stmt}

\noindent\textbf{a)} Let $G=SL_2(K)$ and take up again the root datum in $G$
defined in (\btref{6.1.3})~a). For $u\in K$ put
\[ \varphi_a\!\left(\begin{pmatrix}1&u\\0&1\end{pmatrix}\right)
 =\varphi_{-a}\!\left(\begin{pmatrix}1&0\\u&1\end{pmatrix}\right)=\omega(u). \]
Conditions (V\,0), (V\,1), (V\,3) and (V\,4) are immediately verified. If
$m=\begin{pmatrix}0&x\\y&0\end{pmatrix}\in M_a$, then
$m\begin{pmatrix}1&0\\u&1\end{pmatrix}m^{-1}
 =\begin{pmatrix}1&xuy^{-1}\\0&1\end{pmatrix}$, whence (V\,2). Finally (V\,5)
follows at once from formula (1) of (\btref{6.1.3})~a). We have thus indeed defined a
valuation of the root datum $(T,(U_a))$. Note that
\begin{equation}
M_{a,k}=\left\{\begin{pmatrix}0&u\\-u^{-1}&0\end{pmatrix}\;\middle|\;\omega(u)=k\right\}.
\end{equation}

\medskip
\noindent\textbf{b)} Take up the notation of (\btref{6.1.3})~b). For $a\in\Phi$ and
$u=x_a(t)\in U_a$ ($t\in K$), put $\varphi_a(u)=\omega(t)$. The family
$\varphi=(\varphi_a)$ is a valuation of the root datum $(T,(U_a))$. Here again
(V\,0), (V\,1) and (V\,4) are evident. Condition (V\,3) follows from assertion
(3) of (\btref{6.1.3})~b). Moreover the family
$(\varphi_a\circ\zeta_a,\varphi_{-a}\circ\zeta_{-a})$ is none other than the
valuation of the root datum of $SL_2(K)$ described in example a); one deduces at
once (V\,5), and also (V\,2) when
$m=n_a=\zeta_a\!\left(\begin{smallmatrix}0&-1\\1&0\end{smallmatrix}\right)$.
Since $M_a=Tn_a$, to finish the proof it suffices to show that
$u\mapsto\varphi_a(tut^{-1})-\varphi_a(u)$ is constant on $U_a^{*}$ for every
$t\in T$, which is evident since $tx_a(s)t^{-1}=x_a(a(t)s)$ for $s\in K$,
$t\in T$.

\medskip
\noindent\textbf{c)} One of the goals of this work is to show that if $K$ is a
local field and $\Gsc$ a connected semisimple algebraic group defined over $K$,
then the root datum $(T,(U_a))$ described in (\btref{6.1.3})~c) can be valued: this is
what we shall see in a later chapter (see also \S\,10).

\medskip
\noindent\textbf{d)} Suppose the valued field $K$ commutative, imperfect of
characteristic $2$, and let $L\subsetneq K$ be a $K^{2}$-vector subspace of $K$.
Put $G=SL_2(K)$, $\Phi=\{\pm a,\pm 2a\}$, let $U_{\pm a}$ and $\varphi_{\pm a}$
be defined as in (\btref{6.1.3})~a) and (\btref{6.2.3})~a), put
\[ U_{2a}=\left\{\begin{pmatrix}1&u\\0&1\end{pmatrix}\;\middle|\;u\in L\right\},\qquad
   U_{-2a}=\left\{\begin{pmatrix}1&0\\u&1\end{pmatrix}\;\middle|\;u\in L\right\}, \]
and $\varphi_{\pm2a}=2\varphi_{\pm a}\,|\,U_{\pm2a}$. Then $(T,(U_b)_{b\in\Phi})$
is a root datum of type $\Phi$ (hence of type $BC_1$) in $G$, and
$\varphi=(\varphi_b)$ is a valuation of it. Moreover $U_{2a}$ is dense in $U_a$
(and $\Gamma'_a=\emptyset$) as soon as $L$ is dense in $K$. This is so, for
example, when $K=k((t))$ is a field of formal series over a field $k$ of infinite
degree over $k^{2}$, if one takes for $L$ the group of series $\sum_i a_it^{i}$
such that the $k^{2}$-vector subspace of $k$ generated by the $a_i$ is
finite-dimensional.

\medskip
\noindent\textbf{e)} Take up the notation of (\btref{5.2.30}) and suppose that the root
datum $(T,(U_a)_{a\in\Phi})$ and the system $(N,\nu,(U_\alpha)_{\alpha\in\bSig})$
are \emph{compatible}, in the sense that the following conditions hold:
\begin{enumerate}
\item[(i)] $V=\vA$, $V^{*}=\vv{\mathbf{A}^{*}}$ and $\vv{\bW}=\vW$, which amounts
to saying that there is a surjective map $a\mapsto\bar a$ of $\Phi$ onto $\vSig$
with $\bar a\in\R_{+}a$ for every $a\in\Phi$;
\item[(ii)] the subgroup $N$ coincides with the one also denoted $N$ in (\btref{6.1.2});
\item[(iii)] for every $a\in\Phi$, the subgroup $U_a$ is the union of the
$U_\alpha$ for $\alpha\in\bSig$ with $\vv{\alpha}=a$, i.e.\ is the subgroup
denoted $\Us_a$ in \S\,5.2.
\end{enumerate}
Comparing (\btref{6.1.2})(10) with condition (B\,1) of (\btref{5.2.30}), one sees at once that the
homomorphism $\vnu:N\to\vv{\bW}=\vW$ defined in (\btref{6.1.2}) coincides with the one
also denoted $\vnu$ in \S\,5.2.

Identifying now $\bA$ and $V$ by the choice of an origin $o$ in $\bA$, denote for
$a\in V^{*}$ and $k\in\R$ by $\alpha_{a,k}$ the closed half-space
$\{x\in\bA\mid a(x)+k\ge 0\}$. For $a\in\Phired$ and $u\in U_a$ put
\[ \varphi_a(u)=\sup\{k\in\R\mid u\in U_{\alpha_{a,k}}\}, \]
and for $a\in\Phi-\Phired$ and $u\in U_{2a}$ put $\varphi_{2a}(u)=2\varphi_a(u)$.
Then the family $\varphi=(\varphi_a)_{a\in\Phi}$ so defined is a
\emph{quasi-valuation} of the root datum $(T,(U_a))$, in the sense that it
satisfies (V\,0), (V\,1), (V\,2), (V\,4), (V\,5) and the condition (V\,3 bis),
weaker than (V\,3), obtained by replacing in (V\,3) the hypothesis
``$b\notin-\R_{+}a$'' by ``$b\notin\R a$''. Let us note moreover that part of the
results we shall establish on valuations of a root datum remain valid for
quasi-valuations.

\begin{sketch}
(V\,0), (V\,1), (V\,4) are immediate. For (V\,2), reduce by (V\,4) and
(\btref{6.1.2})(4) to $a\in\Phired$; then $\nu(m)$ is the orthogonal reflection in a
hyperplane $a(x)+\ell=0$, and (B\,1) gives $mU_{a,k}m^{-1}=U_{-a,k-2\ell}$,
whence $\varphi_{-a}(u)=\varphi_a(mum^{-1})-2\ell$.

For (V\,3 bis) one identifies the affine roots containing
$\alpha_{a,k}\cap\alpha_{b,\ell}$ as the $\alpha_{pa+qb,h}$ with
$h\ge pk+q\ell$, applies (\btref{5.2.32}) to get $U_\Omega$ as an ordered product of the
$U_{s,h(s)}$ and (\btref{6.1.6}) to get the analogous product decomposition of the group
$X$ generated by suitable $X_s$, so that
$X\cap U_\Omega=\prod_{pa+qb\in\Phi,\,p,q\in\Nstar}U_{pa+qb,\,pk+q\ell}$;
the commutator $(x,y)$ lies in $U_\Omega$ by (B$'$3 bis) and in $X$ by (DR\,2).

For (V\,5), again with $a\in\Phired$: put $k=\varphi_a(u)$,
$\alpha=\alpha_{a,k}$, use (B\,4) to write $u=u_1'ru_1''$ with
$u_1',u_1''\in U_{-a,-k}$, so $u'=u_1'$ by (\btref{6.1.2})(2) and
$\varphi_{-a}(u')\ge-k$; strict inequality would give
$\alpha_{+}\cap\alpha^{*}=(\alpha^{*})_{+}\cap\alpha$, which is absurd.
\emph{(Original: pp.~119--120.)}
\end{sketch}
\noindent A partial converse will be seen later (\btref{6.5}).

\para{6.2.4}\btanchor{6.2.4} In all the sequel of \S\,6, $\varphi=(\varphi_a)$ denotes a valuation
of the root datum $(T,(U_a))$. When ambiguity about $\varphi$ is to be feared, we
write ${}^{\varphi}U_{a,k}$, ${}^{\varphi}M_{a,k}$, etc.\ for the sets denoted
simply $U_{a,k}$, $M_{a,k}$, etc.\ above; an analogous notation is used without
further warning for the other symbols attached to $\varphi$ defined below.

\para{6.2.5}\btanchor{6.2.5} Let $\lambda:\Phi\to\R^{*}_{+}$ be a function constant on the
irreducible components of $\Phi$, and let $v\in V$. It is easy to see that the
family $\psi=(\psi_a)$ defined by $\psi_a(u)=\lambda(a)\varphi_a(u)+a(v)$ (for
$a\in\Phi$, $u\in U_a$) is a valuation of $(T,(U_a))$, which we denote
$\lambda\varphi+v$. The valuations $\varphi$ and $\psi=\lambda\varphi+v$ are said
to be \emph{equivalent}, and \emph{equipollent} if $\lambda=1$. The map
$(\varphi,v)\mapsto\varphi+v$ is an action of the additive group of $V$ on the set
of valuations, and an equipollence class of valuations is an orbit of $V$ for this
action.

On the other hand let $n\in N$ and $w=\vnu(n)$. For $a\in\Phi$ and $u\in U_a$ one
has $n^{-1}un\in U_{w^{-1}(a)}$ by (\btref{6.1.2})(10). Put this time
\[ \psi_a(u)=\varphi_{w^{-1}(a)}(n^{-1}un). \]
Here again it is immediate --- the inner automorphism defined by $n$ being an
automorphism of the root datum $(T,(U_a))$ --- that the family $\psi=(\psi_a)$ so
defined is a valuation of $(T,(U_a))$, which we denote $n.\varphi$. One thus
obtains an action of the group $N$ on the set of valuations of $(T,(U_a))$, and
one checks at once that for $n\in N$, $v\in V$ and $\lambda$ as above,
\begin{equation}
n.(\lambda\varphi+v)=\lambda(n.\varphi)+\vnu(n)(v).
\end{equation}

\para{6.2.6}\btanchor{6.2.6} Let $A$ be the set of valuations equipollent to $\varphi$: it is an
affine space under $V$, which we equip with the Euclidean distance corresponding
to the given inner product on $V$. For $a\in V^{*}$ and $k\in\R$ we write, as
above, $\alpha_{a,k}$ (or ${}^{\varphi}\alpha_{a,k}$ if confusion is to be feared)
for the closed half-space of $A$ formed by the $x\in A$ with $a(x-\varphi)+k\ge0$.
The boundary of a closed half-space $\alpha$ is denoted $\partial\alpha$. For
$\alpha=\alpha_{a,k}$ with $a\in\Phired$ one puts
\[ U_\alpha=U_{a,k}\qquad\text{and}\qquad U_{\alpha_{+}}=\bigcup_{h>k}U_{a,h}. \]

The \emph{affine roots} of $A$ are the half-spaces $\alpha_{a,k}$ for $a\in\Phi$
and $k\in\Gamma'_a$, and the \emph{walls} of $A$ are the boundaries of the affine
roots. The set of affine roots is denoted $\Sigma$. Finally, $\mathscr{E}$ denotes
the set of pairs $(\alpha,a)\in\Sigma\times\Phi$ such that $\alpha=\alpha_{a,k}$
with $k\in\Gamma'_a$. It is clear that $\Sigma$, $\mathscr{E}$ and the map
$\alpha\mapsto U_\alpha$ remain unchanged if $\varphi$ is replaced by an
equipollent valuation.

\begin{stmt}{Proposition}{6.2.7}\btanchor{6.2.7}
For $i=1,\dots,p$ let $a_i\in\Phi$, $k_i\in\Gamma_{a_i}$ and
$m_i\in M_{a_i,k_i}$. Put $n=m_1\cdots m_p$ and $r_i=r_{a_i}$. Then
\[ n.\varphi=\varphi-v\qquad\text{with}\qquad
   v=\sum_{1\le i\le p}k_i(r_1\cdots r_{i-1})(\co{a_i})\in V. \]
\end{stmt}

\begin{sketch}
Induction on $p$ reduces, via (\btref{6.2.5})(1), to $p=1$ with $a=a_1\in\Phired$; write
$m=u'uu''$, so $\varphi_a(u)=-\varphi_{-a}(u')=-\varphi_{-a}(u'')=k$.
For $b$ not proportional to $a$, the group
$U'=\prod_{c\in\Psi^{\red}}U_{c,h(c)}$ (with $\Psi$ the set of $pa+qb\in\Phi$,
$p\in\Z$, $q\in\Nstar$, and $h(pa+qb)=pk+q\ell$) is normalized by $m$ by (V\,3),
and intersecting with $U_b$, $U_{r(b)}$ gives
$m^{-1}U_{b,\ell}m=U_{r(b),\,\ell-kb(\co{a})}$, i.e.\ (1). The cases $b=-a$ and
$b=a$ are handled from $u''=u^{-1}m(m^{-1}u'^{-1}m)$ and
$u'=(mu''^{-1}m^{-1})mu^{-1}$ together with (V\,2) and $a(\co{a})=2$, giving (2)
and (3); with (V\,4) these say $m.\varphi=\varphi-k\co{a}$.
\emph{(Original: p.~121.)}
\end{sketch}

\begin{stmt}{Corollary}{6.2.8}\btanchor{6.2.8}
Let $\varphi$ and $\psi$ be two valuations of the root datum $(T,(U_a))$. If
$\varphi_a=\psi_a$ for every $a\in\Pi$, then $\varphi=\psi$.
\end{stmt}

\begin{sketch}
One has ${}^{\varphi}M_{a,k}={}^{\psi}M_{a,k}$ for every $a\in\Pi$ and every
$k\in\R$. Let $b\in\Phired$. There is $w\in\vW$ with $c=w^{-1}(b)\in\Pi$, and a
sequence $a_1,\dots,a_p$ of elements of $\Pi$ with $w=r_{a_1}\cdots r_{a_p}$.
Choosing, for $i=1,\dots,p$, an element $k_i\in\Gamma_{a_i}$ and an element
$m_i\in{}^{\varphi}M_{a_i,k_i}={}^{\psi}M_{a_i,k_i}$, and putting
$m=m_1\cdots m_p$, the preceding proposition shows that
\[ \varphi_b(u)-\varphi_c(mum^{-1})=\psi_b(u)-\psi_c(mum^{-1})\qquad(u\in U_b), \]
which proves the assertion.
\end{sketch}

\begin{stmt}{Corollary}{6.2.9}\btanchor{6.2.9}
Let $a,b\in\Phi$ with $b\notin\R a$ and $b-a\notin\Phi$, let
$\Psi=\Phi\cap(\Z a+\Z b)$ be the root subsystem they generate, let
$a=a_1,a_2,\dots,a_s=b,a_{s+1},\dots,a_{2s}$ be the elements of $\Psi^{\red}$
arranged in ``circular order'' (cf.\ (\btref{6.1.8})), let $u\in U_a$, $u'\in U_b$ and
$(u,u')=u_2\cdots u_{s-1}$ with $u_i\in U_{a_i}$. Then, if $i,p,q\in\mathbf{N}$
are such that $2\le i\le s-1$ and $a_i=pa+qb$, one has
\[ \varphi_{a_i}(u_i)=p\,\varphi_a(u)+q\,\varphi_b(u'). \]
\end{stmt}

\begin{sketch}
One may obviously assume $u\ne1\ne u'$. By (\btref{6.1.8}) one has
$u_2=m(u).u'.m(u)^{-1}$, and applying (\btref{6.2.7}) with $n=m(u)$ gives
\begin{equation}
\varphi_{a_2}(u_2)=\varphi_b(u')-b(\co{a})\varphi_a(u).
\end{equation}
Since $a_2=r_a(b)=b-b(\co{a}).a$, this establishes the assertion for $i=2$.

Assume now $i>2$ and argue by induction on $i$. Let $u''\in U^{*}_{-a}$ be defined
by $u\in U_aM_au''$. By axiom (V\,5 bis) of valuations (\btref{6.2.2}),
\begin{equation}
\varphi_{-a}(u'')=-\varphi_a(u).
\end{equation}
On the other hand it follows from (\btref{6.1.8}) that $(u_2^{-1},u'')=u_3\cdots u_{s-1}u'$.
But $a_i=pa+qb=qa_2-(p+q\,b(\co{a})).(-a)$. Hence, by the induction hypothesis and
taking (1) and (2) into account,
\[ \varphi_{a_i}(u_i)=q\,\varphi_{a_2}(u_2)-(p+q\,b(\co{a}))\,\varphi_{-a}(u'')
   =p\,\varphi_a(u)+q\,\varphi_b(u'). \]
\end{sketch}

\begin{stmt}{Proposition}{6.2.10}\btanchor{6.2.10}
Let $A$ be the affine space of valuations equipollent to $\varphi$.
\begin{enumerate}
\item[\rm(i)] $A$ is stable under the action of $N$. For $n\in N$, the map
$\nu(n):\psi\mapsto n.\psi$ of $A$ into itself is an automorphism of the Euclidean
space $A$, whose canonical image in $\operatorname{Aut}V$ equals $\vnu(n)$.
\item[\rm(ii)] For every $a\in\Phi$ and every $k\in\Gamma_a$, the image under
$\nu$ of an element of $M_{a,k}$ is the orthogonal reflection $r_{a,k}$ in the
hyperplane $\partial\alpha_{a,k}=\{x\in A\mid a(x-\varphi)+k=0\}$.
\item[\rm(iii)] For every $n\in N$ and every $\alpha\in\Sigma$ one has
$\nu(n)(\alpha)\in\Sigma$ and $nU_\alpha n^{-1}=U_{\nu(n)(\alpha)}$.
\end{enumerate}
\end{stmt}

\begin{sketch}
(i) By (\btref{6.2.5})(1) and (\btref{6.2.7}), $n.A=A$ for $n$ in the subgroup $N'$ generated by
the $M_{a,k}$; as $N=N'T$ it remains to treat $t\in T$, where (V\,2) gives a
constant $c(t,a)$ with $(t.\psi)_a=\psi_a+c(t,a)$, and the unique $v\in V$ with
$a(v)=c(t,a)$ for $a\in\Pi$ yields $t.\psi=\psi+v$ by (\btref{6.2.8}).
(ii) For $m\in M_{a,k}$ one has $\vnu(m)=r_a$ and, by (\btref{6.2.7}),
$m.(\varphi-\frac12k\co{a})=\varphi-\frac12k\co{a}$.
(iii) Unwinding $u\in nU_\alpha n^{-1}$ gives $u\in U_\beta$ with
$\beta=\{\psi\in A\mid a(n^{-1}.\psi-\varphi)+k\ge0\}=\nu(n)(\alpha)$.
\emph{(Original: pp.~122--123.)}
\end{sketch}

\para{6.2.11}\btanchor{6.2.11} One puts $H=\nu^{-1}(1)$ and $\widehat W=\nu(N)$, and denotes by $W$
the subgroup of $\widehat W$ generated by the reflections $r_{a,k}$ for $a\in\Phi$
and $k\in\Gamma_a$. Since $N$ permutes the $M_{a,k}$, this is a normal subgroup of
$\widehat W$.

Put $N'=\nu^{-1}(W)$, $T'=T\cap N'$, and let $G'$ be the subgroup of $G$ generated
by $N'$ and the $U_a$ for $a\in\Phi$. Since $M_a\cap N'\ne\emptyset$ for
$a\in\Phi$, one sees easily that $(T',(U_a))$ is a generating root datum of type
$\Phi$ in $G'$, for which the group $N$ of (\btref{6.1.2})(10) is none other than $N'$.

\begin{stmt}{Remarks}{6.2.12}\btanchor{6.2.12}
\end{stmt}
\noindent\textbf{a)} Let $\Phi=\bigcup_i\Phi_i$ be the decomposition of $\Phi$
into irreducible components. It is clear that $\varphi_i=(\varphi_a)_{a\in\Phi_i}$
is a valuation of the root datum $(T\cap G_i^{0},(U_a)_{a\in\Phi_i})$ of the group
$G_i^{0}$ defined in (\btref{6.1.5}). On the other hand, to the decomposition of $\Phi$
there corresponds a canonical decomposition of $A$ into a product of affine spaces
$A_i$ under the dual vector space of the subspace generated by $\Phi_i$
(cf.\ (\btref{1.3.9})). It is clear that the action of $N$ is compatible with this product
decomposition of $A$, and that $W$ decomposes as the direct product of the groups
$W_i$ defined from $G_i^{0}$ and $\varphi_i$ as $W$ was defined from $G$ and
$\varphi$. This often allows one to reduce to the case where $\Phi$ is
irreducible.

\medskip
\noindent\textbf{b)} Let $a\in\Phi$ and $u\in U_a$, $u\ne1$. It follows from
(V\,5) and (\btref{6.2.10})(ii) that \emph{the valuation $\varphi_a(u)$ is the unique real
number $k$ such that $\nu(m(u))$ is the orthogonal reflection in the hyperplane
$\partial\alpha_{a,k}$.} One sees therefore that $\varphi$ is completely determined
by the datum of the homomorphism $\nu:N\to\operatorname{Aut}A$.

More precisely, let $A_0$ be an affine space under $V$, let $\nu_0$ be a
homomorphism of $N$ into $\operatorname{Aut}A_0$ whose composite with the canonical
homomorphism $\operatorname{Aut}A_0\to\operatorname{Aut}V$ equals $\vnu$, and let
$\varphi_0$ be a point of $A_0$. For $a$ and $u$ as above, $\nu_0(m(u))$ is an
orthogonal reflection in a hyperplane of $A_0$ of the form
$\{x\in A_0\mid a(x-\varphi_0)+k=0\}$ (with $k\in\R$), since
$\vv{\nu_0}(m(u))=r_a$. Put then $\psi_a(u)=k$ and $\psi_a(1)=\infty$. It is clear
that if there exist a valuation $\varphi$ of $(T,(U_a))$ and an isomorphism $j$ of
${}^{\varphi}A$ onto $A_0$ with $j(\varphi)=\varphi_0$ and
$\nu_0(n)=j\nu(n)j^{-1}$ for all $n\in N$, then $\varphi_a=\psi_a$ for all
$a\in\Phi$.

It is moreover immediate that the family $\psi=(\psi_a)$ always satisfies
conditions (V\,2), (V\,4) and (V\,5); condition (V\,0) is satisfied as soon as
$\nu_0(T)$ generates $V$. To be sure that $\psi$ is indeed a valuation, it remains
only to verify conditions (V\,1) and (V\,3). It is possible that (V\,1) suffices to
entail that $\psi$ is a ``quasi-valuation'' of $(T,(U_a))$, i.e.\ that for a family
$\psi$ of the type considered here, (V\,1) entails condition (V\,3 bis) of
(\btref{6.2.3})~e).

\medskip
\noindent\textbf{Application.} Let $K$ be a field and $G=SL_2(K)$; take up the
notation of (\btref{6.1.3})~a) and suppose the root datum $(T,(U_a,U_{-a}))$ equipped with
a valuation $\varphi$. Replacing $\varphi$ by an equipollent valuation, one may
assume $\varphi_a\!\left(\begin{smallmatrix}1&1\\0&1\end{smallmatrix}\right)=0$.
For $x\in K^{*}$ let $\omega(x)\in\R$ be such that
$\nu\!\left(\begin{smallmatrix}x&0\\0&x^{-1}\end{smallmatrix}\right)=\omega(x)\co{a}$.
One defines in this way a homomorphism $\omega:K^{*}\to\R$. Extend $\omega$ to all
of $K$ by setting $\omega(0)=\infty$. Then \emph{$\omega$ is a non-improper
valuation of $K$.} Indeed, by (\btref{6.1.3})(1) and (\btref{6.2.10})(ii),
$\nu\!\left(\begin{smallmatrix}0&1\\-1&0\end{smallmatrix}\right)$ is the reflection
about $\varphi$; consequently, for $u\in K^{*}$,
\[ \nu\!\left(\begin{pmatrix}0&u\\-u^{-1}&0\end{pmatrix}\right)
 =\nu\!\left(\begin{pmatrix}u&0\\0&u^{-1}\end{pmatrix}\right)\circ
   \nu\!\left(\begin{pmatrix}0&1\\-1&0\end{pmatrix}\right) \]
is the reflection about $\varphi+(1/2)\omega(u)\co{a}$. Using once more (\btref{6.1.3})(1)
and (\btref{6.2.10})(ii), one sees that
$\varphi_a\!\left(\begin{smallmatrix}1&u\\0&1\end{smallmatrix}\right)=\omega(u)$
for all $u\in K$. Condition (V\,1) then shows that
$\omega(u+v)\ge\inf(\omega(u),\omega(v))$ for $u,v\in K$, hence that $\omega$ is a
valuation of $K$, and (V\,0) shows that $\omega$ is not improper.

Taking (\btref{6.2.3})~a) into account, one sees therefore that \emph{the equipollence
classes of valuations of the ``canonical'' root datum of $SL_2(K)$ correspond
bijectively to the non-improper valuations of $K$.} This result will be generalized
in 10.2.

\medskip
\noindent\textbf{c)} Suppose $\Phi$ irreducible and let $\psi$ be another
valuation of the root datum $(T,(U_a))$. \emph{If there exists $a\in\Phi$ with
$\varphi_a=\psi_a$, then the valuations $\varphi$ and $\psi$ are equipollent.}

\begin{sketch}
Reduce to $a\in\Pi\cup2\Pi$ with $u_b\in U_b$ satisfying
$\varphi_b(u_b)=\psi_b(u_b)=0$ for each $b\in\Pi$. By b) it is enough, after
normalizing $\varphi=\psi=0$, to show ${}^{\varphi}\nu={}^{\psi}\nu$. On the
group $N_0$ generated by the $m(u_b)$ both give the reflection $r_{b,0}$, and
$N=N_0.T$; for $t\in T$ one gets $c({}^{\varphi}\nu(t))=c({}^{\psi}\nu(t))$ for
every $c\in{}^{v}\nu(N_0)(a)$, and irreducibility of $\Phi$ makes the $\vW$-orbit
of $a$ span $V^{*}$. \emph{(Original: p.~124.)}
\end{sketch}
\par\endgroup

% ---------------------------------------- bruhattits6.tex
\begingroup\providecommand{\R}{}
\renewcommand{\R}{\mathbf{R}}
\providecommand{\Z}{}
\renewcommand{\Z}{\mathbf{Z}}
\providecommand{\N}{}
\renewcommand{\N}{\mathbf{N}}
\providecommand{\Rt}{}
\renewcommand{\Rt}{\widetilde{\R}}
\providecommand{\Gt}{}
\renewcommand{\Gt}{\widetilde{\Gamma}}
\providecommand{\co}{}
\renewcommand{\co}{^{\vee}}
\providecommand{\vW}{}
\renewcommand{\vW}{{}^{v}W}
\providecommand{\vnu}{}
\renewcommand{\vnu}{{}^{v}\nu}
\providecommand{\vSi}{}
\renewcommand{\vSi}{{}^{v}\Sigma}
\providecommand{\red}{}
\renewcommand{\red}{\mathrm{red}}
\providecommand{\nm}{}
\renewcommand{\nm}{\mathrm{nm}}
\providecommand{\Phired}{}
\renewcommand{\Phired}{\Phi^{\red}}
\providecommand{\Phinm}{}
\renewcommand{\Phinm}{\Phi^{\nm}}
\providecommand{\Ech}{}
\renewcommand{\Ech}{\mathscr{E}}
\providecommand{\Gscr}{}
\renewcommand{\Gscr}{\mathscr{G}}
\providecommand{\Fscr}{}
\renewcommand{\Fscr}{\mathscr{F}}
\providecommand{\Inf}{}
\renewcommand{\Inf}{\operatorname{Inf}}
\providecommand{\Hom}{}
\renewcommand{\Hom}{\operatorname{Hom}}
\providecommand{\cte}{}
\renewcommand{\cte}{\mathrm{const.}}
\providecommand{\para}{}
\renewcommand{\para}[1]{\medskip\noindent\textbf{(#1)}\ \ignorespaces}
\setlength{\parindent}{1.5em}
\bigskip

% =====================================================================
\pmark{125}
\noindent
\textit{d) If $a,b\in\Phi$, there exists a map $g:U_a\to U_b$ such that,
for every valuation $\varphi$ of the root datum $(T,(U_c)_{c\in\Phi})$, one
has}
\[
  \varphi_b(g(u))=-b(a\co)\cdot\varphi_a(u)+\cte
  \qquad(u\in U_a)
\]
(the constant depending of course on $\varphi$).

Indeed, it suffices to choose arbitrarily an element $u_0\in U^{*}_{r_a(b)}$
and to put
\[
  g(u)=\begin{cases}
    1 & \text{if } u=1,\\[2pt]
    m(u)\cdot u_0\cdot m(u)^{-1} & \text{if } u\neq1.
  \end{cases}
\]

\para{6.2.13}\btanchor{6.2.13}
We shall now study the sets $\Gamma_a$ and $\Gamma'_a$ defined in (\btref{6.2.2}).
If $0\in\Gamma_a$, we denote by $\Gt_a$ the subgroup of $\R$ generated by
$\Gamma_a$.

\begin{defn}
The valuation $\varphi$ is said to be \emph{special} if, for every root
$a\in\Phired$, one has $0\in\Gamma_a$.
\end{defn}

\begin{prop}[6.2.14]\btanchor{6.2.14}
In order that $\varphi$ be special it is necessary and sufficient that $0$
belong to $\Gamma_a$ for every $a\in\Pi$. One then has
$\Gamma'_{w(b)}=\Gamma'_b$ for every $w\in\vW$ and every $b\in\Phi$.
\end{prop}

Indeed, (\btref{6.2.7}) implies that, for $a,b\in\Phi$, one has
\begin{equation}
  \Gamma'_{r_a(b)}\supset\Gamma'_b-b(a\co)\,\Gamma_a .
\end{equation}
In particular, $\Gamma'_{r_a(b)}\supset\Gamma'_b$ as soon as
$0\in\Gamma_a$. Since the $r_a$ for $a\in\Pi$ generate $\vW$, and since
every root is the transform of a root belonging to $\Pi\cup2\Pi$ by an
element of $\vW$, this yields the proposition.

\begin{cor}[6.2.15]\btanchor{6.2.15}
There exists a special valuation equipollent to $\varphi$, and even a
valuation $\psi$ equipollent to $\varphi$ such that
$0\in\Gamma_a=\Gamma'_a$ for every root $a$ belonging to the set $\Phinm$
of non-multipliable roots.
\end{cor}

It suffices to take $\psi=\varphi+v$, with $-a(v)\in\Gamma_a$ for every
$a\in(\Pi\cup2\Pi)\cap\Phinm$.

\begin{cor}[6.2.16]\btanchor{6.2.16}
Let $b\in\Phi$. One has $\Gamma'_{-b}=-\Gamma'_b$. If $0\in\Gamma_b$, one
has $\Gamma'_b=\Gamma'_{-b}=\Gamma'_b+2\Gt_b$. If moreover $\Gamma'_b$ is a
discrete subset of $\R$ and if $0\in\Gamma'_b$, then $\Gamma'_b$ is a
subgroup of $\R$ isomorphic to $\Z$.
\end{cor}

The first two assertions follow from (\btref{6.2.14}) and from (\btref{6.2.14})~(1), in
which one replaces $a$ by $b$. If moreover $\Gamma'_b$ is discrete and
contains $0$, then $2\Gt_b$ is contained in $\Gamma'_b$, hence is a
discrete group, isomorphic to $\Z$ in view of (V\,0). The relations
$2\Gt_b\subset\Gamma'_b=\Gamma'_b+2\Gt_b\subset\Gt_b$ then imply
$\Gamma'_b=\Gt_b$ or $\Gamma'_b=2\Gt_b$.

\begin{exam}[6.2.17]\btanchor{6.2.17}
Let $\Gt$ be a subgroup of $\R$ such that $\widetilde{X}=\Gt/2\Gt$ is a
vector space of dimension $>1$ over the field $\Z/2\Z$. Let $X$ be a
generating system of $\widetilde{X}$, containing $0$ and different from the
whole of $\widetilde{X}$, and let $\Gamma$ be the inverse image of $X$ in
$\Gt$. Then $\Gamma$ is not a group, but contains $0$ and is stable under
addition of an element of $2\Gt$, and $\Gt$ is the group generated by
$\Gamma$; it is moreover easy to see that we have just described all the
subsets of $\R$ possessing these properties.
\end{exam}

% =====================================================================
\pmark{126}
On the other hand, let $k$ be a field of characteristic $2$ and let $K_0$ be
the group algebra of the group $\Gt$ with coefficients in $k$. Denote by
$(X^{\gamma})_{\gamma\in\Gt}$ the canonical basis of $K_0$. One obtains a
valuation $\omega$ of $K_0$ with values in $\Gt$ by putting
\begin{equation}
  \omega\Big(\sum_{\gamma}a_{\gamma}X^{\gamma}\Big)
   =\inf\{\gamma\mid a_{\gamma}\neq0\}.
\end{equation}
The completion $K$ of $K_0$ for this valuation is a field, sometimes called
the ``field of formal series over $k$ with exponents in $\Gt$''. The
elements of $K$ are the series $\sum_{\gamma}a_{\gamma}X^{\gamma}$ such
that, for every $\lambda\in\R$, the set of $\gamma\leq\lambda$ with
$a_{\gamma}\neq0$ is finite; the valuation $\omega$ extends to $K$ and (1)
remains valid.

Now let $L$ be the subset of $K$ consisting of the series
$\sum_{\gamma}a_{\gamma}X^{\gamma}$ such that $a_{\gamma}=0$ for
$\gamma\notin\Gamma$. One sees at once that $L$ is a $K^2$-vector subspace
of $K$, containing $K^2$, but is not a subfield of $K$, although $x\in L$,
$x\neq0$, implies $x^{-1}=(x^{-2})x\in L$.

Finally let $G$ be the subset of $\mathrm{SL}_2(K)$ consisting of the
matrices $\begin{pmatrix}a&b\\c&d\end{pmatrix}$ such that the products
$ab$, $ac$, $bd$ and $cd$ belong to $L$. One verifies easily that $G$ is a
subgroup of $\mathrm{SL}_2(K)$. Taking up again the notation of (\btref{6.2.3})~a):
it is immediate that $(T\cap G,(U_a\cap G))$ is a root datum in $G$ and that
$\psi=(\varphi_a\,|\,U_a\cap G)$ is a valuation of it, for which one has
$\Gamma_a=\Gamma'_a=\Gamma$.

We thus see that (\btref{6.2.16}) is ``the best possible result'' in the general
case.

\begin{prop}[6.2.18]\btanchor{6.2.18}
Assume $\varphi$ special. Let $a$ and $b$ be two non-orthogonal roots.
\begin{enumerate}[label=(\roman*),leftmargin=2.2em]
\item If $a$ and $b$ have the same length and $a\neq\pm b$, one has
      $\Gamma_a=\Gamma'_a=\Gamma_b$ and $\Gamma_a$ is a group.
\item If $(b\,|\,b)=2(a\,|\,a)$ and if $0\in\Gamma'_a$, one has
      $\Gamma_b+2\Gt_a\subset\Gamma_b=\Gamma'_b$ and
      $\Gamma'_a+\Gt_b\subset\Gamma'_a$.
\item If $(b\,|\,b)=3(a\,|\,a)$, one has
      $\Gamma_b+3\Gamma_a\subset\Gamma_b=\Gamma'_b\subset\Gamma_a=\Gamma'_a$
      and $\Gamma_a$ and $\Gamma_b$ are groups.
\end{enumerate}
\end{prop}

Under the hypothesis of (i) one has $2a\notin\Phi$, hence
$\Gamma_a=\Gamma'_a$ and $a(b\co)=b(a\co)=\pm1$. Replacing $b$ by $-b$ if
necessary, we may assume that $a(b\co)=-1$. Then $a+b$ is a root of the
same length as $a$ and $b$ and one has $\Gamma_{a+b}=\Gamma_a=\Gamma_b$
(\btref{6.2.14}). On the other hand (\btref{6.2.14})~(1) implies that
$\Gamma_{a+b}\supset\Gamma_a+\Gamma_b$, whence (i).

Under the hypothesis of (ii) one has $2b\notin\Phi$, hence
$\Gamma_b=\Gamma'_b$; moreover, replacing $b$ by $-b$ if necessary, we may
assume that $b(a\co)=2a(b\co)=-2$, and (ii) then follows from (\btref{6.2.14})~(1)
applied to the pairs $(a,b)$ and $(b,a)$.

Under the hypothesis of (iii), the root system generated by $a$ and $b$ is
of type $G_2$ and there exists a root $c$ of the same length as $a$ (resp.
$b$) and not orthogonal to $a$ (resp. $b$). Hence $\Gamma_a=\Gamma'_a$ and
$\Gamma_b=\Gamma'_b$ are groups by (i), and the rest of (iii) is proved as
(ii).

\begin{prop}[6.2.19]\btanchor{6.2.19}
If $\varphi$ is special, the group $W$ (resp.\ $\widehat{W}$) is the
semi-direct product of the stabilizer of $\varphi$ by the subgroup $W\cap V$
(resp.\ $\widehat{W}\cap V$) of the translations of $W$ (resp.\
$\widehat{W}$). Moreover, $W\cap V$ is the group generated by the
translations $ka\co$, with $a\in\Phired$ and $k\in\Gt_a$.
\end{prop}

If $\varphi$ is special, one has $M_{a,0}\neq\emptyset$ for every
$a\in\Phired$, whence $r_{a,0}\in W$. The image in $\vW$
% =====================================================================
\pmark{127}
of the stabilizer of $\varphi$ is therefore the whole of $\vW$, whence the
first assertion. The second follows from (\btref{6.2.14}) and from the fact that
$\vW$ permutes the coroots $a\co$.

\begin{prop}[6.2.20]\btanchor{6.2.20}
Assume $\varphi$ special. One has
\begin{align*}
  W\cap V&=\Big\{\sum_a\gamma_a a\co \;\Big|\; a\in\Pi,\ \gamma_a\in\Gt_a\Big\}\\
  \widehat{W}\cap V&\subset\Big\{\sum_a\gamma_a \varpi_a \;\Big|\; a\in\Pi,\ \gamma_a\in\Gt_a\Big\}
\end{align*}
where $(\varpi_a)$ denotes the basis of $V$ dual to the basis $\Pi$ of
$V^{*}$ (fundamental weights of $(\Phired)\co$).
\end{prop}

For every root $b\in\Phired$ one has $b\co=\sum_{a\in\Pi}n_{b,a}\,a\co$
with $n_{b,a}\in\Z$. Moreover, it is well known that $n_{b,a}$ is even
(resp.\ a multiple of $3$) as soon as $(a\,|\,a)=2(b\,|\,b)$ (resp.\
$(a\,|\,a)=3(b\,|\,b)$). The first assertion then follows from (\btref{6.2.14}),
(\btref{6.2.18}) and (\btref{6.2.19}).

Now let $v\in\widehat{W}\cap V$. Since $W$ is normal in $\widehat{W}$, one
has $v\,r_{a,0}\,v^{-1}r_{a,0}\in V\cap W$ for every $a\in\Pi$, whence
$v-r_a(v)=a(v)a\co\in V\cap W$, which proves the second relation.

\begin{defn}[6.2.21]\btanchor{6.2.21}
One says that $\varphi$ is \emph{discrete} if $\Gamma_a$ is a discrete
subset of $\R$ for every $a\in\Phi$.
\end{defn}

We have already seen (\btref{6.2.16}) that if $\varphi$ is discrete and if
$0\in\Gamma'_a$, then $\Gamma'_a$ is a group isomorphic to $\Z$.

\begin{prop}[6.2.22]\btanchor{6.2.22}
The set $\Phi'=\{a\in\Phi\mid\Gamma'_a\neq\emptyset\}$ is a root system
containing $\Phinm$. Assume $\varphi$ discrete. Then $W$ is an affine Weyl
group, $\Sigma$ is the corresponding affine root system, and $\Ech$ is an
échelonnage of $\Phi'$ by $\Sigma$.
\end{prop}

The first assertion is immediate. Assume $\varphi$ discrete. Replacing it
if necessary by an equipollent valuation, we may assume it special. Each
$\Gt_a$ (for $a\in\Phi'$) is then a group of the form $e_a\Z$ with $e_a>0$
(\btref{6.2.16}), and $V\cap W$ is, by (\btref{6.2.20}), a lattice in $V$. Since $W$ is
generated by orthogonal reflections, it follows that $W$ is an affine Weyl
group ([\btbib{5}], chap.~VI, \S2, no.~5, prop.~8). Moreover, the reflections
$r_{a,k}$ for $a\in\Phi'$ and $k\in\Gamma'_a$ form a generating system of
$W$, invariant under inner automorphisms. Consequently they are all the
reflections of $W$ ([\btbib{5}], chap.~V, \S3, no.~2, cor.\ to th.~1), and $\Sigma$
is exactly the affine root system associated with $W$. That $\Ech$ is an
échelonnage of $\Phi'$ by $\Sigma$ in the sense of (\btref{1.4.1}) is then obvious.

\para{6.2.23}\btanchor{6.2.23}
Assume $\Phi$ irreducible and $\varphi$ discrete, such that $0\in\Gamma'_a$
for every $a\in\Phinm$ (\btref{6.2.15}). One then sees easily that there exists one
and only one real number $e>0$ such that one and only one of the following
conditions is satisfied (the notation is that of \btref{1.4}):

\smallskip
1) The échelonnage $\Ech$ is of type $A_n$ $(n\geq1)$, $B_n$ $(n\geq3)$,
$C_n$ $(n\geq2)$, $D_n$ $(n\geq4)$, $E_6$, $E_7$, $E_8$, $F_4$ or $G_2$;
one has $\Phi'=e\cdot\vSi$ and $\Gamma_a=\Gamma'_a=\Z e$ for every
$a\in\Phi'$.

\smallskip
2) $\Ech$ is of type $B\text{-}C_n$ $(n\geq3)$, $C\text{-}B_n$ $(n\geq2)$ or
$F_4^{\mathrm{I}}$ (resp.\ $G_2^{\mathrm{I}}$); $\Phi'$ is the root system
obtained by multiplying the short roots of $\vSi$ by $e\sqrt2$ (resp.\
$e\sqrt3$) and the
% =====================================================================
\pmark{128}
long roots of $\vSi$ by $e$; one has $\Gamma_a=\Gamma'_a=\Z e$ if $a$ is a
short root of $\Phi'$ and $\Gamma_a=\Gamma'_a=2\Z e$ (resp.\ $3\Z e$) if $a$
is a long root of $\Phi'$.

\smallskip
3) $\Ech$ is of type $B\text{-}BC_n$ $(n\geq3)$, $\Phi'=\Phi$ is of type
$BC_n$, one has $\Phired=e\cdot\vSi$, $\Gamma'_a=\Z e$ for
$a\in\Phired$ and $\Gamma_a=\Gamma'_a=2\Z e$ for $a\in\Phi-\Phired$.

\smallskip
4) $\Ech$ is of type $C\text{-}BC_n^{X}$ $(n\geq1)$ with
$X=\mathrm{I},\mathrm{II},\mathrm{III}$ or $\mathrm{IV}$; $\Phi'=\Phi$ is of
type $BC_n$, $\Phinm=e\cdot\vSi$, one has $\Gamma'_a=\Gamma_a=\Z e$ for
$a\in\Phired\cap\Phinm$; for $a\in\Phired\cap\tfrac12\Phi$ one has

\medskip
\begin{tabular}{@{\hspace{3em}}l@{\hspace{3em}}l@{\hspace{4em}}l}
if $X=\mathrm{I}$   & $\Gamma'_a=(1/2)\Z e$        & $\Gamma_{2a}=2\Z e$\\[4pt]
if $X=\mathrm{II}$  & $\Gamma'_a=(1/2)\Z e$        & $\Gamma_{2a}=\Z e$\\[4pt]
if $X=\mathrm{III}$ & $\Gamma'_a=\Z e+\tfrac12 e$  & $\Gamma_{2a}=2\Z e$\\[4pt]
if $X=\mathrm{IV}$  & $\Gamma'_a=\Z e$ or $\Z e+\tfrac12 e$ & $\Gamma_{2a}=\Z e$.
\end{tabular}
\medskip

\begin{rems}[6.2.24]\btanchor{6.2.24}
1) If $\Ech$ is of type $C\text{-}BC_n^{\mathrm{IV}}$, one can, by replacing
$\varphi$ by an equipollent valuation, reduce to the case where
$\Gamma'_a=\Z e$ for $a\in\Phired\cap\tfrac12\Phi$.

2) One sees that the datum of $\Ech$ determines $\Phi'$ and the system of
the $\Gamma'_a$ up to a multiplicative constant (and up to a
``translation'' on the $\Gamma'_a$); likewise, the datum of $\Phi'$ and of
the $\Gamma'_a$ determines $\Ech$.
\end{rems}

\btsecB{6.3}{6.3. Lemmas concerning rank one}

We keep the hypotheses and notation of 6.2 (see (\btref{6.2.4})). We are going to
prove results which involve only one root $a\in\Phi$ and its integral
multiples (positive or negative), and which therefore in fact express
properties of valued root data of rank one. These results will be
generalized in 6.4: the lemmas (\btref{6.3.2}), (\btref{6.3.3}), (\btref{6.3.6}), (\btref{6.3.9}) and
(\btref{6.3.11}) are particular cases of proposition (\btref{6.4.9}), and lemma (\btref{6.3.7}) is
a particular case of proposition (\btref{6.4.28}), in view of (\btref{6.4.25}).

\begin{lem}[6.3.1]
\btanchor{6.3.1}
Let $a\in\Phi$, $u\in U_a$ and $u'\in U_{-a}$ be such that
$\varphi_a(u)+\varphi_{-a}(u')>0$. There exists one and only one triple
$(u_1,t,u'_1)\in U_a\times T\times U_{-a}$ such that $u'u=u_1tu'_1$.
Moreover one has $t\in H$, $\varphi_a(u_1)=\varphi_a(u)$ and
$\varphi_{-a}(u'_1)=\varphi_{-a}(u')$.
\end{lem}

One cannot have $u'u\in M_aU_{-a}$, for one would then have
$u\in u'^{-1}M_aU_{-a}$ and $\varphi_a(u)+\varphi_{-a}(u')=0$ by (V\,5).
The existence of $(u_1,t,u'_1)$ then follows from (\btref{6.1.2})~(7), and its
uniqueness from (DR\,6). To prove the rest of the lemma it suffices to
treat the case $u\neq1$, $u'\neq1$, which we assume from now on. One then
has $u_1\neq1$ and there exist $u'_2,u'_3\in U_{-a}$ and $m\in M_a$ such
that $u_1=u'_2mu'_3$, with $\varphi_{-a}(u'_2)=-\varphi_a(u_1)$. One then
has
\[
  u=(u'^{-1}u'_2)\cdot mt\cdot(t^{-1}u'_3tu'_1)\in U_{-a}\,mt\,U_{-a}
\]
% =====================================================================
\pmark{129}
whence $\varphi_a(u)=-\varphi_{-a}(u'^{-1}u'_2)$, since $mt\in M_a$. It
follows that
\[
  \varphi_{-a}(u'^{-1}u'_2)<\varphi_{-a}(u')=\varphi_{-a}(u'^{-1}).
\]
Consequently $\varphi_{-a}(u'^{-1}u'_2)=\varphi_{-a}(u'_2)$ and
\[
  \varphi_a(u)=-\varphi_{-a}(u'_2)=\varphi_a(u_1).
\]
Replacing $a$ by $-a$, $u$ by $u'^{-1}$ and $u'$ by $u^{-1}$, one obtains
the equality $\varphi_{-a}(u'_1)=\varphi_{-a}(u')$.

Finally, $m$ and $mt$ both belong to $M_{a,k}$ with
$k=\varphi_a(u)=\varphi_a(u_1)$; one therefore has
$\nu(t)=\nu(m^{-1}mt)=r^2_{a,k}=1$ and $t\in H$.

\begin{lem}[6.3.2]\btanchor{6.3.2}
Let $a\in\Phi$ and $k,\ell\in\Rt$ with $k+\ell>0$. The product
$U_{a,k}HU_{-a,\ell}$ is a group.
\end{lem}

This follows from (\btref{6.3.1}), since $H$ normalizes all the $U_{\alpha}$.

\begin{lem}[6.3.3]\btanchor{6.3.3}
Let $a\in\Phi$, $k\in\Gamma_a$ and $m\in M_{a,k}$. Let $U_{-a,-k+}$ be the
subgroup which is the union of the $U_{-a,\ell}$ for $\ell>-k$. Then
\[
  L_{a,k}=(U_{a,k}\cdot H\cdot U_{-a,-k+})\cup(U_{a,k}\cdot mH\cdot U_{a,k})
\]
is the group generated by $H\cup U_{a,k}\cup U_{-a,-k}$.
\end{lem}

Since $M_{a,k}\subset U_{a,k}U_{-a,-k}U_{a,k}$, it is clear that
$L_{a,k}$ is contained in the group generated by $H$, $U_{a,k}$ and
$U_{-a,-k}$, a group which is also generated by $H$, $U_{a,k}$ and $m$,
since $mU_{a,k}m^{-1}=U_{-a,-k}$ by (\btref{6.2.10}). In view of (\btref{6.3.1}) one has
\[
  L_{a,k}H=L_{a,k}U_{a,k}=L_{a,k}
\]
and it suffices to prove that $L_{a,k}m\subset L_{a,k}$. Since
\[
  U_{a,k}HU_{-a,-k}m=U_{a,k}mHU_{a,k}
\]
it suffices to show that $U_{a,k}mHU_{a,k}m=U_{a,k}HU_{-a,-k}$ is contained
in $L_{a,k}$. But, in view of (\btref{6.2.2})~(3), one has
\[
  U_{-a,-k}\subset U_{-a,-k+}\cup U_{a,k}mHU_{a,k}
\]
whence the lemma.

\begin{rem}[6.3.4]\btanchor{6.3.4}
Let $a\in\Phi$ and let $T_a$ be the inverse image under $\nu$ of the group
of translations with vector proportional to $a\co$. Let $L'_a$ be the group
generated by $U_a$ and $U_{-a}$, and let $L''_a$ be the group generated by
$L'_a$ and $T_a$. Put $T'_a=L'_a\cap T_a$. Let $k\in\Gamma_a$ and let
$m\in M_{a,k}$. One has:
\begin{align}
  L'_a&=U_aT'_a\cup U_amT'_aU_a\\
  L''_a&=U_aT_a\cup U_amT_aU_a.
\end{align}
\end{rem}

Indeed, one has $mT_a=T_am=T_am^{-1}$, since $\nu(m)$ is the reflection
$r_{a,k}$. Consequently
\[
  U_amT_aU_am=U_aT_aU_{-a}\subset U_aT_a\cup
   \Big(\bigcup_{r\in\Gamma_a}U_aT_aM_{a,r}U_a\Big).
\]
But $T_aM_{a,r}=mT_a$, and this shows that the right-hand side of (2) is
stable under right multiplication by $U_a$, $T_a$ and $m$, from which one
deduces at once that it is equal to $L''_a$. Finally, (1) follows
immediately from (2).

% =====================================================================
\pmark{130}
\para{6.3.5}\btanchor{6.3.5}
Let $r,s\in\R$, with $r+s>0$, and let $u\in\varphi_a^{-1}(r)$ and
$u'\in\varphi_{-a}^{-1}(s)$. Put $u=u'_1mu'_2$ and $u'=u_1m'u_2$, with
$m\in M_{a,r}$, $m'\in M_{-a,s}$, $u'_i\in U_{-a,-r}$ and
$u_i\in U_{a,-s}$. We shall evaluate the commutator
$(u',u)=u'uu'^{-1}u^{-1}$. One has
\begin{align}
  (u',u)&\in u_1\cdot m'(u_2,u)m'^{-1}\cdot m'um'^{-1}\cdot U_a\\
  (u',u)&\in U_{-a}\cdot mu'^{-1}m^{-1}\cdot m(u',u'_2)m^{-1}\cdot u'^{-1}_1 .
\end{align}
In view of (\btref{6.2.10}) one has
\[
  m'um'^{-1}\in U_{-a,r+2s}\qquad\text{and}\qquad mu'^{-1}m^{-1}\in U_{a,s+2r}.
\]

Suppose now that $u$ belongs to $U_{2a}$. One then has $(u_2,u)=1$ and one
deduces from (1) that
\[
  (u',u)\in U_{a,-s}\cdot U_{-a,r+2s}\cdot U_a
\]
whence, in view of (\btref{6.3.2}), the relation
\begin{equation}\tag{3}
  (u',u)\in U_{-a,r+2s}\cdot H\cdot U_a .
\end{equation}
On the other hand one also has $u'_2\in U_{-2a}$, whence $(u',u'_2)=1$, and
one deduces likewise from (2) the relation
\begin{equation}\tag{4}
  (u',u)\in U_{-a}\cdot H\cdot U_{a,2r+s}.
\end{equation}
In view of (DR\,6), the two relations (3) and (4) imply ``by intersection''
the relation
\begin{equation}\tag{5}
  (u',u)\in U_{-a,r+2s}\cdot H\cdot U_{a,2r+s}
\end{equation}
and it is immediate that (5) remains valid for $u\in U_{2a,2r}$ and
$u'\in U_{-a,s}$.

\begin{lem}[6.3.6]\btanchor{6.3.6}
Let $k,\ell\in\R$ be such that $k+\ell>0$. The product
\[
  U_{-a,\ell}\cdot H\cdot U_{a,k}\cdot U_{2a,k-\ell}
\]
is a group.
\end{lem}

It is clear that this product is stable under left multiplication by
$U_{-a,\ell}$, $H$ and $U_{a,k}$ (in view of (\btref{6.3.2})). It is also stable
under left multiplication by an element $u$ of $U_{2a,k-\ell}$: indeed, for
$u'\in U_{-a,\ell}$ one has, in view of (\btref{6.3.5})~(5),
\[
  uu'=u'(u'^{-1},u)u\in u'\,U_{-a,\,2\ell+(1/2)(k-\ell)}\,HU_{a,k}u
\]
and $2\ell+(1/2)(k-\ell)>\ell$.

\begin{lem}[6.3.7]\btanchor{6.3.7}
Let $k,\ell\in\R$ with $k+\ell>0$.
\begin{enumerate}[label=(\roman*),leftmargin=2.2em]
\item If $u\in U_{a,k}$ and $u'\in U_{-a,\ell}$, one has
\[
  (u',u)\in U_{-2a,k+3\ell}\cdot U_{-a,k+2\ell}\cdot H\cdot U_{a,2k+\ell}
          \cdot U_{2a,3k+\ell}.
\]
\item If $u\in U_{2a,2k}$ and $u'\in U_{-a,\ell}$, one has
\[
  (u',u)\in U_{-2a,2k+4\ell}\cdot U_{-a,2k+3\ell}\cdot H\cdot U_{a,2k+\ell}.
\]
\end{enumerate}
\end{lem}

Let us prove (i). We may assume $u\neq1\neq u'$. Put $r=\varphi_a(u)$ and
$s=\varphi_{-a}(u')$ and take up again the notation of (\btref{6.3.5}). One has
$(u_2,u)\in U_{2a,r-s}$ and $m'(u_2,u)m'^{-1}\in U_{-2a,r-s+4s}$, and
relation (\btref{6.3.5})~(1) implies
\[
  (u',u)\in U_{a,-s}\cdot U_{-2a,r+3s}\cdot U_{-a,r+2s}\cdot U_a
\]
% =====================================================================
\pmark{131}
whence, in view of lemma (\btref{6.3.6}),
\[
  (u',u)\in U_{-2a,k+3\ell}\cdot U_{-a,k+2\ell}\cdot H\cdot U_a .
\]
One proves in the same way, using (\btref{6.3.5})~(2), that
\[
  (u',u)\in U_{-a}\cdot H\cdot U_{a,2k+\ell}\cdot U_{2a,3k+\ell}
\]
whence (i) (in view of (DR\,6)).

If moreover $u\in U_{2a}$, one deduces from (\btref{6.3.5})~(1), taking into
account the relations $(u_2,u)=1$ and $m'um'^{-1}\in U_{-2a,2r+4s}$, that
$(u',u)\in U_{a,-s}U_{-2a,2r+4s}U_a$, whence, in view of (\btref{6.3.6}),
\[
  (u',u)\in U_{-2a,2k+4\ell}\cdot U_{-a,2k+3\ell}\cdot H\cdot U_a .
\]
On the other hand, (\btref{6.3.5})~(5) implies
\[
  (u',u)\in U_{-a,k+2\ell}\cdot H\cdot U_{a,2k+\ell}
\]
whence (ii).

\para{6.3.8}\btanchor{6.3.8}
Let $a\in\Phi$, $u\in U_a$, $v\in U_{-a}$. By abuse of language, one says
that the commutator $(v,u)$ \emph{admits an $H$-component} if there exists
$h\in H$ such that $(v,u)\in U_{-a}hU_a$; this element $h\in H$ (which is
then unique by (DR\,6)) is called \emph{the $H$-component of the commutator}
$(v,u)$.

\begin{lem}[6.3.9]\btanchor{6.3.9}
Let $k,k',\ell,\ell'\in\R$ be such that $k'\leq2k$, $\ell'\leq2\ell$,
$k\leq k'+\ell$, $\ell\leq\ell'+k$, $k'+\ell'>0$. For
$u\in U_{a,k}\cup U_{2a,k'}$ and $u'\in U_{-a,\ell}\cup U_{-2a,\ell'}$, the
commutator $(u',u)$ admits an $H$-component. If $X$ denotes the group
generated by the $H$-components of these commutators, then the group
generated by
$U_{-2a,\ell'}\cup U_{-a,\ell}\cup U_{a,k}\cup U_{2a,k'}$ is the product
\[
  U_{-2a,\ell'}\cdot U_{-a,\ell}\cdot X\cdot U_{a,k}\cdot U_{2a,k'} .
\]
\end{lem}

The first assertion follows immediately from the hypotheses made on $k$,
$k'$, $\ell$ and $\ell'$ and from lemma (\btref{6.3.2}). To prove the second, it
suffices to show that the product under consideration is stable under left
multiplication by each of its factors, which follows easily from (\btref{6.3.7})
and from the relation $uv=v(v^{-1},u)u$.

\begin{rem}[6.3.10]\btanchor{6.3.10}
Lemma (\btref{6.3.9}) implies in particular that the products considered in
(\btref{6.3.7})~(i) and (ii) are groups, and therefore contain respectively the
commutator groups $(U_{a,k},U_{-a,\ell})$ and $(U_{2a,2k},U_{-a,\ell})$.
\end{rem}

\begin{lem}[6.3.11]\btanchor{6.3.11}
Let $a\in\Phi$ be such that $2a\in\Phi$ and let $k,k'\in\R$ with
$k'\in\Gamma_{2a}$ and $k'<2k$. Let $m\in M_{2a,k'}$ and put
\[
  L_{a,k,k'}=(U_{2a,k'}U_{a,k}HU_{-a,k-k'}U_{-2a,-k'+})\cup              (U_{2a,k'}U_{a,k}mHU_{a,k}U_{2a,k'}).
\]
Then $L_{a,k,k'}$ is the subgroup generated by
$U_{2a,k'}\cup U_{a,k}\cup H\cup U_{-2a,-k'}$.
\end{lem}

One may assume $k'=0$. Taking into account the relation
$mU_{a,k}m^{-1}=U_{-a,k}$ and (\btref{6.3.9}), one sees, arguing as in (\btref{6.3.3}),
that it suffices to show that $L_{a,k,0}m\subset L_{a,k,0}$, and even that
$mHU_{a,k}U_{2a,0}m=HU_{-a,k}U_{-2a,0}$ is contained in $L_{a,k,0}$. Since
one obviously has $HU_{-a,k}U_{-2a,0+}\subset L_{a,k,0}$, it suffices to
show that $HU_{-a,k}u\subset L_{a,k,0}$
% =====================================================================
\pmark{132}
for $u\in\varphi_{2a}^{-1}(0)$. One then has $u\in U_{2a,0}HmU_{2a,0}$. By
(\btref{6.3.9}) (where one takes $k'=0$, $\ell=k$ and $\ell'=2k$) one has
$HU_{-a,k}U_{2a,0}\subset U_{2a,0}U_{a,k}HU_{-a,k}$, whence
\[
  HU_{-a,k}u\subset U_{2a,0}U_{a,k}HU_{-a,k}mU_{2a,0}    =U_{2a,0}U_{a,k}HmU_{a,k}U_{2a,0}\subset L_{a,k,0}
\]
which completes the proof.

\bigskip
\noindent{\large\bfseries 6.4. The groups $U_f$. Extended valuations.}
\medskip

We keep the notation of 6.2 and 6.3.

\para{6.4.1}

\btsecB{6.4}{6.4. The groups \texorpdfstring{$U_{f}$}{U\_f}. Extended valuations}
\btanchor{6.4.1}
Let $\Rt$ be the union of $\R\times\{0,1\}$ and of an element denoted
$\infty$ not belonging to $\R\times\{0,1\}$. We endow $\Rt$ with the
following law of composition: for $r,s\in\R$ one puts
\begin{align*}
  (r,0)+(s,0)&=(r+s,0)\\
  (r,0)+(s,1)&=(s,1)+(r,0)=(r+s,1)\\
  (r,1)+(s,1)&=(r+s,1)\\
  (r,0)+\infty&=(r,1)+\infty=\infty+(r,0)=\infty+(r,1)=\infty .
\end{align*}

We also endow $\Rt$ with the following order relation: for $r,s\in\R$, the
relations $(r,0)\leq(s,0)$, $(r,0)\leq(s,1)$ and $(r,1)\leq(s,1)$ are
equivalent to $r\leq s$, the relation $(r,1)\leq(s,0)$ is equivalent to
$r<s$, and one has $u\leq\infty$ for every $u\in\Rt$. One verifies
immediately that this makes $\Rt$ a \emph{totally ordered commutative
monoid} and that the map $r\mapsto(r,0)$ is an isomorphism of ordered
monoids from $\R$ onto a submonoid of $\Rt$. From now on we shall identify
an element $r\in\R$ with its image $(r,0)$, and we shall denote by $r+$ the
element $(r,1)$ of $\Rt$. Note that one has $r+=\Inf\{s\in\R\mid s>r\}$.
For $t\in\R^{*}_{+}$ and $r\in\R$ one puts $t(r+)=(tr)+$. For $r\in\Rt$ one
puts $0r=0$.

On the other hand, one sees easily that every bounded-below subset of $\Rt$
admits an infimum and that every non-empty subset of $\Rt$ admits a
supremum.

We then extend the definition of $U_{a,k}$ to $k\in\Rt$ (and $a\in\Phi$) by
putting, as in (\btref{6.3.3}),
\[
  U_{a,r+}=\bigcup_{s\in\R,\;s>r}U_{a,s}
\]
for $r\in\R$.

\para{6.4.2}\btanchor{6.4.2}
This paragraph is devoted to the study of certain subgroups of $G$
generalizing the subgroups denoted $P_{\Omega}$ and $U_{\Omega}$ in \S5.2:
in (\btref{5.2.30}), $U_{\Omega}$ was defined, for $\Omega\subset A$, as the group
generated by the $U_{\alpha}$ for $\alpha\supset\Omega$. In other words, in
the case of example (\btref{6.2.3})~e), $U_{\Omega}$ is generated by the
$U_{a,f_{\Omega}(a)}$, where the function $f_{\Omega}:\Phi\to\Rt$ is defined
by
\begin{equation}
  f_{\Omega}(a)=\inf\{k\in\R\mid\Omega\subset\alpha_{a,k}\}\qquad(a\in\Phi).
\end{equation}
As for $P_{\Omega}$, it is equal to $H\cdot U_{\Omega}$.

The generalization is then immediate: let $f$ be a map from $\Phi$ into the
monoid $\Rt$.
% =====================================================================
\pmark{133}
\emph{Throughout the rest of this work we shall denote by $U_f$ the subgroup
of $G$ generated by the union of the subgroups $U_{a,f(a)}$ for $a\in\Phi$.}
Moreover we shall put:
\[
  U_f^{+}=U^{+}\cap U_f,\quad U_f^{-}=U^{-}\cap U_f,\quad   H_f=H\cap U_f,\quad N_f=N\cap U_f,
\]
\[
  U_{f,a}=U_a\cap U_f\qquad(\text{for }a\in\Phi).
\]

When $a\in\Phi$ and $2a\notin\Phi$, it will sometimes be convenient to
assign a value to $f(2a)$: unless expressly stated otherwise, we shall take
$f(2a)=2f(a)$. Recall that we have already made the convention
$U_{2a,k}=\{1\}$ in this case for $k\in\R$; we extend it of course to
$k\in\Rt$.

We denote by $U_f^{(a)}$ the subgroup generated by the union of the
$U_{ia,f(ia)}$ for $i\in\{\pm1,\pm2\}$, and we put, for $a\in\Phi$:
\[
  N_f^{(a)}=N\cap U_f^{(a)}\qquad\text{and}\qquad H_f^{(a)}=H\cap U_f^{(a)}.
\]

It is clear that the subgroups $U_f$, $H_f$ and $N_f$ are normalized by
$H$. Consequently, \emph{for every subgroup $X$ of $H$, the product
$X\cdot U_f$ is a group}; in particular, the subgroups $P_f=H\cdot U_f$ are
the natural generalization of the subgroups $P_{\Omega}$ of \S5.

\para{6.4.3}\btanchor{6.4.3}
If $\Omega$ is a non-empty subset of $A$, the function $f_{\Omega}$ defined
by (\btref{6.4.2})~(1) is ``concave'' in the following sense:

\begin{defn}
A function $f:\Phi\to\Rt$ (resp.\ $f:\Phi\cup\{0\}\to\Rt$) is said to be
\emph{concave} if, for every finite family $(a_i)$ of elements of $\Phi$
(resp.\ $\Phi\cup\{0\}$) such that $\sum_i a_i\in\Phi$ (resp.\
$\sum_i a_i\in\Phi\cup\{0\}$), one has
\begin{equation*}\tag{C}
  f\Big(\sum_i a_i\Big)\leq\sum_i f(a_i).
\end{equation*}
\end{defn}

\para{6.4.4}\btanchor{6.4.4}
Note that every linear function, that is to say every function of the form
$a\mapsto\lambda(a)$ with $\lambda\in V$, is obviously concave: it coincides
moreover with the function $f_{\Omega}$ where $\Omega\subset A$ is the part
reduced to the point of $A$ corresponding to $\lambda$ in the
identification of $V$ with $A$ obtained by the choice of the origin
$\varphi$ in $A$. However, not every concave function, even one with values
in $\R$, is necessarily of the form $f_{\Omega}$ for a subset
$\Omega\subset A$. Besides, the results which we are going to establish
about the groups $U_f$ will apply to functions more general than the
concave functions (cf.\ (\btref{6.4.8})).

\begin{prop}[6.4.5]\btanchor{6.4.5}
In order that a function $f:\Phi\to\Rt$ be concave, it is necessary and
sufficient that the two following conditions be satisfied:
\begin{align*}
  &\text{(C\,1)}\qquad f(a)+f(b)\geq f(a+b)\qquad\text{for }a,b,a+b\in\Phi;\\
  &\text{(C\,2)}\qquad f(a)+f(-a)\geq0\qquad\text{for }a\in\Phi.
\end{align*}
\end{prop}

That (C\,1) and (C\,2) are necessary is obvious. Assume them satisfied and
let $a_1,\dots,a_n\in\Phi$ be such that $a=\sum_i a_i\in\Phi$. Let us show
by induction on $n$ that $f(a)\leq\sum_i f(a_i)$: this follows from (C\,1)
for $n\leq2$. If $n>2$, there exists an index $j\in[1,n]$ such that the
scalar product $(a\,|\,a_j)$ is positive, so that $a-a_j$ is either zero or
a root $a'$.
% =====================================================================
\pmark{134}
In the first case, choose an index $k\in[1,n]$ with $k\neq j$; one then
has, in view of the induction hypothesis and of (C\,2),
\[
  \sum_i f(a_i)=\Big(\sum_{i\neq j,k}f(a_i)\Big)+f(a_k)+f(a_j)
   \geq f(-a_k)+f(a_k)+f(a_j)\geq f(a_j)=f(a).
\]
In the second case, applying the induction hypothesis to the family
$(a_i)_{i\neq j}$, one obtains
\[
  \sum_i f(a_i)\geq f(a')+f(a_j)\geq f(a)
\]
which completes the proof.

\begin{prop}[6.4.6]\btanchor{6.4.6}
Let $f:\Phi\to\Rt$ be a concave function and let
$g:\Phi\to\R\cup\{+\infty\}$ be a function majorising $f$. Put
\[
  f_0=\inf\Big(\sum_i t_i f(a_i)\Big)
\]
where the infimum is extended over the set of pairs consisting of a finite
family $(a_i)$ of elements of $\Phi$ and a finite family $(t_i)$ of
positive real numbers such that $\sum_i t_i a_i=0$ and $\sum_i t_i=1$.
Then one has $f_0\geq0$, and for every real number $c\leq f_0$ there exists
a linear form $\lambda$ on $V^{*}$ such that the restriction of
$\lambda+c$ to $\Phi$ is majorised by $f$. Moreover, if $f(a)+f(-a)>0$ for
every $a\in\Phi$, one has $f_0>0$ and there exist $\mu\in V$ and $k\in\R$,
$k>0$, such that the restriction of $\mu+k$ to $\Phi$ is majorised by $g$.
\end{prop}

Let $(a_i)$ be a finite non-empty family of elements of $\Phi$. Since
$\Phi$ generates a vector subspace over the field of rationals of dimension
equal to the dimension of $V^{*}$, the families $(t_i)$ of rational numbers
such that $\sum_i t_ia_i=0$ are dense in the set of families $(t_i)$ of
real numbers such that $\sum_i t_ia_i=0$. To show that $f_0\geq0$ it
therefore suffices to show that a relation $\sum_i t_ia_i=0$ implies
$\sum_i t_if(a_i)\geq0$ as soon as the $t_i$ are positive rational numbers,
or even only positive integers. In other words, we are reduced to showing
that if $(a_i)$ is a finite non-empty family of elements of $\Phi$ such
that $\sum_i a_i=0$, then $\sum_i f(a_i)\geq0$. Choose an index $i_0$; one
then has
\begin{equation}
  \sum_i f(a_i)=\Big(\sum_{i\neq i_0}f(a_i)\Big)+f(a_{i_0})
   \geq f(-a_{i_0})+f(a_{i_0})\geq0
\end{equation}
by (C) and (C\,2).

To prove the existence of $\lambda$, it suffices to show that the
intersection of the half-spaces $L_a=\{x\in V\mid a(x)+c\leq f(a)\}$, for
$a\in\Phi$ and $f(a)<\infty$, is non-empty. Now it is known ([\btbib{4}],
chap.~II, \S2, exerc.~21) that, in order for the intersection of a finite
family $\Fscr$ of convex subsets of a real vector space of dimension $n$ to
be non-empty, it suffices that the intersection of every subfamily of $n+1$
elements of $\Fscr$ be non-empty. Moreover, if the elements of $\Fscr$ are
half-spaces, repeated application of this result shows that it suffices
that the intersection of every finite subfamily consisting of $r+1$
half-spaces whose corresponding linear forms constitute a system of rank
$r$ (with $1\leq r\leq n$) be non-empty. In other words, it suffices for us
to verify that if $a_1,\dots,a_r$ are linearly
% =====================================================================
\pmark{135}
independent roots and if $a_{r+1}=\sum_{1\leq i\leq r}t_ia_i\in\Phi$ (with
$t_i\in\R$), then $\bigcap_{1\leq i\leq r+1}L_{a_i}$ is non-empty. Now, if
this intersection were empty, that would mean that the relations
$a_i(x)\leq f(a_i)-c$ for $1\leq i\leq r$ imply
$\sum_{1\leq i\leq r}t_ia_i(x)>f(a_{r+1})-c$. This is possible only if
$t_i\leq0$ for $1\leq i\leq r$. For $a\in\Phi$ with $f(a)<\infty$, let
$f^{-}(a)\in\R$ be such that $f(a)=f^{-}(a)$ or $f^{-}(a)+$. Choose
$x\in A$ such that $a_i(x)+c=f^{-}(a_i)$ for $1\leq i\leq r$, which is
legitimate since the $a_i$ are independent. By definition of $f_0$ one has:
\[
  f^{-}(a_{r+1})-\sum_i t_if^{-}(a_i)\geq\Big(1-\sum_i t_i\Big)c
\]
whence
$a_{r+1}(x)+c=\sum_i t_if^{-}(a_i)+(1-\sum_i t_i)c\leq f^{-}(a_{r+1})\leq
f(a_{r+1})$, which completes the proof of the existence of $\lambda$.

Now suppose that $f(a)+f(-a)>0$ for every $a\in\Phi$. In view of (1) one
has $f_0\geq0+$ and there exists $\lambda\in V$ such that $f\geq\lambda$.
Since $f$ is concave, the set $\Psi$ of the $a\in\Phi$ such that
$f(a)=\lambda(a)$ is a closed subset of $\Phi$ such that
$\Psi\cap(-\Psi)=\emptyset$. There therefore exists $\nu\in V$ such that
$\nu(a)>0$ for $a\in\Psi$ ([\btbib{5}], chap.~VI, \S1, no.~7, prop.~22). Put
$h=\inf\{g(a)-\lambda(a)\mid a\notin\Psi\}$,
$d=\sup\{\nu(a)\mid a\in\Phi\}$ and
$\delta=\inf\{\nu(a)\mid a\in\Psi\}>0$. Then $g$ majorises the restriction
to $\Phi$ of the function $\alpha h+\lambda-\beta\nu$ as soon as the real
numbers $\alpha,\beta>0$ satisfy the inequalities
$\alpha h/\delta\leq\beta\leq(1-\alpha)/d$.

\begin{prop}[6.4.7]\btanchor{6.4.7}
Let $f:\Phi\to\Rt$ be a concave function. Then $f$ possesses the two
following properties:
\smallskip

\noindent\textnormal{(QC\,1)} \emph{For every root $a\in\Phi$ one has}
\[
  U_f^{(a)}=U_{-2a,f(-2a)}\cdot U_{-a,f(-a)}\cdot U_{a,f(a)}\cdot             U_{2a,f(2a)}\cdot N_f^{(a)} .
\]

\noindent\textnormal{(QC\,2)} \emph{If $a,b\in\Phi$ are not proportional,
the commutator group $(U_{a,f(a)},U_{b,f(b)})$ is contained in the group
generated by the $U_{pa+qb,\,f(pa+qb)}$ for $p,q\in\N^{*}$ and
$pa+qb\in\Phi$.}
\end{prop}

(QC\,2) follows from axiom (V\,3) of valuations (\btref{6.2.1}), taking into
account condition (C), which implies $f(pa+qb)\leq pf(a)+qf(b)$ for
$a,b,pa+qb\in\Phi$ with $p,q\in\N^{*}$.

Let us prove (QC\,1). Since $f$ is concave, we have, in view of (C\,1), the
relations $f(-2a)\leq2f(-a)$, $f(2a)\leq2f(a)$, $f(a)\leq f(2a)+f(-a)$ and
$f(-a)\leq f(-2a)+f(a)$; from these relations and from (C\,2) one deduces
\[
  0\leq f(2a)+f(-2a)\leq2\big(f(a)+f(-a)\big).
\]

If $f(2a)+f(-2a)>0$, property (QC\,1) follows immediately from (\btref{6.3.9})
(which shows moreover that $N_f^{(a)}=H_f^{(a)}$). If $f(a)+f(-a)=0$, it
follows from the preceding relations that $f(2a)=2f(a)$ and
$f(-2a)=2f(-a)$, and $U_f^{(a)}$ is generated by
$U_{a,f(a)}\cup U_{-a,f(-a)}$; property (QC\,1) then follows at once from
(\btref{6.3.3}) if $f(a)\in\Gamma_a$, and from (\btref{6.3.2}) in the contrary case.

Finally suppose that $f(a)+f(-a)>0$ and that $f(2a)+f(-2a)=0$. To simplify
the notation we shall assume that $f(2a)=f(-2a)=0$, which is legitimate up
to replacing $\varphi$ by an equipollent valuation. One then has
$k=f(a)=f(-a)>0$. If $0\in\Gamma_{2a}$, (QC\,1) follows from (\btref{6.3.11});
otherwise, the group $U_f^{(a)}$ is the union of the groups generated
% =====================================================================
\pmark{136}
by the unions $U_{-2a,r}\cup U_{-a,k}\cup U_{a,k}\cup U_{2a,0}$ for
$0<2r\leq k$, and it suffices once again to apply (\btref{6.3.9}).

\begin{defn}[6.4.8]\btanchor{6.4.8}
A function $f:\Phi\to\Rt$ is said to be \emph{quasi-concave} if it satisfies
conditions (QC\,1) and (QC\,2) of (\btref{6.4.7}).
\end{defn}

We have just seen that every concave function is quasi-concave. In certain
cases the converse is true. Let us take up again, for example, the notation
of example (\btref{6.2.3})~b): $G$ is thus the group of rational points of a
connected semi-simple algebraic group $\Gscr$ defined and split over a
valued field $K$. Suppose moreover that $\Gscr$ is simple. The Chevalley
commutation relations then imply that every quasi-concave function is
concave, except when $\Gscr$ is of type $B_n$, $C_n$ or $F_4$ and the
characteristic of the residue field of $K$ is equal to $2$, or when
$\Gscr$ is of type $G_2$ and the characteristic of the residue field of $K$
is equal to $2$ or to $3$. In these exceptional cases it is easy to
construct quasi-concave functions which are not concave.

\begin{prop}[6.4.9]\btanchor{6.4.9}
Let $f:\Phi\to\Rt$ be a quasi-concave function.
\begin{enumerate}[label=(\roman*),leftmargin=2.2em]
\item $U_{f,a}=U_{a,f(a)}\cdot U_{2a,f(2a)}$ for every $a\in\Phi$.
\item The product map
      $\prod_{a\in\Phi^{\pm\,\red}}U_{f,a}\to U_f^{\pm}$
      is bijective, whatever the order in which the factors of the product
      are arranged.
\item One has $U_f=U_f^{-}\cdot U_f^{+}\cdot N_f$.
\item $N_f$ is the group generated by the $N_f^{(a)}$ for $a\in\Phi$.
\end{enumerate}
\end{prop}

For $a\in\Phi$, put $X_a=U_{a,f(a)}\cdot U_{2a,f(2a)}$. In view of (QC\,2)
and (\btref{6.1.6}), the product of the $X_a$ for $a\in\Phi^{\pm\,\red}$, arranged
in an arbitrary order, is a group, which we shall denote $X^{\pm}$. On the
other hand let $Y$ be the group generated by the $N_f^{(a)}$ for $a\in\Phi$.

We shall first show that the product $X^{-}X^{+}Y$ is independent of the
choice of the Weyl chamber $D$ used to define $\Phi^{+}$ and $\Phi^{-}$.
For this, it suffices to show that this product does not change if one
replaces $D$ by the chamber $r_a(D)$, where $r_a$ is the reflection
associated with a root $a$ simple for $D$. Denote then by $X'_a$ (resp.\
$X'_{-a}$) the product of the $X_b$ for $b\in\Phi^{+\,\red}$, $b\neq a$
(resp.\ $b\in\Phi^{-\,\red}$, $b\neq-a$). In view of (QC\,2) and (\btref{6.1.6}),
$X'_{\pm a}$ is a group and is normalized by $X_a$ and $X_{-a}$, hence by
the group they generate, that is to say by
$U_f^{(a)}=X_{-a}X_aN_f^{(a)}$ ((QC\,1)). Consequently one has
\begin{align*}
  X^{-}X^{+}Y&=X'_{-a}X_{-a}X_aX'_aY=X'_{-a}X'_aX_{-a}X_aY
   =X'_{-a}X'_aX_aX_{-a}N_f^{(a)}Y\\
  &=X'_{-a}X_aX'_aX_{-a}Y
\end{align*}
which proves our assertion.

It follows at once that $X^{-}X^{+}Y$ is stable under left multiplication
by all the $X_a$ for $a\in\Phired$, and consequently that
\[
  U_f=X^{-}X^{+}Y.
\]

Now let $x^{+}\in X^{+}$, $x^{-}\in X^{-}$ and $y\in Y$. If
$x^{-}x^{+}y\in U^{-}$, one has $x^{+}y\in U^{-}$. Since $Y\subset N$, one
deduces that $y=1$ ((\btref{6.1.15})~c)), whence
$x^{+}\in U^{+}\cap U^{-}=\{1\}$ (DR\,6). One
% =====================================================================
\pmark{137}
concludes that $U_f^{-}=X^{-}$, and one shows in the same way that
$U_f^{+}=X^{+}$: this proves (ii), and (i) follows from it by intersection
with $U_a$. Finally, if $x^{-}x^{+}y\in N$, one has, again by (\btref{6.1.15})~c),
$y=x^{-}x^{+}y$, which proves (iv); and (iii) then follows from the
relation $U_f=X^{-}X^{+}Y$.

\begin{prop}[6.4.10]\btanchor{6.4.10}
Let $f:\Phi\to\Rt$ be a quasi-concave function. For $a\in\Phi$ put
\[
  f'(a)=\inf\{k\in\Gamma'_a\mid k\geq f(a)\ \text{ or }\    a/2\in\Phi\ \text{ and }\ k\geq2f(a/2)\}.
\]
Let $\Phi_f$ be the set of the $a\in\Phi$ such that $f'(a)+f'(-a)=0$. Then
$\Phi_f$ is a root system in the space $V^{*}_f$ which it generates, and the
map which to $n\in N_f$ associates the restriction of $\vnu(n)$ to
$V^{*}_f$ is a surjective homomorphism from $N_f$ onto the Weyl group of
$\Phi_f$, with kernel $H_f$.
\end{prop}

In view of (\btref{6.4.9})~(i) one has, for $a\in\Phi$,
\[
  f'(a)=\inf\{k\in\Gamma'_a\mid U_{a,k}\subset U_f\}.
\]

Consequently the function $f'$ is in a certain sense ``covariant'' under
$N_f$. More precisely, let $n\in N_f$, $a\in\Phi$ and $b=\vnu(n)(a)$. Then
\begin{equation}
  f'(b)=f'(a)+a\big(\nu(n)^{-1}(\varphi)-\varphi\big).
\end{equation}

This shows in particular that $\Phi_f$ is stable under $\vnu(N_f)$. It
therefore suffices for us to prove that $\vnu(N_f)$ is generated by the
reflections $r_a$ for $a\in\Phi_f$. In view of (\btref{6.4.9})~(iv), this follows
from the following lemma:

\begin{lem}[6.4.11]\btanchor{6.4.11}
Let $f$ and $f'$ be as in (\btref{6.4.10}) and let $a\in\Phi$.
\begin{enumerate}[label=(\roman*),leftmargin=2.2em]
\item One has $f'(a)+f'(-a)\geq0$ and $2f'(a)+f'(-2a)\geq0$.
\item The following conditions are equivalent:
\begin{enumerate}[label=\alph*),leftmargin=2.2em]
\item $N_f^{(a)}\not\subset H$;
\item there exists $k\in\R$ such that
      $\nu(N_f^{(a)})=\{1,r_{a,k}\}=\{1,r_{2a,2k}\}$;
\item $\vnu(N_f^{(a)})=\{1,r_a\}=\{1,r_{2a}\}$;
\item $a$ or $2a$ belongs to $\Phi_f$.
\end{enumerate}
\end{enumerate}
\end{lem}

Let us first prove the equivalence of a), b) and c). It is clear that b) or
c) implies a). On the other hand, it is known (\btref{6.3.4}) that
$\nu(N_f^{(a)})$ consists of reflections $r_{a,k}$ and of translations
$ka\co$ with $k\in\R$. If there existed $k>0$ and $n\in N_f^{(a)}$ such
that $\nu(n)=ka\co$, one would have, for every integer $p>0$:
\[
  n^{p}U_{a,f(a)}n^{-p}=U_{a,f(a)-2pk}\subset U_f
\]
which is possible only if $f(a)=\infty$; and one sees likewise that one
would have $f(2a)=\infty$. But this is absurd, for it implies
$U_{f,a}\subset U_{-a}$ and $N_f^{(a)}=\{1\}$. Consequently
$\nu(N_f^{(a)})$ contains only reflections (besides $1$), and cannot
contain two distinct reflections whose product would be a non-zero
translation. Whence the equivalence of a), b) and c).

Now suppose $f'(a)+f'(-a)<0$. There then exist $k\in\Gamma'_a$ and
$h\in\Gamma'_{-a}$ such that
\[
  f'(a)\leq k,\qquad f'(-a)\leq h\qquad\text{and}\qquad k+h<0.
\]
% =====================================================================
\pmark{138}
The group $U_f$ then contains $U_{a,k}\cup U_{-a,-k}$ and
$U_{a,-h}\cup U_{-a,h}$, hence $M_{a,k}$ and $M_{a,-h}$, and
$\nu(N_f^{(a)})$ contains the two distinct reflections $r_{a,k}$ and
$r_{a,-h}$. But we have just seen that this is impossible, whence the first
assertion of (i). Likewise, if $2f'(a)+f'(-2a)<0$, there exist
$k\in\Gamma'_a$ and $h\in\Gamma'_{-2a}$ such that $f'(a)\leq k$,
$f'(-2a)\leq h$ and $2k+h<0$, and one sees as above that
$\nu(N_f^{(a)})$ contains the two distinct reflections $r_{2a,2k}$ and
$r_{2a,-h}$, which completes the proof of (i). Note moreover that an
analogous argument shows that if $2f'(a)+f'(-2a)=0$, then $a$ and $2a$
belong to $\Phi_f$.

If $a\in\Phi_f$, one has $k=f'(a)=-f'(-a)\in\R$ by the formulae defining
addition in $\Rt$ (\btref{6.4.1}). Consequently $k\in\Gamma'_a$ and
$M_{a,k}\subset N_f^{(a)}$. One sees likewise that if $2a\in\Phi_f$, there
exists $k\in\Gamma'_{2a}$ such that $M_{2a,k}\subset N_f^{(a)}$. This shows
that d) implies a).

Finally, suppose a), b) and c) satisfied. Up to replacing $\varphi$ by an
equipollent valuation, we may assume that $r_{a,0}\in\nu(N_f^{(a)})$. By
(\btref{6.4.10})~(1) one then has
\[
  f'(a)=f'(-a)\qquad\text{and}\qquad f'(2a)=f'(-2a).
\]
In view of (i) it follows that $f'(a)\geq0$ and $f'(2a)\geq0$. Finally, if
$f'(a)$ and $f'(2a)$ were both strictly positive, $U_{f,a}$ would be
contained in the union of the groups generated by
$U_{a,r}\cup U_{-a,-r}$ for $r\in\R$, $r>0$, and one would have
$N_f^{(a)}\subset H$ by (\btref{6.3.2}), which contradicts c). Consequently d) is
indeed satisfied.

The proof of the lemma and that of proposition (\btref{6.4.10}) are complete.

\begin{rems}[6.4.12]\btanchor{6.4.12}
a) Let us keep the notation of (\btref{6.4.10}). For $a\in\Phi$ put
\[
  f''(a)=\inf\{k\in\Gamma_a\mid k\geq f(a)\ \text{ or }\    a/2\in\Phi\ \text{ and }\ k\geq2f(a/2)\}.
\]
One has $f''(a)\leq f'(a)$, and $f''(a)=f'(a)$ if $\Gamma'_a=\Gamma_a$, for
example if $2a\notin\Phi$. One shows exactly as above that (\btref{6.4.11})~(i) is
also valid for $f''$; in other words one has, for $a\in\Phi$:
\begin{equation}
  f''(a)+f''(-a)\geq0\qquad\text{and}\qquad 2f''(a)+f''(-2a)\geq0.
\end{equation}

On the other hand it is clear that $U_f=U_{f''}$. On the contrary, one does
not always have $U_f=U_{f'}$. More precisely, for $a\in\Phi$ one has
\begin{equation}
  U_{a,f'(a)}\cdot U_{2a,f'(2a)}\subset U_{f,a}\subset
   U_{a,\,\inf\{f'(a),\frac12 f'(2a)\}}
\end{equation}
the first inclusion being an equality as soon as $\Gamma'_a\neq\emptyset$.

b) Suppose $\varphi$ discrete and take up again the notation of (\btref{6.2.22})
sqq. Let $\Omega$ be a subset of the affine space $A$ and consider the
corresponding concave function $f_{\Omega}$ and group
$U_{\Omega}=U_{f_{\Omega}}$ (\btref{6.4.2}). It is immediate that, for every
$a\in\Phi$, $f'_{\Omega}(a)$ is the smallest real number $k$ such that the
half-space $\alpha_{a,k}=\{x\in A\mid a(x-\varphi)+k\geq0\}$ is an affine
root of the system $\Sigma$ associated with $a$ by the échelonnage $\Ech$
and containing $\Omega$. Consequently, \emph{the root system
$\Phi_{\Omega}=\Phi_{f_{\Omega}}$ consists of the roots $a\in\Phi$ for which
there exists $k\in\R$ such that $(\alpha_{a,k},a)\in\Ech$ and
$\Omega\subset\partial\alpha_{a,k}$.}

Suppose in particular that $\Omega$ is a facet $F$ of $A$. Then $\Phi_F$ is
exactly
% =====================================================================
\pmark{139}
\emph{the root system attached to $F$ by $\Ech$ in the sense of (\btref{1.4.2})
and, by (\btref{1.4.5})~c), the Dynkin diagram of $\Phi_F$ is obtained by deleting
from the Dynkin diagram of the échelonnage $\Ech$ the vertices not
belonging to the type $X$ of $F$.}
\end{rems}

\para{6.4.13}\btanchor{6.4.13}
In what follows we shall define certain normal subgroups of $H$ and study
the intersection $H_f$ of $H$ with a subgroup $U_f$. Let $k\in\Rt$; we
shall denote by $H_{(k)}$ the set of the $h\in H$ such that
\[
  (h,U_{a,r})=\{(h,u)\mid u\in U_{a,r}\}\subset U_{a,r+k}\cdot U_{2a,2r+k}
\]
for every $a\in\Phi$ and every $r\in\R$ (or $r\in\Rt$). One verifies
immediately that the $H_{(k)}$ form a decreasing family of normal subgroups
of $H$, and even of $N$. One has $H_{(k)}=H$ for $k\leq0$, and
$H_{(\infty)}$ is, in view of (\btref{6.1.13}), the centralizer of the subgroup of
$G$ generated by the $U_a$ for $a\in\Phi$.

\para{6.4.14}\btanchor{6.4.14}
Let again $k\in\Rt$. We shall define another subgroup of $H$, denoted
$H_{[k]}$. For $k\leq0$, $H_{[k]}$ is the subgroup generated by the union
of the intersections with $H$ of the subgroups generated by
$U_{a,r}\cup U_{-a,-r}$ for $a\in\Phi$ and $r\in\R$. For $0<k<\infty$,
$H_{[k]}$ is the subgroup generated by the $H$-components (\btref{6.3.8}) of the
commutators $(u,u')$ and $(u',u)$ with $u\in U_{a,r}$ and $u'\in U_{-a,s}$,
$r+s=k$, or else $u\in U_{a,r}$ and $u'\in U_{-2a,s}$, $2r+s=k$ (with $a$
running through $\Phi$). Here again, the $H_{[k]}$ form a decreasing family
of normal subgroups of $H$, and even of $N$. Finally, one puts
$H_{[\infty]}=\bigcap_k H_{[k]}$.

\para{6.4.15}\btanchor{6.4.15}
It is not impossible that the condition
\begin{equation*}\tag{Pr}
  H_{[k]}\subset H_{(k)}\qquad\text{for every }k\in\Rt
\end{equation*}
is always satisfied. We shall see later that it is so in very general
cases, for example as soon as all the irreducible components of $\Phi$ are
of rank $>1$, and also in the case of semi-simple algebraic groups over a
local field.

Note that the $H_{[k]}$ are contained in the group $G'$ generated by the
$U_a$ for $a\in\Phi$. Consequently, if $H_{[\infty]}\subset H_{(\infty)}$,
then $H_{[\infty]}$ is contained in the center of $G'$.

\begin{exams}[6.4.16]\btanchor{6.4.16}
a) Let us take up again the notation of (\btref{6.2.3})~a) ($G=\mathrm{SL}_2(K)$),
and let $k>0$. One sees easily that $H_{(k)}$ consists of the matrices
$\begin{pmatrix}x&0\\0&y\end{pmatrix}\in\mathrm{SL}_2(K)$ with
$\omega(xuy^{-1}-u)\geq\omega(u)+k$ for every $u\in K$, and that $H_{[k]}$
is generated by the matrices of the form
$\begin{pmatrix}x&0\\0&ux^{-1}u^{-1}\end{pmatrix}$ with $x,u\in K^{*}$ and
$\omega(x-1)\geq k$. One deduces at once that condition (Pr) is satisfied.
If $K$ is commutative, one has
\begin{align*}
  H_{(k)}&=\left\{\begin{pmatrix}x&0\\0&x^{-1}\end{pmatrix}
    \;\middle|\; x\in K^{*}\ \text{and}\ \omega(x^2-1)\geq k\right\}\\
  H_{[k]}&=\left\{\begin{pmatrix}x&0\\0&x^{-1}\end{pmatrix}
    \;\middle|\; x\in K^{*}\ \text{and}\ \omega(x-1)\geq k\right\}.
\end{align*}

% =====================================================================
\pmark{140}
b) Let us now take up again the notation of (\btref{6.2.3})~b) ($G$ is thus the
group of rational points of a connected semi-simple algebraic group
$\Gscr$ defined and split over the commutative valued field $K$). Let
$k>0$. One verifies immediately that
\[
  H_{(k)}=\{t\in T\mid\omega(a(t)-1)\geq k\ \text{ for every }a\in\Phi\}.
\]

On the other hand, $H_{[k]}$ is generated by the images under the various
homomorphisms $\zeta_a$ (for $a\in\Phi^{+}$) of the analogous subgroups of
$\mathrm{SL}_2(K)$. Now, if $a,b\in\Phi$ and $x\in K^{*}$, one has
\[
  b\left(\zeta_a\left(\begin{pmatrix}x&0\\0&x^{-1}\end{pmatrix}\right)\right)    =x^{n_{a,b}},\qquad\text{with}\quad n_{a,b}=2\frac{(b\,|\,a)}{(a\,|\,a)},
\]
which shows that $H_{[k]}\subset H_{(k)}$.

Suppose moreover that $\Gscr$ is simply connected and let $P$ be the group
of weights of $\Gscr$. It is known that the map which to $t\in T$ associates
the map $p\mapsto p(t)$ from $P$ into $K^{*}$ is an isomorphism of $T$ onto
$\Hom(P,K^{*})$. For $a\in\Phi$, $p\in P$ and $x\in K^{*}$ one has
\[
  p\left(\zeta_a\left(\begin{pmatrix}x&0\\0&x^{-1}\end{pmatrix}\right)\right)    =x^{n}\quad\text{with }n\in\Z,
\]
from which one deduces that $\omega(p(t)-1)\geq k$ for every $t\in H_{[k]}$.
Conversely, let $(p_a)_{a\in\Pi}$ be the family of fundamental weights
associated with the basis $\Pi$ of $\Phi$. One has
\[
  t=\prod_{a\in\Pi}\zeta_a\left(\begin{pmatrix}p_a(t)&0\\0&p_a(t)^{-1}
   \end{pmatrix}\right)\qquad(t\in T)
\]
from which one deduces that
\[
  H_{[k]}=\{t\in T\mid\omega(p(t)-1)\geq k\ \text{ for every weight }p\in P\}.
\]

If $\Gscr$ is not simply connected, $H_{[k]}$ is the canonical image of the
corresponding subgroup of the simply connected covering of $\Gscr$.
\end{exams}

\begin{prop}[6.4.17]\btanchor{6.4.17}
Let $f:\Phi\to\Rt$ be a quasi-concave function. Put
\[
  d(f)=\inf_{a\in\Phi}\big(\inf(f(a)+f(-a),\,f(-2a)+2f(a))\big).
\]
The group $H_f=H\cap U_f$ is contained in $H_{[d(f)]}$. Moreover, $H_f$ is
generated by the $H_f^{(a)}=H\cap N_f^{(a)}=H\cap U_f^{(a)}$ for $a\in\Phi$.
\end{prop}

Let us prove the last assertion. Let $H'$ be the subgroup generated by the
$H_f^{(a)}$. It is obviously a normal subgroup of $N_f$. Let $\Delta$ be a
basis of $\Phi_f$, and let $N_1$ be the subgroup of $N_f$ generated by the
$N_f^{(a)}\cap M_a^{0}$ for $a\in\Delta$. In view of (\btref{6.4.10}) and
(\btref{6.4.9})~(iv) one has $N_f=N_1H'$, and (\btref{6.1.10}) implies $H'=H_f$.

To complete the proof of the proposition, it therefore suffices for us to
show that $H_f^{(a)}\subset H_{[d(f)]}$ for every $a\in\Phi$. Let
$a\in\Phi$; put, with the notation of (\btref{6.4.12})~a),
\[
  r=\inf\big(f''(a),\tfrac12 f''(2a)\big)\qquad\text{and}\qquad   s=\inf\big(f''(-a),\tfrac12 f''(-2a)\big).
\]
In view of (\btref{6.4.12})~(1), one has $r+s\geq0$ and $U_f^{(a)}$ is contained in
the group generated by $U_{a,r}$ and $U_{-a,s}$. One deduces that
$H_f\subset H_{[0]}$, whence the desired result when $d=d(f)\leq0$. Suppose
now $d>0$; for $\varepsilon\in\{\pm1,\pm2\}$ put
\[
  r_{\varepsilon}=\inf\Big\{\sum_i f(\varepsilon_i a)\;\Big|\;
   \varepsilon_i\in\{\pm1,\pm2\},\ \sum_i\varepsilon_i=\varepsilon\Big\}.
\]
% =====================================================================
\pmark{141}
It is easy to see that
\[
  r_{-2}\leq2r_{-1},\quad r_2\leq2r_1,\quad r_1\leq r_2+r_{-1},\quad   r_{-1}\leq r_{-2}+r_1\quad\text{and}\quad r_2+r_{-2}\geq d>0.
\]
One may therefore apply (\btref{6.3.9}), which shows that, if $X$ denotes the group
generated by the $H$-components of the commutators $(u,u')$ for
$u\in U_{a,r_1}\cup U_{2a,r_2}$ and
$u'\in U_{-a,r_{-1}}\cup U_{-2a,r_{-2}}$, the product
\[
  Y=U_{-2a,r_{-2}}\cdot U_{-a,r_{-1}}\cdot X\cdot U_{a,r_1}\cdot U_{2a,r_2}
\]
is a group which contains $U_f^{(a)}$, since
$r_{\varepsilon}\leq f(\varepsilon a)$. Now one sees easily that
$r_1+r_{-1}\geq d$, $r_2+r_{-2}\geq d$, $2r_1+r_{-2}\geq d$ and
$2r_{-1}+r_2\geq d$. Consequently $H_f^{(a)}\subset X\subset H_{[d]}$.

\para{6.4.18}\btanchor{6.4.18}
We shall now study the conditions under which a group of the form $U_f$ (or
more generally $X\cdot U_f$ with $X\subset H$) normalizes another group of
the same form. We shall need the following notation: if $f$ and $g$ are two
functions on $\Phi$ with values in $\Rt$, one denotes by $J(f,g)$ the set of
the triples $(a,\varepsilon,\eta)$ with $a\in\Phi$,
$\varepsilon,\eta\in\{1,2\}$ and
$\varepsilon g(\eta a)+\eta f(-\varepsilon a)>0$, and one denotes by
$H_{f,g}$ the subgroup of $H$ generated by the $H$-components of the
commutators $(u',u)$, with $u'\in U_{-\varepsilon a,f(-\varepsilon a)}$,
$u\in U_{\eta a,g(\eta a)}$ and $(a,\varepsilon,\eta)\in J(f,g)$. It is
immediate that this group is contained in $H_{[k]}$ with
\[
  k=d(f,g)=\inf\{\varepsilon g(\eta a)+\eta f(-\varepsilon a)\mid    (a,\varepsilon,\eta)\in J(f,g)\}.
\]

\begin{prop}[6.4.19]\btanchor{6.4.19}
Let $f,g:\Phi\to\Rt$ be two quasi-concave functions such that
\begin{equation}
  f(pa+qb)\leq pf(a)+qg(b)\qquad\text{for }p,q\in\N^{*}
  \text{ and }a,b,pa+qb\in\Phi.
\end{equation}
Let $X$ be a subgroup of $H$. In order that $U_g$ normalize the group
$X\cdot U_f$ it is necessary and sufficient that the two following
conditions be fulfilled:
\begin{align}
  &(X,U_{g,a})\subset U_{f,a}\qquad\text{for every }a\in\Phi;\\
  &H_{f,g}\subset X\cdot U_f .
\end{align}
\end{prop}

It is clear that these conditions are necessary. Assume them satisfied. It
suffices to show that the commutator of an element of a generating system
of $U_g$ with an element of a generating system of $XU_f$ is contained in
$XU_f$. Taking (2) into account, it therefore suffices to show that
$(U_{a,f(a)},U_{b,g(b)})\subset XU_f$ for $a,b\in\Phi$. If
$b\notin-\R_{+}a$, this follows from axiom (V\,3) of valuations (\btref{6.2.1}) and
from condition (1). Otherwise, there exists $c\in\Phi$ such that
$a=-\varepsilon c$ and $b=\eta c$, with $\varepsilon,\eta\in\{1,2\}$; if
$(c,\varepsilon,\eta)\in J(f,g)$, the inclusion to be established follows
from (1) and (3); if $(c,\varepsilon,\eta)\notin J(f,g)$, in other words if
$\eta f(-\varepsilon c)+\varepsilon g(\eta c)\leq0$, one has, in view of
condition (1),
\[
  f(b)=f\big(\eta(-\varepsilon c)+(\varepsilon+1)(\eta c)\big)
   \leq\eta f(-\varepsilon c)+(\varepsilon+1)g(\eta c)\leq g(\eta c)=g(b)
\]
whence $U_{b,g(b)}\subset U_{b,f(b)}$, which completes the proof.

\begin{cor}[6.4.20]\btanchor{6.4.20}
Let $f,g:\Phi\to\Rt$ be as in (\btref{6.4.19}). Suppose moreover that for every
pair $(a,b)$ of roots such that $a\in-\R_{+}b$ one has $g(a)=\infty$ or
$f(b)=\infty$. Then $U_g$ normalizes $U_f$.
\end{cor}

Indeed, one then has $H_{f,g}=\{1\}$.

% =====================================================================
\pmark{142}
\begin{cor}[6.4.21]\btanchor{6.4.21}
Let $f$ and $g$ be as in (\btref{6.4.19}). Suppose moreover that $f$ and $g$ are
infinite on $\Phi^{-}$. Then $U_g$ normalizes $U_f$.
\end{cor}

This follows from (\btref{6.4.20}).

\begin{cor}[6.4.22]\btanchor{6.4.22}
Let $f,g:\Phi\to\Rt$ be as in (\btref{6.4.19}). Put
\begin{align*}
  r&=\sup\big(\{f(a)-g(a)\mid a\in\Phi\}\cup\{f(2a)-2g(a)\mid a,2a\in\Phi\}\big)\\
  s&=\inf\big(\{f(a)+g(-a)\mid a\in\Phi\}\cup
     \{f(2a)+2g(-a),\,2f(a)+g(-2a)\mid a,2a\in\Phi\}\big).
\end{align*}
If $X$ is a subgroup of $H$ such that $H_{[s]}\subset X\subset H_{(r)}$,
then $U_g$ normalizes $X\cdot U_f$.
\end{cor}

Indeed, one has $H_{f,g}\subset H_{[s]}$ and
\[
  (H_{(r)},U_{a,g(a)})\subset U_{a,g(a)+r}\cdot U_{2a,2g(a)+r}
   \subset U_{a,f(a)}\cdot U_{2a,f(2a)}.
\]

\begin{prop}[6.4.23]\btanchor{6.4.23}
Let $f:\Phi\to\Rt$ be a concave function. For $a\in\Phi$ put
\begin{align*}
  f^{*}(a)&=f(a)   &&\text{if } f(a)+f(-a)>0\\
  f^{*}(a)&=f(a)+  &&\text{if } f(a)+f(-a)=0 .
\end{align*}
Then the function $f^{*}:\Phi\to\Rt$ is concave. Let moreover $X$ be a
subgroup of $H$ and $X^{*}$ a normal subgroup of $X$, satisfying the two
conditions:
\begin{align}
  &H_{f,f^{*}}\subset X^{*};\\
  &(X^{*},U_{f,a})\subset U_{f^{*},a}\qquad\text{for every }a\in\Phi.
\end{align}
Then $X^{*}\cdot U_{f^{*}}$ is a normal subgroup of $X\cdot U_f$.

Finally, let $\overline{G}$ be the quotient group $XU_f/X^{*}U_{f^{*}}$ and
let $\overline{U}_a$ (for $a\in\Phi\cup2\Phi$) be the canonical image of
$U_{a,f(a)}$ in $\overline{G}$, and $\overline{T}$ that of $X\cdot H_f$.
The group $\overline{U}_a$ is not contained in $\overline{U}_{2a}$ if and
only if $a\in\Phi_f$, and
$(\overline{T},(\overline{U}_a)_{a\in\Phi_f})$ is a generating root datum of
type $\Phi_f$ in $\overline{G}$.
\end{prop}

Let us first show that, for $p,q\in\N^{*}$, $a,b,pa+qb\in\Phi$, one has
\begin{equation}
  f^{*}(pa+qb)\leq pf(a)+qf^{*}(b).
\end{equation}
Indeed, if (3) is not satisfied, one necessarily has the following
relations, since $f(pa+qb)\leq pf(a)+qf(b)$:
\begin{align*}
  f^{*}(pa+qb)&\neq f(pa+qb) &&\text{hence } f(pa+qb)=-f(-pa-qb)\in\R;\\
  f^{*}(b)&=f(b) &&\text{hence } f(b)+f(-b)>0;\\
  f(pa+qb)&=pf(a)+qf(b) &&\text{hence } f(b)\in\R.
\end{align*}
One deduces that
\[
  pf(a)+(q-1)f(b)+f(-pa-qb)=-f(b)<f(-b)
\]
which contradicts the concavity of $f$ (condition (C) of (\btref{6.4.3})). Whence
(3).

Since $f^{*}(a)+f^{*}(-a)\geq f(a)+f(-a)\geq0$, one deduces at once from
(3) and from (\btref{6.4.5}) that $f^{*}$ is concave; and conditions (1), (2) and
(3) imply, in view of (\btref{6.4.19}), that $U_f$ normalizes $X^{*}U_{f^{*}}$,
whence the second assertion, since $X$ normalizes $U_{f^{*}}$ and $X^{*}$.

% =====================================================================
\pmark{143}
Suppose that $\overline{U}_a\not\subset\overline{U}_{2a}$, in other words
that
\begin{equation}
  U_{a,f(a)}\not\subset U_{a,f^{*}(a)}\cdot U_{2a,f(2a)}.
\end{equation}
This implies $f(a)\neq f^{*}(a)$, hence $f(a)+f(-a)=0$, whence
$f(2a)=2f(a)$ and
\[
  U_{a,f(a)}\not\subset U_{a,f(a)+}\cdot U_{2a}.
\]
Consequently one has $f(a)=f'(a)\in\Gamma'_a$ and
$f(-a)=-f(a)\in\Gamma'_{-a}$ (\btref{6.2.16}), whence $f(-a)=f'(-a)$ and
$a\in\Phi_f$.

Conversely, if $a\in\Phi_f$, one has $f'(a)+f'(-a)=0$; since
$f'(a)\geq f(a)$ and $f(a)+f(-a)\geq0$, one has
$f(a)=-f(-a)=f'(a)\in\R$, whence $f^{*}(a)>f(a)$ and $f'(a)\in\Gamma'_a$.
One deduces (4), whence
$\overline{U}_a\not\subset\overline{U}_{2a}$.

Finally, the verification of the last assertion is immediate.

\para{6.4.24}\btanchor{6.4.24}
It is clear that $H_{f,f^{*}}\subset H_{[0+]}$ and one has
$(H_{(0+)},U_{f,a})\subset U_{f^{*},a}$ for every $a\in\Phi$. Consequently
conditions (1) and (2) of (\btref{6.4.23}) are satisfied if one has
\[
  H_{[0+]}\subset X^{*}\subset H_{(0+)}.
\]

Of course, this is possible only if $H_{[0+]}\subset H_{(0+)}$, for example
if $\varphi$ satisfies condition (Pr) of (\btref{6.4.15}). But in every case there
exist subgroups $X^{*}$ satisfying the conditions imposed in (\btref{6.4.22}), for
example $H_{f,f^{*}}$ itself. To show this we shall need the following
lemma:

\begin{lem}[6.4.25]\btanchor{6.4.25}
Let $a\in\Phi$, $r\in\Gamma_a$ and $k\in\Rt$, with $k>0$. Let
$u\in U_{-a,-r}$, $v\in U_{a,r+k}\cup U_{2a,2r+k}$ and let $h$ be the
$H$-component of the commutator $(v,u)$.
\begin{enumerate}[label=(\roman*),leftmargin=2.2em]
\item If $b\in\Phi$, $b\notin\R a$, one has
\[
  (h,U_{b,s})\subset U_{b,s+k}\cdot U_{2b,2s+k}\qquad\text{for every }s\in\Rt.
\]
\item $(h,U_{a,r})\subset U_{a,r+k}\cdot U_{2a,2r+k}$.
\item $(h,U_{-a,-r})\subset U_{-a,-r+k}\cdot U_{-2a,-2r+k}$.
\end{enumerate}
\end{lem}

Up to replacing $\varphi$ by an equipollent valuation, we may assume, in
order to simplify the notation, that $r=0$.

Let $b\in\Phi$, $b\notin\R a$, and $s\in\Rt$; let $X$ be the group
generated by $U_{-a,0}\cup U_{a,k}\cup U_{2a,k}$ and $Y$ the group
generated by the $U_{pa+qb,\,t(p,q)}$ with $p\in\Z$, $q\in\N^{*}$,
$pa+qb\in\Phi$ and $t(p,q)=qs$ if $p<0$ and $t(p,q)=k+qs$ if $p\geq0$.
Axiom (V\,3) of valuations shows that $X$ normalizes $Y$ and that the
commutator of an element of $U_{-a,0}\cup U_{a,k}\cup U_{2a,k}$ with an
element of $U_{b,s}$ belongs to $Y$. One therefore has
$(X,U_{b,s})\subset Y$, so that, taking (\btref{6.1.6}) into account,
\[
  (H\cap X,U_{b,s})\subset Y\cap U_b=U_{b,k+s}\cdot U_{2b,k+2s}
\]
which proves (i).

Let us prove (iii). Let $x\in U_{-a,0}$ and let $Z$ be the group
$U_{-2a,k}\cdot U_{-a,k}\cdot H\cdot U_{a,k}\cdot U_{2a,k}$ (\btref{6.3.9}). In
view of (\btref{6.3.7}) there exist $y_2\in U_{-2a,k}$, $y_1\in U_{-a,k}$,
$z_1\in U_{a,k}$ and $z_2\in U_{2a,k}$ such that
% =====================================================================
\pmark{144}
$(v,u)=y_2y_1hz_1z_2$. On the other hand, $v$ and the commutators
$(v^{-1},x)$, $(v^{-1},xu)$, $(x,y_i)$ and $(x,z_i)$ all belong to $Z$.
Consequently one has
\[
  xvu=v(v^{-1},x)xu\in Zxu\qquad\text{and}\qquad xuv=v(v^{-1},xu)xu\in Zxu
\]
whence $x(v,u)x^{-1}\in Z$. But $xy_ix^{-1}=(x,y_i)y_i$ and
$xz_ix^{-1}=(x,z_i)z_i$ also belong to $Z$. Consequently one has
$xhx^{-1}\in Z$ and
$(h,x)\in Z\cap U_{-a}=U_{-a,k}\cdot U_{-2a,k}$, whence (iii).

Finally, (ii) follows from (iii) and from the following lemma:

\begin{lem}[6.4.26]\btanchor{6.4.26}
Let $a\in\Phi$, $h\in H$, $r\in\Gamma_a$ and $k\in\Rt$ with $k>0$. If one
has $(h,U_{-a,-r})\subset U_{-a,-r+k}\cdot U_{-2a,-2r+k}$, then one also
has $(h,U_{a,r})\subset U_{a,r+k}\cdot U_{2a,2r+k}$.
\end{lem}

One may suppose $r=0$. Since $0\in\Gamma_a$, the group $U_{a,0}$ is
generated by $\varphi_a^{-1}(0)$ and, since $h$ normalizes
$U_{a,k}\cdot U_{2a,k}$, it suffices to show that
$(h,x)\in U_{a,k}\cdot U_{2a,k}$ for $x\in\varphi_a^{-1}(0)$. Put
$x=x'mx''$ with $x',x''\in U_{-a,0}$ and $m\in M_{a,0}$. Then
\begin{align*}
  (h,x)&=h\cdot x'\cdot mh^{-1}m^{-1}\cdot m(h,x'')m^{-1}\cdot x'^{-1}\\
   &\in H\cdot U_{-a,0}\cdot H\cdot U_{a,k}\cdot U_{2a,k}\cdot U_{-a,0}
   \subset U_{-a,0}\cdot H\cdot U_{a,k}\cdot U_{2a,k}
\end{align*}
whence the lemma.

\begin{prop}[6.4.27]\btanchor{6.4.27}
Let $f,f^{*}$ be as in (\btref{6.4.23}). The subgroup $H_{f,f^{*}}$ is normal in
$H$ and one has $(H_{f,f^{*}},U_{f,b})\subset U_{f^{*},b}$ for every
$b\in\Phi$.
\end{prop}

The first assertion is obvious, since $H$ normalizes all the $U_{a,k}$.
Since $H_{f,f^{*}}$ and $U_{f,b}$ normalize $U_{f^{*},b}$, it suffices to
prove that, if $h$ is the $H$-component of the commutator of an element
$u\in U_{-\varepsilon a,f(-\varepsilon a)}$ and an element
$v\in U_{\eta a,f^{*}(\eta a)}$ with $a\in\Phired$,
$\varepsilon,\eta\in\{1,2\}$ and
$\varepsilon f^{*}(\eta a)+\eta f(-\varepsilon a)>0$, and if
$x\in U_{b,f(b)}$ with $b\in\Phi$, then $(h,x)\in U_{b,f^{*}(b)}$. This is
obvious if $f^{*}(b)=f(b)$ or if $f(b)\notin\Gamma_b$. If
$b\notin\R a$, it follows from (\btref{6.4.25})~(i). Suppose that
$f(a)\in\Gamma_a$, $f(a)+f(-a)=0$ and $b\in\{\pm a,\pm2a\}$; then
$u\in U_{-a,f(-a)}$ and $v\in U_{a,f(a)+}\cup U_{2a,2f(a)+}$ (note that, if
$2a\in\Phi$, one then has $f(2a)=2f(a)=-f(-2a)$), and our assertion follows
from (\btref{6.4.25})~(ii) and (iii) (taking $r=f(a)$ and $k=0+$).

It remains for us to examine only the case where $f(a)+f(-a)>0$,
$f(2a)+f(-2a)=0$, $f(2a)\in\Gamma_{2a}$ and $b=\pm2a$. Put $f(2a)=2r$. One
has $f(a)>r$, otherwise the concavity of $f$ would imply
$f(-a)\leq f(-2a)+f(a)\leq\tfrac12 f(-2a)\leq-f(a)$. Likewise,
$f(-a)>-r$, and one has $u\in U_{-a,-r}$ and
$v\in U_{a,r+}\cup U_{2a,2r+}$. By (\btref{6.4.25})~(ii) and (iii) one therefore
has
\[
  (h,U_{\pm2a,\pm2r})\subset(h,U_{\pm a,\pm r})\cap U_{\pm2a}\subset    U_{\pm a,\pm r+}\cap U_{\pm2a}=U_{\pm2a,\pm2r+}
\]
which completes the proof.

\begin{prop}[6.4.28]\btanchor{6.4.28}
Let $f,g,h:\Phi\to\Rt$ be three quasi-concave functions. Put
$f_1=\inf(f,h)$ and $g_1=\inf(g,h)$ and suppose that, for $p,q\in\N^{*}$
and $a,b,pa+qb\in\Phi$, one has
\begin{equation}
  h(pa+qb)\leq pf_1(a)+qg_1(b).
\end{equation}
Let moreover $X$ and $Y$ be two subgroups of $H$ such that
\begin{equation}
  (X,U_{g,a})\cup(Y,U_{f,a})\cup(H_{f_1,g_1},U_{f,a}\cup U_{g,a})
   \subset U_{h,a}
\end{equation}
for every $a\in\Phi$.
\end{prop}
\par\endgroup

% ---------------------------------------- bruhattits.tex
\begingroup\linespread{1.06}
\providecommand{\bR}{}
\renewcommand{\bR}{\mathbf{R}}
\providecommand{\bN}{}
\renewcommand{\bN}{\mathbf{N}}
\providecommand{\bZ}{}
\renewcommand{\bZ}{\mathbf{Z}}
\providecommand{\bF}{}
\renewcommand{\bF}{\mathbf{F}}
\providecommand{\Rt}{}
\renewcommand{\Rt}{\widetilde{\bR}}
\providecommand{\Rtp}{}
\renewcommand{\Rtp}{\Rt_{+}}
\providecommand{\Ns}{}
\renewcommand{\Ns}{\bN^{*}}
\providecommand{\Cx}{}
\renewcommand{\Cx}{\mathscr{C}}
\providecommand{\Bldg}{}
\renewcommand{\Bldg}{\mathscr{I}}
\providecommand{\Phired}{}
\renewcommand{\Phired}{\Phi^{\mathrm{red}}}
\providecommand{\Wv}{}
\renewcommand{\Wv}{{}^{v}W}
\providecommand{\nuv}{}
\renewcommand{\nuv}{{}^{v}\nu}
\providecommand{\Sigmav}{}
\renewcommand{\Sigmav}{{}^{v}\Sigma}
\providecommand{\cl}{}
\renewcommand{\cl}{\operatorname{cl}}
\providecommand{\Hcl}{}
\renewcommand{\Hcl}[1]{H_{[#1]}}
\providecommand{\Hop}{}
\renewcommand{\Hop}[1]{H_{(#1)}}
\providecommand{\Bfr}{}
\renewcommand{\Bfr}{\mathfrak{B}}
\providecommand{\Ker}{}\renewcommand{\Ker}{\operatorname{Ker}}
\providecommand{\comm}{}
\renewcommand{\comm}[2]{(#1,#2)}
\providecommand{\cpl}{}
\renewcommand{\cpl}{\complement}
\newenvironment{bt}[2]% {kind}{number}
  {\par\medskip\noindent\textbf{#1 (#2).}\ ---\ \begingroup\itshape}
  {\endgroup\par\medskip}
\newenvironment{btpar}[1]% unnumbered-statement paragraph, e.g. (6.4.33)
  {\par\medskip\noindent\textbf{(#1)}\ }
  {\par\medskip}
\newenvironment{sketch}
  {\par\smallskip\noindent\small\textit{Proof.}\ \ignorespaces}
  {\par\smallskip\normalsize}
\providecommand{\tnote}{}
\renewcommand{\tnote}[1]{\footnote{\emph{Translator's note.} #1}}
\bigskip

% =====================================================================

% =====================================================================

\noindent
Throughout, $\Phi$ is the root system of the root datum
$(T,(U_a))$, $\Rt$ denotes $\bR$ enlarged by the symbols $r+$, and for a
function $f:\Phi\to\Rt$ we write $U_f$ for the subgroup generated by the
$U_{a,f(a)}$, $a\in\Phi$, and $P_f=H\cdot U_f$. Commutator subgroups are
written $\comm{X}{Y}$, as in the original.

\medskip
Page~145 completes the proof of (\btref{6.4.28}), whose conclusion is that
$\comm{X}{Y}\cdot H_{f_1,g_1}\cdot U_h$ is a group containing the
commutator group $\comm{X\cdot U_f}{Y\cdot U_g}$. The argument runs in
three steps: an elementary inequality showing $h(-\varepsilon a)=f_1(-\varepsilon a)$
and $h(\eta a)=g_1(\eta a)$ under the stated sign condition; the
verification, via (\btref{6.4.19}), that $U_f$ and $U_g$ normalize
$H_{f_1,g_1}\cdot U_h$; and the inclusion
$\comm{U_f}{U_g}\subset H_{f_1,g_1}\cdot U_h$, obtained from (V\,3) and
(\btref{6.3.9}). The general case then follows from the commutator identity
\[
  (xu,yv)=xy\cdot(y^{-1},u)\cdot(u,v)\cdot(v,x^{-1})\cdot x^{-1}y^{-1}.
\]

\begin{bt}{Corollary}{6.4.29}\btanchor{6.4.29}
Let $\Delta$ be a subset of the basis $\Pi$ of $\Phi$ and let $\Psi$ be
the set of positive roots that are not linear combinations of elements
of $\Delta$. Let $f,g:\Phi\to\Rt$ be quasi-concave functions, and
suppose $f$ (resp.\ $g$) is infinite on $\cpl\Psi$ (resp.\ on $-\Psi$).
For $a\in\Phi$ put
\[
  h(a)=\inf\Bigl\{\textstyle\sum_{1\le i\le m}f(a_i)         +\sum_{1\le j\le n}g(b_j)\Bigr\},
\]
the infimum being taken over all pairs of non-empty finite sequences
$(a_i)$, $(b_j)$ of elements of $\Phi$ with
$\sum_i a_i+\sum_j b_j=a$. Then $h:\Phi\to\Rt$ is concave and
$\comm{U_f}{U_g}\subset U_h$.
\end{bt}

\begin{sketch}
$h$ is concave and satisfies (\btref{6.4.28})(1) with $f_1=\inf(f,h)$,
$g_1=\inf(g,h)$. Examining the two cases $a\in-\Psi$ and
$a\in(\pm\Phi^{+}\cap\cpl\Psi)$ shows $f_1$ or $g_1$ is infinite at $a$;
hence $H_{f_1,g_1}=\{1\}$ and (\btref{6.4.28}) applies.
\end{sketch}

\begin{bt}{Corollary}{6.4.30}\btanchor{6.4.30}
Let $f_{*},g_{*}:\Phi\cup\{0\}\to\Rt$ be functions whose restrictions to
$\Phi$ are quasi-concave, and let
$h_{*}:\Phi\cup\{0\}\to\Rt\cup\{-\infty\}$ be given by
\[
  h_{*}(a)=\inf\Bigl\{\textstyle\sum_i f_{*}(a_i)+\sum_j g_{*}(b_j)\Bigr\},
\]
the infimum over pairs of non-empty finite sequences $(a_i)$, $(b_j)$ in
$\Phi\cup\{0\}$ with $a=\sum_i a_i+\sum_j b_j$; the order on
$\Rt\cup\{-\infty\}$ is defined by $-\infty<k$ for all $k\in\Rt$.

If $h_{*}(0)\ne-\infty$, then $h_{*}(a)\ne-\infty$ for every $a\in\Phi$
and the restriction $h$ of $h_{*}$ to $\Phi$ is concave. Suppose in
addition that $\Hcl{h_{*}(0)}\subset\Hop{h_{*}(0)}$ (for instance that
$\varphi$ satisfies condition~(Pr)), and let $X$ (resp.\ $Y$) be a
subgroup of $\Hop{f_{*}(0)}$ (resp.\ $\Hop{g_{*}(0)}$). Then
\[
  \comm{X\cdot U_f}{Y\cdot U_g}\subset
  \comm{X}{Y}\cdot \Hcl{h_{*}(0)}\cdot U_h .
\]
\end{bt}

\begin{sketch}
One checks that, in the notation of (\btref{6.4.28}),
$H_{f_1,g_1}\subset \Hcl{h_{*}(0)}$ and that conditions (1) and (2) of
(\btref{6.4.28}) hold.
\end{sketch}

\begin{bt}{Corollary}{6.4.31}\btanchor{6.4.31}
Let $f$ and $g$ be linear functions on $\Phi$ and put
$k=\inf\{g(a)-f(a)\mid a\in\Phi^{+}\}$. Then
$\comm{U^{+}_{f}}{U^{+}_{g}}\subset U^{+}_{f+k}$.
\end{bt}

\begin{sketch}
Apply (\btref{6.4.29}) with $\Delta=\emptyset$ (so $\Psi=\Phi^{+}$ and
$\cpl\Psi=-\Psi=\Phi^{-}$) to the functions $f^{+},g^{+}$ equal to
$\infty$ on $\Phi^{-}$ and to $f,g$ on $\Phi^{+}$. Linearity of $f$ gives
$h(a)=\inf\{f(a)+\sum_j(g(b_j)-f(b_j))\}\ge f(a)+k$ for $a\in\Phi^{+}$.
\end{sketch}

\noindent
The authors remark that $g$ need not be assumed linear: quasi-concavity
of $g$ suffices.

\begin{bt}{Corollary}{6.4.32}\btanchor{6.4.32}
Let $L$ be a half-line with origin $0$ in $V^{*}$, contained in $D_0$.
For $x\in A$ let $U^{+}_{x}$ denote the subgroup generated by the
$U_{a,-a(x-\varphi)}$ for $a\in\Phi^{+}$. There is a constant $k$ with
$0<k\le1$ such that
\[
  \comm{U^{+}_{x}}{U^{+}_{y}}\subset U^{+}_{x+k(y-x)}
\]
for all points $x,y$ of $A$ with $x-y\in L$.
\end{bt}

\begin{sketch}
Take $v\in L$ and set $k=\inf a(v)/\sup a(v)$, the bounds over
$a\in\Phi^{+}$; then $0<k\le1$. Write $x=y+tv$ with $t\in\bR_{+}$ and
apply (\btref{6.4.31}): the resulting function satisfies
$h(b)\ge -b(x)+b(ktv)=-b(x+k(y-x))$.
\end{sketch}

\noindent
If $\Phi$ has rank one and type $BC_1$, then $k=1/2$. If $\Phi$ has type
$A_1$, then $\comm{U^{+}_{x}}{U^{+}_{y}}=\{1\}$ for all $x,y\in A$.

\begin{btpar}{6.4.33}
The subgroups $\Hcl{k}$ and $\Hop{k}$ and condition~(Pr) of (\btref{6.4.15}) are
now studied in more detail.
\end{btpar}

\begin{bt}{Proposition}{6.4.33}\btanchor{6.4.33}
For all $k,\ell\in\Rt$, the commutator group
$\comm{\Hop{k}}{\Hop{\ell}}$ is contained in $\Hop{k+\ell}$. If
$\Hcl{k}\subset\Hop{k}$, then for $a\in\Phi$ and $m\in M^{0}_{a}$ one has
$\comm{m}{\Hop{k}}\subset\Hcl{k}$.
\end{bt}

\begin{sketch}
For the first assertion one may assume $k,\ell>0$; setting
$X_y=U_{a,x+y}\cdot U_{2a,2x+y}\subset U_{a,x}$ for $y\in\Rtp$, (V\,3)
gives $\comm{X_y}{X_z}\subset X_{y+z}$, and a computation with the
commutators $v=(u^{-1},h^{-1})$, $v'=(u^{-1},h'^{-1})$ modulo
$X_{k+\ell}$ yields $((h,h'),u)\in X_{k+\ell}$. For the second, write
$m=u'uu''$ and use the group
$Y=U_{-2a,-2r+k}U_{-a,-r+k}\Hcl{k}U_{a,r+k}U_{2a,2r+k}$ of (\btref{6.3.9}),
normalized by $U_{a,r}$ and $U_{-a,-r}$ by (\btref{6.4.19}); then
$\comm{m}{h}\in Y\cap H=\Hcl{k}$.
\end{sketch}

\begin{bt}{Proposition}{6.4.34}\btanchor{6.4.34}
Let $k\in\Rt$ with $k>0$. An element $h$ of $H$ belongs to $\Hop{k}$ if
and only if
\begin{equation}
  \comm{h}{U_{a,r}}\subset U_{a,r+k}\cdot U_{2a,2r+k}
  \qquad\text{for all } r\in\bR,
  \tag{1}
\end{equation}
for every $a\in\Pi$.
\end{bt}

\begin{sketch}
If (1) holds for a root $a$ it holds for $2a$; so it suffices to show
that the condition ``(1) holds for the elements of $\Pi$'' is
independent of the choice of $\Pi$. This reduces to two assertions:
(2) if (1) holds for $a$ it holds for $-a$ --- this is Lemma (\btref{6.4.26});
and (3) if (1) holds for distinct $a,b\in\Pi$ it holds for $c=r_a(b)$ ---
proved using (\btref{6.1.8}), (\btref{6.4.21}) and (\btref{6.4.9}).
\end{sketch}

\begin{bt}{Proposition}{6.4.35}\btanchor{6.4.35}
If $\Phi$ has no irreducible component of rank one, then condition~(Pr)
is satisfied.
\end{bt}

\begin{sketch}
Lemma (\btref{6.4.25}) gives condition (\btref{6.4.34})(1) for every root
$a\notin\bR b$, and the hypothesis on $\Phi$ allows one to choose the
basis $\Pi$ with $b\notin\bR\Pi$; then apply (\btref{6.4.34}).
\end{sketch}

\begin{bt}{Remark}{6.4.36}\btanchor{6.4.36}
A more elaborate argument shows that (Pr) holds as soon as $\Phi$
contains no irreducible component of type $BC_1$. The authors know of no
example in which (Pr) fails.
\end{bt}

\begin{bt}{Proposition}{6.4.37}\btanchor{6.4.37}
Suppose (Pr) is satisfied and let $k,\ell\in\Rtp$. Then
\begin{align}
  \comm{\Hop{k}}{\Hcl{\ell}} &\subset \Hcl{k+\ell},\tag{1}\\
  \comm{\Hcl{k}}{\Hcl{\ell}} &\subset \Hcl{k+\ell}.\tag{2}
\end{align}
\end{bt}

\begin{sketch}
Only (1) is needed. Assume $k>0$; since $\Hcl{k+\ell}$ is normal in $H$
it suffices to treat commutators with a generating system of
$\Hcl{\ell}$, i.e.\ with $U_f\cap H$ for a suitable $f$ supported on
$\pm a,\pm 2a$. Apply (\btref{6.4.30}) with $X=\{1\}$, $Y=\Hop{k}$, and conclude
with (\btref{6.3.9}), which gives
$U_h\cap H\subset \Hcl{2k+\ell}\subset \Hcl{k+\ell}$.
\end{sketch}

\begin{bt}{Definition}{6.4.38}\btanchor{6.4.38}
An \emph{extension} of the valuation $\varphi$ is a map
$\varphi_0:H\to\bR_{+}\cup\{\infty\}$ satisfying:
\begin{itemize}
\item[(P\,1)] for every $k\in\bR$ the preimage
  $U_{0,k}=\varphi_0^{-1}([k,\infty])$ is a subgroup of $H$, and
  $\Hcl{k}\subset U_{0,k}\subset \Hop{k}$;
\item[(P\,2)] for $k,\ell\in\bR$, $\comm{U_{0,k}}{U_{0,\ell}}\subset U_{0,k+\ell}$.
\end{itemize}
$\varphi$ is \emph{extendable} if it admits an extension. An
\emph{extended valuation} of the root datum $(T,(U_a))$ is a pair
$(\psi,\psi_0)$ consisting of a valuation $\psi$ of $(T,(U_a))$ and an
extension $\psi_0$ of $\psi$.
\end{bt}

\noindent
(P\,2) implies each $U_{0,k}$ is normal in $U_{0,0}=H$; together with
(\btref{6.4.33}), (P\,1) shows $U_{0,k}$ is even normal in $N$. Both conditions
remain valid if $k$ or $\ell$ is allowed the value $\infty$.
Condition (P\,1) may be made explicit: for all $a\in\Phi$ and
$h,k\in\bR$, $U_{0,k}$ is a group and
\begin{itemize}
\item[(P\,1$'$)] $\comm{U_{a,h}}{U_{0,k}}\subset U_{a,h+k}\cdot U_{2a,2h+k}$;
\item[(P\,1$''$)] the $H$-components of the commutators $(u,u')$ and
  $(u',u)$, for $u\in U_{a,r}$ and $u'\in U_{-a,s}$ (resp.\
  $u'\in U_{-2a,s}$) with $r+s=k$ (resp.\ $2r+s=k$), lie in $U_{0,k}$.
\end{itemize}

\begin{bt}{Proposition}{6.4.39}\btanchor{6.4.39}
The valuation $\varphi$ is extendable if and only if condition~(Pr) of
(\btref{6.4.15}) is satisfied.
\end{bt}

\begin{sketch}
Necessity is clear. Conversely, under (Pr) the two functions
\[
  \varphi_{[0]}(h)=\sup\bigl(\{k\in\bR\mid h\in \Hcl{k}\}\cup\{0\}\bigr),
  \qquad
  \varphi_{(0)}(h)=\sup\{k\in\bR\mid h\in \Hop{k}\}
\]
are extensions of $\varphi$, by (\btref{6.4.37}).
\end{sketch}

\begin{btpar}{6.4.40}
From here on the authors set $U_0=H$. This makes the notation $U_{a,k}$
coherent for $a\in\Phi\cup\{0\}$ and $k\in\bR$, and is justified by the
fact that $H$ plays, relative to the element $0$ of $V^{*}$, a role
analogous to that of the groups $U_a$ relative to the roots $a\in\Phi$.
The next proposition exhibits this parallelism between $0$ and the roots
with respect to an extended valuation.
\end{btpar}

\begin{bt}{Proposition}{6.4.41}\btanchor{6.4.41}
For $a\in\Phi\cup\{0\}$ let $\psi_a$ be a map from $U_a$ to
$\bR\cup\{\infty\}$, with $\psi_0(u)\ge0$ for all $u\in U_0$. The pair
$((\psi_a)_{a\in\Phi},\psi_0)$ is an extended valuation of the root group
datum $(T,(U_a))$ if and only if the maps $\psi_a$ satisfy conditions
(V\,0), (V\,1), (V\,2), (V\,4), (V\,5) of (\btref{6.2.1}) (with $\varphi$
replaced by $\psi$) together with:
\begin{itemize}
\item[(V\,1 bis)] for every $k\in\bR$, the set
  $U_{0,k}=\psi_0^{-1}([k,\infty])$ is a group;
\item[(VP)] let $a,b\in\Phi\cup\{0\}$ and $k,\ell\in\bR$; let $\Psi$ be
  the set of elements of $\Phi\cup\{0\}$ of the form $pa+qb$ with
  $p,q\in\Ns$, and for $c\in\Psi$ put
  $k(c)=\inf\{pk+q\ell\mid p,q\in\Ns,\ pa+qb=c\}$. Then the commutator
  group $\comm{U_{a,k}}{U_{b,\ell}}$ is contained in the product of the
  groups $U_{c,k(c)}$ for $c\in\Psi$ and $c/2\in\Psi$, taken in a
  suitable order.
\end{itemize}
\end{bt}

\begin{sketch}
Necessity is straightforward. Conversely (VP) for $b\notin-\bR_{+}a$ is
just (V\,3), so $\psi=(\psi_a)$ is a valuation; (VP) with $a=b=0$ gives
(P\,2), with $a=0,\,b\in\Phi$ gives $U_{0,k}\subset \Hop{k}$, and with
$a\in\Phi$, $b=-a$ or $b=-2a$ gives $\Hcl{r}\subset U_{0,r}$ via (\btref{6.3.9}).
\end{sketch}

\begin{bt}{Definition}{6.4.42}\btanchor{6.4.42}
Let $\varphi_0$ be an extension of $\varphi$. For a function
$f:\Phi\cup\{0\}\to\Rt$, let $U_f$ again denote the subgroup of $G$
generated by the $U_{a,f(a)}$ for $a\in\Phi\cup\{0\}$ (for $k=r+$ with
$r\in\bR$, $U_{0,k}$ denotes the union of the $U_{0,s}$ for
$s\in\bR\cup\{\infty\}$, $s>r$). One says $f$ is \emph{quasi-concave}
(relative to $(\varphi,\varphi_0)$) if its restriction to $\Phi$ is
quasi-concave and $H_{f|\Phi}$ is contained in $U_{0,f(0)}$.
\end{bt}

\noindent
If $f:\Phi\cup\{0\}\to\Rt$ is concave (\btref{6.4.4}) then it is quasi-concave,
by (\btref{6.4.17}), since $f(0)\le d(f|\Phi)$. If $f$ is quasi-concave then
$U_f\cap U_0=U_{0,f(0)}$.

\begin{bt}{Proposition}{6.4.43}\btanchor{6.4.43}
Let $\varphi_0$ be an extension of $\varphi$ and let
$f,g:\Phi\cup\{0\}\to\Rt$ be quasi-concave functions such that
$f(pa+qb)\le pf(a)+qg(b)$ for $p,q\in\Ns$ and
$a,b,pa+qb\in\Phi\cup\{0\}$. Then $U_f$ is normalized by $U_g$.
\end{bt}

\begin{sketch}
One has $H_{f|\Phi,\,g|\Phi}\subset U_{0,f(0)}$ and
\[
  \comm{U_{0,f(0)}}{U_{a,g(a)}}\subset
  \comm{\Hop{f(0)}}{U_{a,g(a)}}\subset
  U_{a,g(a)+f(0)}U_{2a,2g(a)+f(0)}\subset U_{a,f(a)}U_{2a,f(2a)};
\]
now apply (\btref{6.4.19}) with $X=U_{0,f(0)}$.
\end{sketch}

\begin{bt}{Proposition}{6.4.44}\btanchor{6.4.44}
Let $\varphi_0$ be an extension of $\varphi$, let
$f,g:\Phi\cup\{0\}\to\Rt$ be quasi-concave, and let
$h:\Phi\cup\{0\}\to\Rt\cup\{-\infty\}$ be defined by
\[
  h(a)=\inf\Bigl\{\textstyle\sum_i f(a_i)+\sum_j g(b_j)\Bigr\},
\]
the infimum over pairs of non-empty finite sequences $(a_i)$, $(b_j)$ in
$\Phi\cup\{0\}$ with $a=\sum_i a_i+\sum_j b_j$. If $h(0)\ne-\infty$,
then $h$ is concave and $\comm{U_f}{U_g}\subset U_h$.
\end{bt}

\begin{sketch}
A special case of (\btref{6.4.30}), taking $X=U_{0,f(0)}$ and $Y=U_{0,g(0)}$.
\end{sketch}

\begin{bt}{Proposition}{6.4.45}\btanchor{6.4.45}
Suppose $\varphi$ is extendable (i.e.\ (Pr) holds). Let $f:\Phi\to\Rt$
be quasi-concave, and suppose there exist a linear function
$\lambda\in V$ and $k\in\Rt$ with $k>0+$ and $f\ge\lambda+k$
(cf.\ (\btref{6.4.6})). Let $H'$ be a subgroup of $\Hop{k}$. Put $Z=H'\cdot U_f$
and, for $n\in\Ns$,
\[
  Z_n=Z\cap\bigl(\Cx^{n}(H')\cdot \Hcl{nk}\cdot U_{\lambda+nk}\bigr).
\]
Then:
\begin{enumerate}
\item[(1)] $\comm{Z_n}{Z_{n'}}\subset Z_{n+n'}$ for all $n,n'\in\Ns$;
\item[(2)] $\Cx^{n}(Z)\subset Z_n$ for all $n\in\Ns$;
\item[(3)] $\Cx^{\infty}(Z)\subset \Hop{\infty}$ --- recall that this is
  the centralizer of the group generated by the $U^{a}$, $a\in\Phi$
  (\btref{6.4.13});
\item[(4)] if $H'\subset \Hcl{k}$, then $\Cx^{\infty}(Z)$ is contained in
  $\Hcl{\infty}$ and in the center of $Z$. If moreover
  $\Hcl{\infty}=\{1\}$, then $Z$ is pronilpotent.
\end{enumerate}
\end{bt}

\noindent
Here, for a group $X$, $\Cx^{n}(X)$ denotes the $n$-th term of the
descending central series, so $\Cx^{1}(X)=X$ and
$\Cx^{n+1}(X)=(\Cx^{n}(X),X)$, and
$\Cx^{\infty}(X)=\bigcap_{n\ge1}\Cx^{n}(X)$.

\begin{sketch}
Apply (\btref{6.4.30}) with $f_{*},g_{*}$ the restrictions of $\lambda+nk$,
$\lambda+n'k$ and $X=\Cx^{n}(H')\Hcl{nk}$, $Y=\Cx^{n'}(H')\Hcl{n'k}$;
(1) follows from $h_{*}\ge\lambda+(n+n')k$ together with
$(\Cx^{n}(H'),\Cx^{n'}(H'))\subset\Cx^{n+n'}(H')$ and (\btref{6.4.37}). Then (2)
and (3) follow using (\btref{6.4.33}), (\btref{6.4.9}) and (DR\,6), and (4) from
(\btref{6.4.15}).
\end{sketch}

\begin{bt}{Corollary}{6.4.46}\btanchor{6.4.46}
Suppose $\varphi$ is extendable and discrete, and let $e\in\bR^{*}_{+}$
be such that $\Gamma_a\subset\bZ e$ for every $a\in\Phi$. Suppose
moreover $\Hcl{\infty}=\{1\}$, and resume the hypotheses and notation of
(\btref{6.4.23}). Then $H_{f,f^{*}}\subset \Hcl{e}$. If $X^{*}\subset \Hcl{e}$,
then $X^{*}\cdot U_{f^{*}}$ is a pronilpotent normal subgroup of
$X\cdot U_f$.
\end{bt}

\begin{sketch}
By (\btref{6.4.24}), $H_{f,f^{*}}\subset \Hcl{0+}=\Hcl{e}$. Setting
$g(a)=\inf\{t\in\bZ e\mid t\ge f^{*}(a)\}$ gives $U_g=U_{f^{*}}$ with
$g$ quasi-concave, and (\btref{6.4.6}) supplies $\lambda$ and $k\le e$ with
$g>\lambda+k$; apply (\btref{6.4.45}).
\end{sketch}

\begin{bt}{Remarks}{6.4.47}\btanchor{6.4.47}
(1) The hypothesis $\Gamma_a\subset\bZ e$ in (\btref{6.4.46}) is not very
restrictive: by (\btref{6.2.23}), every discrete valuation is equivalent
(\btref{6.2.5}) to one with this property.

(2) With the hypotheses and notation of (\btref{6.4.23}) and (\btref{6.4.46}) one may
take $X=H$ and $X^{*}=\Hcl{e}$. The group $H\cdot U_f$ is then an
extension of $\overline{G}=HU_f/\Hcl{e}U_{f^{*}}$, which carries a
generating root datum $(H/\Hcl{e},(\overline{U}_a))$ of type
$\Phi_f$ --- making it akin to a reductive algebraic group --- by the
pronilpotent group $\Hcl{e}\cdot U_{f^{*}}$. This situation recurs later
for semisimple algebraic groups over a local field: $\overline{G}$ is
then the group of rational points of a reductive algebraic group over
the residue field $k$, and $\Ker(H\cdot U_f\to\overline{G})$ the group of
rational points of a prounipotent proalgebraic group over $k$, the
``prounipotent radical'' of $H\cdot U_f$.
\end{bt}

\begin{bt}{Proposition}{6.4.48}\btanchor{6.4.48}
Let $g:\Phi\cup\{0\}\to\Rt$ be a concave function with $g(0)>0$, let
$f:\Phi\to\Rt$ be a quasi-concave function majorising the restriction of
$g$ to $\Phi$, and let $X$ be a subgroup of $H$ containing $H_g$. Let
$(a_1,\dots,a_m)$ be the elements of $\Phired$ in an arbitrary order and
let $j\in\bN$ with $0\le j\le m$. Let
\[
  \tau_j:\;X\times\prod_{1\le i\le m}U_{f,a_i}\longrightarrow X\cdot U_f
\]
be the map $(x,u_1,\dots,u_m)\mapsto u_1\cdots u_j\,x\,u_{j+1}\cdots u_m$.
\begin{enumerate}
\item[(i)] $\tau_j$ is injective.
\item[(ii)] If, for every $a\in\Phi$, the group $U_a$ --- given the
  topological group structure for which the family $(U_{a,k})_{k\in\bR}$
  is a fundamental system of neighborhoods of $1$ --- is complete, then
  $\tau_j$ is bijective.
\end{enumerate}
\end{bt}

\begin{sketch}
One reduces to $j=m$ with $\Phi$ irreducible. Rank one follows from
(DR\,6), (\btref{6.4.9}), (\btref{6.4.10}) (completeness is not needed there). For rank
$\ge2$, $\varphi$ is extendable by (\btref{6.4.35}), so (\btref{6.4.45}) is available:
with $Z=\Hcl{k}U_{\lambda+k}$, $Z_n=\Hcl{nk}U_{\lambda+nk}$, an
induction on $n$ shows $u_i^{-1}u'_i\in Z_n\cap U_{a_i}$, giving (i) on
letting $n\to\infty$. For (ii) a second induction produces
approximations $u_i^{(n)}$ that are Cauchy by construction; their limits
$v_i$ satisfy $(v_1\cdots v_m)^{-1}u\in X$, so $u$ lies in the image of
$\tau_m$.
\end{sketch}

% =====================================================================

\btsecB{6.5}{6.5. Discrete valuations and double Tits systems}

% =====================================================================

\noindent
In this subsection the notation of 6.2 --- in particular that of (\btref{6.2.4})
--- and of 6.4 (see (\btref{6.4.2}) and (\btref{6.4.3}) for the meaning of $P_f$ and
$f_\Omega$) is retained, and the valuation $\varphi$ is assumed
\emph{discrete}. By (\btref{6.2.22}) one may then also use the notation of 1.3
and assume $A=\mathbf{A}$, $W=\mathbf{W}$, $\Sigma=\boldsymbol{\Sigma}$.

\begin{bt}{Theorem}{6.5}
Let $G'$ be the subgroup of $G$ generated by $N'=\nu^{-1}(\mathbf{W})$
and the $U_a$ $(a\in\Phi)$. Put $T'=T\cap N'$ and $B=P_{f_{\mathbf{C}}}$.
\begin{enumerate}
\item[(i)] $B\cap N'=H$ and $N'/H=\mathbf{W}$. The quadruple
  $(G',B,N',\mathbf{S})$ is a double Tits system. If this system is
  substituted for the one written $(G,B,N,S)$ in \S4, the groups written
  $\Bfr^{0}_{D}$ and $\Bfr_{D}$ in 4.1 are respectively $HU^{+}$ and
  $T'U^{+}$ (for $D=D_0$).
\item[(ii)] The injection of $G'$ into $G$ is a $(B,N')$-adapted
  homomorphism of connected type, for which the objects written
  $\widehat{N}$, $\widehat{\nu}$ and $\widehat{W}$ in 4.1 are
  respectively $N$, $\nu$ and $\widehat{W}$.
\end{enumerate}
\end{bt}

\begin{sketch}
The theorem is an almost immediate consequence of (\btref{5.2.19}) and (\btref{5.2.31});
the authors also give an independent proof based on (\btref{6.4.9}), for readers
who have not read \S5. That (i) implies (ii) follows from
$G=T\cdot G'=N\cdot G'$ together with (\btref{6.2.10})(iii). For (i): the first
two relations come from (\btref{6.4.10}) and the definition of $N'$; axiom
(T\,1) of (\btref{1.2.6}) follows from $U_a\subset T'BT'$, axiom (T\,2) from
$(\mathbf{W},\mathbf{S})$ being a Coxeter system (\btref{1.3.4}), and (T\,4) from
$U_a\cap sBs^{-1}=U_{\alpha_{+}}$, $U_a\cap B=U_\alpha$ by (\btref{6.4.9}).
Axiom (T\,3) follows from the two inclusions $sB\subset BsU_\alpha$ and
$(U_\alpha sHU_\alpha\cup U_{(\alpha^{*})_{+}}HU_\alpha)wB\subset
BswB\cup BwB$. The identification of $\Bfr^{0}_{D}$ and $\Bfr_{D}$ uses
(\btref{4.1.5}) and $B\subset U^{-}HU^{+}$; finally (\btref{6.1.12}) applied to
$(T',(U_a))$ in $G'$ shows $(G',T'U^{+})$ is a Tits system, whence the
double Tits system.
\end{sketch}

\begin{bt}{Remark}{6.5, 1}
Let $\Omega\subset A$, $\Omega\ne\emptyset$. The subgroup written
$P_\Omega$ in 4.1 is none other than $P_{f_\Omega}$: this follows from
(\btref{5.2.19}), and is also a consequence of the more general result proved in
(\btref{7.1.11}).
\end{bt}

\begin{bt}{Remark}{6.5, 2}
Let $(V_n)$ be a decreasing sequence of subgroups of $G$ forming a
fundamental system of neighborhoods of $1$ for a complete topological
group structure on $G$. Suppose moreover that $B$ and the $U_a$ are
closed subgroups and that the topology induced on each $U_a$ is the one
defined by the $U_{a,k}$. Then the condition of (\btref{6.4.48})(ii) is
satisfied, and so are the hypotheses of (\btref{2.8.2}): it suffices to consider
an increasing sequence of bounded convex subsets $(\Omega_n)$ of $A$ with
non-empty interior such that
$f_{\Omega_n}(a)\ge\inf\{k\in\bR\mid U_{a,k}\subset V_n\}$ for all
$a\in\Phi$; one checks at once that $P_{\Omega_n}\subset H\cdot V_n$.
\end{bt}

% =====================================================================

\btsecA{7}{7. The building of a valued root datum}

% =====================================================================

\noindent
The notation of 6.2--6.4 is retained. It was seen in (\btref{6.5}) that when
$\varphi$ is discrete it defines a Tits system of affine type, hence a
building on which $G$ acts. In this section a metric space $\Bldg$ on
which $G$ acts is constructed \emph{without} assuming $\varphi$
discrete; it is identified with the building of $G$ when $\varphi$ is
discrete, and in general retains some of its structures and properties.
Although $\Bldg$ is no longer a polysimplicial complex in the
non-discrete case, its chambers and facets are still defined.

\medskip
Before defining $\Bldg$ in 7.3, certain subgroups of $G$ are studied,
closely related to particular cases of the groups $U_f$ of 6.4, which
generalize the fixers and stabilizers studied in \S4 for the building
of a Tits system of affine type. In particular subgroups attached to the
chambers of $A$, written $B_{x,D}$, are introduced; they play the role of
the group $B$ of such a system.

\medskip
When $\varphi$ is discrete, most results of \S\S7--8 reduce to special
cases of results already established in \S\S3, 4, 5, but are recovered
here by a different method. The main novelty of the approach is that it
applies equally to non-discrete valuations, in particular to
\emph{dense} valuations, i.e.\ those for which $\Gamma_a=\varphi_a(U_a^{*})$
is dense in $\bR$ for every root $a\in\Phi$. (When $\Phi$ is irreducible
this is equivalent to density of $\Gamma_a$ in $\bR$ for a \emph{single}
root $a$.)

\medskip
If $D$ is a vector chamber of $\Phi$, one writes $\Phi^{+}_{D}$
(resp.\ $\Phi^{-}_{D}$) for the set of roots $a\in\Phi$ positive
(resp.\ negative) on $D$, $\Pi_D$ for the corresponding basis of $\Phi$,
and $U^{+}_{D}$ (resp.\ $U^{-}_{D}$) for the subgroup of $G$ generated by
the $U_a$ with $a\in\Phi^{+}_{D}$ (resp.\ $a\in\Phi^{-}_{D}$).

% ---------------------------------------------------------------------

% ---------------------------------------------------------------------

\begin{btpar}{7.1.1}
Let $\Omega$ be a non-empty subset of the affine space $A$ (\btref{6.2.6}). In
(\btref{6.4.3}) it was assigned a concave function $f_\Omega:\Phi\to\bR$ by
\[
  f_\Omega(a)=\inf\{k\in\bR\mid a(x)+k\ge0 \text{ for all } x\in\Omega\},
  \qquad a\in\Phi .
\]
From now on $A$ is identified with the vector space $V$ by the choice of
the origin $\varphi$; this allows the elements of $V^{*}$, in particular
the roots $a\in\Phi$, to be regarded as functions on $A$. Likewise an
element $x\in A$ defines a linear function $a\mapsto-a(x)$ on $\Phi$,
which is sometimes again written $x$. To the concave function
$f_\Omega$ one associates, as in (\btref{6.4.2}), the subgroup $U_{f_\Omega}$
generated by the union of the $U_{a,f_\Omega(a)}$, $a\in\Phi$, and the
subgroup $P_{f_\Omega}=H\cdot U_{f_\Omega}$; these are written more
simply $U_\Omega$ and $P_\Omega$. Since $U_{a,f_\Omega(a)}$ is the union
of the $U_{a,k}$ for $k\in\Gamma_a$, $k\ge f_\Omega(a)$, one sees that
$U_\Omega$ is generated by the union of the subgroups $U_\alpha$, where
$\alpha$ runs over the affine roots (\btref{6.2.6}) containing $\Omega$
(cf.\ (\btref{5.2.30})).
\end{btpar}

\begin{btpar}{7.1.2}
Clearly $U_\Omega$ and $P_\Omega$ depend only on the intersection of the
affine roots containing $\Omega$, that is, on the set
\[
  \cl(\Omega)=\bigcap_{\substack{a\in\Phi,\ k\in\Gamma_a\\ k\ge f_\Omega(a)}}
  \{x\in A\mid a(x)+k\ge0\}.
\]
The set $\cl(\Omega)$, again called the \emph{enclosure} of $\Omega$, is a
closed convex subset of $A$; one says $\Omega$ is \emph{closed} if
$\Omega=\cl(\Omega)$ (cf.\ (\btref{2.4.4})).
\end{btpar}

\begin{btpar}{7.1.3}
It was seen in (\btref{6.4.10}) that the subgroup $N_{f_\Omega}=N\cap U_\Omega$
has the following property: its image $\nu(N_{f_\Omega})$ in
$\widehat{W}=\nu(N)=N/H$ is the subgroup generated by the reflections
$r_{a,k}$ for $a\in\Phi$ and $k\in\Gamma'_a$ with
$\Omega\subset\{x\in A\mid a(x)+k=0\}$, and hence is identified with the
Weyl group of the root system $\Phi_\Omega=\Phi_{f_\Omega}$ consisting of
the roots $a\in\Phi$ for which there exists $k\in\Gamma'_a$ with
$a(\Omega)=\{-k\}$. One sets
\[
  N_\Omega=H\cdot N_{f_\Omega}=H\cdot(N\cap U_\Omega)=N\cap P_\Omega .
\]
When $\Omega$ reduces to a point $x$ one writes $U_x$, $P_x$, etc.
Finally, the map sending $\Omega\subset A$, $\Omega\ne\emptyset$, to
$U_\Omega$ (resp.\ $P_\Omega$, $N_\Omega$, $\Phi_\Omega$) is decreasing
for inclusion.
\end{btpar}

\begin{btpar}{7.1.4}
As in (\btref{1.3.11}), a \emph{sector} of $A$ is a set $x+D$ of translates of a
point $x\in A$ (the \emph{vertex} of the sector) by the elements of a
vector chamber $D$ (the \emph{direction} of the sector). One has
$U_{x+D}\subset U^{+}_{D}$ and $N_{x+D}=H$. If $\Omega\subset A$ contains
a sector of direction $D$ one therefore also has
$U_\Omega\subset U^{+}_{D}$ and $N_\Omega=H$. More generally, for every
non-empty subset $\Omega$ of $A$,
\begin{equation}
  P_\Omega\cap U^{+}_{D}=U_{\Omega+D},\qquad
  P_\Omega\cap U^{-}_{D}=U_{\Omega-D},
  \tag{1}
\end{equation}
whence, by (\btref{6.4.9}),
\begin{equation}
  P_\Omega=N_\Omega\cdot U_{\Omega+D}\cdot U_{\Omega-D}
  \tag{2}
\end{equation}
for every vector chamber $D$.
\end{btpar}

\begin{bt}{Proposition}{7.1.5}

\btsecB{7.1}{7.1. The subgroups \texorpdfstring{$P_{\Omega}$}{P\_Omega} and \texorpdfstring{$\widehat{P}_{\Omega}$}{P\^\_Omega}}
\btanchor{7.1.5}
Let $\Omega$ and $\Omega'$ be non-empty subsets of $A$ and let $D$ be a
vector chamber.
\begin{enumerate}
\item[(i)] If $\Omega'\subset\Omega+\overline{D}$, then
  $P_\Omega\cdot P_{\Omega'}\subset
   N_\Omega\cdot U^{-}_{D}\cdot U_{\Omega'+D}\cdot N_{\Omega'}$.
\item[(ii)] If $\Omega'\subset\Omega+\overline{D}$ and
  $\Omega\subset\Omega'-\overline{D}$, then
  $P_\Omega\cdot P_{\Omega'}=
   N_\Omega\cdot U_{\Omega-D}\cdot U_{\Omega'+D}\cdot N_{\Omega'}$.
\end{enumerate}
\end{bt}

\begin{sketch}
By (7.1.4)(2),
$P_\Omega P_{\Omega'}=N_\Omega U_{\Omega-D}U_{\Omega+D}U_{\Omega'+D}
 U_{\Omega'-D}N_{\Omega'}$.
If $\Omega'\subset\Omega+\overline{D}$ then
$U_{\Omega'+D}\supset U_{\Omega+D}$, which absorbs a factor and gives
(i) since $U_{\Omega-D}U_{\Omega'-D}\subset U^{-}_{D}$. If moreover
$\Omega\subset\Omega'-\overline{D}$ then
$U_{\Omega'-D}\subset U_{\Omega-D}$, giving (ii).
\end{sketch}

\begin{bt}{Corollary}{7.1.6}\btanchor{7.1.6}
Let $x$ be a point of $A$ and $\Omega$ a non-empty subset of $A$. There
is a vector chamber $D$ with
$P_\Omega\cdot P_x\subset N_\Omega\cdot U^{-}_{D}\cdot U_{x+D}\cdot N_x$.
\end{bt}

\begin{sketch}
Choose $y\in\Omega$ and $D$ with $x-y\in\overline{D}$; apply (\btref{7.1.5})(i).
\end{sketch}

\begin{bt}{Corollary}{7.1.7}\btanchor{7.1.7}
Let $x$ and $y$ be two points of $A$. There is a vector chamber $D$ with
$P_y\cdot P_x=N_y\cdot U_{y-D}\cdot U_{x+D}\cdot N_x$.
\end{bt}

\begin{sketch}
Choose $D$ with $x-y\in\overline{D}$; apply (\btref{7.1.5})(ii).
\end{sketch}

\begin{btpar}{7.1.8}
Let again $\Omega\subset A$, $\Omega\ne\emptyset$. Clearly
$\nu(n)(x)=x$ for all $x\in\Omega$ and $n\in N_\Omega$. Write
$\widehat{N}_\Omega$ for the \emph{fixer} of $\Omega$ in $N$:
\[
  \widehat{N}_\Omega=\{n\in N\mid \nu(n)(x)=x \text{ for all } x\in\Omega\}.
\]
One has $N_\Omega\subset\widehat{N}_\Omega$, and $\widehat{N}_\Omega$
normalizes $P_\Omega$ since $N$ permutes the affine roots (\btref{6.2.10}).
Consequently
\[
  \widehat{P}_\Omega=\widehat{N}_\Omega\cdot P_\Omega
                    =\widehat{N}_\Omega\cdot U_\Omega
\]
is a group having $P_\Omega$ and $U_\Omega$ as normal subgroups. For
every vector chamber $D$,
$\widehat{P}_\Omega=\widehat{N}_\Omega\cdot U_{\Omega+D}\cdot U_{\Omega-D}$,
and one deduces immediately from (\btref{6.1.15}) that
\begin{equation}
  \widehat{P}_\Omega\cap U^{+}_{D}=U_{\Omega+D},\quad
  \widehat{P}_\Omega\cap U^{-}_{D}=U_{\Omega-D},\quad
  \widehat{P}_\Omega\cap N=\widehat{N}_\Omega .
  \tag{1}
\end{equation}
The map $\Omega\mapsto\widehat{P}_\Omega$ is decreasing for inclusion.
Moreover, for all $n\in N$ and all non-empty $\Omega\subset A$,
\[
  nP_\Omega n^{-1}=P_{\nu(n)(\Omega)},\qquad   n\widehat{P}_\Omega n^{-1}=\widehat{P}_{\nu(n)(\Omega)} .
\]
One checks at once that in the case of a discrete valuation the
notations $\widehat{N}_\Omega$ and $\widehat{P}_\Omega$ agree with those
of \S4.
\end{btpar}

\begin{bt}{Proposition}{7.1.9}\btanchor{7.1.9}
Let $\Omega$ be a non-empty subset of $A$. If the valuation $\varphi$ is
dense, then $\widehat{N}_{\cl(\Omega)}=\widehat{N}_\Omega$ and
$\widehat{P}_{\cl(\Omega)}=\widehat{P}_\Omega$.
\end{bt}

\begin{sketch}
Since $U_{\cl(\Omega)}=U_\Omega$ (7.1.2), the first equality implies the
second. The fixed-point set $L$ of $\widehat{N}_\Omega$ in $A$ is an
affine subspace whose direction $L^{0}$ is the fixed-point subspace of
$\nuv(\widehat{N}_\Omega)$ in $\Wv$, hence the intersection of the
kernels of the roots containing it. Thus $L$ is an intersection of
closed half-spaces equipollent to affine roots and containing $\Omega$;
density of $\varphi$ makes each such half-space an intersection of
affine roots, so $L$ is closed and $\cl(\Omega)\subset L$.
\end{sketch}

\noindent
When $\varphi$ is discrete the equalities of (\btref{7.1.9}) still hold if
$\Omega$ contains a non-empty open set, but need not hold in general.

\begin{bt}{Remark}{7.1.10}\btanchor{7.1.10}
Let $x\in A$; then $\widehat{N}_x=\nu^{-1}(\widehat{W}_x)$, where
$\widehat{W}_x$ is the stabilizer of $x$ in $\widehat{W}=\nu(N)$. When
$G$ is generated by $H$ and the $U_a$ (the ``simply connected'' case) and
$\varphi$ is discrete, $\widehat{W}=W$ is the affine Weyl group of the
root system $\Sigmav$ (\btref{1.3.8}) and is generated by the reflections in the
walls of $A$. It is known that $\widehat{W}_x$ is then generated by the
reflections in the walls of $A$ passing through $x$ (\btref{1.3.5}). Hence
$\widehat{N}_x=N_x$ and $\widehat{P}_x=P_x$.

This is no longer necessarily correct when $\varphi$ is not discrete,
even assuming $G$ generated by the $U_a$ and all $\Gamma_a$ groups.
\end{bt}

\noindent\textbf{Examples.}
\begin{enumerate}
\item Suppose $\Phi=\{a,-a\}$ of type $A_1$ and $\varphi$ special. Then
  $W$ is the group of transformations $x\mapsto\varepsilon x+\gamma a^{\vee}$
  with $\varepsilon=\pm1$ and $\gamma\in\widetilde{\Gamma}_a$, where
  $\widetilde{\Gamma}_a$ is the group generated by $\Gamma_a$ (\btref{6.2.20}).
  The points of $A$ lying on a wall are those of the form
  $\tfrac12\gamma a^{\vee}$ with $\gamma\in\Gamma_a$, whereas the points
  stabilized by some non-trivial element of $W$ are the
  $\tfrac12\gamma a^{\vee}$ with $\gamma\in\widetilde{\Gamma}_a$. So if
  $\Gamma_a$ is not a group there are points $x\in A$ with
  $\widehat{N}_x\ne N_x$, even assuming $\widehat{W}=W$.
\item Let $K$ be a field with a valuation $\omega$, let
  $\Gamma=\omega(K^{*})$, and suppose $\Gamma/2\Gamma$ is a vector space
  of dimension $>1$ over $\bF_2$. Take for $G$ the group of
  $K$-rational points of a simple algebraic group split over $K$ of type
  $G_2$ (necessarily simply connected), with the valued root datum
  defined in (\btref{6.2.3})\,b). Let $\lambda,\mu\in\Gamma$ have images in
  $\Gamma/2\Gamma$ linearly independent over $\bF_2$, and put
  $x=\tfrac12(\lambda a^{\vee}+\mu b^{\vee})$, where $a,b$ are the two
  roots of a basis of the root system of $G$. One checks easily that the
  symmetry about $x$ lies in the stabilizer of $x$ in $\widehat{W}=W$,
  yet no wall of $A$ passes through $x$. Hence
  $\widehat{N}_x\ne N_x$.
\end{enumerate}

\begin{bt}{Proposition}{7.1.11}\btanchor{7.1.11}
Let $\Omega\subset A$, $\Omega\ne\emptyset$. Then
$\displaystyle\widehat{P}_\Omega=\bigcap_{x\in\Omega}\widehat{P}_x$.
\end{bt}

\begin{sketch}
One first proves
$\widehat{P}_\Omega\cap\widehat{P}_x=\widehat{P}_{\Omega\cup\{x\}}$,
using the decomposition $g=nvu$ of (7.1.8) together with (\btref{6.1.15}),
(DR\,6), (7.1.8)(1) and (\btref{6.4.9}); induction then gives
$\widehat{P}_{\{x_1,\dots,x_k\}}=\widehat{P}_{x_1}\cap\cdots\cap\widehat{P}_{x_k}$.
It remains to treat an increasing directed family $(\Omega_i)$ with union
$\Omega$: one may assume all $\Omega_i$ contain a common point $x$, so
$\widehat{N}_x/H$ is finite; passing to a cofinal subfamily and using
(DR\,6), the components of $g$ are independent of $i$, and lie in $H$,
$\bigcap_i U_{\Omega_i+D}=U_{\Omega+D}$ and
$\bigcap_i U_{\Omega_i-D}=U_{\Omega-D}$ respectively.
\end{sketch}

\begin{bt}{Corollary}{7.1.12}\btanchor{7.1.12}
Let $\Omega$ and $\Omega'$ be non-empty subsets of $A$. Then
$\widehat{P}_\Omega\cap\widehat{P}_{\Omega'}=\widehat{P}_{\Omega\cup\Omega'}$.
\end{bt}

\begin{bt}{Remark}{7.1.13}\btanchor{7.1.13}
For non-empty $\Omega,\Omega'\subset A$ one has
$f_{\cl(\Omega\cup\Omega')}=\sup(f_\Omega,f_{\Omega'})$. One might try to
generalize (\btref{7.1.12}) by showing that $U_f\cap U_g\subset HU_{\sup(f,g)}$
whenever $f$ and $g$ are concave on $\Phi$ --- which would imply that
$\sup(f,g)$ is concave as well. But this assertion is in general false.
For example, take for $G$ the group of $K$-rational points of a split
simple algebraic group of type $G_2$ over a field $K$ with a
non-improper valuation, equipped with the root datum of
(\btref{6.2.3})\,b). Let $(a,b)$ be a basis of $\Phi$ with $b$ the long root, and
put, for $c\in\Phi$,
\[
\begin{aligned}
  f(c)&=0 \ \text{ if } c=\pm a,\ \pm(2b+3a), &\quad f(c)&=0+ \ \text{ otherwise};\\
  g(c)&=0 \ \text{ if } c=\pm(b+a),\ \pm(b+3a), &\quad g(c)&=0+ \ \text{ otherwise}.
\end{aligned}
\]
One checks easily that $f$ and $g$ are concave and that
$U_f\cap U_g\not\subset HU_{\sup(f,g)}$: indeed
$N_{\sup(f,g)}\subset H$, whereas $N_f\cap N_g$ contains an element whose
image in $\Wv$ equals $-1$.
\end{bt}

% ---------------------------------------------------------------------

\btsecB{7.2}{7.2. Chambers and facets of \texorpdfstring{$A$}{A}}

% ---------------------------------------------------------------------

\begin{btpar}{7.2.1}
Let $X$ be a set and $\mathscr{F}$ a filter base on $X$; the
\emph{germ} of $\mathscr{F}$ is the filter generated by $\mathscr{F}$.
If $Y\subset X$, one says the germ $\gamma$ of $\mathscr{F}$ is
\emph{contained in} $Y$ if some element of $\gamma$ (or of
$\mathscr{F}$) is contained in $Y$. More generally, if there is
$Z\in\gamma$ with $Y\cap Z=\emptyset$ one puts $\gamma\cap Y=\emptyset$;
otherwise the intersections $Z\cap Y$, $Z\in\gamma$, form a filter base
on $X$ whose germ is written $\gamma\cap Y$, and it is contained in $Y$.

If $X$ is a topological space, the \emph{closure} of $\gamma$, written
$\overline{\gamma}$, is the germ associated with the filter base formed
by the closures of the elements of $\gamma$ (or of $\mathscr{F}$). One
says $\gamma$ \emph{has non-empty interior} if the interiors of the
elements of $\gamma$ are non-empty; these then form a filter base whose
germ is called the \emph{interior} of $\gamma$. One says $\gamma$ is
\emph{open} if it equals its interior.
\end{btpar}

\begin{btpar}{7.2.2}
Let $\mathscr{F}$ be a filter base on the affine space $A$. The family of
the $U_\Omega$ for $\Omega\in\mathscr{F}$ is increasing and directed for
inclusion, so its union is a subgroup of $G$ depending only on the germ
$\gamma$ of $\mathscr{F}$; it is written $U_\gamma$. The subgroups
$N_\gamma$, $P_\gamma$, $\widehat{N}_\gamma$, $\widehat{P}_\gamma$ and
the root system $\Phi_\gamma$ are defined analogously.
\end{btpar}

\begin{btpar}{7.2.3}
Let $D$ be a vector chamber. The sectors of $A$ of direction $D$ form a
filter base; the associated germ is written $\gamma(D)$ and called the
\emph{sector germ} of direction $D$ of $A$. Immediately,
\[
  U_{\gamma(D)}=U^{+}_{D}
  \qquad\text{and}\qquad
  P_{\gamma(D)}=\widehat{P}_{\gamma(D)}=H\cdot U^{+}_{D}.
\]
\end{btpar}

\begin{btpar}{7.2.4}
Let $x$ be a point of $A$ and $E$ a non-empty convex cone with vertex $0$
in $V$, open in the vector subspace it spans. Write $\gamma(x,E)$ for the
germ of the filter base formed by the subsets $X$ of $A$ with the
following two properties: $X$ is an intersection of affine roots and of
complements of affine roots, and there is a neighborhood $Y$ of $x$ in
$A$ such that $X\cap(x+E)\supset Y\cap(x+E)$. One puts
\[
  U_{x,E}=U_{\gamma(x,E)},\qquad P_{x,E}=P_{\gamma(x,E)},\qquad\text{etc.}
\]
A \emph{facet} of $A$ is a germ of the form $\gamma(x,E)$.\tnote{The
scan of this sentence was not fully legible; the delimiting clause on
the admissible cones $E$ should be checked against the original before
the definition is relied upon.} For a germ $\gamma(x,E)$ to be a facet
it is necessary and sufficient that, for every root $a\in\Phi_x$, the
convex cone $E$ either be contained in the hyperplane
$\{v\in V\mid a(v)=0\}$ or not meet that hyperplane.
\end{btpar}

\begin{btpar}{7.2.4 (continued)}
This is in particular the case when $E$ is a vector facet of $\Phi$; one
then puts $F_{x,E}=\gamma(x,E)$. If $E$ is a vector chamber, the facet
$F_{x,E}$ is called a \emph{chamber} of $A$ and is also written
$C_{x,E}$.

Two facets $F_{x,E}$ and $F_{x',E'}$ are equal if and only if they are
contained in the same affine roots, equivalently if and only if
$P_{x,E}=P_{x',E'}$. When $\varphi$ is dense this forces $x=x'$. In all
cases $F_{x,E}=F_{x,E'}$ if and only if the canonical projections of $E$
and $E'$ into the dual of the vector subspace of $V$ spanned by $\Phi_x$
lie in a single vector facet of $\Phi_x$. The preimage $E''$ of this
vector facet in $V$ --- also called, by abuse of language, a vector facet
of $\Phi_x$ in $V$ --- is therefore the largest convex cone of $V$ with
$F_{x,E}=F_{x,E''}$, and $E''\mapsto F_{x,E''}$ is a bijection from the
set of vector facets of $\Phi_x$ in $V$ onto the set of facets of the
form $F_{x,E}$. The vector facet $E''$ is called the \emph{direction} of
$F_{x,E''}$ at $x$, and the \emph{dimension} of $F_{x,E''}$ is the
dimension of the vector subspace spanned by $E''$. One checks easily
that this dimension does not depend on the choice of $x$, and that the
set of $X\in F_{x,E''}$ which are convex subsets of $x+E''$ all of whose
points are internal is a base of the filter $F_{x,E''}$. A facet $F$ is
a chamber if and only if $\dim F=\dim V$.

Every facet $F$ is contained in the closure of at least one chamber $C$;
one then has $P_C\subset P_F$. More generally, a facet $F$ is contained
in the closure of a facet $F'$ if and only if $P_F\supset P_{F'}$.
\end{btpar}

\begin{btpar}{7.2.5}
When $\varphi$ is discrete one recovers, up to a canonical
identification, the notions of chamber and facet already seen: a facet
(resp.\ chamber) in the sense of (7.2.4) is precisely the filter of all
subsets of $A$ containing a facet (resp.\ chamber) in the sense of \S2.

When $\varphi$ is dense one sees immediately that $\gamma(x,E)$ is the
filter of all subsets $X$ of $A$ for which there is a neighborhood $Y$
of $x$ in $A$ with $X\cap(x+E')\supset Y\cap(x+E')$, where $E'$ is the
direction of $\gamma(x,E)$ at $x$.

In all cases, the \emph{center} of a facet $F_{x,E}$ is the limit along
the filter $\gamma(x,E)$ of the barycenters of the compact subsets
belonging to $\gamma(x,E)$. When $\varphi$ is discrete one recovers the
barycenter of the facet $F_{x,E}$; when $\varphi$ is dense, the center of
$F_{x,E}$ is $x$.
\end{btpar}

\begin{btpar}{7.2.6}
Let $x\in A$ and let $D$ be a vector chamber. Since the chamber
$C_{x,D}$ is open in $A$ (in the sense of (7.2.1)), every element of
$\widehat{N}_{x,D}$ fixes all points of a non-empty open subset of $A$
and hence belongs to $H$. Therefore
$\widehat{P}_{x,D}=P_{x,D}$, and one puts
\[
  B_{x,D}=P_{x,D}=\widehat{P}_{x,D}.
\]
One checks immediately that $U_{x,D}=U_{f_{x,D}}$, where
$f_{x,D}:\Phi\to\Rt$ is defined by
\begin{equation}
  \begin{cases}
    f_{x,D}(a)=-a(x) & \text{if } a\in\Phi^{+}_{D},\\[2pt]
    f_{x,D}(a)=(-a(x))+ & \text{if } a\in\Phi^{-}_{D}.
  \end{cases}
  \tag{1}
\end{equation}
One has $\Phi_{x,D}=\emptyset$, $N_{x,D}=H$ and, taking account of
$f_{x,D}(2a)=2f_{x,D}(a)$ when $a,2a\in\Phi$, (\btref{6.4.9}) gives
\begin{equation}
  B_{x,D}=H\cdot U_{x,D}
   =\prod_{a\in\Phi^{+\mathrm{red}}_{D}}U_{a,-a(x)}\cdot H\cdot
    \prod_{a\in\Phi^{-\mathrm{red}}_{D}}U_{a,-a(x)+}.
  \tag{2}
\end{equation}
More generally, if $\Delta$ is another vector chamber,
\begin{equation}
  B_{x,D}=\prod_{a\in\Phi^{+\mathrm{red}}_{\Delta}}U_{a,f_{x,D}(a)}\cdot H\cdot
          \prod_{a\in\Phi^{-\mathrm{red}}_{\Delta}}U_{a,f_{x,D}(a)}.
  \tag{3}
\end{equation}
Recall (\btref{6.4.9}) that (2) and (3) are in fact decompositions of $B_{x,D}$
as a (set-theoretic) direct product, and that the position of the factor
$H$, placed here ``in the middle'', is immaterial.

More generally still, let $F=\gamma(x,E)$ be a facet. Then $U_F=U_{f_F}$,
where $f_F:\Phi\to\Rt$ is defined by
\[
\begin{aligned}
f_F(a)&=-a(x) &&\text{if }a\text{ is positive or zero on }E,\\
f_F(a)&=-a(x)+ &&\text{if }a\text{ is strictly negative on }E.
\end{aligned}
\]
\end{btpar}

\begin{btpar}{7.2.7}
Let $x\in A$ and consider the function $f^{*}_{x}:\Phi\to\Rt$ associated
as in (\btref{6.4.23}) with the linear function $f_x:a\mapsto-a(x)$. One has
$f^{*}_{x}(a)=-a(x)+$ for all $a\in\Phi_x$. Resuming the notation of
(\btref{6.4.18}) and putting $H^{*}_{x}=H_{f_x,f^{*}_{x}}$ --- a normal subgroup
of $H$ --- it was seen in (\btref{6.4.23}) and (\btref{6.4.27}) that
\[
  P^{*}_{x}=H^{*}_{x}\cdot U_{f^{*}_{x}}
\]
is a normal subgroup of $P_x=H\cdot U_{f_x}$; moreover the quotient group
$\overline{G}_x=P_x/P^{*}_{x}$ carries a root datum
$(\overline{T},(\overline{U}_a)_{a\in\Phi_x})$ of type $\Phi_x$, where
$\overline{T}$ (resp.\ $\overline{U}_a$) is the canonical image of $H$
(resp.\ $U_{a,-a(x)}$) in $\overline{G}_x$. One then sees easily that,
for a vector chamber $D$, \emph{the subgroup $B_{x,D}$ is the preimage in
$P_x$ of the minimal parabolic subgroup
$\overline{T}\,\overline{U}^{+}_{D}$ of $\overline{G}_x$} (\btref{6.1.12}),
generated by $\overline{T}$ and the $\overline{U}_a$ for
$a\in\Phi_x\cap\Phi^{+}_{D}$. More generally the $P_{x,E}$ are the
preimages of the parabolic subgroups of $\overline{G}_x$ associated with
the various vector facets of $\Phi_x$.

Note also that $P^{*}_{x}$ is a normal subgroup of $\widehat{P}_x$ and
that the canonical injection of $P_x/P^{*}_{x}$ into
$\widehat{P}_x/P^{*}_{x}$ is adapted to the Tits system of
$\overline{G}_x$ ((\btref{1.2.13}) and (\btref{6.1.12})).
\end{btpar}

% ---------------------------------------------------------------------

\btsecB{7.3}{7.3. Iwasawa decomposition and Bruhat decomposition}

% ---------------------------------------------------------------------

\noindent
This subsection proves that the subgroups $B_{x,D}$ attached to the
chambers of $A$ give rise to ``Iwasawa decompositions''
$G=U^{+}_{D}\cdot N\cdot B_{x,D}$ and to ``Bruhat decompositions''
$G=B_{x,D}\cdot N\cdot B_{x,D}$. When $\varphi$ is discrete this simply
recovers earlier results ((\btref{1.2.7}) and (\btref{5.1.3})) by another method. Note
however that the ``Schiffmann formula'' (\btref{7.3.3}) has not yet been proved
even in the discrete case.

\begin{bt}{Proposition}{7.3.1}\btanchor{7.3.1}
Let $D$ and $D'$ be two vector chambers and let $x$ be a point of $A$.
\begin{enumerate}
\item[(i)] $G=U^{+}_{D}\cdot N\cdot B_{x,D'}$.
\item[(ii)] More precisely, the natural map from $\widehat{W}=N/H$ to the
  set of double cosets $U^{+}_{D}\backslash G/B_{x,D'}$ is bijective.
\end{enumerate}
\end{bt}

\begin{sketch}
Put $U=U^{+}_{D}$, $B=B_{x,D'}$, $Z=U\cdot N\cdot B$ and
$f(a)=f_{x,D'}(a)$; write $B_a$ for the subgroup generated by
$U_{a,f(a)}$, $U_{-a,f(-a)}$ and $H$, so that
$B_a=HU_{a,f(a)}U_{-a,f(-a)}=HU_{-a,f(-a)}U_{a,f(a)}$ by (\btref{6.3.2}). The
proof of (i) has three steps.
\begin{enumerate}
\item[(1)] For $a\in\Phi$, the group $L_a$ generated by $U_a$, $U_{-a}$
  and $T$ is contained in $Z_a=U_a\cdot T\cdot\{1,m_a\}\cdot B_a$, where
  $m_a$ is any element of $M^{0}_{a}\subset N$ (cf.\ (\btref{6.1.2})(2)). By
  (\btref{6.1.2})(7) it suffices that $m_a u\in Z_a$ for $u\in U_a$, which is
  clear if $\varphi_a(u)\ge f(a)$ and follows from the decomposition
  $u=v'mv''$ and (\btref{6.2.2})(V\,5 bis) otherwise.
\item[(2)] For $a$ a simple root for $D$: $U_{-a}\cdot Z\subset Z$. One
  uses (DR\,2) --- the subgroup $U'_a$ of $U^{+}$ generated by the $U_b$,
  $b\in\Phi^{+}_{D}$, $b\ne a$, is normalized by $U_a$ and $U_{-a}$, with
  $U^{+}=U'_a\cdot U_a$ (\btref{6.1.6}) --- together with step (1), reducing to
  $un\in Z$ for $n\in N$, $u\in U_{-a}$, which follows by applying (1) to
  the root $-b$ where $b=-\nuv(n^{-1})(a)$.
\item[(3)] Hence $Z$ is stable under left multiplication by the $U_{-a}$
  for $a\in\Pi_D$, and evidently also by the $U_a$ for $a\in\Phi^{+}_{D}$
  and by $T$, hence by all of $G$ (\btref{6.1.4}). This proves (i).
\end{enumerate}
\end{sketch}

% =====================================================================
\par\endgroup

% ---------------------------------------- bruhattits7standalone.tex
\begingroup\providecommand{\Imm}{}
\renewcommand{\Imm}{\mathscr{I}}
\providecommand{\hImm}{}
\renewcommand{\hImm}{\widehat{\mathscr{I}}}
\providecommand{\hP}{}
\renewcommand{\hP}{\widehat{P}}
\providecommand{\hN}{}
\renewcommand{\hN}{\widehat{N}}
\providecommand{\hW}{}
\renewcommand{\hW}{\widehat{W}}
\providecommand{\vnu}{}
\renewcommand{\vnu}{{}^{v}\nu}
\providecommand{\vW}{}
\renewcommand{\vW}{{}^{v}W}
\providecommand{\vpi}{}
\renewcommand{\vpi}{{}^{v}\pi}
\providecommand{\red}{}
\renewcommand{\red}{\mathrm{red}}
\providecommand{\cl}{}
\renewcommand{\cl}{\operatorname{cl}}
\providecommand{\Card}{}
\renewcommand{\Card}{\operatorname{Card}}
\providecommand{\bU}{}
\renewcommand{\bU}{\overline{U}}
\providecommand{\bT}{}
\renewcommand{\bT}{\overline{T}}
\providecommand{\bB}{}
\renewcommand{\bB}{\overline{B}}
\providecommand{\bG}{}
\renewcommand{\bG}{\overline{G}}
\providecommand{\bu}{}
\renewcommand{\bu}{\overline{u}}
\providecommand{\bw}{}
\renewcommand{\bw}{\overline{w}}
\providecommand{\bD}{}
\renewcommand{\bD}{\overline{D}}
\providecommand{\tE}{}
\renewcommand{\tE}{\widetilde{E}}
\newenvironment{stmt}[2]%
  {\medskip\par\noindent\textit{#1} \textbf{(#2)}.\ ---\ \begingroup\itshape}%
  {\endgroup\par\medskip}
\providecommand{\num}{}
\renewcommand{\num}[1]{\medskip\par\noindent\textbf{(#1)}\ }
\providecommand{\pf}{}
\renewcommand{\pf}{\par\noindent\emph{Proof.}\ }
\setlength{\parindent}{1.4em}
\bigskip

\num{7.3.1\,(4)} If $n,n'\in N$ satisfy $n'\in U_D^{+}nB$, put
$\Delta=\vnu(n^{-1})(D)$. By (7.2.6)\,(3),
$n'^{\,-1}n\in B\,U_\Delta^{+}=H\,(U_\Delta^{-}\cap B)\,U_\Delta^{+}$, so that
$n'^{\,-1}n\in H$ by (\btref{6.1.15}). This gives (ii).

\begin{stmt}{Corollary}{7.3.2}\btanchor{7.3.2}
(i) The natural map $N/N_x\to U_D^{+}\backslash G/P_x$ (resp.\
$N/\hN_x\to U_D^{+}\backslash G/\hP_x$) is bijective.

(ii) If $x$ is a special point of $A$ \textup{(\btref{6.2.13})}, the natural map
$T/H\to U_D^{+}\backslash G/P_x$ is bijective.
\end{stmt}

\pf For (i) it is enough, given (\btref{7.3.1})\,(i), to see that $n'\in U_D^{+}nP_x$
forces $n\in n'N_x$ (and likewise with hats). With $\Delta=\vnu(n^{-1})(D)$ one has
$n'^{\,-1}n\in P_x\,U_\Delta^{+}=N_x\,(U_\Delta^{-}\cap P_x)\,U_\Delta^{+}$ by
(7.1.4) and (7.1.8), and the claim follows from (\btref{6.1.15}). If $x$ is special, $N$ is
the semi-direct product of $N_x$ by $T$, whence (ii); note that then $\hP_x=P_x$.
\qed

\begin{stmt}{Proposition}{7.3.3}\btanchor{7.3.3}
Let $x\in A$, let $D$ be a vector chamber, $a\in\Phi_D^{-}$, $v\in U_a$ and $n\in N$.
If $v\in U_D^{+}nP_x$, then either
\[
  \varphi_a(v)\geqslant-a(x),\quad v\in\hP_x,\quad n\in\hN_x,
\]
or else
\[
  a\bigl(\nu(n)(x)-x\bigr)>0
  \quad\text{and}\quad
  \varphi_a(v)=-a(x)-\tfrac12\,a\bigl(\nu(n)(x)-x\bigr)
\]
\textup{(``Schiffmann's formula'')}.
\end{stmt}

\pf Suppose $k=\varphi_a(v)<-a(x)$. By (\btref{6.2.2})\,(3) one may write $v=u'mu''$ with
$u',u''\in U_{-a,-k}$ and $m\in M_{a,k}$. As $U_{-a,-k}\subset\hP_x\cap U_D^{+}$,
(\btref{7.3.2}) gives $m\in n\hN_x$, whence $\nu(n)(x)=\nu(m)(x)=x-ka^{\vee}-a(x)a^{\vee}$
and $a(\nu(n)(x)-x)=-2k-2a(x)$. \qed

\smallskip
This is the analogue of a formula of G.~Schiffmann for real semi-simple Lie groups
of real rank one ([\btbib{33}]).

\begin{stmt}{Theorem}{7.3.4}\btanchor{7.3.4}
Let $x,x'\in A$ and let $D,D'$ be vector chambers.

(i) $G=B_{x,D}\,N\,B_{x',D'}$.

(ii) More precisely, the canonical map of $N$ into the set of double cosets modulo
$B_{x,D}$ and $B_{x',D'}$ induces a bijection
$\hW=N/H\;\xrightarrow{\ \sim\ }\;B_{x,D}\backslash G/B_{x',D'}$.
\end{stmt}

\num{7.3.5}\btanchor{7.3.5} \emph{Notation for the proof.} For $x,x'\in A$ and vector chambers
$D,D'$: a hyperplane of $A$ through $x$ and $x'$ parallel to a wall of $\Phi$ cannot
meet the sector $x'+D'$, so there is a unique vector chamber $D(x;x',D')$ such that
the sector $x+D(x;x',D')$ contains the intersection of a neighborhood of $x'$ with
$x'+D'$. Let $w(x,D;x',D')$ be the unique element of $\vW$ with
\begin{equation}
  D(x;x',D')=w(x,D;x',D')\,.\,D \tag{1}
\end{equation}
and let $\lambda(x,D;x',D')$ be its length in $\vW$ with respect to the reflections
in the walls of $D$.

\begin{stmt}{Lemma}{7.3.6}\btanchor{7.3.6}
Let $x,x'\in A$, let $D,D'$ be vector chambers and let $L$ be an open half-line with
origin $o$ in $V$, contained in $D$. Let $g\in G$ and $n\in N$ with
$g\in B_{x,D}\,n\,B_{x',D'}$.

(i) One of the following holds:

\smallskip
\noindent\textup{(i\,1)} $g\in B_{z,D}\,n\,B_{x',D'}$ for every $z\in x+L$;

\smallskip
\noindent\textup{(i\,2)} there are a point $y\in x+L$ and an element $n'\in N$ with
\begin{gather*}
  g\in B_{z,D}\,n\,B_{x',D'}\quad\text{for all }z\in\,]x,y[,\qquad
  g\in B_{y,D}\,n'\,B_{x',D'},\\
  \lambda\bigl(y,D;n'.x',\vnu(n')(D')\bigr)>\lambda\bigl(x,D;n.x',\vnu(n)(D')\bigr).
\end{gather*}

(ii) $g\in B_{z,D}\,N\,B_{x',D'}$ for every $z\in x+L$.
\end{stmt}

\num{7.3.7}\btanchor{7.3.7} \emph{How the lemma gives the theorem.} Let $g\in G$. By (\btref{7.3.1}) there
are $n\in N$, $u\in U_D^{-}$ with $g\in unB_{x',D'}$; write
$u=\prod_{a\in\Phi_D^{+\red}}u_a$, $u_a\in U_{-a}$. Since $a(v)>0$ for
$a\in\Phi_D^{+\red}$, $v\in L$, one finds $y\in x-L$ with
$-a(y)+\varphi_{-a}(u_a)>0$ for all such $a$; then $u_a\in U_{-a,a(y)+}$ and
$u\in B_{y,D}$. Hence $g\in B_{y,D}\,n\,B_{x',D'}$ and (\btref{7.3.6})\,(ii) (applied with
$y$ in place of $x$ and $z=x$) yields $g\in B_{x,D}\,N\,B_{x',D'}$, i.e.\ (i).

For (ii), put $B=B_{x,D}$, $B'=B_{x',D'}$ and suppose $n'\in BnB'$ with $n,n'\in N$.
Then $1\in Bn''B''$ with $n''=nn'^{\,-1}$, $B''=n'B'n'^{\,-1}$, so $n''\in BB''$.
Put $x''=n'.x'$, $D''=\vnu(n')(D')$ (so $B''=B_{x'',D''}$),
$\Delta=D(x;x'',D'')$ and $\gamma=(x+\Delta)\cap C_{x'',D''}$; then $B''=P_\gamma$.
By (\btref{7.1.5}), $BB''\subset U_\Delta^{+}H\,U_\Delta^{-}$, and since
$N\cap U_\Delta^{+}HU_\Delta^{-}=H$ (6.1.15\,(a)) one gets $n''\in H$, i.e.\
$n'\in nH$.

\num{7.3.8}\btanchor{7.3.8} \emph{Proof of the lemma.} Assume (i); then (ii) follows by iterating:
in case (i\,1) there is nothing to do, and in case (i\,2) one applies (i) again with
$y,n'$ in place of $x,n$. Since lengths in $\vW$ are bounded, after finitely many
steps one lands in case (i\,1). One obtains points $y_0=y,y_1,\dots,y_k$ with
$y_{i+1}\in y_i+L$ and elements $n_1=n',n_2,\dots,n_{k+1}$ of $N$ such that
$g\in B_{z,D}\,n\,B_{x',D'}$ for $z\in[x,y_0[$, $g\in B_{z,D}\,n_i\,B_{x',D'}$ for
$z\in[y_{i-1},y_i[$, and $g\in B_{z,D}\,n_{k+1}\,B_{x',D'}$ for $z=y_k$ or
$z\in y_k+L$; whence (ii).

There remains (i). Replacing $g$ by $gn^{-1}$, $x'$ by $n.x'$ and $D'$ by
$\vnu(n).D'$ one reduces to $n=1$, i.e.\ $g\in B_{x,D}B_{x',D'}$. Put
$\Delta=D(x;x',D')$ and
\[
  \Psi=\Phi_D^{-\red}\cap\Phi_\Delta^{-},\qquad
  \Psi'=\Phi_D^{-\red}\cap\Phi_\Delta^{+}.
\]
By (\btref{6.4.9}) every $b\in B_{x,D}$ is uniquely
$b=\prod_{a\in\Phi_D^{+\red}}u_a\cdot\prod_{a\in\Psi}u_a\cdot
   \prod_{a\in\Psi'}u_a\cdot h$
with $u_a\in U_{a,f_{x,D}(a)}$, $h\in H$.

For $a\in\Phi_D^{+}$ and $z\in x+L$ one has $a(z)>a(x)$, hence
$U_{a,f_{x,D}(a)}\subset B_{z,D}$. By definition of $\Delta$ one has
$x'\in x+\bD$ (with $\bD=\overline{\Delta}$) and there is $v\in D'$ with
$x'+v\in x+\Delta$; so for $a\in\Phi_\Delta^{+}$, $a(x')\geqslant a(x)$ and
$a(x')+a(v)>a(x)$, whence $a(x')>a(x)$ when
$a\in\Phi_\Delta^{+}\cap\Phi_{D'}^{-}$. It follows that
$U_{a,f_{x,D}(a)}\subset B_{x',D'}$ for all $a\in\Phi_\Delta^{+}$, in particular for
$a\in\Psi'$. Multiplying $g$ by an element lying in every $B_{z,D}$, $z\in x+L$, one
reduces to $g\in uB_{x',D'}$ with $u=\prod_{a\in\Psi}u_a$, $u_a\in U_{a,-a(x)+}$.

If $u=1$, (i\,1) holds. (This is the case when $\Delta=-D$, i.e.\ when
$\lambda(x,D;x',D')$ is maximal, since then $\Psi=\varnothing$.)

If $u\neq1$, there is a unique $y\in x+L$ with $a(y)\geqslant-\varphi_a(u_a)$ for
all $a\in\Psi$ and with equality for at least one $a\in\Psi$. Then
\begin{equation}
  u\in U_{y-\Delta}\cap U_{y-D}\subset U_y,\qquad u\notin U_{y^{*}}, \tag{1}
\end{equation}
$y^{*}$ being the concave function attached to the linear function $y$ as in
(\btref{6.4.23}). By (\btref{6.4.23}) and (\btref{6.4.27}) there is a subgroup $H^{*}\subset H$ with
$H^{*}U_{y^{*}}$ normal in $HU_y$ and with
$\bG=HU_y/H^{*}U_{y^{*}}$ carrying a generating root datum
$(\bT,(\bU_a)_{a\in\Phi_y})$ of type $\Phi_y$, where $\bT$ is the image of $H$,
$\bU_a$ the image of $U_{a,y(a)}$, and $\Phi_y=\{a\in\Phi: a(y)\in\Gamma'_a\}$.
Putting $\Delta'=D(y;x',D')$, (\btref{6.1.15}) provides $u'\in U_{y+D}$,
$u''\in U_{y+\Delta'}$ and $n'\in N_y$ with
\begin{equation}
  u\in U_{y^{*}}\,u'\,n'\,u''. \tag{2}
\end{equation}

Since $U_{y+D}\cup U_{y^{*}}\subset B_{y,D}$ and
$U_{y+\Delta'}\subset B_{x',D'}$ (because $C_{x',D'}\cap(y+\Delta')\neq\varnothing$),
one gets $u\in B_{y,D}n'B_{x',D'}$ and $g\in B_{y,D}n'B_{x',D'}$: the second
relation of (i\,2). For $z\in\,]x,y[$ one has $a(z)>-\varphi_a(u_a)$ for all
$a\in\Psi$, so $u\in B_{z,D}$ and $g\in B_{z,D}B_{x',D'}$: the first relation.

It remains to prove
\begin{equation}
  \lambda\bigl(y,D;n'.x',\vnu(n')(D')\bigr)>\lambda(x,D;x',D'). \tag{3}
\end{equation}

\emph{Step 1.}
\begin{equation}
  \lambda(y,D;x',D')\geqslant\lambda(x,D;x',D'). \tag{4}
\end{equation}
Indeed, let $\mathscr{M}$ be the set of hyperplanes of $V$ that are kernels of
elements of $\Phi$. The length of $w\in\vW$ is the number of members of
$\mathscr{M}$ separating $D$ and $w(D)$ ([\btbib{5}], ch.~VI, \S1, no.~6). Hence
$\lambda(z,D;x',D')$ is, for $z\in A$, the cardinality of the set $\mathscr{M}_z$ of
$M\in\mathscr{M}$ such that $z+M$ separates $z+D$ from a neighborhood of $x$ in
$x'+\overline{D'}$. For fixed $M$, the set of $z\in A$ with $M\in\mathscr{M}_z$ is
the open half-space $x'+M+D$ or the closed one $x'+M+\bD$, and its intersection with
$x+L$ is a half-line. Hence $\mathscr{M}_y\supset\mathscr{M}_x$, giving (4).
In fact this shows more: $w'=w(y,D;x',D')$ dominates $w=w(x,D;x',D')$, i.e.\
$w'=ww''$ for some $w''\in\vW$ with $\ell(w')=\ell(w)+\ell(w'')$.

\emph{Step 2.}
\begin{equation}
  \lambda(y,D;x',D')>\lambda(x,D;x',D')\quad\text{whenever }n'\in H. \tag{5}
\end{equation}
Equality in (4) forces $w'=w$, i.e.\ $\Delta'=\Delta$. The image $\bu$ of $u$ in
$\bG$ lies in $\bU_D^{-}\cap\bU_\Delta^{-}$, while the images $\bu'$, $\bu''$ of
$u'$, $u''$ lie in $\bU_D^{+}$, $\bU_{\Delta'}^{+}$ respectively. If $n'\in H$ and
$\Delta=\Delta'$ then
$\bu\in\bU_D^{-}\cap\bU_\Delta^{-}\cap\bU_D^{+}\bT\bU_\Delta^{+}$, whence $\bu=1$ by
(DR\,6) and (\btref{6.1.6}); so $u\in U_{y^{*}}$, contradicting (1).

\emph{Step 3.} By (4) and (5), (3) will follow from
\begin{equation}
  \lambda\bigl(y,D;n'(x'),\vnu(n')(D')\bigr)\geqslant\lambda(y,D;x',D'),
  \text{ with equality only if }n'\in H. \tag{6}
\end{equation}
Since $n'\in N_y$ one checks that
$w\bigl(y,D;n'.x',\vnu(n')(D')\bigr)=\vnu(n')\,w(y,D;x',D')$. Setting
$t=\vnu(n')$ and $w'=w(y,D;x',D')$, (6) is equivalent to
\begin{equation}
  \ell(tw')\geqslant\ell(w'),\text{ with equality only if }t=1. \tag{7}
\end{equation}

Identify the Weyl group $\vW'$ of $\Phi_y$ with a subgroup of $\vW$, and identify a
vector chamber of $\Phi_y$ (a simplicial cone in a quotient of $V$) with its inverse
image in $V$. Let $E$ be the vector chamber of $\Phi_y$ containing $D$; let $R$
(resp.\ $R'$) be the reflections in the walls of $D$ (resp.\ $E$), generating $\vW$
(resp.\ $\vW'$), and $w_0$ (resp.\ $w'_0$) the longest element. Then $w_0(D)=-D$,
$w'_0(E)=-E$, so $w'_0(E)\supset w_0(D)$. Let $\bB=\bT\prod_{a\in\Phi_y\cap\Phi_D^{+}}\bU_a$
be the minimal parabolic of $\bG$ attached to $E$. Since $w'(D)=\Delta'$ one has
$\bU_{y+\Delta'}\subset w'\bB w'^{\,-1}$, and (1), (2) give
\begin{align}
  \bu&\in w'_0\,\bB\,w'^{\,-1}_0, \tag{8}\\
  \bu&\in \bB\,t\,w'\,\bB\,w'^{\,-1}. \tag{9}
\end{align}

Take a reduced decomposition $(r_k,\dots,r_1)$ of $w'^{\,-1}w_0$ in $\vW$ w.r.t.\
$R$, so $w_0=w'r_k\cdots r_1$ and $\ell(w')=\ell(w_0)-k$ ([\btbib{5}], ch.~VI, \S1, no.~6,
cor.~3 to prop.~17). For $1\leqslant i\leqslant k$ put $w_i=w'r_k\cdots r_{i+1}$, so
$w_k=w'$ and $w_i=w_{i+1}r_i$; the chambers $w_i(D)$, $w_{i+1}(D)$ are adjacent and
the reflection $s_i$ in their common wall is
$w_{i+1}w_i^{-1}=w_ir_iw_i^{-1}$. Let $w'_i$ be the unique element of $\vW'$ with
$w'_i(E)\supset w_i(D)$; this is consistent with $w'_0$, and $w'_k=w_k=w'$ since
$w'\in\vW'$. The chambers $w'_i(E)$, $w'_{i+1}(E)$ are adjacent or equal according
as $s_i\in\vW'$ or not; in the first case $s_i=w'_{i+1}w'^{\,-1}_i$.

By (6.1.15\,(a)) there is for each $i=0,\dots,k$ a unique $\bw'_i\in\vW'$ with
$\bB\,\bu\,w'_i\,\bB=\bB\,\bw'_i\,\bB$, and $\bw'_0=w'_0$ by (8), $\bw'_k=tw'$ by
(9). Let $I=\{0,\dots,k-1\}$ and $I_1=\{i\in I: w'_{i+1}=w'_i\}$. Clearly
\begin{equation}
  i\in I_1\ \Longrightarrow\ w'_{i+1}=w'_i\ \text{and}\ \bw'_{i+1}=\bw'_i, \tag{10}
\end{equation}
and we saw that
\begin{equation}
  i\notin I_1\ \Longrightarrow\ w'_{i+1}w'^{\,-1}_i=w_{i+1}w_i^{-1}=s_i. \tag{11}
\end{equation}
For $i\notin I_1$, $w'^{\,-1}_{i+1}w'_i=w'^{\,-1}_is_iw'_i$ is a reflection in a
wall of $E$, i.e.\ an element $r'_i\in R'$; by (\btref{6.1.15}) and axiom (T\,3) of Tits
systems,
\[
  \bB\bw'_{i+1}\bB=\bB\bu w'_{i+1}\bB=\bB\bu w'_ir'_i\bB
  \subset\bB\bw'_i\bB r'_i\bB\subset\bB\bw'_i\bB\cup\bB\bw'_ir'_i\bB .
\]
Hence $I-I_1$ splits into $I_2=\{i:\bw'_{i+1}=\bw'_i\}$ and
$I_3=\{i:\bw'_{i+1}=\bw'_ir'_i\}$.

Finally put $\bw_i=\bw'_iw'^{\,-1}_iw_i$ $(i=0,\dots,k)$; then $\bw_0=w_0$ and
$\bw_k=tw'\,w'^{\,-1}w'=tw'$. For $i\in I$ set
$d_i=\bw^{-1}_{i+1}\bw_i=w^{-1}_{i+1}w'_{i+1}\bw'^{\,-1}_{i+1}\bw'_iw'^{\,-1}_iw_i$.
Then $d_i=r_i$ if $i\in I_1$ (by (10)); $d_i=w^{-1}_{i+1}s_iw_i=1$ if $i\in I_2$
(by (11)); $d_i=w^{-1}_{i+1}w_i=r_i$ if $i\in I_3$. Consequently
\[
  w_0=\bw_0=\bw_k(\bw_k^{-1}\bw_{k-1})\cdots(\bw_1^{-1}\bw_0)    =tw'\prod_{i\in I_1\cup I_3}r_i,
\]
whence
$\ell(tw')\geqslant\ell(w_0)-\Card(I_1\cup I_3)\geqslant\ell(w_0)-k=\ell(w')$,
with equality only if $I_2=\varnothing$, which forces $w_0=tw_0$ and $t=1$. This
proves (7), the lemma (\btref{7.3.6}), and therefore the theorem (\btref{7.3.4}). \qed

\begin{stmt}{Proposition}{7.3.9}\btanchor{7.3.9}
Let $F$ be a facet of $A$ and let $L$ be a subgroup of $G$ containing $H$ and such
that $L\cap U_a=P_F\cap U_a$ for every $a\in\Phi$. Let $N_F^{\dagger}$ be the
stabilizer of $F$ in $N$ and put $P_F^{\dagger}=N_F^{\dagger}P_F$. Then
$P_F^{\dagger}$ is a subgroup of $G$ and
$P_F\subset L\subset P_F^{\dagger}$.
\end{stmt}

\pf The first assertion holds because $N_F^{\dagger}$ normalizes $P_F$; and
$P_F\subset L$ since $P_F$ is generated by $H$ and the $U_a\cap P_F$. As $P_F$
contains some $B_{x,D}$, (\btref{7.3.4}) shows that if $L\not\subset P_F^{\dagger}$ there is
$n\in N\cap L$ with $n.F\neq F$. Then some affine root $\alpha$ contains exactly one
of $F$, $\nu(n).F$; replacing if need be $n$ by $n^{-1}$ and $\alpha$ by
$\nu(n^{-1}).\alpha$, assume $\alpha\supset F$ and $\alpha\not\supset\nu(n).F$. Then
$U_\alpha\subset L$ and $U_{\nu(n^{-1}).\alpha}=n^{-1}U_\alpha n\subset L$, which by
(7.2.6) and (\btref{6.4.9}) forces $\nu(n^{-1}).\alpha\supset F$ --- absurd. \qed

\btsecB{7.4}{7.4. The building of \texorpdfstring{$G$}{G}}

\num{7.4.1}\btanchor{7.4.1} Consider the following relation between two elements $(g,x)$ and $(h,y)$
of $G\times A$:
\[
  \text{(R)}\qquad
  \text{there exists }n\in N\text{ such that } y=\nu(n)(x)\text{ and }
  g^{-1}hn\in\hP_x .
\]
Since $\hP_x=P_x\hN_x$ and $\hN_x$ is the stabilizer of $x$ in $N$, (R) is
unchanged if $\hP_x$ is replaced by $P_x$.

\emph{(R) is an equivalence relation.} Let $g,h,k\in G$, $x,y,z\in A$, $m,n\in N$
with $y=\nu(m)(x)$, $z=\nu(n)(y)$, $g^{-1}hm\in\hP_x$, $h^{-1}kn\in\hP_y$. By
(7.1.8), $\hP_y=m\hP_xm^{-1}$; hence
\begin{align*}
  h^{-1}gm^{-1}&=m(m^{-1}h^{-1}g)m^{-1}\in\hP_y, & x&=\nu(m^{-1})(y),\\
  g^{-1}knm&=g^{-1}hm\cdot m^{-1}h^{-1}knm\in\hP_x\,m^{-1}\hP_ym=\hP_x, & z&=\nu(nm)(x).
\end{align*}
The first gives symmetry, the second transitivity; reflexivity is clear.

\begin{stmt}{Definition}{7.4.2}\btanchor{7.4.2}
The \emph{building} of $G$ (equipped with the root datum $(T,(U_a))$ and with the
valuation $\varphi$ of it) is the set ${}^{\varphi}\Imm=\Imm$ obtained as the
quotient of $G\times A$ by the equivalence relation \textup{(R)}.
\end{stmt}

The building ${}^{\varphi}\Imm$ clearly depends only on the equipollence class of
$\varphi$.

The map $j$ sending $x\in A$ to the class of $(1,x)$ is injective: if $j(x)=j(y)$
there is $n\in N$ with $y=\nu(n)(x)$ and $n\in\hP_x$; as $N\cap\hP_x=\hN_x$
(7.1.8), $y=x$. \emph{We identify $A$ with its image in $\Imm$ via $j$ from now on.}

Left translation of $G$ on itself is compatible with (R) and gives an action
$(g,x)\mapsto g.x$ of $G$ on $\Imm$ for which the class of $(g,x)$ equals $g.x$. For
$n\in N$ this extends the action of $n$ on $A$ given by $\nu$: $n.x=\nu(n)(x)$.
Clearly $\Imm=G.A$.

Finally, one checks easily that for $\varphi$ discrete this construction recovers,
up to isomorphism, the building of $G$ attached to the Tits system of $G$ (\btref{6.5}).

\begin{stmt}{Proposition}{7.4.3}\btanchor{7.4.3}
Let $\psi$ be a valuation equivalent to $\varphi$ \textup{(\btref{6.2.5})}. There is one and
only one map $i:{}^{\varphi}\Imm\to{}^{\psi}\Imm$ with the properties:

\quad a) $i(g.x)=g.i(x)$ for $x\in{}^{\varphi}\Imm$, $g\in G$;

\quad b) the restriction of $i$ to $A$ is an affine map of $A=\varphi+V$ onto
$\psi+V$.

\noindent If $\lambda:\Phi\to\mathbf{R}^{*}_{+}$ is defined by
$\psi=\lambda\varphi+V$, then $i(\varphi)=\lambda\varphi$.
\end{stmt}

\pf One has $\lambda(\varphi+V)=\psi+V$. Let $i_1:A\to\psi+V$ be the affine map
$x\mapsto\lambda x$. Clearly
\begin{equation}
  {}^{\varphi}P_x={}^{\psi}P_{i_1(x)}\qquad(x\in A) \tag{1}
\end{equation}
and, by (\btref{6.2.5})\,(1),
\begin{equation}
  i_1(n.x)=n.i_1(x)\qquad(x\in A,\ n\in N), \tag{2}
\end{equation}
whence
\begin{equation}
  {}^{\varphi}\hN_x={}^{\psi}\hN_{i_1(x)}\quad\text{and}\quad
  {}^{\varphi}\hP_x={}^{\psi}\hP_{i_1(x)}\qquad(x\in A). \tag{3}
\end{equation}
So $i_1$ extends uniquely to a map $i$ with (a) and (b).

Conversely let $i$ satisfy (a). By the definition of the building,
\begin{equation}
  {}^{\varphi}\hP_x={}^{\psi}\hP_{i(x)}\qquad\text{for all }x\in A. \tag{4}
\end{equation}
Let $x$ be a special point of $A$ and $y\in A$, $y\neq x$. There is an affine root
$\alpha$ containing $x$ but not $y$; by (7.1.8)\,(1) and (\btref{6.4.9}),
$U_\alpha\subset{}^{\varphi}\hP_x$ and
$U_\alpha\not\subset{}^{\varphi}\hP_y={}^{\psi}\hP_{\lambda y}$. By (4),
$i(x)\neq\lambda y$; hence $i(x)=\lambda x$ for every special point $x$. If $i$ also
satisfies (b), then $i|A=i_1$, whence uniqueness and the last assertion. \qed

\begin{stmt}{Proposition}{7.4.4}\btanchor{7.4.4}
Let $\Omega$ be a non-empty subset or a facet of $A$. The fixer of $\Omega$ in $G$
(for the action of $G$ on $\Imm$) is equal to $\hP_\Omega$.
\end{stmt}

(The fixer of a germ is by definition the union of the fixers of its elements.) This
follows from the definition of $\Imm$ when $\Omega$ is a point, and the general case
follows using (\btref{7.1.11}).

\begin{stmt}{Proposition}{7.4.5}\btanchor{7.4.5}
Let $a\in\Phi$ and $u\in U_a^{*}$; put $k=\varphi_a(u)$. The set of points of $A$
fixed by $u$ is the affine root
\[
  \alpha_{a,k}=\{x\in A\mid a(x)+k\geqslant0\}.
\]
\end{stmt}

Indeed, for every $x\in A$ one has $\hP_x\cap U_a=U_{a,-a(x)}$, and $u\in\hP_x$ is
equivalent to $k\geqslant-a(x)$. Note that there are points of $\Imm$ fixed by $u$
outside $A$ (cf.\ (\btref{7.4.33})).

\begin{stmt}{Corollary}{7.4.6}\btanchor{7.4.6}
For every non-empty closed subset $\Omega$ of $A$, the set of points of $A$ fixed by
all elements of $\hP_\Omega$ is $\Omega$.
\end{stmt}

This follows from (\btref{7.4.5}) when $\Omega$ is an affine root; the general case follows
since $\cl(\Omega)$ is the intersection of the affine roots containing $\Omega$.

\begin{stmt}{Definition}{7.4.7}\btanchor{7.4.7}
An \emph{apartment} of $\Imm$ is a transform of $A$ by an element of $G$.
\end{stmt}

\begin{stmt}{Proposition}{7.4.8}\btanchor{7.4.8}
Let $g\in G$. The intersection $A\cap g^{-1}.A$ is a closed subset of $A$, and there
exists $n\in N$ such that $g.x=n.x$ for all $x\in A\cap g^{-1}.A$.
\end{stmt}

\pf We may assume $\Omega=A\cap g^{-1}.A\neq\varnothing$. Let $\mathscr{X}$ be the
set of subsets $X$ of $\Omega$ with $g^{-1}N\cap\hP_X\neq\varnothing$. By the very
definition of $\Imm$, $\mathscr{X}$ contains all one-point subsets of $\Omega$. Let
$X\in\mathscr{X}$, $y\in\Omega$ and $n_X,n_y\in N$ with $g^{-1}n_X\in\hP_X$,
$g^{-1}n_y\in\hP_y$. By (\btref{7.1.6}) there is a vector chamber $D$ with
\[
  n_X^{-1}n_y\in\hP_X\hP_y\subset\hN_X\,U_D^{-}\,U_D^{+}\,\hN_y
\]
(using $\hP_X=\hN_XP_X$ and $\hP_y=P_y\hN_y$). Hence there are
$n'_X\in\hN_X$, $n'_y\in\hN_y$ with
$n'^{\,-1}_Xn_X^{-1}n_yn'_y\in N\cap U_D^{-}U_D^{+}=\{1\}$. Putting
$n=n_Xn'_X=n_yn'_y$ we get $g^{-1}n\in\hP_X\cap\hP_y=\hP_{X\cup\{y\}}$. By induction
on $k$, every finite subset $\{x_1,\dots,x_k\}$ of $\Omega$ lies in $\mathscr{X}$.

Now fix $x_0\in\Omega$ and write $\Omega$ as the union of an increasing filtering
family of finite subsets $X_i\ni x_0$. There are $n_0,n_i\in N$ with
$g^{-1}n_0\in\hP_{x_0}$ and $g^{-1}n_i\in\hP_{X_i}\subset\hP_{x_0}$. Then
$n_i^{-1}n_0\in\hN_{x_0}$, and since $\hN_{x_0}/H$ is finite we may (passing to a
cofinal subfamily and multiplying the $n_i$ on the right by elements of $H$) assume
$n_i$ independent of $i$. Then $g^{-1}n_i\in\bigcap_i\hP_{X_i}=\hP_\Omega$ (\btref{7.1.11}),
which gives the last assertion. Moreover there is $n'\in\hN_\Omega$ with
$g^{-1}n_in'\in P_\Omega=P_{\cl(\Omega)}$ (7.1.2), so $g.x\in A$ for all
$x\in\cl(\Omega)$, i.e.\ $\Omega$ is closed. \qed

\begin{stmt}{Corollary}{7.4.9}\btanchor{7.4.9}
Let $\Omega$ be a subset or a facet of $A$.

(i) The fixer $\hP_\Omega$ of $\Omega$ acts transitively on the set of apartments
containing $\Omega$.

(ii) Suppose $\Omega$ is a chamber and let $b\in P_\Omega$. Then $b.x=x$ for all
$x\in A\cap b.A$.
\end{stmt}

\pf Let $A'=g.A$ be an apartment containing $\Omega$ and $n\in N$ with
$g^{-1}.x=n.x$ for all $x\in A\cap g.A$. Then $gn\in\hP_\Omega$ and $A'=gn.A$,
whence (i). If $\Omega$ is a chamber and $g\in P_\Omega$, then by (7.2.6)
$n\in N\cap P_C=H$, whence (ii). \qed

\smallskip
Compare (\btref{7.4.9}) with (\btref{2.5.8}) and (\btref{2.5.9}).

\begin{stmt}{Corollary}{7.4.10}\btanchor{7.4.10}
The subgroup $N$ (resp.\ $H$) is the stabilizer (resp.\ the fixer) of $A$ in $G$.
\end{stmt}

\pf Let $g\in G$ with $g.A=A$. By (\btref{7.4.8}) there is $n\in N$ with $g^{-1}n\in\hP_A$.
But $U_A=\{1\}$, since $A$ is contained in no affine root, so
$\hP_A=\hN_A=H$; and $H$ is by definition the kernel of $\nu$. \qed

\smallskip
Compare (\btref{7.4.10}) with (\btref{2.2.5}).

\num{7.4.11}\btanchor{7.4.11} As in (\btref{2.2.8}), it follows from (\btref{7.4.10}) that on every apartment $A'$ of
$\Imm$ there is one and only one structure of Euclidean affine space such that, for
every $g\in G$ with $A'=g.A$, the map $x\mapsto g.x$ is an isomorphism of Euclidean
affine spaces from $A$ onto $A'$. In particular each apartment $A'$ carries a
distance $d_{A'}$.

Moreover (\btref{7.4.8}) shows at once that if $M$ is a subset of an apartment, the set
$g.\cl(g^{-1}.M)$ does not depend on the choice of $g\in G$ with $M\subset g.A$.
This set is called the \emph{enclosure} of $M$ and denoted $\cl(M)$; it is contained
in every apartment containing $M$. One says that $M$ is \emph{closed} if
$\cl(M)=M$.

\begin{stmt}{Definition}{7.4.12}\btanchor{7.4.12}
A \emph{sector} (resp.\ \emph{sector germ}, \emph{chamber}, \emph{facet}) of
$\Imm$ is a transform of a sector (resp.\ sector germ, chamber, facet) of $A$ by
an element of $G$.
\end{stmt}

Of course one identifies a germ of $A$, defined by a filter base $\mathscr{F}$ on
$A$, with the germ defined by $\mathscr{F}$ in $\Imm$. Definition (\btref{7.4.12}) is
justified by:

\begin{stmt}{Proposition}{7.4.13}\btanchor{7.4.13}
(i) If a facet $X$ of $\Imm$ meets an apartment $A'=g.A$ of $\Imm$, there is a facet
$Y$ of $A$ with $X=g.Y$, and $X$ is contained in $A'$.

(ii) If a sector (resp.\ sector germ) $X$ of $\Imm$ is contained in an apartment
$A'=g.A$, there is a sector (resp.\ sector germ) $Y$ of $A$ with $X=g.Y$.
\end{stmt}

\pf (i) Let $x\in A$, let $E$ be a vector facet of $\Phi_x$ and $h\in G$ with
$X=h.F_{x,E}$; suppose $X$ meets $g.A$. For every convex neighborhood $\Omega$ of
$x$ in $A$ one has $h.(\Omega\cap(x+E))\cap g.A\neq\varnothing$. For $\Omega$ small
enough there is an open half-line $E'$ with origin $o$ in $V$, $E'\subset E$, with
$h.(\Omega\cap(x+E'))\subset g.A$. Then $F_{x,E}=F_{x,E'}$, and by (\btref{7.4.8}) there is
$n\in N$ with $g^{-1}h.y=n.y$ for all $y\in A\cap g^{-1}h.A$. Hence
$n^{-1}g^{-1}h\in\hP_{x,E'}=\hP_{x,E}$ and $X=h.F_{x,E}=g.F_{x'',E''}$ with
$x''=n.x$, $E''=\vnu(n)(E)$. The proof of (ii) is analogous and simpler. \qed

\smallskip
In particular, the facets (resp.\ chambers) of $\Imm$ meeting $A$ are exactly the
facets (resp.\ chambers) of $A$.

\begin{stmt}{Corollary}{7.4.14}\btanchor{7.4.14}
Let $F,F'$ be facets of $A$ with $F'$ contained in the closure of $F$ in $A$
\textup{(7.2.1)}, and let $g\in G$. Every apartment of $\Imm$ containing $g.F$ also
contains $g.F'$.
\end{stmt}

This follows from (\btref{7.4.13}), since $A\cap h.A$ is a closed subset of $A$ for every
$h\in G$ (\btref{7.4.8}).

\begin{stmt}{Proposition}{7.4.15}\btanchor{7.4.15}
Let $\Omega,\Omega'$ be subsets or facets of $A$ and let $\hW_\Omega$ (resp.\
$\hW_{\Omega'}$) be the image of $\hN_\Omega$ (resp.\ $\hN_{\Omega'}$) in $\hW$. The
natural map
$\hW_\Omega\backslash\hW/\hW_{\Omega'}\to\hP_\Omega\backslash G/\hP_{\Omega'}$
is bijective.
\end{stmt}

The proof is identical with that of (\btref{4.2.1}), replacing the reference (\btref{2.5.8}) by
(\btref{7.4.9}).

\begin{stmt}{Proposition}{7.4.16}\btanchor{7.4.16}
Suppose the valuation $\varphi$ is dense. Let $x\in A$ and let $E$ be a vector facet.

(i) The only point of $A$ fixed by $P_{x,E}$ is $x$.

(ii) If $g\in G$ is such that $gP_{x,E}g^{-1}\subset\hP_x$, then $g\in\hP_x$.
\end{stmt}

\pf Let $y\in A$, $y\neq x$. As $\varphi$ is dense there is an affine root containing
$x$ but not $y$; by (\btref{7.4.5}) this gives (i). For (ii), by (\btref{7.3.4}) we may assume
$g\in N$; but for $n\in N$ one has $nP_{x,E}n^{-1}=P_{n.x,E'}$ with
$E'=\vnu(n)(E)$, so (ii) follows from (i). \qed

\begin{stmt}{Corollary}{7.4.17}\btanchor{7.4.17}
If $\varphi$ is dense, the normalizer of $P_x$ (resp.\ $\hP_x$) equals $\hP_x$, for
every $x\in A$.
\end{stmt}

\begin{stmt}{Theorem}{7.4.18}\btanchor{7.4.18}
(i) Two chambers (resp.\ facets, points) of $\Imm$ lie in a common apartment.

(ii) A chamber (resp.\ facet, point) of $\Imm$ and a sector germ lie in a common
apartment.

(iii) Two germs of sectors of $\Imm$ lie in a common apartment.
\end{stmt}

\pf These are geometric translations of (\btref{7.3.4}), (\btref{7.3.1}) and (\btref{6.1.15})\,a)
respectively. By (\btref{7.4.14}) it suffices to treat chambers and germs of sectors. Let
$C,C'$ be chambers of $\Imm$; transforming both by one element of $G$ we may assume
$C=C_{x,D}$ and $C'=g.C_{x',D'}$ with $x,x'\in A$, $D,D'$ vector chambers, $g\in G$.
By (\btref{7.3.4}), $g=bnb'$ with $b\in B_{x,D}$, $n\in N$, $b'\in B_{x',D'}$; then $C=b.C$
and $C'=bn.C_{x',D'}$ lie in the common apartment $b.A$. One proves (ii) similarly
from (\btref{7.3.1}) and (iii) from (\btref{6.1.15}); see also (\btref{4.2.5}). \qed

\begin{stmt}{Theorem}{7.4.19}\btanchor{7.4.19}
Let $C$ be a chamber of $\Imm$ and $A'$ an apartment containing $C$. There is one
and only one map $\rho$ from $\Imm$ into $A'$ such that, for every $b\in B_C$, the
restriction of $\rho$ to the apartment $b.A'$ coincides with the restriction of
$x\mapsto b^{-1}.x$. The restriction of $\rho$ to $A'$ is the identity, and if $x$
is a point of $A'$ adherent to $C$ then $\rho^{-1}(x)=\{x\}$.
\end{stmt}

For the proof see (\btref{2.3.2}) and (\btref{2.3.4}). As in (\btref{2.3.5}) this map is denoted
$\rho_{A';C}$ and called the \emph{retraction of $\Imm$ onto $A'$ with center $C$}.

\begin{stmt}{Proposition}{7.4.20}\btanchor{7.4.20}
(i) There is one and only one distance $d$ on $\Imm$ whose restriction to every
apartment $A'$ of $\Imm$ is the Euclidean distance of $A'$ \textup{(\btref{7.4.11})}. It is
$G$-invariant.

(ii) Let $C$ be a chamber of $\Imm$, $x$ a point adherent to $C$ and $A'$ an
apartment containing $C$; let $\rho$ be the retraction of $\Imm$ onto $A'$ with
center $C$. For all $y,z\in\Imm$,
\[
  d(\rho(y),\rho(z))\leqslant d(y,z).
\]
If moreover $C$, $y$ and $z$ lie in one apartment, then
$d(\rho(y),\rho(z))=d(y,z)$. In particular $d(x,\rho(y))=d(x,y)$ for all $y\in\Imm$.

(iii) Let $x,y\in\Imm$ and $S=\{z\in\Imm\mid d(x,y)=d(x,z)+d(z,y)\}$. Then $S$ is
contained in every apartment $A'$ containing $x$ and $y$, and coincides with the
segment $[xy]$ of the affine space $A'$; one writes $S=[xy]$. Every isometry of
$\Imm$ --- in particular every element of $G$ --- fixing $x$ and $y$ fixes every
point of $[xy]$.

(iv) Let $x,y,z,z'\in\Imm$ with $z\in[xy]$, and $\varepsilon\in\mathbf{R}_{+}$ with
$d(x,z')\leqslant d(x,z)+\varepsilon d(x,y)$ and
$d(y,z')\leqslant d(y,z)+\varepsilon d(x,y)$. Then
\[
  d(z,z')^{2}\leqslant4\,d(x,z)\,d(y,z)\,\varepsilon+d(x,y)^{2}\varepsilon^{2}.
\]

(v) Let $x,y\in\Imm$ and $t\in[0,1]$; write $tx+(1-t)y$ for the unique $z\in[xy]$
with $d(y,z)=t\,d(x,y)$. The map $(t,x,y)\mapsto tx+(1-t)y$ from
$[0,1]\times\Imm\times\Imm$ to $\Imm$ is continuous. The metric space $\Imm$ is
contractible.
\end{stmt}

For the proof see nos.\ (\btref{2.5.1})--(\btref{2.5.4}) and (\btref{2.5.14})--(\btref{2.5.16}). In the proof of
(\btref{2.5.3})\,(i) one replaces the use of facets by the following lemma.

\begin{stmt}{Lemma}{7.4.21}\btanchor{7.4.21}
Let $C$ be a chamber, $A'$ an apartment, and $x,y$ two points of $A'$. There is a
monotone sequence $(x_0=x,x_1,\dots,x_m=y)$ of points of the segment $[xy]$ of $A'$
such that for each $i=0,\dots,m-1$ the segment $[x_ix_{i+1}]$ of $A'$ and the
chamber $C$ lie in one apartment.
\end{stmt}

\pf We may assume $x\neq y$. For each $z\in\,]xy[$ there are chambers $C_z^{+}$,
$C_z^{-}$ contained in $A'$ with $]xz[\,\cap\,C_z^{-}\neq\varnothing$ and
$]zy[\,\cap\,C_z^{+}\neq\varnothing$. Let $A_z^{+}$ (resp.\ $A_z^{-}$) be an
apartment containing $C$ and $C_z^{+}$ (resp.\ $C_z^{-}$). There are then
$z^{+}\in\,]zy[$ and $z^{-}\in\,]xz[$ with $[zz^{+}]\subset A_z^{+}$ and
$[z^{-}z]\subset A_z^{-}$. One defines $x^{+},A_x^{+}$ and $y^{-},A_y^{-}$
analogously. There are finitely many points $z_1,\dots,z_k$ of $[xy]$ such that the
$]z_i^{-}z_i^{+}[$ together with $[xx^{+}[$ and $]y^{-}y]$ form an open cover of
$[xy]$. Ordering $x,x^{+},z_i^{-},z_i,z_i^{+},y^{-},y$ monotonically gives the
required sequence. \qed

\begin{stmt}{Remark}{7.4.22}\btanchor{7.4.22}
Let $C,C'$ be chambers of $A$ and $g=bnb'\in G$ with $b\in B_C$, $n\in N$,
$b'\in B_{C'}$. Since the restriction of $\rho_{A;C}$ to the apartment
$b.A\supset g.C'$ coincides with the action of $b^{-1}$, one has
$\rho_{A;C}(g.C')=n.C'$. Likewise $\rho_{A;C}(g.x)=n.x$ as soon as
$g\in B_Cn\hP_x$ \textup{($x\in A$, $n\in N$)}.
\end{stmt}

\num{7.4.23}\btanchor{7.4.23} Let $v$ be an element of length $1$ of $V$, $E=E(v)$ the vector facet
containing it and $\tE=\tE(v)$ the vector subspace generated by $E$. Put
\[
  c(v)=\sin(\alpha/2)
\]
where $\alpha$ is the smallest of the angles that $v$ makes with its transforms
distinct from itself under the various elements of $\vW$.

If $\tE=V$, put $k(v)=+\infty$; if $\tE\neq V$, put
\[
  k(v)=\sin\beta
\]
where $\beta$ is the smallest of the angles that $v$ makes with the vector subspaces
generated by the vector facets not contained in $\tE$ and of dimension
$\leqslant\dim\tE$. Finally let $\delta(v)$ be the distance from $v$ to the
complement of $E$ in $\tE$, and put
\[
  h(v)=\inf\{c(v),k(v),\delta(v)\}.
\]
Note that $h(v)>0$.

\begin{stmt}{Proposition}{7.4.24}\btanchor{7.4.24}
Keep the notation of \textup{(\btref{7.4.23})}. Let $x,y_1,z\in A$ with $x\in[zy_1]$ and
$y_1-z\in\mathbf{R}^{*}_{+}v$. Let $g\in\hP_z$ and $y\in\Imm$ with $g^{-1}.y\in A$
and
\[
  d(y,y_1)+d(g^{-1}.y,y_1)\leqslant2c(v)\,d(x,y_1).
\]
Then $g\in\hP_x$.
\end{stmt}

\pf The set of $t\in[zx]$ with $g\in\hP_t$ is a segment $[zm]$. Suppose $m\neq x$.
Let $D$ be a vector chamber with $v\in\bD$ and $\rho$ the retraction onto $A$ with
center the chamber $C_{m,D}$. Let $n\in N$ with $g\in B_{m,D}nB_{m,D}$. Since
$g.m=m$ we have $n.m=m$; and by (\btref{7.4.22}) there is $m'\in\,]my_1]$ with
$\rho(g.p)=n.p$ for $p\in[mm']$. As $\rho(g.[my_1])$ is the segment
$[m\,\rho(g.y_1)]$, we get $[m\,\rho(g.y_1)]=n.[my_1]$ and
\[
  \rho(g.y_1)-y_1=d(m,y_1)\bigl(\vnu(n)(v)-v\bigr).
\]
Now $\vnu(n)(v)\neq v$: otherwise, since an element of $B_{m,D}$ fixes the points of
a neighborhood of $m$ in $m+\bD$, we would have $g.(m+tv)=m+tv$ for all small
$t\geqslant0$, contradicting the definition of $m$. Hence the length of
$\vnu(n)(v)-v$ is at least $2c(v)$ and
\begin{align*}
  d(\rho(g.y_1),y_1)&\geqslant2c(v)d(m,y_1)>2c(v)d(x,y_1)\geqslant d(y,y_1)+d(g^{-1}.y,y_1),\\
  d(\rho(g.y_1),y_1)&>d(y,y_1)+d(y,g.y_1)\geqslant d(\rho(y),y_1)+d(\rho(y),\rho(g.y_1))\\
  &\geqslant d(y_1,\rho(g.y_1)),
\end{align*}
a contradiction. Hence $m=x$ and $g\in\hP_x$. \qed

\begin{stmt}{Proposition}{7.4.25}\btanchor{7.4.25}
Let $A'$ be an apartment of $\Imm$ and $\gamma$ a sector germ contained in $A'$.
There is one and only one map, denoted $\rho_{A';\gamma}$, from $\Imm$ into $A'$
with the following property: for every bounded subset $M$ of $\Imm$ there is a
sector $\mathfrak{C}$ of $A'$ belonging to $\gamma$ such that, for every chamber $C$
contained in $\mathfrak{C}$, one has $\rho_{A';\gamma}(y)=\rho_{A';C}(y)$ for all
$y\in M$.
\end{stmt}

\pf We may take $A'=A$ and $\gamma=\gamma(D)$, $D$ a vector chamber; uniqueness is
clear. Let $M\subset\Imm$ be bounded, $y_1\in A$, $v\in-D$. Put
$r=\sup\{d(y_1,y)\mid y\in M\}$ and choose $x\in y_1-\mathbf{R}^{*}_{+}v$ so that
the ball of $A$ with center $y_1$ and radius $r$ lies in the sector $x-D$. Let
$y\in M$ and $g\in G$ with $y'=g.y\in A$. Write $g=hnu$ with $h\in\hP_{y'}$,
$n\in N$, $u\in U_D^{+}$ (\btref{7.3.1}); then $u.y\in A$. There is
$z\in x-\mathbf{R}^{*}_{+}v$ with $u\in U_{z+D}$, and by (\btref{7.4.19})
$u.y=\rho_{A;C}(y)$ for every chamber $C\subset z+D$, whence
\[
  d(y_1,u.y)\leqslant d(y_1,y)\leqslant r\leqslant c(v)\,d(y_1,x).
\]
By (\btref{7.4.24}), $u\in\hP_x$, so $u\in U_{x+D}$ and
\begin{equation}
  u.y=\rho_{A;C}(y) \tag{1}
\end{equation}
for every chamber $C\subset x+D$ and every $y\in M$. \qed

\smallskip
As in 2.9, $\rho_{A';\gamma}$ is called the \emph{retraction of $\Imm$ onto $A'$
relative to the sector germ $\gamma$}. It is the identity on $A'$ and decreases
distances.

\begin{stmt}{Proposition}{7.4.26}\btanchor{7.4.26}
Keep the notation of \textup{(\btref{7.4.23})}. Let $z\in A$ and $g\in\hP_z$; put
$\tau=\sup\{t\in\mathbf{R}\mid g.(z+tv)=z+tv\}$. For every $t\in\mathbf{R}$ with
$t\geqslant\tau$,
\[
  2c(v)(t-\tau)\leqslant d\bigl(z+tv,\,g.(z+tv)\bigr)\leqslant2(t-\tau).
\]
\end{stmt}

\pf Let $t\geqslant\tau$. Put $y=y_1=z+tv$, $d=d(y,g.y)$ and
$x=y-\bigl(\tfrac12c(v)^{-1}d\bigr)v$. By (\btref{7.4.24}), $g\in\hP_x$, whence
$\tau\geqslant t-\tfrac12c(v)^{-1}d$: this is the first inequality. For the second,
put $u=z+\tau v$; then
$d(z+tv,g.(z+tv))\leqslant d(u,z+tv)+d(u,g.(z+tv))=2(t-\tau)$. \qed

\smallskip
Note that when $\Phi$ has rank one, $c(v)=1$.

\begin{stmt}{Corollary}{7.4.27}\btanchor{7.4.27}
Let $D$ be a vector chamber with $v\in\bD$; let $y\in A$, $u\in U_D^{+}$ and
$\lambda\in\mathbf{R}_{+}$.

(i) If $d(y,u.y)\leqslant\lambda$, then $u\in U_{y+\frac12c(v)^{-1}\lambda v+D}$.

(ii) If $u\in U_{y+\lambda v+D}$, then $d(y,u.y)\leqslant\sup(0,2\lambda)$.
\end{stmt}

\pf Let $s=\inf\{t\in\mathbf{R}\mid u\in U_{y+tv+D}\}$; since $u\in U_D^{+}$,
$s<\infty$. If $s\leqslant0$ then $u.y=y$ and there is nothing to prove. If $s>0$,
apply (\btref{7.4.26}) with $-v$ in place of $v$, $g=u$, $z=y+sv$, $t=s$: then $\tau=0$ and
$2c(v)s\leqslant d(y,u.y)\leqslant2s$, which gives (i) and (ii). \qed

\begin{stmt}{Remark}{7.4.28}\btanchor{7.4.28}
Let $x\in A$, $u\in U_D^{+}$ and $n\in N$ with $u\in\hP_xn\hP_x$. Then
$d(x,u.x)=d(x,n.x)$, and \textup{(\btref{7.4.27})} gives a two-sided inequality between the
``size'' of $n$, measured by $d(x,n.x)$, and that of $u$, measured by the smallest
$\lambda\in\mathbf{R}_{+}$ with $u\in U_{x+\lambda v+D}$. Note that for $v\in D$ the
subgroups $U_{x+\lambda v+D}$, $\lambda\in\mathbf{R}$, form a decreasing filtration
of $U_D^{+}$ and allow one to define a ``valuation'' $\varphi_v$ of $U_D^{+}$.
\end{stmt}

\begin{stmt}{Proposition}{7.4.29}\btanchor{7.4.29}
Keep the notation of \textup{(\btref{7.4.23})}. Let $x,x',y_1,y'_1\in A$ be such that
$y_1-x$, $x-x'$ and $x'-y'_1$ lie in $\mathbf{R}^{*}_{+}v$. Put
$L=x+\tE=x'+\tE$ and let $y\in\hP_x.L$, $y'\in\hP_{x'}.L$ satisfy
\[
  d(y,y_1)<h(v)\,d(x,y_1)\qquad\text{and}\qquad d(y',y'_1)<h(v)\,d(x',y'_1).
\]
Then $(x-E)\cap(x'+E)\subset\cl(\{y,y'\})$.
\end{stmt}

\pf Let $C$ be a chamber of $A$ to which $y'_1$ is adherent; put $\rho=\rho_{A;C}$
and choose $g\in\hP_C$ with $\rho(y)=g.y$. Then
\[
  d(y,y_1)+d(g.y,y_1)\leqslant2d(y,y_1)<2c(v)\,d(x,y_1),
\]
so $g\in\hP_x$ by (\btref{7.4.24}). Since
$\cl(\{x\}\cup C)\supset(y'_1+E)\cap(x-E)\supset(x'+E)\cap(x-E)$, we get
$g\in\hP_C\cap\hP_{(x'+E)\cap(x-E)}$.

\emph{$\rho(y)\in x+E$.} There are $z\in L$ and $h\in\hP_x$ with $\rho(y)=gh.z$; by
(\btref{7.4.8}) there is $n\in\hN_x$ with $\rho(y)=n.z$, whence
$\rho(y)-x=\vnu(n)(z-x)$. As $z-x\in\tE$, $\rho(y)-x$ lies in a vector facet of
dimension $\leqslant\dim\tE$. If $\rho(y)-x\notin\tE$ then
\[
  d(y_1,y)\geqslant d(y_1,\rho(y))\geqslant k(v)d(y_1,x)\geqslant h(v)d(y_1,x),
\]
contrary to hypothesis. Moreover
$d(y_1,\rho(y))\leqslant d(y_1,y)<\delta(v)d(y_1,x)$, which forces $\rho(y)=g.y$ to
lie in $x+E$.

Now take $t\in\mathbf{R}_{+}$ large enough that $g.y\in x+tv-E$; put $z=x+tv$ and
$y''=g.y'$. Arguing as above but with $x',z,y'_1,y'',-v,-E$ in place of
$x,y'_1,y_1,y,v,E$, and using $d(y'_1,y'')=d(y'_1,y')$ and
$y''\in\hP_{x'}.y'\subset\hP_{x'}.L$, one finds
\[
  g'\in\hP_z\cap\hP_{(x'+E)\cap(z-E)}\subset\hP_{g.y}\cap\hP_{(x'+E)\cap(x-E)}
\]
with $g'g.y=g.y\in x+E$ and $g'g.y'=g'.y''\in x'-E$. Then
\[
  \cl(\{y,y'\})=g^{-1}g'^{\,-1}.\cl(\{g'g.y,g'g.y'\})
  \supset g^{-1}g'^{\,-1}.\bigl((x'+E)\cap(x-E)\bigr)=(x'+E)\cap(x-E). \qed
\]

\begin{stmt}{Corollary}{7.4.30}\btanchor{7.4.30}
Keep the notation of \textup{(\btref{7.4.23})} and let $\lambda$ be a real number with
$0\leqslant\lambda<1$. There is a real number $\varepsilon>0$ such that, for
$x,x'\in A$ and $y,y'\in\Imm$, the relations
$x'-x\in\mathbf{R}^{*}_{+}v$, $d(x,y)<\varepsilon d(x,x')$,
$d(x',y')<\varepsilon d(x,x')$ and $\{y,y'\}\subset\hP_u.(x+\tE)$ for every
$u\in[xx']$ imply the existence of $z,z'\in[yy']\cap A$ with
$d(z,z')\geqslant\lambda d(x,x')$.
\end{stmt}

The derivation from (\btref{7.4.29}) is very close to that of (\btref{2.8.10}) from (\btref{2.8.9}); we
leave it to the reader.

\begin{stmt}{Remark}{7.4.31}\btanchor{7.4.31}
If $E$ is a vector chamber, then $L=x+\tE=A$ and the conditions $y\in\hP_x.L$,
$y'\in\hP_{x'}.L$ (resp.\ $y,y'\in\hP_u.(x+\tE)$) of \textup{(\btref{7.4.29})} and
\textup{(\btref{7.4.30})} hold automatically; indeed $\Imm=\hP_x.A$ for every $x\in A$ by
\textup{(\btref{7.4.9})\,(i)} and \textup{(\btref{7.4.18})\,(i)}.
\end{stmt}

\begin{stmt}{Corollary}{7.4.32}\btanchor{7.4.32}
The apartment $A$ is the smallest convex subset of $\Imm$ stable under $T$.
\end{stmt}

One deduces (\btref{7.4.32}) from (\btref{7.4.30}) exactly as (\btref{2.8.11}) from (\btref{2.8.10}).

\begin{stmt}{Proposition}{7.4.33}\btanchor{7.4.33}
Let $v$ be an element of length $1$ of $V$ contained in the vector chamber $D$.
There is a constant $h>0$ with the following property: if $x\in A$, $u\in U_{x+D}$
and $\lambda\in\mathbf{R}_{+}$, then $u$ fixes every point of $\Imm$ lying in the
ball with center $x+\lambda v$ and radius $h\lambda$. If $\Phi$ has rank one one may
take $h=1/2$; if $\Phi$ is of type $A_1$ one may take $h=1$.
\end{stmt}

\pf Let $x,u,\lambda$ be as stated; put $y=x+\lambda v$ and $c=c(v)^{-1}$. Let
$z\in\Imm$; put $\delta=\lambda^{-1}d(y,z)$ and $t=\rho_{A;\gamma(D)}(z)$. By
(\btref{7.4.25})\,(1) there is $u'\in U_D^{+}$ with $z=u'.t$. Since $d(y,t)\leqslant d(y,z)$,
\[
  d(y,u'.y)\leqslant d(y,z)+d(u'.t,u'.y)=d(y,z)+d(t,y)\leqslant2\lambda\delta,
\]
so $u'\in U_{y+c\lambda\delta v+D}$ by (\btref{7.4.27})\,(i). Put $x'=y+c\lambda\delta v$. By
(\btref{6.4.32}) there is a constant $k$ depending only on $v$, with $0<k\leqslant1$, such
that
\[
  (U_{x+D},U_{x'+D})\subset U_{x'+k(x-x')+D}.
\]
Put $r=\delta(v)$ (\btref{7.4.23}); one checks easily that $t\in x'+k(x-x')+D$ as soon as
\[
  c\delta-k(1+c\delta)\leqslant-r^{-1}\delta .
\]
Since $k\leqslant1$ it suffices for this that $\delta\leqslant kr$. Under that
condition one also has $t\in x+D$ and
\[
  u.z=uu'.t=uu'u^{-1}.t=u'(u'^{\,-1},u).t=u'.t=z .
\]
So one may take $h=kr$. If $\Phi$ has rank one then $r=1$ and one may take $k=1/2$
(\btref{6.4.32}); if $\Phi$ is of type $A_1$ one may take $k=1$. \qed

\begin{stmt}{Corollary}{7.4.34}\btanchor{7.4.34}
Suppose $\varphi$ dense. There exist two constants $\sigma,\tau>0$ such that for any
points $x,y\in\Imm$ there is $g\in G$ satisfying:
\begin{align}
  &g\text{ fixes every point of }\Imm\text{ in the ball with center }x
   \text{ and radius }\sigma d(x,y); \tag{1}\\
  &d(y,g.y)\geqslant\tau d(x,y). \tag{2}
\end{align}
\end{stmt}

\pf For $t\in V$, $t\neq0$, let $\alpha(t)$ be the smallest of the angles formed by
$t$ and the $a^{\vee}$, $a\in\Phi$, and put
$\alpha=\inf_{t\in V,\,t\neq0}\cos\alpha(t)>0$. Let $v$ be an element of length $1$
of $V$ contained in a vector chamber $D$ and $h$ the corresponding constant of
(\btref{7.4.33}). We show that $\sigma=\tfrac13h\alpha$ and $\tau=\tfrac23c(v)\alpha$ work.

It suffices to treat $x,y\in A$ with $x-y\in\bD$. Choose $a\in\Phi$ with
$\cos\angle(x-y,a^{\vee})\geqslant\alpha$; then $a\in\Phi_D^{+}$. Put $d=d(x,y)$ and
choose $z\in[xy]$ with $-a(z)\in\Gamma_a$, $d(x,z)\geqslant\tfrac13d$ and
$d(y,z)\geqslant\tfrac13d$ (possible since $\varphi$ is dense). Let $u\in U_a$ with
$\varphi_a(u)=-a(z)$. By (\btref{7.4.5}) the fixed-point set of $u$ in $A$ is the half-space
$\{t\in A\mid a(t)-a(z)\geqslant0\}$; in particular $u\in U_{x-\lambda v+D}$ with
$\lambda=a(x-z)a(v)^{-1}\geqslant\alpha d(x,z)\geqslant\tfrac{\alpha}{3}d$. By
(\btref{7.4.33}), $u.t=t$ for every $t\in\Imm$ with
$d(x,t)\leqslant h\tfrac{\alpha}{3}d=\sigma d(x,y)$: this is (1).

Similarly $u\notin U_{y+\mu v+D}$ as soon as $\mu<a(z-y)a(v)^{-1}$, so by
(\btref{7.4.27})\,(i)
\[
  d(y,u.y)\geqslant2c(v)a(z-y)a(v)^{-1}\geqslant2c(v)\alpha d(y,z)\geqslant\tau d(x,y),
\]
which is (2). \qed

\btsecB{7.5}{7.5. The completion of the building}

\num{7.5.1}\btanchor{7.5.1} When $\varphi$ is discrete, the building of $G$ is a complete metric
space (\btref{2.5.12}). This is not so in general. We then write $\hImm$ for the completion
of $\Imm$. The operations of $G$, as well as the retractions onto an apartment,
extend by continuity to $\hImm$. From (\btref{7.4.20})\,(iv) one deduces easily that if
$x_i$ (resp.\ $y_i$) in $\Imm$ tend to $x$ (resp.\ $y$) in $\hImm$, then for each
$t\in[0,1]$ the points $tx_i+(1-t)y_i$ tend to a point $z\in\hImm$ depending only on
$t,x,y$; one then puts $z=tx+(1-t)y$ and $[xy]=\{tx+(1-t)y\mid t\in[0,1]\}$. By
continuity, (\btref{7.4.20})\,(iv) and (v) remain valid with $\Imm$ replaced by $\hImm$. It
follows at once that $[xy]$ is the unique geodesic of $\hImm$ joining $x$ to $y$;
more precisely,
\[
  [xy]=\{z\in\hImm\mid d(x,z)+d(z,y)=d(x,y)\}.
\]

\begin{stmt}{Proposition}{7.5.2}\btanchor{7.5.2}
Let $v$ be an element of length $1$ of $V$ contained in a vector chamber $D$. For
every point $x\in\hImm$ there is a triple $\bigl(y,(\lambda_p),(u_p)\bigr)$
satisfying:

\quad\textup{(1)} $y\in A$, $(\lambda_p)_{p\geqslant1}$ is a decreasing sequence of
real numbers tending to $0$, and $(u_p)_{p\geqslant1}$ is a sequence of elements of
$U_D^{+}$;

\quad\textup{(2)} $u_{p+1}\in u_pU_{y+\lambda_pv+D}$ for every $p\geqslant1$;

\quad\textup{(3)} $x=\lim_{p\to\infty}u_p.y$.

\noindent Conversely, for every triple $\bigl(y,(\lambda_p),(u_p)\bigr)$ satisfying
\textup{(1)} and \textup{(2)}, the sequence $u_p.y$ converges to an element $x$ of
$\hImm$. In order that $x\in\Imm$ it is necessary and sufficient that there exist
$u\in U_D^{+}$ with
\begin{equation}
  u_p\in uU_{y+\lambda_pv+D}\quad\text{for every }p\geqslant1. \tag{4}
\end{equation}
One then has $x=u.y$.
\end{stmt}

\pf Let $x\in\hImm$ and let $\rho$ be the retraction onto $A$ relative to the germ
of sector $\gamma(D)$, extended by continuity to $\hImm$. Choose points
$x_p\in\Imm$ with $\varepsilon_p=d(x,x_p)$ decreasing to $0$. Put $y=\rho(x)$,
$y_p=\rho(x_p)$. For each $p$ there is $u_p\in U_D^{+}$ with $x_p=u_p.y_p$
(\btref{7.4.25})\,(1). Since $d(u_p.y,u_p.y_p)=d(y,y_p)\leqslant d(x,x_p)$, after replacing
the $x_p$ by the $u_p.y$ and passing to a subsequence we may assume $x_p=u_p.y$ for
all $p$.

For $q\geqslant p$ we have $d(u_q^{-1}u_p.y,y)=d(u_p.y,u_q.y)\leqslant2\varepsilon_p$,
so by (\btref{7.4.27})\,(i)
\begin{equation}
  u_q^{-1}u_p\in U_{y+\lambda_pv+D}\qquad(q\geqslant p) \tag{2\,bis}
\end{equation}
with $\lambda_p=c(v)^{-1}\varepsilon_p$. Clearly $\bigl(y,(\lambda_p),(u_p)\bigr)$
satisfies (1), (2) and (3).

Conversely let $\bigl(y,(\lambda_p),(u_p)\bigr)$ satisfy (1) and (2). Since
$(\lambda_p)$ decreases, (2\,bis) holds; then (\btref{7.4.27})\,(ii) gives
$d(u_p.y,u_q.y)=d(u_q^{-1}u_p.y,y)\leqslant2\lambda_p$ for $q\geqslant p$, so
$(u_p.y)$ converges to some $x\in\hImm$.

If $x\in\Imm$, write $x=u.y$ with $u\in U_D^{+}$. Then
$d(u.y,u_p.y)=\lim_{q\to\infty}d(u_q.y,u_p.y)\leqslant2\lambda_p$, so by
(\btref{7.4.27})\,(i) $u\in u_pU_{x+c(v)^{-1}\lambda_pv+D}$ for every $p$. But for
$q\geqslant p$ large enough $\lambda_p\geqslant c(v)^{-1}\lambda_q$, whence
$u\in u_qU_{x+\lambda_pv+D}$, which with (2\,bis) gives (4). Conversely, if (4)
holds then (\btref{7.4.27})\,(ii) gives $d(u.y,u_p.y)\leqslant2\lambda_p$, so
$x=u.y\in\Imm$. \qed

\num{7.5.3}\btanchor{7.5.3} For $a\in\Phi$, call \emph{ball with center $u\in U_a$ and radius
$k\in\mathbf{R}$} in $U_a$ the left coset $uU_{a,k}$.

\begin{stmt}{Proposition}{7.5.4}\btanchor{7.5.4}
In order that the building $\Imm$ be complete it is necessary and sufficient that
the following condition hold for every $a\in\Phi^{\red}$:

\smallskip
\noindent\textup{(MC)} \emph{Every decreasing (for inclusion) sequence of balls of
bounded radius of $U_a$ has non-empty intersection.}
\end{stmt}

\pf \emph{Necessity.} Let $a\in\Phi^{\red}$ and, for integers $p\geqslant1$, let
$u_p\in U_a$ and $k_p\in\mathbf{R}$ with $k_q\geqslant k_p$ and
$u_qU_{a,k_q}\subset u_pU_{a,k_p}$ for $q\geqslant p$, and $k=\sup_pk_p<\infty$. Let
$D$ be a vector chamber with $a\in\Phi_D^{+}$ and $v$ an element of length $1$ of
$D$. Let $y\in A$ with $a(y)=-k$ and put $\lambda_p=a(v)^{-1}(k-k_p)$, so that
$U_{a,k_p}=U_a\cap U_{y+\lambda_pv+D}$. By (\btref{7.5.2}) the sequence $(u_p.y)$ converges
to a point $x\in\hImm$. If $\Imm$ is complete then $x\in\Imm$ and, by (\btref{7.5.2}), the
intersection of the $u_pU_{y+\lambda_pv+D}$ is non-empty. Let $u$ be in it. By
(\btref{6.1.6}) there are $u_1\in U_a$ and
$u_2\in\prod_{b\in\Phi_D^{+\red},\,b\neq a}U_b$ with $u=u_1u_2$. Since
$u_p^{-1}u_1u_2\in U_{y+\lambda_pv+D}$, we get
$u_p^{-1}u_1\in U_a\cap U_{y+\lambda_pv+D}=U_{a,k_p}$ (\btref{6.4.9}), whence
$\bigcap_pu_pU_{a,k_p}\neq\varnothing$.

\emph{Sufficiency.} Suppose (MC) holds for every $a\in\Phi^{\red}$. By (\btref{7.5.2}) it
suffices to show that if $y\in A$, $u_p\in U_D^{+}$, $\lambda_p\in\mathbf{R}_{+}$
and $v\in D$ are such that
$u_pU_{y+\lambda_pv+D}\supset u_qU_{y+\lambda_qv+D}$ for $q\geqslant p$, then
$X=\bigcap_pu_pU_{y+\lambda_pv+D}\neq\varnothing$.

Arrange the roots of $\Phi_D^{+\red}$ in a nibbling order $(a_1,\dots,a_m)$
(\btref{1.3.15}). For $i=1,\dots,m$ put $U_i=U_{a_i}$ and let $U'_i$ be the group generated
by the $U_j$ for $i<j\leqslant m$. By (DR\,2) and (\btref{6.1.6}), $U'_i$ is normalized by
$U_i$ and $U'_{i-1}$ is the semi-direct product of $U_i$ by $U'_i$. Put moreover
\[
  U^{p}=U_{y+\lambda_pv+D},\qquad U_i^{p}=U^{p}\cap U_i,\qquad   U'^{p}_i=U^{p}\cap U'_i .
\]
By (\btref{6.4.9}),
\[
  U'^{p}_{i-1}=U_i^{p}\,U'^{p}_i\qquad\text{and}\qquad   U_i^{p}=U_{a_i,\,-a_i(y)-a_i(v)\lambda_p}\supset U_{a_i,\,-a_i(y)} .
\]
We show by induction on $i$ that there exist $t_i\in U_D^{+}$ and
$u_{p,i}\in U'_i$ (for $i=0,\dots,m$ and $p\geqslant1$) with
\begin{align}
  &t_iu_pU^{p}\supset u_{p,i}U'^{p}_i\quad\text{for all }p; \tag{1\,;\,$i$}\\
  &u_{q,i}U'^{q}_i\subset u_{p,i}U'^{p}_i\quad\text{for }q\geqslant p. \tag{2\,;\,$i$}
\end{align}
This holds for $i=0$ (with $U'_0=U_D^{+}$, $U'^{p}_0=U^{p}$): take $t_0=1$,
$u_{p,0}=u_p$. Let $i\geqslant1$ and suppose $t_{i-1}$ and the $u_{p,i-1}$ already
constructed, satisfying (1\,;\,$i-1$) and (2\,;\,$i-1$). Let $f$ be the homomorphism
of $U'_{i-1}$ onto $U_i$ defined by $x\in f(x)U'_i$ for $x\in U'_{i-1}$. Then
$f(u_{p,i-1}U'^{p}_{i-1})=f(u_{p,i-1})U_i^{p}$, and (2\,;\,$i-1$) together with (MC)
and the relation $U_i^{p}\supset U_{a_i,-a_i(y)}$ shows that there exists
$s\in\bigcap_pf(u_{p,i-1})U_i^{p}$. Put $t_i=s^{-1}t_{i-1}$. Since
$U_i^{p}U'_i=U'_iU_i^{p}$, there is exactly one $u_{p,i}\in U'_i$ with
$s^{-1}u_{p,i-1}\in u_{p,i}U_i^{p}$. For every $p$,
\[
  t_iu_pU^{p}\supset s^{-1}u_{p,i-1}U'^{p}_{i-1}=u_{p,i}U'^{p}_{i-1}
  \supset\bigl(u_{p,i}U'^{p}_{i-1}\bigr)\cap U'_i=u_{p,i}U'^{p}_i
\]
and, for $q\geqslant p$,
\[
  u_{q,i}U'^{q}_i=\bigl(s^{-1}u_{q,i-1}U'^{q}_{i-1}\bigr)\cap U'_i
  \subset\bigl(s^{-1}u_{p,i-1}U'^{p}_{i-1}\bigr)\cap U'_i=u_{p,i}U'^{p}_i,
\]
whence (1\,;\,$i$) and (2\,;\,$i$).

Since $U'_m=\{1\}$ we have $u_{p,m}=1$ for all $p$, and (1\,;\,$m$) gives
$t_m^{-1}\in u_pU^{p}$ for all $p$, so $X\neq\varnothing$. \qed

\begin{stmt}{Example}{7.5.5}\btanchor{7.5.5}
Take up the notation of \textup{(\btref{6.2.3})\,a)} (canonical valued root datum of
$SL_2(K)$, $K$ a field, commutative or not, endowed with a non-improper valuation
$\omega$) or of \textup{(\btref{6.2.3})\,b)} (valued root datum of the group of rational
points over $K$ of a semi-simple algebraic group split over $K$, $K$ a commutative
field with a non-improper valuation $\omega$). Each group $U_a$ with the valuation
$\varphi_a$ is then isomorphic to the additive group of $K$ with $\omega$. Condition
\textup{(MC)} therefore says that the completion of $K$ for $\omega$ is maximally
complete \textup{[34]}. Note that this always holds when $\omega$ is discrete, which
matches the fact that the building of $G$ is always complete when $\varphi$ is
discrete.
\end{stmt}

\begin{stmt}{Remark}{7.5.6}\btanchor{7.5.6}
Condition \textup{(MC)} may hold for every $a\in\Phi^{\red}$ without holding for
every $a\in\Phi$. Take up for instance the notation of example
\textup{(\btref{6.2.3})\,d)}, taking for $K$ the field of formal series with well-ordered
exponents in a subgroup $\Gamma$ of $\mathbf{R}$ \textup{([\btbib{6}], ch.~VI, \S3,
exerc.~2)}, dense in $\mathbf{R}$ and with $\Gamma/2\Gamma$ infinite, with
coefficients in a field $k$ of characteristic $2$, and taking for $L$ the additive
subgroup generated by the elements $u^{2}X^{\gamma}$ with $u\in K$ and
$\gamma\in\Gamma$. One sees easily that $L$ is a subfield of $K$. On the other hand
it is well known that $K$ is maximally complete \textup{(ibid., \S5, exerc.~5)},
which shows that \textup{(MC)} holds when $a$ is a short root. But one sees easily
that the completion of $L$ is not maximally complete and that \textup{(MC)} fails
when $a$ is a long root.
\end{stmt}

\num{7.5.7}\btanchor{7.5.7} Let $A$ be an apartment of $\Imm$ and $C$ a chamber of $A$. Clearly the
retraction $\rho_{A;C}$ extends by continuity to a map, still denoted $\rho_{A;C}$,
from $\hImm$ onto $A$, and (\btref{7.4.20})\,(ii) remains valid with $\Imm$ replaced by
$\hImm$. Similarly, a retraction with respect to a sector germ extends by
continuity to $\hImm$.

\num{7.5.8}\btanchor{7.5.8} Suppose the valuation $\varphi$ dense, and let $\sigma,\tau$ be the
constants whose existence (\btref{7.4.34}) guarantees. Let $\varepsilon>0$ be small enough
that $\sigma'=(1-2\varepsilon)\sigma-\varepsilon$ and
$\tau'=(1-2\varepsilon)\tau-2\varepsilon$ are strictly positive. Let $x,y$ be two
distinct points of $\hImm$ and choose $x',y'\in\Imm$ with $d(x,x')$ and $d(y,y')$
less than $\varepsilon d(x,y)$, so that
$d(x',y')\geqslant(1-2\varepsilon)d(x,y)$. By (\btref{7.4.34}) there is $g\in G$ fixing
every point of $\Imm$ in the ball with center $x'$ and radius
$(1-2\varepsilon)\sigma d(x,y)$ and with
$d(y',g.y')\geqslant(1-2\varepsilon)\tau d(x,y)$. By continuity $g$ fixes every
point of $\hImm$ in the ball with center $x$ and radius $\sigma'd(x,y)$, and
$d(y,g.y)\geqslant\tau'd(x,y)$. In other words, after decreasing the constants
$\sigma$ and $\tau$, corollary (\btref{7.4.34}) remains true with $\Imm$ replaced by
$\hImm$.

\btsecB{7.6}{7.6. Restriction to a subgroup}

\begin{stmt}{Definition}{7.6.1}\btanchor{7.6.1}
A subset $\Phi_1$ of $\Phi$ is said to be \emph{quasi-closed} \textup{(cf.\ [\btbib{3}],
3.8)} if, for all $a,b\in\Phi_1$, the commutator group $(U_a,U_b)$ is contained in
the group generated by the $U_{pa+qb}$ for $p,q$ strictly positive integers with
$pa+qb\in\Phi_1$.
\end{stmt}

Clearly a closed subset is quasi-closed.

\num{7.6.2}\btanchor{7.6.2} Throughout the rest of 7.6, $\Phi_1$ denotes a quasi-closed root
subsystem. We write $G_1^{0}$ for the subgroup of $G$ generated by the $U_a$,
$a\in\Phi_1$; $N_1^{0}$ for the subgroup generated by the $M_a^{0}$, $a\in\Phi_1$
((\btref{6.1.2})\,(2)); and we put $T_1^{0}=T\cap N_1^{0}$. We are moreover given a subgroup
$T_1$ of $T$ containing $T_1^{0}$, and we write $G_1$ for the subgroup generated by
$G_1^{0}$ and $T_1$.

\begin{stmt}{Proposition}{7.6.3}\btanchor{7.6.3}
$\bigl(T_1,(U_a,M_a^{0}T_1)_{a\in\Phi_1}\bigr)$ is a generating root datum of type
$\Phi_1$ in $G_1$, and $\varphi_1=(\varphi_a)_{a\in\Phi_1}$ is a valuation of it.
\end{stmt}

\pf The only point that is not completely obvious is that $\varphi_1$ satisfies
(V\,3). Let $a,b\in\Phi_1$ and $h,k\in\mathbf{R}$ with
$b\notin-\mathbf{R}^{*}_{+}a$. For an arbitrary order on the set $E$ of elements of
$\Phi$ of the form $pa+qb$ with $p,q\in\mathbf{N}^{*}$, $p$ and $q$ not both even,
one has
\[
  (U_{a,k},U_{b,h})\subset
  \Bigl(\prod_{c=pa+qb\in E}U_{c,\,pk+qh}\Bigr)\cap
  \Bigl(\prod_{c\in E\cap\Phi_1}U_c\Bigr),
\]
taking into account that if $2pa+2qb\in\Phi_1$ $(p,q\in\mathbf{N}^{*})$ then
$pa+qb\in\Phi_1$, as a glance at the rank-$2$ root systems shows. Since the product
map $\prod_{c\in E}U_c\to G$ is injective (\btref{6.1.6}), (V\,3) follows. \qed

\smallskip
The objects that \S\S6 and 7 attach to the valued root datum of $G_1$ will be
denoted by the same letters as for $G$, but with the index or exponent $1$. For
example $V_1^{*}$ is the subspace of $V^{*}$ generated by $\Phi_1$, and $V_1$ is the
quotient of $V$ by the intersection of the kernels of the $a\in\Phi_1$, an
intersection we denote $L_1$. We write $\vpi$ for the canonical map of $V$ onto
$V_1$ and $\pi$ for the affine map of $A$ onto $A_1$ defined by
$\pi(\varphi+v)=\varphi_1+\vpi(v)$ for $v\in V$. One identifies $\vW_1$ with the
subgroup of $\vW$ generated by the reflections $r_a$, $a\in\Phi_1$. Note that
$\vW_1$ acts trivially on $L_1$.

It is immediate that $N_1=N_1^{0}T_1\subset N$.
\par\endgroup

% ---------------------------------------- bttranslation.tex
\begingroup\providecommand{\headrulewidth}{}
\renewcommand{\headrulewidth}{0pt}
\providecommand{\bui}{}
\renewcommand{\bui}{\mathscr{I}}
\providecommand{\nat}{}
\renewcommand{\nat}{\natural}
\providecommand{\red}{}
\renewcommand{\red}{\mathrm{red}}
\providecommand{\Isom}{}
\renewcommand{\Isom}{\operatorname{Isom}}
\providecommand{\Card}{}
\renewcommand{\Card}{\operatorname{Card}}
\providecommand{\Ker}{}
\renewcommand{\Ker}{\operatorname{Ker}}
\providecommand{\Hom}{}
\renewcommand{\Hom}{\operatorname{Hom}}
\providecommand{\Mult}{}
\renewcommand{\Mult}{\mathrm{Mult}}
\providecommand{\R}{}
\renewcommand{\R}{\mathbf{R}}
\providecommand{\Z}{}
\renewcommand{\Z}{\mathbf{Z}}
\providecommand{\N}{}
\renewcommand{\N}{\mathbf{N}}
\providecommand{\Rp}{}
\renewcommand{\Rp}{\R^{*}_{+}}
\providecommand{\Bor}{}
\renewcommand{\Bor}{\mathscr{B}}
\providecommand{\Nor}{}
\renewcommand{\Nor}{\mathscr{N}}
\providecommand{\Gor}{}
\renewcommand{\Gor}{\mathscr{G}}
\providecommand{\Hor}{}
\renewcommand{\Hor}{\mathscr{H}}
\providecommand{\Uor}{}
\renewcommand{\Uor}{\mathscr{U}}
\providecommand{\Eor}{}
\renewcommand{\Eor}{\mathscr{E}}
\providecommand{\Xor}{}
\renewcommand{\Xor}{\mathscr{X}}
\providecommand{\Yor}{}
\renewcommand{\Yor}{\mathscr{Y}}
\providecommand{\Rel}{}
\renewcommand{\Rel}{\mathscr{R}}
\providecommand{\Pw}{}
\renewcommand{\Pw}{\mathfrak{P}}
\providecommand{\Gr}{}
\renewcommand{\Gr}{\mathscr{G}}
\providecommand{\Tr}{}
\renewcommand{\Tr}{\mathscr{T}}
\providecommand{\vv}{}
\renewcommand{\vv}[1]{{}^{v}\!#1}
\providecommand{\wt}{}
\renewcommand{\wt}{\widetilde}
\providecommand{\ph}{}
\renewcommand{\ph}{\varphi}
\providecommand{\ps}{}
\renewcommand{\ps}{\psi}
\providecommand{\Ph}{}
\renewcommand{\Ph}{\Phi}
\providecommand{\Om}{}
\renewcommand{\Om}{\Omega}
\providecommand{\hP}{}
\renewcommand{\hP}{\widehat{P}}
\providecommand{\hN}{}
\renewcommand{\hN}{\widehat{N}}
\providecommand{\hbui}{}
\renewcommand{\hbui}{\widehat{\bui}}
\providecommand{\dual}{}
\renewcommand{\dual}{{}^{\vee}}
\newenvironment{stmt}[2]
  {\medskip\par\noindent\textit{#1} \textbf{(#2)}. --- \itshape\ignorespaces}
  {\par\medskip}
\newenvironment{stmtup}[2]
  {\medskip\par\noindent\textit{#1} \textbf{(#2)}. --- \ignorespaces}
  {\par\medskip}
\providecommand{\numpar}{}
\renewcommand{\numpar}[1]{\medskip\par\noindent\textbf{(#1)}\ \ignorespaces}
\setlength{\parindent}{1.4em}
\allowdisplaybreaks
\bigskip

\pg{185}
\noindent
If $f$ is a quasi-concave function on $\Ph_1$, we denote by $U^1_f$ the corresponding
subgroup of $G_1$ (\btref{6.4.2}); the notations $U^1_\Om$, $P^1_\Om$ and $\hP^1_\Om$
(for $\Om\subset A_1$) are self-explanatory.

\begin{stmt}{Proposition}{7.6.4}\btanchor{7.6.4}
(i) There exists one and only one map $\wt\pi:G_1.A\to\bui_1$ from the set of the
transforms of the points of $A$ by the elements of $G_1$ into the building of $G_1$,
extending $\pi:A\to A_1$ and commuting with the actions of $G_1$ on $G_1.A$ and on
$\bui_1$.

(ii) One has $\wt\pi^{-1}(A_1)=A$, and the inverse image under $\wt\pi$ of any apartment
(resp.\ half-apartment, wall) of $\bui_1$ is an apartment (resp.\ half-apartment, wall)
of $\bui$.

(iii) There exists one and only one law of operation $(x,v)\mapsto x+v$ of $L_1$ on
$G_1.A$, extending the action of $L_1$ on $A$ and such that $g.(x+v)=g.x+v$ for
$g\in G_1$, $x\in G_1.A$ and $v\in L_1$. The map $\wt\pi$ passes to the quotient by
$L_1$ and defines a bijection of $(G_1.A)/L_1$ onto $\bui_1$.
\end{stmt}

Let us first show that
\begin{equation}\tag{1}
\pi\circ\nu(n)=\nu_1(n)\circ\pi\qquad\text{for all }n\in N_1 .
\end{equation}
This is immediate if there exist $a\in\Ph_1$ and $k\in\R$ such that $n\in M^1_{a,k}$:
indeed, $\nu(n)$ (resp.\ $\nu_1(n)$) is then the orthogonal reflection with respect to
the hyperplane of equation $a(x-\ph)+k=0$ (resp.\ $a(x-\ph_1)+k=0$). If $n\in T_1$,
there exists $v\in V$ (resp.\ $v_1\in V_1$) such that $n.y=y+v$ (resp.\ $n.y_1=y_1+v_1$)
for all $y\in A$ (resp.\ $y_1\in A_1$). One has, for $a\in\Ph_1$ and $k\in\R$:
\[
nU_{a,k}n^{-1}=U_{a,\,k-a(v)}=U_{a,\,k-a(v_1)}
\]
whence $a(v)=a(v_1)$ for all $a\in\Ph_1$. Consequently, $v_1=\vv\pi(v)$ and
$\pi(n.y)=n.\pi(y)$ for all $y\in A$.

Formula (1) then follows from what precedes, since $N_1$ is generated by $T_1$ and by
the union of the $M^1_{a,k}$.

Now let $x\in A$. It is immediate that $U^1_{\pi(x)}\subset U_x$. Since
$G_1=U^1_{\pi(x)}.N_1.U^1_{\pi(x)}$ (\btref{7.3.4}), one has, in view of (1):
\begin{equation}\tag{2}
G_1\cap\hP_x=U^1_{\pi(x)}\cdot(N_1\cap\hP_x)\cdot U^1_{\pi(x)}\subset\hP^1_{\pi(x)} .
\end{equation}

Now let $g\in G_1$ and put $\Om=A\cap g^{-1}.A$. Suppose $\Om$ non-empty and let
$x\in\Om$. Put $\Om'=\Om\cap(x+L_1)$. In view of (\btref{7.4.8}), there exists $n\in N$ such
that $g.y=n.y$ for all $y\in\Om$, whence
\[
g\in n\hP_\Om\cap G_1=n\hP_\Om\cap U^1_{\pi(x),D_1}.N_1.U^1_{\pi(x)}
\subset n\hP_\Om\cap U_{x,D}.N_1.\hP_{\Om'}
\]
denoting by $D$ (resp.\ $D_1$) a vector chamber of $\Ph$ (resp.\ $\Ph_1$) such that
$D\subset\pi^{-1}(D_1)$. In view of (\btref{7.4.15}), there therefore exists $n_1\in N_1$ such
that
\[
n\in(N\cap U_{x,D})\,n_1\,(N\cap\hP_{\Om'})=Hn_1\hN_{\Om'}=n_1\hN_{\Om'},
\]
whence
\begin{equation}\tag{3}
g\in n_1(\hP_u\cap G_1)\subset n_1\hP^1_{\pi(u)}\qquad
\text{for all }u\in A\cap g^{-1}.A\cap(x+L_1).
\end{equation}

This being so, let us prove (i). The uniqueness of $\wt\pi$ is clear. To show its
existence, it suffices to show that if $x\in A$ and $g\in G_1$ are such that $g.x\in A$,
then $\pi(g.x)=g.\pi(x)$. Let then $n_1\in N_1$ satisfy (3); in view of (1) and (2), one
has
\[
\pi(g.x)=\pi(n_1.x)=n_1.\pi(x)=g.\pi(x).
\]

\pg{186}
Let us prove (ii). Let $x\in A$ and $g\in G_1$ be such that $\pi(g.x)\in A_1$. In view
of (\btref{7.4.8}), there exists $n\in N_1$ such that $g\in n\hP^1_{\pi(x)}$. Since
$\hP^1_{\pi(x)}=\hN^1_{\pi(x)}.U^1_{\pi(x)}\subset N_1.\hP_x$, one has
$g.x\in N_1.x\subset A$, whence (ii).

Let us prove (iii). The uniqueness of the sought law of operation is clear. To show its
existence, it suffices to show that if $z=g.y$, with $y,z\in A$, and if $g\in G_1$ and
$v\in L_1$, then $g.(y+v)=z+v$. Now one has in that case
\[
\wt\pi(g.(y+v))=g.\wt\pi(y+v)=g.\wt\pi(y)=\wt\pi(g.y)=\pi(z)\in A_1
\]
and $g.(y+v)\in A$ in view of (ii). Let then $n_1\in N_1$ satisfy (3), with $x=y$. One
has
\[
g.(y+v)=n_1.(y+v)=n_1.y+\vv\nu(n_1).v=z+v
\]
since $\vv\nu(n_1)\in\vv W_1$ is the identity on $L_1$.

\begin{stmt}{Corollary}{7.6.5}\btanchor{7.6.5}
Let $p$ be the orthogonal projection of $V$ onto $L_1$ and let $S_1$ be the subgroup
formed by the $t\in T$ such that $p\circ\nu(t)=0$. Let $\Om$ be a non-empty subset
of $A$.

(i) One has $T^0_1=T\cap G^0_1\subset S_1$.

(ii) If $T_1\subset S_1$, one has $G_1\cap\hP_\Om=\hP^1_{\pi(\Om)}$.

(iii) One has
$G_1\cap\hP_\Om=G_1\cap\hP_{\Om+L_1}=(G_1\cap G^0_1.S_1)\cap\hP_\Om\subset
\hP^1_{\pi(\Om)}$.
\end{stmt}

Assertion (i) follows at once from (\btref{6.1.2})(12), (\btref{6.1.13}) and (\btref{6.2.10})(ii). Let us prove
(ii). In view of (\btref{7.6.4})(i), it suffices to show that $\hP^1_{\pi(\Om)}\subset\hP_\Om$
as soon as $T_1\subset S_1$, or again, in view of (7.1.8), that
$\hN^1_{\pi(\Om)}\subset\hP_\Om$. Let then
$n\in\hN^1_{\pi(\Om)}\subset N_1\subset N^0_1.S_1$. Every element of $N^0_1$ or of $S_1$
leaves invariant the affine subspaces of the form $y+\Ker p$ for $y\in A$; on the other
hand, $\nu(n)$ leaves invariant the affine subspaces of the form $x+L_1$ for $x\in\Om$
((\btref{7.6.4})(i)). It follows that $n\in\hP_\Om$, whence (ii).

Let us prove (iii). The first equality follows from (\btref{7.6.4})(iii) and the last inclusion
from (\btref{7.6.4})(i). It remains to show that $G_1\cap\hP_\Om\subset G^0_1.S_1$. Let
$g\in G_1\cap\hP_\Om$. Since the buildings of $G_1$ and of $G^0_1$ are the same, there
exists $h\in G^0_1\cap\hP^1_{\pi(\Om)}$ such that $h^{-1}.A_1=g.A_1$. In view of (ii),
one has $h\in\hP_\Om$ and $n=hg\in N_1\cap\hP_\Om$. Let us write $n=tn_1$ with
$n_1\in N^0_1$ and $t\in T_1$, and let $x\in\Om$. One has on the one hand
\begin{equation}\tag{1}
\nu(n)(x+\Ker p)=\nu(t)(x+\Ker p)
\end{equation}
and on the other hand
\begin{equation}\tag{2}
\nu(n)(x+\Ker p)=\nu(n).x+\vv\nu(n)(\Ker p)=x+\Ker p .
\end{equation}
Comparing (1) and (2), one sees that $t\in S_1$ and one has indeed
\[
g=h^{-1}n\in G^0_1.N^0_1.S_1=G^0_1.S_1 .
\]

\bigskip
\pg{187}

\btsecA{8}{8. Bornology defined by a valued root datum}

\medskip

\noindent
We keep the notations of \S\S6 and 7.

\medskip

\btsecB{8.1}{8.1. Bornologies compatible with a valued root datum}

\begin{stmt}{Definition}{8.1.1}\btanchor{8.1.1}
A bornology $\Bor$ on $G$ is said to be \upshape weakly compatible\itshape\ with the
valued root datum of $G$ (or, more simply, with the valuation $\ph$) if it is
compatible with the group law of $G$, if the subgroups $U_{a,k}$ are bounded and the
subgroups $U_a$ unbounded for $\Bor$, whatever $a\in\Ph$ and $k\in\R$.

One says that $\Bor$ is \upshape compatible\itshape\ with $\ph$ if moreover the
subgroup $H_{[0]}$ defined in (\btref{6.4.14}) is bounded for $\Bor$.
\end{stmt}

\begin{stmt}{Proposition}{8.1.2}\btanchor{8.1.2}
Let $\Bor$ be a bornology weakly compatible with $\ph$. Let $x$ be a point of $A$ and
$D$ a vector chamber.

(i) For a subset $X$ of $G$ to be bounded for $\Bor$, it is necessary and sufficient
that
\[
(U_{x+D}.U_{x-D}.X.U_{x-D}.U_{x+D})\cap N
\]
be bounded for $\Bor$.

(ii) The set $\Bor\cap\Pw(N)$ of bounded subsets of $N$ is a bornology on $N$,
compatible with the group law of $N$, and satisfying the two following conditions:

\upshape(BCN 1)\itshape\ The image under $\nu$ of any bounded subset of $N$ is a
bounded subset of $\Isom A$ (endowed with its natural bornology) $((3.1.2)\ b))$.

\upshape(BCN 2)\itshape\ Whatever $a\in\Ph$ and $k\in\R$, the intersection
\[
T\cap(U_{a,k}U_{-a,k}U_{a,k}U_{-a,k}U_{a,k})
\]
is bounded.

(iii) Every bounded subset of $U_a$ is contained in some $U_{a,k}$
$(a\in\Ph,\ k\in\R)$.
\end{stmt}

By (\btref{7.3.4}) and (7.2.6)(2), one has $G=U_{x-D}U_{x+D}NU_{x+D}U_{x-D}$, whence
\[
X\subset U_{x-D}.U_{x+D}.\bigl((U_{x+D}U_{x-D}XU_{x-D}U_{x+D})\cap N\bigr).U_{x+D}.U_{x-D}.
\]
On the other hand, $U_{x+D}$ and $U_{x-D}$ are bounded for $\Bor$ since one has, for
example (\btref{6.4.9}),
\[
U_{x+D}=\prod_{a\in\Ph^{+\,\red}_{D}}U_{a,\,-a(x)} .
\]
One deduces (i).

Let us prove (BCN 1) by contradiction. Let $X$ be a bounded subset of $N$ whose image
$\nu(X)$ is not bounded in $\Isom A$. Up to multiplying $X$ by a finite
\pg{188}
subset, one may suppose that $\nu(X)$ is stable under left multiplication by the
stabilizer of a special point, and $\nu(X)$ then contains translations of arbitrarily
large length. More precisely, $X$ contains a sequence $(x_n)$ such that
$t_n=\nu(x_n)$ is a translation of length $|t_n|$ tending to infinity and of direction
$t_n/|t_n|$ tending to an element $v\in V$ of length $1$. Let then $a\in\Ph$ be such
that $a(v)>0$. One has $\lim_{n\to\infty}a(t_n)=+\infty$ and
\[
U_a=\bigcup_n U_{a,\,-a(t_n)}=\bigcup_n x_nU_{a,0}x_n^{-1}\subset X.U_{a,0}.X^{-1}
\]
is bounded, which is absurd.

This proves (ii), for (BCN 2) is obvious.

Let us prove (iii) by contradiction. Let $(u_n)$ be a bounded sequence of points of
$U_a$ such that $\lim_{n\to\infty}\ph_a(u_n)=-\infty$ and let
$k=\sup_n\ph_a(u_n)$. By (V 5), one has $m(u_n)\in U_{-a,-k}u_nU_{-a,-k}$ and the set
of reflections $r_{a,\ph_a(u_n)}=\nu(m(u_n))$ is bounded, which is absurd.

\begin{stmt}{Proposition}{8.1.3}\btanchor{8.1.3}
Let $D$ be a vector chamber and let $\Uor^+_D$ be the smallest bornology on $U^+_D$
compatible with the group law and containing the $U_{a,k}$ for $a\in\Ph^+_D$ and
$k\in\R$. Let $X\subset U^+_D$; the following conditions are equivalent:

(i) $X$ is bounded for $\Uor^+_D$;

(ii) whatever the order in which the roots $a\in\Ph^{+\,\red}_D$ are arranged, there
exists $k\in\R$ such that
$X\subset\prod_{a\in\Ph^{+\,\red}_D}U_{a,k}$;

(iii) there exists $x\in A$ such that $X\subset U_{x+D}$.
\end{stmt}

It is clear that (iii) implies (ii) (\btref{6.4.9}) and that (ii) implies (i). Let
$X\in\Uor^+_D$; then $X$ is contained in a finite product of subgroups
$U_{a_i,k_i}$ with $a_i\in\Ph^+_D$ and $k_i\in\R$. For $a\in\Ph^+_D$, put
$h(a)=\inf\{k_i\mid a_i=a\}$ and let $x\in A$ be such that $a(x)>-h(a)$ for all
$a\in\Ph^+_D$. One then has $X\subset U_{x+D}$.

\begin{stmt}{Proposition}{8.1.4}\btanchor{8.1.4}
Let $\Nor$ be a bornology on $N$, compatible with the group law and satisfying the two
conditions \upshape(BCN 1)\itshape\ and \upshape(BCN 2)\itshape\ of (\btref{8.1.2}). Let $\Gor$
be the smallest bornology on $G$ containing $\Nor$ and the $U_{a,k}$ for $a\in\Ph$ and
$k\in\R$ and such that $M,M'\in\Gor$ implies $MM'\in\Gor$.

(i) $\Gor$ is a bornology weakly compatible with $\ph$ and in particular is compatible
with the group law of $G$.

(ii) $\Gor$ induces on $N$ the given bornology $\Nor$ and, for every vector chamber
$D$, the bornology $\Gor$ induces on $U^+_D$ the bornology $\Uor^+_D$ (\btref{8.1.3}).

(iii) Let $D$ be a vector chamber. For a subset $X$ of $G$ to be bounded for $\Gor$, it
is necessary and sufficient that there exist $Y^+\subset\Uor^+_D$,
$Y^-\subset\Uor^-_D=\Uor^+_{-D}$ and $Y^0\subset\Nor$ such that
\[
X\subset Y^+.Y^-.Y^0 .
\]
\end{stmt}

Let us introduce a few notational conventions: if $\Xor$ and $\Yor$ are two sets of
subsets of $G$, we put $\Xor.\Yor=\{XY\mid X\in\Xor,\ Y\in\Yor\}$ and
$\Xor^{-1}=\{X^{-1}\mid X\in\Xor\}$, and we write $\Xor\prec\Yor$ if every element of
$\Xor$ is contained in a finite union of elements of $\Yor$. Let $\Eor=\Nor\cap\Pw(T)$
be the bornology induced by $\Nor$ on $T$ and, for $a\in\Ph$, let us denote by $\Uor_a$
the set of the $U_{a,k}$ for $k\in\R$.

\pg{189}
Let then $\Hor$ be the set of the subsets $X$ of $G$ satisfying the condition of (iii)
for a given vector chamber $D$. Since $G=U^+_DU^-_DN$ (\btref{6.1.12}), this is a bornology on
$G$. Let us show that $\Hor$ is a bornology compatible with the group law of $G$, or
again that
\begin{equation}\tag{1}
\Hor^{-1}.\Hor\prec\Hor .
\end{equation}

Every element of $\Hor^{-1}$ is contained in a finite union of finite products of sets
belonging either to $\Eor$, or to one of the $\Uor_a$ for $a\in\Ph$, or to
$\{\{n\}\mid n\in N\}$. Since $\Uor_{-a}=n\Uor_an^{-1}$ for $n\in M_a$ and since $N$ is
generated by $T$ and by the union of the $M_b$ for $b$ running through the basis $\Pi_D$
of $\Ph$ corresponding to $D$ (\btref{6.1.11}), one sees that every element of $\Hor^{-1}$ is
contained in a finite union of finite products of sets belonging either to $\Eor$, or to
one of the $\Uor_a$ for $a\in\Ph^+_D$, or to the set of the $\{m\}$ for $b\in\Pi_D$ and
$m\in M_b$.

By (\btref{8.1.3}), one has $\Uor_a.\Hor\prec\Hor$, and condition (BCN 1) implies that
$\Eor.\Hor\prec\Hor$. To prove (1), it therefore suffices to show that:
\begin{equation}\tag{2}
\text{\itshape for $b\in\Pi_D$ and $m\in M_b$, one has $m\Hor\prec\Hor$.}
\end{equation}

Let $U'^+$ (resp.\ $U'^-$) be the group generated by the $U_c$ for $c\in\Ph^+_D$
(resp.\ $\Ph^-_D$) not proportional to $b$. Put $\Uor^\pm=\Uor^\pm_D$,
$\Uor'^+=\Uor^+\cap\Pw(U'^+)$ and $\Uor'^-=\Uor^-\cap\Pw(U'^-)$. The following
relations are obvious:
\begin{align}
&m\Uor_bm^{-1}=\Uor_{-b},\qquad m\Uor_{-b}m^{-1}=\Uor_b;\tag{3}\\
&m\Uor'^+m^{-1}=\Uor'^+,\qquad m\Uor'^-m^{-1}=\Uor'^-;\tag{4}\\
&\Eor.\Uor_b\prec\Uor_b.\Eor,\qquad \Eor.\Uor_{-b}\prec\Uor_{-b}.\Eor;\tag{5}\\
&\Eor.\Uor'^+\prec\Uor'^+.\Eor,\qquad \Eor.\Uor'^-\prec\Uor'^-.\Eor.\tag{6}
\end{align}

We shall show that:
\begin{equation}\tag{7}
\Uor_{-b}.\Uor_b\prec\Uor_b.\Uor_{-b}.\Eor\ \cup\ \Uor_b.\Uor_{-b}.m\Eor .
\end{equation}
For this, it suffices to show that, whatever $h,k\in\R$, one has:
\begin{equation}\tag{8}
\{U_{-b,h}.U_{b,k}\}\prec\Uor_b.\Uor_{-b}.\Eor\ \cup\ \Uor_b.\Uor_{-b}.m\Eor
\end{equation}
and one may suppose $h=k<0$. For $u\in U_{-b}$ with $r=\ph_b(u)\leqslant-2k$, one has:
\[
u\in U_{b,-r}mTU_{b,-r}\subset U_{b,2k}mTU_{b,2k}.
\]
Consequently, there exists a subset $Y$ of $T$ such that:
\begin{equation}\tag{9}
U_{-b,k}-U_{-b,-2k}\subset U_{b,2k}mYU_{b,2k}
\end{equation}
and one may choose $Y$ in such a way that:
\begin{equation}\tag{10}
Y\subset m^{-1}U_{b,2k}U_{-b,k}U_{b,2k}.
\end{equation}
Let $s\in\Gamma_b$, $s>2k$. There exists $t\in T$ such that
$m^{-1}\in tM^0_{b,s}\subset tU_{b,s}U_{-b,-s}U_{b,s}$, and (10) implies that
$Y\in\Eor$ by (BCN 2).

\pg{190}
On the other hand, one has by (\btref{6.3.2})
$U_{-b,-2k}U_{b,k}\subset U_{b,k}U_{-b,-2k}H$, and one sees as above that there exists
$Y'\in\Eor$ such that:
\begin{equation}\tag{11}
U_{-b,-2k}.U_{b,k}\subset U_{b,k}.U_{-b,-2k}.Y' .
\end{equation}
It then follows from (9) and (11) that:
\[
U_{-b,k}.U_{b,k}\subset(U_{b,2k}.mY.U_{b,2k})\cup(U_{b,k}.U_{-b,-2k}.Y').
\]
Taking (3) and (5) into account, one deduces (8), whence (7). Finally, using (3), (4),
(6) and (7), one sees that:
\begin{align*}
m\Hor&\prec m.\Uor'^+.\Uor_b.\Uor_{-b}.\Uor'^-.\Nor
=\Uor'^+.\Uor_{-b}.\Uor_b.\Uor'^-.m\Nor\\
&\prec\Uor'^+.\Uor_b.\Uor_{-b}.\{1,m\}.\Eor.\Uor'^-.\Nor
\prec\Uor^+.\Uor_b.\Uor_{-b}.\Uor'^-.\{1,m\}.\Eor.\Nor\\
&\prec\Uor^+.\Uor^-.\Nor=\Hor .
\end{align*}

It is now clear that $\Hor=\Gor$.

On the other hand, for $Y^+\subset U^+_D$, $Y^-\subset U^-_D$ and $Y^0\subset N$, one
has $((6.1.15)\ c))$
\begin{align*}
U^+_D\cap Y^+Y^-Y^0&\subset Y^+\\
N\cap Y^+Y^-Y^0&\subset Y^0 .
\end{align*}
Consequently, $\Gor=\Hor$ induces on $N$ (resp.\ $U^+_D$) the bornology $\Nor$
(resp.\ $\Uor^+_D$). In view of the invariance of $\Gor$ under inner automorphisms, one
deduces that $\Gor$ is weakly compatible with $\ph$ and satisfies (ii), which completes
the proof.

\begin{stmt}{Corollary}{8.1.5}\btanchor{8.1.5}
The map which, to a bornology on $G$, associates the induced bornology on $N$, is a
bijection from the set of the bornologies weakly compatible with $\ph$ onto the set of
the bornologies on $N$ compatible with the group law of $N$ and satisfying conditions
\upshape(BCN 1)\itshape\ and \upshape(BCN 2)\itshape.
\end{stmt}

This follows from (\btref{8.1.2}) and (\btref{8.1.4}).

The inverse bijection of the one described in the corollary will be denoted
$\Nor\mapsto\Bor(\Nor)$.

\begin{stmtup}{Remarks}{8.1.6}\btanchor{8.1.6}
\textit{a)} A bornology $\Bor$ on $G$ which contains the $U_{a,k}$ for $a\in\Ph$ and
$k\in\R$, which induces on $N$ a bornology $\Nor$ compatible with the group law of $N$
satisfying condition (BCN 1), and which is such that $M,M'\in\Bor$ implies
$MM'\in\Bor$, is weakly compatible with $\ph$: it is indeed clear that $\Nor$ satisfies
(BCN 2) and that $\Bor$ satisfies (\btref{8.1.2})(i); consequently, one has $\Bor=\Bor(\Nor)$.

\textit{b)} Let us take up the notations of (\btref{8.1.4}). It is not true that every bounded
subset of $G$ (for $\Gor$) is contained in a set of the form
\begin{equation}\tag{1}
Y^+.Y^0.Y^-\qquad\text{with }Y^+\in\Uor^+,\ Y^-\in\Uor^-\text{ and }Y^0\in\Nor
\end{equation}
nor even that every bounded subset of $U^+.T.U^-$ is contained in a product of the form
(1) with moreover $Y^0\in\Eor$. A counterexample is furnished by the set
\[
m.(U_{b,k}-\{1\})
\]
for $b\in\Ph^+$, $m\in M_b$ and $k\in\R$.
\end{stmtup}

\pg{191}
\begin{stmt}{Definition}{8.1.7}\btanchor{8.1.7}
One calls \upshape bornology defined by $\ph$\itshape, and denotes $\Bor(\ph)$, the
bornology on $G$ which is the inverse image of the natural bornology of $\Isom\bui$
under the canonical map from $G$ into $\Isom\bui$ (\btref{7.4.20}).
\end{stmt}

One will note that, in the case where $\ph$ is discrete, this bornology coincides with
the one defined by the Tits system defined by $\ph$ (\btref{3.1.8}).

\begin{stmt}{Proposition}{8.1.8}\btanchor{8.1.8}
The bornology $\Bor(\ph)$ is compatible with $\ph$ and coincides with the bornology
$\Bor(\Nor)$, where $\Nor$ is the bornology on $N$ which is the inverse image under
$\nu$ of the natural bornology of $\Isom A$.
\end{stmt}

Let $a\in\Ph$ and $k\in\R$ and let $x\in A$ be such that $a(x)\geqslant-k$. One then has
$U_{a,k}\subset\hP_x$, which shows that $U_{a,k}$ is bounded for $\Bor(\ph)$. On the
other hand, corollary (\btref{7.4.27}) shows that $U_a$ is not bounded. Moreover, all of $H$,
hence \emph{a fortiori} $H_{[0]}$, is bounded, since $H$ fixes every point of $A$.
Finally, the image of a subset $X$ of $N$ is bounded in $\Isom\bui$ if and only if
$\nu(X)$ is bounded in $\Isom A$, whence the last assertion.

\begin{stmt}{Corollary}{8.1.9}\btanchor{8.1.9}
Let $\Bor$ be a bornology weakly compatible with $\ph$. Suppose that, for every subset
$X$ of $\nu(T)$ bounded in $\Isom A$, there exists a subset $Y$ of $T$ belonging to
$\Bor$ such that $X=\nu(Y)$. Then the bornology defined by $\ph$ is generated by
$\Bor.H$.
\end{stmt}

Indeed, the bornology which is the inverse image under $\nu$ of the natural bornology of
$\Isom A$ is generated by $\Nor.H$, where $\Nor$ is the bornology induced by $\Bor$
on $N$.

\begin{stmt}{Theorem}{8.1.10}\btanchor{8.1.10}
Let $\ps$ be another valuation of the root datum of $G$. The following conditions are
equivalent:

\upshape a)\itshape\ the valuations $\ph$ and $\ps$ are equivalent (\btref{6.2.5});

\upshape b)\itshape\ the bornologies $\Bor(\ph)$ and $\Bor(\ps)$ are equal;

\upshape c)\itshape\ the bornologies $\Bor(\ph)$ and $\Bor(\ps)$ induce the same
bornology on $U_a$ for every $a\in\Ph$;

\upshape d)\itshape\ the bornologies $\Bor(\ph)$ and $\Bor(\ps)$ induce the same
bornology on $N$;

\upshape e)\itshape\ for every $a\in\Ph$, there exists a bijection
$f_a:{}^{\ph}\Gamma_a\to{}^{\ps}\Gamma_a$ such that $\ps_a=f_a\circ\ph_a$.
\end{stmt}

Suppose $\ph$ and $\ps$ equivalent and let $i$ be the canonical bijection of the
building ${}^{\ph}\bui$ onto ${}^{\ps}\bui$ (\btref{7.4.3}). In view of $(7.4.3)\ b)$, the
restriction of $i$ to ${}^{\ph}A$ is bounded, hence so is $i$ itself
$((7.4.3)\ a)$, (\btref{7.4.18}) and (\btref{7.4.20})$)$. One deduces at once that
$\Bor(\ph)\subset\Bor(\ps)$. One sees in the same way that
$\Bor(\ps)\subset\Bor(\ph)$, whence the implication a)~$\Rightarrow$~b).

It is clear that b) implies c) and d). Let us show that c) implies d). For this it
suffices to remark that a subset $Y$ of $T$ is bounded for $\Bor(\ph)$ if and only if,
for every $a\in\Ph$ and every subset $X\subset U_a$ bounded for $\Bor(\ph)$, the set of
the $txt^{-1}$ for $t\in Y$ and $x\in X$ is bounded for the bornology induced by
$\Bor(\ph)$ on $U_a$, and that a subset of $N$ is bounded if and only if it is contained
in the union of a finite number of translates of bounded subsets of $T$.

Let us show that d) implies e). For this, let us consider the equivalence relation
${}^{\ph}\Rel_a:\ph_a(u)=\ph_a(v)$ on $U_a$ (for $a\in\Ph$). It is equivalent to the
relation
\[
m(u)\in Hm(v)
\]
\pg{192}
or again to the fact that the two reflections $\nu(m(u))$ and $\nu(m(v))$ are equal.
But, whatever $u$ and $v$ in $U_a$, $\nu(m(u))$ and $\nu(m(v))$ are two reflections with
respect to parallel hyperplanes, and to say that they are equal amounts to saying that
they generate a bounded subgroup of $\Isom A$. In view of (\btref{8.1.8}), this means that
$m(u)$ and $m(v)$ generate a \emph{bounded} subgroup of $N$. Consequently, d) implies
that the equivalence relations ${}^{\ph}\Rel_a$ and ${}^{\ps}\Rel_a$ are the same for
every $a$, whence e).

Let us show finally that e) implies a). Up to replacing $\ph$ and $\ps$ by equipollent
valuations, one may suppose that, for every $a$ belonging to a basis $\Pi$ of $\Ph$,
there exists an element $u_a\in U_a$ such that $\ph_a(u_a)=\ps_a(u_a)=0$. In view of
(\btref{6.2.7})(1), one has, for $u\in U_a$:
\[
\ph_{-a}(m(u_a)um(u_a)^{-1})=\ph_a(u),\qquad
\ps_{-a}(m(u_a)um(u_a)^{-1})=\ps_a(u)
\]
whence $f_{-a}=f_a$.

Now let $u,v\in U_a$, $k=\ph_a(u)$ and $\ell=\ph_a(v)$. Again by (\btref{6.2.8})(1), one has
\begin{align*}
\ph_{-a}(m(u)vm(u)^{-1})&=\ell-2k\\
\ps_{-a}(m(u)vm(u)^{-1})&=f_a(\ell)-2f_a(k)
\end{align*}
whence $f_a(\ell-2k)=f_a(\ell)-2f_a(k)$. Since $0,k\in\Gamma_a$ and
$\Gamma_a\supset\Gamma_a+2\wt\Gamma_a$ (\btref{6.2.16}), one sees that $pk\in\Gamma_a$ and that:
\[
f_a(pk)=pf_a(k)\qquad\text{for all }k\in\Gamma_a\text{ and all }p\in\Z .
\]

For $u\in U_a$, the subgroup ${}^{\ph}U_{a,\ph_a(u)}$ is generated by the $v\in U_a$
such that $\ph_a(v)=\ph_a(u)$. It therefore coincides with the subgroup
${}^{\ps}U_{a,\ps_a(u)}$, whence one deduces easily that $f_a$ is decreasing. Since
$\Gamma_a$ is dense in $\R$, or else is a subgroup of $\R$ isomorphic to $\Z$, one
deduces easily that \emph{there exists a strictly positive constant $\lambda(a)$ such
that $f_a$ is multiplication by $\lambda(a)$}, and this for every $a\in\Pi$.

Now let $a\in\Pi$ and $b\in\Ph$ be such that $f_b$ is also multiplication by a constant
$\lambda(b)$. Using once again (\btref{6.2.7})(1), one has, for $v\in U_b$:
\[
\ph_{r_a(b)}(m(u_a)vm(u_a)^{-1})=\ph_b(v),\qquad
\ps_{r_a(b)}(m(u_a)vm(u_a)^{-1})=\ps_b(v)
\]
whence one deduces that $f_{r_a(b)}$ is multiplication by $\lambda(b)$.

Since the $r_a$ for $a\in\Pi$ generate the Weyl group $\vv W$ of $\Ph$, and since every
$b\in\Ph^{\red}$ is the transform of an element of $\Pi$ by an element of $\vv W$, one
sees that \emph{$f_b$ is multiplication by a constant $\lambda(b)$ for every
$b\in\Ph^{\red}$}, hence also for every $b\in\Ph$ by (V 4).

Finally, for $u\in U_a$, $v\in U_b$, $k=\ph_a(u)$ and $\ell=\ph_b(v)$, one has
\begin{align*}
\ph_{r_a(b)}(m(u)vm(u)^{-1})&=\ell-kb(a\dual)\\
\ps_{r_a(b)}(m(u)vm(u)^{-1})&=\ps_b(v)-\ps_a(u)b(a\dual)
=\lambda(b)\ell-\lambda(a)kb(a\dual)\\
&=\lambda(b)\bigl(\ell-kb(a\dual)\bigr)
\end{align*}
whence one deduces that if $b(a\dual)\neq0$, that is to say \emph{if $a$ and $b$ are not
orthogonal, one has $\lambda(a)=\lambda(b)$}. Consequently, the function $\lambda$ is
indeed constant on the irreducible components of $\Ph$, which completes the proof.

\pg{193}
\begin{stmt}{Corollary}{8.1.11}\btanchor{8.1.11}
Let $\gamma:G\to G$ be an automorphism of bornological group preserving $T$ and the set
of the $U_a$. There exists one and only one map $\gamma_*:\bui\to\bui$ possessing the
following properties:

\upshape a)\itshape\ $\gamma_*(g.x)=\gamma(g).\gamma_*(x)$ for all $g\in G$ and
$x\in\bui$;

\upshape b)\itshape\ $\gamma_*(A)=A$ and the restriction of $\gamma_*$ to $A$ is an
affinity.

The map $\gamma_*$ is bijective and equicontinuous: more precisely, there exist
constants $m,M\in\Rp$ such that
\[
m\,d(x,y)\leqslant d(\gamma_*(x),\gamma_*(y))\leqslant M\,d(x,y)
\]
whatever $x,y\in\bui$.
\end{stmt}

The existence and the uniqueness of $\gamma_*$ are immediate consequences of (\btref{8.1.10})
and (\btref{7.4.3}), applied to the valuations $\ph$ and $\ph\circ\gamma$. The rest of the
statement follows at once from a), b) and (\btref{7.4.20}), taking into account that an affinity
of a Euclidean space onto itself possesses the properties in question.

\begin{stmtup}{Remark}{8.1.12}\btanchor{8.1.12}
One can show that, if $\gamma:G\to G$ is an automorphism of bornological group, there
exists one and only one map $\gamma_*:\bui\to\bui$ possessing property (\btref{8.1.11}) a) and
such that:

\textit{b$'$)} \emph{for every bounded convex subset $\Om$ of $A$, the restriction of
$\gamma_*$ to $\Om$ is an isomorphism of affine sets onto a bounded convex subset of an
apartment.}

If $\ph$ is dense, a) already implies the uniqueness of $\gamma_*$ (this follows from
(\btref{8.2.1})). If $\ph$ is discrete, $\gamma_*$ is characterized by a) and:

\textit{b$''$)} \emph{the restriction of $\gamma_*$ to every facet $F$ of $\bui$ is an
affine map from $F$ onto another facet.}

In the discrete case, these assertions follow from \S\S2 and 3, notably from 3.3
and 3.5.
\end{stmtup}

\begin{stmt}{Proposition}{8.1.13}\btanchor{8.1.13}
Let us take up the hypotheses and notations of 7.6. Let $p$ be the orthogonal projection
of $V$ onto $L_1$ and let $S_1$ be the subgroup of $T$ formed by the elements $t\in T$
such that $p\circ\nu(t)=0$.

(i) $\Bor(\ph)$ induces on $G_1$ a bornology compatible with $\ph_1$.

(ii) For $\Bor(\ph)$ to induce on $G_1$ the bornology $\Bor(\ph_1)$, it is necessary and
sufficient that $G_1$ be contained in $G^0_1S_1$.

(iii) For a subset $M$ of $G_1$ to be bounded for $\Bor(\ph)$, it is necessary and
sufficient that it be bounded for $\Bor(\ph_1)$ and contained in a subset of the form
$G^0_1.(p\circ\nu)^{-1}(X)$, where $X$ is a bounded subset of $L_1$.

(iv) For a subgroup of $G_1$ to be bounded for $\Bor(\ph)$, it is necessary and
sufficient that it be bounded for $\Bor(\ph_1)$ and contained in $G^0_1.S_1$.
\end{stmt}

Assertion (i) is immediate. Let us prove (iii). Let $M\subset G_1$, bounded for
$\Bor(\ph_1)$ (resp.\ $\Bor(\ph)$). There exist subsets $Y^+\subset U^+_1$,
$Y^-\subset U^-_1$ and $Z\subset T_1$ bounded for $\Bor(\ph_1)$ (resp.\ $\Bor(\ph)$),
and a finite subset $F\subset N_1\cap G^0_1$, such that $M\subset Y^+.Y^-.F.Z$. Since
$Y^+$, $Y^-$ and $F$ are bounded for $\Bor(\ph)$ (resp.\ $\Bor(\ph_1)$), one sees
easily that one is reduced to proving that the image under $\nu$ of a subset $Z$ of
$T_1$ in $\Isom A$ is bounded if and only if $\nu_1(Z)$ is bounded in $\Isom A_1$ and
$p\circ\nu(Z)$ is bounded in $L_1$, which is immediate.

\pg{194}
Finally, in view of (\btref{7.6.5})(i), the homomorphism $p\circ\nu:T\to L_1$ extends in a
unique way to a homomorphism from $G^0_1.T$ into $L_1$ with kernel $G^0_1.S_1$.
Assertions (ii) and (iv) then follow from (iii).

\medskip

\btsecB{8.2}{8.2. Maximal bounded subgroups}

\medskip
In this subsection, we suppose that \emph{the valuation $\ph$ of the root datum of $G$
is dense}: for the discrete case, see \S3. We endow $G$ with the bornology $\Bor(\ph)$
defined by $\ph$ (\btref{8.1.7}).

\begin{stmt}{Proposition}{8.2.1}\btanchor{8.2.1}
(i) The map $x\mapsto\hP_x$ is a bijection from the completion $\hbui$ of the building
of $G$ onto the set of the maximal bounded subgroups of $G$.

(ii) Let $x$ and $y$ be two points of $\hbui$. The maximal bounded subgroups $\hP_x$ and
$\hP_y$ are conjugate if and only if $x$ and $y$ belong to the same orbit of $G$ in
$\hbui$.
\end{stmt}

It follows at once from the definition of the bornology of $G$ (\btref{8.1.7}) that the
subgroups $\hP_x$ are bounded. On the other hand, the complete metric space $\hbui$
satisfies condition (CN) of (\btref{3.2.3}): indeed, one shows exactly as in (\btref{3.2.1}) (taking
(\btref{7.4.20}) into account) that relation (\btref{3.2.1})(1) is exact whatever $x,y,z\in\bui$, on
taking $m=(x/2)+(y/2)$ (cf.\ (\btref{7.4.20})(v)), and this remains true by continuity for
$x,y,z\in\hbui$ (cf.\ (\btref{7.5.1})). Consequently, \emph{every bounded subgroup of
$\Isom\hbui$, hence also every bounded subgroup of $G$, possesses a fixed point in
$\hbui$} (\btref{3.2.3}), hence is contained in a subgroup $\hP_x$ with $x\in\hbui$. To complete
the proof of (i), it suffices to show that if $x$ and $y$ are two distinct points of
$\bui$, one has $\hP_x\not\subset\hP_y$; now this follows from (\btref{7.4.34}) and (\btref{7.5.8}).
Assertion (ii) is then obvious.

\numpar{8.2.2}\btanchor{8.2.2}
One will compare (\btref{8.2.1}) with the results of \S3; when $\ph$ is dense, all the
stabilizers of the various points of $\bui$ are maximal bounded subgroups (and there are
still others when $\bui$ is not complete), whereas, when $\ph$ is discrete, only certain
among them are: those which correspond to the vertices of the polysimplicial complex
$\bui$ when $G$ is generated by $H$ and the $U_a$, and, in the general case, those which
correspond to the barycenters of certain facets. Moreover, when $\ph$ is discrete, there
is only a finite number, in general greater than one, of conjugacy classes of maximal
bounded subgroups. When $\ph$ is dense, there are in general infinitely many conjugacy
classes; however, when $\bui$ is complete (condition (MC) of (\btref{7.5.4})) and
$\Gamma_a=\R$ for every $a\in\Ph$, there is only one conjugacy class!

\begin{stmt}{Lemma}{8.2.3}\btanchor{8.2.3}
Let $x$ be a point of $A$, $D$ a vector chamber and $L$ a subgroup of $G$ containing
$B_{x,D}$ and not contained in $\hP_x$. There exists a root $a\in\Ph$ such that
$L\cap U_a\not\subset U_{a,-a(x)}=U_a\cap\hP_x$.
\end{stmt}

In view of (\btref{7.3.4}), one has $L=B_{x,D}.(N\cap L).B_{x,D}$ and there exists $n\in N\cap L$
with $n\notin\hN_x$. Put
\pg{195}
$v=\nu(n).x-x\in V$; since $v\neq0$, there exists $a\in\Ph^{\red}$ such that $a(v)>0$.
Put $b=\vv\nu(n^{-1}).a$. Since $U_{b,-b(x)+}\subset B_{x,D}\subset L$, one has:
\[
U_{a,-a(x+v)+}=n.U_{b,-b(x)+}.n^{-1}\subset L
\]
and since $\ph$ is dense, there exists $h\in\Gamma_a$ such that $-a(x+v)<h<-a(x)$,
whence the lemma, since $U_{a,h}\subset L$ and $U_{a,h}\not\subset U_{a,-a(x)}$.

\begin{stmt}{Proposition}{8.2.4}\btanchor{8.2.4}
Let $x\in A$ and let $D$ be a vector chamber.

(i) Every bounded subgroup of $G$ containing $B_{x,D}$ is contained in $\hP_x$.

(ii) If $\Ph$ is irreducible, every subgroup of $G$ containing $B_{x,D}$ either is
contained in $\hP_x$, or contains the subgroup $G'$ of $G$ generated by $H$ and the
$U_a$ for $a\in\Ph$.
\end{stmt}

Let $L$ be a subgroup containing $B_{x,D}$ and not contained in $\hP_x$. In view of
lemma (\btref{8.2.3}), there exists a root $a\in\Ph^{\red}$ such that
$L\cap U_a\not\subset U_{a,-a(x)}$. Taking into account that $\ph$ is dense, one
deduces that there exist $h,k\in\Gamma_a$ such that $h<k<-a(x)$ and $U_{a,h}\subset L$.
The subgroup $L$ then contains on the one hand the two subgroups $U_{a,h}$ and
$U_{-a,-h}$, on the other hand the two subgroups $U_{a,k}$ and $U_{-a,-k}$.
Consequently, $\nu(N\cap L)$ contains the two reflections $r_{a,k}$ and $r_{a,h}$, hence
the translation $t=r_{a,h}\circ r_{a,k}$ of non-zero vector $(k-h)a\dual$. It follows
that $L$ is not bounded, whence (i).

Now let $c$ be a root not orthogonal to $a$, and put $m=\langle c,a\dual\rangle$; one
has $m\neq0$. For $q\in\Z$, one has
\[
L\supset t^qU_{c,-c(x)+}t^{-q}=U_{c,\,-c(x)-qm(k-h)+}
\]
which shows that $L$ contains all of $U_c$, and the same reasoning shows that $L$
contains $U_{-c}$. Consequently, $\nu(N\cap L)$ contains a translation of the form
$\ell.c\dual$, with $\ell\neq0$. One shows thus step by step that $L$ contains all the
groups $U_c$ for the roots $c\in\Ph$ which can be joined to $a$ by a sequence
$(c_0=a,c_1,\dots,c_s=c)$ of roots such that $c_i$ and $c_{i+1}$ are not orthogonal for
$0\leqslant i<s$, in other words for the roots $c$ of the irreducible component of $\Ph$
which contains $a$. This proves (ii).

Let us recall (\btref{6.2.20}) that the quotient group $G/G'$ is a torsion group every element
of which is annihilated by a power of the connection index of $\Ph$.

\begin{stmtup}{Remark}{8.2.5}\btanchor{8.2.5}
Proposition (\btref{8.2.4}), together with the results recalled in (7.2.7) and in (\btref{1.2.19}),
allows one to determine all the subgroups of $G$ containing $B_{x,D}$.
\end{stmtup}

\bigskip
\pg{196}

\btsecA{9}{9. Descent of a valued root datum}

\medskip

\noindent
We keep the notations of \S\S6, 7 and 8. The aim of this section is to show that, under
certain conditions, the valued root datum of $G$ ``descends'' to a subgroup and
provides a valued root datum of that subgroup.

\medskip

\btsecB{9.1}{9.1. Induced valuations}

\numpar{9.1.1}\btanchor{9.1.1}
Throughout \S9, we denote by $V^\nat$ a vector subspace of $V$. For $a\in V^*$, we
denote by $a^\nat$ the restriction of $a$ to $V^\nat$. For every $b\in(V^\nat)^*$, we
denote by $\Ph_b$ (resp.\ $\Ph^+_b$) the set of the $a\in\Ph$ such that
$a^\nat\in\R.b$ (resp.\ $a^\nat\in\Rp.b$). It is clear that $\Ph_b$ is a rationally
closed root subsystem of $\Ph$. We denote by $Z$ the subgroup generated by $T$ and the
$U_a$ for $a\in\Ph_0$ (i.e.\ for $a\in\Ph$ with $a^\nat=0$). For $b\in(V^\nat)^*$, we
denote:

\begin{itemize}[leftmargin=2em,itemsep=1pt]
\item by $U_b$ the subgroup generated by the $U_a$ for $a\in\Ph^+_b$;
\item by $L_b$ the subgroup generated by $Z$, $U_b$ and $U_{-b}$;
\item by $M_b$ the set of the $x\in L_b$ such that $xZx^{-1}=Z$, $xU_bx^{-1}=U_{-b}$
and $xU_{-b}x^{-1}=U_b$.
\end{itemize}

\begin{stmt}{Lemma}{9.1.2}\btanchor{9.1.2}
For every $b\in(V^\nat)^*$, the pair $(T,(U_a,M_a)_{a\in\Ph_b})$ is a generating root
datum of type $\Ph_b$ in $L_b$. Moreover, $Z$ normalizes $U_b$ and $U_{-b}$; the
subgroups $Z.U_b$ and $Z.U_{-b}$ are parabolic subgroups of $L_b$ and one has
$ZU_b\cap ZU_{-b}=Z$.
\end{stmt}

The first assertion is immediate (cf.\ (\btref{7.6.3})). The others are so as well when $b=0$.
Let us therefore suppose $b\neq0$. That $Z$ normalizes $U_b$ and $U_{-b}$ follows at
once from condition (DR 2) of (\btref{6.1.1}) and from (\btref{6.1.2})(3). Moreover, if one chooses the
vector chamber $D_0$ in such a way that $\Ph^+_b\subset\Ph^+$, which is permissible, and
if one puts $U^+_Z=Z\cap U^+$ and $U^-_Z=Z\cap U^-$, one sees that $TU^+_ZU_b$ and
$TU^-_ZU_{-b}$ are minimal parabolic subgroups of $L_b$, which implies that $ZU_b$ and
$ZU_{-b}$ are parabolic subgroups of $L_b$. Finally, let $u\in U_b$, $v\in U_{-b}$ and
$z\in Z$ be such that $v=zu$. In view of (\btref{6.1.15}), there exists one and only one element
$n\in Z\cap N$ such that $z=v'nu'$, with $v'\in U^-_Z$ and $u'\in U^+_Z$. One then has
$v=v'nu'u\in U^-nU^+\cap U^-$, whence $n=1$ by (\btref{6.1.15}) and $v=v'$ by (DR 6). One
therefore has $v\in U_{-b}\cap U^-_Z=\{1\}$ (\btref{6.1.6}), which proves the last assertion.

\begin{stmt}{Lemma}{9.1.3}\btanchor{9.1.3}
Let $b\in(V^\nat)^*$; if $M_b\neq\emptyset$, there exists an element $n\in N\cap L_b$
such that $M_b=Zn$.
\end{stmt}

This is obvious if $b=0$. Suppose $b\neq0$ and let us take up the notations of (\btref{9.1.2}).
In view of $(6.1.15)\ c)$, there exists $n\in N\cap L_b$ such that
\[
n^{-1}TU^+_ZU_bn=TU^-_ZU_{-b}\qquad\text{and}\qquad n^{-1}TU^-_ZU_{-b}n=TU^+_ZU_b .
\]
\pg{197}
If $m\in M_b$, the subgroups $n^{-1}ZU_bn$ and $m^{-1}ZU_bm=ZU_{-b}$ are two conjugate
parabolic subgroups both containing the same minimal parabolic subgroup
$TU^-_ZU_{-b}$. They are therefore equal, and $mn^{-1}$ normalizes $ZU_b$, hence belongs
to $ZU_b$. One sees in the same way that $mn^{-1}$ belongs to $ZU_{-b}$, hence finally
to $Z$ by (\btref{9.1.2}); whence the lemma.

\numpar{9.1.4}\btanchor{9.1.4}
For $b\in(V^\nat)^*$, $b\neq0$, and $k\in\R$, we denote by $U_{b,k}$ the subgroup of $G$
generated by the $U_{a,rk}$ for $a\in\Ph^+_b$, $r\in\Rp$ and $a^\nat=rb$. With the
notations of 6.4, $U_{b,k}$ is the subgroup $U_f$ associated with the function
$f=f(b,k)$ on $\Ph$ defined by
\begin{equation}\tag{1}
\left\{
\begin{aligned}
f(b,k)(a)&=rk &&\text{if }a=rb\text{ with }r\in\Rp\\
f(b,k)(a)&=+\infty &&\text{if }a\notin\Ph^+_b .
\end{aligned}
\right.
\end{equation}
One verifies at once that $f(b,k)$ is a \emph{concave} function.

\begin{stmt}{Lemma}{9.1.5}\btanchor{9.1.5}
Let $b\in(V^\nat)^*$, $b\neq0$, and $k\in\R$. For $r\in\Rp$, let $\beta(r)$ be the set
of the $a\in\Ph^{\red}$ such that $a^\nat=rb$. The product map (for an arbitrary fixed
order) is a bijection from
$\prod_{r\in\Rp}\bigl(\prod_{a\in\beta(r)}U_a\bigr)$
$\bigl($resp.\
$\prod_{r\in\Rp}\bigl(\prod_{a\in\beta(r)}U_{a,rk}\bigr)\bigr)$
onto $U_b$ (resp.\ $U_{b,k}$).
\end{stmt}

This follows at once from (\btref{6.1.6}) and (\btref{6.4.9}).

\begin{stmt}{Lemma}{9.1.6}\btanchor{9.1.6}
Let $b\in(V^\nat)^*$, $b\neq0$. For $u\in U_b$, put
\[
\ph_b(u)=\sup\{k\in\R\cup\{\infty\}\mid u\in U_{b,k}\}.
\]
The inverse image $\ph_b^{-1}([k,+\infty])$ is the subgroup $U_{b,k}$. For every
$s\in\Rp$, one has $U_{sb}=U_b$ and $\ph_{sb}=s\ph_b$.
\end{stmt}

This is evident.

\numpar{9.1.7}\btanchor{9.1.7}
Let us take up the notations of \S7: let $\bui$ be the building of $G$, $A$ the
canonical apartment, identified with the set $\ph+V$ of the valuations equipollent to
$\ph$. Put $A_\nat=\ph+V^\nat$.

\begin{stmt}{Lemma}{9.1.8}\btanchor{9.1.8}
Let $b\in(V^\nat)^*$, $b\neq0$, $k\in\R$ and $u\in U_b$. The two following conditions
are equivalent:

(i) $\ph_b(u)=k$;

(ii) the set of the fixed points of $u$ in $A_\nat$ is the half-space $\alpha_{b,k}$ of
$A_\nat$ formed by the $x\in A_\nat$ such that $b(x-\ph)+k\geqslant0$.

Moreover, if $k\in\ph_b(U_b-\{1\})$, the half-space $\alpha_{b,k}$ is the intersection
with $A_\nat$ of at least one affine root of $A$.
\end{stmt}

One may suppose that $\Ph^+_b$ is not empty.

Let us denote by $D_b$ the intersection of the closed half-spaces of $V$ of the form
$\{x\in V\mid a(x)\geqslant0\}$ for $a\in\Ph^+_b$. This is a convex cone with non-empty
interior, containing a Weyl chamber $D$ of $\Ph$, and whose intersection with $V^\nat$
is the closed half-space $R_b=\{x\in V\mid b(x)\geqslant0\}$.

Let $a\in\Ph$ and $h\in\R$. Every element of $U_{a,h}$ fixes the points of the half-space
\pg{198}
$\alpha_{a,h}=\{x\in A\mid a(x-\ph)+h\geqslant0\}$ of $A$ (\btref{7.4.5}); if $a^\nat=rb$, with
$r\in\Rp$, one has $\alpha_{a,rk}\cap A_\nat=\alpha_{b,k}$. Consequently, \emph{every
element of $U_{b,k}$ fixes all the points of $\alpha_{b,k}$.}

On the other hand, if $a\in\Ph^+_b$ and $h\in\R$, the half-space $\alpha_{a,h}$ contains
a translate of $D_b$. Since the intersection of a finite number of translates of the
convex cone $D_b$ still contains a translate of $D_b$, and since the set $F(u)$ of the
fixed points of $u$ in $A$ is convex, one sees that $F(u)+D_b=F(u)$. Let $x\in A_\nat$ be
such that $u.x=x$; put $h=-b(x-\ph)$. Then $u$ fixes all the points of $\Om=x+D_b$ and
belongs to the subgroup $\hP_\Om\cap U^+_D$, which is generated by the $U_\alpha$ for
$\alpha\supset\Om$ and $\alpha\in\Sigma$ (7.1.8); now, the only affine roots containing
$\Om$ are the $\alpha_{a,\ell}$ for $a\in\Ph^+_b$, $s\in\Rp$, $a^\nat=sb$ and
$\ell\geqslant sh$. One therefore has $u\in U_{b,h}$ and $\ph_b(u)\geqslant h$. This,
together with the preceding result, shows the equivalence of (i) and (ii).

Finally, if $\ph_b(u)=k$, there exist $s\in\Rp$ and $a\in\Ph^+_b$ such that $a^\nat=sb$
and $sk\in\ph_a(U_a)$, whence the last assertion.

\begin{stmt}{Definition}{9.1.9}\btanchor{9.1.9}
Let $G^\nat$ be a subgroup of $G$, $\Ph^\nat$ a root system in
$(V^\nat)^*$\footnote{\upshape One will take care that, contrary to what the notation
might lead one to believe, one does \emph{not} assume that
$\Ph^\nat=\{a^\nat\mid a\in\Ph\}$.} and $(T^\nat,(U^\nat_b)_{b\in\Ph^\nat})$ a
generating root datum of type $\Ph^\nat$ in $G^\nat$. One says that this root datum is
\upshape compatible\itshape\ with that of $G$ if one has
\[
\text{\upshape(DDR 1)}\qquad U^\nat_b\subset U_b\quad\text{for every }b\in\Ph^\nat .
\]
\end{stmt}

Throughout the rest of 9.1, we are given the vector subspace $V^\nat$ of $V$, a subgroup
$G^\nat$ of $G$, a root system $\Ph^\nat$ in $(V^\nat)^*$ and a generating root datum
$(T^\nat,(U^\nat_b))$ of type $\Ph^\nat$ in $G^\nat$, compatible with that of $G$. The
objects which \S6 allows one to associate with this root datum will be denoted by the
same letter as in \S6, but affected with the exponent $\nat$: for example, one denotes
by $\vv W^\nat$ the Weyl group of $\Ph^\nat$ and one denotes by $D^\nat$ a vector
chamber of $\Ph^\nat$ in $V^\nat$. We assume, which is permissible, that the vector
chamber $D$ is such that $\overline D\cap V^\nat\subset\overline D^\nat$. The notations
$M^\nat_b$, $N^\nat$, $\Ph^{\nat+}$, $U^{\nat+}$, etc.\ then explain themselves.

For $b\in\Ph^\nat$, we denote by $\ph^\nat_b$ the restriction of $\ph_b$ to $U^\nat_b$
and we put $\ph^\nat=(\ph^\nat_b)_{b\in\Ph^\nat}$.

\begin{stmt}{Proposition}{9.1.10}\btanchor{9.1.10}
The family $\ph^\nat$ satisfies conditions \upshape(V 1)\itshape\ and
\upshape(V 4)\itshape\ of (\btref{6.2.1}). It also satisfies \upshape(V 3)\itshape\ except
perhaps when the two roots $a$ and $b$ of $\Ph^\nat$ are equal and $2a\in\Ph^\nat$.
However, this last restriction can be lifted if one supposes that the following
condition is fulfilled:

\upshape(DDR 2)\itshape\ for every $a\in\Ph^\nat$ such that $2a\in\Ph^\nat$, one has
$\Card\{c^\nat\mid c\in\Ph^+_a\}\leqslant2$.
\end{stmt}

That (V 1) and (V 4) are satisfied follows from lemma (\btref{9.1.6}). Now let
$a,b\in\Ph^\nat$. Let us put
\begin{align*}
\Psi^\nat&=\{pa+qb\mid p,q\in\N^*\}\cap\Ph^\nat\\
\Psi&=\{c\in\Ph\mid c^\nat\in\Rp a+\Rp b\}
\end{align*}
\pg{199}
Let $Y$ be the group generated by the $U_c$ for $c\in\Psi$ and let $Y^\nat$ be the group
generated by the $U^\nat_d$ for $d\in\Psi^\nat$. In view of (\btref{6.1.6}), the following
product maps are bijections:
\begin{align*}
\theta&:\Bigl(\prod_{d\in\Psi^{\nat\,\red}}U_d\Bigr)\times
\prod_{c\in\Psi^{\red},\ c^\nat\notin\Rp.\Psi^\nat}U_c\longrightarrow Y\\
\theta&:\Bigl(\prod_{d\in\Psi^{\nat\,\red}}U^\nat_d\Bigr)\times\{1\}\longrightarrow
Y^\nat .
\end{align*}

Let us then apply (\btref{6.4.29}) to the two concave functions $f=f(a,k)$ and $g=f(b,\ell)$
(with $k,\ell\in\R$): one has $(U_f,U_g)\subset U_h$, where the function
$h:\Ph\to\wt\R$ is defined by
\begin{equation}\tag{1}
h(c)=\inf\Bigl\{\sum_i f(a_i)+\sum_j g(b_j)\Bigr\}
\end{equation}
the infimum being taken over the set of the decompositions
$c=\sum_{i\in I}a_i+\sum_{j\in J}b_j$, where $a_i,b_j\in\Ph$ and $\Card I\geqslant1$,
$\Card J\geqslant1$.

Let us suppose $a$ and $b$ non-proportional. One then verifies immediately that
$h(c)=+\infty$ if $c\notin\Psi$ and that $h(c)=rk+s\ell$ if $c^\nat=ra+sb$. In view of
(\btref{6.4.9}), the restriction of $\theta$ to
\[
\Bigl(\prod_{d\in\Psi^{\nat\,\red}}U_{d,h(d)}\Bigr)\times
\prod_{c\in\Psi^{\red},\ c^\nat\notin\Rp.\Psi^\nat}U_{c,h(c)}
\]
(one has put $h(ra+sb)=rk+s\ell$) is a bijection onto $U_h$. Now, by (DR 2), one has
$(U^\nat_a,U^\nat_b)\subset Y^\nat$. Consequently,
$(U^\nat_{a,k},U^\nat_{b,\ell})\subset Y^\nat\cap U_h$ and one deduces condition (V 3)
for the roots $a$ and $b$ by taking the inverse image of $Y^\nat\cap U_h$ under
$\theta$.

If $a$ and $b$ are proportional, but $a+b\notin\Ph^\nat$, there is nothing to prove
since $U^\nat_a$ and $U^\nat_b$ commute.

There remains to examine the case where $b=a$, with $2a\in\Ph^\nat$, supposing that
$\Card\{c^\nat\mid c\in\Psi\}\leqslant2$. One then sees that $U_a$ is commutative, and
that there is nothing to prove, except when there exists $s\in\Rp$ such that
$\{c^\nat\mid c\in\Psi\}=\{sa,2sa\}$. But in this last case one verifies easily that the
function $h$ given by (1) is infinite except on the roots $c\in\Psi$ of the form
$c=c_1+c_2$ with $c_1,c_2\in\Psi$ and $c_1^\nat=c_2^\nat=sa$, in which case
$h(c)=s(k+\ell)$. One deduces that
$(U_{a,k},U_{a,\ell})\subset U_{a,(1/2)(k+\ell)}$, whence
\[
(U_{a,k},U_{a,\ell})\subset U_{a,(1/2)(k+\ell)}\cap U_{2a}=U_{2a,k+\ell}
\]
which completes the proof.

\begin{stmt}{Definition}{9.1.11}\btanchor{9.1.11}
If the family $\ph^\nat$ is a valuation of the root datum
$(T^\nat,(U^\nat_b))$ of $G^\nat$, one says that $\ph$ \upshape descends\itshape\ to
$G^\nat$ and that $\ph^\nat$ is the valuation \upshape induced\itshape\ by $\ph$.
\end{stmt}

If $\ph$ descends to $G^\nat$, we extend to the objects associated with the valued
root datum of $G^\nat$ the notational conventions of (\btref{9.1.9}).

\begin{stmtup}{Remarks}{9.1.12}\btanchor{9.1.12}
\textit{a)} Under the hypotheses of (\btref{9.1.10}) (including (DDR 2)), the valuation $\ph$
descends if and only if the family $\ph^\nat$ satisfies conditions (V 0), (V 2) and
(V 5) of (\btref{6.2.1}). One will note that these are all relative ``to rank $1$'', in the
sense that they involve simultaneously only the integer multiples of one and the
\pg{200}
same root. We shall give later (\btref{9.2}) conditions for them to be satisfied.

\textit{b)} Let $V^\sharp$ be a vector subspace of $V^\nat$ and let $G^\sharp$ be a
subgroup of $G^\nat$, endowed with a root datum compatible with that of $G^\nat$. Then
the root datum of $G^\sharp$ is compatible with that of $G$. If $\ph$ descends to
$G^\nat$, one verifies easily that the valuation $\ph^\nat$ descends to $G^\sharp$ if
and only if $\ph$ descends to $G^\sharp$.
\end{stmtup}

\numpar{9.1.13}\btanchor{9.1.13}
For $b\in(V^\nat)^*$, $b\neq0$, and $k\in\R$, let us denote by $U^e_b$ the subgroup of
$G$ generated by the $U_a$ for $a\in\Ph$ and $a^\nat\in\N^*.b$, and by $U^e_{b,k}$ the
one generated by the $U_{a,nk}$ for $a\in\Ph$, $n\in\N^*$ and $a^\nat=nb$. Using
(\btref{6.1.6}) and (\btref{6.4.9}), one sees easily that
\[
U^e_{b,k}=U_{b,k}\cap U^e_b .
\]

In practice, the root datum of $G^\nat$ will most of the time satisfy the following
condition, stronger than (DDR 1):
\[
\text{\upshape(DDR 1 bis)}\qquad U^\nat_b\subset U^e_b\quad\text{for every }
b\in\Ph^\nat .
\]

\numpar{9.1.14}\btanchor{9.1.14}
Let us suppose $\ph$ extendable and let $\ph_0:H\to\R_+\cup\{\infty\}$ be an extension
of $\ph$ (\btref{6.4.38}). For $u\in Z$ (\btref{9.1.1}), let $\ph_Z(u)$ be the supremum, finite or
$+\infty$, of the set of the $k\in\R$ such that $u$ belongs to the subgroup generated by
$U_{0,k}=\ph_0^{-1}([k,+\infty])$ and the $U_{a,k}$ for $a\in\Ph_0$.

Let us suppose also that $\ph$ descends to $G^\nat$. One may then consider the subgroup
$H^\nat=\nu^{\nat\,-1}(1)$ and put, for $u\in H^\nat$:
\[
\ph^\nat_0(u)=\sup\{0,\ph_Z(u)\}.
\]

\begin{stmt}{Proposition}{9.1.15}\btanchor{9.1.15}
Let us keep the hypotheses and notations of (\btref{9.1.14}). Suppose moreover that
$\Ph^\nat=\{a^\nat\mid a\in\Ph\}$ and that there exists a constant $t\geqslant0$ such
that, if $b$ and $2b$ belong to $\Ph^\nat$, one has
\[
\text{\upshape(DP)}\qquad
(U_{b,k+t}.U^e_{2b,2k-h})\cap U^\nat_b=U^\nat_{b,k}.U^\nat_{2b,2k-h}
\]
whatever $k\in\R$ and $h\in\R_+$. Then $\ph^\nat_0$ is an extension of $\ph^\nat$.
\end{stmt}

For $\ell\geqslant0$, put $U^\nat_{0,\ell}=\ph^{\nat\,-1}_0([\ell,+\infty])$. We must
show (\btref{6.4.38}) the following assertions:

\medskip
\noindent
\textit{a)} \emph{One has $(U^\nat_{0,k},U^\nat_{0,\ell})\subset U^\nat_{0,k+\ell}$ for
$k,\ell\geqslant0$.} For this, it suffices to apply (\btref{6.4.44}) to the functions
$f,g:\Ph\cup\{0\}\to\R\cup\{\infty\}$ such that $f(a)=k$ and $g(a)=\ell$ if
$a\in\Ph_0\cup\{0\}$, and $f(a)=g(a)=\infty$ otherwise.

\medskip
\noindent
\textit{b)} \emph{For $b\in\Ph^\nat$ and $k,\ell\in\R$ such that $k+\ell>0$, the
$H^\nat$-components (\btref{6.3.8}) of the commutators $(u,u')$ with $u\in U^\nat_{b,k}$ and
$u'\in U^\nat_{-b,\ell}$ belong to $U^\nat_{0,k+\ell}$.} For this one applies (\btref{6.4.44})
to the functions $f,g:\Ph\cup\{0\}\to\R\cup\{\infty\}$ defined by $f(a)=k$ if
$a^\nat=b$, $f(a)=2k$ if $a^\nat=2b$, $f(a)=\infty$ otherwise, and $g(a)=\ell$ if
$a^\nat=-b$, $g(a)=2\ell$ if $a^\nat=-2b$, $g(a)=\infty$ otherwise. One draws from this
that
\[
(U_{b,k},U_{-b,\ell})\subset U^e_{2b,3k+\ell}.U_{b,2k+\ell}.Y.U_{-b,2\ell+k}. U^e_{-2b,3\ell+k}
\]
\pg{201}
where $Y$ is the subgroup generated by $U_{0,k+\ell}$ and the $U_{a,k+\ell}$ for
$a\in\Ph_0$, whence the result. One sees in the same way that the $H^\nat$-components of
the commutators $(u,u')$ and $(u',u)$ with $u\in U^\nat_{b,k/2}$ and
$u'\in U^\nat_{-2b,\ell}$ belong to $U^\nat_{0,k+\ell}$. This shows that
$H^\nat_{[k+\ell]}\subset U^\nat_{0,k+\ell}$.

\medskip
\noindent
\textit{c)} \emph{For $b\in\Ph$, $k\in\R$ and $\ell\in\Rp$, one has
$(U^\nat_{b,k},U^\nat_{0,\ell})\subset U^\nat_{b,k+\ell}.U^\nat_{2b,2k+\ell}$;} in other
words $U^\nat_{0,\ell}\subset H^\nat_{(\ell)}$. A reasoning of the same type as the
preceding ones shows that
\[
(U_{b,k},U^\nat_{0,\ell})\subset U_{b,k+\ell}.U^e_{2b,2k+\ell}\qquad\text{for all } k\in\R
\]
whence the result if $2b\notin\Ph^\nat$. If $2b\in\Ph^\nat$, one uses (DP) and what
precedes to show, taking into account that $(U_b,U^e_{2b})=\{1\}$, that
\begin{align*}
(U^\nat_{b,k},U^\nat_{0,\ell})&\subset(U_{b,k+t},U^\nat_{0,\ell}).U^e_{2b,2k+\ell}\cap
U^\nat_b\\
&\subset(U_{b,k+\ell+t}.U^e_{2b,2k+\ell})\cap U^\nat_b
=U^\nat_{b,k+\ell}.U^\nat_{2b,2k+\ell}
\end{align*}
which completes the proof.

\begin{stmtup}{Remark}{9.1.16}\btanchor{9.1.16}
One cannot always take $t=0$ in (DP). Let for example $K$ be a commutative field,
endowed with an involution $\sigma$ and with a non-improper valuation $\omega$ invariant
under $\sigma$. Suppose $K$ of residue characteristic $2$ and let $\lambda\in K$ be such
that $\lambda+\lambda^\sigma=1$. With the notations that we shall employ in \S10
(cf.\ (\btref{10.1.1})), let us consider the hermitian form on $K^3$ defined by
$q(xe_1+ze_0+ye_{-1})=\lambda^\sigma zz+x^\sigma y+K_{\sigma,1}$ and let us endow
$G^\nat=SU_3(q)$ with the root datum defined in (10.1.6). Let us endow $G=SL_3(K)$ with
the root datum defined in (\btref{10.2.2}) and with the valuation $\ph$ defined in (\btref{10.2.3})
(relative to the basis $(e_1,e_0,e_{-1})$ of $K^3$). One verifies easily that these root
data are compatible and that $\ph$ descends to $G^\nat$ and induces there the valuation
defined in (\btref{10.1.13}). Moreover, we shall see that $\ph$ and $\ph^\nat$ are extendable
((\btref{10.1.24}) and (\btref{10.2.4})). But it is easy to see that condition (DP) is satisfied if and
only if one has $t=-(1/2)\sup\{\omega(\mu)\mid\mu\in K,\ \mu+\mu^\sigma=1\}$. If $K$ and
the field of invariants of $\sigma$ have the same residue field, or if the residue
extension is inseparable, one has $t>0$.
\end{stmtup}

\begin{stmt}{Proposition}{9.1.17}\btanchor{9.1.17}
Suppose that $\ph$ descends to $G^\nat$ and that there exists a subgroup $S^\nat$ of
$T^\nat$, normal in $N^\nat$, such that $G^\nat$ is generated by $S^\nat$ and the
$U^\nat_a$ for $a\in\Ph^\nat$, that $\nu^\nat(S^\nat)$ generates the vector space
$V^\nat$ and that $S^\nat.A_\nat=A_\nat$. Then the scalar product induced on $V^\nat$ by
that of $V$ is invariant under the Weyl group of $\Ph^\nat$. If one endows $\bui^\nat$
with the corresponding distance, there exists one and only one isometric map
$j:\bui^\nat\to\bui$ such that $j(\ph^\nat)=\ph$ and such that $j(g.x)=g.j(x)$ for all
$g\in G^\nat$ and all $x\in\bui^\nat$. Moreover, $j(\bui^\nat)$ is convex,
$j(A^\nat)=A_\nat$ and $A^\nat=j^{-1}(A)$.
\end{stmt}

Let us denote by $j$ the bijection of $A^\nat=\ph^\nat+V^\nat$ onto $A_\nat$ defined by
$j(\ph^\nat+v)=\ph+v$. Let $s\in S^\nat$, $\ps\in A_\nat$ and $v\in V^\nat$ be such that
$s.\ps=\ps+v$. It follows at once from the definitions that
$(\ps+v)^\nat=\ps^\nat+v$. On the other hand, the inner automorphism defined by $s$
preserves the root datum $(T^\nat,(U^\nat_b))$ and transforms the root datum
$(T,(U_a))$ and the apartment $A$ into $(sTs^{-1},(sU_as^{-1}))$ and $s.A$. Since
$s.A_\nat=A_\nat$, one deduces from (\btref{9.1.8}) that the valuation $\ps^\nat+v$ of
$(T^\nat,(U^\nat_b))$ is equal to the valuation obtained by descent from the valuation
$s.\ps$ of $(sTs^{-1},(sU_as^{-1}))$, which is none other than $s.\ps^\nat$ (transport
of structure!). One therefore has:
\begin{equation}\tag{1}
j(s.x)=s.j(x)\qquad\text{for }x\in A^\nat\text{ and }s\in S^\nat .
\end{equation}

\pg{202}
Let $b\in\Ph^\nat$, $k\in\R$, $m\in M^\nat_{b,k}$ (\btref{6.2.2}) and $y\in A^\nat$ be such that
$m.y=y$. One has:
\[
b(j(y)-\ph)+k=b(y-\ph^\nat)+k=0 .
\]
Since $m$ belongs to the group generated by $U_{b,k}$ and $U_{-b,-k}$, this implies that
$m.j(y)=j(y)$. For $s\in S^\nat$, one then has, in view of (1):
\[
j(ms.y)=j(msm^{-1}.y)=msm^{-1}.j(y)=ms.j(y)=m.j(s.y)
\]
whence one deduces at once that $j(m.x)=m.j(x)$ for all $x\in A^\nat$. The group
$N^\nat$ being generated by $S^\nat$ and the $M^\nat_{b,k}$, one concludes that
$j(n.x)=n.j(x)$ for $n\in N^\nat$ and $x\in A^\nat$. This implies that the scalar
product induced on $V^\nat$ by that of $V$ is invariant under $\vv W^\nat$. If one
endows $A^\nat$ with the corresponding distance, $j$ is isometric.

Let still $x\in A^\nat$. One knows that the stabilizer of $x$ in $G^\nat$ is generated
by its stabilizer $\hN^\nat_x$ in $N^\nat$ and the $U^\nat_{a,-a(x-\ph^\nat)}$ for
$a\in\Ph^\nat$. But the very definition of $\ph^\nat$ shows that these last are contained
in the stabilizer of $j(x)$, and what precedes shows that $\hN^\nat_x$ is too. One
therefore has $\hP^\nat_x\subset\hP_{j(x)}$ and $j$ extends in one and only one way to a
map, still denoted $j$, from $\bui^\nat$ into $\bui$ such that $j(g.x)=g.j(x)$ for
$g\in G^\nat$ and $x\in\bui^\nat$. Since the restriction of $j$ to $A^\nat$ is isometric
and since any two points of $\bui^\nat$ are transformed into two points of $A^\nat$ by
an operation of $G^\nat$, one sees that $j$ is isometric and that $j(\bui^\nat)$ is
convex. On the other hand, the uniqueness of $j$ is immediate.

There remains finally to show that $j^{-1}(A)=A^\nat$. Let $x\in\bui^\nat$,
$x\notin A^\nat$, be such that $j(x)\in A$. Let $C$ be a chamber of $A^\nat$ and let
$A'$ be an apartment of $\bui^\nat$ containing $C$ and $x$. Since $j(C)\subset A_\nat$,
there exists a non-empty open subset $\Om$ of $A'$ such that $j(\Om)$ is a non-empty
open subset of $A_\nat$. The restriction of $j$ to the convex hull $\Om'$ of
$\Om\cup\{x\}$ is then an isometry onto a subset with non-empty interior of the affine
subspace $L$ of $A$ generated by $A_\nat\cup\{j(x)\}$, whence one draws
$\dim A'\geqslant\dim\Om'=\dim L>\dim A_\nat$, which is absurd.

\begin{stmtup}{Remark}{9.1.18}\btanchor{9.1.18}
Conversely, if $\ph_1$ is a valuation of the root datum of $G^\nat$ and $j$ an isometry
of the corresponding building $\bui_1$ into $\bui$, such that $j(g.x)=g.j(x)$ for
$x\in\bui_1$ and $g\in G^\nat$, that $j(\ph_1)=\ph$ and that $j(A_1)\subset A$ (where
$A_1=\ph_1+V^\nat$), then one can show that $\ph$ descends, that $\ph^\nat=\ph_1$ and
that the hypotheses of (\btref{9.1.17}) are satisfied with $S^\nat=T^\nat$.
\end{stmtup}

\begin{stmtup}{Examples}{9.1.19}\btanchor{9.1.19}
\textit{a)} If $V^\nat=V$, $T^\nat\subset T$ and $N^\nat\subset N$, one verifies easily
that (DDR 2) is satisfied and that $\ph^\nat$ satisfies conditions (V 2) and (V 5) of
(\btref{6.2.1}). Suppose moreover that $\Card\ph^\nat_b(U^\nat_b)\geqslant3$ for every
$b\in\Ph^\nat$: then \emph{$\ph$ descends to $G^\nat$ and the conditions of (\btref{9.1.17}) are
satisfied} (with $S^\nat=T^\nat$).

For example, let $K$ be a commutative field endowed with a valuation $\omega$, and let
$L$ be a subfield of $K$ such that the restriction of $\omega$ to $L$ is non-improper.
Take for $G$ (resp.\ $G^\nat$) the group of the rational points over $K$ (resp.\ $L$) of
a connected semi-simple algebraic group $\Gr$ defined and split over $L$, endowed with
the root datum defined in $(6.1.3)\ b)$. Then (DDR 1 bis) and (DDR 2) are satisfied. If
$\ph$ (resp.\ $\ph_1$) is the
\pg{203}
valuation of $G$ (resp.\ $G^\nat$) defined in $(6.2.3)\ b)$, one sees that $\ph$
descends to $G^\nat$ and that $\ph^\nat=\ph_1$.

Another example is that where $K$ is a commutative field complete for a non-improper
discrete valuation, with perfect residue field, where $G$ is the group of the rational
points over $K$ of a connected semi-simple algebraic group $\Gr$ defined over $K$,
endowed with the root datum defined in $(6.1.3)\ c)$ and with the valuation $\ph$ that
we shall construct later, and where $G^\nat$ is the group of the rational points over
$K$ of a semi-simple subgroup $\Gr^\nat$ of $\Gr$, defined over $K$ and of relative rank
over $K$ equal to that of $\Gr$ (for example one of the maximal split subgroups
constructed in [\btbib{3}], \S7). Then (DDR 1 bis) and (DDR 2) are satisfied ([\btbib{3}], 3.12) and
$\Card\ph^\nat_b(U^\nat_b)=\infty$ for every $b\in\Ph^\nat$. Hence $\ph$ descends to
$G^\nat$ and one may apply (\btref{9.1.17}).

\medskip
\textit{b)} Under the conditions of the first paragraph of \textit{a)}, and with the
notations of (\btref{9.1.17}), one has $A_\nat=A$ and one may identify $A^\nat$ with $A$ thanks
to the isometry $j$. The set of the walls (resp.\ affine roots) of $A^\nat$ is then
contained in that of $A$ (\btref{9.1.8}), but is not necessarily equal to it. For example, under
the hypotheses of the second paragraph of \textit{a)}, the affine roots of $A$ are the
closed half-spaces $\alpha_{a,k}$ for $a\in\Ph$ and $k\in\omega(K^*)$, whereas those of
$A^\nat$ are the $\alpha_{a,k}$ for $a\in\Ph$ and $k\in\omega(L^*)$.

One can also give examples where \emph{no special point of $A$ is a special point of
$A^\nat$}.

\medskip
\textit{c)} Let $\Gr$ be a Chevalley scheme (over $\Z$), $\Tr$ a maximal split torus of
$\Gr$ and $\Gr^\nat$ a closed group subscheme of $\Gr$ which is also a Chevalley scheme
and such that $\Tr^\nat=\Tr\cap\Gr^\nat$ is a maximal split torus of $\Gr^\nat$. Put
$V=\Hom(\Mult,\Tr)\otimes\R$, $V^\nat=\Hom(\Mult,\Tr^\nat)\otimes\R$ and let
$\Ph\subset V^*$ and $\Ph^\nat\subset(V^\nat)^*$ be the root systems of $\Gr$ and
$\Gr^\nat$. Let $K$ be a field endowed with a non-improper valuation $\omega$. Let
$G=\Gr(K)$ and $G^\nat=\Gr^\nat(K)$, endowed with the root data of types $\Ph$ and
$\Ph^\nat$ defined as in $(6.1.3)\ b)$ (with $T=\Tr(K)$ and $T^\nat=\Tr^\nat(K)$); it is
easy to see that the compatibility condition (DDR 1) of (\btref{9.1.9}) is satisfied. Finally,
let $\ph$ and $\ph_1$ be the valuations of the root data in question defined as in
$(6.2.3)\ b)$. Let us denote by $\bui$ (resp.\ $\bui_1$) the building of $G$
(resp.\ $G^\nat$) relative to $\ph$ (resp.\ $\ph_1$). Then:

\emph{The valuation $\ph$ descends to $G^\nat$ and induces there a valuation $\ph^\nat$
equivalent to $\ph_1$ (\btref{6.2.5}). Moreover, the hypotheses of (\btref{9.1.17}) are satisfied and
there exists (\btref{7.4.3}) one and only one map $\gamma:\bui_1\to\bui$, compatible with the
actions of $G^\nat$ on $\bui_1$ and $\bui$, and such that $\gamma(\ph_1)=\ph$ and such
that the restriction of $\gamma$ to $A_1=\ph_1+V^\nat$ is an affine injection of $A_1$
into $A$. If $G$ is almost simple, $\gamma$ multiplies distances by a constant.}

We shall only sketch the proof of these assertions.

It is first of all easy to reduce to the case where $G^\nat$ is almost simple, which we
shall suppose from now on.

Let $\overline K$ be a field containing $K$, endowed with a valuation
$\overline\omega$ extending $\omega$ and whose value group is $\R$. Applying remark
$(9.1.12)\ b)$ to the inclusions
\begin{align*}
&\Gr(\overline K)\supset\Gr(K)\supset\Gr^\nat(K)\\
&\Gr(\overline K)\supset\Gr^\nat(\overline K)\supset\Gr^\nat(K)
\end{align*}
\pg{204}
and taking into account example \textit{a)} above, one sees that one may, without loss
of generality, substitute $\overline K$ for $K$. We shall therefore suppose from now on
that $\omega(K^*)=\R$.

Let $\mathcal{O}$ be the ring of integers of $K$. The stabilizer $\hP_\ph$ of $\ph$
(considered as a point of $\bui$) in $G$ is none other than $\Gr(\mathcal{O})$.
Similarly, $\hP_{\ph_1}=\Gr^\nat(\mathcal{O})$. Consequently
$\hP_{\ph_1}=\hP_\ph\cap G^\nat$ and, since $G^\nat$ is transitive on $\bui_1$, there
exists one and only one map $\gamma:\bui_1\to\bui$ compatible with the action of
$G^\nat$, that is to say such that
\begin{equation}\tag{1}
\gamma(gx)=g.\gamma(x)\qquad(x\in\bui_1;\ g\in G^\nat).
\end{equation}
Let us take $g$ in $T^\nat$. In view of $(6.1.3)\ b)$, one has
$\nu(g)=\nu^\nat(g)\in V^\nat$ and $\nu^\nat(T^\nat)=V^\nat$, whence
\[
\gamma(\ph_1+v)=\ph+v\qquad(v\in V^\nat).
\]
In particular, $\gamma$ induces an affine injection of $A_1$ into $A$. Let $N^\nat$ be
the normalizer of $T^\nat$ in $G^\nat$. The metric induced on $A_1$
(resp.\ $\gamma(A_1)$) by the metric of $\bui_1$ (resp.\ of $\bui$) is, up to a factor,
the only Euclidean metric invariant under $N^\nat$. In view of (1), it follows that
$\gamma|_{A_1}$ multiplies distances by a constant. Again by virtue of (1), and taking
(\btref{7.4.18}) into account, the same holds for $\gamma$. A reference to remark (\btref{9.1.18})
completes the proof.
\end{stmtup}

\medskip

\btsecB{9.2}{9.2. A descent theorem}

\numpar{9.2.1}\btanchor{9.2.1}
Let us take up for a moment the hypotheses and notations of (\btref{9.1.17}). Put
$\bui_\nat=j(\bui^\nat)$, $A_\nat=j(A^\nat)$ and let $F$ be a facet of $\bui$ meeting
$\bui_\nat$ of the largest possible dimension. Since $\bui_\nat=G^\nat.A_\nat$, one may
suppose that $F$ meets $A_\nat$, hence is contained in $A$. One then sees immediately
that the following conditions are satisfied:

\medskip
\noindent
(DI 1) \emph{$\bui_\nat$ is a convex subset of $\bui$ stable under $G^\nat$};

\noindent
(DI 2) \emph{$A_\nat=\bui_\nat\cap A$ is an affine subspace of $A$};

\noindent
(DI 3) \emph{there exists a facet $F$ of $\bui$ meeting $A_\nat$ (hence contained in
$A$) which is contained in the closure of no other facet $F'\neq F$ of $\bui$ meeting
$\bui_\nat$.}

\medskip
\noindent
(DV 1) \emph{$\ph\in A_\nat$.}

\numpar{9.2.2}\btanchor{9.2.2}
Throughout 9.2, we are given a subgroup $G^\nat$ of $G$ and a subset $\bui_\nat$ of the
building $\bui$ of $G$ satisfying conditions (DI 1), (DI 2) and (DI 3) above. The aim of
this subsection is to show that these conditions, joined to conditions (DDR 1) (\btref{9.1.9}),
(DDR 2) (\btref{9.1.10}), (DV 1) (\btref{9.2.1}) and to two other conditions (DDR 3) (\btref{9.2.8}) and
(DV 2) (\btref{9.2.8}), imply that the valuation $\ph$ descends to $G^\nat$ and that the
hypotheses of (\btref{9.1.17}) are satisfied.

\begin{stmtup}{Example}{9.2.3}\btanchor{9.2.3}
Let $\Gamma$ be a group of automorphisms of $G$ preserving the root datum and the
bornology of $G$. We have seen (\btref{8.1.11}) that $\Gamma$ then operates on the building
$\bui$ of $G$ in such a way that $\gamma(g.x)=\gamma(g).\gamma(x)$ whatever
$\gamma\in\Gamma$, $g\in G$ \dots
\end{stmtup}
\par\endgroup

% ---------------------------------------- bruhattitsen_3.tex
\begingroup\providecommand{\nat}{}
\renewcommand{\nat}{^{\natural}}
\providecommand{\Imm}{}
\renewcommand{\Imm}{\mathscr{I}}
\providecommand{\Fil}{}
\renewcommand{\Fil}{\mathscr{F}}
\providecommand{\AFil}{}
\renewcommand{\AFil}{\mathscr{A}\mathscr{F}}
\providecommand{\R}{}
\renewcommand{\R}{\mathbf{R}}
\providecommand{\N}{}
\renewcommand{\N}{\mathbf{N}}
\providecommand{\Card}{}
\renewcommand{\Card}{\operatorname{Card}}
\providecommand{\Aut}{}
\renewcommand{\Aut}{\operatorname{Aut}}
\providecommand{\Hom}{}
\renewcommand{\Hom}{\operatorname{Hom}}
\providecommand{\Is}{}
\renewcommand{\Is}{\operatorname{Is}}
\providecommand{\Sm}{}
\renewcommand{\Sm}{\operatorname{Sim}}
\providecommand{\id}{}
\renewcommand{\id}{\mathrm{id.}}
\providecommand{\Kse}{}
\renewcommand{\Kse}{\mathrm{K}_{\sigma,\varepsilon}}
\providecommand{\Khse}{}
\renewcommand{\Khse}{\widehat{\mathrm{K}}_{\sigma,\varepsilon}}
\providecommand{\op}{}
\renewcommand{\op}[1]{\mathopen{]}#1\mathclose{[}}
\providecommand{\st}{}
\renewcommand{\st}[1]{\medskip\noindent\textbf{#1}\ ---\ }
\providecommand{\num}{}
\renewcommand{\num}[1]{\medskip\noindent\textbf{(#1)}\ }
\providecommand{\pf}{}
\renewcommand{\pf}{\smallskip\noindent\textit{Argument (outline).}\ }
\bigskip

The setting at the top of p.~205: a group $\Gamma$ acts on the building $\Imm$,
preserving the apartment $A$ and acting on it by affine automorphisms. More
generally $\gamma\in\Gamma$ carries an apartment $A'$ onto an apartment $A''$ by
an affine isomorphism. Taking
\[
\Imm\nat=\Imm^{\Gamma},\qquad G\nat=G^{\Gamma}
\]
(the $\Gamma$-fixed points) yields data satisfying (DI~1) and (DI~2). If in
addition $\Gamma$ preserves a chamber $C$ of $A$, then $A\nat$ contains the
center of $C$ (7.2.5) and (DI~3) holds with $F=C$ --- the situation that occurs
for quasi-split algebraic groups over a valued field.

% ---------------------------------------------------------------- 9.2.4
\num{9.2.4}\btanchor{9.2.4} \textbf{Dimension of a filter base.}
The dimension of a facet was defined in (7.2.4). More generally, let $\Fil$ be a
filter base in a finite-dimensional affine space. The affine subspaces spanned
by the members of $\Fil$ form a filter base $\AFil$; the \emph{affine subspace
generated by} $\Fil$ is the intersection $L$ of the members of $\AFil$ --- equivalently,
the smallest member of $\AFil$ --- and one sets $\dim\Fil=\dim L$. Note that
$\Fil$ is contained in $L$ in the sense of (7.2.1).

% ---------------------------------------------------------------- 9.2.5
\st{Proposition (9.2.5).}
\emph{Let $M$ be a convex subset of $\Imm$, and let $F$ (resp.\ $F'$) be a facet
of $\Imm$ that meets $M$ and that lies in the closure of no other facet of
$\Imm$ meeting $M$. Then:}
\begin{itemize}
\item[(i)] $\dim F=\dim F'$;
\item[(ii)] $\dim F\cap M=\dim F'\cap M$;
\item[(iii)] \emph{if $M$ lies in an apartment $A'$ of $\Imm$, then
$\dim F\cap M=\dim M$, $M$ is contained in the affine subspace $L$ of $A'$
generated by $F\cap M$, and $F\cap M$ has non-empty interior relative to $L$
(cf.\ (7.2.1));}
\item[(iv)] \emph{let $A'$ be an apartment containing $F$ and let $F''$ be a
facet of $A'$ meeting $M$. Suppose that at least one of the following holds:
$\dim F''=\dim F$; or $\dim F''\cap M=\dim F\cap M$; or $F''\cap M$ has
non-empty interior relative to the affine subspace $L$ generated by $F\cap M$ in
$A'$. Then $F''$ lies in the closure of no other facet of $\Imm$ meeting $M$.}
\end{itemize}

\pf
For (iii) one reduces to $A'=A$ and writes $F=F_{x,E}$ with $x\in\overline M$
and $E$ the direction of $F$ at $x$ (7.2.4). Let $\widetilde E$ be the span of
the convex cone $E$, and choose a convex $X\in F_{x,E}$ with $X\subset x+E$ and
$X\cap M\subset L$ (when $\varphi$ is discrete one may take for $X$ the facet of
\S\,2 corresponding to $F$). Picking $y$ interior to $X\cap M$ relative to $L$,
the open segment $\op{xy}$ consists of interior points of $X\cap M$. Suppose
some $z\in M$ lies outside $L$. The interior of the triangle $xyz$ lies in $M$.
If $z\in x+\widetilde E$, a point $u\in\op{yz}\cap(x+E)$ produces an open
segment inside $X\cap M$, forcing $u\in L$ and $z\in L$ --- contradiction. Hence
$z\notin x+\widetilde E$ and $[yz]\cap(x+\widetilde E)=\{y\}$, so there is a
facet direction $E'$ at $x$ with $E\subset\overline{E}'$ and
$\op{yz}\cap(x+E')\neq\emptyset$; this gives $F_{x,E'}\cap M\neq\emptyset$,
contradicting the maximality hypothesis on $F$ since $F\subset\overline
F_{x,E'}$. Therefore $M\subset L$, and (iii) follows.

For (i) and (ii): translating by an element of $G$ one may assume
$F,F'\subset A$, and replacing $M$ by $M\cap A$ one may assume $M\subset A$.
Then (ii) is immediate from (iii). For (i), write $F'=F_{x',E'}$ with
$x'\in\overline M$. By (iii), $M\subset L\subset x+\widetilde E=x'+\widetilde E$.
Since $E$ is a facet direction at $x$, the affine subspace $x+\widetilde E$ is
an intersection of walls of $A$, and every facet direction at $x'$ meeting
$\widetilde E$ lies in $\widetilde E$. If $\dim E'>\dim E$ one would get
$E'\cap\widetilde E=\emptyset$ and hence
$M\cap F_{x',E'}\subset L\cap F_{x',E'}=\emptyset$, which is false. So
$\dim E'\le\dim E$.

For (iv): choose a facet $F'''$ meeting $M$, maximal in the above sense, with
$F''\subset\overline{F}'''$; reduce to $F'''=F'$, $A'=A$, $M\subset A$, so
$M\subset L$. Each of the three hypotheses implies the next in the chain
\[
\dim F''\cap M=\dim F\cap M
\ \Longrightarrow\
\operatorname{int}_L(F''\cap M)\neq\emptyset
\ \Longrightarrow\
\dim F''=\dim F,
\]
the middle implication because all facets meeting $L$ in a germ of non-empty
interior have the same dimension, namely that of the intersection of the affine
roots containing $L$. Finally $\dim F''=\dim F=\dim F'$ gives $F''=F'$.

% ---------------------------------------------------------------- 9.2.6
\st{Corollary (9.2.6).}
\emph{(i) Let $F$ be a facet satisfying (DI~3), $A'$ an apartment containing
$F$, and $L$ the affine subspace of $A'$ generated by $F\cap A\nat$. Then
$A'\cap\Imm\nat\subset L$.}

\emph{(ii) If $A'$ is an apartment containing $A\nat$, then
$A'\cap\Imm\nat=A\nat$.}

\pf (i) is (9.2.5)(iii) applied to $M=A'\cap\Imm\nat$; (ii) follows from (i).

% ---------------------------------------------------------------- 9.2.7
\st{Proposition (9.2.7).}
\emph{Let $F$ be a facet satisfying (DI~3) and $C$ a chamber of $A$ with
$F\subset\overline C$. The retraction $\rho_{A;C}$ of $\Imm$ onto $A$ maps
$\Imm\nat$ onto $A\nat$; moreover its restriction to $\Imm\nat$ is independent
of the choice of such a chamber $C\subset A$.}

\pf For $x\in\Imm\nat$ take an apartment $A'$ containing $x$ and $F$. Then
$F\cap\Imm\nat=F\cap\Imm\nat\cap A\cap A'$ consists of fixed points of
$\rho_{A;C}$ (7.2.1), so $\rho_{A;C}$ carries the affine span $L$ of
$F\cap\Imm\nat$ in $A'$ onto its span in $A$, i.e.\ onto $A\nat$ by
(9.2.5)(iii); with (9.2.6) this gives $\rho_{A;C}(x)\in A\nat$. Since
$F\cap\Imm\nat$ has non-empty interior in $L$ (resp.\ $A\nat$), the restriction
of $\rho_{A;C}$ to $L$ is the \emph{unique} affine map $L\to A\nat$ restricting
to the identity on $F\cap\Imm\nat$ --- whence the independence of $C$.

% ---------------------------------------------------------------- 9.2.8
\num{9.2.8}\btanchor{9.2.8} Let $V\nat$ be the direction of $A\nat$ in $V$; notation as in (\btref{9.1.1}).

\st{Lemma.} $Z.A\cap\Imm\nat=A\nat$.

\pf Apply 7.6 with $\Phi_1=\Phi_0$ (a closed subsystem) and $T_1=T$, so that
$G_1=Z$. Let $\pi:Z.A\to\Imm_1$ be the canonical projection onto the building of
$Z$ ((\btref{7.6.4})(i)). On $A$, $\pi$ is the quotient by the subspace $L_1\subset V$
cut out by the kernels of the $a\in\Phi_0$; since $V\nat\subset L_1$, the image
$\pi(A\nat)$ is a single point. Given $x\in A$, $z\in Z$, choose an apartment
$A_1$ of $\Imm_1$ containing $\pi(A\nat)$ and $\pi(z.x)$; then
$A'=\pi^{-1}(A_1)$ is an apartment of $\Imm$ ((\btref{7.6.4})(ii)) containing $A\nat$
and $z.x$. So if $z.x\in\Imm\nat$ then $z.x\in A\nat$ by (9.2.6).

% ---------------------------------------------------------------- 9.2.9
\num{9.2.9}\btanchor{9.2.9} \textbf{Standing hypotheses for the rest of \S\,9.2.}
In addition to $G\nat$ and $\Imm\nat$ satisfying (DI~1), (DI~2), (DI~3), we are
given a root system $\Phi\nat$ in $(V\nat)^{*}$ and a generating root datum
$(T\nat,(U_b\nat,M_b\nat)_{b\in\Phi\nat})$ in $G\nat$ satisfying (DDR~1) (\btref{9.1.9})
and (DDR~2) (\btref{9.1.10}). Keeping the notation of (\btref{9.1.9}), we assume:
\begin{align*}
&\text{(DDR 3)} && T\nat\subset Z \quad\text{and}\quad M_b\nat\subset M_b
                   \quad\text{for all } b\in\Phi\nat;\\[1mm]
&\text{(DV 1)}  && \varphi\in A\nat;\\[1mm]
&\text{(DV 2)}  && \Card\bigl(\varphi_b\nat(U_b\nat)\bigr)\geq 3
                   \quad\text{for all } b\in\Phi\nat.
\end{align*}

% ---------------------------------------------------------------- 9.2.10
\st{Theorem (9.2.10).}
\emph{Under the hypotheses of (\btref{9.2.9}) the valuation $\varphi$ descends to
$G\nat$: that is, the family $\varphi\nat$ is a valuation of the root datum of
$G\nat$. If $\varphi$ is discrete, so is $\varphi\nat$.}

\pf One must verify (V~0)--(V~5) of (\btref{6.2.1}). Condition (V~0) is exactly (DV~2),
and (V~1), (V~3), (V~4) are already known from (\btref{9.1.10}). Discreteness is
immediate from the definition of $\varphi_b\nat$ (\btref{9.1.6}), which gives
\[
\varphi_b\nat(U_b\nat)\ \subset\!\!
\bigcup_{a\in\Phi,\ r\in\R^{*}_{+},\ a\nat=rb}\!\! r^{-1}\varphi_a(U_a)
\qquad (b\in\Phi\nat).
\]
There remain (V~2) and (V~5), proved in (\btref{9.2.12}) and (\btref{9.2.13}).

% ---------------------------------------------------------------- 9.2.11
\st{Lemma (9.2.11).}
\emph{Let $N\nat$ be the group generated by $T\nat$ and the $M_b\nat$,
$b\in\Phi\nat$ (cf.\ (\btref{6.1.2})(10)). For every $n\in N\nat$:}
\begin{equation}
n.A\nat=A\nat. \tag{1}
\end{equation}
\emph{If $\nu\nat(n)\in\Aut A\nat$ denotes the restriction of $n$ to $A\nat$,
then}
\begin{equation}
\lambda\circ\nu\nat(n)={}^{v}\nu\nat(n) \tag{2}
\end{equation}
\emph{where $\lambda:\Aut A\nat\to\Aut V\nat$ is the canonical homomorphism and
${}^{v}\nu\nat:N\nat\to{}^{v}W\nat$ is the homomorphism given by the root datum
of $G\nat$ ((\btref{6.1.2})(10)).}

\pf For (1) it is enough to treat $n\in T\nat$ or $n\in M_b\nat$; in both cases
$n=zn'$ with $n'\in N$, $z\in Z$ (by (DDR~3) and (\btref{9.1.3})), and (\btref{9.2.8}) gives
\[
n.A\nat\subset (n.A)\cap\Imm\nat=(z.A)\cap\Imm\nat\subset A\nat .
\]
For (2): if $\Phi\nat=\emptyset$ then $V\nat=\{0\}$ and there is nothing to
prove. Otherwise $N\nat$ is generated by the $M_b\nat$, so it suffices to see
that $\lambda\circ\nu\nat(m)=r_b$ for $m\in M_b\nat$. By (DR~5),
$mU_c\nat m^{-1}=U^{\natural}_{r_b(c)}$ for $c\in\Phi\nat$; by (\btref{9.1.8}),
$\nu\nat(m)$ carries a closed half-space
$\{x\in A\nat\mid c(x-\varphi)+k\ge 0\}$, $k\in\Gamma_c\nat$, onto one of the
form $\{x\in A\nat\mid r_b(c)(x-\varphi)+h\ge 0\}$, $h\in\R$. As $\Phi\nat$
spans $(V\nat)^{*}$ and $\Gamma_c\nat\neq\emptyset$, this forces
$\lambda\circ\nu\nat(m)=r_b$.

% ---------------------------------------------------------------- 9.2.12
\num{9.2.12}\btanchor{9.2.12} \textbf{Proof of (V~2).}
Let $b\in\Phi\nat$, $m\in M_b\nat$. By (9.2.11)(2), $\nu\nat(m)$ is a
translation $t\in V\nat$ followed by the orthogonal reflection in the hyperplane
$\{x\in A\nat\mid b(x-\varphi)=0\}$; hence it carries
$\{x\in A\nat\mid b(x-\varphi)+k\ge 0\}$ $(k\in\R)$ onto
$\{x\in A\nat\mid -b(x-\varphi)-b(t)+k\ge 0\}$. With (\btref{9.1.8}) this yields
\[
\varphi^{\natural}_{-b}(mum^{-1})=\varphi_b\nat(u)-b(t)
\qquad\text{for all } u\in U_b\nat,
\]
which is (V~2).

% ---------------------------------------------------------------- 9.2.13
\num{9.2.13}\btanchor{9.2.13} \textbf{Proof of (V~5).}
Let $b\in\Phi\nat$, $u\in U_b\nat$, $u\neq 1$, and put $k=\varphi_b\nat(u)$.
Choose $u',u''\in U^{\natural}_{-b}$ and $m\in M_b\nat$ with $u=u'mu''$, and set
$k'=\varphi^{\natural}_{-b}(u')$, $k''=\varphi^{\natural}_{-b}(u'')$. Let
$\alpha,\alpha',\alpha''$ be the fixed-point sets in $A\nat$ of $u,u',u''$
respectively; by (\btref{9.1.8}),
\begin{align*}
\alpha  &=\{x\in A\nat\mid  b(x-\varphi)+k\ \ge 0\},\\
\alpha' &=\{x\in A\nat\mid -b(x-\varphi)+k'\ge 0\},\\
\alpha''&=\{x\in A\nat\mid -b(x-\varphi)+k''\ge 0\}.
\end{align*}
It suffices to show that $\partial\alpha=\partial\alpha'$, equivalently that
$\alpha\cap\alpha'$ is neither empty nor of non-empty interior.

\smallskip
\emph{Case 1: $\alpha\cap\alpha'$ has interior.} Let $x$ be an interior point.
Then $u''.x=m^{-1}u'^{-1}u.x=m^{-1}.x\in A\nat$, and for every $y\in\alpha''$,
$d(u''.x,y)=d(u''.x,u''.y)=d(x,y)$. Since $\alpha''$ contains an open half-space
of $A\nat$, this forces $u''.x=x$. Hence $m$ fixes the non-empty open set
$\operatorname{int}(\alpha\cap\alpha')$ pointwise and induces the identity on
$A\nat$, contradicting (9.2.11).

\smallskip
\emph{Case 2: $\alpha\cap\alpha'=\emptyset$.} Take $x\in\alpha''$ and put
$x'=u.x=u'm.x$, so $x'\in u(A\nat)\cap u'(A\nat)\subset\Imm\nat$. The set
$X=A\nat\smallsetminus(\alpha\cup\alpha')$ is open and non-empty in $A\nat$, so
some facet $F$ of $A$ has $F\cap X$ of non-empty interior relative to $A\nat$;
by (9.2.5)(iv) this $F$ satisfies (DI~3). Let $C$ be a chamber of $A$ with
$F\subset\overline C$ and $\rho$ the retraction of $\Imm$ onto $A$ centred at
$C$; by (9.2.7), $x''=\rho(x')\in A\nat$. Choose an apartment $A'$ containing
$C$ and $x'$, and a point $y\in A'\cap X$ with $y\neq x''$ and $[yx'']$ not
parallel to $\partial\alpha$ or $\partial\alpha'$. The line through this segment
meets both $\alpha$ and $\alpha'$, which cut out on it two half-lines of
opposite directions; hence either there is $z\in\alpha$ with $y\in[zx'']$, or
there is $z'\in\alpha'$ with $y\in[z'x'']$.

In the first case, since $\rho$ does not increase distances and restricts to an
isometry $A'\to A$ (\btref{7.4.19}),
\[
d(z,x')\ \ge\ d(z,x'')=d(z,y)+d(y,x'')=d(z,y)+d(y,x')\ \ge\ d(z,x'),
\]
so $d(z,x')=d(z,y)+d(y,x')$ and $y\in[zx']$ (\btref{7.4.20}). As $z$ and $x'$ both lie
in the convex set $u(A\nat)$, we get $y\in u(A\nat)$. But
$d(u.y,t)=d(u.y,u.t)=d(y,t)$ for every $t$ in the closed half-space
$\alpha=u(\alpha)$ of the affine space $u(A\nat)$, which forces $u.y=y$; i.e.\
$y\in\alpha$, contradicting $y\notin\alpha\cup\alpha'$. The second case gives
$y\in u'(A\nat)$ and $y=u'.y$, the same contradiction. Hence
$\alpha\cap\alpha'\neq\emptyset$, and the proofs of (V~5) and of Theorem
(9.2.10) are complete.

% ---------------------------------------------------------------- 9.2.14
\num{9.2.14}\btanchor{9.2.14} \textbf{The canonical embedding $j$.}
The hypotheses of (\btref{9.1.17}) hold with $S\nat=T\nat$: indeed
$T\nat.\varphi\subset A\nat=\varphi+V\nat$ by (9.2.11)(1). Consequently:

\emph{there exists exactly one isometry $j$ of the building $\Imm\nat$ of
$G\nat$ (relative to $\varphi\nat$) into $\Imm$ such that $j(g.x)=g.j(x)$ for
$x\in\Imm\nat$, $g\in G\nat$, and $j(\varphi\nat)=\varphi$.}
If $A\nat=\varphi\nat+V\nat$ is the model apartment of $\Imm\nat$, then
$j(A\nat)=A\nat$ and $A\nat=A\cap j(\Imm\nat)$. Moreover
$j(\Imm\nat)\subset\Imm\nat$, \emph{but in general $j(\Imm\nat)\neq\Imm\nat$.}

Furthermore, \emph{the building $\Imm\nat$ of $G\nat$ does not depend (up to
isomorphism) on the choice of $\Imm\nat$ nor on the choice of $\varphi$ within
its equivalence class} (subject to the conditions of (\btref{9.2.9})): this follows from
(\btref{8.1.10}), since the bornology of $G\nat$ defined by $\varphi\nat$ is, by
(\btref{8.1.7}), the one induced by the bornology of $G$. In fact one can show that
$\Imm\nat$ depends only on $\Imm$ and $G\nat$ (cf.\ (\btref{8.1.12})). The next
proposition shows that $j$ depends only on $\Imm\nat$.

% ---------------------------------------------------------------- 9.2.15
\st{Proposition (9.2.15).}
\emph{(i) $A\nat$ is the smallest non-empty convex subset of $\Imm\nat$ stable
under $T\nat$.}

\emph{(ii) $j(\Imm\nat)$ is the smallest non-empty convex subset of $\Imm\nat$
stable under $G\nat$.}

\emph{(iii) Let $i$ be an isometry of $\Imm\nat$ into $\Imm\nat$ with
$i(g.x)=g.i(x)$ for $g\in G\nat$, $x\in\Imm\nat$. Then $i=j$.}

\pf For (i), argue as in (\btref{2.8.11}), taking for $g$ an element of $T\nat$ such
that $v={}^{v}\nu\nat(g)$ lies in the kernel of no vector root outside
$\Phi_0$, and applying (\btref{7.4.30}) in place of (\btref{2.8.10}), after noting that (9.2.7)
gives $\Imm\nat\subset\widehat P_x.A\nat$ for every $x\in A\nat$. Assertion (ii)
is an immediate consequence of (i). For (iii): $i(\Imm\nat)$ is a non-empty
convex subset of $\Imm\nat$, so $A\nat\subset i(\Imm\nat)$; since the
stabilizers in $G\nat$ of two distinct chambers of $\Imm\nat$ are distinct, one
deduces $i(A\nat)=A\nat$ and $i|A\nat=j|A\nat$, whence (iii).

\bigskip
\hrule
\bigskip

% ======================================================================

\btsecA{10}{10. Valuations of the classical groups}

The aim of \S\,10 is to show that the ``natural'' root datum of the classical
groups over a complete valued field --- possibly of infinite rank over its
center --- can be equipped with a valuation, making explicit along the way
several notions introduced in the earlier sections. Contrary to custom, the
authors begin with the orthogonal, unitary and symplectic groups; the (much
easier) linear group is treated more briefly in 10.2.

\btsecB{10.1}{10.1. Orthogonal, unitary and symplectic groups}

% ---------------------------------------------------------------- 10.1.1
\num{10.1.1}\btanchor{10.1.1} \textbf{Sesquilinear and pseudo-quadratic forms.}
(For the notions and properties used here, cf.\ [\btbib{38}], [\btbib{39}].)

Let $K$ be a field, not necessarily commutative, $\sigma$ an involution of $K$,
and $\varepsilon\in K$ with $\varepsilon=\pm 1$. Put
$\Kse=\{t-\varepsilon t^{\sigma}\mid t\in K\}$. Let $X$ be a right vector space
over $K$, let $f:X\times X\to K$ be a sesquilinear form (relative to $\sigma$)
such that
\begin{align}
f(y,x)&=\varepsilon f(x,y)^{\sigma} && (x,y\in X) \tag{1}\\
f(x,x)&=0 \quad\text{if } (\sigma,\varepsilon)=(\id,-1) && (x\in X), \tag{2}
\end{align}
and let $q:X\to K/\Kse$ be a \emph{pseudo-quadratic form} associated with $f$,
i.e.\ a function such that
\begin{align}
q(xk)&=k^{\sigma}q(x)k && (k\in K;\ x\in X) \tag{3}\\
q(x+y)&=q(x)+q(y)+\bigl(f(x,y)+\Kse\bigr) && (x,y\in X). \tag{4}
\end{align}
If $(\sigma,\varepsilon)=(\id,1)$, then $q$ is a quadratic form in the usual
sense and $f$ is the associated symmetric bilinear form.

Since a non-zero sesquilinear form has all of $K$ as image, $q$ determines $f$
except when $\Kse=K$, which is equivalent to $\sigma=\id$ and
$\varepsilon\neq 1$; in that last case $q$ is zero and $f$ is alternating. In
all cases $f$ is \emph{trace-valued}:
\begin{equation}
f(x,x)=k+\varepsilon k^{\sigma}\qquad\text{for } k\in q(x). \tag{5}
\end{equation}
Indeed, (4) gives
\begin{equation}
q(x(1+t))=q(x)+q(xt)+f(x,xt)+\Kse \tag{6}
\end{equation}
whence $(f(x,x)-k-\varepsilon k^{\sigma})t\in\Kse$ for every $t\in K$.

% ------- p. 212
If $\operatorname{char}K\neq 2$, then $f$ determines $q$. The same holds if
$\operatorname{char}K=2$ and the set of $\sigma$-invariant elements of $K$
equals $\Kse$: for if $f=0$ then $k=k^{\sigma}$ for every $k\in q(x)$ by (5),
whence $q=0$.

In particular, suppose there is a central element $\lambda$ of $K$ with
$\lambda+\lambda^{\sigma}=1$. This happens if $\operatorname{char}K\neq2$
(take $\lambda=1/2$), or if $\operatorname{char}K=2$ and $\sigma$ restricted to
the center of $K$ is not the identity. Taking $t=-\lambda$ in (6) gives
\begin{equation}
q(x)=\lambda f(x,x)+\Kse . \tag{7}
\end{equation}
Moreover
\begin{equation}
\Kse=\{t\in K\mid t^{\sigma}=-\varepsilon t\}, \tag{8}
\end{equation}
since $t^{\sigma}=-\varepsilon t$ implies $t=(\lambda t)-\varepsilon(\lambda t)^{\sigma}$.

\smallskip
\textbf{Extension of scalars.} Let $K'\supset K$ be a field and $\sigma'$ an
involution of $K'$ extending $\sigma$. Then $f$ and $q$ extend naturally to
$X\otimes_K K'$. For $f$ this is the familiar operation; the extension $q'$ of
$q$ is obtained as follows. Define $\mathrm{K}'_{\sigma',\varepsilon}$ as $\Kse$
was defined, and for
\begin{equation}
x=\sum_{i=1}^{m}x_i\otimes k_i \qquad (x_i\in X;\ k_i\in K') \tag{9}
\end{equation}
put
\begin{equation}
q'(x)=\sum_{i=1}^{m}k_i^{\sigma'}\bigl(q(x_i)+\mathrm{K}'_{\sigma',\varepsilon}\bigr)k_i
 \;+\!\!\sum_{\substack{i,j=1\\ i<j}}^{m}\!\!
   k_i^{\sigma'}\bigl(f(x_i,x_j)+\mathrm{K}'_{\sigma',\varepsilon}\bigr)k_j . \tag{10}
\end{equation}
This is well defined, since one checks easily that the right-hand side of (10)
does not depend on the decomposition (9) chosen for $x$. Note that when
$\mathrm{K}'_{\sigma',\varepsilon}\cap K\neq\Kse$ --- which can happen if (8)
fails --- the extension $q'$ does not determine $q$.

\smallskip
\textbf{Standing assumptions.} From now on the pair $(f,q)$ is assumed
\emph{non-degenerate}, i.e.:
\begin{equation}
\text{for }x\in X,\quad f(x,X)=\{0\}\ \text{ and }\ q(x)=0
\ \Longrightarrow\ x=0. \tag{11}
\end{equation}
We also assume $(f,q)$ has \emph{finite Witt index} $r$ (recall the Witt index
is the maximal dimension of a totally singular subspace $X'\subset X$, i.e.\ one
with $f(X',X')=\{0\}$ and $q(X')=\{0\}$).

Put $I=\{\pm1,\dots,\pm r\}$ and set $\varepsilon(i)=1$ (resp.\
$\varepsilon(i)=\varepsilon$) for $i\in I$ with $i>0$ (resp.\ $i<0$). Then $X$
admits a decomposition
\[
X=\coprod_{i\in I}e_i.K\;+\;X_0
\]
where the $e_i$ are elements of $X$ and $X_0$ a vector subspace, such that for
$i,j\in I$:
\begin{equation}
\left\{
\begin{aligned}
f(e_i,e_j)&=q(e_i)=0 &&\text{if } i\neq -j,\\
f(e_i,e_{-i})&=\varepsilon(i),\\
f(e_i,X_0)&=\{0\},\\
0&\notin q(X_0-\{0\}).
\end{aligned}\right. \tag{12}
\end{equation}

% ------- p. 213
Such a decomposition being fixed once and for all, put $X_i=e_i.K$; for an
endomorphism $g$ of $X$ let $c_{ij}(g)$ $(i,j\in I\cup\{0\})$ denote the element
of $\Hom(X_j,X_i)$ obtained by composing the canonical injection $X_j\to X$,
$g$, and the canonical projection $X\to X_i$ --- i.e.\ the $(i,j)$ entry of the
matrix of $g$ for this decomposition. Since a basis element $e_i$ has been
singled out in $X_i$ for $i\neq0$, we may, where convenient, identify
$\Hom(X_j,X_i)$ with $K$ when $i\neq 0\neq j$, with $X_0$ when $i=0\neq j$, and
with the dual $X_0^{*}$ of $X_0$ when $i\neq0=j$.

\smallskip
\emph{In the remainder of 10.1, $Z$ denotes the set
$\{(z,k)\mid z\in X_0,\ k\in K,\ q(z)=k+\Kse\}$. Unless explicitly stated
otherwise, $r$ is always assumed non-zero, and $\neq1$ if
$(\sigma,X_0)=(\id,\{0\})$.}

% ---------------------------------------------------------------- 10.1.2
\num{10.1.2}\btanchor{10.1.2} For $i\in I$ and $(z,k)\in Z$ (resp.\ $Z-\{(0,0)\}$), let
$u_i(z,k)$ (resp.\ $m_i(z,k)$) be the linear transformation of $X$ defined by
\[
u_i(z,k):\ \left\{
\begin{aligned}
x&\mapsto x-e_{-i}.f(z,x).\varepsilon(i) && (x\in X_0)\\
e_i&\mapsto e_i+z-e_{-i}.k.\varepsilon(i)\\
e_j&\mapsto e_j && (j\in I;\ j\neq i)
\end{aligned}\right.
\]
\[
\left(\text{resp.}\quad
m_i(z,k):\ \left\{
\begin{aligned}
x&\mapsto x-z.k^{-1}.f(z,x) && (x\in X_0)\\
e_i&\mapsto -e_{-i}.k.\varepsilon(i)\\
e_{-i}&\mapsto -e_i.(k^{\sigma})^{-1}.\varepsilon(-i)\\
e_j&\mapsto e_j && (j\in I;\ j\neq\pm i)
\end{aligned}\right.\right)
\]
For $i,j\in I$ with $j\neq\pm i$, and $k\in K$ (resp.\ $K-\{0\}$), let
$u_{ij}(k)$ (resp.\ $m_{ij}(k)$) be the linear transformation of $X$ defined by
\[
u_{ij}(k):\ \left\{
\begin{aligned}
x&\mapsto x && (x\in X_0)\\
e_i&\mapsto e_i+e_{-j}.k^{\sigma}.\varepsilon(-j)\\
e_j&\mapsto e_j-e_{-i}.k.\varepsilon(i)\\
e_h&\mapsto e_h && (h\in I-\{i,j\})
\end{aligned}\right.
\]
\[
\left(\text{resp.}\quad
m_{ij}(k):\ \left\{
\begin{aligned}
x&\mapsto x && (x\in X_0)\\
e_i&\mapsto e_{-j}.k^{\sigma}.\varepsilon(-j)\\
e_j&\mapsto -e_{-i}.k.\varepsilon(i)\\
e_{-i}&\mapsto e_j.k^{-1}.\varepsilon(i)\\
e_{-j}&\mapsto -e_i.(k^{\sigma})^{-1}.\varepsilon(-j)\\
e_h&\mapsto e_h && (h\in I-\{\pm i,\pm j\})
\end{aligned}\right.\right)
\]

% ------- p. 214
Consider the space $V^{*}=\R^{r}$ with its usual Euclidean norm, let
$(a_h)_{h=1,\dots,r}$ be its canonical basis, and set $a_{-h}=-a_h$ and
$a_{ij}=a_i+a_j$ for $i,j\in I$, $j\neq\pm i$. For $i\in I$ put
\begin{align*}
U_{a_i}   &=\{u_i(z,k)\mid (z,k)\in Z\},\\
U_{2a_i}  &=\{u_i(0,k)\mid k\in\Kse\},\\
M^{0}_{a_i}  &=\{m_i(z,k)\mid (z,k)\in Z-\{(0,0)\}\},\\
M^{0}_{2a_i} &=\{m_i(0,k)\mid k\in\Kse-\{0\}\},
\end{align*}
for $i,j\in I$ with $j\neq\pm i$ put
\[
U_{a_{ij}}=\{u_{ij}(k)\mid k\in K\},\qquad M^{0}_{a_{ij}}=\{m_{ij}(k)\mid k\in K-\{0\}\},
\]
and for $a\in V^{*}-\{a_i,2a_i,a_{ij}\mid i\in I,\ j\in I-\{\pm i\}\}$ put
$U_a=\{1\}$.

The set $\Phi$ of those $a\in V^{*}$ with $U_a-U_{2a}\neq\emptyset$ is a root
system. Precisely, with $i,j$ ranging over $I$ and $j\neq\pm i$:
\begin{align*}
\text{(B)}\quad  &\Phi=\{a_i,a_{ij}\}      &&\text{when } X_0\neq\{0\}\text{ and }(\sigma,\varepsilon)=(\id,1);\\
\text{(BC)}\quad &\Phi=\{a_i,2a_i,a_{ij}\} &&\text{when } X_0\neq\{0\}\text{ and }\sigma\neq\id;\\
\text{(C)}\quad  &\Phi=\{2a_i,a_{ij}\}     &&\text{when } X_0=\{0\}\text{ and }(\sigma,\varepsilon)\neq(\id,1);\\
\text{(D)}\quad  &\Phi=\{a_{ij}\}          &&\text{when } X_0=\{0\}\text{ and }(\sigma,\varepsilon)=(\id,1).
\end{align*}

% ---------------------------------------------------------------- 10.1.3
\st{Remark (10.1.3) (``Change of coordinates'').}
Let $\eta\in K$ be equal to $1$ if $\sigma=\id$ and to $\pm1$ if
$\sigma\neq\id$, and let $k$ be an invertible element of $K$ with
$k^{\sigma}=\eta k$. The conditions imposed on $\varepsilon,\sigma,f,q$ are
preserved, and the groups $U_{a_i},U_{a_{ij}}$ are unchanged, if one replaces
$\sigma,\varepsilon,f,q,e_i,X_0$ by $\sigma',\varepsilon',f'$, etc., where
\begin{gather*}
\varepsilon'=\varepsilon\eta,\\
t^{\sigma'}=kt^{\sigma}k^{-1}\quad\text{for } t\in K
 \quad(\text{whence } \mathrm{K}_{\sigma',\varepsilon'}=k.\Kse),\\
f'=kf,\qquad q'=kq,\\
e_i'=e_i.k^{-1}\ \text{ or }\ e_i \quad\text{according as } i<0 \text{ or } i>0,\\
X_0'=X_0 .
\end{gather*}

If $\sigma\neq\id$ and $\varepsilon\neq1$, one may take for $k$ a non-zero
anti-invariant element, e.g.\ one of the form $h-h^{\sigma}$; then
$\varepsilon'=1$.

If $(\sigma,\varepsilon)\neq(\id,1)$ and one takes for $k$ an invertible element
of $\Kse$ (whose inverse then also lies in $\Kse$), one gets
$\varepsilon'=-1$ and $1\in\mathrm{K}_{\sigma',\varepsilon'}$. On the other
hand, if $(\sigma,\varepsilon)=(\id,1)$ and $X_0\neq\{0\}$, it is possible to
choose $k$ so that $q'$ takes the value $1$ at an arbitrarily prescribed point
of $X_0$. Hence, for the study of the system of groups $(U_a)_{a\in\Phi}$, one
may always reduce to one of the following cases:
\begin{align}
&\varepsilon=-1 \ \text{ and }\ 1\in\Kse
   &&(\text{whence } \Phi=C_r \text{ or } BC_r); \tag{1}\\
&(\sigma,\varepsilon)=(\id,1) \ \text{ and }\ 1\in q(X_0)
   &&(\text{whence } \Phi=B_r); \tag{2}\\
&(\sigma,\varepsilon,X_0)=(\id,1,\{0\})
   &&(\text{whence } \Phi=D_r). \tag{3}
\end{align}

% ---------------------------------------------------------------- 10.1.4
\num{10.1.4}\btanchor{10.1.4} Let $g\in\Hom(X,X)$. One says $g$ is a \emph{similitude} if there
is an element $c$ of the center of $K$, invariant under $\sigma$, such that for
$x,y\in X$
\[
f(gx,gy)=cf(x,y),\qquad q(gx)=cq(x).
\]
This element $c$ is then unique; it is called the \emph{ratio} (multiplier) of
the similitude $g$ and is written $c(g)$. An \emph{isometry} is by definition a
similitude of ratio $1$.

We write $G^{0}$ for the group generated by the $U_a$ $(a\in\Phi)$,
$\Is(f,q)$ for the group of isometries, $\Sm(f,q)$ for the group of similitudes,
and $T^{0}$ (resp.\ $T(f,q)$; resp.\ $\widetilde T(f,q)$) for the group of those
elements $t$ of $G^{0}$ (resp.\ $\Is(f,q)$; resp.\ $\Sm(f,q)$) stabilizing each
of the $X_i$ $(i\in[-r,+r])$, i.e.\ represented by a diagonal matrix
$((c_{ij}(t)))$. If $\Phi$ is of type $D$, it is known that the linear subspaces
of $X$ of dimension $r$ on which $q$ vanishes (the maximal totally singular
subspaces) split into two families, two such subspaces belonging to the same
family if and only if their intersection has even codimension in each of them;
we then write $\Is^{+}(f,q)$ (resp.\ $\Sm^{+}(f,q)$) for the subgroup of index
$2$ of $\Is(f,q)$ (resp.\ $\Sm(f,q)$) consisting of those elements preserving
each of these two families.

% ---------------------------------------------------------------- 10.1.5
\st{Proposition (10.1.5).}
\emph{(i) For a linear transformation $t$ of $X$ stabilizing each $X_i$
$(i\in I\cup\{0\})$ to be a similitude of ratio $c$, it is necessary and
sufficient that}
\begin{align}
c_{ii}(t)^{\sigma}.c_{-i,-i}(t)&=c &&\text{for } i\in I \tag{1}\\
\intertext{\emph{and}}
q\bigl(c_{00}(t)(z)\bigr)&=c.q(z) &&\text{for } z\in X_0. \tag{2}
\end{align}
\emph{(ii) The group $\widetilde T(f,q)$ normalizes each of the groups $U_a$
$(a\in\Phi)$ and $G^{0}$.}

\emph{(iii) If $\Phi$ is not of type $D$, then $\Is(f,q)=T(f,q).G^{0}$ and
$\Sm(f,q)=\widetilde T(f,q).G^{0}$. If $\Phi$ is of type $D$, then
$\Is^{+}(f,q)=T(f,q).G^{0}$ and $\Sm^{+}(f,q)=\widetilde T(f,q).G^{0}$.}

% ---------------------------------------------------------------- 10.1.6
\st{Proposition (10.1.6).}
\emph{Let $T$ be a subgroup of $\widetilde T(f,q)$ containing $T^{0}$, let $G$
be the group $T.G^{0}$, and for $a\in\Phi$ put $M_a=T.M^{0}_a$. Then the pair
$(T,(U_a,M_a)_{a\in\Phi})$ is a generating root datum of type $\Phi$ in $G$.
With the notation of (\btref{6.1.2})(2) one has, for this datum,
$m(u_i(z,k))=m_i(z,k)$ and $m(u_{ij}(k))=m_{ij}(k)$; in particular the set
$M^{0}_a$ $(a\in\Phi)$ coincides with the set so denoted in (\btref{6.1.2})(2).}

\smallskip
These two propositions are proved in (\btref{10.1.10}) and (\btref{10.1.12}) respectively.

% ---------------------------------------------------------------- 10.1.7
\st{Proposition (10.1.7).}
\emph{(i) The group $N^{0}$ generated by the $M^{0}_a$ $(a\in\Phi)$ stabilizes
$X_0$ and permutes the $X_i$ $(i\in I)$. The group ${}^{v}W$ of permutations of
the index set $I$ that it induces is the group $W_C$ of all permutations
commuting with multiplication by $-1$, or the group $W_D$ of all even
permutations with that property, according as the root system $\Phi$ is not, or
is, of type $D$.}

\emph{(ii) The group $N'_1$ generated by the $m_{ij}(1)$ $(i,j\in I;\
j\neq\pm i)$ is a finite subgroup of $N^{0}$ preserving the set
$\{\pm e_i\mid i\in I\}$, fixing $X_0$, and whose image in ${}^{v}W$ is the
group $W_D$.}

% ------- p. 216
\emph{(iii) Suppose one of the conditions (1), (2), (3) of (10.1.3) holds, and
let $e_0$ be an element of $X_0$ equal to $0$ in case (1) and such that
$q(e_0)=1$ in case (2). Then the group $N'$ generated by $N'_1$ and, if $\Phi$
is not of type $D$, by the $m_i(e_0,1)$ $(i\in I)$, is a finite group preserving
the set $\{\pm e_i\mid i\in I\}$, whose image in ${}^{v}W$ is all of
${}^{v}W$.}

\pf It is clear that $N'_1$ preserves $\{\pm e_i\mid i\in I\}$ and fixes $X_0$;
hence it is a finite group. Finiteness of $N'$ follows because every element of
the given generating system --- hence every element of the group itself ---
preserves $\{\pm e_i\mid i\in I\}$ and induces on $X_0$ either the identity or
the transformation $x\mapsto x-e_0.f(e_0,x)$, which has order $1$ or $2$. The
remaining assertions are obvious.

% ---------------------------------------------------------------- 10.1.8
\st{Lemma (10.1.8).}
\emph{(i) For $i\in I$, $U_{a_i}$ (resp.\ $U_{2a_i}$) is the group of all
isometries $1+n$ of $X$ such that}
\begin{align*}
n(X_i)&\subset X_0+X_{-i} &&(\text{resp.\ } X_{-i}),\\
n(X_0)&\subset X_{-i}     &&(\text{resp.\ } =\{0\}),\\
n(X_j)&=\{0\} &&\text{for } j\neq i,0.
\end{align*}
\emph{This group acts simply transitively on the set of singular lines of $q$
(i.e.\ one-dimensional subspaces on which $q$ vanishes) contained in
$X_{-i}+X_0+X_i$ (resp.\ $X_{-i}+X_i$) and distinct from $X_{-i}$.}

\emph{(ii) For $i,j\in I$ with $j\neq\pm i$, $U_{a_{ij}}$ is the group of all
isometries $1+n$ of $X$ such that}
\begin{align*}
n(X_i)&\subset X_{-j},\\
n(X_j)&\subset X_{-i},\\
n(X_h)&=\{0\} &&\text{for } h\in[-r,r]-\{i,j\}.
\end{align*}
\emph{This group acts simply transitively on the set of lines of $X_i+X_{-j}$
distinct from $X_{-j}$. It is contained in $\Is^{+}(f,q)$ if $\Phi$ is of type
$D$.}

\smallskip
The verification is easy.

% ---------------------------------------------------------------- 10.1.9
\st{Lemma (10.1.9).}
\emph{(i) For every $i\in I$, the stabilizer of $e_i$ in $G^{0}$ acts
transitively on the lines $Y$ of $X$ with $f(e_i,Y)\neq\{0\}$ and $q(Y)=\{0\}$.}

\emph{(ii) The group $G^{0}$ acts transitively on the singular lines of $q$.}

\pf (i) is proved by induction on $r$; one may assume $i=1$. For $r=1$ the
assertion is (10.1.8)(i). Let $r\ge2$ and let $x\in X$ satisfy $f(e_1,x)\neq0$
and $q(x)=0$; it suffices to find $g\in G^{0}$ with $g(e_1)=e_1$ and
$g(x)\in X_{-1}$. The set of $y\in X_1+X_r$ with $f(y,x)=0$ is a line distinct
from $X_1$; by (10.1.8)(ii), after transforming $x$ by a suitable element of
$U_{a_{r,-1}}$, one may assume this line is $X_r$, i.e.\ $f(e_r,x)=0$. Then
$x\in\sum_{i=-r+1}^{r}X_i$; write $x=x'+e_r.k$ with
$x'\in X'=\sum_{i=-r+1}^{r-1}X_i$ and $k\in K$. By the inductive hypothesis
applied to $X'$ and the restrictions of $f$ and $q$, there is an element of
$G^{0}$ fixing $e_1,e_r,e_{-r}$ and carrying $x'$ into a point of $X_{-1}$. One
may therefore assume directly that $x\in X_{-1}+X_r$; but then $U_{a_{-1,-r}}$
contains an element $g$ with the required properties, by (10.1.8)(ii).

For (ii) it suffices to show that if $x\in q^{-1}(0)$ there is $g\in G^{0}$ with
$g(x)\in X_1$. Assume $x\neq0$; there is $i\in I$ with $f(e_i,x)\neq0$. If
$i\neq-1$, put $y=e_{-1}+e_{-i}$; if $i=-1$ and $r\ge2$, put
$y=e_1-e_2+e_{-2}-e_{-1}$; finally, if $i=-1$ and $r=1$ --- in which case
$\dim X_0>0$ by hypothesis --- take
$y=u_1(k,z)(e_1)=e_1+z-e_{-1}k\varepsilon(1)$ with $(z,k)\in Z-\{(0,0)\}$. In
each case $q(y)=0$, $f(e_1,y)\neq0$ and $f(e_i,y)\neq0$. By (i) there is an
element $g_1$ (resp.\ $g_i$) of the stabilizer of $e_1$ (resp.\ $e_i$) in
$G^{0}$ with $g_1(y)\in X_{-1}$ (resp.\ $g_i(y)\in xK$), whence the assertion.

% ---------------------------------------------------------------- 10.1.10
\num{10.1.10}\btanchor{10.1.10} \textbf{Proof of Proposition (10.1.5).}
Assertion (i) results from a straightforward computation, and (ii) is an
immediate consequence of (10.1.8).

For (iii) we restrict to the case where $\Phi$ is not of type $D$, the case
$\Phi=D_r$ being handled analogously. By (10.1.8) one has
$G^{0}\subset\Is(f,q)\subset\Sm(f,q)$, so it suffices to prove
$\Sm(f,q)\subset\widetilde T(f,q).G^{0}$. The proof is by induction on $r$, the
assertion being obvious for $r=0$ (when $\Phi=D_r$ one starts the induction at
$r=1$). Let $g\in\Sm(f,q)$. Applying (10.1.9)(ii) and then (10.1.9)(i), one
finds $g^{0}\in G^{0}$ with $g^{0}(g(e_1))=e_1$ and
$g^{0}(g(e_{-1}))=e_{-1}$, whence $g^{0}(g(X'))=X'$, where
\[
X'=\{x\mid x\in X,\ f(e_1,x)=f(e_{-1},x)=0\}=\sum_{i\neq\pm1}X_i .
\]
By the inductive hypothesis applied to the restriction of $g^{0}g$ to $X'$, one
gets $g^{0}g\in G^{0}.\widetilde T(f,q)$, whence
$g\in G^{0}.\widetilde T(f,q)$. \hfill q.e.d.

% ---------------------------------------------------------------- 10.1.11
\num{10.1.11}\btanchor{10.1.11} \textbf{Relations.}
In the relations below, $h,i,j,i',j'$ denote elements of $I$, $k,k'$ elements of
$K$, and $z,z'$ elements of $X_0$. These elements are in each case assumed to
satisfy the conditions needed for the relations written to make sense (thus in
relation (3) one must have $(z,k)\in Z-\{(0,0)\}$; in relation (15),
$j\neq\pm i$ and $j'\neq\pm i'$; etc.).
\begin{align}
&u_i(z,k).u_i(z',k')=u_i(z+z',\,k+k'+f(z,z')), \tag{1}\\
&u_i(z,k)^{-1}=u_i(-z,\,\varepsilon k^{\sigma}), \tag{2}\\
&u_i(z,k)=u_{-i}\bigl(-z.k^{-1}.\varepsilon(i),\,k^{-1}.\varepsilon\bigr)
   .m_i(z,k).
   u_{-i}\bigl(-z.(k^{\sigma})^{-1}.\varepsilon(-i),\,k^{-1}.\varepsilon\bigr), \tag{3}\\
&u_{ij}(k).u_{ij}(k')=u_{ij}(k+k'), \tag{4}\\
&u_{ij}(k)^{-1}=u_{ij}(-k), \tag{5}\\
&u_{ij}(k)=u_{-i,-j}\bigl(-(k^{\sigma})^{-1}.\varepsilon(i).\varepsilon(j)\bigr)
   .m_{ij}(k).
   u_{-i,-j}\bigl(-(k^{\sigma})^{-1}.\varepsilon(i).\varepsilon(j)\bigr), \tag{6}\\
&u_{ji}(k)=u_{ij}(-\varepsilon k^{\sigma}), \tag{7}\\
&\bigl(u_i(z,k),\,u_i(z',k')\bigr)=u_i\bigl(0,\,f(z,z')-f(z',z)\bigr), \tag{8}
\end{align}
% ------- p. 218
\begin{align}
&\bigl(u_i(z,k),\,u_{i'}(z',k')\bigr)=u_{ii'}(f(z,z'))
   \qquad\text{for } i'\neq\pm i, \tag{9}\\
&\bigl(u_{-i}(z,k),\,u_{ij}(k')\bigr)
   =u_j\bigl(-z.k'.\varepsilon(i),\,\varepsilon k'^{\sigma}.k^{\sigma}.k'\bigr)
    .u_{-i,j}\bigl(-k.k'.\varepsilon(i)\bigr), \tag{10}\\
&\bigl(u_i(z,k),\,u_{i',j'}(k')\bigr)=1
   \qquad\text{for } i\notin\{-i',-j'\}, \tag{11}\\
&\bigl(u_{ih}(k),\,u_{-h,j}(k')\bigr)=u_{ij}\bigl(-k.k'.\varepsilon(-h)\bigr)
   \qquad\text{for } j\neq\pm i, \tag{12}\\
&\bigl(u_{ih}(k),\,u_{-h,i}(k')\bigr)
   =u_i\bigl(0,\,(k'^{\sigma}.k^{\sigma}-k.k'.\varepsilon).\varepsilon(h)\bigr), \tag{13}\\
&\bigl(u_{ij}(k),\,u_{ij'}(k')\bigr)=1 \qquad\text{for } j'\neq-j, \tag{14}\\
&\bigl(u_{ij}(k),\,u_{i',j'}(k')\bigr)=1
   \qquad\text{for } i',j'\notin\{\pm i,\pm j\}. \tag{15}
\end{align}
The verifications are easy.

% ---------------------------------------------------------------- 10.1.12
\num{10.1.12}\btanchor{10.1.12} \textbf{Proof of Proposition (10.1.6).}
For $a\in\Phi$, the quotient of two elements of $M^{0}_a$ stabilizes each
$X_i$; hence it lies in $T^{0}$, and $M_a=T.M^{0}_a$ is a right coset of $T$.
Granting this, the fact that the system $(T,(U_a,M_a)_{a\in\Phi})$ satisfies
axioms (DR~1), (DR~2) and (DR~4) of root data (\btref{6.1.1}) follows from the relations
of (\btref{10.1.11}). Axiom (DR~3) is obvious. To check (DR~5) it clearly suffices to
treat the case where $n_a$ belongs to $M^{0}_a$; the required relation then
follows from the shape of the $m_i$ and $m_{ij}$, from the fact that they lie in
$\Is(f,q)$ (cf.\ (10.1.5)), and from Lemma (10.1.8). Axiom (DR~6) follows
because the matrix $((c_{ij}(g)))$ (cf.\ (\btref{10.1.1})) representing an element $g$
of $TU^{+}$ (resp.\ of $U^{-}$) is lower triangular (resp.\ upper triangular
with only $1$'s on the diagonal). Finally, the assertion of the statement
concerning the $m(u_i)$ and $m(u_{ij})$ is a consequence of relations (3) and
(6) of (\btref{10.1.11}).

% ---------------------------------------------------------------- 10.1.13
\num{10.1.13}\btanchor{10.1.13} For any function $\omega:K\to\R\cup\{\infty\}$, we write
$\varphi^{\omega}=(\varphi^{\omega}_a)_{a\in\Phi}$ for the system of functions
$\varphi^{\omega}_a:U_a\to\R\cup\{\infty\}$ defined by
\begin{align}
\varphi^{\omega}_{a_i}\bigl(u_i(z,k)\bigr) &=\tfrac{1}{2}\omega(k)
   &&\text{if } a_i\in\Phi \tag{1}\\
\varphi^{\omega}_{2a_i}\bigl(u_i(0,k)\bigr) &=\omega(k)
   &&\text{if } 2a_i\in\Phi \tag{2}\\
\varphi^{\omega}_{a_{ij}}\bigl(u_{ij}(k)\bigr) &=\omega(k). \tag{3}
\end{align}

% ---------------------------------------------------------------- 10.1.14
\num{10.1.14}\btanchor{10.1.14} Let $\omega:K\to\R\cup\{\infty\}$ be a non-improper valuation of
$K$, invariant under $\sigma$, and let $\widehat K$ be the completion of $K$ for
$\omega$. The continuous extensions of $\omega$ and $\sigma$ to $\widehat K$ are
again written $\omega$ and $\sigma$. Put
$\widehat X=X\otimes_K\widehat K$, $\widehat X_i=X_i\otimes_K\widehat K$, write
$\widehat f:\widehat X\times\widehat X\to\widehat K$ and
$\widehat q:\widehat X\to\widehat K/\Khse$ for the natural extensions of $f$ and
$q$ (\btref{10.1.1}), and set
\[
\widehat Z=\{(z,k)\mid z\in\widehat X_0,\ k\in\widehat K,\
\widehat q(z)=k+\Khse\}.
\]
We equip $K$ and $\widehat K$ with the topology defined by $\omega$, and every
finite-dimensional vector space over $K$ or $\widehat K$ with the topology
naturally associated with it (the topology transported from the product topology
on $K^{m}$ or $\widehat K^{m}$). Closure for these topologies is denoted by a
bar. One has
\begin{equation}
\Khse\subset\overline{\Kse}. \tag{1}
\end{equation}
Put
\[
X^{\perp}=\{x\in X\mid f(x,X)=\{0\}\},\qquad
\widehat X^{\perp}=\{x\in\widehat X\mid f(x,\widehat X)=\{0\}\}.
\]
One also has $X^{\perp}=\{z\in X_0\mid f(z,X_0)=\{0\}\}$ and
\begin{equation}
\widehat X^{\perp}=\{z\in\widehat X_0\mid f(z,X_0)=\{0\}\}
 =X^{\perp}\otimes_K\widehat K. \tag{2}
\end{equation}
Note further that, for $(z,k)\in Z$,
\begin{equation}
\omega(f(z,z))=\omega(k+\varepsilon k^{\sigma})\ \ge\ \omega(k). \tag{3}
\end{equation}

% ---------------------------------------------------------------- 10.1.15
\st{Theorem (10.1.15).}
\emph{With the notation and conventions of (\btref{10.1.14}), the following conditions
are equivalent.}
\begin{itemize}
\item[(i)] \emph{For $(z,k),(z',k')\in Z$,}
 $\ \omega(f(z,z'))\ \ge\ \tfrac12\bigl(\omega(k)+\omega(k')\bigr).$
\item[(ii)] \emph{For $(z,k),(z',k')\in Z$,}
 $\ \omega(f(z,z'))\ \ge\ \inf\{\omega(k),\omega(k')\}.$
\item[(iii)] $\widehat q(\widehat X_0-\widehat X^{\perp})\cap\overline{\Kse}=\emptyset$.
\item[(iv)] \emph{$\varphi^{\omega}$ is a valuation of the root datum
$(T,(U_a))$.}
\end{itemize}

% ---------------------------------------------------------------- 10.1.16
\st{Remarks (10.1.16).}
\emph{a)} If $\Khse$ is closed in $\widehat K$ (for instance if $K$ is
finite-dimensional over its center, or if $\varepsilon=1$ and there is a central
element $\lambda$ with $\lambda+\lambda^{\sigma}=1$, cf.\ (\btref{10.1.1})(8)), then
condition (iii) means that the restriction of $q$ to $\widehat X_0$ does not
vanish outside $\widehat X^{\perp}$.

\emph{b)} If $\operatorname{char}K=2$, it may happen that $q$ has non-trivial
zeros in $\widehat X^{\perp}$. For example, suppose $K$ is a non-perfect
commutative field whose completion is perfect, and let $\sigma$ be the identity,
$X_0$ the field $K^{1/2}$ regarded as a vector space over $K$, $q|_{X_0}$ the
quadratic form $z\mapsto z^{2}$, and $f|_{X_0\times X_0}$ the zero bilinear
form. Then $\widehat q$ is the square of a linear form and hence has a kernel of
codimension $1$ in $\widehat X_0$.

% ---------------------------------------------------------------- 10.1.17
\st{Lemma (10.1.17).}
\emph{Let $Y$ be a finite-dimensional subspace of $X_0$ and
$\widehat Y=Y\otimes_K\widehat K$; identify $Y$ with $Y\otimes1$. Then:}

\emph{(i) There is a continuous function $\overline q:\widehat Y\to\widehat K$
such that $\overline q(Y)\subset K$ and}
\begin{align*}
q(z)&=\overline q(z)+\Kse  &&\text{for } z\in Y,\\
\widehat q(z)&=\overline q(z)+\Khse &&\text{for } z\in\widehat Y.
\end{align*}

\emph{(ii) Let $(z,k)\in\widehat Y$ and $z_n\in Y$ $(n\in\N)$ with
$\lim_{n\to\infty}z_n=z$. Then there exist $k_n\in K$ with
$\lim_{n\to\infty}k_n=k$ and $q(z_n)=k_n+\Kse$ for every $n$.}

\pf In what follows $s,s'$ denote integers running through $[1,\dim Y]$. Let
$(x_s)$ be a basis of $Y$ and, for each $s$, let $d_s$ be a representative of
$q(x_s)$ in $K$. Then the function $\overline q$ defined by
\[
\overline q\Bigl(\sum_s x_s\xi_s\Bigr)
 =\sum_s \xi_s^{\sigma}d_s\xi_s+\sum_{s<s'}\xi_s^{\sigma}f(x_s,x_{s'})\xi_{s'}
\]
satisfies the conditions of (i).

For (ii), let $(z,k)$ and $z_n$ be as in the statement. By (\btref{10.1.14})(1) there
are $k'_n\in\Kse$ with $\lim_{n\to\infty}k'_n=k-\overline q(z)$. The sequence
$k_n=\overline q(z_n)+k'_n$ then has the required property.

% ---------------------------------------------------------------- 10.1.18
\st{Lemma (10.1.18).}
\emph{Let $(z,k),(z_0,k_0)\in Z-\{(0,0)\}$ and}
\[
\eta=\omega(k)+\omega(k_0)-2\omega(f(z,z_0)).
\]
\emph{Suppose $\eta$ is strictly positive. For $n\in\N^{*}$, define
$z_n\in X_0$, $k_n\in K$ and $\lambda_n\in K^{*}$ inductively by}
\[
\lambda_n=-f(z_{n-1},z)^{-1}.k_{n-1},\qquad
z_n=z_{n-1}+z.\lambda_n,\qquad k_n=\lambda_n^{\sigma}.k.\lambda_n .
\]
\emph{Then:}

\emph{(i) For every $n\in\N$,}
\begin{align}
&(z_n,k_n)\in Z \tag{1}\\
&\omega(\lambda_{n+1})=(2^{n}-1)\eta+\omega(k_0)-\omega(f(z_0,z)); \tag{2}\\
&\omega(f(z_n,z))=\omega(f(z_0,z)); \tag{3}\\
&\omega(k_n)=(2^{n}-1)\eta+\omega(k_0). \tag{4}
\end{align}

\emph{(ii) The sequence $(z_n)$ converges in $z.\widehat K+z_0.\widehat K$ to a
point $\widehat z$ such that}
\[
\omega(f(\widehat z,z))=\omega(f(z_0,z))
\qquad\text{and}\qquad
\widehat q(\widehat z)\subset\overline{\Kse}/\Khse .
\]

\emph{(iii) For every real number $A$ there exist $(z',k'),(z'_0,k'_0)\in Z$
with $z',z'_0\in z.K+z_0.K$, $\omega(f(z'_0,z'))=\omega(f(z_0,z))$,
$\omega(k')\ge A$ and $\omega(k'_0)\ge A$.}

\pf The verification of (i) by induction on $n$ is immediate, taking
(\btref{10.1.14})(3) into account. Assertion (ii) follows at once from (i). For (iii),
note first that by (i) there is an integer $n$ with $\omega(k_n)\ge A$; put
$(z',k')=(z_n,k_n)$ and apply (i) again, substituting $(z',k')$ and $(z,k)$ for
$(z,k)$ and $(z_0,k_0)$ respectively.

% ---------------------------------------------------------------- 10.1.19
\num{10.1.19}\btanchor{10.1.19} \textbf{Proof of Theorem (10.1.15).}
It is clear that (i) implies (ii). The implications (ii)$\Rightarrow$(i) and
(iii)$\Rightarrow$(i) are immediate consequences of (10.1.18)(iii) and
(10.1.18)(ii) respectively.

Suppose there is $\widehat z\in\widehat X_0-\widehat X^{\perp}$ with
$\widehat q(\widehat z)\in\overline{\Kse}/\Kse$, and let $z'\in X_0$ and
$\widehat k\in\overline{\Kse}$ be such that $f(\widehat z,z')\neq0$ and
$\widehat q(\widehat z)=\widehat k+\Khse$. By (10.1.17)(ii),
$(\widehat z,\widehat k)$ is a limit of a sequence $((z_n,k_n))_{n\in\N}$ of
elements of $Z$. Let $(\ell_n)_{n\in\N}$ be a sequence in $\Kse$ tending to
$\widehat k$, and put $k'_n=k_n-\ell_n$. Then $(z_n,k'_n)\in Z$,
$\lim_{n\to\infty}k'_n=0$ and
$\lim_{n\to\infty}f(z_n,z')=f(\widehat z,z')\neq0$, so that
$\omega(f(z_n,z'))<\tfrac12(\omega(k)+\omega(k'_n))$ for $n$ large enough. This
shows (i)$\Rightarrow$(iii).

The implication (iv)$\Rightarrow$(ii) follows from (\btref{10.1.11})(1) and axiom (V~1)
of valuations (\btref{6.2.1}).

It remains to show that if conditions (i), (ii), (iii) hold --- which we assume
from now on --- then $\varphi^{\omega}$ is a valuation of $(T,(U_a))$.

Axiom (V~0) is a consequence of the assumption that $\omega$ is non-improper
(when $a=a_i$, note that $(z,k)\in Z$ implies
$(z\ell,\ell^{\sigma}k\ell)\in Z$ for every $\ell\in K$).

Axiom (V~1) holds by (ii) together with relations (1) and (4) of (\btref{10.1.11}).

Axiom (V~2) follows from the observations: if $a=a_{ij}$ and $m\in M_a$, then
\[
m.u_{ij}(k).m^{-1}  =u_{-i,-j}\bigl(c_{i,-i}(m).k.c_{j,-j}(m^{-1}).\varepsilon\bigr);
\]
if $a=a_i$ or $2a_i$ and $m\in M_a$, then
\[
m.u_i(z,k).m^{-1}  =u_i\bigl(m(z),\ c_{i,-i}(m).k.c_{i,-i}(m^{-1})\bigr).
\]

Axiom (V~3) follows from relations (8)--(15) of (\btref{10.1.11}), from (i), and from
the fact that $\omega$ is a valuation.

Finally the validity of (V~4) is obvious, and that of (V~5) follows from
(\btref{10.1.11})(3) and (6).

% ---------------------------------------------------------------- 10.1.20
\num{10.1.20}\btanchor{10.1.20} Let $X_0^{*}=\Hom(X_0,K)$ be the dual of $X_0$. For any function
$\omega:K\to\R\cup\{\infty\}$ and any function
$\omega_1:X_0\to\R\cup\{\infty\}$, let
$\omega_q:X_0\to\R\cup\{\infty\}$ and
$\omega_1^{*}:X_0^{*}\to\R\cup\{\infty\}\cup\{-\infty\}$ be the functions
defined by the following relations, where one agrees that
$\infty-\infty=\infty$:
\begin{align*}
\omega_q(z)&=\tfrac12\sup\{\omega(k)\mid k\in q(z)\} && (z\in X_0)\\
\omega_1^{*}(x)&=\inf\{\omega(x(z))-\omega_1(z)\mid z\in X_0\} && (x\in X_0^{*}).
\end{align*}
A function $\alpha:Y\to\R\cup\{\infty\}\cup\{-\infty\}$ defined on a
$K$-vector space $Y$ is called an \emph{($\omega$-)additive semi-norm} if, for
$\lambda\in K$ and $y,y'\in Y$,
\begin{align}
\alpha(y\lambda)&=\alpha(y)+\omega(\lambda), \tag{1}\\
\alpha(y+y')&\ge\inf\{\alpha(y),\alpha(y')\}. \tag{2}
\end{align}
For $z\in X_0$ we write $z^{*}$ for the element of $X_0^{*}$ defined by
$z^{*}(z')=f(z,z')$.

% ---------------------------------------------------------------- 10.1.21
\st{Proposition (10.1.21).}
\emph{Conditions (i) to (iv) of Theorem (10.1.15) are further equivalent to the
following:}
\begin{itemize}
\item[(v)] \emph{for $(z,k)\in Z$ one has
$\omega_q^{*}(z^{*})\ge\tfrac12\omega(k)$;}
\item[(vi)] \emph{there exists a function $\omega_1:X_0\to\R\cup\{\infty\}$
such that, for $(z,k)\in Z$, one has $\omega_1(z)\ge\tfrac12\omega(k)$ and
$\omega_1^{*}(z^{*})\ge\tfrac12\omega(k)$;}
\item[(vii)] \emph{$\omega_q:X_0\to\R\cup\{\infty\}$ is an additive
semi-norm.}
\end{itemize}
\emph{These conditions imply that the set of $x\in X_0$ with
$\omega_q(x)=\infty$ is a vector subspace of $X^{\perp}$ (\btref{10.1.14}).}

\pf Condition (v) is merely a restatement of (10.1.15)(i). If (v) holds, so does
(vi) (take $\omega_1=\omega_q$). Conversely, if $\omega_1$ has the properties in
(vi), then $\omega_1\ge\omega_q$, whence $\omega_1^{*}\le\omega_q^{*}$, and for
every $(z,k)\in Z$
\[
\tfrac12\omega(k)\le\omega_1^{*}(z^{*})\le\omega_q^{*}(z^{*}).
\]

Let us show that (ii) implies (vii). The function $\omega_q$ obviously satisfies
relation (1) of (\btref{10.1.20}), so it suffices to establish (\btref{10.1.20})(2). Let
$(z,k),(z',k')\in Z$. By (ii), for all $x,x'\in\Kse$,
\[
\omega(f(z,z'))\ \ge\ \inf\{\omega(k+x),\ \omega(k'+x')\},
\]
whence
\begin{align*}
2\omega_q(z+z')
 &=\sup\{\omega(k+k'+f(z,z')+x)\mid x\in\Kse\}\\
 &\ge\sup\bigl\{\inf\{\omega(k+x),\omega(k'+x')\}\ \big|\ x,x'\in\Kse\bigr\}
   =\inf\{2\omega_q(z),\,2\omega_q(z')\}.
\end{align*}

To finish, we show that if (ii) fails then $\omega_q$ is not an additive
semi-norm. The set $\Kse$ is not dense in $K$; otherwise one would have
$x^{\sigma}=-\varepsilon x$ for all $x\in K$, whence $\sigma=\id$ and
$\varepsilon\neq1$, which would force $X_0=\{0\}$ and condition (ii) would hold
vacuously. Let $\lambda\in K-\overline{\Kse}$ and
$A=\sup\{\omega(\lambda+x)\mid x\in\Kse\}$. By (10.1.18)(iii) there exist
$(z,k),(z_1,k_1)\in Z$ with $\omega(k)>A$ and
$\omega(k_1)>A+2\omega(f(z,z_1))-2\omega(\lambda)$. Put
$\lambda_1=f(z,z_1)^{-1}.\lambda$, $z'=z_1.\lambda_1$ and
$k'=\lambda_1^{\sigma}.k_1.\lambda_1$. Then $(z',k')\in Z$,
$f(z,z')=\lambda$ and $\omega(k')>A$, whence
\begin{align*}
\omega_q(z+z')&=\tfrac12\sup\{\omega(\lambda+k+k'+x)\mid x\in\Kse\}\\
&=\tfrac12 A<\inf\Bigl\{\tfrac12\omega(k),\tfrac12\omega(k')\Bigr\}
\le\inf\{\omega_q(z),\omega_q(z')\},
\end{align*}
contradicting (vii). Finally, the last assertion follows from (10.1.15)(i),
taking into account that $f(X_0,X_i)=\{0\}$ for $i\neq0$.

% ---------------------------------------------------------------- 10.1.22
\st{Remarks (10.1.22).}
\emph{a)} Let $\omega_f:X_0\to\R\cup\{\infty\}$ be defined by
\[
\omega_f(z)=\tfrac12\omega(f(z,z))\qquad (z\in X_0).
\]
By (\btref{10.1.14})(3) one always has
\begin{equation}
\omega_f\ \ge\ \omega_q. \tag{1}
\end{equation}

\emph{b)} Write $Z(K)$ for the center of $K$. Suppose there is $\lambda\in Z(K)$
with $\lambda+\lambda^{\sigma}=1$. Then $\lambda f(z,z)\in q(z)$ for every
$z\in X_0$ and
\begin{equation}
\omega_q\ \ge\ \omega_f+\tfrac12\omega(\lambda). \tag{2}
\end{equation}
If $\lambda$ is integral --- which means either that the residue characteristic
of $K$ is different from $2$, or that $Z(K)$ is an unramified extension of the
field of fixed points of $\sigma$ in $Z(K)$ and the corresponding residue
extension is separable --- then (1) and (2) imply $\omega_q=\omega_f$.

\emph{c)} Suppose one of the following conditions holds:
\begin{itemize}
\item[(I)] $\dim X_0=1$;
\item[(II)] $K$ is commutative and $(\sigma,\operatorname{char}K)\neq(\id,2)$;
\item[(III)] $\varepsilon=1$, $K$ is a quaternion algebra of characteristic
$\neq2$, and $\sigma(k)=\overline k$ (the conjugate of $k$) for $k\in K$.
\end{itemize}
\emph{Then there exists $c\in\R\cup\{\infty\}$ with $\omega_f=\omega_q+c$, and
the function $\omega_f$ is an additive semi-norm if and only if $\omega_q$ is
one.}

In case (I) this is clear. In the other cases one may assume $\varepsilon=1$ by
(10.1.3); one then observes that
\begin{align}
&\Kse[\,\sigma,1\,]=\{k\mid k\in K,\ k+k^{\sigma}=0\} \tag{4}\\
\intertext{(whence, for $x\in X$, $f(x,x)=0$ implies $q(x)=0$), and that}
&f(z,z)\in Z(K)\qquad\text{for all } z\in X_0, \tag{5}
\end{align}
so that, for $(z,k)\in Z-\{(0,0)\}$,
\begin{align*}
\omega_q(z)-\omega_f(z)
&=\tfrac12\sup\{\omega(k+h)\mid h\in\mathrm{K}_{\sigma,1}\}
 -\tfrac12\omega(f(z,z))\\
&=\tfrac12\sup\{\omega(k.f(z,z)^{-1}+h)\mid h\in\mathrm{K}_{\sigma,1}\}
 =\tfrac12\sup\{\omega(\ell)\mid \ell\in K,\ \ell+\ell^{\sigma}=1\}.
\end{align*}
It is easy to see that if $\varepsilon=1$ and $X_0\neq0$, relations (4) and (5),
needed for the preceding argument, force one of the conditions (II), (III) to
hold.

\emph{d)} Suppose condition (10.1.15)(ii) fails and that there is $z\in X_0$
with $f(z,z)\neq0$. For such a $z$ let $k\in q(z)$ (so
$k+\varepsilon k^{\sigma}\neq0$) and choose $z_1,z_2\in X_0$ with
\[
\omega_f(z_1)>\tfrac12\omega(k+\varepsilon k^{\sigma})-\omega(k)+\omega(f(z_1,z_2))
\quad\text{and}\quad
\omega_f(z_2)>\tfrac12\omega(k+\varepsilon k^{\sigma})
\]
(the existence of such $z_i$ is guaranteed by (10.1.18)(iii)). Put
$z'_1=z_1.f(z_2,z_1)^{-1}.k$. A simple computation then shows that
$\omega_f(z'_1+z_2)<\inf\{\omega_f(z'_1),\omega_f(z_2)\}$. One concludes:

\emph{If the function $z\mapsto f(z,z)$ does not vanish identically on $X_0$ and
$\omega_f$ is an additive semi-norm, then the equivalent conditions (i) to (vii)
of (10.1.15) and (10.1.21) hold.}

\smallskip
The converse is false, as the following example shows.

\smallskip
Let $C$ be a local field of characteristic $0$ and residue characteristic $2$,
$\pi$ a uniformiser of $C$, and $L$ an unramified quadratic extension of $C$
whose residue extension is separable. Let $\ell\mapsto\overline\ell$ be the
$C$-automorphism of order $2$ of $L$, and let $d$ be an integer of $L$ with
$d+\overline d=1$ (such integers exist by the hypothesis on the residue
extension of $L/C$). Let $K=L(j)$ be the quaternion algebra defined by the
relations
\[
j\ell j^{-1}=\overline\ell\quad\text{for }\ell\in L,
\qquad j^{2}=\pi,
\]
let $\varepsilon=1$, let $\sigma:K\to K$ be the involution
$x+yj\mapsto\overline x+yj$ $(x,y\in L)$, let $X_0$ be the vector space $K^{2}$,
and let $f$ be the hermitian form defined by (\btref{10.1.1})(1) and
\[
f(z,z)=\zeta_1^{\sigma}\zeta_1-\zeta_2^{\sigma}(1+\pi j)\zeta_2
\qquad\text{for } z=(\zeta_1,\zeta_2)\in X_0 .
\]
Put $z_1=(\pi,0)$ and $z_2=(d,d)$. One finds
\[
f(z_1,z_1)=\pi^{2},\quad f(z_2,z_2)=-\overline d\,\pi j d,\quad f(z_1+z_2,z_1+z_2)=\pi+\pi^{2}-\overline d\,\pi j d,
\]
whence $\omega_f(z_1)=\omega(\pi)$, $\omega_f(z_2)=\tfrac34\omega(\pi)$ and
$\omega_f(z_1+z_2)=\tfrac12\omega(\pi)$. Consequently $\omega_f$ is not an
additive semi-norm. On the other hand, for condition (iii) of (10.1.15) to fail,
the function $z\mapsto f(z,z)$ must have a non-trivial zero in $X_0$, i.e.\
there must exist $\zeta\in K$ with
\begin{equation}
1+\pi j=\zeta^{\sigma}\zeta. \tag{6}
\end{equation}
Writing $\zeta=\xi+\eta j$ with $\xi,\eta\in L$, relation (6) is equivalent to
the pair
\begin{align}
\xi\overline\xi+\eta\overline\eta\,\pi&=1, \tag{7}\\
2\overline\xi\eta&=\pi. \tag{8}
\end{align}
From (7) it follows that $\xi$ and $\eta$ must be integral; one then deduces
from (8) that if $\zeta$ exists one must have $\omega(\pi)\ge\omega(2)$.
Choosing $C$ so that $\omega(\pi)<\omega(2)$, one obtains the announced example.

% ---------------------------------------------------------------- 10.1.23
\num{10.1.23}\btanchor{10.1.23} We shall see later ((\btref{10.1.34}), (\btref{10.1.36}) and (\btref{10.1.37})) that,
under certain conditions (for example when $r\ge2$), the valuations
$\varphi^{\omega}$ described by Theorem (10.1.15) are essentially (up to
equipollence) the only valuations of the root datum
$(T,(U_a)_{a\in\Phi})$. In the meantime we establish a few properties of the
valuations $\varphi^{\omega}$ in question.

\emph{Up to and including no.\ (\btref{10.1.33}), $\omega$ will denote a valuation of the
field $K$ satisfying conditions (i) to (iv) of Theorem (10.1.15), and
$\varphi=\varphi^{\omega}$ the corresponding valuation of the root datum
$(T,(U_a))$.} Note that a ``change of coordinates'' of the type described in
(10.1.3) has the effect of replacing $\varphi$ by an equipollent valuation.
\par\endgroup

% ---------------------------------------- bruhattits10.tex
\begingroup\setlength{\parindent}{1.5em}
\allowdisplaybreaks
\providecommand{\R}{}
\renewcommand{\R}{\mathbf{R}}
\providecommand{\Z}{}
\renewcommand{\Z}{\mathbf{Z}}
\providecommand{\Kse}{}
\renewcommand{\Kse}{\mathrm{K}_{\sigma,\varepsilon}}
\providecommand{\K}{}
\renewcommand{\K}{\mathrm{K}}
\providecommand{\Is}{}
\renewcommand{\Is}{\mathrm{Is}}
\providecommand{\Sim}{}
\renewcommand{\Sim}{\mathrm{Sim}}
\providecommand{\Isom}{}
\renewcommand{\Isom}{\mathrm{Isom}}
\providecommand{\Hom}{}
\renewcommand{\Hom}{\operatorname{Hom}}
\providecommand{\End}{}
\renewcommand{\End}{\operatorname{End}}
\providecommand{\Tr}{}
\renewcommand{\Tr}{\operatorname{Tr}}
\providecommand{\Ker}{}
\renewcommand{\Ker}{\operatorname{Ker}}
\providecommand{\id}{}
\renewcommand{\id}{\mathrm{id.}}
\providecommand{\ov}{}
\renewcommand{\ov}{\overline{\omega}}
\providecommand{\Bor}{}
\renewcommand{\Bor}{\mathscr{B}}
\providecommand{\Nor}{}
\renewcommand{\Nor}{\mathscr{N}}
\providecommand{\Eor}{}
\renewcommand{\Eor}{\mathscr{E}}
\providecommand{\Ech}{}
\renewcommand{\Ech}{\mathscr{E}}
\providecommand{\Build}{}
\renewcommand{\Build}{\mathscr{I}}
\providecommand{\Jcal}{}
\renewcommand{\Jcal}{\mathscr{J}}
\providecommand{\Ncal}{}
\renewcommand{\Ncal}{\mathscr{N}}
\providecommand{\Phat}{}
\renewcommand{\Phat}{\widehat{\mathrm{P}}}
\providecommand{\Nhat}{}
\renewcommand{\Nhat}{\widehat{\mathrm{N}}}
\providecommand{\What}{}
\renewcommand{\What}{\widehat{\mathrm{W}}}
\providecommand{\stmt}{}
\renewcommand{\stmt}[2]{\medskip\noindent\emph{#1} \textbf{(#2)}.\ ---\ }
\providecommand{\num}{}
\renewcommand{\num}[1]{\medskip\noindent\textbf{(#1)}\ }
\providecommand{\pf}{}
\renewcommand{\pf}{\par\medskip\noindent}
\begin{center}

\end{center}

\bigskip
\begin{center}\rule{0.5\textwidth}{0.4pt}\end{center}

% =====================================================================

% =====================================================================

\stmt{Proposition}{10.1.24}\btanchor{10.1.24}
\emph{The valuation $\varphi$ is extendable.}

\pf
If $r\neq 1$, this follows from (\btref{6.4.39}) and (\btref{6.4.35}). So suppose $r=1$. Let
$\mathrm{X}_{\pm 2}=e_{\pm 2}.\K$ be two one-dimensional vector spaces. Put
$\mathrm{X}'=\mathrm{X}_{-2}\oplus \mathrm{X}\oplus \mathrm{X}_{2}$ and let $f'$
and $q'$ be the forms obtained by extending $f$ and $q$ to
$\mathrm{X}'\times \mathrm{X}'$ and $\mathrm{X}'$ in such a way that the
relations (\btref{10.1.1})\,(12) remain satisfied. Let $\Phi'$,
$\widetilde{\mathrm{T}}{}'=\widetilde{\mathrm{T}}{}'(f',q')$, $\mathrm{U}'_a$
($a\in\Phi'$) and $\varphi'^{\,\omega}$ be defined in the same way as $\Phi$,
$\mathrm{T}(f,q)$, $\mathrm{U}_a$, $\varphi^\omega$, but substituting
$\mathrm{X}'$, $f'$, $q'$ for $\mathrm{X}$, $f$, $q$. Then
$\varphi'^{\,\omega}$ is, by what we have just seen, an extendable valuation of
the root datum $(\mathrm{T}',(\mathrm{U}'_a)_{a\in\Phi'})$; the assertion of the
statement follows from this one by restriction.

\num{10.1.25}\btanchor{10.1.25} \emph{The $\Gamma_a$.} --- It follows easily from the definitions
of (\btref{6.2.2}) that one has
\begin{align*}
\Gamma_{a_{ij}} &= \Gamma'_{a_{ij}} = \omega(\K^{*})\\[1ex]
\Gamma'_{a_i} &= \tfrac{1}{2}\bigl\{\omega(k)\ \big|\ k\in \mathrm{pr}_2 \mathrm{Z},\
   \omega(k)=\sup \omega(k+\Kse)\bigr\}
   \qquad\text{if } a_i\in\Phi\\[1ex]
\Gamma_{2a_i} &= \Gamma'_{2a_i} = \omega(\Kse-\{0\})
   \qquad\text{if } 2a_i\in\Phi .
\end{align*}

In particular, one sees that $\Gamma'_{a_i}$ is empty when $q|_{\mathrm{X}_0}$
takes its values in $\overline{\Kse}/\Kse$, and that if one of the conditions
(1), (2) or (3) of (10.1.3) is fulfilled, the valuation $\varphi$ is
\emph{special}.

If $\omega$ is \emph{discrete}, or if $\Kse$ is maximally complete, and if
$a_i\in\Phi$, one has
\[
\Gamma'_{a_i}=\omega_q(\mathrm{X}_0)-\{\infty\}.
\]

\num{10.1.26}\btanchor{10.1.26} \emph{Échelonnage.} --- If $\omega$ is a discrete valuation, the
discussion of no.~(\btref{6.2.23}) allows one to determine the échelonnage $\Ech$
associated with the valuation $\varphi^\omega$. We shall see, by means of
examples, that \emph{each of the échelonnages} $\mathrm{B}_r$, $\mathrm{C}_r$,
$\mathrm{D}_r$, $\mathrm{B}\text{-}\mathrm{C}_r$,
$\mathrm{C}\text{-}\mathrm{B}_r$, $\mathrm{B}\text{-}\mathrm{BC}_r$
\emph{and} $\mathrm{C}\text{-}\mathrm{BC}_r^{\mathrm{J}}$
($\mathrm{J}=\mathrm{I},\mathrm{II},\mathrm{III},\mathrm{IV}$)
\emph{can occur}.

Let us begin by examining the case where $\K$ \emph{is commutative}. We assume
that one of the conditions (1), (2), (3) of (10.1.3) is fulfilled, and we
normalize $\omega$ in such a way that $\omega(\K^{*})=\Z$.

\medskip
\emph{a)} Consider first the case where $\sigma=\id$. If $\Phi$ is of type
$\mathrm{C}$ (resp.\ $\mathrm{D}$), $\Ech$ is of type $\mathrm{C}_r$ (resp.\
$\mathrm{D}_r$). Suppose $\Phi$ of type $\mathrm{B}$; then $\Ech$ is of type
$\mathrm{B}_r$ if $\omega(q(\mathrm{X}_0-\{0\}))=2\Z$ (for example if
$\dim \mathrm{X}_0=1$) and of type $\mathrm{C}\text{-}\mathrm{B}_r$ if
$\omega(q(\mathrm{X}_0-\{0\}))=\Z$ (for example if $\dim \mathrm{X}_0=2$ and if
$q|_{\mathrm{X}_0}$ is the quadratic form $x^2+\pi y^2$, where $\pi$ is a
uniformiser of $\K$).

\medskip
\emph{b)} Suppose now that $\sigma\neq\id$, and let $\mathrm{L}$ be the field of
fixed points of $\sigma$. Since we are in case (1) of (10.1.3), we have
$\varepsilon=-1$ and $1\in\Kse$, whence $\Kse=\mathrm{L}$ ((\btref{10.1.1})\,(8)). Let
$\pi$ be a uniformiser of $\mathrm{L}$.

\smallskip
\emph{b\,1)} If $\mathrm{X}_0=\{0\}$, $\Ech$ is of type $\mathrm{C}_r$ or
$\mathrm{B}\text{-}\mathrm{C}_r$ according as the extension $\K/\mathrm{L}$ is
unramified ($\omega(\mathrm{L}^{*})=\Z$) or ramified
($\omega(\mathrm{L}^{*})=2\Z$).

\smallskip
\emph{b\,2)} Suppose $\mathrm{X}_0\neq\{0\}$. Since $\Kse=\mathrm{L}$ is closed
in $\K$, one has
\[
\omega_q(\mathrm{X}_0)-\{\infty\}=\Z,\qquad (1/2)\Z\qquad\text{or}\qquad (1/2)+\Z .
\]

If $\K/\mathrm{L}$ is unramified and if
$\omega_q(\mathrm{X}_0)-\{\infty\}=\Z$ or $(1/2)+\Z$ (for example if
$\dim \mathrm{X}_0=1$), $\Ech$ is of type
$\mathrm{C}\text{-}\mathrm{BC}_r^{\mathrm{IV}}$.

If $\K/\mathrm{L}$ is unramified and if
$\omega_q(\mathrm{X}_0)-\{\infty\}=(1/2)\Z$ (for example if
$\mathrm{X}_0=\K^2$ and if $q((x,y))=t(xx^\sigma+\pi yy^\sigma)+\mathrm{L}$,
with $t^\sigma=-t$ when $\operatorname{char}\K\neq 2$ and $t+t^\sigma=1$ when
$\operatorname{char}\K=2$), $\Ech$ is of type
$\mathrm{C}\text{-}\mathrm{BC}_r^{\mathrm{II}}$.

If $\K/\mathrm{L}$ is ramified, one necessarily has
$\omega_q(\mathrm{X}_0)-\{\infty\}=(1/2)+\Z$ and $\Ech$ is of type
$\mathrm{C}\text{-}\mathrm{BC}_r^{\mathrm{III}}$.

The examination of the commutative case is thus finished, and it remains for us
to give examples of échelonnages of types $\mathrm{B}\text{-}\mathrm{BC}_r$ and
$\mathrm{C}\text{-}\mathrm{BC}_r^{\mathrm{I}}$. Let $\K$ be a quaternion algebra
over a local field $\mathrm{C}(\K)$ of residue characteristic different from
$2$, $\sigma$ the involution $k\mapsto \bar k$, $\varepsilon=-1$, $\pi$ a
uniformiser of $\mathrm{C}(\K)$, and $c,d\in\K$ pure quaternions (elements
antisymmetric for $\sigma$) such that $\omega(c^2)=0$ and $d^2=\pi$ (such
elements always exist, for example, if $\mathrm{C}(\K)$ is locally compact). In
view of (10.1.22)\,\emph{b)}, one has
\[
\Gamma'_{a_i}=\omega_f(\mathrm{X}_0)-\{\infty\}.
\]

One thus sees that if $\dim \mathrm{X}_0=2$ (resp.\ $1$) and if $z\mapsto
f(z,z)$ is the antihermitian form $\bar x c x+\bar y d y$ (resp.\ $\bar x c x$,
resp.\ $\bar x d x$), the échelonnage $\Ech$ is of type
$\mathrm{C}\text{-}\mathrm{BC}_r^{\mathrm{I}}$ (resp.\
$\mathrm{B}\text{-}\mathrm{BC}_r$, resp.\
$\mathrm{C}\text{-}\mathrm{BC}_r^{\mathrm{III}}$).

\num{10.1.27}\btanchor{10.1.27} For $i,j\in[-r,+r]$, let
$\omega_i:\mathrm{X}_i\to\R\cup\{\infty\}$ and
\[
\omega_{ij}:\Hom(\mathrm{X}_j,\mathrm{X}_i)\to\R\cup\{\infty\}\cup\{-\infty\}
\]
be defined by
\begin{align*}
\omega_0 &= \omega_q,\\
\omega_i(e_i k) &= \omega(k)\qquad\text{for } i\in\mathrm{I}\text{ and } k\in\K,\\
\omega_{ij}(\alpha) &= \inf_{x\in \mathrm{X}_j}\bigl(\omega_i(\alpha(x))-\omega_j(x)\bigr)
   \qquad\text{for } \alpha\in\Hom(\mathrm{X}_j,\mathrm{X}_i),
\end{align*}
always with the convention $\infty-\infty=\infty$.

By means of the identifications of certain
$\Hom(\mathrm{X}_i,\mathrm{X}_j)$ with $\K$, $\mathrm{X}_0$ or
$\mathrm{X}_0^{*}$ introduced in no.~(\btref{10.1.1}), one has, for
$i,j\in\mathrm{I}$, $\omega_{ij}=\omega$, $\omega_{0i}=\omega_q$ and
$\omega_{i0}=\omega_q^{*}$. It follows that only the $\omega_{i0}$, for
$i\in[-r,+r]$, can take the value $-\infty$.

If $h,i,j\in[-r,+r]$ and $\alpha\in\Hom(\mathrm{X}_h,\mathrm{X}_i)$,
$\beta\in\Hom(\mathrm{X}_i,\mathrm{X}_j)$, and if $\omega_{ih}(\alpha)$,
$\omega_{ji}(\beta)\neq-\infty$, one has
\begin{equation*}
\omega_{jh}(\beta\circ\alpha)\ \geqslant\ \omega_{ji}(\beta)+\omega_{ih}(\alpha).
\tag{1}
\end{equation*}

For $g\in\Hom(\mathrm{X},\mathrm{X})$ we set furthermore
\[
\ov_{ij}(g)=\omega_{ij}(c_{ij}(g))
\qquad\text{and}\qquad
\ov(g)=\inf_{i,j\in[-r,+r]}\{\ov_{ij}(g)\}.
\]

It is clear that
$\ov:\Hom(\mathrm{X},\mathrm{X})\to\R\cup\{\infty\}\cup\{-\infty\}$ is an
additive semi-norm. For $g,g'\in\Hom(\mathrm{X},\mathrm{X})$ with $\ov(g)$,
$\ov(g')\neq-\infty$, one has, by virtue of (1),
\begin{equation*}
\ov(gg')\ \geqslant\ \ov(g)+\ov(g').
\tag{2}
\end{equation*}

\medskip
\noindent
\emph{Henceforth, and up to and including no.~(\btref{10.1.33}), we identify the affine
space $\mathrm{A}$ of the valuations equipollent to $\varphi$ with the dual
$\mathrm{V}$ of $\mathrm{V}^{*}$ (\btref{10.1.2}) in such a way that the valuation
$\varphi$ coincides with the point $0$.}

\num{10.1.28}\btanchor{10.1.28} \emph{The homomorphism $\nu$.}

\btunnum{Proposition}\emph{Let $n$ be an element of
$\mathrm{N}=\mathrm{N}^{0}.\mathrm{T}$, let $c=c(n)$ be its similitude ratio and
let $\nu$ be the homomorphism of $\mathrm{N}$ into the affine group of
$\mathrm{V}$ defined in (\btref{6.2.10}). Then:}

\begin{enumerate}[label=(\roman*),leftmargin=2.5em]
\item \emph{For every $i\in[-r,+r]$ there exists one and only one element $s(i)$
of $[-r,+r]$ such that $c_{s(i),i}(n)\neq 0$. One has $s(0)=0$ and
$s(-i)=-s(i)$.}
\item \emph{One has}
\begin{equation*}
\ov_{s(i),i}(n)+\ov_{s(-i),-i}(n)=\omega(c)\qquad\text{\emph{for } } i\in\mathrm{I};
\tag{1}
\end{equation*}
\begin{equation*}
\ov_{00}(n)=\tfrac{1}{2}\omega(c)\ \text{ \emph{if} }\ \omega_q(\mathrm{X}_0)\neq\{\infty\}
\qquad\text{\emph{and}}\qquad
\ov_{00}(n)=\infty\ \text{ \emph{if} }\ \omega_q(\mathrm{X}_0)=\{\infty\}.
\tag{2}
\end{equation*}
\item \emph{For $v\in\mathrm{V}$ and $i\in\mathrm{I}$ one has:}
\begin{align*}
a_{s(i)}\bigl(\nu(n)(v)\bigr)
&= a_i(v)+\ov_{s(i),i}(n)-\tfrac{1}{2}\omega(c)\\
&= a_i(v)+\tfrac{1}{2}\omega(c)-\ov_{s(-i),-i}(n).
\end{align*}
\end{enumerate}

\pf
(i) is an immediate consequence of (10.1.7).

If relation (1) is true for two elements of $\mathrm{N}$, it is also true for
their product. Since it is manifestly true for the $m_{ij}(k)$, for the
$m_i(z,k)$ and, by virtue of (10.1.5)\,(i), for the elements of $\mathrm{T}$, it
is true for every $n\in\mathrm{N}$. Relation (2) is obvious.

Finally, (iii) follows from a simple computation, taking (ii) into account.

\num{10.1.29}\btanchor{10.1.29} \emph{The group $\mathrm{H}$.} --- Let us denote by
$\Is_\omega(f,q)$ the group of all $(f,q)$-similitudes of $\mathrm{X}$ whose
ratio $c$ is a unit (i.e.\ such that $\omega(c)=0$).

\btunnum{Corollary}\emph{For $t\in\mathrm{T}$, the following
properties are equivalent:}
\begin{enumerate}[label=(\roman*),leftmargin=2.5em]
\item \emph{$t\in \mathrm{H}=\nu^{-1}(1)$;}
\item \emph{$\ov_{ii}(t)$ does not depend on $i$, for $i\in\mathrm{I}$.}
\end{enumerate}
\emph{If $\mathrm{G}\subset \Is_\omega(f,q)$, these conditions are furthermore
equivalent to:}
\begin{enumerate}[label=(\roman*),leftmargin=2.5em,start=3]
\item \emph{$\ov_{ii}(t)=0$ for every $i\in\mathrm{I}$.}
\end{enumerate}

\num{10.1.30}\btanchor{10.1.30} \emph{The groups $\mathrm{W}$ and $\What$.} --- Let us suppose the
valuation $\varphi$ special, which is not, as we have seen (\btref{10.1.25}), an
essential restriction. By (\btref{6.2.19}), each of the groups $\mathrm{W}$ and $\What$
is then the semi-direct product of its intersection with $\mathrm{V}$ and of the
Weyl group of $\Phi$ (regarded as a group of linear transformations of
$\mathrm{V}$). Let $v$ be an element of $\mathrm{V}$ and $v_i=a_i(v)$ its
coordinates. Then it again follows from (\btref{6.2.19}) that $v$ (or, more precisely,
the corresponding translation) belongs to $\mathrm{W}$ if and only if the
following conditions are fulfilled:
\[
v_i\in\omega(\K^{*})\qquad\text{for } i\in[1,r];
\]
$\sum_{i=1}^{r} v_i$ belongs to the group generated by the set
\[
\bigl\{\omega(k)\ \big|\ (z,k)\in \mathrm{Z}-\{(0,0)\}\bigr\}\ \cup\ 2\omega(\K^{*}).
\]

On the other hand, $v\in\What$ if and only if there exists $t\in\mathrm{T}$ such
that
\[
v_i=a_i(\nu(t))=\tfrac{1}{2}\bigl(\omega(c_{ii}(t))-\omega(c_{-i,-i}(t))\bigr).
\]
In particular, one sees that
\begin{equation*}
2(\What\cap \mathrm{V})\subset \mathrm{W}\cap \mathrm{V}.
\tag{1}
\end{equation*}

\num{10.1.31}\btanchor{10.1.31} \emph{Bornology.}

\btunnum{Proposition}\emph{Let $\Bor$ be the set of subsets
$\mathrm{M}$ of $\mathrm{G}$ such that $\ov(\mathrm{M})$ and
$\{\omega(c(g)^{-1})\mid g\in \mathrm{M}\}$ are bounded below. Then $\Bor$ is a
bornology on $\mathrm{G}$ compatible with $\varphi$ (\btref{8.1.1}). In particular,
$\ov(g)\neq-\infty$ for every $g\in\mathrm{G}$. The bornology defined by
$\varphi$ (\btref{8.1.7}) is the bornology generated by $\Bor.\mathrm{H}$. It coincides
with $\Bor$ if and only if $\mathrm{G}\subset \Is_\omega(f,q)$ (\btref{10.1.29}).}

\pf
It follows from (\btref{10.1.27})\,(2) that $\mathrm{M},\mathrm{M}'\in\Bor$ implies
$\mathrm{M}\mathrm{M}'\in\Bor$. On the other hand, one verifies at once, taking
(10.1.21)\,(v) into account, that the $\mathrm{U}_{a,k}$ for $a\in\Phi$ and
$k\in\R$ are bounded for $\Bor$. In view of (\btref{10.1.28}), every one-element subset
of $\mathrm{N}$ belongs to $\Bor$. Consequently $\Bor$ induces on $\mathrm{N}$
(resp.\ $\mathrm{T}$) a bornology $\Nor$ (resp.\ $\Eor$). For
$t\in\mathrm{T}$ and $i\in\mathrm{I}$ one has
$\ov_{00}(t)\geqslant\frac{1}{2}\omega(c(t))$ ((\btref{10.1.28})\,(2)) and
$c_{ii}(t^{-1})=c(t)^{-1}c_{-i,-i}(t)^{\sigma}$ (10.1.5). One deduces easily
that $\Eor$ is compatible with the group law of $\mathrm{T}$ and, since
$\mathrm{N}/\mathrm{T}$ is finite, that $\Nor$ is compatible with the group law
of $\mathrm{N}$. Moreover, (\btref{10.1.28})\,(iii) shows that the image under $\nu$ of
every bounded subset of $\mathrm{N}$ is bounded in $\Isom \mathrm{A}$, that is,
that $\Nor$ satisfies condition (BNC~1) of (\btref{8.1.2}). Consequently $\Bor$ is
weakly compatible with $\varphi$ ((\btref{8.1.6})\,\emph{a)}). Finally, one has
$\mathrm{H}_{[0]}\subset \mathrm{H}\cap \mathrm{G}^{0}\in\Bor$ by (\btref{10.1.29}).
This proves the first assertion.

Suppose that $\mathrm{H}\in\Bor$; let $t\in\mathrm{T}$ and $c=c(t)$. By virtue
of (\btref{10.1.30})\,(1) there exists $t_0\in\mathrm{T}^{0}$ (\btref{10.1.4}) such that
$t^{4}t_0^{-1}\in \mathrm{H}$. The set of powers of $h=t^{4}t_0^{-1}$ being
bounded for $\Bor$, one has, for every $i\in\mathrm{I}$, $\ov_{ii}(h)=0$, whence:
\[
h\in \Is_\omega(f,q),\qquad t\in \Is_\omega(f,q)
\qquad\text{and}\qquad
\mathrm{G}=\mathrm{T}.\mathrm{G}^{0}\cap \Is_\omega(f,q).\mathrm{G}^{0}=\Is_\omega(f,q).
\]

It then follows from the definition of $\Bor$ and from (\btref{10.1.28}) that the
bornology induced by $\Bor$ on $\mathrm{N}$ is the inverse image under $\nu$ of
the natural bornology of $\Isom \mathrm{A}$; consequently $\Bor$ is the
bornology $\Bor(\varphi)$ defined by $\varphi$.

Suppose now that $\mathrm{H}\notin\Bor$. In view of (\btref{10.1.29}) one then has
$\mathrm{G}\not\subset \Is_\omega(f,q)$. Moreover, by virtue of (\btref{10.1.28}),
$\omega(c(\mathrm{H}))\neq\{0\}$ and there exists a bounded subset $\mathrm{Y}$
of $\R$ such that $\mathrm{Y}+\omega(c(\mathrm{H}))=\R$. Let $\mathrm{L}$ be an
arbitrary subset of $\nu(\mathrm{T})$. There exists $\mathrm{M}\subset
\mathrm{T}$ such that $\nu(\mathrm{M})=\mathrm{L}$ and
$\omega(c(\mathrm{M}))\subset \mathrm{Y}$. If moreover $\mathrm{L}$ is bounded
in $\Isom \mathrm{A}$, it follows from (\btref{10.1.28}) that $\mathrm{M}\in\Bor$.
Corollary (\btref{8.1.9}) then shows that $\Bor(\varphi)$ is generated by
$\Bor.\mathrm{H}$, which completes the proof.

\num{10.1.32}\btanchor{10.1.32} \emph{The groups $\Phat_{x,\mathrm{E}}$.}

\btunnum{Proposition}\emph{Let $\mathrm{E}$ be a vectorial facet of
$\Phi$ and $x$ a point of $\mathrm{A}$. Then an element $g$ of $\mathrm{G}$
belongs to $\Phat_{x,\mathrm{E}}$ ((7.1.8), (7.2.4)) if and only if one has, for
$i,j\in[-r,+r]$,}
\begin{equation*}
\ov_{ij}(g)-(1/2)\omega(c(g))\ \geqslant\ a_i(x)-a_j(x)
\tag{1}
\end{equation*}
\emph{and}
\begin{equation*}
\ov_{ij}(g)-(1/2)\omega(c(g))\ >\ a_i(x)-a_j(x)
\qquad\text{\emph{when} } (a_i-a_j)(\mathrm{E})\subset\R^{*}_{+}.
\tag{2}
\end{equation*}

\pf
It follows immediately from (\btref{10.1.27})\,(1) that the set $\mathrm{L}$ of elements
$g\in\mathrm{G}$ which satisfy relations (1) and (2) is stable under
multiplication. On the other hand, it follows from (\btref{10.1.28})\,(2), (\btref{10.1.29}) and
(10.1.5)\,(i) that $\mathrm{H}\subset \mathrm{L}$, and one verifies easily that,
for every $a\in\Phi$, one has
\[
\mathrm{U}_a\cap \mathrm{L}=\mathrm{U}_a\cap \mathrm{P}_{x,\mathrm{E}}.
\]
In view of (7.2.6) we thus have $\mathrm{B}_{x,\mathrm{D}}\subset \mathrm{L}$
for every vector chamber $\mathrm{D}$ whose closure contains $\mathrm{E}$
and, in view of (\btref{7.3.4}), we have
\[
\mathrm{L}=\mathrm{B}_{x,\mathrm{D}}.(\mathrm{L}\cap \mathrm{N}).\mathrm{B}_{x,\mathrm{D}}.
\]
It therefore suffices to show that $\mathrm{L}\cap
\mathrm{N}=\Nhat_{x,\mathrm{E}}$. So let $n\in \mathrm{N}\cap \mathrm{L}$ and
let $i\in\mathrm{I}$. One has, with the notations of (\btref{10.1.28}),
\begin{equation*}
\ov_{s(i),i}(n)-(1/2)\omega(c(n))\ \geqslant\ a_{s(i)}(x)-a_i(x)
\tag{3}
\end{equation*}
\begin{equation*}
\ov_{s(-i),-i}(n)-(1/2)\omega(c(n))\ \geqslant\ a_{-s(i)}(x)-a_{-i}(x).
\tag{4}
\end{equation*}
Since $s(-i)=-s(i)$ (\btref{10.1.28}) and $a_{-j}=-a_j$ for every $j$, one deduces that
\begin{equation*}
\ov_{s(i),i}(n)-(1/2)\omega(c(n))\ =\ a_{s(i)}(x)-a_i(x)
\tag{5}
\end{equation*}
for every $i\in\mathrm{I}$, that is, by (\btref{10.1.28})\,(iii), that
\[
\nu(n)(x)=x
\]
or again $n\in\Nhat_{x}$. Moreover each of the inequalities (3) and (4) is an
equality, and one has, by (2),
\begin{equation*}
(a_{s(i)}-a_i)(\mathrm{E})=\{0\}
\tag{6}
\end{equation*}
whence, by (10.2.28)\,(iii) and (5), $n\in\Nhat_{x,\mathrm{E}}$.
% [TN] "(10.2.28)" is printed thus in the original; (10.1.28) is meant.
Conversely, if $n\in\Nhat_{x,\mathrm{E}}$, relations (3) to (6) are satisfied,
which shows that $n$ satisfies (1) and (2) when $i=s(j)$, with $j\in\mathrm{I}$.
When $i\neq s(j)$ one has $\ov_{ij}(n)=+\infty$, and when $i=j=0$ one has
$\ov_{00}(n)=(1/2)\omega(c(n))$ or $\infty$ ((\btref{10.1.28})\,(2)). Thus $n$ indeed
belongs to $\mathrm{L}$, which completes the proof.

\stmt{Corollary}{10.1.33}\btanchor{10.1.33}
\emph{Let $\Omega$ be a convex subset of $\mathrm{A}$. Then, in order that an
element $g$ of $\mathrm{G}$ belong to $\Phat_{\Omega}$, it is necessary and
sufficient that one have, for $i,j\in[-r,+r]$,}
\[
\ov_{ij}(g)\ \geqslant\ \sup\{a_i(x)-a_j(x)\mid x\in\Omega\}+(1/2)\omega(c(g)).
\]

% =====================================================================
%  From here on: statements translated in full, proofs given in
%  condensed/summary form rather than as a continuous translation.
% =====================================================================

\btanchor{10.1.34}

\noindent\small
From this point the file gives the \emph{statements} in full translation
together with a condensed account of each proof (the structure of the argument
and the relations it establishes), rather than a continuous rendering of the
original prose. See the note at the end of this document.
\normalsize

\stmt{Proposition}{10.1.34}
\emph{Suppose $r\geqslant 2$ and let $\mathrm{A}$ be an equipollence class of
valuations of the root datum $(\mathrm{T},(\mathrm{U}_a)_{a\in\Phi})$. Then
there exists one and only one function $\omega:\K\to\R\cup\{\infty\}$ such that
$\varphi^{\omega}$ (defined as in (\btref{10.1.13})) belongs to $\mathrm{A}$. This
function is a valuation of $\K$, invariant under $\sigma$.}

\pf
\emph{Summary of the proof (pp.~230).} The finite group $\mathrm{N}'_1$ of
(10.1.7)\,(ii) acts on $\mathrm{A}$ through $\nu$, and by (\btref{6.2.10})\,(ii) a point
$\varphi$ of $\mathrm{A}$ is $\mathrm{N}'_1$-invariant precisely when
\[
\varphi_{a_{ij}}(u_{ij}(1))=0\qquad (i,j\in\mathrm{I},\ j\neq\pm i).
\tag{$*$}
\]
Let $\varphi$ denote the (unique) fixed point. Uniqueness of $\omega$ follows
from ($*$) together with (\btref{6.2.1})\,(V\,1), (V\,5) and (\btref{10.1.11})\,(5),(6), which
force $\omega(-k)=\omega(k)$ and $\omega(1)=0$, hence $\varphi^{\omega}=\varphi$;
and $\omega$ is determined by $\varphi^{\omega}$ by (\btref{10.1.13})\,(3).

For existence one sets $\Psi_{ij}(k)=\varphi_{a_{ij}}(u_{ij}(k))$ for
$j\neq\pm i$, so that $\Psi_{ij}(1)=0$ by ($*$). Letting the group
$\mathrm{G}_{ij}=\langle \mathrm{U}_{a_{ij}},\mathrm{U}_{-a_{ij}}\rangle$ act on
$\mathrm{X}_{-i}+\mathrm{X}_j$ produces a homomorphism
$\mathrm{G}_{ij}\to \mathrm{SL}_2(\K)$ with
\[
u_{ij}(k)\mapsto\begin{pmatrix}1&0\\-\varepsilon(i)k&1\end{pmatrix},
\qquad
u_{-j,-i}(k)\mapsto\begin{pmatrix}1&-\varepsilon(-j)k\\0&1\end{pmatrix},
\]
and (\btref{6.2.12})\,\emph{b)} then shows that $\Psi_{ij}$ \emph{is} a valuation of
$\K$ and that $\Psi_{ij}=\Psi_{-j,-i}$. Comparing indices via (\btref{6.2.9}) and
(\btref{10.1.11})\,(12) (case $r\geqslant 3$) or (\btref{10.1.11})\,(10) (case $r=2$, where
$\mathrm{Z}\neq\{(0,0)\}$) gives $\Psi_{ih}=\Psi_{ij}$ and
$\Psi_{-i,j}=\Psi_{ij}$, so that $\Psi_{ij}$ is independent of the pair $(i,j)$;
by (\btref{10.1.11})\,(7) it is $\sigma$-invariant. Setting $\omega=\Psi_{ij}$, one
computes from (\btref{6.2.9}) and (\btref{10.1.11})\,(10) that, for $(z,k)\in\mathrm{Z}$,
\[
\varphi_{-a_i}(u_{-i}(z,k))=\tfrac12\omega(k)\ \ (a_i\in\Phi),
\qquad
\varphi_{-2a_i}(u_{-i}(0,k))=\omega(k)\ \ (2a_i\in\Phi),
\]
so that $\varphi=\varphi^{\omega}$. \qed

\stmt{Lemma}{10.1.35}\btanchor{10.1.35}
\emph{Suppose $r=1$ and $(\sigma,\varepsilon,\dim \mathrm{X}_0)\neq(\id,1,2)$.
Let $\varphi$ be a valuation of the root datum $(\mathrm{T},(\mathrm{U}_a))$ and
let $s\in\{1,2\}$ be such that $sa_1\in\Phi$. Then
$\frac{1}{s}\varphi_{sa_1}(u_1(z,k))$ depends only on $k$.}

\pf
\emph{Summary.} If $(\sigma,\varepsilon)=(\id,1)$: for proportional $z,z'$ one
has $u_1(z',k)=u_1(z,k)^{\pm1}$ and the claim follows from (\btref{6.2.1})\,(V\,1); for
non-proportional $z,z'$ one has $\dim \mathrm{X}_0>2$, so one may choose
$z''\in\mathrm{X}_0-\{0\}$ with $f(z,z'')=f(z',z'')=0$, and the identity
\[
m_1(z,k)\,u_1(z'',k'')\,m_1(z,k)^{-1} = u_{-1}(-z''k^{-1},\,k''k^{-2}) = m_1(z',k)\,u_1(z'',k'')\,m_1(z',k)^{-1}
\]
(with $k''=q(z'')$) gives $\nu(m_1(z,k))=\nu(m_1(z',k))$, whence the claim by
(\btref{6.2.10})\,(ii). If $(\sigma,\varepsilon)\neq(\id,1)$, the conjugation formula
$m_1(z,k)u_1(0,k')m_1(z,k)^{-1}=u_{-1}(0,(k^{\sigma})^{-1}k'k^{-1})$ expresses
$\frac1s\varphi_{sa_1}(u_1(z,k))$ in terms of $k$ alone. \qed

\stmt{Proposition}{10.1.36}\btanchor{10.1.36}
\emph{Let $r=1$. Suppose one of the conditions} (1), (2) \emph{of} (10.1.3)
\emph{is fulfilled. Put}
\[
\K_0=\bigcup_{z\in\mathrm{X}_0} q(z)=\{k\mid k\in\K,\ (z,k)\in\mathrm{Z}\},
\qquad
\K_1^{*}=c_{11}(\mathrm{T}\cap \Is(f,q)),
\qquad
\K_1=\K_1^{*}\cup\{0\}.
\]

\begin{enumerate}[label=(\roman*),leftmargin=2.5em]
\item \emph{One has $\K_0\subset\K_1$. If $(\sigma,\varepsilon)\neq(\id,1)$,
$\K_1$ contains the invariant elements of the center $\mathrm{Z}(\K)$ of $\K$.
If $(\sigma,\varepsilon)=(\id,1)$, $\K_1$ contains the elements of the form
$x^{\sigma}x=x^{2}$, for $x\in\K$. If $\mathrm{G}\supset \Is(f,q)$, then
$\K_1=\K$.}

\item \emph{Suppose $(\sigma,\varepsilon,\dim \mathrm{X}_0)\neq(\id,1,2)$ and
let $\mathrm{A}$ be an equipollence class of valuations of the root datum
$(\mathrm{T},(\mathrm{U}_a)_{a\in\Phi})$. Then there exists one and only one
function $\omega:\K\to\R\cup\{\infty\}$ such that}
\begin{equation*}
\omega(x^{\sigma}x)=2\omega(x)\qquad\text{\emph{for all} } x\in\K
\tag{1}
\end{equation*}
\emph{and such that $\varphi^{\omega}\in\mathrm{A}$. The valuation
$\varphi^{\omega}$ is special.}

\item \emph{Let $\omega:\K\to\R\cup\{\infty\}$ be a function satisfying relation}
(1)\emph{. Then, in order that $\varphi^{\omega}$ be a valuation of the root
datum $(\mathrm{T},(\mathrm{U}_a))$, it is necessary and sufficient that one
have, for $(z,k),(z',k')\in\mathrm{Z}$ and $t\in\mathrm{T}$,}
\begin{align}
\omega^{-1}(\infty)&=\{0\},\tag{2}\\
\omega(k+k'+f(z,z'))&\geqslant\inf(\omega(k),\omega(k')),\tag{3}\\
\omega(f(z,z')-f(z',z))&\geqslant\omega(k)+\omega(k')\tag{4}
\end{align}
\emph{and}
\begin{equation}
\omega\bigl(c_{-1,-1}(t)\,k\,c_{11}(t^{-1})\bigr)
=\omega(k)+\omega\bigl(c_{-1,-1}(t)\,c_{11}(t^{-1})\bigr).\tag{5}
\end{equation}
\emph{Under these conditions,}
\begin{equation}
\omega|_{\K_1^{*}}:\K_1^{*}\to\R \text{ \emph{is a homomorphism.}}\tag{5$'$}
\end{equation}

\item \emph{If $\mathrm{G}\subset \Is(f,q)$ or $\mathrm{G}\supset \Is(f,q)$,}
(5$'$) \emph{implies} (5)\emph{.}
\end{enumerate}

\pf
\emph{Summary of the proof (pp.~232--233).} Part (i) rests on $\K_1^{*}$ being a
group together with the three explicit similitude computations
\[
\begin{aligned}
c_{11}\bigl(m_1(e_0,1)m_1(z,k)\bigr)&=\varepsilon k,\\
c_{11}\bigl(m_1(0,1)m_1(0,\varepsilon\lambda)\bigr)&=\lambda,\\
c_{11}\bigl(m_1(e_0,1)m_1(e_0x,x^{2})\bigr)&=x^{2},
\end{aligned}
\]
the last assertion following from (10.1.5)\,(i).

For (ii) and the necessity in (iii): let $\varphi$ be the fixed point of
$\nu(\mathrm{N}')$ in $\mathrm{A}$ and suppose there is a (necessarily unique)
$\Psi:\K_0\to\R$ with $\frac1s\varphi_{sa_1}(u_1(z,k))=\Psi(k)$ --- guaranteed by
Lemma (\btref{10.1.35}) under the hypothesis on
$(\sigma,\varepsilon,\dim\mathrm{X}_0)$. Invariance of $\varphi$ under
$\nu(\mathrm{N}')$ extends this to $s\in\{-1,-2\}$ and gives $\Psi(1)=0$ and
$\Psi(\varepsilon k^{\sigma})=\Psi(k)$. Conjugating $u_1(z,k)$ by $t\in\mathrm{T}$
yields, with $\tau(t)=a_1(\nu(t))$,
\[
\Psi\bigl(c_{-1,-1}(t^{-1})\,k\,c_{11}(t)\bigr)=\Psi(k)+\tau(t),
\qquad
\tau(t)=\Psi\bigl(c_{-1,-1}(t^{-1})c_{11}(t)\bigr),
\]
and hence $\Psi(k'^{\sigma}kk')=\Psi(k)+\Psi(k'^{\sigma}k')$ for $k'\in\K_1$.
Defining $\omega(k'')=\frac12\Psi(k''^{\sigma}k'')$ and substituting
$k=k''^{\sigma}k''$ gives multiplicativity $\omega(k''k')=\omega(k'')+\omega(k')$
for $k'\in\K_1$, $k''\in\K$. Using $\Psi(\varepsilon k^{-1})=-\Psi(k)$ (from
(\btref{6.2.1})\,(V\,5) and (\btref{10.1.11})\,(3)) one checks $\Psi=\omega|_{\K_0}$, so
$\varphi=\varphi^{\omega}$. Relations (3), (4) are then the axioms (V\,1),
(V\,3); (5) comes from the two displayed formulas above; (5$'$) from
multiplicativity; and $\varphi^{\omega}$ is special because $\Psi(1)=0$.
Uniqueness follows as in (\btref{10.1.34}), using
$\{k^{\sigma}k\mid k\in\K\}\subset\K_0$. Sufficiency in (iii) is a direct
verification. Finally, $\mathrm{T}$ being normalized by $m_1(e_0,1)$ gives
$c_{-1,-1}(\mathrm{T})=c_{11}(\mathrm{T})$, and when
$\Is(f,q)\subset\mathrm{G}$ or $\supset\mathrm{G}$ one has
$c_{11}(\mathrm{T})=\K_1^{*}$, whence (5$'$)$\Rightarrow$(5). \qed

\stmt{Corollary}{10.1.37}\btanchor{10.1.37}
\emph{Keep the hypotheses of} (\btref{10.1.36})\,(ii)\emph{. Then $\omega$ is a
valuation of $\K$, invariant under $\sigma$, in each of the following cases:}
\begin{enumerate}[label=(\roman*),leftmargin=2.5em]
\item \emph{$\K$ is commutative and $\mathrm{X}_0\neq\{0\}$;}
\item \emph{$\K$ is a quaternion field, $k^{\sigma}=\Tr k-k$ for every $k\in\K$,
and $\mathrm{X}_0\neq\{0\}$;}
\item \emph{$\K=\K_0.\Kse$.}
\end{enumerate}

\pf
\emph{Summary.} Case (iii): the computation
$c_{11}\bigl(m_1(zk^{-1},(k^{\sigma})^{-1})\,m_1(0,\varepsilon k')\bigr)=kk'$
gives $\K_1=\K$, so $\omega|_{\K^{*}}$ is a homomorphism; taking $z'=0$ in
(\btref{10.1.36})\,(3) gives $\omega(k+k')\geqslant\inf(\omega(k),\omega(k'))$, so
$\omega$ is a valuation. Case (i): if
$(\sigma,\varepsilon)\neq(\id,1)$ then $\sigma\neq\id$ (else $f$ would be
alternating and $\mathrm{X}_0=\{0\}$), so $\varepsilon=-1$ and
$\Kse=\{t\in\K\mid t^{\sigma}=t\}$; since $q$ is not zero on $\mathrm{X}_0$,
$\K_0\not\subset\Kse$, whence $\K=\K_0.\Kse$ and case (iii) applies. If
$(\sigma,\varepsilon)=(\id,1)$ one applies (1) and (3) with $z=e_0x$,
$z'=e_0y$. Case (ii) reduces to $\dim\mathrm{X}_0=1$ and then to case (i) via
the classical isomorphism between a quaternionic antihermitian group in three
variables and a hermitian group in four variables over a commutative field.
\qed

\stmt{Remarks}{10.1.38}\btanchor{10.1.38}

\emph{a)} The conclusion of (\btref{10.1.36})\,(ii) is in general \emph{false} when
$(\sigma,\varepsilon,\dim \mathrm{X}_0)=(\id,1,2)$. Indeed, take
$\mathrm{G}=\mathrm{G}^{0}$ (\btref{10.1.4}) with $f$ nondegenerate and let $\mathrm{L}$
be the quadratic extension of $\K$ splitting $f$ (i.e.\ over which
$q|_{\mathrm{X}_0}$ acquires nontrivial zeros). There is then an isomorphism of
$\mathrm{G}$ onto $\mathrm{PSL}_2(\mathrm{L})$ carrying
$(\mathrm{T},(\mathrm{U}_a))$ onto the ``natural'' root datum of
$\mathrm{PSL}_2(\mathrm{L})$. The equipollence classes $\mathrm{A}$ therefore
correspond bijectively and naturally to the non-improper valuations $\omega_1$
of $\mathrm{L}$ (by (\btref{6.2.12})\,\emph{b)}), and $\mathrm{A}$ satisfies the
conclusion of (\btref{10.1.36})\,(ii) only when $\omega_1$ is invariant under the
automorphism of order $2$ of $\mathrm{L}$ fixing $\K$.

\emph{b)} Keep the hypotheses of (\btref{10.1.36})\,(ii), with $\K$ commutative and
$\mathrm{X}_0=\{0\}$. If $\sigma=\id$ then $\varepsilon\neq1$ by the conventions
of (\btref{10.1.1}), $\mathrm{G}^{0}$ with its root datum is isomorphic to
$\mathrm{SL}_2(\K)$, and (\btref{6.2.12})\,\emph{b)} shows that $\omega$ is a valuation
of $\K$. If on the other hand $\sigma\neq\id$, then $\mathrm{G}^{0}$ with its
root datum is isomorphic to $\mathrm{SL}_2(\mathrm{L})$, $\mathrm{L}$ being the
field of invariants of $\sigma$; the restriction $\omega_1$ of $\omega$ to
$\mathrm{L}$ is a valuation of $\mathrm{L}$ which may be chosen arbitrarily
((\btref{6.2.12})\,\emph{b)}), and $\omega$ is a valuation of $\K$ only if $\omega_1$
extends in a unique way to a valuation of $\K$ ((\btref{10.1.36})\,(1)).

\smallskip
The authors add that they know no example with $\mathrm{X}_0\neq\{0\}$ in which
$\omega$ fails to be a valuation of $\K$.

% =====================================================================

\btsecB{10.2}{10.2. The groups \texorpdfstring{$\mathrm{GL}_{n}(\mathrm{K})$}{GL\_n(K)} and \texorpdfstring{$\mathrm{SL}_{n}(\mathrm{K})$}{SL\_n(K)}}

% =====================================================================

\num{10.2.1}\btanchor{10.2.1} \emph{Setting.} Let $\K$ be a field, not necessarily commutative,
and let $\mathrm{X}$ be a right vector space of finite dimension $\geqslant 2$
over $\K$, equipped with a basis $(e_i)_{1\leqslant i\leqslant r+1}$. For an
endomorphism $g$ of $\mathrm{X}$, write $(g_{ij})_{1\leqslant i,j\leqslant r+1}$
for its matrix in this basis. For $1\leqslant i,j\leqslant r+1$, $i\neq j$, let
$\zeta_{ij}$ send $m=\left(\begin{smallmatrix}a&b\\c&d\end{smallmatrix}\right)\in
\mathrm{M}_2(\K)$ to the endomorphism $\zeta_{ij}(m)$ of $\mathrm{X}$ given by
\[
\zeta_{ij}(m).e_i=e_i a+e_j c,\qquad
\zeta_{ij}(m).e_j=e_i b+e_j d,\qquad
\zeta_{ij}(m).e_h=e_h\ \ (h\neq i,j).
\]
For $k\in\K$ one puts
\[
u_{ij}(k)=\zeta_{ij}\begin{pmatrix}1&k\\0&1\end{pmatrix} =\zeta_{ji}\begin{pmatrix}1&0\\k&1\end{pmatrix},
\qquad
m_{ij}(k)=\zeta_{ij}\begin{pmatrix}0&k\\-k^{-1}&0\end{pmatrix}=m_{ji}(-k^{-1})
\ \ (k\in\K^{*}).
\]
In $\R^{r+1}$ with its usual euclidean norm and canonical basis
$(a_i)_{1\leqslant i\leqslant r+1}$, put $a_{ij}=a_j-a_i$; the set $\Phi$ of the
$a_{ij}$ is an irreducible root system of type $\mathrm{A}_r$ in the space
$\mathrm{V}^{*}$ they span ([\btbib{5}], ch.~VI, p.~251). For $a_{ij}\in\Phi$ set
\[
\mathrm{U}_{a_{ij}}=\{u_{ij}(k)\mid k\in\K\},\qquad
\mathrm{M}^{0}_{a_{ij}}=\{m_{ij}(k)\mid k\in\K^{*}\}.
\]
The map $k\mapsto u_{ij}(k)$ is an injective homomorphism $\K\to\End(\mathrm{X})$,
and $\mathrm{U}_{a_{ij}}$ is the set of elements of $\End(\mathrm{X})$ of the
form $1+n$ with $n(e_h)=0$ for $h\neq j$ and $n(e_j)\in e_i\K$.

\stmt{Proposition}{10.2.2}\btanchor{10.2.2}
\emph{Let $\mathrm{G}$ be a subgroup of $\mathrm{GL}(\mathrm{X})$ containing the
subgroups $\mathrm{U}_{a_{ij}}$ for $1\leqslant i,j\leqslant r+1$, $i\neq j$,
and let $\mathrm{T}$ be the subgroup formed of the elements of $\mathrm{G}$
whose matrix relative to the basis $(e_i)$ is diagonal. Put
$\mathrm{M}_{a_{ij}}=\mathrm{T}.\mathrm{M}^{0}_{a_{ij}}$. Then
$(\mathrm{T},(\mathrm{U}_{a_{ij}},\mathrm{M}_{a_{ij}})_{a_{ij}\in\Phi})$ is a
generating root datum of type $\Phi$ in $\mathrm{G}$.}

\pf
\emph{Summary.} (DR~1) and (DR~3) are clear. (DR~2) follows from the commutator
formulas
\begin{equation*}
(u_{ij}(x),u_{h\ell}(y))=
\begin{cases}
1 & \text{if } i\neq\ell \text{ and } h\neq j,\\
u_{i\ell}(xy) & \text{if } i\neq\ell \text{ and } h=j,\\
u_{hj}(-xy) & \text{if } i=\ell \text{ and } h\neq j,
\end{cases}
\tag{1}
\end{equation*}
(the case $i=\ell$, $h=j$ being automatic). (DR~4) follows from
(\btref{6.1.3})\,\emph{a)}. Since $m_{ij}(k)$ exchanges the lines $e_i\K$, $e_j\K$ and
fixes the others, $m_{ij}(k)\mathrm{U}_{a_{h\ell}}m_{ij}(k)^{-1}
=\mathrm{U}_{a_{\tau(h)\tau(\ell)}}$ with $\tau=(i\,j)$, matching
$r_{a_{ij}}(a_{h\ell})=a_{\tau(h)\tau(\ell)}$; together with $\mathrm{T}$-invariance
of the $\mathrm{U}_{a_{h\ell}}$ this gives (DR~5). Taking $\Phi^{+}$ to be the
$a_{ij}$ with $i<j$, the groups $\mathrm{U}^{\pm}$ consist of unipotent
upper/lower triangular matrices, whence (DR~6). Consequently $\mathrm{G}^{0}$ is
normal in $\mathrm{GL}(\mathrm{X})$ and equals $\mathrm{SL}(\mathrm{X})$, the
group generated by all transvections (= kernel of the Dieudonné determinant)
[\btbib{21}]; since $\mathrm{GL}(\mathrm{X})=\widetilde{\mathrm{T}}.\mathrm{SL}(\mathrm{X})$
the datum is generating. In particular $\mathrm{G}$ may be \emph{any} subgroup
of $\mathrm{GL}(\mathrm{X})$ containing $\mathrm{SL}(\mathrm{X})$. \qed

\num{10.2.3}\btanchor{10.2.3} Let $\omega$ be a non-improper valuation of $\K$. For
$1\leqslant i,j\leqslant r+1$, $i\neq j$, and $k\in\K$, put
\[
\varphi^{\omega}_{a_{ij}}(u_{ij}(k))=\omega(k).
\]

\btunnum{Proposition}\emph{The family
$\varphi^{\omega}=(\varphi^{\omega}_a)_{a\in\Phi}$ is a special valuation of the
root datum $(\mathrm{T},(\mathrm{U}_a)_{a\in\Phi})$.}

\pf
\emph{Summary.} (V\,0), (V\,1), (V\,4) are immediate; (V\,3) follows from the
formulas (1) of (\btref{10.2.2}); (V\,5) and (V\,2) for $m\in\mathrm{M}^{0}_a$ follow
from (\btref{6.2.3})\,\emph{a)} via the homomorphisms
$\zeta_{ij}:\mathrm{SL}_2(\K)\to\mathrm{G}$. For (V\,2) with $t\in\mathrm{T}$ one
uses
\begin{equation*}
t\,u_{ij}(k)\,t^{-1}=u_{ij}(t_{ii}\,k\,t_{jj}^{-1}),
\tag{1}
\end{equation*}
which gives
$\varphi^{\omega}_{a_{ij}}(tut^{-1})=\varphi^{\omega}_{a_{ij}}(u)+\omega(t_{ii})-\omega(t_{jj})$,
together with $\varphi^{\omega}_{a_{ij}}(u_{ij}(1))=0$. \qed

\num{10.2.4}\btanchor{10.2.4} One proves as in (\btref{10.1.24}) that the valuation $\varphi^{\omega}$ is
\emph{extendable}. Moreover $\Gamma_{a_{ij}}=\Gamma'_{a_{ij}}=\omega(\K^{*})$
for every $a_{ij}\in\Phi$, and when $\omega$ is discrete the échelonnage $\Ech$
associated with $\varphi^{\omega}$ is of type $\mathrm{A}_r$.

\stmt{Proposition}{10.2.5}\btanchor{10.2.5}
\begin{enumerate}[label=(\roman*),leftmargin=2.5em]
\item \emph{The group $\mathrm{N}$ is the subgroup of $\mathrm{G}$ formed of the
elements whose matrix $(n_{ij})$ is monomial, and there is a surjective
homomorphism $s$ from $\mathrm{N}$ onto the symmetric group $\mathfrak{S}_{r+1}$
such that the only nonzero coefficients of the matrix of an element
$n\in\mathrm{N}$ are the coefficients $n_{s(n)(i),\,i}$.}
\item \emph{Let $\mathrm{A}$ be the affine space under the dual $\mathrm{V}$ of
$\mathrm{V}^{*}$ of the valuations equipollent to $\varphi^{\omega}$, and let
$\nu$ be the homomorphism from $\mathrm{N}$ into the affine group of
$\mathrm{A}$ defined in} (\btref{6.2.10})\emph{. Identify $\mathrm{A}$ and $\mathrm{V}$
by taking $\varphi^{\omega}$ as origin. Let $n\in\mathrm{N}$ and $\sigma=s(n)$.
For $x\in\mathrm{A}$ and $1\leqslant i,j\leqslant r+1$ one has}
\[
(a_{\sigma(i)\sigma(j)})(\nu(n).x)=(a_{ij})(x)+\omega(n_{\sigma(j),j})-\omega(n_{\sigma(i),i}).
\]
\item \emph{The subgroup $\mathrm{H}=\Ker\nu$ of $\mathrm{N}$ is the subgroup of
the $t\in\mathrm{T}$ such that $\omega(t_{ii})$ is independent of $i$ for
$1\leqslant i\leqslant r+1$.}
\end{enumerate}

\pf
\emph{Summary.} $\mathrm{N}$ is generated by $\mathrm{T}$ and the $m_{ij}(k)$,
and $s(m_{ij}(k))$ is the transposition of $i$ and $j$; (ii) is a direct
computation from (\btref{10.2.3})\,(1), and (iii) is an immediate consequence. \qed

\num{10.2.6}\btanchor{10.2.6} Let $\mathbf{W}$ be the affine Weyl group of $\Phi$ acting on
$\mathrm{A}$ (identified with $\mathrm{V}$ via $\varphi^{\omega}$), and
$\widetilde{\mathbf{W}}$ its normalizer in the affine group of $\mathrm{A}$
(cf.~(\btref{1.3.18})). Using (\btref{6.2.19}) one checks that the translation group
$\mathrm{W}\cap\mathrm{V}$ contained in
$\mathrm{W}=\nu(\mathrm{N}\cap\mathrm{G}^{0})$ is
$(\mathbf{W}\cap\mathrm{V})\otimes\omega(\K^{*})$; identifying $\mathrm{V}$ and
$\mathrm{V}^{*}$ by the scalar product induced from $\R^{r+1}$, it is the
restriction to $\mathrm{V}^{*}$ of the group of translations
$\sum_i\gamma_i a_i$ with $\gamma_i\in\omega(\K^{*})$ and $\sum_i\gamma_i=0$.

If $\mathrm{G}=\mathrm{GL}(\mathrm{X})$, then by (\btref{10.2.5})\,(ii) the translation
group of $\What=\nu(\mathrm{N})$ is identified with
$(\widetilde{\mathbf{W}}\cap\mathrm{V})\otimes\omega(\K^{*})$, i.e.\ with the
restrictions to $\mathrm{V}^{*}$ of the translations
$\sum_i\bigl(\gamma_i-(r+1)^{-1}\sum_j\gamma_j\bigr)a_i$ with
$\gamma_i\in\omega(\K^{*})$. One deduces that for any $\mathrm{G}$ with
$\mathrm{SL}(\mathrm{X})\subset\mathrm{G}\subset\mathrm{GL}(\mathrm{X})$,
\begin{equation*}
(r+1)(\What\cap\mathrm{V})\subset \mathrm{W}\cap\mathrm{V}.
\tag{1}
\end{equation*}

\num{10.2.7}\btanchor{10.2.7} Let $\mathrm{C}$ be the commutator group of $\K^{*}$ and
$\det:\mathrm{GL}(\mathrm{X})\to\K^{*}/\mathrm{C}$ the Dieudonné determinant
[\btbib{21}]. Since $\omega(\mathrm{C})=\{0\}$, the homomorphism $\omega:\K^{*}\to\R$
induces $\omega':\K^{*}/\mathrm{C}\to\R$; one sets $\delta=\omega'\circ\det$.

\btunnum{Proposition}\emph{Let $\Bor$ be the set of subsets
$\mathrm{M}$ of $\mathrm{G}$ such that the set of the $\omega(g_{ij})$ for
$g\in\mathrm{M}$, $1\leqslant i,j\leqslant r+1$, and the set of the $-\delta(g)$
for $g\in\mathrm{M}$ are both bounded below. Then $\Bor$ is a bornology
compatible with $\varphi^{\omega}$. The bornology defined by $\varphi^{\omega}$
is generated by $\Bor.\mathrm{H}$. It coincides with $\Bor$ if and only if
$\mathrm{G}\subset\Ker\delta$.}

\pf
\emph{Summary.} Stability under products and boundedness of the
$\mathrm{U}_{a,k}$ are clear (the Dieudonné determinant of a triangular matrix
is the image of the product of its diagonal coefficients). For
$\mathrm{M}\in\Bor$ with $\mathrm{M}\subset\mathrm{T}$, setting
$p=\inf\{\omega(t_{ii})\}$ and $q=\sup\{\delta(t)\}$ over $\mathrm{M}$, the
relations $-\delta(t^{-1})=\sum_i\omega(t_{ii})\geqslant(r+1)p$ and
$\omega((t^{-1})_{ii})=-\omega(t_{ii})\geqslant -q+rp$ give
$\mathrm{M}^{-1}\in\Bor$ and
$(a_i-a_j)(\nu(t))=\omega(t_{ii})-\omega(t_{jj})\leqslant q-(r+1)p$ by
(\btref{10.2.5})\,(ii). Hence $\Bor$ induces on $\mathrm{T}$, and so on $\mathrm{N}$, a
bornology compatible with the group law for which $\nu(\mathrm{M})$ is bounded
in $\Isom\mathrm{A}$; by (\btref{8.1.6})\,\emph{a)} this makes $\Bor$ weakly compatible
with $\varphi^{\omega}$. Finally
\[
\mathrm{H}_{[0]}\subset\mathrm{H}\cap\mathrm{SL}(\mathrm{X})\subset
\mathrm{H}\cap\Ker\delta=\{t\in\mathrm{T}\mid \omega(t_{ii})=0,\ 1\leqslant i\leqslant r+1\}\in\Bor .
\]
The proof is finished as in (\btref{10.1.31}), replacing $\Is(f,q)$ by
$\mathrm{SL}(\mathrm{X})$ and $\omega\circ c$ by $\delta$. \qed

\stmt{Proposition}{10.2.8}\btanchor{10.2.8}
\emph{Let $\mathrm{E}$ be a vectorial facet of $\Phi$ and $x$ a point of
$\mathrm{A}$. An element $g\in\mathrm{G}$ belongs to $\Phat_{x,\mathrm{E}}$}
(7.2.4) \emph{if and only if one has, for $1\leqslant i,j\leqslant r+1$:}
\begin{equation*}
\omega(g_{ij})+a_j(x)-a_i(x)\ \geqslant\ (r+1)^{-1}\delta(g)
\tag{1}
\end{equation*}
\emph{and}
\begin{equation*}
\omega(g_{ij})+a_j(x)-a_i(x)\ >\ (r+1)^{-1}\delta(g)
\qquad\text{\emph{when} } (a_i-a_j)(\mathrm{E})\subset\R^{*}_{+}.
\tag{2}
\end{equation*}

\pf
\emph{Summary.} As in (\btref{10.1.32}) it suffices to show
$\mathrm{N}\cap\mathrm{L}=\Nhat_{x,\mathrm{E}}$, where $\mathrm{L}$ is the set
of $g$ satisfying (1) and (2). For $n\in\mathrm{N}\cap\mathrm{L}$ with
$\sigma=s(n)$, summing (1)/(2) over $i=\sigma(j)$ gives
$\sum_j\omega(n_{\sigma(j),j})\geqslant\delta(n)$, which is an \emph{equality};
hence every one of those inequalities is an equality. By (\btref{10.2.5})\,(ii) this
gives $\nu(n).x=x$, and one must have
$(a_{\sigma(j)}-a_j)(\mathrm{E})\not\subset\R^{*}_{+}$ for every $j$, whence
$(a_{\sigma(j)}-a_j)(\mathrm{E})=\{0\}$ and $n\in\Nhat_{x,\mathrm{E}}$.
Conversely, for $n\in\Nhat_{x,\mathrm{E}}$ the numbers
$\omega(n_{\sigma(j),j})+a_j(x)-a_{\sigma(j)}(x)$ are all equal, giving (1); and
for $i\neq\sigma(j)$ one has $\omega(n_{ij})=+\infty$. \qed

\stmt{Corollary}{10.2.9}\btanchor{10.2.9}
\emph{Let $\Omega$ be a convex subset of $\mathrm{A}$. In order that an element
$g\in\mathrm{G}$ belong to $\Phat_{\Omega}$ it is necessary and sufficient that}
\[
\omega(g_{ij})\ \geqslant\ \sup\{a_i(x)-a_j(x)\mid x\in\Omega\}+(r+1)^{-1}\delta(g)
\qquad (1\leqslant i,j\leqslant r+1).
\]

\stmt{Proposition}{10.2.10}\btanchor{10.2.10}
\emph{Let $\psi$ be a valuation of the root datum
$(\mathrm{T},(\mathrm{U}_a)_{a\in\Phi})$ of $\mathrm{G}$. There exists one and
only one valuation $\omega$ of the field $\K$ such that $\psi$ is equipollent to
$\varphi^{\omega}$.}

\pf
The proof is entirely analogous to, and simpler than, that of (\btref{10.1.34}).

% =====================================================================
\btsecA{}{Note added in proof}
% =====================================================================

\noindent\small
\emph{The following is a condensed account of the note on pp.~238--239, not a
sentence-by-sentence translation.}
\normalsize

\medskip
\noindent
\textbf{(a) $\mathrm{GL}_{r+1}$ and the Goldman--Iwahori space of norms.}
When $\K$ is a complete, locally compact valued field, the building of
$\mathrm{SL}_{r+1}(\K)$ for the valuation $\varphi^{\omega}$ of (\btref{10.2.3}) is
closely related to the space of norms of O.~Goldman and N.~Iwahori
[\emph{Acta Math.} \textbf{109} (1963), 137--177]. In general, with the notation
of 10.2 and no extra hypothesis on $\K$, call an \emph{additive norm} $\gamma$
on $\mathrm{X}$ (i.e.\ the logarithm of the inverse of a norm in the usual
sense) \emph{decomposable} if there is a basis
$(u_i)_{1\leqslant i\leqslant r+1}$ of $\mathrm{X}$ and reals $c_i$ with
\begin{equation*}
\gamma\Bigl(\sum_{i=1}^{r+1}u_i x_i\Bigr)=\inf_i\{\omega(x_i)-c_i\}
\qquad (x_i\in\K),
\tag{1}
\end{equation*}
(every norm has this property if $\K$ is maximally complete), and call two norms
\emph{homothetic} if they differ by a constant. Then there is exactly one
$\mathrm{G}$-invariant bijection from the building $\Build$ of $\mathrm{G}$ (for
$\varphi^{\omega}$) onto the set of homothety classes of additive norms on
$\mathrm{X}$ carrying the class of the norm $\gamma$ of (1) to the point of
$\Build$ representing $\varphi^{\omega}+v$, where $v\in\mathrm{V}$ is given by
$a_i(v)=c_i$. In particular, for $\K$ locally compact, $\Build$ is the quotient
of the Goldman--Iwahori space $\Ncal(\mathrm{X})$ by homothety. The topology on
$\Build$ from the metric of (\btref{7.4.20}) (or (\btref{2.5.4})) is the quotient of the
Goldman--Iwahori topology, although their metric is unrelated to the one used
here. The authors note that it is sometimes useful to attach to a reductive
group over a local field a building which is the direct product of a building of
the present chapter with a euclidean affine space; for $\mathrm{GL}_{r+1}$ this
modified building is the space of norms itself, rather than its quotient by
homothety.

\medskip
\noindent
\textbf{(b) The other classical groups (after a suggestion of A.~Weil).}
Resume the notation and conventions of (\btref{10.1.1}) and (\btref{10.1.14}), assume conditions
(i)--(iv) of (10.1.15), and let $\Jcal$ be the set of additive norms $\alpha$ on
$\mathrm{X}$ satisfying
\begin{align}
\omega(q(x))&\geqslant 2\alpha(x) && (x\in\mathrm{X}),\tag{2}\\
\omega(f(x,y))&\geqslant \alpha(x)+\alpha(y) && (x,y\in\mathrm{X}),\tag{3}
\end{align}
where for $\bar t\in\K/\Kse$ one puts $\omega(\bar t)=\sup\{\omega(t)\mid t\in\bar t\}$.
Then there is exactly one $\mathrm{G}$-invariant injection $\iota$ of the
building $\Build$ of $\mathrm{G}$ (for $\varphi^{\omega}$) into $\Jcal$ such
that, for $v\in\mathrm{V}$ given by $a_i(v)=c_i$, $\iota(\varphi^{\omega}+v)$ is
the norm
\begin{equation*}
\gamma\Bigl(\sum_{i=1}^{r} e_i x_i+x_0\Bigr)
=\inf\Bigl\{\omega(x_i)-c_i,\ \tfrac12\omega(q(x_0))\Bigr\}.
\tag{4}
\end{equation*}

If moreover $\mathrm{X}$ is finite-dimensional and $\K$, $\Kse$ are maximally
complete for $\omega$, $\omega|_{\Kse}$, then $\iota(\Build)$ is exactly the set
of maximal elements of $\Jcal$, and every element of $\Jcal$ is minorised by an
element of $\iota(\Build)$.

\emph{Sketch of the argument.} Call subspaces $\mathrm{Y},\mathrm{Y}'$ of
$\mathrm{X}$ \emph{orthogonal for $\alpha$} if
$\alpha(y+y')=\inf\{\alpha(y),\alpha(y')\}$ for $y\in\mathrm{Y}$,
$y'\in\mathrm{Y}'$. Put $\mathrm{X}_{+}=\sum_{i=1}^{r}\mathrm{X}_i$ and choose
$\alpha$-orthogonal supplements $\mathrm{X}'_{-}$ of
$\mathrm{X}_0+\mathrm{X}_{+}$ and $\mathrm{X}''$ of $\mathrm{X}_{+}$ in
$\mathrm{X}$. Setting
$\alpha^{*}(x_{+})=\sup\{f(x_{+},x')-\alpha(x')\mid x'\in\mathrm{X}'_{-}\}$,
the norm $x_{+}+x''\mapsto\inf\{\alpha^{*}(x_{+}),\alpha(x'')\}$ majorises
$\alpha$ and lies in $\Jcal$, so one may assume
$\alpha^{*}=\alpha|_{\mathrm{X}_{+}}$. Take an $\alpha$-orthogonal basis
$(u_i)_{1\leqslant i\leqslant r}$ of $\mathrm{X}_{+}$ and the dual basis
$(u'_{-i})$ of $\mathrm{X}'_{-}$ with $f(u_i,u'_{-j})=\delta_{ij}$, and for each
$i$ pick $\lambda_i\in q(u'_{-i})$ with $\omega(\lambda_i)=\omega(q(u'_{-i}))$
(possible by the hypothesis on $\Kse$). Then
$u_{-i}=u'_{-i}-\varepsilon u_i\lambda_i$ spans a space
$\mathrm{X}_{-}=\sum_i u_{-i}\K$ also orthogonal to $\mathrm{X}_{+}+\mathrm{X}_0$;
choosing it as $\mathrm{X}'_{-}$ one may, after transforming $\alpha$ and the
data by a suitable element of $\mathrm{G}$, assume $e_i=u_i$ for
$i\in\mathrm{I}$. Writing (2) for $x=x_{+}+x_{-}+x_0$ and (3) for that $x$
against $\mathrm{X}_{+}$ and $\mathrm{X}_{-}$ shows that $\alpha$ minorises the
norm $\gamma$ of (4) with $c_i=-\omega(e_i)$, and this $\gamma$ is easily seen
to be a maximal element of $\Jcal$.
% [TN] The original prints "on a u'_{-i}=u_i" and "pour y in Y_+ ou Y_-";
% X_+ / X_- are evidently intended in the latter.

% =====================================================================
\btsecA{}{Bibliography}
% =====================================================================
\small
\begin{enumerate}[leftmargin=3.4em,labelwidth=3em,labelsep=0.4em,align=left,itemsep=1pt]
\item[{[1]}]\hypertarget{bib:1}{} A.~\textsc{Borel}, Groupes algébriques linéaires, \emph{Ann.\ of Math.}, \textbf{2} (1956), 20--80.
\item[{[2]}]\hypertarget{bib:2}{} ---, \emph{Linear algebraic groups}, New York, Benjamin, 1969.
\item[{[2 bis]}]\hypertarget{bib:2bis}{} --- et J.-P.~\textsc{Serre}, Cohomologie à support compact des immeubles de Bruhat--Tits; applications à la cohomologie des groupes S-arithmétiques, \emph{C.\ R.\ Acad.\ Sci.}, \textbf{272} (1971), 110--113.
\item[{[3]}]\hypertarget{bib:3}{} --- et J.~\textsc{Tits}, Groupes réductifs, \emph{Publ.\ I.H.É.S.}, \textbf{27} (1965), 55--150.
\item[{[4]}]\hypertarget{bib:4}{} N.~\textsc{Bourbaki}, \emph{Espaces vectoriels topologiques}, chapitres I et II, Paris, Hermann, 1966.
\item[{[5]}]\hypertarget{bib:5}{} ---, \emph{Groupes et algèbres de Lie}, chapitres IV, V et VI, Paris, Hermann, 1968.
\item[{[6]}]\hypertarget{bib:6}{} ---, \emph{Algèbre commutative}, chapitres V et VI, Paris, Hermann, 1964.
\item[{[7]}]\hypertarget{bib:7}{} F.~\textsc{Bruhat}, Sur les représentations des groupes classiques $\mathfrak{p}$-adiques, \emph{Amer.\ J.\ of Math.}, \textbf{83} (1961), 321--338 et 343--368.
\item[{[8]}]\hypertarget{bib:8}{} ---, Sur une classe de sous-groupes compacts maximaux des groupes de Chevalley sur un corps $\mathfrak{p}$-adique, \emph{Publ.\ I.H.É.S.}, \textbf{23} (1964), 46--74.
\item[{[9]}]\hypertarget{bib:9}{} ---, $\mathfrak{p}$-adic groups, in \emph{Algebraic groups and discontinuous subgroups}, Proc.\ Symp.\ in Pure Math., IX, A.M.S., 1966, 63--70.
\item[{[10]}]\hypertarget{bib:10}{} ---, Groupes semi-simples sur un corps local, \emph{Actes du Congrès International des Mathématiciens} (Nice, 1970), Paris, Gauthier-Villars, 1971.
\item[{[11]}]\hypertarget{bib:11}{} --- et J.~\textsc{Tits}, Un théorème de point fixe, \emph{Publ.\ I.H.É.S.}, 1966 (multigraphié).
\item[{[12]}]\hypertarget{bib:12}{} --- et ---, BN-paires de type affine et données radicielles affines, \emph{C.\ R.\ Acad.\ Sci.}, \textbf{263} (1966), 598--601.
\item[{[13]}]\hypertarget{bib:13}{} --- et ---, Groupes simples résiduellement déployés sur un corps local, \emph{C.\ R.\ Acad.\ Sci.}, \textbf{263} (1966), 766--768.
\item[{[14]}]\hypertarget{bib:14}{} --- et ---, Groupes algébriques simples sur un corps local, \emph{C.\ R.\ Acad.\ Sci.}, \textbf{263} (1966), 822--825.
\item[{[15]}]\hypertarget{bib:15}{} --- et ---, Groupes algébriques simples sur un corps local, cohomologie galoisienne (décompositions d'Iwasawa et de Cartan), \emph{C.\ R.\ Acad.\ Sci.}, \textbf{263} (1966), 867--869.
\item[{[16]}]\hypertarget{bib:16}{} --- et ---, Groupes algébriques simples sur un corps local, \emph{Proc.\ Conf.\ on Local Fields} (Driebergen, 1966), Springer, 1967, 23--36.
\item[{[17]}]\hypertarget{bib:17}{} C.~\textsc{Chevalley}, Sur certains groupes simples, \emph{Tôhoku Math.\ J.}, \textbf{7} (1955), 14--66.
\item[{[18]}]\hypertarget{bib:18}{} ---, Classification des groupes de Lie algébriques, \emph{Séminaire E.N.S.}, Paris, 1956--1958 (multigraphié).
\item[{[19]}]\hypertarget{bib:19}{} M.~\textsc{Demazure}, Schémas en groupes réductifs, \emph{Bull.\ Soc.\ Math.\ Fr.}, \textbf{93} (1965), 369--414.
\item[{[20]}]\hypertarget{bib:20}{} --- et A.~\textsc{Grothendieck}, Schémas en groupes, \emph{Sém.\ Géom.\ Alg.\ I.H.É.S.}, 1963--64.
\item[{[21]}]\hypertarget{bib:21}{} J.~\textsc{Dieudonné}, \emph{La géométrie des groupes classiques}, Springer, 1963.
\item[{[22]}]\hypertarget{bib:22}{} \textsc{Euclide}, \emph{\OE uvres}, Paris, Patris, 1814.
\item[{[23]}]\hypertarget{bib:23}{} \textsc{Harish-Chandra}, Harmonic analysis on $p$-adic reductive groups, \emph{Lecture Notes}, \textbf{162}, Springer, 1970.
\item[{[24]}]\hypertarget{bib:24}{} S.~\textsc{Helgason}, \emph{Differential geometry and symmetric spaces}, New York, Academic Press, 1962.
\item[{[25]}]\hypertarget{bib:25}{} H.~\textsc{Hijikata}, Maximal compact subgroups of $p$-adic classical groups (in Japanese), \emph{Sugaku no Ayumi}, \textbf{10-2} (1963), 12--23.
\item[{[26]}]\hypertarget{bib:26}{} ---, \emph{On the arithmetic of $p$-adic Steinberg groups}, Yale University, 1964 (mimeographed).
\item[{[27]}]\hypertarget{bib:27}{} ---, \emph{Maximal compact subgroups of some $p$-adic classical groups}, Yale University, 1964 (mimeographed).
\item[{[28]}]\hypertarget{bib:28}{} N.~\textsc{Iwahori} et H.~\textsc{Matsumoto}, On some Bruhat decomposition and the structure of the Hecke ring of $p$-adic Chevalley groups, \emph{Publ.\ I.H.É.S.}, \textbf{25} (1965), 5--48.
\item[{[29]}]\hypertarget{bib:29}{} I.~\textsc{Macdonald}, Harmonic analysis on semi-simple $p$-adic groups, \emph{Actes du Congrès International des Mathématiciens} (Nice, 1970), Paris, Gauthier-Villars, 1971.
\item[{[30]}]\hypertarget{bib:30}{} H.~\textsc{Matsumoto}, Fonctions sphériques sur un groupe semi-simple $p$-adique, \emph{C.\ R.\ Acad.\ Sci.}, \textbf{269} (1969), 829--832.
\item[{[31]}]\hypertarget{bib:31}{} \textsc{Pappus d'Alexandrie}, \emph{Mathematicae collectiones}, Bologne, H.~H.~de Ducciis, 1660.
\item[{[32]}]\hypertarget{bib:32}{} I.~\textsc{Satake}, Theory of spherical functions on reductive algebraic groups over $p$-adic fields, \emph{Publ.\ I.H.É.S.}, \textbf{18} (1963), 5--69.
\item[{[33]}]\hypertarget{bib:33}{} G.~\textsc{Schiffmann}, Intégrales d'entrelacement et fonctions de Whittaker, \emph{Bull.\ Soc.\ Math.\ Fr.}, \textbf{99} (1971), 3--72.
\item[{[34]}]\hypertarget{bib:34}{} O.~\textsc{Schilling}, Theory of valuations, \emph{Math.\ Surveys}, IV, New York, 1950.
\item[{[35]}]\hypertarget{bib:35}{} J.-P.~\textsc{Serre}, Cohomologie des groupes discrets, \emph{C.\ R.\ Acad.\ Sci.}, \textbf{268} (1969), 268--271.
\item[{[36]}]\hypertarget{bib:36}{} J.~\textsc{Tits}, Algebraic and abstract simple groups, \emph{Ann.\ of Math.}, \textbf{80} (1964), 313--329.
\item[{[37]}]\hypertarget{bib:37}{} ---, Structures et groupes de Weyl, \emph{Séminaire Bourbaki}, \textbf{288} (février 1965).
\item[{[38]}]\hypertarget{bib:38}{} ---, Formes quadratiques, groupes orthogonaux et algèbres de Clifford, \emph{Inv.\ Math.}, \textbf{5} (1968), 19--41.
\item[{[39]}]\hypertarget{bib:39}{} ---, Buildings of spherical type and finite BN-pairs, \emph{Lecture Notes}, Springer (to appear).
\end{enumerate}
\normalsize

\medskip
\noindent\emph{Manuscript received 7 June 1971.}

% =====================================================================
\btsecA{}{Index of notations}
% =====================================================================
\small
\begin{multicols}{2}
\raggedright
(\btref{1.1.5}) $\dim \mathrm{A}$, $\operatorname{codim}\mathrm{F}$.\\
(\btref{1.2.2}) $\ell(w)$, $\mathrm{S}_w$.\\
(\btref{1.2.4}) $\mathrm{W}_{\mathrm{X}}$.\\
(\btref{1.2.5}) $\mathrm{Cox}(\mathrm{W},\mathrm{S})$, $\mathrm{Cox}\,\mathrm{W}$.\\
(\btref{1.2.9}) $\mathrm{B}_{\mathrm{X}}$.\\
(\btref{1.2.12}) $\mathrm{H}$, $\widetilde{\mathrm{N}}$.\\
(\btref{1.2.13}) $\widehat{\mathrm{G}}$.\\
(\btref{1.2.15}) $^{g}\mathrm{P}$, $\operatorname{Stab}\mathrm{P}$.\\
(\btref{1.2.16}) $\xi$, $\Xi$.\\
(\btref{1.2.17}) $\widehat{\mathrm{G}}_0$, $\widehat{\mathrm{B}}$, $\widehat{\mathrm{P}}$.\\
(\btref{1.2.18}) $\Xi_{\mathrm{X}}$.\\
(\btref{1.3.1}) $\mathbf{A}$ (cf.\ \btref{2.8} Introd.); $^{v}\mathbf{A}$, $\mathbf{W}$, $^{v}w$.\\
(\btref{1.3.2}) $\mathbf{A}_i$, $\mathbf{W}_i$, $^{v}\mathbf{A}_i$.\\
(\btref{1.3.3}) $s_{\mathrm{L}}$, $\mathrm{L}_s$; $\boldsymbol{\Sigma}$; $r_\alpha$, $\partial\alpha$, $\alpha^{*}$, $\alpha_{+}$.\\
(\btref{1.3.4}) $\mathbf{C}$, $\mathbf{S}$, $\mathbf{C}_i$, $\mathbf{S}_i$.\\
(\btref{1.3.5}) $\mathbf{C}_{\mathrm{X}}$, $\mathbf{T}$.\\
(\btref{1.3.6}) $s(\Gamma)$ (cf.\ (\btref{2.1.3})).\\
(\btref{1.3.7}) $\mathbf{W}_x$, $\mathrm{V}$, $^{v}\mathbf{W}$.\\
(\btref{1.3.8}) $^{v}\boldsymbol{\Sigma}$, $^{v}\alpha$; $\mathrm{L}_{a,k}$, $\alpha_{a,k}$, $r_{a,k}$; $r_a$, $a^{\vee}$.\\
(\btref{1.3.12}) $\mathrm{B}(\mathrm{D})$, $^{v}\boldsymbol{\Sigma}^{+}_{\mathrm{D}}$, $^{v}\boldsymbol{\Sigma}^{+}(\mathrm{D})$.\\
(\btref{1.3.16}) $\leqslant_{\mathrm{D}}$.\\
(\btref{1.3.18}) $\widetilde{\mathbf{W}}$, $\widetilde{\mathbf{V}}$, $\widetilde{\mathbf{W}}^{\dagger}$, $\widetilde{\mathbf{W}}_{\mathrm{int}}$.\\
(\btref{1.4.1}) $\Ech$.\\
(\btref{1.4.2}) $\Phi_x$, $\Phi_{\mathrm{F}}$.\\
(\btref{1.4.4}) $\operatorname{Dyn}\Phi$.\\
(\btref{1.4.5}) $\operatorname{Dyn}\Ech$; $\mathrm{A}_n$, $\mathrm{B}_n,\dots$, $\mathrm{B}\text{-}\mathrm{C}_n$, $\mathrm{C}\text{-}\mathrm{B}_n^{\mathrm{X}}$, etc.\\
(\btref{1.5.2}) $\mathrm{G}$, $\mathrm{B}$, $\mathrm{N}$, $\mathrm{S}$.\\
(\btref{2.1.1}) $\mathscr{P}$, $\tau(\mathrm{P})$, $\mathrm{F}(\mathrm{P})$; $\Build$ (cf.\ (\btref{7.4.2})).\\
(\btref{2.1.7}) $w(\mathrm{C},\mathrm{C}')$.\\
(\btref{2.2.1}) $j$.\\
(\btref{2.2.4}) $\partial \mathrm{M}$.\\
(\btref{2.2.8}) $d_{\mathrm{A}}$ (cf.\ (\btref{7.4.11})).\\
(\btref{2.3.2}) $\rho_{\mathrm{A}',\mathrm{A};\mathfrak{C}}$.\\
(\btref{2.3.5}) $\rho_{\mathrm{A};\mathrm{C}}$ (cf.\ (\btref{7.4.19}), (\btref{7.5.7})).\\
(\btref{2.3.10}) $\mathrm{T}_w$.\\
(\btref{2.4.2}) $\operatorname{cl}(\mathrm{M})$ (cf.\ (\btref{2.4.6}), (\btref{2.5.7}), (7.1.2) and (\btref{7.4.11})).\\
(\btref{2.5.6}) $[xy]$ (cf.\ (\btref{7.4.20}), (\btref{7.5.1})).\\
(\btref{2.5.13}) $tx+(1-t)y$ (cf.\ (\btref{1.1.1}), (\btref{7.4.20}) and (\btref{7.5.1})).\\
(\btref{2.6.2}) $\mathrm{B}^{\mathrm{J}}$.\\
(\btref{2.7.4}) $\operatorname{Aut}_{\mathrm{B}}\mathrm{G}$, $\operatorname{Aut}_{(\mathrm{B},\mathrm{N})}\mathrm{G}$.\\
(\btref{2.9.1}) $\rho_{\mathrm{A};\mathfrak{C}}$.\\
(\btref{3.1.2})\,b) $\Isom \mathrm{E}$.\\
4.\ Introd.\ $\varphi:\mathrm{G}\to\widehat{\mathrm{G}}$.\\
(\btref{4.1.1}) $\Phat_{\Omega}$, $\Phat^{\dagger}_{\Omega}$, $\mathrm{P}_{\Omega}$, $\mathrm{P}^{\dagger}_{\Omega}$ (cf.\ (7.1.1), (7.1.8) and (\btref{7.3.9})); $\mathrm{P}_x$, etc.\ (cf.\ (7.1.3)).\\
(\btref{4.1.2}) $\Nhat$, $\widehat{\nu}$, $\widehat{\mathbf{W}}$, $\widehat{\mathbf{V}}$, $\widehat{\mathrm{H}}$.\\
(\btref{4.1.4}) $\mathrm{N}^{\dagger}_{\Omega}$, $\Nhat_{\Omega}$ (cf.\ (7.1.8)).\\
(\btref{4.1.5}) $\widehat{\mathrm{V}}_{\mathrm{D}}$, $\mathrm{E}_{\mathrm{D}}$; $\widehat{\mathfrak{P}}^{0}_{\mathrm{D}}$, $\widehat{\mathfrak{P}}_{\mathrm{D}}$, $\mathfrak{B}^{0}_{\mathrm{D}}$, $\mathfrak{B}_{\mathrm{D}}$; $\mathbf{D}$, $\widehat{\mathfrak{P}}^{0}$, $\widehat{\mathfrak{P}}$, $\mathfrak{B}^{0}$, $\mathfrak{B}$.\\
4.2.\ Introd.\ $\widehat{\mathbf{W}}_{\mathrm{Q}}$.\\
(\btref{5.1.26}) $\mathrm{U}_\alpha$ (cf.\ (\btref{6.2.6})).\\
(\btref{5.1.33}) $^{v}\Build$.\\
(\btref{5.2.4}) $\widetilde{\boldsymbol{\Sigma}}(\Omega)$, $\boldsymbol{\Sigma}(\Omega)$, $^{v}\boldsymbol{\Sigma}(\Omega)$, $a(\Omega)$.\\
6.\ Introd.\ $\mathrm{V}$, $\mathrm{V}^{*}$, $\Phi$, $^{v}\mathrm{W}$, $r_a$, $a^{\vee}$, $\mathrm{D}_0$, $\Pi$, $\Phi^{+}$, $\Phi^{-}$, $\Phi^{\text{réd}}$, $\Phi^{\text{nm}}$, $\Phi^{+\text{réd}}$, $\Phi^{-\text{réd}}$; $\Ncal_{\mathrm{X}}(\mathrm{Y})$, $\Ncal(\mathrm{Y})$.\\
(\btref{6.1.1}) $\mathrm{T}$, $\mathrm{U}_a$, $\mathrm{M}_a$, $\mathrm{U}^{+}$, $\mathrm{U}^{-}$.\\
(\btref{6.1.2}) $\mathrm{U}^{*}_a$.\\
(\btref{6.1.2})(2) $m(u)$, $\mathrm{M}^{0}_a$.\\
(\btref{6.1.2})(7) $\mathrm{L}_a$.\\
(\btref{6.1.2})(10) $^{v}\nu$.\\
(\btref{6.1.2})(12) $\mathrm{T}^{0}$, $\mathrm{G}'$.\\
(\btref{6.1.4}) $\mathrm{U}_{2a}$ (for $2a\notin\Phi$).\\
(\btref{6.1.5}) $\mathrm{G}^{0}$, $\mathrm{G}^{0}_i$, $\Phi_i$.\\
(\btref{6.2.1}) $\varphi$, $\varphi_a$, $\mathrm{U}_{a,k}$.\\
(\btref{6.2.2}) $\mathrm{U}_{2a,k}$ (for $2a\notin\Phi$), $\mathrm{M}_{a,k}$; $\Gamma_a$, $\Gamma'_a$.\\
(\btref{6.2.4}) $^{\varphi}\mathrm{U}_{a,k}$, $^{\varphi}\mathrm{M}_{a,k}$, etc.\\
(\btref{6.2.5}) $\lambda\varphi+v$, $n.\varphi$.\\
(\btref{6.2.6}) $\mathrm{A}=\varphi+\mathrm{V}$; $\alpha_{a,k}$, $\mathrm{U}_\alpha$, $\mathrm{U}_{\alpha+}$; $\Sigma$, $\Ech$.\\
(\btref{6.2.10}) $\nu$, $r_{a,k}$.\\
(\btref{6.2.11}) $\mathrm{H}$, $\widehat{\mathbf{W}}$, $\mathbf{W}$, $\mathrm{N}'$, $\mathrm{T}'$, $\mathrm{G}'$.\\
(\btref{6.2.13}) $\widetilde{\Gamma}_a$.\\
(\btref{6.3.3}) $\mathrm{U}_{a,k+}$.\\
(\btref{6.4.1}) $\widetilde{\R}$, $r+$.\\
(\btref{6.4.2}) $f_{\Omega}$, $\mathrm{U}_f$ (cf.\ (\btref{6.4.42})), $\mathrm{U}^{+}_f$, $\mathrm{U}^{-}_f$, $\mathrm{H}_f$, $\mathrm{N}_f$, $\mathrm{U}_{f,a}$, $\mathrm{U}^{(a)}_f$, $\mathrm{N}^{(a)}_f$, $\mathrm{H}^{(a)}_f$, $\mathrm{P}_f$.\\
(\btref{6.4.10}) $f'$, $\Phi_f$, $\mathrm{V}^{*}_f$.\\
(\btref{6.4.13}) $\mathrm{H}_{(k)}$.\\
(\btref{6.4.14}) $\mathrm{H}_{[k]}$.\\
(\btref{6.4.18}) $\mathrm{H}_{f,g}$.\\
(\btref{6.4.23}) $f^{*}$.\\
(\btref{6.4.38}) $\mathrm{U}_{0,k}$.\\
(6.4.40) $\mathrm{U}_0$.\\
(\btref{6.4.42}) $\mathrm{U}_f$ (for $f:\Phi\cup\{0\}\to\widetilde{\R}$); $\mathrm{U}_{0,k}$ (for $k\in\widetilde{\R}$).\\
(\btref{6.4.45})$^{(1)}$ $\mathscr{C}^{n}$.\\
7.\ Introd.\ $\Phi^{+}_{\mathrm{D}}$, $\Phi^{-}_{\mathrm{D}}$, $\Pi_{\mathrm{D}}$, $\mathrm{U}^{+}_{\mathrm{D}}$, $\mathrm{U}^{-}_{\mathrm{D}}$.\\
(7.1.1) $\mathrm{U}_{\Omega}$, $\mathrm{P}_{\Omega}$.\\
(7.1.2) $\operatorname{cl}(\Omega)$ (cf.\ (\btref{2.4.2}) and (\btref{7.4.11})).\\
(7.1.3) $\mathrm{N}_{\Omega}$, $\Phi_{\Omega}$; $\mathrm{U}_x$, $\mathrm{P}_x$, $\mathrm{N}_x$, etc.\\
(7.1.8) $\Nhat_{\Omega}$, $\Phat_{\Omega}$.\\
(\btref{7.1.10}) $\widehat{\mathbf{W}}_x$.\\
(7.2.1) $\gamma\cap \mathrm{Y}$, $\bar\gamma$ ($\gamma$ a filter base).\\
(7.2.2) $\mathrm{U}_\gamma$, $\mathrm{N}_\gamma$, $\mathrm{P}_\gamma$, $\Nhat_\gamma$, $\Phat_\gamma$, $\Phi_\gamma$ ($\gamma$ a filter base).\\
(7.2.3) $\gamma(\mathrm{D})$.\\
(7.2.4) $\gamma(x,\mathrm{E})$, $\mathrm{U}_{x,\mathrm{E}}$, $\mathrm{P}_{x,\mathrm{E}}$, $\mathrm{N}_{x,\mathrm{E}}$, \dots, $\mathrm{F}_{x,\mathrm{E}}$, $\mathrm{C}_{x,\mathrm{E}}$.\\
(7.2.6) $\mathrm{B}_{x,\mathrm{D}}$, $f_{x,\mathrm{D}}$; $f_{\mathrm{F}}$.\\
(7.2.7) $\overline{\mathrm{G}}_x$, $\mathrm{P}^{*}_x$, $\mathrm{H}^{*}_x$.\\
(\btref{7.3.9}) $\mathrm{N}^{\dagger}_{\mathrm{F}}$, $\mathrm{P}^{\dagger}_{\mathrm{F}}$.\\
(\btref{7.4.2}) $^{\varphi}\Build$, $\Build$ (cf.\ (\btref{2.1.1})).\\
(\btref{7.4.23}) $\mathrm{E}(v)$, $\widetilde{\mathrm{E}}(v)$, $c(v)$, $k(v)$, $\delta(v)$, $h(v)$.\\
(\btref{7.4.25}) $\rho_{\mathrm{A}';\gamma}$.\\
(\btref{7.5.1}) $\widehat{\Build}$.\\
(\btref{8.1.3}) $\mathscr{U}^{+}_{\mathrm{D}}$.\\
(\btref{8.1.5}) $\Bor(\Nor)$.\\
(\btref{8.1.7}) $\Bor(\varphi)$.\\
(\btref{9.1.1}) $\mathrm{V}^{\natural}$, $a^{\natural}$, $\Phi_b$, $\Phi^{+}_b$, $\mathrm{Z}$, $\mathrm{U}_b$, $\mathrm{L}_b$, $\mathrm{M}_b$.\\
(\btref{9.1.4}) $\mathrm{U}_{b,k}$, $f(b,k)$.\\
(\btref{9.1.6}) $\varphi_b$.\\
(\btref{9.1.7}) $\mathrm{A}^{\natural}$.\\
(\btref{9.1.8}) $\alpha_{b,k}$.\\
(\btref{9.1.9}) $\mathrm{G}^{\natural}$, $\Phi^{\natural}$, $\mathrm{T}^{\natural}$, $\mathrm{U}^{\natural}_b$, $^{v}\mathrm{W}^{\natural}$, $\mathrm{D}^{\natural}$, $\mathrm{M}^{\natural}_b$, $\mathrm{N}^{\natural}$, $\Phi^{\natural+}$, $\mathrm{U}^{\natural+}$, $\varphi^{\natural}_b$, $\varphi^{\natural}$.\\
(\btref{9.1.11}) $\nu^{\natural}$, $\Build^{\natural}$, etc.\\
(\btref{9.1.13}) $\mathrm{U}^{\varepsilon}_b$, $\mathrm{U}^{\varepsilon}_{b,k}$.\\
(\btref{9.1.14}) $\mathrm{H}^{\natural}$, $\varphi^{\natural}_0$.\\
(\btref{9.2.2}) $\Build^{\natural}$.\\
(\btref{10.1.1}) $\K$, $\sigma$, $\varepsilon$, $\Kse$; $\mathrm{X}$, $f$, $q$; $\mathrm{I}$, $\varepsilon(i)$, $e_i$, $\mathrm{X}_0$, $\mathrm{X}_i$, $c_{ij}(g)$, $\mathrm{Z}$.\\
(\btref{10.1.2}) $u_i(z,k)$, $m_i(z,k)$, $u_{ij}(k)$, $m_{ij}(k)$, $a_h$, $a_{ij}$, $\mathrm{U}_{a_i}$, $\mathrm{U}_{a_{ij}}$, $\mathrm{U}_{2a_i}$, $\mathrm{M}^{0}_{a_i}$, $\mathrm{M}^{0}_{2a_i}$, $\mathrm{M}^{0}_{a_{ij}}$.\\
(\btref{10.1.4}) $c(g)$, $\Is(f,q)$, $\Sim(f,q)$, $\Is^{+}(f,q)$, $\Sim^{+}(f,q)$, $\mathrm{T}(f,q)$, $\widetilde{\mathrm{T}}(f,q)$, $\mathrm{T}^{0}$, $\mathrm{G}^{0}$.\\
(10.1.6) $\mathrm{G}$.\\
(\btref{10.1.13}) $\varphi^{\omega}$, $\varphi^{\omega}_a$.\\
(\btref{10.1.14}) $\widehat{\K}$, $\widehat{\mathrm{X}}$, $\widehat{\mathrm{X}}_i$, $\widehat{f}$, $\widehat{q}$, $\widehat{\mathrm{Z}}$, $\mathrm{X}^{\perp}$, $\widehat{\mathrm{X}}^{\perp}$.\\
(\btref{10.1.20}) $\mathrm{X}^{*}_0$, $\omega_q$, $\omega^{*}_1$, $z^{*}$.\\
(10.1.22) $\omega_f$.\\
(\btref{10.1.27}) $\omega_i$, $\omega_{ij}$, $\ov_i$, $\ov_{ij}$, $\ov$.\\
(\btref{10.1.28}) $s(i)$.\\
(\btref{10.1.29}) $\Is_\omega(f,q)$.\\
(\btref{10.2.1}) $\mathrm{X}$, $e_i$, $g_{ij}$, $\zeta_{ij}$, $u_{ij}(k)$, $m_{ij}(k)$, $a_i$, $a_{ij}$, $\mathrm{V}^{*}$, $\Phi$, $\mathrm{U}_{a_{ij}}$, $\mathrm{M}^{0}_{a_{ij}}$.\\
(\btref{10.2.2}) $\mathrm{T}$, $\mathrm{M}_{a_{ij}}$.\\
(\btref{10.2.3}) $\varphi^{\omega}$.\\
(\btref{10.2.7}) $\det$, $\delta$.
\end{multicols}
\normalsize

% =====================================================================
\btsecA{}{Index of conditions}
% =====================================================================
\small
\begin{multicols}{2}
\raggedright
(T\,1) to (T\,4) : (\btref{1.2.6}).\\
(E\,1) to (E\,3) : (\btref{1.4.1}).\\
(CN) : (\btref{3.2.3}).\\
(A\,1) to (A\,6) : (\btref{5.2.1}).\\
(A\,1 \emph{bis}) : (\btref{5.2.2}).\\
(B\,1) to (B\,6) : (\btref{5.2.30}).\\
(DR\,1) to (DR\,6) : (\btref{6.1.1}).\\
(V\,0) to (V\,5) : (\btref{6.2.1}).\\
(V\,5 \emph{bis}) : (\btref{6.2.2}).\\
(C) : (\btref{6.4.3}).\\
(C\,1) and (C\,2) : (\btref{6.4.5}).\\
(QC\,1) and (QC\,2) : (\btref{6.4.7}).\\
(Pr) : (\btref{6.4.15}).\\
(P\,1), (P\,2), (P\,1$'$), (P\,1$''$) : (\btref{6.4.38}).\\
(V\,1 \emph{bis}), (VP) : (\btref{6.4.41}).\\
(MC) : (\btref{7.5.4}).\\
(BCN\,1) and (BCN\,2) : (\btref{8.1.2}).\\
(DDR\,1) : (\btref{9.1.9}).\\
(DDR\,2) : (\btref{9.1.10}).\\
(DDR\,1 \emph{bis}) : (\btref{9.1.13}).\\
(DP) : (\btref{9.1.15}).\\
(DI\,1) to (DI\,3) : (\btref{9.2.1}).\\
(DV\,1) : (\btref{9.2.1}).\\
(DDR\,3) : (\btref{9.2.9}).\\
(DV\,2) : (\btref{9.2.9}).
\end{multicols}
\normalsize

% =====================================================================
\btsecA{}{Terminological index}
% =====================================================================
\noindent\small
Headwords are given in English and re-alphabetized accordingly; the French
original is added in parentheses where the correspondence is not transparent.
Section references are those of the original index.
\normalsize

\medskip
\small
\begin{multicols}{2}
\raggedright
\textbf{A}\\
Adapted (B-, B-N-) : (\btref{1.2.13}).\\
Additive semi-norm : (\btref{10.1.20}).\\
Adjacent chambers (\emph{mitoyennes}) : (\btref{1.1.5}).\\
Admissible order : (\btref{5.2.13}).\\
Affine : see Datum, Set, Space, Weyl group, Root, Structure, Type.\\
Affine root : (\btref{1.3.3}), (\btref{6.2.6}).\\
Affine root datum : (\btref{5.2.36}) (footnote).\\
Affine root system : (\btref{1.3.3}).\\
Affine set : (\btref{1.1.1}).\\
Affine structure : (\btref{1.1.1}).\\
Affine subspace generated by a filter base : (\btref{9.2.4}).\\
Affine Weyl group : (\btref{1.3.2}).\\
Alcove (\emph{alcôve}) : (\btref{1.3.8}).\\
Apartment (\emph{appartement}) : (\btref{2.2.2}), (\btref{7.4.7}).\\
Apex of a sector : (\btref{1.3.11}), (7.1.4).\\
Associated : see Échelonnage, Pseudo-quadratic form, Reflection, Order relation, Vertex, Saturated Tits system.\\
Associated roots : (\btref{1.4.1}).\\
Attached root system : (\btref{1.4.2}).\\
\textbf{B}\\
B-adapted homomorphism : (\btref{1.2.13}).\\
B-N-adapted homomorphism : (\btref{1.2.13}), (\btref{4.1.3}).\\
Base of a root system : (\btref{1.3.12}).\\
Bornological group : (\btref{3.1.1}).\\
Bornology : (\btref{3.1.1}).\\
Bornology compatible with a group law : (\btref{3.1.1}).\\
Bornology compatible with a valued root datum : (\btref{8.1.1}).\\
Bornology defined by a Tits system : (\btref{3.1.4}), (\btref{3.1.9}).\\
Bornology defined by a valuation : (\btref{8.1.7}).\\
Bornology weakly compatible with a valued root datum : (\btref{8.1.1}).\\
Bruhat decomposition : (\btref{1.2.7}), 7.3 introd., (\btref{7.3.4}).\\
\textbf{C}\\
Canonical : see Map, Decomposition.\\
Canonical map of $\Build$ into $\mathbf{C}$ : (\btref{2.1.2}).\\
Canonical map of $\mathbf{A}$ into $\Build$ : (\btref{2.2.2}).\\
Canonical map of $\mathbf{C}_{\mathrm{X}}$ onto a facet of $\Build$ : (\btref{2.1.2}).\\
Cartan decomposition : (\btref{4.4.3}).\\
Center of a facet : (7.2.5).\\
Center of a retraction : (\btref{2.3.5}), (\btref{7.4.19}).\\
Chamber : (\btref{1.1.5}).\\
Chamber of a building : (\btref{2.1.2}), (\btref{7.4.12}).\\
Chamber of an affine space relative to an affine Weyl group : (\btref{1.3.3}).\\
Chamber of an equipollence class of valuations : (7.2.4).\\
Chambered morphism (\emph{chambré}) : (\btref{1.1.7}).\\
Closed part of a building : (\btref{2.4.1}), (\btref{2.5.7}).\\
Closed part of an apartment : (\btref{2.4.1}), (\btref{2.4.6}), (7.1.2), (\btref{7.4.11}).\\
Closed, open polysimplex : (\btref{1.1.3}).\\
Closed, open simplex : (\btref{1.1.2}).\\
Closure (\emph{enclos}) of a part of a building : (\btref{2.4.2}), (\btref{2.5.7}).\\
Closure (\emph{enclos}) of a part of an apartment : (\btref{2.4.2}), (\btref{2.4.6}), (7.1.2), (\btref{7.4.11}).\\
Closure of a germ (\emph{adhérence}) : (7.2.1).\\
Codimension of a facet : (\btref{2.1.1}).\\
Combinatorial building of a Tits system : (\btref{5.1.33}).\\
Compatible : see Bornology, Root data.\\
Completion of a building : (\btref{7.5.1}).\\
Concave function : (\btref{6.4.3}).\\
Connected type : (\btref{4.1.3}).\\
Contained in a set (germ) : (7.2.1).\\
Convex hull : (\btref{2.5.6}).\\
Convex part : (\btref{2.5.6}).\\
Coroot (\emph{coracine}) : \S\,6, introd.\\
Coxeter : see Graph, System.\\
Coxeter graph : (\btref{1.2.5}).\\
Coxeter graph underlying a system $(\mathrm{G},f,\mathrm{M},\mathrm{F})$ : (\btref{1.4.4}).\\
Coxeter system : (\btref{1.2.1}).\\
\textbf{D}\\
Dense valuation : \S\,7 introd.\\
Descent of a valuation : (\btref{9.1.11}).\\
Dimension of a facet : (7.2.4).\\
Dimension of a filter base : (\btref{9.2.4}).\\
Dimension of a polysimplicial complex ($\dim \mathrm{A}$) : (\btref{1.1.5}).\\
Dimension of a simplex : (\btref{1.1.2}).\\
Direction of a facet at a point : (7.2.4).\\
Direction of a sector : (\btref{1.3.11}), (7.1.4).\\
Discrete valuation : (\btref{6.2.21}).\\
Distance in a building : (\btref{2.5.4}), (\btref{7.4.20}).\\
Distance in an apartment : (\btref{2.2.8}), (\btref{7.4.11}).\\
Double Tits system : (\btref{5.1.1}).\\
Dynkin : see Graph.\\
Dynkin graph : (\btref{1.4.4}).\\
Dynkin graph of a root system : (\btref{1.4.4}).\\
Dynkin graph of an échelonnage : (\btref{1.4.5}).\\
\textbf{E}\\
Échelonnage : (\btref{1.4.1}).\\
Échelonnage associated with a discrete valuation : (\btref{6.2.22}).\\
Échelonnages, homothetic : (\btref{1.4.3}).\\
Échelonnages, similar : (\btref{1.4.3}).\\
End of a gallery (\emph{extrémité}) : (\btref{1.1.5}).\\
Endpoints of a segment of a building : (\btref{2.5.6}).\\
Equipollence class of valuations : (\btref{6.2.5}).\\
Equipollent parts of an affine space : (\btref{1.3.1}).\\
Equipollent valuations : (\btref{6.2.5}).\\
Equivalent Tits systems : (\btref{1.2.12}).\\
Equivalent valuations : (\btref{6.2.5}).\\
Euclidean affine space : (\btref{1.3.1}).\\
Extendable valuation (\emph{prolongeable}) : (\btref{6.4.38}).\\
Extended valuation (\emph{prolongée}) : (\btref{6.4.38}).\\
Extension of a valuation : (\btref{6.4.38}).\\
Extremal element of a subset of a root system : (\btref{1.3.15}).\\
\textbf{F}\\
Facet of a building : (\btref{2.1.1}), (\btref{7.4.12}).\\
Facet of a facet : (\btref{1.1.4}), (\btref{2.1.1}).\\
Facet of a polysimplicial complex : (\btref{1.1.4}).\\
Facet of an affine set : (\btref{1.1.1}).\\
Facet of an affine space, of an affine Weyl group, of an affine root system : (\btref{1.3.3}).\\
Facet of an equipollence class of valuations : (7.2.4).\\
Facet, vectorial : (\btref{1.3.10}).\\
\textbf{G}\\
Gallery : (\btref{1.1.5}).\\
Gallery joining two facets : (\btref{1.1.5}).\\
Gallery without stammering (\emph{sans bégaiement}) : (\btref{1.1.5}).\\
Gallery, minimal : (\btref{1.1.5}).\\
Generating root datum : (\btref{6.1.1}).\\
Geodesic : (\btref{2.5.13}).\\
Germ of a filter base : (7.2.1).\\
Germ with nonempty interior : (7.2.1).\\
Germ, open : (7.2.1).\\
Good maximal bounded subgroup : (\btref{4.4.1}).\\
\textbf{H}\\
H-component : (\btref{6.3.8}).\\
Half-apartment : (\btref{2.2.2}).\\
Homothetic échelonnages : (\btref{1.4.3}).\\
\textbf{I}\\
Induced valuation : (\btref{9.1.11}).\\
Interior of a germ : (7.2.1).\\
Inverse image of a bornology under a map : (\btref{3.1.2})\,c).\\
Inverse root : (\btref{1.3.8}).\\
Irreducible affine Weyl group : (\btref{1.3.2}).\\
Irreducible component of a Tits system : (\btref{2.6.5}).\\
Isometry : (\btref{10.1.4}).\\
Iwasawa decomposition : (\btref{4.4.3}), 7.3 introd., (\btref{7.3.1}).\\
\textbf{L}\\
Length of a gallery : (\btref{1.1.5}).\\
Length of an element of a Coxeter group : (\btref{1.2.2}).\\
\textbf{M}\\
Maximal bounded subgroup, good : (\btref{4.4.1}).\\
Maximal bounded subgroup, special : (\btref{4.4.1}).\\
Midpoint of two points of a building : (\btref{3.2.1}), (\btref{7.4.20}).\\
Morphism of affine sets : (\btref{1.1.1}).\\
Morphism of polysimplicial complexes : (\btref{1.1.7}).\\
Multipliable vertex : (\btref{1.4.4}).\\
\textbf{N}\\
Nibbling order (\emph{ordre grignotant}) : (\btref{1.3.15}).\\
Nondegenerate pair $(f,q)$ : (\btref{10.1.1}).\\
\textbf{O}\\
Open germ : (7.2.1).\\
Opposite sectors : (\btref{1.3.11}).\\
Order relation associated with a vector chamber : (\btref{1.3.16}).\\
Origin of a gallery : (\btref{1.1.5}).\\
\textbf{P}\\
Panel (\emph{cloison}) : (\btref{1.1.5}).\\
Panel of an affine space relative to an affine Weyl group : (\btref{1.3.3}).\\
Parabolic subgroup : (\btref{1.2.10}).\\
Parahoric subgroup : (\btref{1.5.1}).\\
Part : see Closed, Convex, Equipollent, Quasi-closed, Transverse.\\
Polysimplicial complex : (\btref{1.1.6}).\\
Positive root : (\btref{1.3.13}).\\
Product of affine sets : (\btref{1.1.1}).\\
Product of polysimplicial complexes : (\btref{1.1.8}).\\
Pseudo-quadratic form associated with a sesquilinear form : (\btref{10.1.1}).\\
\textbf{Q}\\
Quasi-closed part : (\btref{7.6.1}).\\
Quasi-concave function : (\btref{6.4.8}), (\btref{6.4.42}).\\
Quasi-valuation of a root datum : (\btref{6.2.3})\,e).\\
\textbf{R}\\
Rank of a Coxeter system : (\btref{1.3.4}).\\
Rank of a Tits system of affine type : (\btref{1.5.1}).\\
Rank of an affine Weyl group : (\btref{1.3.4}).\\
Reduced decomposition of an element of a Coxeter group : (\btref{1.2.2}).\\
Reflection associated with an element of a Weyl group : (\btref{2.3.10}).\\
Retraction of a building onto an apartment : (\btref{2.3.5}), (\btref{7.4.19}).\\
Retraction of a completed building : (\btref{7.5.7}).\\
Retraction relative to a sector : (\btref{2.9.2}).\\
Retraction relative to a sector germ : (\btref{7.4.25}).\\
Root datum (\emph{donnée radicielle}) : (\btref{6.1.1}).\\
Root datum compatible with another : (\btref{9.1.9}).\\
Root datum, valued : (\btref{6.2.1}).\\
Root system attached to a facet, to a point : (\btref{1.4.2}).\\
Root weight (\emph{poids radiciel}) : (\btref{1.3.8}).\\
Root, vectorial : (\btref{1.3.8}).\\
\textbf{S}\\
Saturated Tits system : (\btref{1.2.12}).\\
Saturated Tits system associated with a Tits system : (\btref{1.2.12}).\\
Schiffmann's formula : (\btref{7.3.3}).\\
Sector (\emph{quartier}) : (\btref{1.3.11}), (7.1.4).\\
Sector germ : (\btref{2.9.2}), (7.2.3), (\btref{7.4.12}).\\
Sector of a building : (\btref{2.2.2}), (\btref{7.4.12}).\\
Segment with given endpoints in a building : (\btref{2.5.6}), (\btref{7.4.20}).\\
Similar échelonnages : (\btref{1.4.3}).\\
Similitude : (\btref{10.1.4}).\\
Similitude ratio (\emph{rapport de similitude}) : (\btref{10.1.4}).\\
Simple root : (\btref{1.3.12}).\\
Simple root system : (\btref{1.3.12}).\\
Simplicial complex : (\btref{1.1.6}).\\
Special maximal bounded subgroup : (\btref{4.4.1}).\\
Special point of an affine space equipped with an affine Weyl group : (\btref{1.3.7}).\\
Special valuation : (\btref{6.2.13}).\\
Special vertex of a Coxeter graph : (\btref{1.3.7}).\\
Strongly equivalent Tits systems : (\btref{1.2.12}).\\
Structural map : (\btref{2.2.2}).\\
Support of a facet : (\btref{1.3.3}).\\
Support of an element of a Coxeter group : (\btref{1.2.2}).\\
\textbf{T}\\
Taut gallery between two facets, two points (\emph{tendue}) : (\btref{1.1.5}).\\
Tits : see System.\\
Tits system : (\btref{1.2.6}).\\
Tits system of affine type : (\btref{1.5.1}).\\
Tits systems, equivalent : (\btref{1.2.11}).\\
Totally singular subspace : (\btref{10.1.1}).\\
Trace form (\emph{forme tracique}) : (\btref{10.1.1}).\\
Transverse closed parts : (\btref{5.1.2}).\\
Type of a facet : (\btref{1.3.5}), (\btref{2.1.1}).\\
Type of a gallery : (\btref{1.3.6}), (\btref{2.1.3}).\\
Type of a panel : (\btref{1.3.6}).\\
Type of a parabolic subgroup : (\btref{1.2.10}).\\
Type of a root datum : (\btref{6.1.1}).\\
\textbf{U}\\
Underlying Coxeter graph : (\btref{1.4.4}).\\
\textbf{V}\\
Valuation of a root datum : (\btref{6.2.1}).\\
Vector chamber : (\btref{1.3.10}).\\
Vectorial : see Chamber, Facet, Root.\\
Vertex of a building associated with a parahoric subgroup : (\btref{2.1.2}).\\
\textbf{W}\\
Wall of a building : (\btref{2.2.2}).\\
Wall of a half-apartment : (\btref{2.2.4}).\\
Wall of an affine space relative to a Weyl group : (\btref{1.3.3}).\\
Wall of an equipollence class of valuations : (\btref{6.2.6}).\\
Weight (\emph{poids}) : (\btref{1.3.8}).\\
Weyl group : (\btref{1.2.6}), (\btref{1.3.2}).\\
Weyl group of a Tits system : (\btref{1.2.6}).\\
Witt index : (\btref{10.1.1}).\\
\end{multicols}
\normalsize
\par\endgroup

\clearpage
\btsecA{}{Appendix: French--English glossary}
\noindent\small
The renderings of the French technical vocabulary used in this translation.
\par\medskip
\raggedright
\begin{longtable}{@{}p{0.44\linewidth}p{0.48\linewidth}@{}}
\'echelonnage & \'echelonnage (retained)\\
\'epinglage & pinning\\
\'equipollent, \'equivalent & equipollent, equivalent\\
\'equipollentes & equipollent\\
adhérence & closure\\
alc\^ove & alcove\\
appartement & apartment\\
applications structurales & structural maps\\
base de filtre & filter base\\
borne inférieure / supérieure & infimum / supremum\\
bornologie & bornology\\
c\^one simplicial & simplicial cone\\
centre de gravité & barycenter\\
chambre & chamber\\
chambre vectorielle & vector chamber\\
cloison & panel\\
close & closed\\
complexe polysimplicial & polysimplicial complex\\
corps & field (not assumed commutative unless stated)\\
corps gauche & skew field\\
corps local & local field\\
corps r\'esiduel & residue field\\
corps résiduel & residue field\\
corps valué & valued field\\
d\'ecomposition r\'eduite & reduced decomposition\\
d\'eplacements & displacements\\
d\'eploy\'e / quasi-d\'eploy\'e & split / quasi-split\\
descente & descent\\
donn\'ee radicielle & root datum\\
donn\'ee radicielle g\'en\'eratrice & generating root datum\\
donn\'ee radicielle valu\'ee & valued root datum\\
donnée radicielle (valuée) & (valued) root datum\\
donnée radicielle génératrice & generating root datum\\
double syst\`eme de Tits & double Tits system\\
déployé & split\\
ensemble affine & affine set\\
enveloppe convexe & convex hull\\
espace affine euclidien & Euclidean affine space\\
facette & facet\\
facette / facette vectorielle & facet / vector facet\\
facette vectorielle & vector facet\\
famille filtrante croissante & increasing directed family\\
fixateur & fixer (pointwise stabilizer)\\
fortement \'equivalents & strongly equivalent\\
fortement orthogonales & strongly orthogonal\\
galerie & gallery\\
germe de quartier & sector germ\\
graphe de Coxeter & Coxeter graph\\
groupe des commutateurs & commutator group\\
homomorphisme B-adapt\'e & B-adapted homomorphism\\
immeuble & building\\
mitoyennes (chambres) & adjacent (chambers)\\
morphisme chambr\'e & chambered morphism\\
mur & wall\\
mur, bord & wall, boundary\\
ordre grignotant & nibbling order\\
par récurrence sur & by induction on\\
partie close & closed subset\\
parties closes & closed subsets\\
poids radiciels & radical weights\\
point extr\'emal & extreme point\\
point interne & internal point\\
point sp\'ecial & special point\\
polysimplexe & polysimplex\\
prolongement / prolongeable & extension / extendable\\
pronilpotent / prounipotent & pronilpotent / prounipotent\\
quartier & sector\\
quitte à remplacer & after replacing, if necessary\\
r\'eseau & lattice\\
r\'etraction & retraction\\
racine affine & affine root\\
racine vectorielle & vector root\\
racines non divisibles / non multipliables & non-divisible / non-multipliable roots\\
radical prounipotent & prounipotent radical\\
sans b\'egaiement & non-stuttering\\
satur\'e & saturated\\
similitude & similarity\\
sommet (d'un quartier) & apex (of a sector); vertex (of a sector)\\
sous-groupe born\'e & bounded subgroup\\
sous-groupe distingu\'e & normal subgroup\\
sous-groupe distingué & normal subgroup\\
sous-groupe parabolique & parabolic subgroup\\
sous-groupe parahorique & parahoric subgroup\\
stabilisateur & stabilizer\\
syst\`eme de racines inverse & inverse (dual) root system\\
syst\`eme de Tits & Tits system (= BN-pair)\\
série centrale descendante & descending central series\\
tendue (galerie) & taut (gallery)\\
valuation & valuation\\
valuation dense & dense valuation\\
valuation impropre & improper valuation\\
valuation prolongée & extended valuation\\
équipollent & equipollent (parallel and equally directed)\\
\end{longtable}
\normalsize

\end{document}
